\documentclass[UTF8,fontset=none,11pt,oneside,openany]{ctexbook}

\usepackage{iftex}
\usepackage{fontspec}
\ifXeTeX
  \usepackage{xeCJK}
\fi
\defaultfontfeatures{Ligatures=TeX}
\setmainfont{texgyretermes-regular.otf}[
  BoldFont=texgyretermes-bold.otf,
  ItalicFont=texgyretermes-italic.otf
]
\setsansfont{texgyreheros-regular.otf}
\setmonofont{lmmono10-regular.otf}[ItalicFont=lmmono10-italic.otf]
\setCJKmainfont{FandolSong-Regular.otf}[
  BoldFont=FandolSong-Bold.otf,
  ItalicFont=FandolKai-Regular.otf
]
\setCJKsansfont{FandolHei-Regular.otf}
\setCJKmonofont{FandolFang-Regular.otf}
\ifXeTeX
  \xeCJKsetup{PunctStyle=kaiming,CheckSingle=true}
  \XeTeXlinebreaklocale "zh"
  \XeTeXlinebreakskip=0pt plus 1pt
\fi

\usepackage{amsmath,amssymb,mathrsfs}
\usepackage{graphicx}
\usepackage{tikz-cd}
\ifLuaTeX
  \usepackage{luatex85}
\fi
\usepackage[all]{xy}
\usepackage[shortlabels]{enumitem}
\usepackage{longtable}
\usepackage{multicol}
\usepackage{marginnote}
\ifLuaTeX
\else
  \usepackage{microtype}
\fi
\usepackage{fancyhdr}
\usepackage[
  a4paper,
  inner=26mm,
  outer=24mm,
  top=24mm,
  bottom=24mm,
  headheight=14pt,
  headsep=8mm,
  footskip=12mm
]{geometry}
\usepackage[
  unicode=true,
  colorlinks=true,
  linkcolor=black,
  citecolor=black,
  urlcolor=blue,
  hypertexnames=false
]{hyperref}

\hypersetup{
  pdftitle={代数几何原理 I：概形语言（简体中文工作版）},
  pdfauthor={Alexander Grothendieck；Jean Dieudonné（协作编写）},
  pdfsubject={依据 NUMDAM 法文印本与法文外交式 TeX 制作的大陆简体中文版本},
  pdfkeywords={代数几何；概形；EGA；简体中文}
}

\AtBeginDocument{\fontsize{10.95pt}{16.5pt}\selectfont}
\setlength{\parindent}{2em}
\setlength{\parskip}{0pt}
\setlength{\emergencystretch}{2em}
\flushbottom

\pagestyle{fancy}
\fancyhf{}
\fancyhead[L]{\small 代数几何原理 I}
\fancyhead[R]{\small 概形语言}
\fancyfoot[C]{\thepage}
\renewcommand{\headrulewidth}{0.4pt}
\renewcommand{\footrulewidth}{0pt}
\fancypagestyle{plain}{%
  \fancyhf{}%
  \fancyhead[L]{\small 代数几何原理 I}%
  \fancyhead[R]{\small 概形语言}%
  \fancyfoot[C]{\thepage}%
  \renewcommand{\headrulewidth}{0.4pt}%
  \renewcommand{\footrulewidth}{0pt}%
}

\renewcommand{\contentsname}{目录}
\renewcommand{\bibname}{参考文献}
\renewcommand{\figurename}{图}
\renewcommand{\tablename}{表}
% EGA numbers its sections independently of the LaTeX book chapter counter:
% §1, §2, ... and subsections 1.1, 1.2, ... rather than 0.1 / 0.1.1.
\renewcommand{\thesection}{\arabic{section}}
\renewcommand{\thesubsection}{\thesection.\arabic{subsection}}

\newenvironment{env}[1][]{%
  \par\medskip\noindent\textbf{(#1)}\ %
}{\par}

\newenvironment{definition}[1][]{%
  \par\medskip\noindent\textbf{定义 (#1).}\ %
}{\par}
\newenvironment{lemma}[1][]{%
  \par\medskip\noindent\textbf{引理 (#1).}\ %
}{\par}
\newenvironment{theorem}[1][]{%
  \par\medskip\noindent\textbf{定理 (#1).}\ %
}{\par}
\newenvironment{proposition}[1][]{%
  \par\medskip\noindent\textbf{命题 (#1).}\ %
}{\par}
\newenvironment{corollary}[1][]{%
  \par\medskip\noindent\textbf{推论 (#1).}\ %
}{\par}
\newenvironment{remark}[1][]{%
  \par\medskip\noindent\textbf{注 (#1).}\ %
}{\par}
\newenvironment{notation}[1][]{%
  \par\medskip\noindent\textbf{记号 (#1).}\ %
}{\par}
\newenvironment{example}[1][]{%
  \par\medskip\noindent\textbf{例 (#1).}\ %
}{\par}
\newenvironment{examples}[1][]{%
  \par\medskip\noindent\textbf{例 (#1).}\ %
}{\par}
\newenvironment{proof}{%
  \par\medskip\noindent\textbf{证明.}\ %
}{\par}
\newenvironment{namedtheorem}[2][]{%
  \par\medskip\noindent\textbf{#2 (#1).}\ %
}{\par}
\newenvironment{namedlemma}[2][]{%
  \par\medskip\noindent\textbf{#2 (#1).}\ %
}{\par}

\newcommand{\oldpage}[2][]{%
  \hypertarget{ega-source-page-#1-#2}{}%
  \marginnote{\scriptsize\textbf{#1}~|~#2}\ignorespaces}

\let\EGAOriginalLabel\label
\renewcommand{\label}[1]{\phantomsection\EGAOriginalLabel{#1}}

\begin{document}
% Authority: EGA I NUMDAM physical PDF p.1; no diplomatic-French-TeX counterpart.
% Physical-page anchor: EGA-I-PDF-001.
\begin{titlepage}
\thispagestyle{empty}
\begin{center}
{\Large 法国高等科学研究所《数学出版物》\par}
\end{center}

\vspace*{7em}

{\large ALEXANDER GROTHENDIECK\par}
\medskip
{\bfseries 《代数几何原理》第一卷：概形语言\par}
\medskip
法国高等科学研究所《数学出版物》，第 4 卷（1960），第 5--228 页

\medskip
\url{http://www.numdam.org/item?id=PMIHES_1960__4__5_0}

\vfill

\begin{minipage}{0.76\textwidth}
\small
\copyright\ 法国高等科学研究所《数学出版物》，1960。保留所有权利。

访问期刊《法国高等科学研究所数学出版物》的档案
（\url{http://www.ihes.fr/IHES/Publications/Publications.html}），
即表示同意 NUMDAM 的一般使用条件
（\url{http://www.numdam.org/conditions}）。任何商业使用或系统性打印均构成
刑事违法行为。对本文件制作的任何副本或打印件，均须载有本版权声明。
\end{minipage}

\vfill

{\large NUMDAM\par}
\smallskip
本文由“古典数学文献数字化”计划数字化

\smallskip
\url{http://www.numdam.org/}
\end{titlepage}
\clearpage
% French editable witness: source/ega1/frontmatter-fr.tex
% Physical-page anchors: EGA-I-PDF-002 and EGA-I-PDF-003.
\begin{titlepage}
\centering

{\Large 法国高等科学研究所\par}
\vspace{0.4em}
{\large INSTITUT DES HAUTES ÉTUDES SCIENTIFIQUES\par}

\includegraphics[width=42mm]{03_zh_hans_cn_tex/assets/ihes-emblem-ega1.png}
\vfill

{\Large\bfseries 代数几何原理\par}
\vspace{1.5em}
{\large A. 格罗滕迪克 著\par}
\vspace{0.7em}
{J. 迪厄多内协作编写\par}

\vspace{2em}
I

\vspace{1.5em}
{\Large\bfseries 概形语言\par}

\vfill
{\Large 1960\par}
\vspace{1.2em}
{\Large 《数学出版物》第 4 号\par}
\vspace{0.5em}
{\large 5, ROND-POINT BUGEAUD --- PARIS (XVI\textsuperscript{e})\par}
\end{titlepage}

\thispagestyle{empty}
\vspace*{\fill}
\begin{center}
\fbox{%
  \begin{minipage}{0.50\textwidth}
  \centering
  法定送存

  \smallskip
  第 1 版 \dotfill 1960 年第 3 季度

  \smallskip
  版权所有

  各国权利均予保留

  \smallskip
  \copyright\ 1960，\emph{法国高等科学研究所}
  \end{minipage}%
}
\end{center}
\vspace*{\fill}
\clearpage
\setcounter{page}{5}
\pdfbookmark[0]{导言}{ega1-intro}
\section*{导言}

\begin{center}
\emph{献给奥斯卡·扎里斯基和安德烈·韦伊。}
\end{center}

\oldpage[I]{5}

本篇论文以及日后承续它的许多篇论文，旨在共同构成一部关于代数几何基础的论著。
原则上，它们不预设读者对这门学科具有任何专门知识；甚至事实表明，尽管这种知识显然有益，
但它有时也会妨碍希望熟悉本文所述观点与方法的读者，因为这种知识往往使人过度习惯于双有理观点。
反之，我们将假定读者充分掌握下列各门学科：

\begin{enumerate}
  \item[a)] \emph{交换代数}，例如 N.~Bourbaki 尚在编写的《数学原理》各卷中所述的内容
  （在这些卷出版以前，可参阅 Samuel--Zariski~\cite{I-13} 以及 Samuel~\cite{I-11},~\cite{I-12}）。
  \item[b)] \emph{同调代数}，可参阅 Cartan--Eilenberg~\cite{I-2}（下引作（M））、
  Godement~\cite{I-4}（下引作（G）），以及 A.~Grothendieck 最近的论文~\cite{I-6}（下引作（T））。
  \item[c)] \emph{层论}；我们的主要参考文献是（G）与（T）。层论提供了不可缺少的语言，
  用以把交换代数中的基本概念解释成“几何的”概念，并把它们“整体化”。
  \item[d)] 最后，读者若对\emph{函子语言}已有一定熟悉，将会很有帮助。本论著将不断使用这种语言；
  读者可参阅（M）、（G），尤其是（T）。这种语言的原理以及一般函子理论的主要结果，
  将在本论著两位作者正在编写的另一部著作中作更详细的阐述。
\end{enumerate}

\bigskip
\begin{center}*\quad*\quad*\end{center}
\bigskip

在这篇导言中，不宜对代数几何中的“概形”观点作或详或略的说明，也不宜罗列促使人们必须采用这种观点的漫长理由；
尤其不宜在此详述，为何必须系统地接纳我们所考虑的“簇”的局部环中的幂零元
（这必然使有理映射的概念退居次位，而让位于正则映射或“态射”的概念）。
本论著的目的，正是系统地建立“概形”语言，并将证明——我们希望如此——这种语言的必要性。
即使做到这一点并不困难，我们也

\oldpage[I]{6}
不会在此对第一章所展开的概念作“直观的”介绍。
希望预先概览本论著内容的读者，可参阅 A.~Grothendieck 于 1958 年在爱丁堡国际数学家大会所作的报告~\cite{I-7}，
以及同一作者的讲稿~\cite{I-8}。
J.-P.~Serre 的论文~\cite{I-14}（下引作（FAC））也可视为经典观点与代数几何概形观点之间的一种过渡性阐述；
因此，阅读它可为阅读我们的《代数几何原理》作极好的准备。

\bigskip
\begin{center}*\quad*\quad*\end{center}
\bigskip

为供参考，现列出本论著预定的总体纲目；不过它日后仍可能修改，尤其是后面几章：

\medskip
\begin{tabular}{rl}
  第一章。 & --- 概形语言。\\
  --- \textbf{II}. & --- 若干类态射的初等整体研究。\\
  --- \textbf{III}. & --- 凝聚代数层的上同调及其应用。\\
  --- \textbf{IV}. & --- 态射的局部研究。\\
  --- \textbf{V}. & --- 构造概形的初等方法。\\
  --- \textbf{VI}. & --- 下降技巧与构造概形的一般方法。\\
  --- \textbf{VII}. & --- 群概形与主纤维丛。\\
  --- \textbf{VIII}. & --- 纤维空间的微分研究。\\
  --- \textbf{IX}. & --- 基本群。\\
  --- \textbf{X}. & --- 留数与对偶。\\
  --- \textbf{XI}. & --- 相交理论、Chern 类与 Riemann--Roch 定理。\\
  --- \textbf{XII}. & --- 阿贝尔概形与 Picard 概形。\\
  --- \textbf{XIII}. & --- Weil 上同调。\\
\end{tabular}
\medskip

原则上，所有章节都视为开放的，日后随时可以增添新的节；为减轻这种出版方式带来的不便，新增的节将以单独分册出版。
如果某一节在一章出版时已经列入计划或正在编写，即使由于轻重缓急的安排，其实际出版时间会晚得多，
也仍将在该章目录中列出。
为方便读者，我们在“第 0 章”中汇集本论著各章所用的若干交换代数、同调代数和层论补充材料；
这些材料或多或少已为人所知，却无法找到便于引用的文献。
建议读者只在阅读本论著正文的过程中，并且在对所引结果还不

\oldpage[I]{7}
够熟悉时，才回头查阅第 0 章。
我们还认为，按照这种方式阅读本论著，对于初学者会是熟悉交换代数与同调代数的一种良好途径；
相当多的人认为，如果学习这两门学科时没有具体应用相伴，过程就会枯燥，甚至令人沮丧。

\bigskip
\begin{center}*\quad*\quad*\end{center}
\bigskip

我们无力在这篇导言中对所述概念和结果作哪怕简略的历史综述。
正文只会列出对理解内容特别有用的参考文献，也只会说明最重要结果的渊源。
至少在形式上，本书所处理的主题相当新；因此，对 19 世纪及 20 世纪初代数几何先辈的引用很少，
我们也只是间接听闻他们的工作。
不过，仍应在此简述几部最直接影响作者并促成概形观点发展的著作。
首先必须提到 J.-P.~Serre 的奠基性论文（FAC）。它曾把不止一位年轻研究者
（包括本论著的一位作者）引入代数几何；这些研究者原先都因 A.~Weil 的经典著作
\emph{Foundations}~\cite{I-18} 过于艰涩而却步。
正是在（FAC）中，人们第一次证明：一个“抽象”代数簇的“Zariski 拓扑”完全适合应用代数拓扑的某些方法，
并且尤其能够产生一种上同调理论。
此外，（FAC）所给的代数簇定义，也是最自然地适合推广为本文所建立概念的定义\footnote{正如 J.-P.~Serre 向我们指出的，
用一个环层来规定簇的结构这一思想应归于 H.~Cartan；Cartan 正是以此为出发点建立了解析空间理论。
当然，与代数几何中的情形完全一样，在“解析几何”中也应当让幂零元在解析空间的局部环中取得正式地位。
H.~Cartan 与 J.-P.~Serre 的定义的这种推广，最近已由 H.~Grauert~\cite{I-5} 着手研究；
可以期待，在这一一般框架内对解析几何作系统阐述的著作不久将会问世。
显然，本论著所建立的概念和方法在解析几何中仍有意义，尽管在后一理论中必须预期更为严重的技术困难。
代数几何的方法较为简明，因而可以预见，它能够成为解析空间理论未来发展的某种形式模型。}。
Serre 本人也已注意到，把域上的仿射代数换成任意交换环以后，仿射代数簇的上同调理论仍可毫无困难地转述。
因此，本论著第一、二章以及第三章最初两节，就其本质而言，可以看作是在这一扩大框架中对（FAC）
及同一作者后来一篇论文~\cite{I-15} 的主要结果所作的直接转述。
C.~Chevalley 的\emph{代数几何讨论班}~\cite{I-1} 也使我们获益良多；尤其是，
系统使用他所引入的“可构造集”，在概形理论中被证明极为有用（参见第四章）。
我们还从他那里借用了关于态射的维数研究

\oldpage[I]{8}
（第四章）；这项研究可在概形框架内基本原样转述。
还应注意，Chevalley 所引入的“局部环概形”概念，自然适合用来推广代数几何，
不过它不具备我们希望在此赋予理论的全部灵活性和一般性；关于该概念与本理论的关系，见第一章第 8 节。
M.~Nagata 已在一系列论文~\cite{I-9} 中发展了这种推广，其中包含许多关于 Dedekind 环上的代数几何的专门结果\footnote{在接近我们代数几何观点的工作中，
还应提到 E.~Kähler 的重要论文~\cite{I-22}，以及 Chow 与 Igusa 最近的一则札记~\cite{I-3}；
后二人在 Nagata--Chevalley 理论的框架内重新讨论了（FAC）的若干结果，并给出了一个 Künneth 公式。}。

\bigskip
\begin{center}*\quad*\quad*\end{center}
\bigskip

最后，不言而喻，任何一本关于代数几何的书，尤其是关于其基础的书，
都必然会受到 O.~Zariski、A.~Weil 等数学家的影响，哪怕这种影响是间接传来的。
特别是，Zariski 的\emph{全纯函数理论}~\cite{I-20} 经上同调方法适当软化，
再由一个存在性定理补足（第三章第 4、5 节），便与第六章所述下降技巧一道，
构成本论著所用的主要工具之一；在我们看来，这也是代数几何现有工具中最强有力的工具之一。

它所归属的一般方法可以概述如下（第九章研究基本群时将给出一个典型例子）。
设有从一个代数簇到另一个代数簇的固有态射（第二章）$f:X\to Y$
（更一般地，是从一个概形到另一个概形的态射）；我们想在某一点 $y\in Y$ 附近研究它，
以解决关于 $y$ 的某个邻域的问题 P。依次进行下列步骤：
\begin{enumerate}
  \item[$1^\circ$] 可以假设 $Y$ 是仿射的，于是 $X$ 成为定义在仿射环 $A$（即 $Y$ 的仿射环）上的概形；
  甚至还可以把 $A$ 换成 $y$ 的局部环。这种约化在实践中总是容易的（第五章），
  因而我们归结到 $A$ 为\emph{局部环}的情形。
  \item[$2^\circ$] 当 $A$ 为\emph{Artin 局部环}时，研究所考虑的问题。
  为使问题在不假定 $A$ 为整环时仍真正有意义，有时必须重新表述问题 P；
  这样做往往会使我们对问题获得更好的理解，而问题在此阶段具有“无穷小”性质。
  \item[$3^\circ$] 形式概形理论（第三章第 3、4、5 节）使我们能够从 Artin 环的情形过渡到\emph{完备局部环}的情形。
  \item[$4^\circ$] 最后，若 $A$ 是任意局部环，考察 $X$ 上适当概形的“多值截面”，
  使其逼近给定的“形式”截面（第四章），往往可以把从 $X$ 经标量扩张到 $A$ 的完备化所得概形

  \oldpage[I]{9}
  的已知结果，过渡为 $A$ 的某个相当简单的有限扩张（例如非分歧扩张）上的类似结果。
\end{enumerate}

这一概述表明，系统研究定义在 Artin 环 $A$ 上的概形十分重要。
Serre 在局部类域论表述中所采取的观点以及 Greenberg 最近的工作似乎提示：
可以对这样的概形 $X$ 进行研究，办法是函子性地给它配上一个概形 $X'$；
该概形定义在剩余域 $k$（即 $A$ 的剩余域，假定为完美域）上，并且在良好情形下其维数等于 $n\dim X$，
其中 $n$ 是 $A$ 的长度。

关于 A.~Weil 的影响，只需指出：为了以所需的全部一般性表述“Weil 上同调”的定义，
并着手证明\footnote{为避免误解，需要说明：在撰写这篇导言时，这项工作才刚刚开始，
因此尚未导出 Weil 猜想的证明。}建立他在丢番图几何中的著名猜想~\cite{I-19} 所需的一切形式性质，
必须发展相应的工具；这种需要是撰写本论著的主要动机之一。
同样重要的动机，是为代数几何中惯用的概念与方法寻找自然框架，并借此使作者有机会真正理解这些概念和方法。

\bigskip
\begin{center}*\quad*\quad* \end{center}
\bigskip

最后，我们认为有必要提醒读者：他们和作者本人一样，在习惯概形语言之前，
以及在确信几何直觉所提示的通常构造基本上只能以一种合理方式转写为这种语言之前，
大概都会经历一段困难时期。
与现代数学的许多分支一样，初始直觉在表面上越来越远离能够以所需精确性和一般性表达它的语言。
这里的心理困难，在于必须把集合中熟悉的概念——笛卡儿积、群律、环律、模律、纤维丛、齐性主丛等——
搬到一个已经与集合范畴颇为不同的范畴中的对象上
（即预概形范畴，或给定预概形上的预概形范畴）。
未来的数学家恐怕难以回避这一新的抽象训练；归根结底，
同我们的先辈为熟悉集合论所付出的努力相比，这点训练也许相当有限。

\bigskip
\begin{center}*\quad*\quad*\end{center}
\bigskip

引用将采用十进制编号。例如在 III, 4.9.3 中，III 表示章，4 表示节，9 表示该节内的款。
在同一章内引用时，将省略章号。

\thispagestyle{fancy}
\clearpage
\setcounter{page}{11}
% French authority: source/ega1/ega0-1-fr.tex
% Sequential translation cursor: complete source file through 7.8.3 and （待续。） (source lines 1--5978; printed pages 11--78).

\section{分式环}
\label{section:0.1-fr}

\setcounter{subsection}{-1}
\subsection{环与代数}
\label{subsection:0.1.0-fr}

\oldpage[0\textsubscript{I}]{11}
\begin{env}[1.0.1]
\label{0.1.0.1-fr}
本论著所考虑的所有环都具有\emph{单位元}；所有这样的环上的模均假定为\emph{幺模}；环同态一律假定\emph{把单位元映到单位元}；除非明确说明相反情形，环 $A$ 的子环均假定\emph{包含 $A$ 的单位元}。
我们主要考虑\emph{交换环}；以后若只说“环”而不加限定，均默认所指的是交换环。
若 $A$ 是不必交换的环，则除非明确说明相反情形，所谓 $A$-模一律指\emph{左}模。
\end{env}

\begin{env}[1.0.2]
\label{0.1.0.2-fr}
设 $A$、$B$ 是两个不必交换的环，$\varphi:A\to B$ 是一个同态。
任一左（相应地，右）$B$-模 $M$ 都可按 $a.m=\varphi(a).m$（相应地，$m.a=m.\varphi(a)$）赋予左（相应地，右）$A$-模结构；需要区分 $M$ 上的 $A$-模结构与 $B$-模结构时，把这样定义的左（相应地，右）$A$-模记作 $M_{[\varphi]}$。
若 $L$ 是一个 $A$-模，则同态 $u:L\to M_{[\varphi]}$ 就是满足 $u(a.x)=\varphi(a).u(x)$（$a\in A$，$x\in L$）的交换群同态；也称它为一个\emph{$\varphi$-同态} $L\to M$，并称二元组 $(\varphi,u)$（或滥用语言时简称 $u$）为从 $(A,L)$ 到 $(B,M)$ 的一个\emph{双同态}（di-homomorphism）。
因此，由二元组 $(A,L)$——其中 $A$ 是环，$L$ 是 $A$-模——构成一个\emph{范畴}，其态射就是双同态。
\end{env}

\begin{env}[1.0.3]
\label{0.1.0.3-fr}
在 (1.0.2) 的假设下，若 $\mathfrak{J}$ 是 $A$ 的一个左（相应地，右）理想，则以 $B\mathfrak{J}$（相应地，$\mathfrak{J}B$）表示 $B$ 中由 $\varphi(\mathfrak{J})$ 生成的左（相应地，右）理想 $B\varphi(\mathfrak{J})$（相应地，$\varphi(\mathfrak{J})B$）；它也是左（相应地，右）$B$-模的典范同态 $B\otimes_A\mathfrak{J}\to B$（相应地，$\mathfrak{J}\otimes_A B\to B$）的像。
\end{env}

\begin{env}[1.0.4]
\label{0.1.0.4-fr}
若 $A$ 是一个（交换）环，$B$ 是一个不必交换的环，则在 $B$ 上给定一个\emph{$A$-代数}结构，等价于给定一个环同态 $\varphi:A\to B$，使得 $\varphi(A)$ 包含在 $B$ 的中心内。
此时，对 $A$ 的任一理想 $\mathfrak{J}$，$\mathfrak{J}B=B\mathfrak{J}$ 是 $B$ 的双边理想；对任一 $B$-模 $M$，$\mathfrak{J}M$ 是一个 $B$-模，并且等于 $(B\mathfrak{J})M$。
\end{env}

\begin{env}[1.0.5]
\label{0.1.0.5-fr}
关于\emph{有限生成模}与\emph{有限生成}（交换）\emph{代数}的概念，以下不再赘述；称 $A$-模 $M$ 有限生成，是指存在

\oldpage[0\textsubscript{I}]{12}
正合列 $A^p\to M\to 0$。
若 $A$-模 $M$ 同构于某个同态 $A^p\to A^q$ 的余核，换言之，若存在正合列 $A^p\to A^q\to M\to 0$，便称其有一个\emph{有限表示}。
应注意，在\emph{Noether} 环 $A$ 上，每个有限生成 $A$-模都有有限表示。

回忆：若 $A$-代数 $B$ 在 $A$ 上满足如下条件：$B$ 的每个元素都是 $B$ 中某个以 $A$ 中元素为系数的首一多项式的根，就称它为\emph{整代数}；这也等价于说，$B$ 的每个元素都包含在 $B$ 的某个作为\emph{有限生成 $A$-模}的子代数中。
在这种情形下，若 $B$ 交换，则 $B$ 中由 $B$ 的任一有限子集生成的子代数都是有限生成 $A$-模；因此，交换代数 $B$ 在 $A$ 上既整又有限生成，当且仅当 $B$ 是有限生成 $A$-模；此时也称 $B$ 为\emph{有限整 $A$-代数}（若不致混淆，则简称\emph{有限代数}）。
还应注意，这些定义并不假定同态 $A\to B$ 所规定的 $A$-代数结构是单射的。
\end{env}

\begin{env}[1.0.6]
\label{0.1.0.6-fr}
\emph{整环}是这样的环：任一由元素 $\neq 0$ 组成的有限族，其元素的乘积仍 $\neq 0$；等价地，在这样的环中 $0\neq 1$，并且任意两个元素 $\neq 0$ 的乘积非零。
环 $A$ 的\emph{素理想}是满足 $A/\mathfrak{p}$ 为整环的理想 $\mathfrak{p}$；因此必有 $\mathfrak{p}\neq A$。
环 $A$ 至少有一个素理想，当且仅当 $A\neq\{0\}$。
\end{env}

\begin{env}[1.0.7]
\label{0.1.0.7-fr}
\emph{局部环}是恰有一个极大理想的环 $A$；这个极大理想就是可逆元集合的补集，并包含所有理想 $\neq A$。
若 $A$、$B$ 是两个局部环，$\mathfrak{m}$、$\mathfrak{n}$ 分别是它们的极大理想，则当 $\varphi(\mathfrak{m})\subset\mathfrak{n}$（等价地，$\varphi^{-1}(\mathfrak{n})=\mathfrak{m}$）时，称同态 $\varphi:A\to B$ 为\emph{局部同态}。
取商以后，这样的同态给出从剩余域 $A/\mathfrak{m}$ 到剩余域 $B/\mathfrak{n}$ 的单同态。
两个局部同态的复合仍是局部同态。
\end{env}

\subsection{理想的根。环的幂零根与 Jacobson 根}
\label{subsection:0.1.1-fr}

\begin{env}[1.1.1]
\label{0.1.1.1-fr}
设 $\mathfrak{a}$ 是环 $A$ 的一个理想；$\mathfrak{a}$ 的\emph{根}，记作 $\mathfrak{r}(\mathfrak{a})$，是所有满足 $x^n\in\mathfrak{a}$（至少对某个整数 $n>0$ 成立）的 $x\in A$ 所成的集合；它是包含 $\mathfrak{a}$ 的一个理想。
有 $\mathfrak{r}(\mathfrak{r}(\mathfrak{a}))=\mathfrak{r}(\mathfrak{a})$；关系 $\mathfrak{a}\subset\mathfrak{b}$ 蕴含 $\mathfrak{r}(\mathfrak{a})\subset\mathfrak{r}(\mathfrak{b})$；有限多个理想之交的根，等于这些理想的根之交。
若 $\varphi$ 是从环 $A'$ 到 $A$ 的一个同态，则对任一理想 $\mathfrak{a}\subset A$，有 $\mathfrak{r}(\varphi^{-1}(\mathfrak{a}))=\varphi^{-1}(\mathfrak{r}(\mathfrak{a}))$。
一个理想是某个理想的根，当且仅当它是一些素理想的交。
理想 $\mathfrak{a}$ 的根，是所有包含 $\mathfrak{a}$ 的素理想中那些\emph{极小}素理想的交；若 $A$ 是 Noether 环，则这样的极小素理想只有有限多个。

理想 $(0)$ 的根也称为 $A$ 的\emph{幂零根}；它就是由 $A$ 的幂零元组成的集合 $\mathfrak{N}$。
若 $\mathfrak{N}=(0)$，就称环 $A$ 是\emph{约化的}；对任一环 $A$，以其幂零根取商所得的商环 $A/\mathfrak{N}$，即 $A$ 对该幂零根的商，是约化环。
\end{env}

\begin{env}[1.1.2]
\label{0.1.1.2-fr}
回忆：一个（不必交换的）环 $A$ 的\emph{Jacobson 根} $\mathfrak{R}(A)$，是 $A$ 的所有极大左理想的交（也等于所有极大右理想的交）。
$A/\mathfrak{R}(A)$ 的 Jacobson 根为 $(0)$。
\end{env}

\subsection{分式模与分式环}
\label{subsection:0.1.2-fr}

\oldpage[0\textsubscript{I}]{13}
\begin{env}[1.2.1]
\label{0.1.2.1-fr}
若环 $A$ 的一个子集 $S$ 满足 $1\in S$，且 $S$ 中任意两个元素的乘积仍在 $S$ 中，就称该子集为\emph{乘法子集}。
以下最重要的例子是：$1^\circ$ 对元素 $f\in A$，由其各次幂 $f^n$（$n\geq 0$）组成的集合 $S_f$；$2^\circ$ 素理想 $\mathfrak{p}$ 在 $A$ 中的补集 $A-\mathfrak{p}$。
\end{env}

\begin{env}[1.2.2]
\label{0.1.2.2-fr}
设 $S$ 是环 $A$ 的一个乘法子集，$M$ 是一个 $A$-模；在集合 $M\times S$ 上，二元组 $(m_1,s_1)$ 与 $(m_2,s_2)$ 之间的关系
\[
  \text{“存在 $s\in S$，使得 $s(s_1m_2-s_2m_1)=0$”}
\]
是一个等价关系。
以 $S^{-1}M$ 表示 $M\times S$ 对此关系的商集，以 $m/s$ 表示二元组 $(m,s)$ 在 $S^{-1}M$ 中的典范像；把从 $M$ 到 $S^{-1}M$ 的映射 $i_M^S:m\mapsto m/1$（也记作 $i^S$）称为\emph{典范映射}。
一般说来，此映射既非单射也非满射；它的核是所有满足如下条件的 $m\in M$ 所成的集合：存在 $s\in S$ 使得 $sm=0$。

在 $S^{-1}M$ 上按下式定义加法群运算：
\[
  (m_1/s_1)+(m_2/s_2)=(s_2m_1+s_1m_2)/(s_1s_2)
\]
（可以验证，此定义确实与所考虑的 $S^{-1}M$ 中元素的表示无关）。
此外，在 $S^{-1}A$ 上按 $(a_1/s_1)(a_2/s_2)=(a_1a_2)/(s_1s_2)$ 定义乘法运算；最后，在 $S^{-1}M$ 上按 $(a/s)(m/s')=(am)/(ss')$ 定义一个以 $S^{-1}A$ 为算子集的外部运算。
由此可验证，$S^{-1}A$ 具有环结构（称为\emph{$A$ 的以 $S$ 中元素为分母的分式环}），而 $S^{-1}M$ 具有 $S^{-1}A$-模结构（称为\emph{$M$ 的以 $S$ 中元素为分母的分式模}）；对任一 $s\in S$，$s/1$ 在 $S^{-1}A$ 中可逆，其逆元为 $1/s$。
典范映射 $i_A^S$（相应地，$i_M^S$）是环同态（相应地，是 $A$-模同态；这里借助同态 $i_A^S:A\to S^{-1}A$，把 $S^{-1}M$ 看作 $A$-模）。
\end{env}

\begin{env}[1.2.3]
\label{0.1.2.3-fr}
若对某个 $f\in A$ 有 $S_f=\{f^n\}_{n\geq 0}$，则以 $A_f$、$M_f$ 分别代替 $S_f^{-1}A$、$S_f^{-1}M$；把 $A_f$ 看作 $A$ 上的代数时，可以写成 $A_f=A[1/f]$。
$A_f$ 同构于商代数 $A[T]/(fT-1)A[T]$。
当 $f=1$ 时，$A_f$、$M_f$ 分别与 $A$、$M$ 典范地等同；若 $f$ 是幂零元，则 $A_f$、$M_f$ 都退化为 $0$。

当 $S=A-\mathfrak{p}$，其中 $\mathfrak{p}$ 是 $A$ 的一个素理想时，以 $A_{\mathfrak{p}}$、$M_{\mathfrak{p}}$ 分别代替 $S^{-1}A$、$S^{-1}M$；$A_{\mathfrak{p}}$ 是一个\emph{局部环}，其极大理想 $\mathfrak{q}$ 由 $i_A^S(\mathfrak{p})$ 生成，并且有 $(i_A^S)^{-1}(\mathfrak{q})=\mathfrak{p}$；取商以后，$i_A^S$ 给出从整环 $A/\mathfrak{p}$ 到域 $A_{\mathfrak{p}}/\mathfrak{q}$ 的单同态，而后者可与 $A/\mathfrak{p}$ 的分式域等同。
\end{env}

\begin{env}[1.2.4]
\label{0.1.2.4-fr}
分式环 $S^{-1}A$ 与典范同态 $i_A^S$ 给出一个\emph{泛映射问题}的解：设 $u$ 是从 $A$ 到环 $B$ 的同态，并且 $u(S)$ 由 $B$ 中的\emph{可逆}元素组成，则它唯一地分解为
\[
  u:A\xrightarrow{i_A^S}S^{-1}A\xrightarrow{u^*}B
\]

\oldpage[0\textsubscript{I}]{14}
其中 $u^*$ 是环同态。
在同样的假设下，设 $M$ 是 $A$-模，$N$ 是 $B$-模，$v:M\to N$ 是 $A$-模同态（$N$ 上的 $B$-模结构由 $u:A\to B$ 规定）；则 $v$ 唯一地分解为
\[
  v:M\xrightarrow{i_M^S}S^{-1}M\xrightarrow{v^*}N
\]
其中 $v^*$ 是 $S^{-1}A$-模同态（$N$ 上的 $S^{-1}A$-模结构由 $u^*$ 规定）。
\end{env}

\begin{env}[1.2.5]
\label{0.1.2.5-fr}
定义 $S^{-1}A$-模的典范同构 $S^{-1}A\otimes_A M\xrightarrow{\sim}S^{-1}M$：把元素 $(a/s)\otimes m$ 对应到元素 $(am)/s$；其逆同构把 $m/s$ 映到 $(1/s)\otimes m$。
\end{env}

\begin{env}[1.2.6]
\label{0.1.2.6-fr}
对 $S^{-1}A$ 的任一理想 $\mathfrak{a}'$，$\mathfrak{a}=(i_A^S)^{-1}(\mathfrak{a}')$ 是 $A$ 的理想；而 $\mathfrak{a}'$ 是 $S^{-1}A$ 中由 $i_A^S(\mathfrak{a})$ 生成的理想，它可与 $S^{-1}\mathfrak{a}$ 等同 (1.3.2)。
映射 $\mathfrak{p}'\mapsto(i_A^S)^{-1}(\mathfrak{p}')$ 是一个保序同构：它把 $S^{-1}A$ 的\emph{素}理想集合映到 $A$ 的满足 $\mathfrak{p}\cap S=\varnothing$ 的素理想 $\mathfrak{p}$ 所成的集合上。
此外，此时局部环 $A_{\mathfrak{p}}$ 与 $(S^{-1}A)_{S^{-1}\mathfrak{p}}$ 典范地同构 (1.5.1)。
\end{env}

\begin{env}[1.2.7]
\label{0.1.2.7-fr}
当 $A$ 是\emph{整环}，并以 $K$ 表示其分式域时，对任一不含 $0$ 的乘法子集 $S$，典范映射 $i_A^S:A\to S^{-1}A$ 都是单射；此时 $S^{-1}A$ 可与 $K$ 的一个包含 $A$ 的子环典范地等同。
特别地，对 $A$ 的任一素理想 $\mathfrak{p}$，$A_{\mathfrak{p}}$ 是一个包含 $A$ 的局部环，其极大理想为 $\mathfrak{p}A_{\mathfrak{p}}$，并且有 $\mathfrak{p}A_{\mathfrak{p}}\cap A=\mathfrak{p}$。
\end{env}

\begin{env}[1.2.8]
\label{0.1.2.8-fr}
若 $A$ 是\emph{约化环} (1.1.1)，则 $S^{-1}A$ 也是约化环：事实上，若对 $x\in A$、$s\in S$ 有 $(x/s)^n=0$，这意味着存在 $s'\in S$ 使 $s'x^n=0$，从而 $(s'x)^n=0$；依假设可得 $s'x=0$，故 $x/s=0$。
\end{env}

\subsection{函子性质}
\label{subsection:0.1.3-fr}

\begin{env}[1.3.1]
\label{0.1.3.1-fr}
设 $M$、$N$ 是两个 $A$-模，$u$ 是一个 $A$-同态 $M\to N$。
若 $S$ 是 $A$ 的一个乘法子集，则可定义一个 $S^{-1}A$-同态 $S^{-1}M\to S^{-1}N$，记作 $S^{-1}u$，规定 $(S^{-1}u)(m/s)=u(m)/s$；若把 $S^{-1}M$ 与 $S^{-1}N$ 分别同 $S^{-1}A\otimes_A M$ 与 $S^{-1}A\otimes_A N$ 典范地等同 (1.2.5)，则 $S^{-1}u$ 与 $1\otimes u$ 等同。
若 $P$ 是第三个 $A$-模，$v$ 是一个 $A$-同态 $N\to P$，则有 $S^{-1}(v\circ u)=(S^{-1}v)\circ(S^{-1}u)$；换言之，$S^{-1}M$ 是关于 $M$ 的\emph{协变函子}，它从 $A$-模范畴取值于 $S^{-1}A$-模范畴（这里 $A$ 与 $S$ 固定）。
\end{env}

\begin{env}[1.3.2]
\label{0.1.3.2-fr}
函子 $S^{-1}M$ 是\emph{正合的}；换言之，若序列
\[
  M\xrightarrow{u}N\xrightarrow{v}P
\]
正合，则序列
\[
  S^{-1}M\xrightarrow{S^{-1}u}S^{-1}N\xrightarrow{S^{-1}v}S^{-1}P.
\]
也正合。
特别地，若 $u:M\to N$ 是单射（相应地，满射），则 $S^{-1}u$ 亦然；

\oldpage[0\textsubscript{I}]{15}
若 $N$ 与 $P$ 是 $M$ 的两个子模，则 $S^{-1}N$ 与 $S^{-1}P$ 可典范地等同于 $S^{-1}M$ 的子模，并且有
\[
  S^{-1}(N+P)=S^{-1}N+S^{-1}P
  \quad\text{且}\quad
  S^{-1}(N\cap P)=(S^{-1}N)\cap(S^{-1}P).
\]
\end{env}

\begin{env}[1.3.3]
\label{0.1.3.3-fr}
设 $(M_\alpha,\varphi_{\beta\alpha})$ 是一个 $A$-模归纳系，则 $(S^{-1}M_\alpha,S^{-1}\varphi_{\beta\alpha})$ 是一个 $S^{-1}A$-模归纳系。
把 $S^{-1}M_\alpha$ 与 $S^{-1}\varphi_{\beta\alpha}$ 表成张量积 (1.2.5 与 1.3.1)，由张量积运算与归纳极限运算的可交换性，得到典范同构
\[
  S^{-1}\varinjlim M_\alpha\xrightarrow{\sim}\varinjlim S^{-1}M_\alpha,
\]
这也表述为：函子 $S^{-1}M$（关于 $M$）\emph{与归纳极限可交换}。
\end{env}

\begin{env}[1.3.4]
\label{0.1.3.4-fr}
设 $M$、$N$ 是两个 $A$-模；存在关于 $M$ 与 $N$ 的典范\emph{函子性}同构
\[
  (S^{-1}M)\otimes_{S^{-1}A}(S^{-1}N)
  \xrightarrow{\sim}
  S^{-1}(M\otimes_A N)
\]
它把 $(m/s)\otimes(n/t)$ 变换为 $(m\otimes n)/st$。
\end{env}

\begin{env}[1.3.5]
\label{0.1.3.5-fr}
同样，存在一个关于 $M$ 与 $N$ 的\emph{函子性}同态
\[
  S^{-1}\operatorname{Hom}_A(M,N)
  \longrightarrow
  \operatorname{Hom}_{S^{-1}A}(S^{-1}M,S^{-1}N)
\]
它把 $u/s$ 对应到同态 $m/t\mapsto u(m)/st$。
当 $M$ 有\emph{有限表示}时，上述同态是\emph{同构}：$M$ 形如 $A^r$ 时这立即成立；一般情形则从正合列 $A^p\to A^q\to M\to 0$ 出发，利用函子 $S^{-1}M$ 的正合性以及函子 $\operatorname{Hom}_A(M,N)$ 关于 $M$ 的左正合性即可得到。
应注意，当 $A$ 是\emph{Noether 环}且 $A$-模 $M$ \emph{有限生成}时，总是属于这种情形。
\end{env}

\subsection{乘法子集的变换}
\label{subsection:0.1.4-fr}

\begin{env}[1.4.1]
\label{0.1.4.1-fr}
设 $S$、$T$ 是环 $A$ 的两个乘法子集，满足 $S\subset T$；存在典范同态 $\rho_A^{T,S}$（或简称 $\rho^{T,S}$）从 $S^{-1}A$ 到 $T^{-1}A$，它把 $S^{-1}A$ 中记作 $a/s$ 的元素对应到 $T^{-1}A$ 中记作 $a/s$ 的元素；并且有
\[
  i_A^T=\rho_A^{T,S}\circ i_A^S.
\]
对任一 $A$-模 $M$，同样存在一个从 $S^{-1}M$ 到 $T^{-1}M$ 的 $S^{-1}A$-线性映射（这里借助同态 $\rho_A^{T,S}$，把后者看作 $S^{-1}A$-模）；它把 $S^{-1}M$ 中的元素 $m/s$ 对应到 $T^{-1}M$ 中的元素 $m/s$；把此映射记作 $\rho_M^{T,S}$，或简称 $\rho^{T,S}$，并且仍有
\[
  i_M^T=\rho_M^{T,S}\circ i_M^S;
\]
在典范等同 (1.2.5) 下，$\rho_M^{T,S}$ 与 $\rho_A^{T,S}\otimes 1$ 等同。
同态 $\rho_M^{T,S}$ 是从函子 $S^{-1}M$ 到函子 $T^{-1}M$ 的一个\emph{函子性态射}（或自然变换）；换言之，图表
\[
\begin{tikzcd}[column sep=large,row sep=large]
S^{-1}M \arrow[r,"S^{-1}u"] \arrow[d,"\rho_M^{T,S}"'] & S^{-1}N \arrow[d,"\rho_N^{T,S}"] \\
T^{-1}M \arrow[r,"T^{-1}u"] & T^{-1}N
\end{tikzcd}
\]

\oldpage[0\textsubscript{I}]{16}
对任一同态 $u:M\to N$ 都可交换；还应注意，$T^{-1}u$ 完全由 $S^{-1}u$ 确定，因为对 $m\in M$ 与 $t\in T$，有
\[
  (T^{-1}u)(m/t)=(t/1)^{-1}\rho^{T,S}\bigl((S^{-1}u)(m/1)\bigr).
\]
\end{env}

\begin{env}[1.4.2]
\label{0.1.4.2-fr}
沿用上述记号，对两个 $A$-模 $M$、$N$，下列图表（参见 (1.3.4) 与 (1.3.5)）
\[
\begin{tikzcd}[column sep=large,row sep=large]
(S^{-1}M)\otimes_{S^{-1}A}(S^{-1}N) \arrow[r,"\sim"] \arrow[d]
  & S^{-1}(M\otimes_A N) \arrow[d]
  \\
(T^{-1}M)\otimes_{T^{-1}A}(T^{-1}N) \arrow[r,"\sim"]
  & T^{-1}(M\otimes_A N)
\end{tikzcd}
\]
\[
\begin{tikzcd}[column sep=large,row sep=large]
S^{-1}\operatorname{Hom}_A(M,N) \arrow[r] \arrow[d]
  & \operatorname{Hom}_{S^{-1}A}(S^{-1}M,S^{-1}N) \arrow[d] \\
T^{-1}\operatorname{Hom}_A(M,N) \arrow[r]
  & \operatorname{Hom}_{T^{-1}A}(T^{-1}M,T^{-1}N)
\end{tikzcd}
\]
都是可交换的。
\end{env}

\begin{env}[1.4.3]
\label{0.1.4.3-fr}
在一种重要情形下，同态 $\rho^{T,S}$ 是\emph{双射}：即 $T$ 的每个元素都是 $S$ 的某个元素的因子；此时借助 $\rho^{T,S}$ 把模 $S^{-1}M$ 与 $T^{-1}M$ 等同起来。
若 $S$ 中某元素在 $A$ 中的每个因子都属于 $S$，就称 $S$ 是\emph{饱和的}；以 $S$ 中各元素的所有因子组成的集合 $T$（此集合是乘法的且饱和）替换 $S$，可见若有需要，总可以只考虑 $S$ 饱和时的分式模 $S^{-1}M$。
\end{env}

\begin{env}[1.4.4]
\label{0.1.4.4-fr}
若 $S$、$T$、$U$ 是 $A$ 的三个乘法子集，满足 $S\subset T\subset U$，则有
\[
  \rho^{U,S}=\rho^{U,T}\circ\rho^{T,S}.
\]
\end{env}

\begin{env}[1.4.5]
\label{0.1.4.5-fr}
考虑 $A$ 的乘法子集所成的一个\emph{递增滤过族} $(S_\alpha)$（用 $\alpha\leq\beta$ 表示 $S_\alpha\subset S_\beta$），并令 $S$ 为乘法子集 $\bigcup_\alpha S_\alpha$；规定
\[
  \rho_{\beta\alpha}=\rho_A^{S_\beta,S_\alpha}\quad\text{当 }\alpha\leq\beta;
\]
由 (1.4.4)，同态 $\rho_{\beta\alpha}$ 定义了一个环 $A'$，它是环的归纳系 $(S_\alpha^{-1}A,\rho_{\beta\alpha})$ 的\emph{归纳极限}。
设 $\rho_\alpha$ 是典范映射 $S_\alpha^{-1}A\to A'$，并规定 $\varphi_\alpha=\rho_A^{S,S_\alpha}$；由 (1.4.4)，当 $\alpha\leq\beta$ 时有 $\varphi_\alpha=\varphi_\beta\circ\rho_{\beta\alpha}$，故可唯一地定义一个同态 $\varphi:A'\to S^{-1}A$，使图表
\[
\begin{tikzcd}[column sep=huge,row sep=large]
  & S_\alpha^{-1}A
      \arrow[ddl,"\rho_\alpha"']
      \arrow[d,"\rho_{\beta\alpha}"]
      \arrow[ddr,"\varphi_\alpha"] & \\
  & S_\beta^{-1}A
      \arrow[dl,"\rho_\beta"]
      \arrow[dr,"\varphi_\beta"'] & \\
A' \arrow[rr,"\varphi"'] & & S^{-1}A
\end{tikzcd}
\qquad(\alpha\leq\beta)
\]
可交换。
事实上，$\varphi$ 是\emph{同构}：由构造立即可见 $\varphi$ 是满射。
另一方面，若 $\rho_\alpha(a/s_\alpha)\in A'$ 满足 $\varphi(\rho_\alpha(a/s_\alpha))=0$，这意味着 $a/s_\alpha=0$ 在 $S^{-1}A$ 中成立，即存在 $s\in S$ 使 $sa=0$；但存在某个 $\beta\geq\alpha$ 使 $s\in S_\beta$，因而由 $\rho_\alpha(a/s_\alpha)=\rho_\beta(sa/ss_\alpha)=0$ 可知 $\varphi$ 是单射。
对 $A$-模 $M$ 的情形作同样处理，于是定义了典范同构
\[
  \varinjlim S_\alpha^{-1}A\xrightarrow{\sim}
  (\varinjlim S_\alpha)^{-1}A,
  \qquad
  \varinjlim S_\alpha^{-1}M\xrightarrow{\sim}
  (\varinjlim S_\alpha)^{-1}M,
\]
其中第二个关于 $M$ 是\emph{函子性的}。
\end{env}

\oldpage[0\textsubscript{I}]{17}
\begin{env}[1.4.6]
\label{0.1.4.6-fr}
设 $S_1$、$S_2$ 是 $A$ 的两个乘法子集；则 $S_1S_2$ 也是 $A$ 的乘法子集。
以 $S_2'$ 表示 $S_2$ 在环 $S_1^{-1}A$ 中的典范像；它是该环的一个乘法子集。
对任一 $A$-模 $M$，存在函子性同构
\[
  S_2'^{-1}(S_1^{-1}M)\xrightarrow{\sim}(S_1S_2)^{-1}M
\]
它把 $(m/s_1)/(s_2/1)$ 对应到元素 $m/(s_1s_2)$。
\end{env}

\subsection{环的变换}
\label{subsection:0.1.5-fr}

\begin{env}[1.5.1]
\label{0.1.5.1-fr}
设 $A$、$A'$ 是两个环，$\varphi$ 是一个同态 $A'\to A$，$S$（相应地，$S'$）是 $A$（相应地，$A'$）的一个乘法子集，且满足 $\varphi(S')\subset S$；复合同态
\[
  A'\xrightarrow{\varphi}A\longrightarrow S^{-1}A
\]
分解为
\[
  A'\longrightarrow S'^{-1}A'\xrightarrow{\varphi^{S'}}S^{-1}A
\]
这是 (1.2.4) 的结论；并且有 $\varphi^{S'}(a'/s')=\varphi(a')/\varphi(s')$。
若 $A=\varphi(A')$ 且 $S=\varphi(S')$，则 $\varphi^{S'}$ 是\emph{满射}。
若 $A'=A$ 且 $\varphi$ 是恒等映射，则 $\varphi^{S'}$ 正是 (1.4.1) 中定义的同态 $\rho_A^{S,S'}$。
\end{env}

\begin{env}[1.5.2]
\label{0.1.5.2-fr}
在 (1.5.1) 的假设下，设 $M$ 是一个 $A$-模。
存在典范函子性同态
\[
  \sigma:S'^{-1}(M_{[\varphi]})\longrightarrow
  (S^{-1}M)_{[\varphi^{S'}]}
\]
它是 $S'^{-1}A'$-模同态，并把 $S'^{-1}(M_{[\varphi]})$ 的每个元素 $m/s'$ 对应到 $(S^{-1}M)_{[\varphi^{S'}]}$ 的元素 $m/\varphi(s')$；事实上，立即可验证此定义与所考察元素的表示 $m/s'$ 无关。
当 $S=\varphi(S')$ 时，同态 $\sigma$ 是\emph{双射}。
当 $A'=A$ 且 $\varphi$ 是恒等映射时，$\sigma$ 正是 (1.4.1) 中定义的同态 $\rho_M^{S,S'}$。

特别取 $M=A$ 时，同态 $\varphi$ 在 $A$ 上定义一个 $A'$-代数结构；此时 $S'^{-1}(A_{[\varphi]})$ 具有环结构，在此结构下可与 $(\varphi(S'))^{-1}A$ 等同，并且同态
\[
  \sigma:S'^{-1}(A_{[\varphi]})\longrightarrow S^{-1}A
\]
是 $S'^{-1}A'$-代数同态。
\end{env}

\begin{env}[1.5.3]
\label{0.1.5.3-fr}
设 $M$、$N$ 是两个 $A$-模；复合 (1.3.4) 与 (1.5.2) 中定义的同态，得到同态
\[
  (S^{-1}M\otimes_{S^{-1}A}S^{-1}N)_{[\varphi^{S'}]}
  \longleftarrow
  S'^{-1}((M\otimes_A N)_{[\varphi]})
\]
当 $\varphi(S')=S$ 时，它是同构。
同样，复合 (1.3.5) 与 (1.5.2) 中的同态，得到同态
\[
  S'^{-1}((\operatorname{Hom}_A(M,N))_{[\varphi]})
  \longrightarrow
  (\operatorname{Hom}_{S^{-1}A}(S^{-1}M,S^{-1}N))_{[\varphi^{S'}]}
\]
当 $\varphi(S')=S$ 且 $M$ 有有限表示时，它是同构。
\end{env}

\begin{env}[1.5.4]
\label{0.1.5.4-fr}
现在考虑一个 $A'$-模 $N'$，并作张量积 $N'\otimes_{A'}A_{[\varphi]}$；按下式规定后，可把它看作 $A$-模：
\[
  a.(n'\otimes b)=n'\otimes(ab).
\]
存在 $S^{-1}A$-模的函子性同构
\[
  \tau:(S'^{-1}N')\otimes_{S'^{-1}A'}(S^{-1}A)_{[\varphi^{S'}]}
  \xrightarrow{\sim}
  S^{-1}(N'\otimes_{A'}A_{[\varphi]})
\]
\oldpage[0\textsubscript{I}]{18}
它把元素 $(n'/s')\otimes(a/s)$ 对应到元素
$(n'\otimes a)/(\varphi(s')s)$；事实上，可以分别验证，当把 $n'/s'$（相应地，$a/s$）换成同一元素的另一种表示时，
$(n'\otimes a)/(\varphi(s')s)$ 不变；另一方面，可以定义 $\tau$ 的逆同态，把 $(n'\otimes a)/s$ 对应到元素 $(n'/1)\otimes(a/s)$：这里利用了
$S^{-1}(N'\otimes_{A'}A_{[\varphi]})$ 与
$(N'\otimes_{A'}A_{[\varphi]})\otimes_A S^{-1}A$ 典范同构 (1.2.5)，因而也与
$N'\otimes_{A'}(S^{-1}A)_{[\psi]}$ 典范同构；其中以 $\psi$ 表示从 $A'$ 到 $S^{-1}A$ 的复合同态 $a'\mapsto\varphi(a')/1$。
\end{env}

\begin{env}[1.5.5]
\label{0.1.5.5-fr}
若 $M'$、$N'$ 是两个 $A'$-模，复合 (1.3.4) 与 (1.5.4) 中的同构，得到同构
\[
  S'^{-1}M'\otimes_{S'^{-1}A'}S'^{-1}N'
  \otimes_{S'^{-1}A'}S^{-1}A
  \xrightarrow{\sim}
  S^{-1}(M'\otimes_{A'}N'\otimes_{A'}A).
\]
同样，若 $M'$ 有有限表示，则由 (1.3.5) 与 (1.5.4) 得到同构
\[
  \operatorname{Hom}_{S'^{-1}A'}(S'^{-1}M',S'^{-1}N')
  \otimes_{S'^{-1}A'}S^{-1}A
  \xrightarrow{\sim}
  S^{-1}(\operatorname{Hom}_{A'}(M',N')\otimes_{A'}A).
\]
\end{env}

\begin{env}[1.5.6]
\label{0.1.5.6-fr}
在 (1.5.1) 的假设下，设 $T$（相应地，$T'$）是 $A$（相应地，$A'$）的第二个乘法子集，满足 $S\subset T$（相应地，$S'\subset T'$）且 $\varphi(T')\subset T$。则图表
\[
\begin{tikzcd}[column sep=huge,row sep=large]
S'^{-1}A' \arrow[r,"\varphi^{S'}"] \arrow[d,"\rho^{T',S'}"']
  & S^{-1}A \arrow[d,"\rho^{T,S}"] \\
T'^{-1}A' \arrow[r,"\varphi^{T'}"']
  & T^{-1}A
\end{tikzcd}
\]
可交换。若 $M$ 是一个 $A$-模，则图表
\[
\begin{tikzcd}[column sep=huge,row sep=large]
S'^{-1}(M_{[\varphi]}) \arrow[r,"\sigma"] \arrow[d,"\rho^{T',S'}"']
  & (S^{-1}M)_{[\varphi^{S'}]} \arrow[d,"\rho^{T,S}"] \\
T'^{-1}(M_{[\varphi]}) \arrow[r,"\sigma"']
  & (T^{-1}M)_{[\varphi^{T'}]}
\end{tikzcd}
\]
可交换。最后，若 $N'$ 是一个 $A'$-模，则图表
\[
\begin{tikzcd}[column sep=huge,row sep=large]
(S'^{-1}N')\otimes_{S'^{-1}A'}(S^{-1}A)_{[\varphi^{S'}]}
  \arrow[r,"\tau","{\sim}"'] \arrow[d]
  & S^{-1}(N'\otimes_{A'}A_{[\varphi]}) \arrow[d,"\rho^{T,S}"] \\
(T'^{-1}N')\otimes_{T'^{-1}A'}(T^{-1}A)_{[\varphi^{T'}]}
  \arrow[r,"\tau"',"{\sim}"]
  & T^{-1}(N'\otimes_{A'}A_{[\varphi]})
\end{tikzcd}
\]
可交换；其中左侧竖直箭头是把 $\rho_{N'}^{T',S'}$ 作用于 $S'^{-1}N'$，并把 $\rho_A^{T,S}$ 作用于 $S^{-1}A$ 而得到的。
\end{env}

\oldpage[0\textsubscript{I}]{19}

\begin{env}[1.5.7]
\label{0.1.5.7-fr}
设 $A''$ 是第三个环，$\varphi':A''\to A'$ 是环同态，$S''$ 是 $A''$ 的一个乘法子集，满足 $\varphi'(S'')\subset S'$。规定 $\varphi''=\varphi\circ\varphi'$；则有
\[
  \varphi''^{S''}=\varphi^{S'}\circ\varphi'^{S''}.
\]
设 $M$ 是一个 $A$-模；显然有 $M_{[\varphi'']}=(M_{[\varphi]})_{[\varphi']}$；若 $\sigma'$ 与 $\sigma''$ 分别由 $\varphi'$ 与 $\varphi''$ 按照 (1.5.2) 中由 $\varphi$ 定义 $\sigma$ 的方式定义，则有传递性公式
\[
  \sigma''=\sigma\circ\sigma'.
\]
最后，设 $N''$ 是一个 $A''$-模；$A$-模
$N''\otimes_{A''}A_{[\varphi'']}$ 可典范地等同于
\[
  (N''\otimes_{A''}A'_{[\varphi']})\otimes_{A'}A_{[\varphi]},
\]
同样，$S^{-1}A$-模
\[
  (S''^{-1}N'')\otimes_{S''^{-1}A''}
  (S^{-1}A)_{[\varphi''^{S''}]}
\]
可典范地等同于
\[
  \bigl((S''^{-1}N'')\otimes_{S''^{-1}A''}
  (S'^{-1}A')_{[\varphi'^{S''}]}\bigr)
  \otimes_{S'^{-1}A'}(S^{-1}A)_{[\varphi^{S'}]}.
\]
在这些等同下，若 $\tau'$ 与 $\tau''$ 分别由 $\varphi'$ 与 $\varphi''$ 按照 (1.5.4) 中由 $\varphi$ 定义 $\tau$ 的方式定义，则有传递性公式
\[
  \tau''=\tau\circ(\tau'\otimes 1).
\]
\end{env}

\begin{env}[1.5.8]
\label{0.1.5.8-fr}
设 $A$ 是环 $B$ 的一个子环；对 $A$ 的任一\emph{极小}素理想 $\mathfrak p$，存在 $B$ 的一个极小素理想 $\mathfrak q$，使得 $\mathfrak p=A\cap\mathfrak q$。事实上，$A_{\mathfrak p}$ 是 $B_{\mathfrak p}$ 的一个子环 (1.3.2)，并且只有一个素理想 $\mathfrak p'$ (1.2.6)；由于 $B_{\mathfrak p}$ 不退化为 $0$，它至少有一个素理想 $\mathfrak q'$，而且必有 $\mathfrak q'\cap A_{\mathfrak p}=\mathfrak p'$；因此，$B$ 的素理想 $\mathfrak q_1$——即 $\mathfrak q'$ 的原像——满足 $\mathfrak q_1\cap A=\mathfrak p$；\emph{更进一步}，对 $B$ 的任一包含于 $\mathfrak q_1$ 的极小素理想 $\mathfrak q$，都有 $\mathfrak q\cap A=\mathfrak p$。
\end{env}

\subsection{模 $M_f$ 与一个归纳极限的等同}
\label{subsection:0.1.6-fr}

\begin{env}[1.6.1]
\label{0.1.6.1-fr}
设 $M$ 是一个 $A$-模，$f$ 是 $A$ 的一个元素。考虑一列 $A$-模 $(M_n)$，它们都与 $M$ 相同；对任意一对整数 $m\leq n$，令 $\varphi_{nm}$ 为从 $M_m$ 到 $M_n$ 的同态 $z\mapsto f^{n-m}z$；立即可见，$((M_n),(\varphi_{nm}))$ 是一个 $A$-模\emph{归纳系}；令 $N=\varinjlim M_n$ 为该系的归纳极限。我们将定义从 $N$ 到 $M_f$ 的一个典范\emph{函子性} $A$-同构。为此，注意对任一 $n$，$\theta_n:z\mapsto z/f^n$ 是从 $M=M_n$ 到 $M_f$ 的一个 $A$-同态；由定义可知，当 $m\leq n$ 时有 $\theta_n\circ\varphi_{nm}=\theta_m$。因此，存在一个 $A$-同态 $\theta:N\to M_f$，使得：若 $\varphi_n$ 表示典范同态 $M_n\to N$，则对任一 $n$ 都有 $\theta_n=\theta\circ\varphi_n$。依假设，$M_f$ 的每个元素至少可对某个 $n$ 写成 $z/f^n$，故 $\theta$ 显然是满射。另一方面，若 $\theta(\varphi_n(z))=0$，即 $z/f^n=0$，则存在整数 $k>0$ 使 $f^kz=0$，因而 $\varphi_{n+k,n}(z)=0$，这又推出 $\varphi_n(z)=0$。所以可借助 $\theta$ 把 $M_f$ 与 $\varinjlim M_n$ 等同起来。
\end{env}

\begin{env}[1.6.2]
\label{0.1.6.2-fr}
现在以 $M_{f,n}$、$\varphi_{nm}^f$、$\varphi_n^f$ 分别代替 $M_n$、$\varphi_{nm}$、$\varphi_n$。设 $g$ 是 $A$ 的第二个元素。由于 $f^n$ 整除 $f^ng^n$，存在函子性同态
\[
  \rho_{fg,f}:M_f\longrightarrow M_{fg}\qquad\text{(1.4.1 与 1.4.3)};
\]

\oldpage[0\textsubscript{I}]{20}

若把 $M_f$ 与 $M_{fg}$ 分别同 $\varinjlim M_{f,n}$ 与 $\varinjlim M_{fg,n}$ 等同，则 $\rho_{fg,f}$ 与映射 $\rho_{fg,f}^n:M_{f,n}\to M_{fg,n}$ 的\emph{归纳极限}等同，其中 $\rho_{fg,f}^n(z)=g^nz$。事实上，这立即来自图表
\[
\begin{tikzcd}[column sep=huge,row sep=large]
M_{f,n} \arrow[r,"\rho_{fg,f}^n"] \arrow[d,"\varphi_n^f"']
  & M_{fg,n} \arrow[d,"\varphi_n^{fg}"] \\
M_f \arrow[r,"\rho_{fg,f}"]
  & M_{fg}.
\end{tikzcd}
\]
的可交换性。
\end{env}

\subsection{模的支撑}
\label{subsection:0.1.7-fr}

\begin{env}[1.7.1]
\label{0.1.7.1-fr}
给定 $A$-模 $M$，把 $M$ 的\emph{支撑}定义为 $A$ 的所有满足 $M_{\mathfrak p}\neq0$ 的素理想 $\mathfrak p$ 所成的集合，并记作 $\operatorname{Supp}(M)$。$M=0$ 的充要条件是 $\operatorname{Supp}(M)=\varnothing$；因为若对每个 $\mathfrak p$ 都有 $M_{\mathfrak p}=0$，则元素 $x\in M$ 的零化子不可能包含于 $A$ 的任何素理想中，因而必为整个 $A$。
\end{env}

\begin{env}[1.7.2]
\label{0.1.7.2-fr}
若 $0\to N\to M\to P\to0$ 是 $A$-模正合列，则有
\[
  \operatorname{Supp}(M)=\operatorname{Supp}(N)\cup\operatorname{Supp}(P)
\]
因为对 $A$ 的每个素理想 $\mathfrak p$，序列 $0\to N_{\mathfrak p}\to M_{\mathfrak p}\to P_{\mathfrak p}\to0$ 都正合 (1.3.2)，而 $M_{\mathfrak p}=0$ 的充要条件是 $N_{\mathfrak p}=P_{\mathfrak p}=0$。
\end{env}

\begin{env}[1.7.3]
\label{0.1.7.3-fr}
若 $M$ 是一族子模 $(M_\lambda)$ 之和，则对 $A$ 的每个素理想 $\mathfrak p$，$M_{\mathfrak p}$ 是各 $(M_\lambda)_{\mathfrak p}$ 之和 (1.3.3 与 1.3.2)，因此
\[
  \operatorname{Supp}(M)=\bigcup_\lambda\operatorname{Supp}(M_\lambda).
\]
\end{env}

\begin{env}[1.7.4]
\label{0.1.7.4-fr}
若 $M$ 是一个\emph{有限生成} $A$-模，则 $\operatorname{Supp}(M)$ 是所有\emph{包含 $M$ 的零化子的}素理想所成的集合。事实上，若 $M$ 是由 $x$ 生成的循环模，则说 $M_{\mathfrak p}=0$，就等价于存在 $s\notin\mathfrak p$ 使 $s.x=0$，从而等价于 $\mathfrak p$ 不包含 $x$ 的零化子。现在，若 $M$ 有一个有限生成元系 $(x_i)_{1\leq i\leq n}$，且 $\mathfrak a_i$ 是 $x_i$ 的零化子，则由 (1.7.3)，$\operatorname{Supp}(M)$ 是所有包含某个 $\mathfrak a_i$ 的 $\mathfrak p$ 所成的集合；等价地，它是所有包含 $\mathfrak a=\bigcap_i\mathfrak a_i$ 的 $\mathfrak p$ 所成的集合，而后者正是 $M$ 的零化子。
\end{env}

\begin{env}[1.7.5]
\label{0.1.7.5-fr}
若 $M$、$N$ 是两个\emph{有限生成} $A$-模，则有
\[
  \operatorname{Supp}(M\otimes_A N)
  =\operatorname{Supp}(M)\cap\operatorname{Supp}(N).
\]
需要证明的是：若 $\mathfrak p$ 是 $A$ 的一个素理想，则条件 $M_{\mathfrak p}\otimes_{A_{\mathfrak p}}N_{\mathfrak p}\neq0$ 等价于“$M_{\mathfrak p}\neq0$ 且 $N_{\mathfrak p}\neq0$”（计及 (1.3.4)）。换言之，需要证明：若 $P$、$Q$ 是局部环 $B$ 上两个不退化为 $0$ 的有限生成模，则 $P\otimes_B Q\neq0$。设 $\mathfrak m$ 是 $B$ 的极大理想。由 Nakayama 引理，向量空间 $P/\mathfrak mP$ 与 $Q/\mathfrak mQ$ 都不退化为 $0$，故其张量积
\[
  (P/\mathfrak mP)\otimes_{B/\mathfrak m}(Q/\mathfrak mQ)
  =(P\otimes_B Q)\otimes_B(B/\mathfrak m),
\]
也非零，结论随即得到。

特别地，若 $M$ 是一个有限生成 $A$-模，$\mathfrak a$ 是 $A$ 的一个理想，则 $\operatorname{Supp}(M/\mathfrak aM)$ 是所有同时包含 $\mathfrak a$ 与 $M$ 的零化子 $\mathfrak n$ 的素理想所成的集合 (1.7.4)；换言之，它是所有包含 $\mathfrak a+\mathfrak n$ 的素理想所成的集合。
\end{env}

\oldpage[0\textsubscript{I}]{21}

\section{不可约空间；Noether 空间}
\label{section:0.2-fr}

\subsection{不可约空间}
\label{subsection:0.2.1-fr}

\begin{env}[2.1.1]
\label{0.2.1.1-fr}
称拓扑空间 $X$ 是\emph{不可约的}，如果它非空，并且不能表示为 $X$ 的两个均异于 $X$ 的闭子空间之并。等价地，也可以说 $X\neq\varnothing$，且 $X$ 中任意两个（从而任意有限多个）非空开集的交非空；或者说每个非空开集都稠密；或者说每个闭子集 $\neq X$ 都\emph{无处稠密}；又或者说 $X$ 的每个开集都\emph{连通}。
\end{env}

\begin{env}[2.1.2]
\label{0.2.1.2-fr}
拓扑空间 $X$ 的子空间 $Y$ 不可约，当且仅当其闭包 $\overline Y$ 不可约。特别地，凡是某个单点子空间的闭包 $\overline{\{x\}}$，都是不可约的；我们把关系 $y\in\overline{\{x\}}$（等价于 $\overline{\{y\}}\subset\overline{\{x\}}$）表述为：$y$ 是 $x$ 的\emph{特化}，或 $x$ 是 $y$ 的\emph{一般化}。若不可约空间 $X$ 中存在一点 $x$，使得 $X=\overline{\{x\}}$，则称 $x$ 为 $X$ 的\emph{泛点}。这时，$X$ 的每个非空开集都包含 $x$，而每个包含 $x$ 的子空间也以 $x$ 为泛点。
\end{env}

\begin{env}[2.1.3]
\label{0.2.1.3-fr}
回忆：若拓扑空间 $X$ 满足下列分离公理，就称它为 \emph{Kolmogorov 空间}：

\smallskip
\noindent
\textup{(T$_0$)} 若 $x\neq y$ 是 $X$ 中任意两个点，则存在一个开集，它包含 $x$、$y$ 中的一点而不包含另一点。

不可约 Kolmogorov 空间若有泛点，就\emph{只能有一个}，因为每个非空开集都包含所有泛点。

回忆：若 $X$ 的任一开覆盖都能抽出一个有限子覆盖（等价地，任一由非空闭集组成的滤过递减族都有非空交），则称拓扑空间 $X$ \emph{拟紧}。若 $X$ 拟紧，则 $X$ 的每个非空闭子集 $A$ 都包含一个\emph{极小}非空闭集 $M$，因为 $A$ 的非空闭子集全体对关系 $\supset$ 构成归纳有序集；若 $X$ 还是 Kolmogorov 空间，则 $M$ 必定只含一个点（借用术语，也称它为一个\emph{闭点}）。
\end{env}

\begin{env}[2.1.4]
\label{0.2.1.4-fr}
不可约空间 $X$ 的每个非空开子空间 $U$ 都不可约；若 $X$ 有泛点 $x$，则 $x$ 也是 $U$ 的泛点。

设 $(U_\alpha)$ 是拓扑空间 $X$ 的一个由非空开集组成的覆盖，其指标集非空。则 $X$ 不可约，当且仅当对每个 $\alpha$，$U_\alpha$ 都不可约，并且对任意 $\alpha$、$\beta$ 都有 $U_\alpha\cap U_\beta\neq\varnothing$。条件的必要性显然。为证明充分性，只需证明：若 $V$ 是 $X$ 的非空开集，则对每个 $\alpha$，$V\cap U_\alpha$ 都非空；这样，对每个 $\alpha$，$V\cap U_\alpha$ 都在 $U_\alpha$ 中稠密，从而 $V$ 在 $X$ 中稠密。事实上，至少有一个指标 $\gamma$ 使 $V\cap U_\gamma\neq\varnothing$，所以 $V\cap U_\gamma$ 在 $U_\gamma$ 中稠密；又因对每个 $\alpha$ 都有 $U_\alpha\cap U_\gamma\neq\varnothing$，故也有 $V\cap U_\alpha\cap U_\gamma\neq\varnothing$。
\end{env}

\oldpage[0\textsubscript{I}]{22}

\begin{env}[2.1.5]
\label{0.2.1.5-fr}
设 $X$ 是不可约空间，$f$ 是从 $X$ 到拓扑空间 $Y$ 的连续映射。则 $f(X)$ 不可约；若 $x$ 是 $X$ 的泛点，则 $f(x)$ 是 $f(X)$ 的泛点，因而也是 $\overline{f(X)}$ 的泛点。特别地，若 $Y$ 也不可约且只有一个泛点 $y$，则 $f(X)$ 稠密，当且仅当 $f(x)=y$。
\end{env}

\begin{env}[2.1.6]
\label{0.2.1.6-fr}
拓扑空间 $X$ 的每个不可约子空间，都包含于一个\emph{极大}不可约子空间中，而后者必为闭集。$X$ 的极大不可约子空间称为 $X$ 的\emph{不可约分支}。若 $Z_1$、$Z_2$ 是 $X$ 的两个不同不可约分支，则 $Z_1\cap Z_2$ 在子空间 $Z_1$、$Z_2$ 中分别都是闭的\emph{无处稠密}集；特别地，若 $X$ 的一个不可约分支有泛点 (2.1.2)，则该点不可能属于其他任何不可约分支。若 $X$ 只有\emph{有限}个不可约分支 $Z_i$（$1\leq i\leq n$），并且对每个 $i$ 置
\[
  U_i=Z_i\cap\complement\Bigl(\bigcup_{j\neq i}Z_j\Bigr),
\]
则各 $U_i$ 都开且不可约，两两不交，并且它们的并在 $X$ 中稠密。

设 $U$ 是拓扑空间 $X$ 的一个开子集。若 $Z$ 是 $X$ 中与 $U$ 相交的不可约子集，则 $Z\cap U$ 在 $Z$ 中开且稠密，因而不可约；反过来，对 $U$ 的任一闭不可约子集 $Y$，$Y$ 在 $X$ 中的闭包 $\overline Y$ 不可约，且 $\overline Y\cap U=Y$。由此可知，$U$ 的不可约分支与 $X$ 中同 $U$ 相交的不可约分支之间存在\emph{双射对应}。
\end{env}

\begin{env}[2.1.7]
\label{0.2.1.7-fr}
若拓扑空间 $X$ 是有限多个闭不可约子空间 $Y_i$ 之并，则 $X$ 的不可约分支正是各 $Y_i$ 中的极大元。事实上，若 $Z$ 是 $X$ 的闭不可约子集，则 $Z$ 是各 $Z\cap Y_i$ 之并，立即可见 $Z$ 必包含于某个 $Y_i$。设 $Y$ 是拓扑空间 $X$ 的一个子空间，并假设 $Y$ 只有有限多个不可约分支 $Y_i$（$1\leq i\leq n$）；则它们在 $X$ 中的各闭包 $\overline{Y_i}$，正是 $\overline Y$ 的不可约分支。
\end{env}

\begin{env}[2.1.8]
\label{0.2.1.8-fr}
设 $Y$ 是只有一个泛点 $y$ 的不可约空间，$X$ 是拓扑空间，$f$ 是从 $X$ 到 $Y$ 的连续映射。则对 $X$ 中每个同 $f^{-1}(y)$ 相交的不可约分支 $Z$，$f(Z)$ 都在 $Y$ 中稠密。反命题未必成立；但是，若 $Z$ 有泛点 $z$，并且 $f(Z)$ 在 $Y$ 中稠密，则必有 $f(z)=y$ (2.1.5)；此外，$Z\cap f^{-1}(y)$ 是 $\{z\}$ 在 $f^{-1}(y)$ 中的闭包，因而不可约；又因为 $f^{-1}(y)$ 中任何包含 $z$ 的不可约子集必包含于 $Z$ (2.1.6)，所以 $z$ 是 $Z\cap f^{-1}(y)$ 的泛点。由于 $f^{-1}(y)$ 的每个不可约分支都包含于 $X$ 的一个不可约分支中，可知：若 $X$ 中同 $f^{-1}(y)$ 相交的每个不可约分支 $Z$ 都有泛点，则这些分支与 $f^{-1}(y)$ 的各不可约分支 $Z\cap f^{-1}(y)$ 之间存在\emph{双射对应}；$Z$ 与 $Z\cap f^{-1}(y)$ 的泛点相同。
\end{env}

\oldpage[0\textsubscript{I}]{23}

\subsection{Noether 空间}
\label{subsection:0.2.2-fr}

\begin{env}[2.2.1]
\label{0.2.2.1-fr}
若拓扑空间 $X$ 的开集全体满足\emph{极大条件}，或者等价地，$X$ 的闭集全体满足\emph{极小条件}，则称 $X$ 为 \emph{Noether 空间}。若每个 $x\in X$ 都有一个作为 Noether 子空间的邻域，则称 $X$ \emph{局部 Noether}。
\end{env}

\begin{env}[2.2.2]
\label{0.2.2.2-fr}
设 $E$ 是满足\emph{极小条件}的有序集，$P$ 是 $E$ 的元素所具有的一个性质，并满足下列条件：若 $a\in E$，且对每个 $x<a$，$P(x)$ 都成立，则 $P(a)$ 成立。在这些条件下，\emph{对每个 $x\in E$，$P(x)$ 都成立}（“Noether 归纳原理”）。事实上，设 $F$ 是使 $P(x)$ 不成立的 $x\in E$ 所成的集合；若 $F$ 非空，它就有一个极小元 $a$，于是对每个 $x<a$，$P(x)$ 都成立，故 $P(a)$ 也成立，这就产生了矛盾。

特别地，我们会在 $E$ 是一个 \emph{Noether 空间的闭子集所成的集合}时应用这一原理。
\end{env}

\begin{env}[2.2.3]
\label{0.2.2.3-fr}
Noether 空间的每个子空间都是 Noether 空间。反过来，若一个拓扑空间是有限多个 Noether 子空间之并，则它是 Noether 空间。
\end{env}

\begin{env}[2.2.4]
\label{0.2.2.4-fr}
每个 Noether 空间都拟紧；反过来，若一个拓扑空间的每个开集都拟紧，则它是 Noether 空间。
\end{env}

\begin{env}[2.2.5]
\label{0.2.2.5-fr}
由 Noether 归纳可知，Noether 空间只有\emph{有限}个不可约分支。
\end{env}

\section{关于层的补充}
\label{section:0.3-fr}

\subsection{取值于一个范畴的层}
\label{subsection:0.3.1-fr}

\begin{env}[3.1.1]
\label{0.3.1.1-fr}
设 $\mathbf{K}$ 是一个范畴，$(A_\alpha)_{\alpha\in I}$、$(A_{\alpha\beta})_{(\alpha,\beta)\in I\times I}$ 是 $\mathbf{K}$ 中的两族对象，满足 $A_{\beta\alpha}=A_{\alpha\beta}$；又设 $(\rho_{\alpha\beta})_{(\alpha,\beta)\in I\times I}$ 是一族态射，其中 $\rho_{\alpha\beta}:A_\alpha\to A_{\alpha\beta}$。若由 $\mathbf{K}$ 的一个对象 $A$ 与一族态射 $\rho_\alpha:A\to A_\alpha$ 组成的二元组具有下列性质，就称它是由各族 $(A_\alpha)$、$(A_{\alpha\beta})$、$(\rho_{\alpha\beta})$ 所定义的\emph{泛问题的解}：对 $\mathbf{K}$ 的每个对象 $B$，把任一 $f\in\operatorname{Hom}(B,A)$ 映到族 $(\rho_\alpha\circ f)\in\prod_\alpha\operatorname{Hom}(B,A_\alpha)$ 的映射，是从 $\operatorname{Hom}(B,A)$ 到所有满足
$\rho_{\alpha\beta}\circ f_\alpha=\rho_{\beta\alpha}\circ f_\beta$（对任意指标对 $(\alpha,\beta)$）的族 $(f_\alpha)$ 所成集合的一个\emph{双射}。立即可见，这样的解若存在，就在相差一个同构的意义下唯一。
\end{env}

\begin{env}[3.1.2]
\label{0.3.1.2-fr}
这里不再回忆拓扑空间 $X$ 上取值于范畴 $\mathbf{K}$ 的\emph{预层} $U\to\mathcal{F}(U)$ 的定义 (G, I, 1.9)；若这样的预层满足下列公理，就称它为\emph{取值于 $\mathbf{K}$ 的层}：

\smallskip
\noindent
\textup{(F)} \emph{对 $X$ 的开集 $U$ 的任一开覆盖 $(U_\alpha)$，其中各 $U_\alpha$ 都包含于 $U$；若以 $\rho_\alpha$（相应地，以 $\rho_{\alpha\beta}$）表示限制态射}
\[
  \mathcal{F}(U)\to\mathcal{F}(U_\alpha)
  \qquad
  \bigl(\text{相应地，}\mathcal{F}(U_\alpha)
  \to\mathcal{F}(U_\alpha\cap U_\beta)\bigr),
\]

\oldpage[0\textsubscript{I}]{24}

\emph{则由 $\mathcal{F}(U)$ 与族 $(\rho_\alpha)$ 组成的二元组，是关于 $(\mathcal{F}(U_\alpha))$、$(\mathcal{F}(U_\alpha\cap U_\beta))$ 与 $(\rho_{\alpha\beta})$ 的泛问题的解 (3.1.1)}%
\footnote{这是\emph{逆极限}（不必滤过）这一一般概念的特例（\emph{见} (T, I, 1.8) 以及引言中提到的在编著作）。}。

等价地说，对 $\mathbf{K}$ 的每个对象 $T$，族 $U\to\operatorname{Hom}(T,\mathcal{F}(U))$ 是一个\emph{集合层}。
\end{env}

\begin{env}[3.1.3]
\label{0.3.1.3-fr}
假设 $\mathbf{K}$ 是由一类“带态射的结构” $\Sigma$ 所定义的范畴；于是，$\mathbf{K}$ 的对象是赋有 $\Sigma$ 类结构的集合，态射则是 $\Sigma$ 中的态射。再假设范畴 $\mathbf{K}$ 满足下列条件：

\smallskip
\noindent
\textup{(E)} 若 $(A,(\rho_\alpha))$ 在\emph{范畴 $\mathbf{K}$ 中}是各族 $(A_\alpha)$、$(A_{\alpha\beta})$、$(\rho_{\alpha\beta})$ 的一个泛映射问题的解，那么它在\emph{集合范畴中}也是这些族的泛映射问题的解（也就是说，把 $A$、各 $A_\alpha$ 与 $A_{\alpha\beta}$ 看作集合，把各 $\rho_\alpha$ 与 $\rho_{\alpha\beta}$ 看作映射）%
\footnote{可以证明，这也意味着典范函子 $\mathbf{K}\to(\mathrm{Ens})$ \emph{与逆极限可交换}（这些逆极限不必滤过）。}。

在这些条件下，(F) 蕴含：若把 $U\to\mathcal{F}(U)$ 看作\emph{集合预层}，它就是一个\emph{层}。此外，由 (F)，映射 $u:T\to\mathcal{F}(U)$ 是 $\mathbf{K}$ 中的态射，当且仅当每个映射 $\rho_\alpha\circ u$ 都是态射 $T\to\mathcal{F}(U_\alpha)$；这意味着 $\mathcal{F}(U)$ 上的 $\Sigma$ 类结构，是关于各态射 $\rho_\alpha$ 的\emph{初始结构}。反过来，假设 $X$ 上取值于 $\mathbf{K}$ 的预层 $U\to\mathcal{F}(U)$ 是一个\emph{集合层}，并满足前述条件；显然，它就满足 (F)，因而是一个\emph{取值于 $\mathbf{K}$ 的层}。
\end{env}

\begin{env}[3.1.4]
\label{0.3.1.4-fr}
当 $\Sigma$ 是群结构或环结构这一类结构时，取值于 $\mathbf{K}$ 的预层 $U\to\mathcal{F}(U)$ 若是一个\emph{集合层}，就\emph{自动}是一个\emph{取值于 $\mathbf{K}$ 的层}（换言之，是 (G) 意义下的群层或环层）%
\footnote{这是因为，在范畴 $\mathbf{K}$ 中，任何作为集合映射是\emph{双射}的态射都是\emph{同构}。例如，当 $\mathbf{K}$ 是拓扑空间范畴时，这一点便不再成立。}。但若 $\mathbf{K}$ 是例如\emph{拓扑环}范畴（其态射为连续映射），结论就不再如此：取值于 $\mathbf{K}$ 的层，是这样的环层 $U\to\mathcal{F}(U)$：对每个开集 $U$ 以及由开集 $U_\alpha\subset U$ 组成的 $U$ 的每个覆盖，环 $\mathcal{F}(U)$ 的拓扑都是使各映射 $\mathcal{F}(U)\to\mathcal{F}(U_\alpha)$ 连续的\emph{最粗}拓扑。在这种情形下，把 $U\to\mathcal{F}(U)$ 视为不带拓扑的环层时，称它为拓扑环层 $U\to\mathcal{F}(U)$ 的\emph{底层}环层。因此，拓扑环层的态射 $u_V:\mathcal{F}(V)\to\mathcal{G}(V)$（$V$ 为 $X$ 的任意开集）就是底层环层的同态，并且对每个开集 $V\subset X$，$u_V$ 都\emph{连续}；为了同底层环层的任意同态相区别，称它们为拓扑环层的\emph{连续}同态。对拓扑空间层或拓扑群层，也采用类似的定义与约定。
\end{env}

\oldpage[0\textsubscript{I}]{25}

\begin{env}[3.1.5]
\label{0.3.1.5-fr}
显然，对任意范畴 $\mathbf{K}$，若 $\mathcal{F}$ 是 $X$ 上取值于 $\mathbf{K}$ 的预层（相应地，层），而 $U$ 是 $X$ 的开集，则当各开集 $V\subset U$ 遍历时，各 $\mathcal{F}(V)$ 构成一个取值于 $\mathbf{K}$ 的预层（相应地，层）；称它为 $\mathcal{F}$ 在 $U$ 上的\emph{诱导}预层（相应地，诱导层），并记作 $\mathcal{F}|U$。

对 $X$ 上取值于 $\mathbf{K}$ 的预层的任一态射 $u:\mathcal{F}\to\mathcal{G}$，以 $u|U$ 表示由所有 $V\subset U$ 时的 $u_V$ 所组成的态射 $\mathcal{F}|U\to\mathcal{G}|U$。
\end{env}

\begin{env}[3.1.6]
\label{0.3.1.6-fr}
现在假设范畴 $\mathbf{K}$ 有\emph{归纳极限} (T, 1.8)。于是，对 $X$ 上任一取值于 $\mathbf{K}$ 的预层（特别地，对任一层）$\mathcal{F}$ 以及任一点 $x\in X$，可以把\emph{茎} $\mathcal{F}_x$ 定义为 $\mathbf{K}$ 中的如下归纳极限对象：当开邻域 $U$ 遍历 $x$ 在 $X$ 中的各开邻域所成的滤过集合（次序为 $\supset$）时，对各 $\mathcal{F}(U)$ 及各态射 $\rho_U^V:\mathcal{F}(V)\to\mathcal{F}(U)$ 取归纳极限。若 $u:\mathcal{F}\to\mathcal{G}$ 是取值于 $\mathbf{K}$ 的预层态射，则对每个 $x\in X$，把态射 $u_x:\mathcal{F}_x\to\mathcal{G}_x$ 定义为各 $u_U:\mathcal{F}(U)\to\mathcal{G}(U)$ 关于 $x$ 的各开邻域所成集合的归纳极限；这样，对每个 $x\in X$，$\mathcal{F}_x$ 就被定义为关于 $\mathcal{F}$ 的、取值于 $\mathbf{K}$ 的协变函子。

若 $\mathbf{K}$ 还由一类带态射的结构 $\Sigma$ 所定义，则仍把取值于 $\mathbf{K}$ 的\emph{层 $\mathcal{F}$} 的 $\mathcal{F}(U)$ 中各元素称为\emph{$U$ 上的截面}，并写作 $\Gamma(U,\mathcal{F})$ 而不写 $\mathcal{F}(U)$；对 $s\in\Gamma(U,\mathcal{F})$ 以及包含于 $U$ 的开集 $V$，写作 $s|V$ 而不写 $\rho_V^U(s)$；对每个 $x\in U$，$s$ 在 $\mathcal{F}_x$ 中的典范像称为 $s$ 在点 $x$ 处的\emph{芽}，记作 $s_x$（\emph{我们绝不在这个意义下使用记号 $s(x)$}；该记号将留给本论著所考虑的特殊层的另一概念 (5.5.1)）。

若 $u:\mathcal{F}\to\mathcal{G}$ 是取值于 $\mathbf{K}$ 的层态射，则对每个 $s\in\Gamma(U,\mathcal{F})$，写作 $u(s)$ 而不写 $u_V(s)$。

若 $\mathcal{F}$ 是交换群层、环层或模层，则把所有满足 $\mathcal{F}_x\neq\{0\}$ 的 $x\in X$ 所成的集合称为 $\mathcal{F}$ 的\emph{支撑}，记作 $\operatorname{Supp}(\mathcal{F})$；此集合在 $X$ 中未必闭。

当 $\mathbf{K}$ 由一类带态射的结构所定义时，对取值于 $\mathbf{K}$ 的层，\emph{我们将始终避免引入“espace étalé”空间的观点}；换言之，绝不把一个层看作拓扑空间（甚至不把它看作其各茎之并所成的集合），也不把 $X$ 上这类层的态射 $u:\mathcal{F}\to\mathcal{G}$ 看作拓扑空间之间的连续映射。
\end{env}

\subsection{拓扑基上的预层}
\label{subsection:0.3.2-fr}

\begin{env}[3.2.1]
\label{0.3.2.1-fr}
下文将只考虑有\emph{逆极限}的范畴 $\mathbf{K}$（广义的逆极限，即对应于不一定滤过的预序集，见 (T, 1.8)）。设 $X$ 是拓扑空间，$\mathfrak{B}$ 是 $X$ 的一个拓扑基。所谓\emph{$\mathfrak{B}$ 上取值于 $\mathbf{K}$ 的预层}，是指：对每个 $U\in\mathfrak{B}$ 配以一个对象 $\mathcal{F}(U)\in\mathbf{K}$，并且对 $\mathfrak{B}$ 中满足 $U\subset V$ 的每一对元素 $(U,V)$，给定一个态射
$\rho_U^V:\mathcal{F}(V)\to\mathcal{F}(U)$，

\oldpage[0\textsubscript{I}]{26}

使得 $\rho_U^U=\text{恒等态射}$，而当 $\mathfrak{B}$ 中的 $U$、$V$、$W$ 满足 $U\subset V\subset W$ 时，$\rho_U^W=\rho_U^V\circ\rho_V^W$。可以按通常意义给它配上一个\emph{取值于 $\mathbf{K}$ 的预层} $U\to\mathcal{F}'(U)$：对每个开集 $U$，令
\[
  \mathcal{F}'(U)=\varprojlim_{V\subset U}\mathcal{F}(V),
\]
其中 $V$ 遍历由满足 $V\subset U$ 的各 $V\in\mathfrak{B}$ 所成的、按 $\subset$ 排序的集合（它\emph{一般并不滤过}）；事实上，各 $\mathcal{F}(V)$ 关于 $\rho_V^W$ 构成一个逆系统（$V\subset W\subset U$，$V\in\mathfrak{B}$，$W\in\mathfrak{B}$）。若 $U$、$U'$ 是 $X$ 的两个开集且 $U\subset U'$，则把 $\rho'_U{}^{U'}$ 定义为各典范态射 $\mathcal{F}'(U')\to\mathcal{F}(V)$ 关于 $V\subset U$ 的逆极限；换言之，它是唯一的态射 $\mathcal{F}'(U')\to\mathcal{F}'(U)$，使得它与各典范态射 $\mathcal{F}'(U)\to\mathcal{F}(V)$ 的复合等于相应的典范态射 $\mathcal{F}'(U')\to\mathcal{F}(V)$。于是，$\rho'_U{}^{U'}$ 的传递性立即可得。此外，若 $U\in\mathfrak{B}$，则典范态射 $\mathcal{F}'(U)\to\mathcal{F}(U)$ 是同构，因而可以把这两个对象等同起来%
\footnote{若 $X$ 是 \emph{Noether 空间}，则只须假设 $\mathbf{K}$ 有\emph{有限}逆系统的逆极限，仍可定义 $\mathcal{F}'(U)$ 并证明它是通常意义下的预层。事实上，若 $U$ 是 $X$ 的任意开集，则存在一个由 $\mathfrak{B}$ 中集合组成的 $U$ 的\emph{有限}覆盖 $(V_i)$；对每一对指标 $(i,j)$，取一个由 $\mathfrak{B}$ 中集合组成的 $V_i\cap V_j$ 的有限覆盖 $(V_{ijk})$。令 $I$ 为各 $i$ 及各三元组 $(i,j,k)$ 所成的集合，仅以关系 $i>(i,j,k)$、$j>(i,j,k)$ 排序；然后把 $\mathcal{F}'(U)$ 取为由各 $\mathcal{F}(V_i)$ 与 $\mathcal{F}(V_{ijk})$ 组成的系统的逆极限。容易验证，这一对象与覆盖 $(V_i)$、$(V_{ijk})$ 的选择无关，并且 $U\to\mathcal{F}'(U)$ 是一个预层。}。
\end{env}

\begin{env}[3.2.2]
\label{0.3.2.2-fr}
上述定义的预层 $\mathcal{F}'$ 是一个\emph{层}，当且仅当 $\mathfrak{B}$ 上的预层 $\mathcal{F}$ 满足下列条件：

\smallskip
\noindent
\textup{(F\textsubscript{0})} \emph{对 $U\in\mathfrak{B}$ 的任一覆盖 $(U_\alpha)$，其中各 $U_\alpha\in\mathfrak{B}$ 且包含于 $U$，以及任一对象 $T\in\mathbf{K}$，把每个 $f\in\operatorname{Hom}(T,\mathcal{F}(U))$ 映到族
\[
  (\rho_{U_\alpha}^{U}\circ f)
  \in\prod_\alpha\operatorname{Hom}(T,\mathcal{F}(U_\alpha))
\]
所得的映射，是从 $\operatorname{Hom}(T,\mathcal{F}(U))$ 到所有满足下列条件的族 $(f_\alpha)$ 所成集合的双射：对任意一对指标 $(\alpha,\beta)$ 以及任意满足 $V\subset U_\alpha\cap U_\beta$ 的 $V\in\mathfrak{B}$，都有 $\rho_V^{U_\alpha}\circ f_\alpha=\rho_V^{U_\beta}\circ f_\beta$}%
\footnote{这也就是说，由 $\mathcal{F}(U)$ 与各 $\rho_\alpha=\rho_{U_\alpha}^{U}$ 组成的二元组，是由下列数据所定义的泛问题 (3.1.1) 的\emph{解}：$A_\alpha=\mathcal{F}(U_\alpha)$，$A_{\alpha\beta}=\prod\mathcal{F}(V)$（其中 $V\in\mathfrak{B}$ 且 $V\subset U_\alpha\cap U_\beta$），以及 $\rho_{\alpha\beta}=(\rho_V''):\mathcal{F}(U_\alpha)\to
\prod\mathcal{F}(V)$；后一个态射由下列条件规定：若 $V\in\mathfrak{B}$、$V'\in\mathfrak{B}$、$W\in\mathfrak{B}$，$V\cup V'\subset U_\alpha\cap U_\beta$ 且 $W\subset V\cap V'$，则 $\rho_W^V\circ\rho_V''=\rho_W^{V'}\circ\rho_{V'}''$。}。

条件的必要性显然。为证明充分性，先取 $X$ 的另一个\emph{包含于 $\mathfrak{B}$ 中的}拓扑基 $\mathfrak{B}'$；记 $\mathcal{F}''$ 为由子族 $(\mathcal{F}(V))_{V\in\mathfrak{B}'}$ 导出的预层，我们来证明 $\mathcal{F}''$ 与 $\mathcal{F}'$ \emph{典范同构}。首先，对每个开集 $U$，把各典范态射 $\mathcal{F}'(U)\to\mathcal{F}(V)$ 关于 $V\in\mathfrak{B}'$、$V\subset U$ 取逆极限，得到一个态射 $\mathcal{F}'(U)\to\mathcal{F}''(U)$。若 $U\in\mathfrak{B}$，此态射是同构；因为按照假设，对 $V\in\mathfrak{B}'$、$V\subset U$，各典范态射 $\mathcal{F}''(U)\to\mathcal{F}(V)$ 都分解为
\[
  \mathcal{F}''(U)\to\mathcal{F}(U)\to\mathcal{F}(V),
\]
而立即可见，这样定义的态射 $\mathcal{F}(U)\to\mathcal{F}''(U)$ 与 $\mathcal{F}''(U)\to\mathcal{F}(U)$ 的两个复合都是恒等态射。于是，对每个开集 $U$，当 $W\in\mathfrak{B}$ 且 $W\subset U$ 时，各态射 $\mathcal{F}''(U)\to\mathcal{F}''(W)=\mathcal{F}(W)$ 满足刻画各 $\mathcal{F}(W)$（$W\in\mathfrak{B}$，$W\subset U$）之逆极限的条件。结合逆极限在相差一个同构意义下的唯一性，这就证明了所断言的结论。

现在设 $U$ 是 $X$ 的任意开集，$(U_\alpha)$ 是由包含于 $U$ 的开集组成的 $U$ 的一个覆盖；令 $\mathfrak{B}'$ 为 $\mathfrak{B}$ 中各满足下述条件的集合所成的子族：

\oldpage[0\textsubscript{I}]{27}

这些 $\mathfrak{B}$ 中的集合至少包含于某个 $U_\alpha$。显然，$\mathfrak{B}'$ 仍是 $X$ 的一个拓扑基；所以 $\mathcal{F}'(U)$（相应地，$\mathcal{F}''(U_\alpha)$）是各 $\mathcal{F}(V)$ 关于 $V\in\mathfrak{B}'$、$V\subset U$（相应地，$V\subset U_\alpha$）的逆极限；由逆极限的定义，公理 (F) 随即得到验证。

当 (F\textsubscript{0}) 成立时，我们沿用一种语言上的滥用，把拓扑基 $\mathfrak{B}$ 上的预层 $\mathcal{F}$ 称为一个\emph{层}。
\end{env}

\begin{env}[3.2.3]
\label{0.3.2.3-fr}
设 $\mathcal{F}$、$\mathcal{G}$ 是拓扑基 $\mathfrak{B}$ 上取值于 $\mathbf{K}$ 的两个预层。所谓\emph{态射} $u:\mathcal{F}\to\mathcal{G}$，是指一族态射 $(u_V)_{V\in\mathfrak{B}}$，其中 $u_V:\mathcal{F}(V)\to\mathcal{G}(V)$，并且它们与各限制态射 $\rho_V^W$ 满足通常的相容条件。采用 (3.2.1) 的记号，对 $V\in\mathfrak{B}$、$V\subset U$ 时的各 $u_V$ 取逆极限，即得到相应的通常预层之间的一个态射 $u':\mathcal{F}'\to\mathcal{G}'$，其中 $u_U'$ 是该逆极限；它与 $\rho'_U{}^{U'}$ 的相容条件由逆极限的函子性立即得出。
\end{env}

\begin{env}[3.2.4]
\label{0.3.2.4-fr}
若范畴 $\mathbf{K}$ 有归纳极限，而 $\mathcal{F}$ 是拓扑基 $\mathfrak{B}$ 上取值于 $\mathbf{K}$ 的预层，则对每个 $x\in X$，$\mathfrak{B}$ 中属于 $x$ 的各邻域，在按 $\supset$ 排序的 $x$ 的全体邻域中组成一个共尾集。因此，若 $\mathcal{F}'$ 是与 $\mathcal{F}$ 对应的通常预层，则茎 $\mathcal{F}_x'$ 等于
\[
  \varinjlim_{\mathfrak{B}}\mathcal{F}(V),
\]
其指标遍历所有包含 $x$ 的 $V\in\mathfrak{B}$。若 $u:\mathcal{F}\to\mathcal{G}$ 是 $\mathfrak{B}$ 上取值于 $\mathbf{K}$ 的预层态射，而 $u':\mathcal{F}'\to\mathcal{G}'$ 是相应的通常预层态射，则 $u_x'$ 同样是各态射 $u_V:\mathcal{F}(V)\to\mathcal{G}(V)$ 关于 $V\in\mathfrak{B}$、$x\in V$ 的归纳极限。
\end{env}

\begin{env}[3.2.5]
\label{0.3.2.5-fr}
回到 (3.2.1) 的一般条件。若 $\mathcal{F}$ 是一个通常的取值于 $\mathbf{K}$ 的\emph{层}，而 $\mathcal{F}_1$ 是把 $\mathcal{F}$ 限制到 $\mathfrak{B}$ 后得到的\emph{$\mathfrak{B}$ 上的层}，则按 (3.2.1) 的方法从 $\mathcal{F}_1$ 得到的通常层 $\mathcal{F}_1'$，由条件 (F) 及逆极限的唯一性，与 $\mathcal{F}$ 典范同构。通常把 $\mathcal{F}$ 与 $\mathcal{F}_1'$ 等同起来。

若 $\mathcal{G}$ 是 $X$ 上另一个通常的取值于 $\mathbf{K}$ 的层，而 $u:\mathcal{F}\to\mathcal{G}$ 是态射，则上述说明表明，\emph{仅对 $V\in\mathfrak{B}$} 给出各 $u_V:\mathcal{F}(V)\to\mathcal{G}(V)$，便已完全决定 $u$。反过来，若对 $V\in\mathfrak{B}$ 已给定各 $u_V$，则只须对 $V\in\mathfrak{B}$、$W\in\mathfrak{B}$ 且 $V\subset W$ 验证它们与限制态射 $\rho_V^W$ 构成的图交换，便存在从 $\mathcal{F}$ 到 $\mathcal{G}$ 的唯一态射 $u'$，使对每个 $V\in\mathfrak{B}$ 都有 $u_V'=u_V$ (3.2.3)。
\end{env}

\begin{env}[3.2.6]
\label{0.3.2.6-fr}
仍假设 $\mathbf{K}$ 有逆极限。于是，\emph{$X$ 上取值于 $\mathbf{K}$ 的层所成的范畴}也有\emph{逆极限}。若 $(\mathcal{F}_\lambda)$ 是 $X$ 上取值于 $\mathbf{K}$ 的层所成的逆系统，则
\[
  \mathcal{F}(U)=\varprojlim_\lambda\mathcal{F}_\lambda(U)
\]
确实定义了一个取值于 $\mathbf{K}$ 的预层；公理 (F) 由逆极限的传递性得到验证，而 $\mathcal{F}$ 随即就是各 $\mathcal{F}_\lambda$ 的逆极限。

当 $\mathbf{K}$ 是集合范畴时，对任何满足下述条件的逆系统 $(\mathcal{H}_\lambda)$：

\oldpage[0\textsubscript{I}]{28}

每个 $\lambda$ 下，$\mathcal{H}_\lambda$ 都是 $\mathcal{F}_\lambda$ 的一个\emph{子层}，则 $\varprojlim_\lambda\mathcal{H}_\lambda$ 典范地等同于 $\varprojlim_\lambda\mathcal{F}_\lambda$ 的一个\emph{子层}。

若 $\mathbf{K}$ 是交换群范畴，则协变函子 $\varprojlim_\lambda\mathcal{F}_\lambda$ 是\emph{可加且左正合}的。
\end{env}

\subsection{层的黏合}
\label{subsection:0.3.3-fr}

\begin{env}[3.3.1]
\label{0.3.3.1-fr}
仍假设范畴 $\mathbf{K}$ 有（广义）逆极限。设 $X$ 是拓扑空间，$\mathfrak{U}=(U_\lambda)_{\lambda\in L}$ 是 $X$ 的一个开覆盖；对每个 $\lambda\in L$，设 $\mathcal{F}_\lambda$ 是 $U_\lambda$ 上取值于 $\mathbf{K}$ 的一个层。对每一对指标 $(\lambda,\mu)$，假设给定一个\emph{同构}
\[
  \theta_{\lambda\mu}:
  \mathcal{F}_\mu|(U_\lambda\cap U_\mu)
  \xrightarrow{\sim}
  \mathcal{F}_\lambda|(U_\lambda\cap U_\mu);
\]
此外，假设对每个三元组 $(\lambda,\mu,\nu)$，以 $\theta'_{\lambda\mu}$、$\theta'_{\mu\nu}$、$\theta'_{\lambda\nu}$ 分别表示 $\theta_{\lambda\mu}$、$\theta_{\mu\nu}$、$\theta_{\lambda\nu}$ 在 $U_\lambda\cap U_\mu\cap U_\nu$ 上的限制，则有
\[
  \theta'_{\lambda\nu}
  =\theta'_{\lambda\mu}\circ\theta'_{\mu\nu}
\]
（各 $\theta_{\lambda\mu}$ 的\emph{黏合条件}）。于是，存在 $X$ 上取值于 $\mathbf{K}$ 的一个层 $\mathcal{F}$，并且对每个 $\lambda$ 存在一个同构
\[
  \eta_\lambda:\mathcal{F}|U_\lambda\xrightarrow{\sim}\mathcal{F}_\lambda,
\]
使得对每一对 $(\lambda,\mu)$，若以 $\eta'_\lambda$、$\eta'_\mu$ 分别表示 $\eta_\lambda$、$\eta_\mu$ 在 $U_\lambda\cap U_\mu$ 上的限制，则
\[
  \theta_{\lambda\mu}=\eta'_\lambda\circ{\eta'_\mu}^{-1};
\]
此外，$\mathcal{F}$ 与各 $\eta_\lambda$ 在唯一同构意义下由这些条件确定。唯一性立即由 (3.2.5) 得出。为证明 $\mathcal{F}$ 的存在性，令 $\mathfrak{B}$ 为由至少包含于某个 $U_\lambda$ 的开集组成的拓扑基；对每个 $U\in\mathfrak{B}$，借助 Hilbert 的 $\tau$ 函数，在满足 $U\subset U_\lambda$ 的各 $\lambda$ 中选取一个，并取相应的一个 $\mathcal{F}_\lambda(U)$。若把这个对象记为 $\mathcal{F}(U)$，则当 $U\subset V$、$U\in\mathfrak{B}$、$V\in\mathfrak{B}$ 时，各 $\rho_U^V$ 可借助 $\theta_{\lambda\mu}$ 以显然的方式定义，而传递性条件由黏合条件推出；此外，(F\textsubscript{0}) 的验证是立即的。因此，这样定义的 $\mathfrak{B}$ 上的预层确为一个层，再由 (3.2.1) 的一般方法得到一个通常的层，仍记作 $\mathcal{F}$，它便是所求。称 $\mathcal{F}$ 是\emph{借助各 $\theta_{\lambda\mu}$ 将各 $\mathcal{F}_\lambda$ 黏合而得的}；通常借助 $\eta_\lambda$ 把 $\mathcal{F}_\lambda$ 与 $\mathcal{F}|U_\lambda$ 等同起来。

显然，$X$ 上任一取值于 $\mathbf{K}$ 的层 $\mathcal{F}$，都可视为借助恒等同构 $\theta_{\lambda\mu}$，由各层 $\mathcal{F}_\lambda=\mathcal{F}|U_\lambda$ 黏合而得（其中 $(U_\lambda)$ 是 $X$ 的任意开覆盖）。
\end{env}

\begin{env}[3.3.2]
\label{0.3.3.2-fr}
沿用上述记号，对每个 $\lambda\in L$，设 $\mathcal{G}_\lambda$ 是 $U_\lambda$ 上另一个取值于 $\mathbf{K}$ 的层；并且对每一对 $(\lambda,\mu)$，给定一个同构
\[
  \omega_{\lambda\mu}:
  \mathcal{G}_\mu|(U_\lambda\cap U_\mu)
  \xrightarrow{\sim}
  \mathcal{G}_\lambda|(U_\lambda\cap U_\mu),
\]
这些同构满足黏合条件。最后，假设对每个 $\lambda$ 给定一个态射 $u_\lambda:\mathcal{F}_\lambda\to\mathcal{G}_\lambda$，并且各图
\begin{equation}
\begin{tikzcd}[column sep=huge,row sep=large]
\mathcal{F}_\mu|(U_\lambda\cap U_\mu)
  \arrow[r,"u_\mu"] \arrow[d]
  & \mathcal{G}_\mu|(U_\lambda\cap U_\mu) \arrow[d] \\
\mathcal{F}_\lambda|(U_\lambda\cap U_\mu)
  \arrow[r,"u_\lambda"']
  & \mathcal{G}_\lambda|(U_\lambda\cap U_\mu)
\end{tikzcd}
\tag{3.3.2.1}\label{0.3.3.2.1-fr}
\end{equation}
交换。若 $\mathcal{G}$ 是借助各 $\omega_{\lambda\mu}$ 将各 $\mathcal{G}_\lambda$ 黏合而得的，则存在唯一的态射 $u:\mathcal{F}\to\mathcal{G}$，使得各图

\oldpage[0\textsubscript{I}]{29}

\[
\begin{tikzcd}[column sep=huge,row sep=large]
\mathcal{F}|U_\lambda
  \arrow[r,"u|U_\lambda"] \arrow[d]
  & \mathcal{G}|U_\lambda \arrow[d] \\
\mathcal{F}_\lambda
  \arrow[r,"u_\lambda"']
  & \mathcal{G}_\lambda
\end{tikzcd}
\]
交换；这立即由 (3.2.3) 得出。族 $(u_\lambda)$ 与 $u$ 之间的对应，是从 $\prod_\lambda\operatorname{Hom}(\mathcal{F}_\lambda,\mathcal{G}_\lambda)$ 中满足条件 (3.3.2.1) 的部分到 $\operatorname{Hom}(\mathcal{F},\mathcal{G})$ 的一个函子性双射。
\end{env}

\begin{env}[3.3.3]
\label{0.3.3.3-fr}
沿用 (3.3.1) 的记号，设 $V$ 是 $X$ 的开集。立即可见，各 $\theta_{\lambda\mu}$ 在 $V\cap U_\lambda\cap U_\mu$ 上的限制，对各诱导层 $\mathcal{F}_\lambda|(V\cap U_\lambda)$ 满足黏合条件；并且，将这些诱导层黏合而得的 $V$ 上的层，典范地等同于 $\mathcal{F}|V$。
\end{env}

\subsection{预层的直像}
\label{subsection:0.3.4-fr}

\begin{env}[3.4.1]
\label{0.3.4.1-fr}
设 $X$、$Y$ 是两个拓扑空间，$\psi:X\to Y$ 是连续映射。设 $\mathcal{F}$ 是 $X$ 上取值于范畴 $\mathbf{K}$ 的一个预层；对每个开集 $U\subset Y$，令
\[
  \mathcal{G}(U)=\mathcal{F}(\psi^{-1}(U)),
\]
并且，若 $U$、$V$ 是 $Y$ 的两个开子集且 $U\subset V$，令 $\rho_U^V$ 为态射
\[
  \mathcal{F}(\psi^{-1}(V))\to\mathcal{F}(\psi^{-1}(U)).
\]
显然，各 $\mathcal{G}(U)$ 与 $\rho_U^V$ 定义了 $Y$ 上取值于 $\mathbf{K}$ 的一个\emph{预层}，称为 $\mathcal{F}$ 在 $\psi$ 下的\emph{直像}，记作 $\psi_*(\mathcal{F})$。若 $\mathcal{F}$ 是一个层，则容易验证预层 $\mathcal{G}(U)$ 满足公理 (F)，因而 $\psi_*(\mathcal{F})$ 是一个层。
\end{env}

\begin{env}[3.4.2]
\label{0.3.4.2-fr}
设 $\mathcal{F}_1$、$\mathcal{F}_2$ 是 $X$ 上取值于 $\mathbf{K}$ 的两个预层，而 $u:\mathcal{F}_1\to\mathcal{F}_2$ 是一个态射。当 $U$ 遍历 $Y$ 的全体开子集时，态射族
\[
  u_{\psi^{-1}(U)}:
  \mathcal{F}_1(\psi^{-1}(U))\to\mathcal{F}_2(\psi^{-1}(U))
\]
满足与限制态射的相容条件，因而定义了一个态射
\[
  \psi_*(u):\psi_*(\mathcal{F}_1)\to\psi_*(\mathcal{F}_2).
\]
若 $v:\mathcal{F}_2\to\mathcal{F}_3$ 是从 $\mathcal{F}_2$ 到 $X$ 上第三个取值于 $\mathbf{K}$ 的预层的态射，则
\[
  \psi_*(v\circ u)=\psi_*(v)\circ\psi_*(u).
\]
换言之，$\psi_*(\mathcal{F})$ 关于 $\mathcal{F}$ 是一个\emph{协变函子}：它从 $X$ 上取值于 $\mathbf{K}$ 的预层（相应地，层）范畴，到 $Y$ 上取值于 $\mathbf{K}$ 的预层（相应地，层）范畴。
\end{env}

\begin{env}[3.4.3]
\label{0.3.4.3-fr}
设 $Z$ 是第三个拓扑空间，$\psi':Y\to Z$ 是连续映射，并令 $\psi''=\psi'\circ\psi$。显然，对 $X$ 上任一取值于 $\mathbf{K}$ 的预层 $\mathcal{F}$，都有
\[
  \psi_*''(\mathcal{F})=\psi_*'(\psi_*(\mathcal{F}));
\]
此外，对这类预层的任一态射 $u:\mathcal{F}\to\mathcal{G}$，都有
\[
  \psi_*''(u)=\psi_*'(\psi_*(u)).
\]
换言之，$\psi_*''$ 是函子 $\psi_*'$ 与 $\psi_*$ 的\emph{复合}，可以写成
\[
  (\psi'\circ\psi)_*=\psi_*'\circ\psi_*.
\]

此外，对 $Y$ 的任一开集 $U$，诱导预层 $\mathcal{F}|\psi^{-1}(U)$ 在限制映射 $\psi|\psi^{-1}(U)$ 下的直像，正是诱导预层 $\psi_*(\mathcal{F})|U$。
\end{env}

\begin{env}[3.4.4]
\label{0.3.4.4-fr}
假设范畴 $\mathbf{K}$ 有归纳极限，并设 $\mathcal{F}$ 是 $X$ 上取值于 $\mathbf{K}$ 的一个预层；对每个 $x\in X$，各态射
\[
  \Gamma(\psi^{-1}(U),\mathcal{F})\to\mathcal{F}_x
\]
（其中 $U$ 是 $Y$ 中 $\psi(x)$ 的开邻域）组成一个归纳系统，取其归纳极限得到一个茎之间的态射

\oldpage[0\textsubscript{I}]{30}

\[
  \psi_x:(\psi_*(\mathcal{F}))_{\psi(x)}\to\mathcal{F}_x.
\]
一般而言，这个态射\emph{既非单射也非满射}。它具有函子性；事实上，若 $u:\mathcal{F}_1\to\mathcal{F}_2$ 是 $X$ 上取值于 $\mathbf{K}$ 的预层态射，则图
\[
\begin{tikzcd}[column sep=huge,row sep=large]
(\psi_*(\mathcal{F}_1))_{\psi(x)}
  \arrow[r,"\psi_x"] \arrow[d,"{(\psi_*(u))_{\psi(x)}}"']
  & (\mathcal{F}_1)_x \arrow[d,"u_x"] \\
(\psi_*(\mathcal{F}_2))_{\psi(x)}
  \arrow[r,"\psi_x"']
  & (\mathcal{F}_2)_x
\end{tikzcd}
\]
交换。若 $Z$ 是第三个拓扑空间，$\psi':Y\to Z$ 是连续映射，而 $\psi''=\psi'\circ\psi$，则对 $x\in X$ 有
\[
  \psi_x''=\psi_x\circ\psi'_{\psi(x)}.
\]
\end{env}

\begin{env}[3.4.5]
\label{0.3.4.5-fr}
在 (3.4.4) 的假设下，再假设 $\psi$ 是从 $X$ 到 $Y$ 的子空间 $\psi(X)$ 上的一个\emph{同胚}。于是，对每个 $x\in X$，$\psi_x$ 都是\emph{同构}。特别地，这适用于 $Y$ 的一个子集 $X$ 到 $Y$ 中的典范嵌入 $j$。
\end{env}

\begin{env}[3.4.6]
\label{0.3.4.6-fr}
假设 $\mathbf{K}$ 是群范畴、环范畴等。若 $\mathcal{F}$ 是 $X$ 上取值于 $\mathbf{K}$、支撑为 $S$ 的一个层，而 $y\notin\overline{\psi(S)}$，则由 $\psi_*(\mathcal{F})$ 的定义可得
\[
  (\psi_*(\mathcal{F}))_y=\{0\}.
\]
换言之，$\psi_*(\mathcal{F})$ 的支撑包含于 $\overline{\psi(S)}$，但不一定包含于 $\psi(S)$。在同样的假设下，若 $j$ 是 $Y$ 的一个子集 $X$ 到 $Y$ 中的典范嵌入，则层 $j_*(\mathcal{F})$ 在 $X$ 上诱导出 $\mathcal{F}$；若 $X$ 还是 $Y$ 中的\emph{闭集}，则 $j_*(\mathcal{F})$ 是 $Y$ 上这样的层：它在 $X$ 上诱导出 $\mathcal{F}$，而在 $Y-X$ 上诱导出 $0$ (G, II, 2.9.2)；但当只假设 $X$ 局部闭而非闭时，它一般不同于后一层。
\end{env}

\subsection{预层的逆像}

\begin{env}[3.5.1]
\label{0.3.5.1-fr}
在 (3.4.1) 的假设下，若 $\mathcal{F}$（相应地，$\mathcal{G}$）是 $X$（相应地，$Y$）上取值于 $\mathbf{K}$ 的预层，则 $Y$ 上预层的任一态射
$u:\mathcal{G}\to\psi_*(\mathcal{F})$ 也称为从 $\mathcal{G}$ 到 $\mathcal{F}$ 的一个 $\psi$-\emph{态射}，并也记作 $\mathcal{G}\to\mathcal{F}$。从 $\mathcal{G}$ 到 $\mathcal{F}$ 的全体 $\psi$-态射所成的集合
$\operatorname{Hom}_Y(\mathcal{G},\psi_*(\mathcal{F}))$ 也记作
$\operatorname{Hom}_\psi(\mathcal{G},\mathcal{F})$。

对每一对 $(U,V)$，其中 $U$ 是 $X$ 的开集，$V$ 是 $Y$ 的开集，且
$\psi(U)\subset V$，将限制态射
$\mathcal{F}(\psi^{-1}(V))\to\mathcal{F}(U)$ 与态射
$u_V:\mathcal{G}(V)\to\psi_*(\mathcal{F})(V)=\mathcal{F}(\psi^{-1}(V))$
复合，便得到一个态射
$u_{U,V}:\mathcal{G}(V)\to\mathcal{F}(U)$；显然，对
$U'\subset U$、$V'\subset V$、$\psi(U')\subset V'$，这些态射使下图交换：
\begin{equation*}
\begin{tikzcd}[column sep=huge,row sep=large]
\mathcal{G}(V) \arrow[r,"u_{U,V}"] \arrow[d]
  & \mathcal{F}(U) \arrow[d] \\
\mathcal{G}(V') \arrow[r,"u_{U',V'}"']
  & \mathcal{F}(U')
\end{tikzcd}
\tag{3.5.1.1}\label{0.3.5.1.1-fr}
\end{equation*}
反过来，给定一个使图 (3.5.1.1) 交换的态射族 $(u_{U,V})$，便定义了一个
$\psi$-态射 $u$，因为只须取
\[
  u_V=u_{\psi^{-1}(V),V}.
\]

\oldpage[0\textsubscript{I}]{31}

若范畴 $\mathbf{K}$ 有（广义）逆极限，而 $\mathfrak{B}$、$\mathfrak{B}'$
分别是 $X$、$Y$ 的拓扑基，则为定义层的一个 $\psi$-态射 $u$，只须给出
$U\in\mathfrak{B}$、$V\in\mathfrak{B}'$、$\psi(U)\subset V$ 时的
$u_{U,V}$，并要求它们对 $\mathfrak{B}$ 中的 $U$、$U'$ 和
$\mathfrak{B}'$ 中的 $V$、$V'$ 满足相容条件 (3.5.1.1)。事实上，对任一开集
$W\subset Y$，只须把 $u_W$ 定义为所有 $u_{U,V}$ 的逆极限，其中
$V\in\mathfrak{B}'$、$V\subset W$，且 $U\in\mathfrak{B}$、
$\psi(U)\subset V$。

当范畴 $\mathbf{K}$ 有归纳极限时，对每个 $x\in X$ 以及 $Y$ 中
$\psi(x)$ 的每个开邻域 $V$，都有态射
\[
  \mathcal{G}(V)\to\mathcal{F}(\psi^{-1}(V))\to\mathcal{F}_x
\]
这些态射组成一个归纳系统；取其极限便得到态射
\[
  \mathcal{G}_{\psi(x)}\to\mathcal{F}_x.
\]
\end{env}

\begin{env}[3.5.2]
\label{0.3.5.2-fr}
在 (3.4.3) 的假设下，设 $\mathcal{F}$、$\mathcal{G}$、$\mathcal{H}$
分别是 $X$、$Y$、$Z$ 上取值于 $\mathbf{K}$ 的预层，并设
$u:\mathcal{G}\to\psi_*(\mathcal{F})$、
$v:\mathcal{H}\to\psi'_*(\mathcal{G})$ 分别是一个 $\psi$-态射和一个
$\psi'$-态射。由此得到一个 $\psi''$-态射
\[
  w:\mathcal{H}\xrightarrow{v}\psi'_*(\mathcal{G})
  \xrightarrow{\psi'_*(u)}
  \psi'_*(\psi_*(\mathcal{F}))=\psi''_*(\mathcal{F}),
\]
依定义称它为 $u$ 与 $v$ 的\emph{复合}。因此，可以把由拓扑空间 $X$ 与
$X$ 上的预层 $\mathcal{F}$（取值于 $\mathbf{K}$）组成的偶
$(X,\mathcal{F})$ 看成一个\emph{范畴}；其中的态射
$(\psi,\theta):(X,\mathcal{F})\to(Y,\mathcal{G})$ 由连续映射
$\psi:X\to Y$ 和 $\psi$-态射 $\theta:\mathcal{G}\to\mathcal{F}$ 组成。
\end{env}

\begin{env}[3.5.3]
\label{0.3.5.3-fr}
设 $\psi:X\to Y$ 是连续映射，$\mathcal{G}$ 是 $Y$ 上取值于
$\mathbf{K}$ 的一个\emph{预层}。所谓 $\mathcal{G}$ 在 $\psi$ 下的
\emph{逆像}，是指一个偶 $(\mathcal{G}',\rho)$，其中 $\mathcal{G}'$ 是
$X$ 上取值于 $\mathbf{K}$ 的一个\emph{层}，而
$\rho:\mathcal{G}\to\mathcal{G}'$ 是一个 $\psi$-态射（换言之，是一个同态
$\mathcal{G}\to\psi_*(\mathcal{G}'))$，并且对 $X$ 上取值于
$\mathbf{K}$ 的每个\emph{层} $\mathcal{F}$，映射
\begin{equation*}
  \operatorname{Hom}_X(\mathcal{G}',\mathcal{F})
  \longrightarrow \operatorname{Hom}_\psi(\mathcal{G},\mathcal{F})
  =\operatorname{Hom}_Y(\mathcal{G},\psi_*(\mathcal{F}))
\tag{3.5.3.1}\label{0.3.5.3.1-fr}
\end{equation*}
把 $v$ 映到 $\psi_*(v)\circ\rho$，且是一个\emph{双射}；由于这个映射关于
$\mathcal{F}$ 具有函子性，它于是定义了关于 $\mathcal{F}$ 的一个函子同构。
偶 $(\mathcal{G}',\rho)$ 是一个泛问题的解，因此，只要它存在，我们知道它在
\emph{唯一同构意义下确定}。此时写作
$\mathcal{G}'=\psi^*(\mathcal{G})$、$\rho=\rho_{\mathcal{G}}$；也约定把
$\psi^*(\mathcal{G})$ 称为 $\mathcal{G}$ 在 $\psi$ 下的\emph{逆像层}，这里总把
$\psi^*(\mathcal{G})$ 看成带有\emph{典范} $\psi$-态射
$\rho_{\mathcal{G}}:\mathcal{G}\to\psi^*(\mathcal{G})$，即带有 $Y$ 上预层的
\emph{典范同态}
\begin{equation*}
  \rho_{\mathcal{G}}:\mathcal{G}\to\psi_*(\psi^*(\mathcal{G})).
\tag{3.5.3.2}\label{0.3.5.3.2-fr}
\end{equation*}
对任一同态 $v:\psi^*(\mathcal{G})\to\mathcal{F}$（其中
$\mathcal{F}$ 是 $X$ 上取值于 $\mathbf{K}$ 的层），记
$v^b=\psi_*(v)\circ\rho_{\mathcal{G}}:
\mathcal{G}\to\psi_*(\mathcal{F})$。依定义，预层的\emph{每个}态射
$u:\mathcal{G}\to\psi_*(\mathcal{F})$ 都能写成 $v^b$；其中的 $v$ 唯一，
并记作 $u^\sharp$。换言之，预层的每个态射
$u:\mathcal{G}\to\psi_*(\mathcal{F})$ 都唯一地分解为
\begin{equation*}
  u:\mathcal{G}\xrightarrow{\rho_{\mathcal{G}}}
  \psi_*(\psi^*(\mathcal{G}))
  \xrightarrow{\psi_*(u^\sharp)}\psi_*(\mathcal{F}).
\tag{3.5.3.3}\label{0.3.5.3.3-fr}
\end{equation*}
\end{env}

\oldpage[0\textsubscript{I}]{32}

\begin{env}[3.5.4]
\label{0.3.5.4-fr}
现在假设范畴 $\mathbf{K}$ 具有如下性质\footnote{我们将在引言所引的著作中，给出范畴 $\mathbf{K}$ 所应满足的十分一般的条件，以保证取值于 $\mathbf{K}$ 的预层存在逆像。}：
$Y$ 上取值于 $\mathbf{K}$ 的\emph{每个}预层 $\mathcal{G}$ 都有在
$\psi$ 下的逆像，记作 $\psi^*(\mathcal{G})$。

我们将看到，可以把 $\psi^*(\mathcal{G})$ 定义成关于
$\mathcal{G}$ 的\emph{协变函子}：它从 $Y$ 上取值于 $\mathbf{K}$ 的预层范畴，到
$X$ 上取值于 $\mathbf{K}$ 的层范畴，并使同构 $v\to v^b$ 成为关于
$\mathcal{G}$ 和 $\mathcal{F}$ 的一个\emph{双函子同构}
\begin{equation*}
  \operatorname{Hom}_X(\psi^*(\mathcal{G}),\mathcal{F})
  \xrightarrow{\sim}
  \operatorname{Hom}_Y(\mathcal{G},\psi_*(\mathcal{F}))
\tag{3.5.4.1}\label{0.3.5.4.1-fr}
\end{equation*}

事实上，对 $Y$ 上取值于 $\mathbf{K}$ 的预层的任一态射
$w:\mathcal{G}_1\to\mathcal{G}_2$，考虑复合态射
\[
  \mathcal{G}_1\xrightarrow{w}\mathcal{G}_2
  \xrightarrow{\rho_{\mathcal{G}_2}}
  \psi_*(\psi^*(\mathcal{G}_2));
\]
与它对应的态射
$(\rho_{\mathcal{G}_2}\circ w)^\sharp:
\psi^*(\mathcal{G}_1)\to\psi^*(\mathcal{G}_2)$ 记作 $\psi^*(w)$。
因此，由 (3.5.3.3) 有
\begin{equation*}
  \psi_*(\psi^*(w))\circ\rho_{\mathcal{G}_1}
  =\rho_{\mathcal{G}_2}\circ w.
\tag{3.5.4.2}\label{0.3.5.4.2-fr}
\end{equation*}
对任一态射 $u:\mathcal{G}_2\to\psi_*(\mathcal{F})$，其中
$\mathcal{F}$ 是 $X$ 上取值于 $\mathbf{K}$ 的层，由 (3.5.3.3)、
(3.5.4.2) 以及 $v^b$ 的定义，有
\[
  (u^\sharp\circ\psi^*(w))^b
  =\psi_*(u^\sharp)\circ\psi_*(\psi^*(w))\circ\rho_{\mathcal{G}_1}
  =\psi_*(u^\sharp)\circ\rho_{\mathcal{G}_2}\circ w
  =u\circ w
\]
也就是说
\begin{equation*}
  (u\circ w)^\sharp=u^\sharp\circ\psi^*(w).
\tag{3.5.4.3}\label{0.3.5.4.3-fr}
\end{equation*}
特别地，若取 $u$ 为态射
\[
  \mathcal{G}_2\xrightarrow{w'}\mathcal{G}_3
  \xrightarrow{\rho_{\mathcal{G}_3}}
  \psi_*(\psi^*(\mathcal{G}_3)),
\]
便有
\[
  \psi^*(w'\circ w)
  =(\rho_{\mathcal{G}_3}\circ w'\circ w)^\sharp
  =(\rho_{\mathcal{G}_3}\circ w')^\sharp\circ\psi^*(w)
  =\psi^*(w')\circ\psi^*(w),
\]
从而得到上述断言。

最后，对 $X$ 上取值于 $\mathbf{K}$ 的任一层 $\mathcal{F}$，令
$i_{\mathcal{F}}$ 为 $\psi_*(\mathcal{F})$ 的恒等态射，并将态射
$(i_{\mathcal{F}})^\sharp$ 记作
\[
  \sigma_{\mathcal{F}}:
  \psi^*(\psi_*(\mathcal{F}))\to\mathcal{F}
\]
对任一态射 $u:\mathcal{G}\to\psi_*(\mathcal{F})$，公式 (3.5.4.3)
特别给出如下分解：
\begin{equation*}
  u^\sharp:\psi^*(\mathcal{G})
  \xrightarrow{\psi^*(u)}
  \psi^*(\psi_*(\mathcal{F}))
  \xrightarrow{\sigma_{\mathcal{F}}}\mathcal{F}
\tag{3.5.4.4}\label{0.3.5.4.4-fr}
\end{equation*}
称态射 $\sigma_{\mathcal{F}}$ 为\emph{典范}态射。
\end{env}

\begin{env}[3.5.5]
\label{0.3.5.5-fr}
设 $\psi':Y\to Z$ 是连续映射，并假设 $Z$ 上取值于 $\mathbf{K}$ 的每个预层
$\mathcal{H}$ 都有在 $\psi'$ 下的逆像 $\psi'^*(\mathcal{H})$。于是（在
(3.5.4) 的假设下）$Z$ 上取值于 $\mathbf{K}$ 的每个预层 $\mathcal{H}$
都有在 $\psi''=\psi'\circ\psi$ 下的逆像，并且有一个具有函子性的典范同构
\begin{equation*}
  \psi''^*(\mathcal{H})\xrightarrow{\sim}
  \psi^*(\psi'^*(\mathcal{H}))
\tag{3.5.5.1}\label{0.3.5.5.1-fr}
\end{equation*}

\oldpage[0\textsubscript{I}]{33}

事实上，注意到 $\psi''_*=\psi'_*\circ\psi_*$，这立即由定义得出。此外，若
$u:\mathcal{G}\to\psi_*(\mathcal{F})$ 是一个 $\psi$-态射，
$v:\mathcal{H}\to\psi'_*(\mathcal{G})$ 是一个 $\psi'$-态射，而
$w=\psi'_*(u)\circ v$ 是它们的复合 (3.5.2)，则立即可见 $w^\sharp$ 是复合态射
\[
  w^\sharp:\psi^*(\psi'^*(\mathcal{H}))
  \xrightarrow{\psi^*(v^\sharp)}
  \psi^*(\mathcal{G})
  \xrightarrow{u^\sharp}\mathcal{F}.
\]
\end{env}

\begin{env}[3.5.6]
\label{0.3.5.6-fr}
特别取 $\psi$ 为恒等映射 $1_X:X\to X$。若 $X$ 上取值于
$\mathbf{K}$ 的预层 $\mathcal{F}$ 在 $\psi$ 下的逆像存在，则称这一逆像为预层
$\mathcal{F}$ 的\emph{层化}。从 $\mathcal{F}$ 到取值于 $\mathbf{K}$ 的一个
\emph{层} $\mathcal{F}'$ 的任一态射 $u:\mathcal{F}\to\mathcal{F}'$，因而都唯一地分解为
\[
  \mathcal{F}\xrightarrow{\rho_{\mathcal{F}}}
  1_X^*(\mathcal{F})\xrightarrow{u^\sharp}\mathcal{F}'.
\]
\end{env}

\subsection{常值层和局部常值层}

\begin{env}[3.6.1]
\label{0.3.6.1-fr}
称 $X$ 上取值于 $\mathbf{K}$ 的一个\emph{预层} $\mathcal{F}$ 是
\emph{常值的}，如果对每个非空开集 $U\subset X$，典范态射
$\mathcal{F}(X)\to\mathcal{F}(U)$ 都是\emph{同构}；应注意，
$\mathcal{F}$ 不一定是层。称一个\emph{层}为\emph{常值层}，如果它由一个常值预层
层化而来 (3.5.6)。称一个层 $\mathcal{F}$ 为\emph{局部常值层}，如果每个
$x\in X$ 都有一个开邻域 $U$，使得 $\mathcal{F}|U$ 是常值层。
\end{env}

\begin{env}[3.6.2]
\label{0.3.6.2-fr}
假设 $X$ 是\emph{不可约的} (2.1.1)；于是，下列性质等价：
\begin{enumerate}
  \item[a)] \emph{$\mathcal{F}$ 是 $X$ 上的一个常值预层}；
  \item[b)] \emph{$\mathcal{F}$ 是 $X$ 上的一个常值层}；
  \item[c)] \emph{$\mathcal{F}$ 是 $X$ 上的一个局部常值层}。
\end{enumerate}

事实上，设 $\mathcal{F}$ 是 $X$ 上的一个常值预层；若 $U$、$V$ 是
$X$ 中两个非空开集，则 $U\cap V$ 非空。因此，由
$\mathcal{F}(X)\to\mathcal{F}(U)\to\mathcal{F}(U\cap V)$ 和
$\mathcal{F}(X)\to\mathcal{F}(U)$ 都是同构可知，
$\mathcal{F}(U)\to\mathcal{F}(U\cap V)$ 也是同构；同样，
$\mathcal{F}(V)\to\mathcal{F}(U\cap V)$ 也是同构。由此立即推出
(3.1.2) 的公理 (F) 确实成立，$\mathcal{F}$ 同构于它的层化；因而 a)
蕴含 b)。

现在设 $(U_\alpha)$ 是 $X$ 的一个由非空开集组成的开覆盖，而
$\mathcal{F}$ 是 $X$ 上的一个层，并且对每个 $\alpha$，
$\mathcal{F}|U_\alpha$ 都是常值层。由于 $U_\alpha$ 不可约，由前面的结论，
$\mathcal{F}|U_\alpha$ 是常值预层。又因
$U_\alpha\cap U_\beta$ 非空，态射
$\mathcal{F}(U_\alpha)\to\mathcal{F}(U_\alpha\cap U_\beta)$ 和
$\mathcal{F}(U_\beta)\to\mathcal{F}(U_\alpha\cap U_\beta)$ 都是同构，
故对每一对指标都有典范同构
$\theta_{\alpha\beta}:\mathcal{F}(U_\alpha)\to\mathcal{F}(U_\beta)$。
此时，对 $U=X$ 应用条件 (F)，便看出，对每个指标 $\alpha_0$，
$\mathcal{F}(U_{\alpha_0})$ 与各 $\theta_{\alpha_0\alpha}$ 都是这个泛问题的解；
由唯一性可知，$\mathcal{F}(X)\to\mathcal{F}(U_{\alpha_0})$ 是同构，
从而证明 c) 蕴含 a)。
\end{env}

\subsection{群预层和环预层的逆像}

\oldpage[0\textsubscript{I}]{34}

\begin{env}[3.7.1]
\label{0.3.7.1-fr}
我们将证明，当 $\mathbf{K}$ 取为集合范畴时，取值于 $\mathbf{K}$ 的每个预层
$\mathcal{G}$ 在 $\psi$ 下的逆像\emph{总是存在}（关于 $X$、$Y$、$\psi$
的记号和假设同 (3.5.3)）。事实上，对每个开集 $U\subset X$，如下定义
$\mathcal{G}'(U)$：$\mathcal{G}'(U)$ 的一个元素 $s'$ 是族
$(s'_x)_{x\in U}$，其中对每个 $x\in U$ 都有
$s'_x\in\mathcal{G}_{\psi(x)}$，并且对每个 $x\in U$ 满足下述条件：
存在 $Y$ 中 $\psi(x)$ 的一个开邻域 $V$、$x$ 的一个邻域
$W\subset\psi^{-1}(V)\cap U$ 以及一个元素 $s\in\mathcal{G}(V)$，使得对每个
$z\in W$ 都有 $s'_z=s_{\psi(z)}$。立即可以验证，
$U\mapsto\mathcal{G}'(U)$ 确实满足\emph{层}的公理。

现在设 $\mathcal{F}$ 是 $X$ 上的一个集合层，并设
$u:\mathcal{G}\to\psi_*(\mathcal{F})$、
$v:\mathcal{G}'\to\mathcal{F}$ 是态射。如下定义 $u^\sharp$ 与 $v^b$：
若 $s'$ 是 $\mathcal{G}'$ 在 $x\in X$ 的一个邻域 $U$ 上的截面，而
$V$ 是 $\psi(x)$ 的一个开邻域，并且 $s\in\mathcal{G}(V)$ 满足：
存在 $x$ 的一个包含于 $\psi^{-1}(V)\cap U$ 的邻域，使得对其中每个
$z$ 都有 $s'_z=s_{\psi(z)}$，则取
$u_x^\sharp(s'_x)=u_{\psi(x)}(s_{\psi(x)})$。同样，若
$s\in\mathcal{G}(V)$（$V$ 是 $Y$ 中的开集），则 $v^b(s)$ 是
$\mathcal{F}$ 在 $\psi^{-1}(V)$ 上的截面；它是 $\mathcal{G}'$ 的截面
$s'$ 在 $v$ 下的像，其中对每个 $x\in\psi^{-1}(V)$ 都有
$s'_x=s_{\psi(x)}$。此外，典范同态 (3.5.3)
$\rho:\mathcal{G}\to\psi_*(\psi^*(\mathcal{G}))$ 定义如下：对每个开集
$V\subset Y$ 和每个截面 $s\in\Gamma(V,\mathcal{G})$，$\rho(s)$ 是
$\psi^*(\mathcal{G})$ 在 $\psi^{-1}(V)$ 上的截面
$(s_{\psi(x)})_{x\in\psi^{-1}(V)}$。关系
$(u^\sharp)^b=u$、$(v^b)^\sharp=v$ 和
$v^b=\psi_*(v)\circ\rho$ 立即可得，从而证明了上述断言。

可以验证，若 $w:\mathcal{G}_1\to\mathcal{G}_2$ 是 $Y$ 上集合预层的同态，
则 $\psi^*(w)$ 可具体写成如下形式：若
$s'=(s'_x)_{x\in U}$ 是 $\psi^*(\mathcal{G}_1)$ 在 $X$ 的一个开集
$U$ 上的截面，则 $(\psi^*(w))(s')$ 是族
$(w_{\psi(x)}(s'_x))_{x\in U}$。最后，显然，对 $Y$ 的每个开集 $V$，
$\mathcal{G}|V$ 在 $\psi$ 到 $\psi^{-1}(V)$ 上的限制映射下的逆像，
就是诱导层 $\psi^*(\mathcal{G})|\psi^{-1}(V)$。

当 $\psi$ 是恒等映射 $1_X$ 时，便回到集合预层的层化定义
(G, II, 1.2)。当 $\mathbf{K}$ 是群范畴或环范畴（不必交换）时，
上述论证不加改变仍然适用。

设 $X$ 是拓扑空间 $Y$ 的任一子集，而 $j$ 是典范嵌入 $X\to Y$。
对 $Y$ 上取值于范畴 $\mathbf{K}$ 的任一层 $\mathcal{G}$，称其逆像
$j^*(\mathcal{G})$（如果存在）为 $\mathcal{G}$ 在 $X$ 上的\emph{诱导层}；
对集合层（或群层、环层），这就给出了通常的定义 (G, II, 1.5)。
\end{env}

\begin{env}[3.7.2]
\label{0.3.7.2-fr}
沿用 (3.5.3) 的记号与假设，并设 $\mathcal{G}$ 是 $Y$ 上的一个群层
（相应地，环层）。由 $\psi^*(\mathcal{G})$ 的截面的定义 (3.7.1)，
并考虑到 (3.4.4)，可知茎同态
$\psi_x\circ\rho_{\psi(x)}:\mathcal{G}_{\psi(x)}\to
(\psi^*(\mathcal{G}))_x$ 是一个\emph{关于 $\mathcal{G}$ 具有函子性的同构}，
由此可将这两个茎等同起来；在这一等同下，$u_x^\sharp$ 与 (3.5.1)
中定义的同态相同。特别地，有
$\operatorname{Supp}(\psi^*(\mathcal{G}))=
\psi^{-1}(\operatorname{Supp}(\mathcal{G}))$。

这一结果的一个直接推论是：在交换群层所成的阿贝尔范畴中，
\emph{函子 $\psi^*(\mathcal{G})$ 关于 $\mathcal{G}$ 是正合的}。
\end{env}

\subsection{伪离散空间层}

\oldpage[0\textsubscript{I}]{35}

\begin{env}[3.8.1]
\label{0.3.8.1-fr}
设 $X$ 是一个拓扑空间，它的拓扑有一个由\emph{拟紧}开集组成的基
$\mathfrak{B}$。设 $\mathcal{F}$ 是 $X$ 上的一个\emph{集合层}；给每个
$\mathcal{F}(U)$ 赋予\emph{离散}拓扑以后，
$U\mapsto\mathcal{F}(U)$ 就是一个\emph{拓扑空间预层}。我们将看到，
拓扑空间预层 $\mathcal{F}$ 的层化 $\mathcal{F}'$ 存在 (3.5.6)，
使得对每个\emph{拟紧}开集 $U$，$\Gamma(U,\mathcal{F}')$ 都是离散空间
$\mathcal{F}(U)$。为此，只须证明拓扑基 $\mathfrak{B}$ \emph{上的}
离散拓扑空间预层 $U\mapsto\mathcal{F}(U)$ 满足 (3.2.2) 的条件
$(\mathrm{F}_0)$；更一般地，只须证明：若 $U$ 是拟紧开集，而
$(U_\alpha)$ 是由 $\mathfrak{B}$ 中的集合构成的 $U$ 的一个覆盖，则使各映射
$\Gamma(U,\mathcal{F})\to\Gamma(U_\alpha,\mathcal{F})$ 连续的
$\Gamma(U,\mathcal{F})$ 上的最粗拓扑 $\mathcal{T}$，就是\emph{离散}拓扑。
事实上，存在有限多个指标 $\alpha_i$，使得
$U=\bigcup_i U_{\alpha_i}$。设 $s\in\Gamma(U,\mathcal{F})$，并令 $s_i$
为它在 $\Gamma(U_{\alpha_i},\mathcal{F})$ 中的像；依定义，各单点集
$\{s_i\}$ 的逆像之交是 $s$ 在 $\mathcal{T}$ 下的一个邻域。但是，由于
$\mathcal{F}$ 是集合层，而各 $U_{\alpha_i}$ 覆盖 $U$，这个交恰由
$s$ 这一点组成，从而得到上述断言。

应注意，若 $U$ 是 $X$ 的一个非拟紧开集，则拓扑空间
$\Gamma(U,\mathcal{F}')$ 的底集仍是 $\Gamma(U,\mathcal{F})$，但其拓扑一般
不是离散的：它是使所有映射
$\Gamma(U,\mathcal{F})\to\Gamma(V,\mathcal{F})$ 连续的最粗拓扑，其中
$V\in\mathfrak{B}$ 且 $V\subset U$（各 $\Gamma(V,\mathcal{F})$ 都取离散拓扑）。

上述论证不加改变也适用于群层或环层（不必交换），并分别给出
\emph{拓扑群}层或\emph{拓扑环}层。为简便起见，称层 $\mathcal{F}'$ 为与
集合层（相应地，群层、环层）$\mathcal{F}$ 对应的\emph{伪离散空间}
（相应地，\emph{伪离散群}、\emph{伪离散环}）层。
\end{env}

\begin{env}[3.8.2]
\label{0.3.8.2-fr}
设 $\mathcal{F}$、$\mathcal{G}$ 是 $X$ 上的两个集合层（相应地，群层、
环层），而 $u:\mathcal{F}\to\mathcal{G}$ 是一个同态。以
$\mathcal{F}'$、$\mathcal{G}'$ 分别表示与 $\mathcal{F}$、$\mathcal{G}$
对应的伪离散层，则 $u$ 也是一个\emph{连续}同态
$\mathcal{F}'\to\mathcal{G}'$；这由 (3.2.5) 得出。
\end{env}

\begin{env}[3.8.3]
\label{0.3.8.3-fr}
设 $\mathcal{F}$ 是集合层，$\mathcal{H}$ 是 $\mathcal{F}$ 的一个子层，而
$\mathcal{F}'$、$\mathcal{H}'$ 分别是与 $\mathcal{F}$、$\mathcal{H}$
对应的伪离散层。于是，对每个开集 $U\subset X$，
$\Gamma(U,\mathcal{H}')$ 在 $\Gamma(U,\mathcal{F}')$ 中是\emph{闭的}：
事实上，它是所有 $\Gamma(V,\mathcal{H})$（其中 $V\in\mathfrak{B}$、
$V\subset U$）在连续映射
$\Gamma(U,\mathcal{F})\to\Gamma(V,\mathcal{F})$ 下的逆像之交，而
$\Gamma(V,\mathcal{H})$ 在离散空间 $\Gamma(V,\mathcal{F})$ 中是闭的。
\end{env}

\section{赋环空间}
\label{section:0.4-fr}

\subsection[赋环空间、A-模层、A-代数层]{赋环空间、$\mathcal{A}$-模层、$\mathcal{A}$-代数层}
\label{subsection:0.4.1-fr}

\begin{env}[4.1.1]
\label{0.4.1.1-fr}
一个\emph{赋环空间}（相应地，\emph{拓扑赋环空间}）是一个偶
$(X,\mathcal{A})$，其中 $X$ 是拓扑空间，$\mathcal{A}$ 是 $X$ 上的一个
环层（不必交换）（相应地，拓扑环层）；称 $X$ 为赋环空间
$(X,\mathcal{A})$ 的
\oldpage[0\textsubscript{I}]{36}
\emph{底层}拓扑空间，称 $\mathcal{A}$ 为它的\emph{结构层}。后者记作
$\mathcal{O}_X$，它在点 $x\in X$ 处的茎记作 $\mathcal{O}_{X,x}$；
在不致混淆时，也简记作 $\mathcal{O}_x$。

用 $1$ 或 $e$ 表示 $\mathcal{O}_X$ 在 $X$ 上的\emph{单位截面}
（即 $\Gamma(X,\mathcal{O}_X)$ 的单位元）。

本书所考察的主要是\emph{交换}环层，因此，在不另加说明地谈到赋环空间
$(X,\mathcal{A})$ 时，总约定 $\mathcal{A}$ 是一个交换环层。

若把一个\emph{态射} $(X,\mathcal{A})\to(Y,\mathcal{B})$ 定义为一个偶
$(\psi,\theta)=\Psi$，其中 $\psi:X\to Y$ 是连续映射，而这里的环层
（相应地，拓扑环层）的一个 $\psi$-\emph{态射}
% EGA1-ERR-0002: French print literally has G->F; EGA II errata corrects this to B->A. Canonical translation preserves the printed formula.
$\theta:\mathcal{G}\to\mathcal{F}$ (3.5.1)，则结构层不必交换的赋环空间
（相应地，拓扑赋环空间）构成一个\emph{范畴}。设另一个态射为
$\Psi'=(\psi',\theta'):(Y,\mathcal{B})\to(Z,\mathcal{C})$，则它与
$\Psi$ 的\emph{复合}记作 $\Psi''=\Psi'\circ\Psi$，并定义为态射
$(\psi'',\theta'')$，其中 $\psi''=\psi'\circ\psi$，而 $\theta''$ 是
$\theta$ 与 $\theta'$ 的复合（即
$\psi'_*(\theta)\circ\theta'$，参见 3.5.2）。对赋环空间，回忆此时有
$\theta''^\sharp=\theta^\sharp\circ\psi^*(\theta'^\sharp)$ (3.5.5)；
因此，若 $\theta'^\sharp$ 和 $\theta^\sharp$ 都是\emph{单射}同态
（相应地，\emph{满射}同态），则考虑到对每个 $x\in X$，
$\psi_x\circ\rho_{\psi(x)}$ 都是同构 (3.7.2)，可知
$\theta''^\sharp$ 也具有同样的性质。由此立即验证：当 $\psi$ 是\emph{单射}
连续映射而 $\theta^\sharp$ 是环层的\emph{满射}同态时，$(\psi,\theta)$
是赋环空间范畴中的一个\emph{单态射} (T, 1.1)。

按照一种术语上的简便约定，常在记号中以 $\Psi$ 代替 $\psi$；例如，对
$Y$ 的一个子集 $U$，在不会引起混淆时写 $\Psi^{-1}(U)$ 而不写
$\psi^{-1}(U)$。
\end{env}

\begin{env}[4.1.2]
\label{0.4.1.2-fr}
对 $X$ 的任一子集 $M$，偶 $(M,\mathcal{A}|M)$ 显然是一个赋环空间，称为
由赋环空间 $(X,\mathcal{A})$ 在 $M$ 上\emph{诱导}的赋环空间（也称
$(X,\mathcal{A})$ 在 $M$ 上的\emph{限制}）。若 $j$ 是典范嵌入
$M\to X$，而 $\omega$ 是 $\mathcal{A}|M$ 的恒等映射，则
$(j,\omega^b)$ 是赋环空间的一个单态射
$(M,\mathcal{A}|M)\to(X,\mathcal{A})$，称为\emph{典范嵌入}。态射
$\Psi:(X,\mathcal{A})\to(Y,\mathcal{B})$ 与这个嵌入的复合，称为
$\Psi$ 在 $M$ 上的\emph{限制}。
\end{env}

\begin{env}[4.1.3]
\label{0.4.1.3-fr}
我们不再重复赋环空间 $(X,\mathcal{A})$ 上的 $\mathcal{A}$-模层或代数层
的定义 (G, II, 2.2)。当 $\mathcal{A}$ 是不必交换的环层时，除非明确
另有说明，$\mathcal{A}$-模层总是指“左 $\mathcal{A}$-模层”。
$\mathcal{A}$ 自身的子-$\mathcal{A}$-模层称为 $\mathcal{A}$ 中的理想层
（左理想层、右理想层或双边理想层），也称为 $\mathcal{A}$-理想层。

当 $\mathcal{A}$ 是交换环层时，把 $\mathcal{A}$-模层定义中各处的模结构
换成代数结构，就得到 $X$ 上 $\mathcal{A}$-代数层的定义。等价地，一个
（不必交换的）$\mathcal{A}$-代数层是一个 $\mathcal{A}$-模层
$\mathcal{C}$，配有一个 $\mathcal{A}$-模层同态
$\varphi:\mathcal{C}\otimes_{\mathcal{A}}\mathcal{C}\to\mathcal{C}$
以及 $X$ 上的一个截面 $e$，并满足：1\textsuperscript{o} 图
\oldpage[0\textsubscript{I}]{37}
\[
\begin{tikzcd}[column sep=large,row sep=large]
\mathcal{C}\otimes_{\mathcal{A}}\mathcal{C}\otimes_{\mathcal{A}}\mathcal{C}
  \arrow[r,"\varphi\otimes 1"]
  \arrow[d,"1\otimes\varphi"']
&
\mathcal{C}\otimes_{\mathcal{A}}\mathcal{C}
  \arrow[d,"\varphi"]
\\
\mathcal{C}\otimes_{\mathcal{A}}\mathcal{C}
  \arrow[r,"\varphi"']
&
\mathcal{C}
\end{tikzcd}
\]
交换；2\textsuperscript{o} 对每个开集 $U\subset X$ 和每个截面
$s\in\Gamma(U,\mathcal{C})$，都有
$\varphi((e|U)\otimes s)=\varphi(s\otimes(e|U))=s$。再要求图
\[
\begin{tikzcd}[column sep=large,row sep=large]
\mathcal{C}\otimes_{\mathcal{A}}\mathcal{C}
  \arrow[rr,"\sigma"]
  \arrow[dr,"\varphi"']
&&
\mathcal{C}\otimes_{\mathcal{A}}\mathcal{C}
  \arrow[dl,"\varphi"]
\\
& \mathcal{C} &
\end{tikzcd}
\]
交换，就等价于说 $\mathcal{C}$ 是交换 $\mathcal{A}$-代数层；这里
$\sigma$ 表示张量积 $\mathcal{C}\otimes_{\mathcal{A}}\mathcal{C}$ 的
典范对称同构。

$\mathcal{A}$-代数层同态也如 (G, II, 2.2) 中的
$\mathcal{A}$-模层同态那样定义，不过它们自然不再构成阿贝尔群。

若 $\mathcal{M}$ 是 $\mathcal{A}$-代数层 $\mathcal{C}$ 的一个
子-$\mathcal{A}$-模层，则由 $\mathcal{M}$ 生成的 $\mathcal{C}$ 的
子-$\mathcal{A}$-代数层，是各同态
$\bigotimes^n\mathcal{M}\to\mathcal{C}$（$n\geqslant0$）的像之和。
它也等于代数预层 $U\mapsto\mathcal{B}(U)$ 的层化，其中
$\mathcal{B}(U)$ 是 $\Gamma(U,\mathcal{C})$ 中由子模
$\Gamma(U,\mathcal{M})$ 生成的子代数。
\end{env}

\begin{env}[4.1.4]
\label{0.4.1.4-fr}
设 $\mathcal{A}$ 是拓扑空间 $X$ 上的一个环层。若茎
$\mathcal{A}_x$ 是\emph{约化}环 (1.1.1)，则称 $\mathcal{A}$ 在
$X$ 的点 $x$ 处\emph{约化}；若它在 $X$ 的每一点都约化，则称
$\mathcal{A}$ \emph{约化}。回忆，一个环 $A$ 称为\emph{正则}的，是指
它的每个局部环 $A_{\mathfrak p}$（其中 $\mathfrak p$ 遍历 $A$ 的所有
素理想）都是正则局部环。若茎 $\mathcal{A}_x$ 是正则环（相应地，
$X$ 的每一点处的茎都正则），则称 $\mathcal{A}$ 在点 $x$ 处
\emph{正则}（相应地，称 $\mathcal{A}$ \emph{正则}）。最后，若茎
$\mathcal{A}_x$ 是整环且整闭（相应地，每个茎都如此），
则称 $\mathcal{A}$ 在点 $x$ 处\emph{正规}（相应地，称
其\emph{正规}）。若赋环空间 $(X,\mathcal{A})$ 的环层具有上述某一性质，
就说这个赋环空间具有该性质。

一个\emph{分次环层} $\mathcal{A}$，按定义是阿贝尔群层族
$(\mathcal{A}_n)_{n\in\mathbf Z}$ 的直和 (G, II, 2.7)，并且满足
$\mathcal{A}_m\mathcal{A}_n\subset\mathcal{A}_{m+n}$。一个
$\mathcal{A}$-\emph{分次模层}，是形如阿贝尔群层族
$(\mathcal{F}_n)_{n\in\mathbf Z}$ 之直和的 $\mathcal{A}$-模层
$\mathcal{F}$，并满足
$\mathcal{A}_m\mathcal{F}_n\subset\mathcal{F}_{m+n}$。事实上，这分别
等价于在每个点 $x$ 都有
$(\mathcal{A}_m)_x(\mathcal{A}_n)_x\subset(\mathcal{A}_{m+n})_x$
和
$(\mathcal{A}_m)_x(\mathcal{F}_n)_x\subset(\mathcal{F}_{m+n})_x$。
\end{env}

\begin{env}[4.1.5]
\label{0.4.1.5-fr}
给定一个赋环空间 $(X,\mathcal{A})$（不必交换），这里不再重述左或右
$\mathcal{A}$-模层范畴（视情形而定）上的双函子
$\mathcal{F}\otimes_{\mathcal{A}}\mathcal{G}$、
$\mathcal{H}\!om_{\mathcal{A}}(\mathcal{F},\mathcal{G})$ 和
$\operatorname{Hom}_{\mathcal{A}}(\mathcal{F},\mathcal{G})$
(G, II, 2.8 和 2.2) 的定义；它们取值于阿贝尔群层范畴（或者更一般地，

\oldpage[0\textsubscript{I}]{38}
取值于 $\mathcal{C}$-模层范畴，其中 $\mathcal{C}$ 是 $\mathcal{A}$ 的
中心）。在每个点 $x\in X$，茎
$(\mathcal{F}\otimes_{\mathcal{A}}\mathcal{G})_x$ 典范地等同于
$\mathcal{F}_x\otimes_{\mathcal{A}_x}\mathcal{G}_x$，并有一个具有
函子性的典范同态
\[
  (\mathcal{H}\!om_{\mathcal{A}}(\mathcal{F},\mathcal{G}))_x
  \longrightarrow
  \operatorname{Hom}_{\mathcal{A}_x}(\mathcal{F}_x,\mathcal{G}_x)
\]
它一般既不是单射也不是满射。上述双函子是可加的，特别地，与有限直和
交换；$\mathcal{F}\otimes_{\mathcal{A}}\mathcal{G}$ 关于
$\mathcal{F}$ 和 $\mathcal{G}$ 都是右正合的，与归纳极限交换，并且
$\mathcal{A}\otimes_{\mathcal{A}}\mathcal{G}$（相应地，
$\mathcal{F}\otimes_{\mathcal{A}}\mathcal{A}$）典范地等同于
$\mathcal{G}$（相应地，$\mathcal{F}$）。函子
$\mathcal{H}\!om_{\mathcal{A}}(\mathcal{F},\mathcal{G})$ 和
$\operatorname{Hom}_{\mathcal{A}}(\mathcal{F},\mathcal{G})$ 关于
$\mathcal{F}$ 和 $\mathcal{G}$ 都是\emph{左正合}的；确切地说，若有
形如 $0\to\mathcal{G}'\to\mathcal{G}\to\mathcal{G}''$ 的正合列，则
序列
\[
  0\longrightarrow
  \mathcal{H}\!om_{\mathcal{A}}(\mathcal{F},\mathcal{G}')
  \longrightarrow
  \mathcal{H}\!om_{\mathcal{A}}(\mathcal{F},\mathcal{G})
  \longrightarrow
  \mathcal{H}\!om_{\mathcal{A}}(\mathcal{F},\mathcal{G}'')
\]
正合；若有形如
$\mathcal{F}'\to\mathcal{F}\to\mathcal{F}''\to0$ 的正合列，则序列
\[
  0\longrightarrow
  \mathcal{H}\!om_{\mathcal{A}}(\mathcal{F}'',\mathcal{G})
  \longrightarrow
  \mathcal{H}\!om_{\mathcal{A}}(\mathcal{F},\mathcal{G})
  \longrightarrow
  \mathcal{H}\!om_{\mathcal{A}}(\mathcal{F}',\mathcal{G})
\]
正合；函子 $\operatorname{Hom}$ 也有类似的性质。此外，
$\mathcal{H}\!om_{\mathcal{A}}(\mathcal{A},\mathcal{G})$ 典范地等同于
$\mathcal{G}$；最后，对每个开集 $U\subset X$，有
\[
  \Gamma\bigl(U,\mathcal{H}\!om_{\mathcal{A}}(\mathcal{F},\mathcal{G})\bigr)
  =
  \operatorname{Hom}_{\mathcal{A}|U}(\mathcal{F}|U,\mathcal{G}|U).
\]

对任一左 $\mathcal{A}$-模层 $\mathcal{F}$（相应地，右
$\mathcal{A}$-模层），称右 $\mathcal{A}$-模层（相应地，左
$\mathcal{A}$-模层）
$\mathcal{H}\!om_{\mathcal{A}}(\mathcal{F},\mathcal{A})$ 为
$\mathcal{F}$ 的\emph{对偶}，并记作 $\check{\mathcal{F}}$。

最后，若 $\mathcal{A}$ 是交换环层，$\mathcal{F}$ 是
$\mathcal{A}$-模层，则
$U\mapsto\bigwedge^p\Gamma(U,\mathcal{F})$ 是一个预层，它的层化是
一个 $\mathcal{A}$-模层，记作 $\bigwedge^p\mathcal{F}$，称为
$\mathcal{F}$ 的\emph{第 $p$ 外幂}。容易验证，从预层
$U\mapsto\bigwedge^p\Gamma(U,\mathcal{F})$ 到它的层化
$\bigwedge^p\mathcal{F}$ 的典范映射是\emph{单射}，而且对每个
$x\in X$ 都有
$(\bigwedge^p\mathcal{F})_x=\bigwedge^p(\mathcal{F}_x)$。显然，
$\bigwedge^p\mathcal{F}$ 关于 $\mathcal{F}$ 是协变函子。
\end{env}

\begin{env}[4.1.6]
\label{0.4.1.6-fr}
设 $\mathcal{A}$ 是不必交换的环层，$\mathcal{J}$ 是
$\mathcal{A}$ 的一个左理想层，$\mathcal{F}$ 是左 $\mathcal{A}$-模层。
记 $\mathcal{J}\mathcal{F}$ 为 $\mathcal{F}$ 的如下
子-$\mathcal{A}$-模层：它是 $\mathcal{J}\otimes_{\mathbf Z}\mathcal{F}$
在典范映射
$\mathcal{J}\otimes_{\mathbf Z}\mathcal{F}\to\mathcal{F}$ 下的像，
其中 $\mathbf Z$ 是常值预层 $U\to\mathbf Z$ 的层化。显然，对每个
$x\in X$，都有
$(\mathcal{J}\mathcal{F})_x=\mathcal{J}_x\mathcal{F}_x$。当
$\mathcal{A}$ 交换时，$\mathcal{J}\mathcal{F}$ 也就是
$\mathcal{J}\otimes_{\mathcal{A}}\mathcal{F}\to\mathcal{F}$ 的典范像。
立即可见，$\mathcal{J}\mathcal{F}$ 也等于预层
$U\to\Gamma(U,\mathcal{J})\Gamma(U,\mathcal{F})$ 的层化。若
$\mathcal{J}_1$、$\mathcal{J}_2$ 是 $\mathcal{A}$ 的两个左理想层，则
有
$\mathcal{J}_1(\mathcal{J}_2\mathcal{F})=(\mathcal{J}_1\mathcal{J}_2)\mathcal{F}$。
\end{env}

\begin{env}[4.1.7]
\label{0.4.1.7-fr}
设 $(X_\lambda,\mathcal{A}_\lambda)_{\lambda\in L}$ 是一族赋环空间。
对每个偶 $(\lambda,\mu)$，设给定 $X_\lambda$ 的一个开子集
$V_{\lambda\mu}$ 以及赋环空间的一个同构
\[
  \varphi_{\lambda\mu}:
  (V_{\mu\lambda},\mathcal{A}_\mu|V_{\mu\lambda})
  \xrightarrow{\sim}
  (V_{\lambda\mu},\mathcal{A}_\lambda|V_{\lambda\mu}),
\]
并且 $V_{\lambda\lambda}=X_\lambda$、$\varphi_{\lambda\lambda}$ 是恒等映射。
再假设，对每个三元组 $(\lambda,\mu,\nu)$，若以
$\varphi'_{\mu\lambda}$ 表示 $\varphi_{\mu\lambda}$ 在
$V_{\lambda\mu}\cap V_{\lambda\nu}$ 上的限制，则
$\varphi'_{\mu\lambda}$ 是从
\[
  (V_{\lambda\mu}\cap V_{\lambda\nu},
   \mathcal{A}_\lambda|(V_{\lambda\mu}\cap V_{\lambda\nu}))
\]
到
\[
  (V_{\mu\nu}\cap V_{\mu\lambda},
   \mathcal{A}_\mu|(V_{\mu\nu}\cap V_{\mu\lambda}))
\]
的一个同构，并且有
$\varphi'_{\lambda\nu}=\varphi'_{\lambda\mu}\circ\varphi'_{\mu\nu}$
（各 $\varphi_{\lambda\mu}$ 的\emph{黏合条件}）。首先，可以考虑通过
各 $\varphi_{\lambda\mu}$ 沿各 $V_{\lambda\mu}$ 将各 $X_\lambda$
黏合而得的拓扑空间 $X$。

\oldpage[0\textsubscript{I}]{39}
若将 $X_\lambda$ 与 $X$ 中相应的开子集 $X'_\lambda$ 等同起来，则上述
假设蕴含三个集合
$V_{\lambda\mu}\cap V_{\lambda\nu}$、
$V_{\mu\nu}\cap V_{\mu\lambda}$、
$V_{\nu\lambda}\cap V_{\nu\mu}$ 都等同于
$X'_\lambda\cap X'_\mu\cap X'_\nu$。现在可以把 $X_\lambda$ 的赋环空间
结构搬到 $X'_\lambda$ 上；若 $\mathcal{A}'_\lambda$ 是由
$\mathcal{A}_\lambda$ 搬来的环层，则各 $\mathcal{A}'_\lambda$ 满足
黏合条件 (3.3.1)，因而在 $X$ 上定义一个环层 $\mathcal{A}$。称
$(X,\mathcal{A})$ 是借助各 $\varphi_{\lambda\mu}$，沿各
$V_{\lambda\mu}$ \emph{黏合各
$(X_\lambda,\mathcal{A}_\lambda)$ 而得的赋环空间}。
\end{env}

\subsection[一个 A-模层的直像]{一个 $\mathcal{A}$-模层的直像}
\label{subsection:0.4.2-fr}

\begin{env}[4.2.1]
\label{0.4.2.1-fr}
设 $(X,\mathcal{A})$、$(Y,\mathcal{B})$ 是两个赋环空间，
$\Psi=(\psi,\theta)$ 是一个态射
$(X,\mathcal{A})\to(Y,\mathcal{B})$；于是 $\psi_*(\mathcal{A})$ 是
$Y$ 上的一个环层，而 $\theta$ 是环层同态
$\mathcal{B}\to\psi_*(\mathcal{A})$。再设 $\mathcal{F}$ 是一个
$\mathcal{A}$-模层；它的直像 $\psi_*(\mathcal{F})$ 是 $Y$ 上的一个
阿贝尔群层。此外，对每个开集 $U\subset Y$，
\[
  \Gamma(U,\psi_*(\mathcal{F}))
  =\Gamma(\psi^{-1}(U),\mathcal{F})
\]
带有关于环
$\Gamma(U,\psi_*(\mathcal{A}))=\Gamma(\psi^{-1}(U),\mathcal{A})$ 的
模结构；定义这些结构的双线性映射与限制运算相容，因而在
$\psi_*(\mathcal{F})$ 上定义一个 $\psi_*(\mathcal{A})$-模层结构。
同态 $\theta:\mathcal{B}\to\psi_*(\mathcal{A})$ 又使
$\psi_*(\mathcal{F})$ 带有一个 $\mathcal{B}$-模层结构；称这个
$\mathcal{B}$-模层为\emph{$\mathcal{F}$ 在态射 $\Psi$ 下的直像}，并记作
$\Psi_*(\mathcal{F})$。若 $\mathcal{F}_1$、$\mathcal{F}_2$ 是 $X$ 上的
两个 $\mathcal{A}$-模层，而 $u$ 是一个 $\mathcal{A}$-模层同态
$\mathcal{F}_1\to\mathcal{F}_2$，则立即可见（考察 $Y$ 的开集上的截面）
$\psi_*(u)$ 是一个 $\psi_*(\mathcal{A})$-模层同态
$\psi_*(\mathcal{F}_1)\to\psi_*(\mathcal{F}_2)$，\emph{从而}也是一个
$\mathcal{B}$-模层同态
$\Psi_*(\mathcal{F}_1)\to\Psi_*(\mathcal{F}_2)$；把它视为
$\mathcal{B}$-模层同态时，记作 $\Psi_*(u)$。因此，$\Psi_*$ 是从
$\mathcal{A}$-模层范畴到 $\mathcal{B}$-模层范畴的一个\emph{协变函子}。
此外，这个函子显然是\emph{左正合}的 (G, II, 2.12)。

$\psi_*(\mathcal{A})$ 上的 $\mathcal{B}$-模层结构与环层结构共同定义
一个 $\mathcal{B}$-代数层结构；把这个 $\mathcal{B}$-代数层记作
$\Psi_*(\mathcal{A})$。
\end{env}

\begin{env}[4.2.2]
\label{0.4.2.2-fr}
设 $\mathcal{M}$、$\mathcal{N}$ 是两个 $\mathcal{A}$-模层。对 $Y$ 的
每个开集 $U$，有典范映射
\[
  \Gamma(\psi^{-1}(U),\mathcal{M})
  \times\Gamma(\psi^{-1}(U),\mathcal{N})
  \longrightarrow
  \Gamma(\psi^{-1}(U),\mathcal{M}\otimes_{\mathcal{A}}\mathcal{N})
\]
它关于环
$\Gamma(\psi^{-1}(U),\mathcal{A})=\Gamma(U,\psi_*(\mathcal{A}))$
是双线性的，因而关于 $\Gamma(U,\mathcal{B})$ \emph{当然}也是双线性的；所以它
定义一个同态
\[
  \Gamma(U,\Psi_*(\mathcal{M}))
  \otimes_{\Gamma(U,\mathcal{B})}
  \Gamma(U,\Psi_*(\mathcal{N}))
  \longrightarrow
  \Gamma(U,\Psi_*(\mathcal{M}\otimes_{\mathcal{A}}\mathcal{N}))
\]
而且立即可验，这些同态与限制运算相容，因而给出一个具有函子性的
$\mathcal{B}$-模层典范同态
\[
  \Psi_*(\mathcal{M})\otimes_{\mathcal{B}}\Psi_*(\mathcal{N})
  \longrightarrow
  \Psi_*(\mathcal{M}\otimes_{\mathcal{A}}\mathcal{N}).
  \tag{4.2.2.1}\label{0.4.2.2.1-fr}
\]

\oldpage[0\textsubscript{I}]{40}
这个同态一般既不是单射也不是满射。若 $\mathcal{P}$ 是第三个
$\mathcal{A}$-模层，则立即可验图
\[
\begin{tikzcd}[column sep=large,row sep=large]
\Psi_*(\mathcal{M})\otimes_{\mathcal{B}}\Psi_*(\mathcal{N})
  \otimes_{\mathcal{B}}\Psi_*(\mathcal{P})
  \arrow[r]
  \arrow[d]
&
\Psi_*(\mathcal{M}\otimes_{\mathcal{A}}\mathcal{N})
  \otimes_{\mathcal{B}}\Psi_*(\mathcal{P})
  \arrow[d]
\\
\Psi_*(\mathcal{M})\otimes_{\mathcal{B}}
  \Psi_*(\mathcal{N}\otimes_{\mathcal{A}}\mathcal{P})
  \arrow[r]
&
\Psi_*(\mathcal{M}\otimes_{\mathcal{A}}\mathcal{N}
  \otimes_{\mathcal{A}}\mathcal{P})
\end{tikzcd}
\tag{4.2.2.2}\label{0.4.2.2.2-fr}
\]
交换。
\end{env}

\begin{env}[4.2.3]
\label{0.4.2.3-fr}
设 $\mathcal{M}$、$\mathcal{N}$ 是两个 $\mathcal{A}$-模层。按定义，
对每个开集 $U\subset Y$ 都有
\[
  \Gamma\bigl(\psi^{-1}(U),
    \mathcal{H}\!om_{\mathcal{A}}(\mathcal{M},\mathcal{N})\bigr)
  =
  \operatorname{Hom}_{\mathcal{A}|V}
    (\mathcal{M}|V,\mathcal{N}|V),
\]
这里为简化记号置 $V=\psi^{-1}(U)$；映射 $u\to\Psi_*(u)$ 是关于
$\Gamma(U,\mathcal{B})$-模结构的一个同态
\[
  \operatorname{Hom}_{\mathcal{A}|V}(\mathcal{M}|V,\mathcal{N}|V)
  \longrightarrow
  \operatorname{Hom}_{\mathcal{B}|U}
    (\Psi_*(\mathcal{M})|U,\Psi_*(\mathcal{N})|U)
\]
这些同态与限制运算相容，因而定义一个具有函子性的
$\mathcal{B}$-模层典范同态
\[
  \Psi_*\bigl(\mathcal{H}\!om_{\mathcal{A}}(\mathcal{M},\mathcal{N})\bigr)
  \longrightarrow
  \mathcal{H}\!om_{\mathcal{B}}
    (\Psi_*(\mathcal{M}),\Psi_*(\mathcal{N})).
  \tag{4.2.3.1}\label{0.4.2.3.1-fr}
\]
\end{env}

\begin{env}[4.2.4]
\label{0.4.2.4-fr}
若 $\mathcal{C}$ 是一个 $\mathcal{A}$-代数层，则复合同态
\[
  \Psi_*(\mathcal{C})\otimes_{\mathcal{B}}\Psi_*(\mathcal{C})
  \longrightarrow
  \Psi_*(\mathcal{C}\otimes_{\mathcal{A}}\mathcal{C})
  \longrightarrow
  \Psi_*(\mathcal{C})
\]
在 $\Psi_*(\mathcal{C})$ 上定义一个 $\mathcal{B}$-代数层结构，这由
(4.2.2.2) 即得。同理可见，若 $\mathcal{M}$ 是一个
$\mathcal{C}$-模层，则 $\Psi_*(\mathcal{M})$ 典范地带有一个
$\Psi_*(\mathcal{C})$-模层结构。
\end{env}

\begin{env}[4.2.5]
\label{0.4.2.5-fr}
特别考察 $X$ 是 $Y$ 的\emph{闭}子空间且 $\psi$ 是典范嵌入
$j:X\to Y$ 的情形。若
$\mathcal{B}'=\mathcal{B}|X=j^*(\mathcal{B})$ 是环层 $\mathcal{B}$ 在
$X$ 上的限制，则借助同态 $\theta^\sharp:\mathcal{B}'\to\mathcal{A}$，
可以把一个 $\mathcal{A}$-模层 $\mathcal{M}$ 看作
$\mathcal{B}'$-模层；此时 $\Psi_*(\mathcal{M})$ 是这样一个
$\mathcal{B}$-模层：它在 $X$ 上诱导出 $\mathcal{M}$，而在别处为 $0$。
若 $\mathcal{N}$ 是另一个 $\mathcal{A}$-模层，则
$\Psi_*(\mathcal{M})\otimes_{\mathcal{B}}\Psi_*(\mathcal{N})$ 可等同于
$\Psi_*(\mathcal{M}\otimes_{\mathcal{B}'}\mathcal{N})$，而
$\mathcal{H}\!om_{\mathcal{B}}(\Psi_*(\mathcal{M}),\Psi_*(\mathcal{N}))$
可等同于
$\Psi_*\bigl(\mathcal{H}\!om_{\mathcal{B}'}(\mathcal{M},\mathcal{N})\bigr)$。
\end{env}

\begin{env}[4.2.6]
\label{0.4.2.6-fr}
设 $(Z,\mathcal{C})$ 是第三个赋环空间，$\Psi'=(\psi',\theta')$ 是一个
态射 $(Y,\mathcal{B})\to(Z,\mathcal{C})$；若 $\Psi''$ 是复合态射
$\Psi'\circ\Psi$，则显然有 $\Psi''_*=\Psi'_*\circ\Psi_*$。
\end{env}

\oldpage[0\textsubscript{I}]{41}
\subsection[一个 B-模层的逆像]{一个 $\mathcal{B}$-模层的逆像}
\label{subsection:0.4.3-fr}

\begin{env}[4.3.1]
\label{0.4.3.1-fr}
沿用 (4.2.1) 的假设和记号，设 $\mathcal{G}$ 是一个
$\mathcal{B}$-模层，而 $\psi^*(\mathcal{G})$ 是它的逆像 (3.7.1)，
因而是 $X$ 上的一个阿贝尔群层。$\psi^*(\mathcal{G})$ 与
$\psi^*(\mathcal{B})$ 的截面之定义 (3.7.1) 表明，
$\psi^*(\mathcal{G})$ 典范地带有一个 $\psi^*(\mathcal{B})$-模层结构。
另一方面，同态 $\theta^\sharp:\psi^*(\mathcal{B})\to\mathcal{A}$
使 $\mathcal{A}$ 带有一个 $\psi^*(\mathcal{B})$-模层结构；在需要避免
混淆时，将后者记作 $\mathcal{A}_{[\theta]}$。于是张量积
$\psi^*(\mathcal{G})\otimes_{\psi^*(\mathcal{B})}\mathcal{A}_{[\theta]}$
带有一个 $\mathcal{A}$-模层结构。称这个 $\mathcal{A}$-模层为
$\mathcal{G}$ 在态射 $\Psi$ 下的逆像，并记作 $\Psi^*(\mathcal{G})$。
若 $\mathcal{G}_1$、$\mathcal{G}_2$ 是 $Y$ 上的两个
$\mathcal{B}$-模层，而 $v$ 是一个 $\mathcal{B}$-模层同态
$\mathcal{G}_1\to\mathcal{G}_2$，则容易验证，$\psi^*(v)$ 是从
$\psi^*(\mathcal{G}_1)$ 到 $\psi^*(\mathcal{G}_2)$ 的一个
$\psi^*(\mathcal{B})$-模层同态；所以 $\psi^*(v)\otimes1$ 是一个
$\mathcal{A}$-模层同态
$\Psi^*(\mathcal{G}_1)\to\Psi^*(\mathcal{G}_2)$，将它记作
$\Psi^*(v)$。由此把 $\Psi^*$ 定义为从 $\mathcal{B}$-模层范畴到
$\mathcal{A}$-模层范畴的一个\emph{协变函子}。这里，这个函子
（不同于 $\psi^*$）一般不再正合，而只是\emph{右正合}的，因为同
$\mathcal{A}$ 作张量积是 $\psi^*(\mathcal{B})$-模层范畴中的一个
右正合函子。

对每个 $x\in X$，由 (3.7.2) 有
\[
  (\Psi^*(\mathcal{G}))_x
  =\mathcal{G}_{\psi(x)}\otimes_{\mathcal{B}_{\psi(x)}}\mathcal{A}_x,
\]
所以 $\Psi^*(\mathcal{G})$ 的支撑包含于
$\psi^{-1}(\operatorname{Supp}\mathcal{G})$。
\end{env}

\begin{env}[4.3.2]
\label{0.4.3.2-fr}
设 $(\mathcal{G}_\lambda)$ 是 $\mathcal{B}$-模层的一个归纳系，而
$\mathcal{G}=\varinjlim\mathcal{G}_\lambda$ 是它的归纳极限。各典范同态
$\mathcal{G}_\lambda\to\mathcal{G}$ 定义了
$\psi^*(\mathcal{B})$-模层同态
$\psi^*(\mathcal{G}_\lambda)\to\psi^*(\mathcal{G})$，它们给出一个
典范同态
\[
  \varinjlim\psi^*(\mathcal{G}_\lambda)
  \longrightarrow
  \psi^*(\mathcal{G}).
\]
由于层的归纳极限在一点的茎，是各层在同一点的茎之归纳极限
(G, II, 1.11)，所以前一个典范同态是\emph{双射} (3.7.2)。此外，
张量积与层的归纳极限交换，因而有 $\mathcal{A}$-模层的一个
具有函子性的\emph{典范同构}
\[
  \varinjlim\Psi^*(\mathcal{G}_\lambda)
  \xrightarrow{\sim}
  \Psi^*\bigl(\varinjlim\mathcal{G}_\lambda\bigr)
\]

另一方面，对 $\mathcal{B}$-模层的一个有限直和
$\bigoplus_i\mathcal{G}_i$，显然有
\[
  \psi^*\bigl(\bigoplus_i\mathcal{G}_i\bigr)
  =\bigoplus_i\psi^*(\mathcal{G}_i)~;
\]
所以再同 $\mathcal{A}_{[\theta]}$ 作张量积，便有
\[
  \Psi^*\bigl(\bigoplus_i\mathcal{G}_i\bigr)
  =\bigoplus_i\Psi^*(\mathcal{G}_i).
  \tag{4.3.2.1}\label{0.4.3.2.1-fr}
\]
根据上述结论，取归纳极限便知，前一个等式对\emph{任意}直和仍然成立。
\end{env}

\begin{env}[4.3.3]
\label{0.4.3.3-fr}
设 $\mathcal{G}_1$、$\mathcal{G}_2$ 是两个 $\mathcal{B}$-模层；由
阿贝尔群层逆像的定义 (3.7.1)，立即得到
$\psi^*(\mathcal{B})$-模层的一个典范同态
\[
  \psi^*(\mathcal{G}_1)\otimes_{\psi^*(\mathcal{B})}\psi^*(\mathcal{G}_2)
  \longrightarrow
  \psi^*(\mathcal{G}_1\otimes_{\mathcal{B}}\mathcal{G}_2)
\]
而层的张量积在一点的茎，就是各层在该点的茎之张量积
(G, II, 2.8)，所以由 (3.7.2) 可知，前一个同态实际上是一个
\emph{同构}。再同 $\mathcal{A}$ 作张量积，便得到一个具有函子性的
\emph{典范同构}
\[
  \Psi^*(\mathcal{G}_1)\otimes_{\mathcal{A}}\Psi^*(\mathcal{G}_2)
  \xrightarrow{\sim}
  \Psi^*(\mathcal{G}_1\otimes_{\mathcal{B}}\mathcal{G}_2).
  \tag{4.3.3.1}\label{0.4.3.3.1-fr}
\]
\end{env}

\begin{env}[4.3.4]
\label{0.4.3.4-fr}
设 $\mathcal{C}$ 是一个 $\mathcal{B}$-代数层；给定
$\mathcal{C}$ 上的代数层结构，等价于给定一个 $\mathcal{B}$-模层同态
\[
  \mathcal{C}\otimes_{\mathcal{B}}\mathcal{C}\longrightarrow\mathcal{C}
\]
满足结合性和交换性条件（这些条件可以逐茎验证）。前一个同构使我们可以
把这个同态搬成一个 $\mathcal{A}$-模层同态
\[
  \Psi^*(\mathcal{C})\otimes_{\mathcal{A}}\Psi^*(\mathcal{C})
  \longrightarrow\Psi^*(\mathcal{C})
\]
并满足同样的条件，从而使 $\Psi^*(\mathcal{C})$ 带有一个
$\mathcal{A}$-代数层结构。特别地，由定义立即可知，
$\mathcal{A}$-代数层 $\Psi^*(\mathcal{B})$ 在典范同构意义下等于
$\mathcal{A}$。

同样，若 $\mathcal{M}$ 是一个 $\mathcal{C}$-模层，则给定这种模层结构
等价于
\oldpage[0\textsubscript{I}]{42}
给定一个满足结合性条件的 $\mathcal{B}$-模层同态
$\mathcal{C}\otimes_{\mathcal{B}}\mathcal{M}\to\mathcal{M}$；把它搬过去，
便在 $\Psi^*(\mathcal{M})$ 上得到一个
$\Psi^*(\mathcal{C})$-模层结构。
\end{env}

\begin{env}[4.3.5]
\label{0.4.3.5-fr}
设 $\mathcal{J}$ 是 $\mathcal{B}$ 的一个理想层；由于函子
$\psi^*$ 正合，$\psi^*(\mathcal{B})$-模层
$\psi^*(\mathcal{J})$ 典范地等同于 $\psi^*(\mathcal{B})$ 的一个
理想层；于是典范单射
$\psi^*(\mathcal{J})\to\psi^*(\mathcal{B})$ 给出一个
$\mathcal{A}$-模层同态
\[
  \Psi^*(\mathcal{J})
  =\psi^*(\mathcal{J})\otimes_{\psi^*(\mathcal{B})}\mathcal{A}_{[\theta]}
  \longrightarrow\mathcal{A}~;
\]
把 $\Psi^*(\mathcal{J})$ 在这个同态下的像记作
$\Psi^*(\mathcal{J})\mathcal{A}$，在不致混淆时也记作
$\mathcal{J}\mathcal{A}$。所以按定义有
\[
  \mathcal{J}\mathcal{A}
  =\theta^\sharp\bigl(\psi^*(\mathcal{J})\bigr)\mathcal{A}
\]
特别地，考虑到 $\psi^*(\mathcal{J})$ 的茎与 $\mathcal{J}$ 的相应茎
之间的典范等同 (3.7.2)，对每个 $x\in X$ 都有
\[
  (\mathcal{J}\mathcal{A})_x
  =\theta_x^\sharp(\mathcal{J}_{\psi(x)})\mathcal{A}_x,
\]
若 $\mathcal{J}_1$、$\mathcal{J}_2$ 是 $\mathcal{B}$ 的两个理想层，
则有
\[
  (\mathcal{J}_1\mathcal{J}_2)\mathcal{A}
  =\mathcal{J}_1(\mathcal{J}_2\mathcal{A})
  =(\mathcal{J}_1\mathcal{A})(\mathcal{J}_2\mathcal{A}).
\]
若 $\mathcal{F}$ 是一个 $\mathcal{A}$-模层，则置
$\mathcal{J}\mathcal{F}=(\mathcal{J}\mathcal{A})\mathcal{F}$。
\end{env}

\begin{env}[4.3.6]
\label{0.4.3.6-fr}
设 $(Z,\mathcal{C})$ 是第三个赋环空间，$\Psi'=(\psi',\theta')$ 是一个
态射 $(Y,\mathcal{B})\to(Z,\mathcal{C})$；若 $\Psi''$ 是复合态射
$\Psi'\circ\Psi$，则由定义 (4.3.1) 和 (4.3.3.1) 可知
${\Psi''}^*=\Psi^*\circ{\Psi'}^*$。
\end{env}

\subsection{直像与逆像之间的关系}
\label{subsection:0.4.4-fr}

\begin{env}[4.4.1]
\label{0.4.4.1-fr}
沿用 (4.2.1) 的假设和记号，设 $\mathcal{G}$ 是一个
$\mathcal{B}$-模层。按定义，$\mathcal{B}$-模层同态
$u:\mathcal{G}\to\Psi_*(\mathcal{F})$ 仍称为从 $\mathcal{G}$ 到
$\mathcal{F}$ 的一个 $\Psi$-\emph{态射}，也简称为从
$\mathcal{G}$ 到 $\mathcal{F}$ 的一个\emph{同态}；在不致混淆时写作
$u:\mathcal{G}\to\mathcal{F}$。给出这样的同态，等价于对每一对
$(U,V)$ 给出一个\emph{同态}
\[
  u_{U,V}:\Gamma(V,\mathcal{G})\longrightarrow\Gamma(U,\mathcal{F})
\]

这里 $U$ 是 $X$ 的开集，$V$ 是 $Y$ 的开集，且
$\psi(U)\subset V$；这个同态是
$\Gamma(V,\mathcal{B})$-\emph{模同态}，其中借助环同态
$\theta_{U,V}:\Gamma(V,\mathcal{B})\to\Gamma(U,\mathcal{A})$，把
$\Gamma(U,\mathcal{F})$ 看作 $\Gamma(V,\mathcal{B})$-模。此外，各
$u_{U,V}$ 还须使图 (3.5.1.1) 交换。事实上，为定义 $u$，只须在
$U$（相应地，$V$）遍历 $X$（相应地，$Y$）的一个拓扑基
$\mathfrak{B}$（相应地，$\mathfrak{B}'$）时给出各 $u_{U,V}$，并对
这些限制验证 (3.5.1.1) 的交换性。
\end{env}

\begin{env}[4.4.2]
\label{0.4.4.2-fr}
在 (4.2.1) 和 (4.2.6) 的假设下，设 $\mathcal{H}$ 是一个
$\mathcal{C}$-模层，而
$v:\mathcal{H}\to\Psi'_*(\mathcal{G})$ 是一个 $\Psi'$-态射；则
\[
  w:\mathcal{H}\xrightarrow{v}\Psi'_*(\mathcal{G})
  \xrightarrow{\Psi'_*(u)}\Psi'_*(\Psi_*(\mathcal{F}))
\]
是一个 $\Psi''$-态射，称为 $u$ 与 $v$ 的\emph{复合}。
\end{env}

\begin{env}[4.4.3]
\label{0.4.4.3-fr}
现在说明，可以定义关于 $\mathcal{F}$ 和 $\mathcal{G}$ 的\emph{双函子}
之间的一个典范\emph{同构}
\[
  \operatorname{Hom}_{\mathcal{A}}(\Psi^*(\mathcal{G}),\mathcal{F})
  \xrightarrow{\sim}
  \operatorname{Hom}_{\mathcal{B}}(\mathcal{G},\Psi_*(\mathcal{F}))
  \tag{4.4.3.1}\label{0.4.4.3.1-fr}
\]
将其记作 $v\mapsto v_\theta^\flat$（在不致混淆时简记作
$v\mapsto v^\flat$）；其逆同构记作 $u\mapsto u_\theta^\sharp$，
或简记作 $u\mapsto u^\sharp$。定义如下：把
$v:\Psi^*(\mathcal{G})\to\mathcal{F}$ 与典范映射
$\psi^*(\mathcal{G})\to\Psi^*(\mathcal{G})$ 复合，得到群层同态
$v':\psi^*(\mathcal{G})\to\mathcal{F}$，它也是一个
$\psi^*(\mathcal{B})$-模层同态。由此得到 (3.7.1) 一个同态
${v'}^\flat:\mathcal{G}\to\psi_*(\mathcal{F})=\Psi_*(\mathcal{F})$；
容易验证，它也是一个 $\mathcal{B}$-模层同态；
\oldpage[0\textsubscript{I}]{43}
置 $v_\theta^\flat={v'}^\flat$。同样，由
$u:\mathcal{G}\to\Psi_*(\mathcal{F})$ 这个 $\mathcal{B}$-模层同态，
得到 (3.7.1) 一个 $\psi^*(\mathcal{B})$-模层同态
$u^\sharp:\psi^*(\mathcal{G})\to\mathcal{F}$，再同
$\mathcal{A}$ 作张量积，便得到一个 $\mathcal{A}$-模层同态
$\Psi^*(\mathcal{G})\to\mathcal{F}$，将它记作
$u_\theta^\sharp$。立即可验
$(u_\theta^\sharp)_\theta^\flat=u$ 和
$(v_\theta^\flat)_\theta^\sharp=v$，以及同构
$v\mapsto v_\theta^\flat$ 关于 $\mathcal{F}$ 的函子性。于是
$u\mapsto u_\theta^\sharp$ 关于 $\mathcal{G}$ 的函子性，也如 (3.5.4)
那样形式地推出（这套论证也可以证明在 (4.3.1) 中直接建立的
$\Psi^*$ 的函子性）。

若取 $v$ 为 $\Psi^*(\mathcal{G})$ 的恒等同态，则
$v_\theta^\flat$ 是一个同态
\[
  \rho_{\mathcal{G}}:\mathcal{G}\longrightarrow
  \Psi_*(\Psi^*(\mathcal{G}))~;
  \tag{4.4.3.2}\label{0.4.4.3.2-fr}
\]
若取 $u$ 为 $\Psi_*(\mathcal{F})$ 的恒等同态，则
$u_\theta^\sharp$ 是一个同态
\[
  \sigma_{\mathcal{F}}:\Psi^*(\Psi_*(\mathcal{F}))
  \longrightarrow\mathcal{F}~;
  \tag{4.4.3.3}\label{0.4.4.3.3-fr}
\]
称这些同态为\emph{典范}同态。它们一般既不是单射也不是满射。还有与
(3.5.3.3) 和 (3.5.4.4) 类似的典范分解。

注意，若 $s$ 是 $\mathcal{G}$ 在 $Y$ 的开集 $V$ 上的一个截面，则
$\rho_{\mathcal{G}}(s)$ 是 $\Psi^*(\mathcal{G})$ 在
$\psi^{-1}(V)$ 上的截面 $s'\otimes1$，这里 $s'$ 对每个
$x\in\psi^{-1}(V)$ 都有 $s'_x=s_{\psi(x)}$。又若
$u:\mathcal{G}\to\psi_*(\mathcal{F})$ 是一个同态，则它对每个
$x\in X$ 都在茎上定义一个同态
$u_x:\mathcal{G}_{\psi(x)}\to\mathcal{F}_x$；它由
$(u^\sharp)_x:(\Psi^*(\mathcal{G}))_x\to\mathcal{F}_x$ 与典范同态
$s_x\mapsto s_x\otimes1$（从 $\mathcal{G}_{\psi(x)}$ 到
$(\Psi^*(\mathcal{G}))_x
=\mathcal{G}_{\psi(x)}\otimes_{\mathcal{B}_{\psi(x)}}\mathcal{A}_x$）
复合而得。也可以从同态
\[
  \Gamma(V,\mathcal{G})\xrightarrow{u}
  \Gamma(\psi^{-1}(V),\mathcal{F})\longrightarrow\mathcal{F}_x,
\]
取归纳极限得到 $u_x$，这里 $V$ 遍历 $\psi(x)$ 的各邻域。
\end{env}

\begin{env}[4.4.4]
\label{0.4.4.4-fr}
设 $\mathcal{F}_1$、$\mathcal{F}_2$ 是 $\mathcal{A}$-模层，
$\mathcal{G}_1$、$\mathcal{G}_2$ 是 $\mathcal{B}$-模层，而 $u_i$
（$i=1,2$）是从 $\mathcal{G}_i$ 到 $\mathcal{F}_i$ 的一个同态。
用 $u_1\otimes u_2$ 表示同态
\[
  u:\mathcal{G}_1\otimes_{\mathcal{B}}\mathcal{G}_2
  \longrightarrow
  \mathcal{F}_1\otimes_{\mathcal{A}}\mathcal{F}_2
\]
使得 $u^\sharp=(u_1)^\sharp\otimes(u_2)^\sharp$（这里利用
(4.3.3.1)）。可以验证，$u$ 也等于复合
\[
  \mathcal{G}_1\otimes_{\mathcal{B}}\mathcal{G}_2
  \longrightarrow
  \Psi_*(\mathcal{F}_1)\otimes_{\mathcal{B}}\Psi_*(\mathcal{F}_2)
  \longrightarrow
  \Psi_*(\mathcal{F}_1\otimes_{\mathcal{A}}\mathcal{F}_2),
\]
其中第一个箭头是通常的张量积 $u_1\otimes_{\mathcal{B}}u_2$，第二个
箭头是典范同态 (4.2.2.1)。
\end{env}

\begin{env}[4.4.5]
\label{0.4.4.5-fr}
设 $(\mathcal{G}_\lambda)_{\lambda\in L}$ 是 $\mathcal{B}$-模层的一个
归纳系；对每个 $\lambda\in L$，设 $u_\lambda$ 是一个同态
$\mathcal{G}_\lambda\to\Psi_*(\mathcal{F})$，
并且这些同态组成一个归纳系。置
$\mathcal{G}=\varinjlim\mathcal{G}_\lambda$ 和
$u=\varinjlim u_\lambda$；则各 $(u_\lambda)^\sharp$ 组成同态
$\Psi^*(\mathcal{G}_\lambda)\to\mathcal{F}$ 的一个归纳系，而
这个归纳系的归纳极限正是 $u^\sharp$。
\end{env}

\begin{env}[4.4.6]
\label{0.4.4.6-fr}
设 $\mathcal{M}$、$\mathcal{N}$ 是两个 $\mathcal{B}$-模层，$V$ 是
$Y$ 的一个开集，并置 $U=\psi^{-1}(V)$；映射
$v\mapsto\Psi^*(v)$ 是关于 $\Gamma(V,\mathcal{B})$-模结构的一个同态
\[
  \operatorname{Hom}_{\mathcal{B}|V}(\mathcal{M}|V,\mathcal{N}|V)
  \longrightarrow
  \operatorname{Hom}_{\mathcal{A}|U}
  (\Psi^*(\mathcal{M})|U,\Psi^*(\mathcal{N})|U)
\]
\begin{sloppypar}
（$\operatorname{Hom}_{\mathcal{A}|U}
(\Psi^*(\mathcal{M})|U,\allowbreak\Psi^*(\mathcal{N})|U)$ 本来带有一个
$\Gamma(U,\psi^*(\mathcal{B}))$-模结构，而借助典范同态 (3.7.2)
\oldpage[0\textsubscript{I}]{44}
$\Gamma(V,\mathcal{B})\to\Gamma(U,\psi^*(\mathcal{B}))$，它也成为一个
$\Gamma(V,\mathcal{B})$-模）。
\end{sloppypar}
立即可见，这些同态与限制相容，因而定义一个具有函子性的典范同态
\[
  \gamma:\mathcal{H}\!om_{\mathcal{B}}(\mathcal{M},\mathcal{N})
  \longrightarrow
  \Psi_*\bigl(\mathcal{H}\!om_{\mathcal{A}}
  (\Psi^*(\mathcal{M}),\Psi^*(\mathcal{N}))\bigr)~;
\]
因而，与这个同态对应的还有同态
\[
  \gamma^\sharp:
  \Psi^*\bigl(\mathcal{H}\!om_{\mathcal{B}}(\mathcal{M},\mathcal{N})\bigr)
  \longrightarrow
  \mathcal{H}\!om_{\mathcal{A}}
  (\Psi^*(\mathcal{M}),\Psi^*(\mathcal{N}))
\]
而且这些典范同态关于 $\mathcal{M}$ 和 $\mathcal{N}$ 都具有函子性。
\end{env}

\begin{env}[4.4.7]
\label{0.4.4.7-fr}
设 $\mathcal{F}$（相应地，$\mathcal{G}$）是一个
$\mathcal{A}$-代数层（相应地，$\mathcal{B}$-代数层）。若
$u:\mathcal{G}\to\Psi_*(\mathcal{F})$ 是 $\mathcal{B}$-代数层同态，
则 $u^\sharp$ 是一个 $\mathcal{A}$-代数层同态
$\Psi^*(\mathcal{G})\to\mathcal{F}$；这由图
\[
\begin{tikzcd}[column sep=huge,row sep=large]
\mathcal{G}\otimes_{\mathcal{B}}\mathcal{G}
  \arrow[r] \arrow[d]
  & \mathcal{G} \arrow[d,"u"] \\
\Psi_*(\mathcal{F}\otimes_{\mathcal{A}}\mathcal{F})
  \arrow[r]
  & \Psi_*(\mathcal{F})
\end{tikzcd}
\]
的交换性与 (4.4.4) 即得。同样，若
$v:\Psi^*(\mathcal{G})\to\mathcal{F}$ 是一个 $\mathcal{A}$-代数层同态，
则 $v^\flat:\mathcal{G}\to\Psi_*(\mathcal{F})$ 是一个
$\mathcal{B}$-代数层同态。
\end{env}

\begin{env}[4.4.8]
\label{0.4.4.8-fr}
设 $(Z,\mathcal{C})$ 是第三个赋环空间，$\Psi'=(\psi',\theta')$ 是一个
态射 $(Y,\mathcal{B})\to(Z,\mathcal{C})$，而
$\Psi'':(X,\mathcal{A})\to(Z,\mathcal{C})$ 是复合态射
$\Psi'\circ\Psi$。设 $\mathcal{H}$ 是一个 $\mathcal{C}$-模层，
% EGA1-ERR-0003: French print introduces u' here, but the composite and both following arrows use v'. Canonical translation preserves the printed variables.
$u'$ 是从 $\mathcal{H}$ 到 $\mathcal{G}$ 的一个同态；按定义，复合
$v''=v\circ v'$ 是由
\[
  \mathcal{H}\xrightarrow{v'}\Psi'_*(\mathcal{G})
  \xrightarrow{\Psi'_*(v)}\Psi'_*(\Psi_*(\mathcal{F}))~;
\]
定义的从 $\mathcal{H}$ 到 $\mathcal{F}$ 的同态。可以验证，
${v''}^\sharp$ 是同态
\[
  \Psi^*({\Psi'}^*(\mathcal{H}))
  \xrightarrow{\Psi^*({v'}^\sharp)}
  \Psi^*(\mathcal{G})
  \xrightarrow{v^\sharp}\mathcal{F}.
\]
\end{env}

\section{拟凝聚层与凝聚层}
\label{section:0.5-fr}

\subsection{拟凝聚层}
\label{subsection:0.5.1-fr}

\begin{env}[5.1.1]
\label{0.5.1.1-fr}
设 $(X,\mathcal{O}_X)$ 是一个赋环空间，$\mathcal{F}$ 是一个
$\mathcal{O}_X$-模层。给定一个同态
$u:\mathcal{O}_X\to\mathcal{F}$ 的数据，其中它是一个
$\mathcal{O}_X$-模层同态，等价于给定截面
$s=u(1)\in\Gamma(X,\mathcal{F})$。事实上，给定 $s$ 后，对于任意截面
$t\in\Gamma(U,\mathcal{O}_X)$，必有 $u(t)=t\cdot(s|U)$；称 $u$
\emph{由截面 $s$ 定义}。现在设 $I$ 是任意指标集，考虑直和层
$\mathcal{O}_X^{(I)}$；对于每个 $i\in I$，设 $h_i$ 是第 $i$ 个因子到
$\mathcal{O}_X^{(I)}$ 的典范嵌入。已知 $u\mapsto(u\circ h_i)$ 是从
$\operatorname{Hom}_{\mathcal{O}_X}(\mathcal{O}_X^{(I)},\mathcal{F})$
到乘积
$\bigl(\operatorname{Hom}_{\mathcal{O}_X}(\mathcal{O}_X,\mathcal{F})\bigr)^I$
的一个同构。因此，同态
$u:\mathcal{O}_X^{(I)}\to\mathcal{F}$ 与
\emph{截面族 $(s_i)_{i\in I}$，它们是 $\mathcal{F}$ 在 $X$ 上的截面}
之间有一个典范双射。同态 $u$（与 $(s_i)$ 对应）把元素
$(a_i)\in(\Gamma(U,\mathcal{O}_X))^{(I)}$ 映为
$\sum_{i\in I}a_i\cdot(s_i|U)$。

称 $\mathcal{F}$ \emph{由族 $(s_i)$ 生成}，如果由这个族定义的同态
\oldpage[0\textsubscript{I}]{45}
$\mathcal{O}_X^{(I)}\to\mathcal{F}$ 是\emph{满射}（换言之，对每个 $x\in X$，
$\mathcal{F}_x$ 是一个 $\mathcal{O}_x$-模，由各 $(s_i)_x$ 生成）。
又称 $\mathcal{F}$ \emph{由其在 $X$ 上的截面生成}，如果它由所有这些
截面组成的族（或其中某个子族）生成；换言之，
存在一个满同态 $\mathcal{O}_X^{(I)}\to\mathcal{F}$，其中 $I$ 适当。

应当注意，可能存在一个 $\mathcal{O}_X$-模层 $\mathcal{F}$ 和一点
$x_0\in X$，使得 $\mathcal{F}|U$ 不由其在 $U$ 上的截面生成，
\emph{无论 $U$ 是 $x_0$ 的哪一个开邻域}：只需取
$X=\mathbb{R}$，取 $\mathcal{O}_X$ 为常值层 $\mathbb{Z}$，并取
$\mathcal{F}$ 为 $\mathcal{O}_X$ 的代数子层，使得
$\mathcal{F}_0=\{0\}$、$\mathcal{F}_x=\mathbb{Z}$（$x\ne0$），最后取
$x_0=0$；$\mathcal{F}|U$ 在 $U$ 上唯一的截面是 $0$，其中 $U$ 是
$0$ 的一个邻域。
\end{env}

\begin{env}[5.1.2]
\label{0.5.1.2-fr}
设 $f:X\to Y$ 是赋环空间的一个态射。若 $\mathcal{F}$ 是一个
$\mathcal{O}_X$-模层，且由其在 $X$ 上的截面生成，则典范同态
$f^*(f_*(\mathcal{F}))\to\mathcal{F}$ (4.4.3.3) 是\emph{满射}：
事实上，采用 (5.1.1) 的记号，$s_i\otimes1$ 是
$f^*(f_*(\mathcal{F}))$ 在 $X$ 上的一个截面，而它在
$\mathcal{F}$ 中的像是 $s_i$。在 (5.1.1) 的例子中令 $f$ 为恒等态射，
便表明这条命题的逆命题一般并不成立。
\end{env}

\begin{env}[5.1.3]
\label{0.5.1.3-fr}
称一个 $\mathcal{O}_X$-模层 $\mathcal{F}$ 为\emph{拟凝聚的}，如果
对每个 $x\in X$，都存在一个开邻域 $U$，它含有 $x$，使
% EGA1-ERR-0004: French print introduces U here but uses V in both following restrictions. Canonical translation preserves the printed variables.
$\mathcal{F}|V$ 同构于一个形如
$\mathcal{O}_X^{(I)}|V\to\mathcal{O}_X^{(J)}|V$ 的同态的\emph{余核}，
其中 $I$、$J$ 是任意指标集。显然，$\mathcal{O}_X$ 本身是一个
拟凝聚 $\mathcal{O}_X$-模层，而拟凝聚 $\mathcal{O}_X$-模层的任意直和
仍是拟凝聚 $\mathcal{O}_X$-模层。若 $\mathcal{O}_X$-代数层
$\mathcal{A}$ 作为 $\mathcal{O}_X$-模层是拟凝聚的，则称
它是\emph{拟凝聚的}。
\end{env}

\begin{env}[5.1.4]
\label{0.5.1.4-fr}
设 $f:X\to Y$ 是赋环空间的一个态射。若 $\mathcal{G}$ 是一个拟凝聚
$\mathcal{O}_Y$-模层，则 $f^*(\mathcal{G})$ 是一个拟凝聚
$\mathcal{O}_X$-模层。事实上，对每个 $x\in X$，都存在一个开邻域
$V$，它是 $f(x)$ 在 $Y$ 中的邻域，使得 $\mathcal{G}|V$ 是一个同态
$\mathcal{O}_Y^{(I)}|V\to\mathcal{O}_Y^{(J)}|V$ 的余核。若
$U=f^{-1}(V)$，而 $f_U$ 是 $f$ 在 $U$ 上的限制，则
$f^*(\mathcal{G})|U=f_U^*(\mathcal{G}|V)$；由于 $f_U^*$ 右正合且与
直和交换，$f_U^*(\mathcal{G}|V)$ 是一个同态
$\mathcal{O}_X^{(I)}|U\to\mathcal{O}_X^{(J)}|U$ 的余核。
\end{env}

\subsection{有限型层}
\label{subsection:0.5.2-fr}

\begin{env}[5.2.1]
\label{0.5.2.1-fr}
称一个 $\mathcal{O}_X$-模层 $\mathcal{F}$ 是\emph{有限型的}，如果对每个
$x\in X$，都存在一个开邻域 $U$，它含有 $x$，使得
$\mathcal{F}|U$ 由它在 $U$ 上的一个\emph{有限}截面族生成；等价地，
它同构于一个形如 $(\mathcal{O}_X|U)^p$ 的层的商层，其中 $p$ 有限。
有限型层的每个商层仍是有限型层；有限型层的任意有限直和与任意有限
张量积也都是有限型层。一个有限型 $\mathcal{O}_X$-模层未必拟凝聚；
$\mathcal{O}_X$-模层 $\mathcal{O}_X/\mathcal{F}$ 就是一个例子，其中
$\mathcal{F}$ 是 (5.1.1) 的例子。若 $\mathcal{F}$ 是有限型的，则
$\mathcal{F}_x$ 是一个有限型 $\mathcal{O}_x$-模（对每个 $x\in X$）；但 (5.1.1)
的例子表明，这个必要条件一般并不充分。
\end{env}

\begin{env}[5.2.2]
\label{0.5.2.2-fr}
设 $\mathcal{F}$ 是一个\emph{有限型} $\mathcal{O}_X$-模层。若 $s_i$
$(1\leq i\leq n)$ 是 $\mathcal{F}$ 在一个开邻域 $U$ 上的截面，该邻域
含有 $x\in X$，且各 $(s_i)_x$ 生成 $\mathcal{F}_x$，则存在一个开邻域
$V\subset U$，它含有 $x$，使得各 $(s_i)_y$ 都生成 $\mathcal{F}_y$
（对每个 $y\in V$）（FAC, I, 2, 12, prop.~1）。特别地，由此可知
$\mathcal{F}$ 的支撑是\emph{闭的}。

\oldpage[0\textsubscript{I}]{46}

同样，若 $u:\mathcal{F}\to\mathcal{G}$ 是一个满足 $u_x=0$ 的同态，
则存在一个邻域 $U$，它含有 $x$，使得 $u_y=0$ 对每个 $y\in U$ 成立。
\end{env}

\begin{env}[5.2.3]
\label{0.5.2.3-fr}
设 $X$ \emph{拟紧}，而 $\mathcal{F}$、$\mathcal{G}$ 是两个
$\mathcal{O}_X$-模层，其中 $\mathcal{G}$ \emph{有限型}，且
$u:\mathcal{F}\to\mathcal{G}$ 是一个\emph{满}同态。再设
$\mathcal{F}$ 是一个归纳系 $(\mathcal{F}_\lambda)$ 的归纳极限，而该
归纳系由 $\mathcal{O}_X$-模层组成。则存在一个指标 $\mu$，使同态
$\mathcal{F}_\mu\to\mathcal{G}$ 是\emph{满射}。事实上，对每个
$x\in X$，存在一个有限截面族 $s_i$，它们是 $\mathcal{G}$ 在一个开邻域
$U(x)$ 上的截面，该邻域含有 $x$；各 $(s_i)_y$ 都生成
$\mathcal{G}_y$（对每个 $y\in U(x)$）。因而存在一个开邻域
$V(x)\subset U(x)$，它含有 $x$，以及 $n$ 个截面 $t_i$，它们是
$\mathcal{F}$ 在 $V(x)$ 上的截面，使得 $s_i|V(x)=u(t_i)$ 对每个
$i$ 都成立；还可以假定各 $t_i$
都是同一个层 $\mathcal{F}_{\lambda(x)}$ 在 $V(x)$ 上的截面的典范像。
现在覆盖 $X$：取有限多个邻域 $V(x_k)$；并令 $\mu$ 是一个大于所有
$\lambda(x_k)$ 的指标；显然，这个指标符合要求。

仍设 $X$ 拟紧，并设 $\mathcal{F}$ 是一个有限型
$\mathcal{O}_X$-模层，由其在 $X$ 上的截面生成 (5.1.1)；则 $\mathcal{F}$ 由这些截面的
一个\emph{有限}子族生成：事实上，只需覆盖 $X$：取有限多个开邻域
$U_k$；使得对每个 $k$，存在有限多个截面 $s_{ik}$，它们是
$\mathcal{F}$ 在 $X$ 上的截面，其在 $U_k$ 上的限制生成
$\mathcal{F}|U_k$；显然，这些
$s_{ik}$ 便生成 $\mathcal{F}$。
\end{env}

\begin{env}[5.2.4]
\label{0.5.2.4-fr}
设 $f:X\to Y$ 是赋环空间的一个态射。若 $\mathcal{G}$ 是有限型
$\mathcal{O}_Y$-模层，则 $f^*(\mathcal{G})$ 是有限型
$\mathcal{O}_X$-模层。事实上，对每个 $x\in X$，存在一个开邻域 $V$，
它是 $f(x)$ 在 $Y$ 中的邻域，并存在一个满同态
$v:\mathcal{O}_Y^p|V\to\mathcal{G}|V$。若 $U=f^{-1}(V)$，且 $f_U$ 是
$f$ 在 $U$ 上的限制，则
$f^*(\mathcal{G})|U=f_U^*(\mathcal{G}|V)$；由于
% Published EGA II erratum: the French print has f_U here; the correction is f_U^*. Canonical translation preserves the printed expression.
$f_U$ 右正合 (4.3.1) 且与直和交换 (4.3.2)，$f_U^*(v)$ 是满同态
$\mathcal{O}_X^p|U\to f^*(\mathcal{G})|U$。
\end{env}

\begin{env}[5.2.5]
\label{0.5.2.5-fr}
称一个 $\mathcal{O}_X$-模层 $\mathcal{F}$ \emph{有有限表示}，如果对
每个 $x\in X$，都存在一个开邻域 $U$，它含有 $x$，使得
$\mathcal{F}|U$ 同构于一个
\emph{$(\mathcal{O}_X|U)$-同态的余核}
$\mathcal{O}_X^p|U\to\mathcal{O}_X^q|U$，其中 $p$、$q$ 是两个
$>0$ 的整数。这样的 $\mathcal{O}_X$-模层因而既有限型又拟凝聚。
若 $f:X\to Y$ 是赋环空间的一个态射，而 $\mathcal{G}$ 是一个有有限
表示的 $\mathcal{O}_Y$-模层，则由 (5.1.4) 的论证可知，
$f^*(\mathcal{G})$ 有有限表示。
\end{env}

\begin{env}[5.2.6]
\label{0.5.2.6-fr}
设 $\mathcal{F}$ 是一个有有限表示 (5.2.5) 的 $\mathcal{O}_X$-模层；
则对任意 $\mathcal{O}_X$-模层 $\mathcal{H}$，具有函子性的典范同态
\[
  \bigl(\mathcal{H}\!om_{\mathcal{O}_X}(\mathcal{F},\mathcal{H})\bigr)_x
  \longrightarrow
  \operatorname{Hom}_{\mathcal{O}_x}(\mathcal{F}_x,\mathcal{H}_x)
\]
是\emph{双射} (T, 4.1.1)。
\end{env}

\begin{env}[5.2.7]
\label{0.5.2.7-fr}
设 $\mathcal{F}$、$\mathcal{G}$ 是两个有有限表示的
$\mathcal{O}_X$-模层。若对某个 $x\in X$，$\mathcal{F}_x$ 与
$\mathcal{G}_x$ 是两个\emph{同构的} $\mathcal{O}_x$-模，则存在一个
开邻域 $U$，它含有 $x$，使得 $\mathcal{F}|U$ 与 $\mathcal{G}|U$
\emph{同构}。事实上，若
$\varphi:\mathcal{F}_x\to\mathcal{G}_x$ 与
$\psi:\mathcal{G}_x\to\mathcal{F}_x$ 是两个互逆同构，则由 (5.2.6)，
存在一个开邻域 $V$，它含有 $x$，以及两个截面 $u$（相应地，$v$），
前者属于 $\mathcal{H}\!om_{\mathcal{O}_X}(\mathcal{F},\mathcal{G})$
（后者属于 $\mathcal{H}\!om_{\mathcal{O}_X}(\mathcal{G},\mathcal{F})$），
它们定义在 $V$ 上，使得

\oldpage[0\textsubscript{I}]{47}

$u_x=\varphi$（相应地，$v_x=\psi$）。由于 $(u\circ v)_x$ 与
$(v\circ u)_x$ 都是恒等自同构，存在一个开邻域 $U\subset V$，它含有
$x$，使得 $(u\circ v)|U$ 与 $(v\circ u)|U$ 都是恒等
自同构，命题得证。
\end{env}

\subsection{凝聚层}
\label{subsection:0.5.3-fr}

\begin{env}[5.3.1]
\label{0.5.3.1-fr}
称一个 $\mathcal{O}_X$-模层 $\mathcal{F}$ 是\emph{凝聚的}，如果它满足
下列两个条件：
\begin{enumerate}
  \item[$a)$] $\mathcal{F}$ 是有限型的。
  \item[$b)$] 对任意开集 $U\subset X$、任意整数 $n>0$ 以及任意同态
    $u:\mathcal{O}_X^n|U\to\mathcal{F}|U$，$u$ 的核都是有限型的。
\end{enumerate}
应当注意，这两个条件都具有\emph{局部}性质。

下文所回顾的凝聚层性质，其大多数证明可参见 (FAC), I, 2。
\end{env}

\begin{env}[5.3.2]
\label{0.5.3.2-fr}
每个凝聚 $\mathcal{O}_X$-模层都有有限表示 (5.2.5)；逆命题未必成立，
因为 $\mathcal{O}_X$ 本身未必是一个凝聚 $\mathcal{O}_X$-模层。

凝聚 $\mathcal{O}_X$-模层的每个\emph{有限型} $\mathcal{O}_X$-子模层
都是凝聚的；凝聚 $\mathcal{O}_X$-模层的任意\emph{有限}直和仍是凝聚
$\mathcal{O}_X$-模层。
\end{env}

\begin{env}[5.3.3]
\label{0.5.3.3-fr}
若 $0\to\mathcal{F}\to\mathcal{G}\to\mathcal{H}\to0$ 是
$\mathcal{O}_X$-模层的一个正合列，而其中两个 $\mathcal{O}_X$-模层
凝聚，则第三个也凝聚。
\end{env}

\begin{env}[5.3.4]
\label{0.5.3.4-fr}
若 $\mathcal{F}$、$\mathcal{G}$ 是两个凝聚 $\mathcal{O}_X$-模层，而
$u:\mathcal{F}\to\mathcal{G}$ 是一个同态，则 $\operatorname{Im}(u)$、
$\operatorname{Ker}(u)$ 与 $\operatorname{Coker}(u)$ 都是凝聚
$\mathcal{O}_X$-模层。特别地，若 $\mathcal{F}$、$\mathcal{G}$ 是某个
凝聚 $\mathcal{O}_X$-模层的凝聚 $\mathcal{O}_X$-子模层，则
$\mathcal{F}+\mathcal{G}$ 与 $\mathcal{F}\cap\mathcal{G}$ 都凝聚。
\end{env}

\begin{env}[5.3.5]
\label{0.5.3.5-fr}
若 $\mathcal{F}$、$\mathcal{G}$ 是两个凝聚 $\mathcal{O}_X$-模层，则
$\mathcal{F}\otimes_{\mathcal{O}_X}\mathcal{G}$ 与
$\mathcal{H}\!om_{\mathcal{O}_X}(\mathcal{F},\mathcal{G})$ 也凝聚。
\end{env}

\begin{env}[5.3.6]
\label{0.5.3.6-fr}
设 $\mathcal{F}$ 是一个凝聚 $\mathcal{O}_X$-模层，$\mathcal{J}$ 是
$\mathcal{O}_X$ 的一个凝聚理想层。则 $\mathcal{O}_X$-模层
$\mathcal{J}\mathcal{F}$ 是凝聚的，因为它是
$\mathcal{J}\otimes_{\mathcal{O}_X}\mathcal{F}$ 在典范同态
$\mathcal{J}\otimes_{\mathcal{O}_X}\mathcal{F}\to\mathcal{F}$ 下的像
（5.3.4 与 5.3.5）。
\end{env}

\begin{env}[5.3.7]
\label{0.5.3.7-fr}
称一个 $\mathcal{O}_X$-代数层 $\mathcal{A}$ 是\emph{凝聚的}，如果
$\mathcal{A}$ 作为 $\mathcal{O}_X$-模层是凝聚的。特别地，
$\mathcal{O}_X$ 是一个\emph{凝聚环层}，当且仅当对任意开集
$U\subset X$ 以及任意形如
$u:\mathcal{O}_X^p|U\to\mathcal{O}_X|U$ 的同态，$u$ 的核都是一个
有限型 $(\mathcal{O}_X|U)$-模层。

若 $\mathcal{O}_X$ 是凝聚环层，则每个有有限表示 (5.2.5) 的
$\mathcal{O}_X$-模层 $\mathcal{F}$ 都凝聚，这是 (5.3.4) 的结果。

一个 $\mathcal{O}_X$-模层 $\mathcal{F}$ 的\emph{零化子}，是核
$\mathcal{J}$，即典范同态
$\mathcal{O}_X\to\mathcal{H}\!om_{\mathcal{O}_X}(\mathcal{F},\mathcal{F})$
的核；
这个同态把任意截面 $s\in\Gamma(U,\mathcal{O}_X)$ 映为乘法映射 $s$，
它属于 $\operatorname{Hom}(\mathcal{F}|U,\mathcal{F}|U)$。
若 $\mathcal{O}_X$ 凝聚，且 $\mathcal{F}$ 是凝聚
$\mathcal{O}_X$-模层，则 $\mathcal{J}$ 凝聚（5.3.4 与 5.3.5）；
对每个 $x\in X$，$\mathcal{J}_x$ 是 $\mathcal{F}_x$ 的零化子 (5.2.6)。
\end{env}

\oldpage[0\textsubscript{I}]{48}

\begin{env}[5.3.8]
\label{0.5.3.8-fr}
设 $\mathcal{O}_X$ 凝聚；设 $\mathcal{F}$ 是一个凝聚
$\mathcal{O}_X$-模层，$x$ 是 $X$ 的一点，$M$ 是 $\mathcal{F}_x$ 的
一个有限型子模。则存在一个开邻域 $U$，它含有 $x$，以及一个凝聚
$(\mathcal{O}_X|U)$-模层 $\mathcal{G}$，它是 $\mathcal{F}|U$ 的子模层，
并使得 $\mathcal{G}_x=M$（T, 4.1, lemme 1）。

把这个结果同凝聚 $\mathcal{O}_X$-模层的 $\mathcal{O}_X$-子模层的
性质结合起来，便得到环 $\mathcal{O}_x$ 为使 $\mathcal{O}_X$ 凝聚所
必须满足的条件。例如，由 (5.3.4)，$\mathcal{O}_x$ 的两个有限型理想
之交仍须是有限型理想。
\end{env}

\begin{env}[5.3.9]
\label{0.5.3.9-fr}
设 $\mathcal{O}_X$ 凝聚，而 $M$ 是一个有有限表示的
$\mathcal{O}_x$-模，因而同构于一个同态
$\varphi:\mathcal{O}_x^p\to\mathcal{O}_x^q$ 的余核；则存在
% EGA1-ERR-0006: French print says an open neighbourhood U of X; the following sentence and construction require a neighbourhood of x. Canonical translation preserves X here.
一个开邻域 $U$，它是 $X$ 的邻域，以及一个凝聚 $(\mathcal{O}_X|U)$-模层
$\mathcal{F}$，使 $\mathcal{F}_x$ 同构于 $M$。事实上，由 (5.2.6)，
存在一个截面 $u$，它属于
$\mathcal{H}\!om_{\mathcal{O}_X}(\mathcal{O}_X^p,\mathcal{O}_X^q)$，
定义在一个开邻域 $U$ 上，该邻域
含有 $x$，并使得 $u_x=\varphi$；余核 $\mathcal{F}$，即同态
$u:\mathcal{O}_X^p|U\to\mathcal{O}_X^q|U$ 的余核，
符合要求 (5.3.4)。
\end{env}

\begin{env}[5.3.10]
\label{0.5.3.10-fr}
设 $\mathcal{O}_X$ 凝聚，而 $\mathcal{J}$ 是 $\mathcal{O}_X$ 的一个
凝聚理想层。一个 $(\mathcal{O}_X/\mathcal{J})$-模层 $\mathcal{F}$
凝聚，当且仅当它作为 $\mathcal{O}_X$-模层凝聚。特别地，
$\mathcal{O}_X/\mathcal{J}$ 是凝聚环层。
\end{env}

\begin{env}[5.3.11]
\label{0.5.3.11-fr}
设 $f:X\to Y$ 是赋环空间的一个态射，并设 $\mathcal{O}_X$ 凝聚；
则对任意凝聚 $\mathcal{O}_Y$-模层 $\mathcal{G}$，
$f^*(\mathcal{G})$ 都是凝聚 $\mathcal{O}_X$-模层。事实上，采用
(5.2.4) 的记号，可以假设 $\mathcal{G}|V$ 是同态
$v:\mathcal{O}_Y^q|V\to\mathcal{O}_Y^p|V$ 的余核；由于 $f_U^*$
右正合，$f^*(\mathcal{G})|U=f_U^*(\mathcal{G}|V)$ 是同态
$f_U^*(v):\mathcal{O}_X^q|U\to\mathcal{O}_X^p|U$ 的余核，故结论成立。
\end{env}

\begin{env}[5.3.12]
\label{0.5.3.12-fr}
设 $Y$ 是 $X$ 的一个闭子集，$j:Y\to X$ 是典范嵌入，
$\mathcal{O}_Y$ 是 $Y$ 上的一个环层，并置
$\mathcal{O}_X=j_*(\mathcal{O}_Y)$。一个 $\mathcal{O}_Y$-模层
$\mathcal{G}$ 是有限型的（相应地，拟凝聚的、凝聚的），当且仅当
$j_*(\mathcal{G})$ 是有限型 $\mathcal{O}_X$-模层（相应地，拟凝聚的、
凝聚的）。
\end{env}

\subsection{局部自由层}
\label{subsection:0.5.4-fr}

\begin{env}[5.4.1]
\label{0.5.4.1-fr}
设 $X$ 是一个赋环空间。称一个 $\mathcal{O}_X$-模层 $\mathcal{F}$
\emph{局部自由}，如果对每个 $x\in X$，都存在一个开邻域 $U$，它含有
$x$，使得 $\mathcal{F}|U$ 同构于一个 $(\mathcal{O}_X|U)$-模层，
而该模层形如 $\mathcal{O}_X^{(I)}|U$，其中 $I$ 可以依赖于
$U$。若对每个 $U$，$I$ 都有限，则称 $\mathcal{F}$ \emph{秩有限}；若
对每个 $U$，$I$ 都有相同的有限元素个数 $n$，则称 $\mathcal{F}$
\emph{秩为 $n$}。秩为 1 的局部自由 $\mathcal{O}_X$-模层也称为
\emph{可逆的}（参见 (5.4.3)）。若 $\mathcal{F}$ 是秩有限的局部自由
$\mathcal{O}_X$-模层，则对每个 $x\in X$，$\mathcal{F}_x$ 是一个有限秩
自由 $\mathcal{O}_x$-模，其秩为 $n(x)$，而且存在一个邻域 $U$，它含有 $x$，
使 $\mathcal{F}|U$ 的秩为 $n(x)$；若 $X$ 连通，则 $n(x)$ 因而是
\emph{常数}。

显然，每个局部自由层都拟凝聚；若 $\mathcal{O}_X$ 是凝聚环层，则每个
秩有限的局部自由 $\mathcal{O}_X$-模层都凝聚。

若 $\mathcal{L}$ 局部自由，则
$\mathcal{L}\otimes_{\mathcal{O}_X}\mathcal{F}$ 是模层范畴中关于
$\mathcal{F}$ 的一个\emph{正合}函子，该模层范畴由
$\mathcal{O}_X$-模层组成。

我们主要考虑秩有限的局部自由 $\mathcal{O}_X$-模层，

\oldpage[0\textsubscript{I}]{49}

以后若只说局部自由层而不加限定，均默认它们\emph{秩有限}。
\end{env}

\begin{env}[5.4.2]
\label{0.5.4.2-fr}
若 $\mathcal{L}$、$\mathcal{F}$ 是两个 $\mathcal{O}_X$-模层，则有一个
具有函子性的典范同态
\[
  \check{\mathcal{L}}\otimes_{\mathcal{O}_X}\mathcal{F}
  =\mathcal{H}\!om_{\mathcal{O}_X}(\mathcal{L},\mathcal{O}_X)
   \otimes_{\mathcal{O}_X}\mathcal{F}
  \longrightarrow
  \mathcal{H}\!om_{\mathcal{O}_X}(\mathcal{L},\mathcal{F})
  \tag{5.4.2.1}\label{0.5.4.2.1-fr}
\]
定义如下：对任意开集 $U$，给任意一对 $(u,t)$，其中
\[
  u\in\Gamma\bigl(U,
    \mathcal{H}\!om_{\mathcal{O}_X}(\mathcal{L},\mathcal{O}_X)\bigr)
  =\operatorname{Hom}(\mathcal{L}|U,\mathcal{O}_X|U)
  \quad\text{且}\quad
  t\in\Gamma(U,\mathcal{F}),
\]
对应 $\operatorname{Hom}(\mathcal{L}|U,\mathcal{F}|U)$ 中的一个元素；
对每个 $x\in U$，这个元素把 $s_x\in\mathcal{L}_x$ 映为元素
$u_x(s_x)t_x$，它属于 $\mathcal{F}_x$。若 $\mathcal{L}$
\emph{局部自由且秩有限}，则这个同态是\emph{双射}；由于这个性质是
局部的，可以约化到 $\mathcal{L}=\mathcal{O}_X^n$ 的情形。又因为对
任意 $\mathcal{O}_X$-模层 $\mathcal{G}$，
$\mathcal{H}\!om_{\mathcal{O}_X}(\mathcal{O}_X^n,\mathcal{G})$ 典范同构于
$\mathcal{G}^n$，故问题归结为 $\mathcal{L}=\mathcal{O}_X$ 的情形，
这时结论立即成立。
\end{env}

\begin{env}[5.4.3]
\label{0.5.4.3-fr}
若 $\mathcal{L}$ 可逆，则它的对偶
$\check{\mathcal{L}}=
\mathcal{H}\!om_{\mathcal{O}_X}(\mathcal{L},\mathcal{O}_X)$ 也可逆；
事实上，由于问题是局部的，立即约化到 $\mathcal{L}=\mathcal{O}_X$
的情形。此外，有一个典范同构
\[
  \mathcal{H}\!om_{\mathcal{O}_X}(\mathcal{L},\mathcal{O}_X)
  \otimes_{\mathcal{O}_X}\mathcal{L}
  \xrightarrow{\sim}
  \mathcal{O}_X
  \tag{5.4.3.1}\label{0.5.4.3.1-fr}
\]
事实上，由 (5.4.2)，只须定义一个典范同构
$\mathcal{H}\!om_{\mathcal{O}_X}(\mathcal{L},\mathcal{L})
\xrightarrow{\sim}\mathcal{O}_X$。而对\emph{任意}
$\mathcal{O}_X$-模层 $\mathcal{F}$，都有一个典范同态
% EGA1-ERR-0007: French print marks the canonical map as an isomorphism before proving bijectivity in the invertible case. Canonical translation preserves the printed arrow.
$\mathcal{O}_X\xrightarrow{\sim}
\mathcal{H}\!om_{\mathcal{O}_X}(\mathcal{F},\mathcal{F})$ (5.3.7)。
还需证明，当 $\mathcal{F}=\mathcal{L}$ 可逆时，这个同态是双射；由于
问题是局部的，问题归结为 $\mathcal{L}=\mathcal{O}_X$ 的情形，
这时结论立即成立。

根据以上结果，置
$\mathcal{L}^{-1}=
\mathcal{H}\!om_{\mathcal{O}_X}(\mathcal{L},\mathcal{O}_X)$，并称
$\mathcal{L}^{-1}$ 为 $\mathcal{L}$ 的\emph{逆}。当 $X$ 只有一点，且
$\mathcal{O}_X$ 是一个\emph{局部}环 $A$，其极大理想为 $\mathfrak m$
时，“可逆层”这一术语可作如下解释。若 $M$、$M'$ 是两个 $A$-模
（$M$ 有限型），并且 $M\otimes_A M'$ 同构于 $A$，则
$(A/\mathfrak m)\otimes_A(M\otimes_A M')$ 可与
$(M/\mathfrak mM)\otimes_{A/\mathfrak m}(M'/\mathfrak mM')$ 等同。
后一个张量积是域 $A/\mathfrak m$ 上向量空间的张量积，并同构于
$A/\mathfrak m$；这迫使 $M/\mathfrak mM$ 与 $M'/\mathfrak mM'$ 都是
一维的。对任意元素 $z\in M$，只要它不属于 $\mathfrak mM$，就有
$M=Az+\mathfrak mM$；由 Nakayama 引理即得 $M=Az$，这里 $M$
有限型。此外，由于 $z$ 的零化子零化
$M\otimes_A M'$，而后者同构于 $A$，这个零化子就是 $\{0\}$，因此
$M$ \emph{同构于 $A$}。在一般情形下，这表明：若 $\mathcal{L}$ 是一个
有限型 $\mathcal{O}_X$-模层，且存在一个 $\mathcal{O}_X$-模层
$\mathcal{F}$，使 $\mathcal{L}\otimes_{\mathcal{O}_X}\mathcal{F}$
同构于 $\mathcal{O}_X$，又若各环 $\mathcal{O}_x$ 都是局部环，则
$\mathcal{L}_x$ 都是一个 $\mathcal{O}_x$-模，同构于 $\mathcal{O}_x$
（对每个 $x\in X$）。若还假设 $\mathcal{O}_X$ 与 $\mathcal{L}$
\emph{凝聚}，则由 (5.2.7) 可知 $\mathcal{L}$ 可逆。
\end{env}

\begin{env}[5.4.4]
\label{0.5.4.4-fr}
若 $\mathcal{L}$、$\mathcal{L}'$ 是两个可逆 $\mathcal{O}_X$-模层，则
$\mathcal{L}\otimes_{\mathcal{O}_X}\mathcal{L}'$ 也可逆；事实上，
由于问题是局部的，可以假设 $\mathcal{L}=\mathcal{O}_X$，这时结论显然。
对每个整数 $n\geq1$，以 $\mathcal{L}^{\otimes n}$ 表示 $n$ 个都等于

\oldpage[0\textsubscript{I}]{50}

$\mathcal{L}$ 的层的张量积；约定
$\mathcal{L}^{\otimes0}=\mathcal{O}_X$，并对 $n\geq1$ 置
$\mathcal{L}^{\otimes(-n)}=(\mathcal{L}^{-1})^{\otimes n}$。采用这些记号，
便对任意有理整数 $m,n$ 有一个\emph{具有函子性的典范同构}
\[
  \mathcal{L}^{\otimes m}\otimes_{\mathcal{O}_X}\mathcal{L}^{\otimes n}
  \xrightarrow{\sim}
  \mathcal{L}^{\otimes(n+m)}
  \tag{5.4.4.1}\label{0.5.4.4.1-fr}
\]
事实上，根据定义，立即约化到 $m=-1$、$n=1$ 的情形，而所论同构已在
(5.4.3) 中定义。
\end{env}

\begin{env}[5.4.5]
\label{0.5.4.5-fr}
设 $f:Y\to X$ 是赋环空间的一个态射。若 $\mathcal{L}$ 是局部自由的
（相应地，可逆的）$\mathcal{O}_X$-模层，则 $f^*(\mathcal{L})$ 是局部
自由的（相应地，可逆的）$\mathcal{O}_Y$-模层。这立即来自如下事实：
两个局部同构 $\mathcal{O}_X$-模层的逆像局部同构；$f^*$ 与有限直和
交换；并且 $f^*(\mathcal{O}_X)=\mathcal{O}_Y$ (4.3.4)。此外，已知有
一个具有函子性的典范同态
\[
  f^*(\check{\mathcal{L}})
  \longrightarrow
  \check{\bigl(f^*(\mathcal{L})\bigr)}
\]
(4.4.6)；当
% EGA1-ERR-0008: French print uses bare L here although the surrounding sheaf is mathcal L. Canonical translation preserves the printed symbol.
$L$ 局部自由时，这个同态是\emph{双射}：事实上，问题又归结为
$\mathcal{L}=\mathcal{O}_X$ 这个显然的情形。由此可知，若
$\mathcal{L}$ 可逆，则 $f^*(\mathcal{L}^{\otimes n})$ 可典范地与
$(f^*(\mathcal{L}))^{\otimes n}$ 等同，对任意有理整数 $n$ 都如此。
\end{env}

\begin{env}[5.4.6]
\label{0.5.4.6-fr}
设 $\mathcal{L}$ 是一个可逆 $\mathcal{O}_X$-模层；以
$\Gamma_*(X,\mathcal{L})$ 或简记 $\Gamma_*(\mathcal{L})$ 表示直和
Abel 群 $\bigoplus_{n\in\mathbb{Z}}\Gamma(X,\mathcal{L}^{\otimes n})$。
如下赋予它一个\emph{分次环}结构：对一对 $(s_n,s_m)$，其中
$s_n\in\Gamma(X,\mathcal{L}^{\otimes n})$、
$s_m\in\Gamma(X,\mathcal{L}^{\otimes m})$，令它对应于
$\mathcal{L}^{\otimes(n+m)}$ 在 $X$ 上的截面；这个截面通过 (5.4.4.1)
典范地对应于截面 $s_n\otimes s_m$，后者属于
$\mathcal{L}^{\otimes n}\otimes_{\mathcal{O}_X}\mathcal{L}^{\otimes m}$。
这个乘法的结合性立即可验。显然，
$\Gamma_*(X,\mathcal{L})$ 关于 $\mathcal{L}$ 是一个取值于分次环范畴的
协变函子。

现在若 $\mathcal{F}$ 是任意 $\mathcal{O}_X$-模层，置
\[
  \Gamma_*(\mathcal{L},\mathcal{F})
  =\bigoplus_{n\in\mathbb{Z}}
   \Gamma\bigl(X,
     \mathcal{F}\otimes_{\mathcal{O}_X}\mathcal{L}^{\otimes n}\bigr).
\]
如下赋予这个 Abel 群一个分次环 $\Gamma_*(\mathcal{L})$ 上的
\emph{分次模}结构：对一对 $(s_n,u_m)$，其中
$s_n\in\Gamma(X,\mathcal{L}^{\otimes n})$ 且
$u_m\in\Gamma(X,
\mathcal{F}\otimes_{\mathcal{O}_X}\mathcal{L}^{\otimes m})$，令它对应于
$\mathcal{F}\otimes_{\mathcal{O}_X}\mathcal{L}^{\otimes(m+n)}$ 的截面；
这个截面通过 (5.4.4.1) 典范地对应于 $s_n\otimes u_m$。模公理立即
可验。当 $X$、$\mathcal{L}$ 固定时，
$\Gamma_*(\mathcal{L},\mathcal{F})$ 关于 $\mathcal{F}$ 是一个取值于
分次 $\Gamma_*(\mathcal{L})$-模范畴的协变函子；当 $X$、$\mathcal{F}$
固定时，它关于 $\mathcal{L}$ 是一个取值于 Abel 群范畴的协变函子。

若 $f:Y\to X$ 是赋环空间的一个态射，则典范同态 (4.4.3.2)
\[
  \rho:\mathcal{L}^{\otimes n}
  \longrightarrow
  f_*\bigl(f^*(\mathcal{L}^{\otimes n})\bigr)
\]
定义一个 Abel 群同态
\[
  \Gamma(X,\mathcal{L}^{\otimes n})
  \longrightarrow
  \Gamma\bigl(Y,f^*(\mathcal{L}^{\otimes n})\bigr),
\]
而由于 $f^*(\mathcal{L}^{\otimes n})=(f^*(\mathcal{L}))^{\otimes n}$，
由典范同态 (4.4.3.2) 与 (5.4.4.1) 的定义可知，以上同态定义了一个
\emph{具有函子性的分次环同态}
$\Gamma_*(\mathcal{L})\to\Gamma_*(f^*(\mathcal{L}))$。同一个典范同态
(4.4.3) 也定义一个 Abel 群同态
\[
  \Gamma\bigl(X,
    \mathcal{F}\otimes_{\mathcal{O}_X}\mathcal{L}^{\otimes n}\bigr)
  \longrightarrow
  \Gamma\bigl(Y,
    f^*(\mathcal{F}\otimes_{\mathcal{O}_X}\mathcal{L}^{\otimes n})\bigr),
\]
又因为
\[
  f^*(\mathcal{F}\otimes_{\mathcal{O}_X}\mathcal{L}^{\otimes n})
  =f^*(\mathcal{F})\otimes_{\mathcal{O}_Y}
   (f^*(\mathcal{L}))^{\otimes n}
  \tag{4.3.3.1}
\]

\oldpage[0\textsubscript{I}]{51}

这些同态（$n$ 取遍各值）定义一个\emph{分次模双同态}
\[
  \Gamma_*(\mathcal{L},\mathcal{F})
  \longrightarrow
  \Gamma_*(f^*(\mathcal{L}),f^*(\mathcal{F})).
\]
\end{env}

\begin{env}[5.4.7]
\label{0.5.4.7-fr}
可以证明，存在一个\emph{集合} $\mathfrak{M}$（也记作
$\mathfrak{M}(X)$），它由可逆 $\mathcal{O}_X$-模层组成，使得每个可逆
$\mathcal{O}_X$-模层都恰与 $\mathfrak{M}$ 中的一个元素同构
\footnote{见《引言》中所引的在编著作。}。在 $\mathfrak{M}$ 上如下定义
合成律：令两个元素 $\mathcal{L}$、$\mathcal{L}'$ 属于
$\mathfrak{M}$，并令它们对应于 $\mathfrak{M}$ 中唯一一个同构于
$\mathcal{L}\otimes_{\mathcal{O}_X}\mathcal{L}'$ 的元素。在这个合成律下，
$\mathfrak{M}$ 是一个\emph{同构于上同调群}
$H^1(X,\mathcal{O}_X^*)$ 的群；这里 $\mathcal{O}_X^*$ 是
$\mathcal{O}_X$ 的子层，使得 $\Gamma(U,\mathcal{O}_X^*)$ 是环
$\Gamma(U,\mathcal{O}_X)$ 的可逆元群，对每个开集 $U\subset X$ 都如此
（所以 $\mathcal{O}_X^*$ 是
\emph{乘法} Abel 群层）。

为此应当注意，对任意开集 $U\subset X$，截面群
$\Gamma(U,\mathcal{O}_X^*)$ 典范地等同于 $(\mathcal{O}_X|U)$-模层
$\mathcal{O}_X|U$ 的\emph{自同构群}。这个等同把截面 $\varepsilon$
（属于 $\mathcal{O}_X^*$ 且定义在 $U$ 上）对应到自同构 $u$，它作用于
$\mathcal{O}_X|U$，使得
% EGA1-ERR-0010: French print quantifies x over X although the sections and automorphism are restricted to U. Canonical translation preserves X.
$u_x(s_x)=\varepsilon_xs_x$ 对每个 $x\in X$ 以及每个
$s_x\in\mathcal{O}_x$ 成立。现在设 $\mathfrak{U}=(U_\lambda)$ 是 $X$ 的
一个开覆盖；对每一对指标 $(\lambda,\mu)$ 给定一个自同构
$\theta_{\lambda\mu}$，它作用于
$\mathcal{O}_X|(U_\lambda\cap U_\mu)$，等价于给定一个
$1$-\emph{上链}，它属于覆盖 $\mathfrak{U}$ 并取值于
$\mathcal{O}_X^*$；说各 $\theta_{\lambda\mu}$ 满足
黏合条件 (3.3.1)，就意味着相应的上链是一个\emph{上循环}。同样，
对每个 $\lambda$ 给定一个自同构 $\omega_\lambda$，它作用于
$\mathcal{O}_X|U_\lambda$，等价于给定一个 $0$-上链，它属于覆盖
$\mathfrak{U}$ 并取值于 $\mathcal{O}_X^*$；它的\emph{上边界}对应于自同构族
$(\omega_\lambda|U_\lambda\cap U_\mu)\circ
 (\omega_\mu|U_\lambda\cap U_\mu)^{-1}$。

对每个 $1$-上循环（属于 $\mathfrak{U}$ 并取值于 $\mathcal{O}_X^*$），
可令它对应于 $\mathfrak{M}$ 中的一个元素：这个元素同构于一个可逆
$\mathcal{O}_X$-模层，而该模层由与这个上循环对应的自同构族
$(\theta_{\lambda\mu})$ 黏合而得；两个同属一个上同调类的上循环对应于
$\mathfrak{M}$ 中两个相等的元素 (3.3.2)。换言之，由此定义了一个映射
$\varphi_{\mathfrak{U}}:H^1(\mathfrak{U},\mathcal{O}_X^*)\to\mathfrak{M}$。
此外，若 $\mathfrak{B}$ 是 $X$ 的第二个开覆盖，且
比 $\mathfrak{U}$ 更细，则图
\[
\begin{tikzcd}[column sep=huge,row sep=large]
H^1(\mathfrak{U},\mathcal{O}_X^*)
  \arrow[dd]
  \arrow[rrd,"\varphi_{\mathfrak{U}}"] & & \\
& & \mathfrak{M} \\
H^1(\mathfrak{B},\mathcal{O}_X^*)
  \arrow[rru,swap,"\varphi_{\mathfrak{B}}"] & &
\end{tikzcd}
\]
交换，其中竖直箭头是典范同态 (G, II, 5.7)；这由 (3.3.3) 即得。
取归纳极限，便得到映射 $H^1(X,\mathcal{O}_X^*)\to\mathfrak{M}$，
因为已知 \v{C}ech 上同调群 $\check{H}^1(X,\mathcal{O}_X^*)$
可与第一上同调群 $H^1(X,\mathcal{O}_X^*)$ 等同
（G, II, 5.9, cor. du th.~5.9.1）。这个映射是\emph{满射}：事实上，
按照定义，对每个可逆 $\mathcal{O}_X$-模层 $\mathcal{L}$，存在一个
开覆盖 $(U_\lambda)$，它覆盖 $X$，使得 $\mathcal{L}$ 由各层
$\mathcal{O}_X|U_\lambda$ 黏合而得 (3.3.1)。它也是\emph{单射}，因为
只须对映射
% EGA1-ERR-0009: French print omits the star on O_X in this H^1 source, although the entire construction is valued in O_X^*. Canonical translation preserves the printed formula.
$H^1(\mathfrak{U},\mathcal{O}_X)\to\mathfrak{M}$ 证明这一点，而这由
(3.3.2) 即得。还需证明

\oldpage[0\textsubscript{I}]{52}

这样定义的双射是一个群同态。事实上，给定两个可逆
$\mathcal{O}_X$-模层 $\mathcal{L}$、$\mathcal{L}'$，存在一个开覆盖
$(U_\lambda)$，使 $\mathcal{L}|U_\lambda$ 与
$\mathcal{L}'|U_\lambda$ 都同构于 $\mathcal{O}_X|U_\lambda$，对每个
$\lambda$ 都如此。因此，对每个指标 $\lambda$，存在元素
$a_\lambda$（相应地，$a'_\lambda$），
它属于 $\Gamma(U_\lambda,\mathcal{L})$（相应地，属于
$\Gamma(U_\lambda,\mathcal{L}')$），使得
$\Gamma(U_\lambda,\mathcal{L})$（相应地，
$\Gamma(U_\lambda,\mathcal{L}')$）的元素都是
$s_\lambda\cdot a_\lambda$（相应地，$s_\lambda\cdot a'_\lambda$），
其中 $s_\lambda$ 遍历 $\Gamma(U_\lambda,\mathcal{O}_X)$。对应的上循环
$(\varepsilon_{\lambda\mu})$、$(\varepsilon'_{\lambda\mu})$ 具有如下
性质：$s_\lambda\cdot a_\lambda=s_\mu\cdot a_\mu$（相应地，
$s_\lambda\cdot a'_\lambda=s_\mu\cdot a'_\mu$）在
$U_\lambda\cap U_\mu$ 上成立，等价于
$s_\lambda=\varepsilon_{\lambda\mu}s_\mu$（相应地，
$s_\lambda=\varepsilon'_{\lambda\mu}s_\mu$）在
$U_\lambda\cap U_\mu$ 上成立。由于
$\mathcal{L}\otimes_{\mathcal{O}_X}\mathcal{L}'$ 在 $U_\lambda$ 上的
截面是各式
$s_\lambda s'_\lambda\cdot(a_\lambda\otimes a'_\lambda)$ 的有限和，
其中 $s_\lambda$、
$s'_\lambda$ 遍历 $\Gamma(U_\lambda,\mathcal{O}_X)$，显然，上循环
$(\varepsilon_{\lambda\mu}\varepsilon'_{\lambda\mu})$ 对应于
$\mathcal{L}\otimes_{\mathcal{O}_X}\mathcal{L}'$，证明完毕
\footnote{这个结果的一般形式见第 51 页注释所引的著作。}。
\end{env}

\begin{env}[5.4.8]
\label{0.5.4.8-fr}
设 $f=(\psi,\omega)$ 是赋环空间的一个态射 $Y\to X$。函子
$f^*(\mathcal{L})$ 在自由 $\mathcal{O}_X$-模层范畴中定义了一个映射
（仍滥用记号写作 $f^*$），它从集合 $\mathfrak{M}(X)$ 映到集合
$\mathfrak{M}(Y)$。另一方面，有一个典范同态 (T, 3.2.2)
\[
  H^1(X,\mathcal{O}_X^*)\longrightarrow H^1(Y,\mathcal{O}_Y^*).
  \tag{5.4.8.1}\label{0.5.4.8.1-fr}
\]
当典范地等同 (5.4.7) $\mathfrak{M}(X)$ 与
$H^1(X,\mathcal{O}_X^*)$（相应地，等同 $\mathfrak{M}(Y)$ 与
$H^1(Y,\mathcal{O}_Y^*)$）时，同态 (5.4.8.1)
\emph{就等同于映射 $f^*$}。事实上，若 $\mathcal{L}$ 来自一个上循环
$(\varepsilon_{\lambda\mu})$，后者对应于一个开覆盖 $(U_\lambda)$，
它覆盖 $X$，则只须证明 $f^*(\mathcal{L})$ 来自一个上循环，且这个
上循环的上同调类是 $(\varepsilon_{\lambda\mu})$ 的上同调类在
(5.4.8.1) 下的像。若 $\theta_{\lambda\mu}$ 是
$\mathcal{O}_X|(U_\lambda\cap U_\mu)$ 的自同构，并对应于
$\varepsilon_{\lambda\mu}$，那么显然，$f^*(\mathcal{L})$ 是将各
$\mathcal{O}_Y|\psi^{-1}(U_\lambda)$ 借助自同构
$f^*(\theta_{\lambda\mu})$ 黏合而得。因此只须验证，这些自同构
对应于上循环 $(\omega^\#(\varepsilon_{\lambda\mu}))$；这由定义立即
得到（这里把 $\varepsilon_{\lambda\mu}$ 等同于它通过 $\rho$
(3.7.2) 得到的典范像；该像是 $\psi^*(\mathcal{O}_X^*)$ 在
$\psi^{-1}(U_\lambda\cap U_\mu)$ 上的截面）。
\end{env}

\begin{env}[5.4.9]
\label{0.5.4.9-fr}
设 $\mathcal{E}$、$\mathcal{F}$ 是两个 $\mathcal{O}_X$-模层，并假设
$\mathcal{F}$ \emph{局部自由}；再设 $\mathcal{G}$ 是一个
$\mathcal{O}_X$-模层，并且是一个
\emph{$\mathcal{F}$ 被 $\mathcal{E}$ 的扩张}，也就是说，存在一个正合列
\[
  0\longrightarrow\mathcal{E}\xrightarrow{i}\mathcal{G}
  \xrightarrow{p}\mathcal{F}\longrightarrow 0.
\]
那么，对每个 $x\in X$，都存在一个开邻域 $U$，它是 $x$ 的邻域，使
$\mathcal{G}|U$ 同构于\emph{直和}
$\mathcal{E}|U\oplus\mathcal{F}|U$。事实上，可以归结到
$\mathcal{F}=\mathcal{O}_X^n$ 的情形。设 $e_i$（$1\leq i\leq n$）
是 $\mathcal{O}_X^n$ 的典范截面 (5.5.5)；于是存在一个开邻域
$U$（它是 $x$ 的邻域），以及 $n$ 个截面 $s_i$，它们是
$\mathcal{G}$ 在 $U$ 上的截面，使得
$p(s_i|U)=e_i|U$，其中 $1\leq i\leq n$。现在令 $f$ 为同态
$\mathcal{F}|U\to\mathcal{G}|U$；它由截面 $s_i|U$ 定义 (5.1.1)。
立即可见，对任意开集 $V\subset U$ 及任意截面
$s\in\Gamma(V,\mathcal{G})$，都有
$s-f(p(s))\in\Gamma(V,\mathcal{E})$；所断言的结论由此成立。
\end{env}

\begin{env}[5.4.10]
\label{0.5.4.10-fr}
设 $f:X\to Y$ 是赋环空间的一个态射，$\mathcal{F}$ 是一个
$\mathcal{O}_X$-模层，$\mathcal{L}$ 是一个有限秩局部自由
$\mathcal{O}_Y$-模层。那么存在一个典范同构
\[
  f_*(\mathcal{F})\otimes_{\mathcal{O}_Y}\mathcal{L}
  \xrightarrow{\sim}
  f_*\bigl(\mathcal{F}\otimes_{\mathcal{O}_X}f^*(\mathcal{L})\bigr).
  \tag{5.4.10.1}\label{0.5.4.10.1-fr}
\]

\oldpage[0\textsubscript{I}]{53}

事实上，对任意 $\mathcal{O}_Y$-模层 $\mathcal{L}$，都有典范同态
\[
  f_*(\mathcal{F})\otimes_{\mathcal{O}_Y}\mathcal{L}
  \xrightarrow{1\otimes\rho}
  f_*(\mathcal{F})\otimes_{\mathcal{O}_Y}f_*(f^*(\mathcal{L}))
  \xrightarrow{\alpha}
  f_*\bigl(\mathcal{F}\otimes_{\mathcal{O}_X}f^*(\mathcal{L})\bigr),
\]
其中 $\rho$ 是同态 (4.4.3.2)，而 $\alpha$ 是同态 (4.2.2.1)。
为证明当 $\mathcal{L}$ 局部自由时这个同态是双射，只须注意问题关于
$Y$ 是局部的，从而考虑 $\mathcal{L}=\mathcal{O}_Y^n$ 的情形。此外，
$f_*$ 与 $f^*$ 都与有限直和可交换，故可以假设 $n=1$；此时命题立即
由定义以及关系 $f^*(\mathcal{O}_Y)=\mathcal{O}_X$ 得到。
\end{env}

\subsection{局部赋环空间上的层}
\label{subsection:0.5.5-fr}

\begin{env}[5.5.1]
\label{0.5.5.1-fr}
我们称赋环空间 $(X,\mathcal{O}_X)$ 是一个\emph{局部赋环空间}，如果对每个
$x\in X$，$\mathcal{O}_x$ 都是一个局部环；在本书所考虑的赋环空间中，
这类空间远为常见。此时，我们以 $\mathfrak{m}_x$ 表示
$\mathcal{O}_x$ 的\emph{极大理想}，以 $k(x)$ 表示\emph{剩余域}
$\mathcal{O}_x/\mathfrak{m}_x$。对任意 $\mathcal{O}_X$-模层
$\mathcal{F}$、任意开集 $U$（属于 $X$）、任意点 $x\in U$，以及任意截面
$f\in\Gamma(U,\mathcal{F})$，我们以 $f(x)$ 表示茎
$f_x\in\mathcal{F}_x\pmod{\mathfrak{m}_x\mathcal{F}_x}$ 的\emph{类}，
并称它为 $f$ 在点 $x$ 的\emph{值}。关系 $f(x)=0$ 因而意味着
$f_x\in\mathfrak{m}_x\mathcal{F}_x$；当这个关系成立时，我们（滥用语言）
称 $f$ \emph{在 $x$ 处为零}。应当注意，不要把这个关系与
$f_x=0$ 混淆。
\end{env}

\begin{env}[5.5.2]
\label{0.5.5.2-fr}
设 $X$ 是一个局部赋环空间，$\mathcal{L}$ 是一个可逆
$\mathcal{O}_X$-模层，而 $f$ 是 $\mathcal{L}$ 在 $X$ 上的一个截面。
那么，在一点 $x\in X$ 处，以下三个性质彼此\emph{等价}：
\begin{enumerate}
  \item[$a)$] \emph{$f_x$ 是 $\mathcal{L}_x$ 的一个生成元}；
  \item[$b)$] \emph{$f_x\notin\mathfrak{m}_x\mathcal{L}_x$}（也就是说，
    $f(x)\neq0$）；
  \item[$c)$] \emph{存在一个截面 $g$，它是 $\mathcal{L}^{-1}$ 在一个
    开邻域 $V$ 上的截面，该邻域是 $x$ 的邻域，并使 $f\otimes g$ 在
    $\Gamma(V,\mathcal{O}_X)$ 中的典范像 (5.4.3) 是单位截面。}
\end{enumerate}

事实上，问题是局部的，故可以归结到 $\mathcal{L}=\mathcal{O}_X$ 的
情形；此时 $a)$ 与 $b)$ 的等价性显然，并且 $c)$ 蕴含 $b)$ 也是显然的。
另一方面，如果 $f_x\notin\mathfrak{m}_x$，则 $f_x$ 在
$\mathcal{O}_x$ 中可逆；设 $f_xg_x=1_x$。按照截面之茎的定义，这意味着
存在一个邻域 $V$（它是 $x$ 的邻域），以及一个截面 $g$，它是
$\mathcal{O}_X$ 在 $V$ 上的截面，使得 $fg=1$ 在 $V$ 上成立；从而
$c)$ 成立。

由条件 $c)$ 立即可知，集合 $X_f$（其元素是点 $x$，并且这些点满足等价
条件 $a)$、$b)$、$c)$）在 $X$ 中是\emph{开集}；按照 (5.5.1) 中引入的
术语，这就是各点 $x$ 所成的集合，而 $f$ \emph{在这些点处不为零}。
\end{env}

\begin{env}[5.5.3]
\label{0.5.5.3-fr}
在 (5.5.2) 的假设下，设 $\mathcal{L}'$ 是第二个可逆
$\mathcal{O}_X$-模层；若 $f\in\Gamma(X,\mathcal{L})$ 且
$g\in\Gamma(X,\mathcal{L}')$，则
\[
  X_f\cap X_g=X_{f\otimes g}.
\]
事实上，立即可以归结到
$\mathcal{L}=\mathcal{L}'=\mathcal{O}_X$ 的情形（因为问题是局部的）；
此时 $f\otimes g$ 典范地等同于乘积 $fg$，命题显然。
\end{env}

\oldpage[0\textsubscript{I}]{54}

\begin{env}[5.5.4]
\label{0.5.5.4-fr}
设 $\mathcal{F}$ 是一个局部自由 $\mathcal{O}_X$-模层，秩为 $n$；
立即可见，$\bigwedge^p\mathcal{F}$ 是一个局部自由
$\mathcal{O}_X$-模层，其秩为 $\binom{n}{p}$（如果 $p\leq n$），而它
等于 $0$（如果 $p>n$）。这是因为问题是局部的，因而归结到
$\mathcal{F}=\mathcal{O}_X^n$ 的情形。此外，对每个 $x\in X$，
\[
  (\bigwedge^p\mathcal{F})_x/
  \mathfrak{m}_x(\bigwedge^p\mathcal{F})_x
\]
是一个维数为 $\binom{n}{p}$ 的 $k(x)$-向量空间，并典范地等同于
\[
  \bigwedge^p(\mathcal{F}_x/\mathfrak{m}_x\mathcal{F}_x).
\]
设 $s_1,\ldots,s_p$ 是
% EGA1-ERR-0011: French print uses bare F here although the ambient sheaf is mathcal F. Canonical translation preserves the printed symbol.
$F$ 在一个开集 $U$ 上的截面，该开集属于 $X$；并置
$s=s_1\wedge\cdots\wedge s_p$；它是
$\bigwedge^p\mathcal{F}$ 在 $U$ 上的一个截面 (4.1.5)。我们有
$s(x)=s_1(x)\wedge\cdots\wedge s_p(x)$；因而，说
$s_1(x),\ldots,s_p(x)$ \emph{线性相关}，就意味着 $s(x)=0$。由此可知，
\emph{满足如下性质的点 $x\in X$ 所成的集合：各
$s_1(x),\ldots,s_p(x)$ 线性无关；这个集合在 $X$ 中是开集}。事实上，
只须归结到
$\mathcal{F}=\mathcal{O}_X^n$ 的情形，并把 (5.5.2) 应用于 $s$ 在各个
投影下的像；这些投影从
$\bigwedge^p\mathcal{F}=\mathcal{O}_X^{\binom{n}{p}}$ 映到它的
$\binom{n}{p}$ 个因子之一。

特别地，若 $s_1,\ldots,s_n$ 是 $n$ 个截面，它们是 $\mathcal{F}$ 在
$U$ 上的截面，
并且 $s_1(x),\ldots,s_n(x)$ 对每个 $x\in U$ 都线性无关，则由这些
同态 $u:\mathcal{O}_X^n|U\to\mathcal{F}|U$（由这些 $s_i$ 定义）
(5.1.1) 是一个\emph{同构}。事实上，可以把问题限制到
$\mathcal{F}=\mathcal{O}_X^n$ 且将 $\bigwedge^n\mathcal{F}$ 与
$\mathcal{O}_X$ 典范等同的情形；此时
$s=s_1\wedge\cdots\wedge s_n$ 是 $\mathcal{O}_X$ 在 $U$ 上的一个
\emph{可逆}截面，并可用 Cramer 公式定义 $u$ 的逆同态。
\end{env}

\begin{env}[5.5.5]
\label{0.5.5.5-fr}
设 $\mathcal{E}$、$\mathcal{F}$ 是两个（有限秩）局部自由
$\mathcal{O}_X$-模层，并设 $u:\mathcal{E}\to\mathcal{F}$ 是一个同态。
为了存在一个
% EGA1-ERR-0012: French print names the neighbourhood u although the restriction immediately uses U. Canonical translation preserves lowercase u.
邻域 $u$，它是 $x\in X$ 的邻域，并使 $u|U$ \emph{单射}，且使
$\mathcal{F}|U$ \emph{为 $u(\mathcal{E})|U$ 与一个局部自由
子-$(\mathcal{O}_X|U)$-模层 $\mathcal{G}$ 的直和}，其充要条件是：
$u_x:\mathcal{E}_x\to\mathcal{F}_x$ 在取商之后给出向量空间间的
\emph{单射}同态
\[
  \mathcal{E}_x/\mathfrak{m}_x\mathcal{E}_x
  \longrightarrow
  \mathcal{F}_x/\mathfrak{m}_x\mathcal{F}_x.
\]
这个条件确实是\emph{必要的}，因为此时 $\mathcal{F}_x$ 是两个自由
$\mathcal{O}_x$-模 $u_x(\mathcal{E}_x)$ 与 $\mathcal{G}_x$ 的直和；
所以 $\mathcal{F}_x/\mathfrak{m}_x\mathcal{F}_x$ 是
$u_x(\mathcal{E}_x)/\mathfrak{m}_x u_x(\mathcal{E}_x)$ 与
$\mathcal{G}_x/\mathfrak{m}_x\mathcal{G}_x$ 的直和。这个条件也是
\emph{充分的}，因为可以归结到
$\mathcal{E}=\mathcal{O}_X^m$ 的情形。设 $s_1,\ldots,s_m$ 是各截面在
$u$ 下的像；这些截面记为 $e_i$，属于 $\mathcal{O}_X^m$，并使
$(e_i)_y$ 等于第 $i$ 个元素，该元素属于 $\mathcal{O}_y^m$ 的典范基，
对每个 $y\in X$ 都如此（即 $\mathcal{O}_X^m$ 的\emph{典范截面}）。
按照假设，
$s_1(x),\ldots,s_m(x)$ 线性无关；因此，若 $\mathcal{F}$ 的秩为 $n$，
则存在 $n-m$ 个截面 $s_{m+1},\ldots,s_n$；它们是 $\mathcal{F}$ 在
某个邻域 $V$ 上的截面，而这个邻域是 $x$ 的邻域；并且各 $s_i(x)$
$(1\leq i\leq n)$ 构成
$\mathcal{F}_x/\mathfrak{m}_x\mathcal{F}_x$ 的一个基。于是存在一个
邻域 $U\subset V$，它是 $x$ 的邻域，使得各 $s_i(y)$
$(1\leq i\leq n)$
构成 $\mathcal{F}_y/\mathfrak{m}_y\mathcal{F}_y$ 的一个基，对每个
% EGA1-ERR-0013: French print quantifies y over V after shrinking to U. Canonical translation preserves V.
$y\in V$ 都如此 (5.5.4)。再由 (5.5.4) 可知，存在一个从
$\mathcal{F}|U$ 到 $\mathcal{O}_X^n|U$ 的同构，它把各
$s_i|U$ $(1\leq i\leq m)$ 映到 $e_i|U$；证明完毕。
\end{env}

\section{平坦性}
\label{section:0.6-fr}

\begin{env}[6.0]
\label{0.6.0.0-fr}
平坦性的概念源于 J.-P.~Serre [16]；下文省略那些已在
N.~Bourbaki 的\emph{《交换代数》}中给出的结果之证明，并请读者参阅该书。
我们假设所有环都是交换的。\footnote{关于其中大多数结果推广到非交换情形，
见所引 N.~Bourbaki 的论述。}

\oldpage[0\textsubscript{I}]{55}

若 $M$、$N$ 是两个 $A$-模，$M'$（相应地，$N'$）是 $M$
（相应地，$N$）的一个子模，则以
$\operatorname{Im}(M'\otimes_A N')$ 表示 $M\otimes_A N$ 的如下子模：
典范映射 $M'\otimes_A N'\to M\otimes_A N$ 的像。
\end{env}

\subsection{平坦模}
\label{subsection:0.6.1-fr}

\begin{env}[6.1.1]
\label{0.6.1.1-fr}
设 $M$ 是一个 $A$-模。以下性质彼此等价：
\begin{enumerate}
  \item[$a)$] 在 $A$-模范畴中，关于 $N$ 的函子
    $M\otimes_A N$ 是正合的；
  \item[$b)$] 对每个 $i>0$ 以及每个 $A$-模 $N$，都有
    $\operatorname{Tor}_i^A(M,N)=0$；
  \item[$c)$] 对每个 $A$-模 $N$，都有
    $\operatorname{Tor}_1^A(M,N)=0$。
\end{enumerate}

当 $M$ 满足这些条件时，称 $M$ 是一个\emph{平坦 $A$-模}。显然，每个自由
$A$-模都是平坦的。

为了使 $M$ 成为一个平坦 $A$-模，只要对 $A$ 的每个\emph{有限型}理想
$\mathfrak{J}$，典范映射
$M\otimes_A\mathfrak{J}\to M\otimes_A A=M$ 都是\emph{单射}即可。
\end{env}

\begin{env}[6.1.2]
\label{0.6.1.2-fr}
平坦 $A$-模的任意归纳极限仍是平坦 $A$-模。$A$-模的直和
$\bigoplus_{\lambda\in L}M_\lambda$ 是平坦 $A$-模，当且仅当每个
$A$-模 $M_\lambda$ 都是平坦的。特别地，每个投射 $A$-模都是平坦的。

设
\[
  0\longrightarrow M'\longrightarrow M\longrightarrow M''\longrightarrow 0
\]
是一个 $A$-模正合列，并且 $M''$ \emph{平坦}。那么，对每个 $A$-模
$N$，序列
\[
  0\longrightarrow M'\otimes N\longrightarrow M\otimes N
  \longrightarrow M''\otimes N\longrightarrow 0
\]
都是正合的。此外，在这种情形下，$M$ 平坦的充要条件是 $M'$ 平坦
（但可能出现 $M$ 与 $M'$ 都平坦，而 $M''=M/M'$ 不平坦的情形）。
\end{env}

\begin{env}[6.1.3]
\label{0.6.1.3-fr}
设 $M$ 是一个平坦 $A$-模，$N$ 是任意 $A$-模；对 $N$ 的两个子模
$N'$、$N''$，有
\[
  \operatorname{Im}(M\otimes(N'+N''))
  =\operatorname{Im}(M\otimes N')+\operatorname{Im}(M\otimes N''),
\]
\[
  \operatorname{Im}(M\otimes(N'\cap N''))
  =\operatorname{Im}(M\otimes N')\cap\operatorname{Im}(M\otimes N'')
\]
（这里各像都取在 $M\otimes N$ 中）。
\end{env}

\begin{env}[6.1.4]
\label{0.6.1.4-fr}
设 $M$、$N$ 是两个 $A$-模，$M'$（相应地，$N'$）是 $M$
（相应地，$N$）的一个子模，并假设 $M/M'$、$N/N'$ 这两个模中有一个
是平坦的。那么
\[
  \operatorname{Im}(M'\otimes N')
  =\operatorname{Im}(M'\otimes N)\cap\operatorname{Im}(M\otimes N')
\]
（各像都取在 $M\otimes N$ 中）。特别地，若 $\mathfrak{J}$ 是
$A$ 的一个理想且 $M/M'$ 平坦，则
$\mathfrak{J}M'=M'\cap\mathfrak{J}M$。
\end{env}

\subsection{环的变换}
\label{subsection:0.6.2-fr}

当一个加法群 $M$ 关于环 $A$、$B$、……带有多个模结构时，我们除了说
$M$ 作为 $A$-模、$B$-模、……是平坦的，有时也说 $M$
\emph{在 $A$ 上平坦}、\emph{在 $B$ 上平坦}、……。

\begin{env}[6.2.1]
\label{0.6.2.1-fr}
设 $A$、$B$ 是两个环，$M$ 是一个 $A$-模，$N$ 是一个
$(A,B)$-双模。若 $M$ 平坦且 $N$ 在 $B$ 上平坦，则
$M\otimes_A N$ 在 $B$ 上平坦。特别地，若 $M$、$N$ 是两个平坦
$A$-模，则 $M\otimes_A N$ 是一个平坦 $A$-模。若 $B$ 是一个
$A$-代数，而 $M$ 是
\oldpage[0\textsubscript{I}]{56}
一个平坦 $A$-模，则 $B$-模 $M_{(B)}=M\otimes_A B$ 是平坦的。
最后，若 $B$ 是一个作为 $A$-模平坦的 $A$-代数，而 $N$ 是一个平坦
$B$-模，则 $N$ 也在 $A$ 上平坦。
\end{env}

\begin{env}[6.2.2]
\label{0.6.2.2-fr}
设 $A$ 是一个环，$B$ 是一个作为 $A$-模平坦的 $A$-代数。再设
$M$、$N$ 是两个 $A$-模，并且 $M$ 有一个有限表示；那么典范同态
\[
  \operatorname{Hom}_A(M,N)\otimes_A B
  \longrightarrow
  \operatorname{Hom}_B(M\otimes_A B,N\otimes_A B)
  \tag{6.2.2.1}
\]
（它把 $u\otimes b$ 变为同态
$m\otimes b'\longmapsto u(m)\otimes b'b$）是一个同构。
\end{env}

\begin{env}[6.2.3]
\label{0.6.2.3-fr}
设 $(A_\lambda,\varphi_{\mu\lambda})$ 是一个环的滤过归纳系，并令
$A=\varinjlim A_\lambda$。另一方面，对每个 $\lambda$，设
$M_\lambda$ 是一个 $A_\lambda$-模；当 $\lambda\leq\mu$ 时，设
$\theta_{\mu\lambda}:M_\lambda\to M_\mu$ 是一个
$\varphi_{\mu\lambda}$-同态，并使
$(M_\lambda,\theta_{\mu\lambda})$ 成为一个归纳系；那么
$M=\varinjlim M_\lambda$ 是一个 $A$-模。在这种情形下，若对每个
$\lambda$，$M_\lambda$ 都是\emph{平坦} $A_\lambda$-模，则 $M$
是一个\emph{平坦} $A$-模。事实上，设 $\mathfrak{J}$ 是 $A$
的一个\emph{有限型}理想；由归纳极限的定义，存在一个指标
$\lambda$ 以及 $A_\lambda$ 的一个理想 $\mathfrak{J}_\lambda$，使
$\mathfrak{J}=\mathfrak{J}_\lambda A$。若对
$\mu\geq\lambda$ 置
$\mathfrak{J}'_\mu=\mathfrak{J}_\lambda A_\mu$，则也有
$\mathfrak{J}=\varinjlim\mathfrak{J}'_\mu$（这里 $\mu$ 遍历所有
指标 $\geq\lambda$），从而（因为函子 $\varinjlim$ 正合并且
与张量积可交换）
\[
  M\otimes_A\mathfrak{J}
  =\varinjlim(M_\mu\otimes_{A_\mu}\mathfrak{J}'_\mu)
  =\varinjlim\mathfrak{J}'_\mu M_\mu
  =\mathfrak{J}M.
\]
\end{env}

\subsection{平坦性的局部化}
\label{subsection:0.6.3-fr}

\begin{env}[6.3.1]
\label{0.6.3.1-fr}
若 $A$ 是一个环，$S$ 是 $A$ 的一个乘法子集，则 $S^{-1}A$ 是一个
\emph{平坦} $A$-模。事实上，对每个 $A$-模 $N$，
$N\otimes_A S^{-1}A$ 可与 $S^{-1}N$ 等同 (1.2.5)，而我们已知
(1.3.2)，$S^{-1}N$ 关于 $N$ 是一个正合函子。

现在若 $M$ 是一个平坦 $A$-模，则
$S^{-1}M=M\otimes_A S^{-1}A$ 是一个平坦 $S^{-1}A$-模 (6.2.1)，
因而根据前述结论与 (6.2.1)，它在 $A$ 上也平坦。特别地，若 $P$
是一个 $S^{-1}A$-模，则可把它视为一个同构于 $S^{-1}P$ 的 $A$-模；
$P$ 在 $A$ 上平坦的充要条件是它在 $S^{-1}A$ 上平坦。
\end{env}

\begin{env}[6.3.2]
\label{0.6.3.2-fr}
设 $A$ 是一个环，$B$ 是一个 $A$-代数，$T$ 是 $B$ 的一个乘法子集。
若 $P$ 是一个在 $A$ 上平坦的 $B$-模，则 $T^{-1}P$ 在 $A$ 上平坦。
事实上，对每个 $A$-模 $N$，有
\[
  (T^{-1}P)\otimes_A N
  =(T^{-1}B\otimes_B P)\otimes_A N
  =T^{-1}B\otimes_B(P\otimes_A N)
  =T^{-1}(P\otimes_A N)~;
\]
而 $T^{-1}(P\otimes_A N)$ 关于 $N$ 是一个正合函子，因为它是两个正合
函子 $P\otimes_A N$（关于 $N$）与 $T^{-1}Q$（关于 $Q$）的复合。
若 $S$ 是 $A$ 的一个乘法子集，并且它在 $B$ 中的像
\emph{包含于 $T$}，则 $T^{-1}P$ 等于 $S^{-1}(T^{-1}P)$，所以由
(6.3.1)，它在 $S^{-1}A$ 上也平坦。
\end{env}

\begin{env}[6.3.3]
\label{0.6.3.3-fr}
设 $\varphi:A\to B$ 是一个环同态，$M$ 是一个 $B$-模。以下性质彼此
等价：
\begin{enumerate}
  \item[$a)$] $M$ 是一个平坦 $A$-模。
  \item[$b)$] 对 $B$ 的每个极大理想 $\mathfrak{n}$，
    $M_\mathfrak{n}$ 都是一个平坦 $A$-模。
  \item[$c)$] 对 $B$ 的每个极大理想 $\mathfrak{n}$，置
    $\mathfrak{m}=\varphi^{-1}(\mathfrak{n})$ 后，
    $M_\mathfrak{n}$ 是一个平坦 $A_\mathfrak{m}$-模。
\end{enumerate}

事实上，因为 $M_\mathfrak{n}=(M_\mathfrak{n})_\mathfrak{m}$，
$b)$ 与 $c)$ 的等价性来自 (6.3.1)，而 $a)$ 蕴含 $b)$ 这件事是
(6.3.2) 的一个特例。还须证明 $b)$ 蕴含 $a)$，
\oldpage[0\textsubscript{I}]{57}
也就是说，对 $A$-模的每个单同态 $u:N'\to N$，同态
$v=1\otimes u:M\otimes_A N'\to M\otimes_A N$ 都是单射。可是，
$v$ 也是一个 $B$-模同态；我们知道，为使它是单射，只要对 $B$ 的
每个极大理想 $\mathfrak{n}$，同态
\[
  v_\mathfrak{n}:(M\otimes_A N')_\mathfrak{n}
  \longrightarrow(M\otimes_A N)_\mathfrak{n}
\]
都是单射即可。但由于
\[
  (M\otimes_A N)_\mathfrak{n}
  =B_\mathfrak{n}\otimes_B(M\otimes_A N)
  =M_\mathfrak{n}\otimes_A N,
\]
$v_\mathfrak{n}$ 正是同态
\[
  1\otimes u:M_\mathfrak{n}\otimes_A N'
  \longrightarrow M_\mathfrak{n}\otimes_A N,
\]
它是单射，因为 $M_\mathfrak{n}$ 在 $A$ 上平坦。

特别地（令 $B=A$），$A$-模 $M$ 平坦的充要条件是：对 $A$ 的每个
极大理想 $\mathfrak{m}$，$M_\mathfrak{m}$ 都在
$A_\mathfrak{m}$ 上平坦。
\end{env}

\begin{env}[6.3.4]
\label{0.6.3.4-fr}
设 $M$ 是一个 $A$-模；若 $M$ 平坦，而 $f\in A$ 不是 $A$ 中
$0$ 的因子，则 $f$ 不零化 $M$ 中任何 $\neq0$ 的元素。事实上，同态
$m\mapsto f\cdot m$ 可写成 $1\otimes u$，其中 $u$ 是 $A$ 中的乘法
$a\mapsto f\cdot a$，并把 $M$ 与 $M\otimes_A A$ 等同；若 $u$
是单射，则因 $M$ 平坦，$1\otimes u$ 也是单射。特别地，若 $A$
是\emph{整环}，则 $M$ \emph{无挠}。

反过来，假设 $A$ 是整环，$M$ 无挠，并且对 $A$ 的每个极大理想
$\mathfrak{m}$，$A_\mathfrak{m}$ 都是一个\emph{离散赋值环}；
那么 $M$ 在 $A$ 上平坦。事实上，由 (6.3.3)，只须证明
$M_\mathfrak{m}$ 在 $A_\mathfrak{m}$ 上平坦，因而可以假设 $A$
本身就是一个离散赋值环。又因为 $M$ 是其有限型子模的归纳极限，而这些
子模都无挠，所以还可以归结到 $M$ 为有限型的情形 (6.1.2)。此时命题
来自 $M$ 是一个自由 $A$-模这一事实。

特别地，若 $A$ 是一个\emph{整环}，而 $\varphi:A\to B$ 是一个环同态，
它使 $B$ 成为一个\emph{平坦}且 $\neq\{0\}$ 的 $A$-模，则
$\varphi$ 必为\emph{单射}。反过来，若 $B$ 是整环，$A$ 是 $B$
的一个子环，并且对 $A$ 的每个极大理想 $\mathfrak{m}$，
$A_\mathfrak{m}$ 都是一个离散赋值环，则 $B$ 在 $A$ 上平坦。
\end{env}

\subsection{忠实平坦模}
\label{subsection:0.6.4-fr}

\begin{env}[6.4.1]
\label{0.6.4.1-fr}
对一个 $A$-模 $M$，以下四个性质彼此等价：
\begin{enumerate}
  \item[$a)$] $A$-模序列 $N'\to N\to N''$ 正合的充要条件是序列
    $M\otimes N'\to M\otimes N\to M\otimes N''$ 正合；
  \item[$b)$] $M$ 平坦，并且对每个 $A$-模 $N$，关系
    $M\otimes N=0$ 蕴含 $N=0$；
  \item[$c)$] $M$ 平坦，并且对 $A$-模的每个同态 $v:N\to N'$，
    关系 $1_M\otimes v=0$ 蕴含 $v=0$，其中 $1_M$ 是 $M$ 的恒等
    自同构；
  \item[$d)$] $M$ 平坦，并且对 $A$ 的每个极大理想
    $\mathfrak{m}$，都有 $\mathfrak{m}M\neq M$。
\end{enumerate}

当 $M$ 满足这些条件时，称 $M$ 是一个\emph{忠实平坦 $A$-模}；
此时 $M$ 必为一个\emph{忠实}模。此外，若 $u:N\to N'$ 是一个
$A$-模同态，则 $u$ 是单射（相应地，满射、双射）的充要条件是
$1\otimes u:M\otimes N\to M\otimes N'$ 具有同一性质。
\end{env}

\oldpage[0\textsubscript{I}]{58}

\begin{env}[6.4.2]
\label{0.6.4.2-fr}
一个自由模 $\neq\{0\}$ 是忠实平坦的；一个平坦模与一个忠实
平坦模的直和也是如此。若 $S$ 是 $A$ 的一个乘法子集，则只有当 $S$
由可逆元素组成时，$S^{-1}A$ 才是忠实平坦 $A$-模（因此
$S^{-1}A=A$）。
\end{env}

\begin{env}[6.4.3]
\label{0.6.4.3-fr}
设
\[
  0\longrightarrow M'\longrightarrow M\longrightarrow M''\longrightarrow 0
\]
是一个 $A$-模正合列；若 $M'$ 与 $M''$ 都平坦，并且其中一个忠实平坦，
则 $M$ 忠实平坦。
\end{env}

\begin{env}[6.4.4]
\label{0.6.4.4-fr}
设 $A$、$B$ 是两个环，$M$ 是一个 $A$-模，$N$ 是一个
$(A,B)$-双模。若 $M$ 忠实平坦，而 $N$ 是一个忠实平坦 $B$-模，
则 $M\otimes_A N$ 是一个忠实平坦 $B$-模。特别地，若 $M$、$N$
是两个忠实平坦 $A$-模，则 $M\otimes_A N$ 也是忠实平坦的。若
$B$ 是一个 $A$-代数，而 $M$ 是一个忠实平坦 $A$-模，则 $B$-模
$M_{(B)}$ 忠实平坦。
\end{env}

\begin{env}[6.4.5]
\label{0.6.4.5-fr}
若 $M$ 是一个忠实平坦 $A$-模，而 $S$ 是 $A$ 的一个乘法子集，则
$S^{-1}M$ 是一个忠实平坦 $S^{-1}A$-模，因为
$S^{-1}M=M\otimes_A(S^{-1}A)$ (6.4.4)。反过来，若对 $A$ 的每个
极大理想 $\mathfrak{m}$，$M_\mathfrak{m}$ 都是一个忠实平坦
$A_\mathfrak{m}$-模，则 $M$ 是一个忠实平坦 $A$-模，因为 $M$
在 $A$ 上平坦 (6.3.3)，并且
\[
  M_\mathfrak{m}/\mathfrak{m}M_\mathfrak{m}
  =(M\otimes_A A_\mathfrak{m})
  \otimes_{A_\mathfrak{m}}
  (A_\mathfrak{m}/\mathfrak{m}A_\mathfrak{m})
  =M\otimes_A(A/\mathfrak{m})
  =M/\mathfrak{m}M,
\]
所以假设蕴含：对 $A$ 的每个极大理想 $\mathfrak{m}$，
$M/\mathfrak{m}M\neq0$；这就证明了所断言的结论 (6.4.1)。
\end{env}

\subsection{标量限制}
\label{subsection:0.6.5-fr}

\begin{env}[6.5.1]
\label{0.6.5.1-fr}
设 $A$ 是一个环，$\varphi:A\to B$ 是一个使 $B$ 成为 $A$-代数的环同态。
假设存在一个 $B$-模 $N$，它是一个\emph{忠实平坦} $A$-模。那么，对每个
$A$-模 $M$，从 $M$ 到 $B\otimes_A M=M_{(B)}$ 的同态
$x\mapsto1\otimes x$ 是\emph{单射}。特别地，$\varphi$ 是单射；
对 $A$ 的每个理想 $\mathfrak{a}$，有
$\varphi^{-1}(\mathfrak{a}B)=\mathfrak{a}$；对 $A$ 的每个极大
（相应地，素）理想 $\mathfrak{m}$，存在 $B$ 的一个极大（相应地，素）
理想 $\mathfrak{n}$，使
$\varphi^{-1}(\mathfrak{n})=\mathfrak{m}$。
\end{env}

\begin{env}[6.5.2]
\label{0.6.5.2-fr}
当 (6.5.1) 的条件满足时，我们通过 $\varphi$ 把 $A$ 等同于 $B$ 的一个
子环；更一般地，对每个 $A$-模 $M$，把 $M$ 等同于 $M_{(B)}$ 的一个
$A$-子模。应注意，此时若 $B$ 是\emph{Noether 环}，则 $A$ 也是
Noether 环，因为映射 $\mathfrak{a}\mapsto\mathfrak{a}B$ 是从 $A$
的理想集合到 $B$ 的理想集合的保序单射；因此，$A$ 的理想若存在一个
无穷严格递增序列，就会蕴含 $B$ 的理想也存在一个同类序列。
\end{env}

\subsection{忠实平坦环}
\label{subsection:0.6.6-fr}

\begin{env}[6.6.1]
\label{0.6.6.1-fr}
设 $\varphi:A\to B$ 是一个使 $B$ 成为 $A$-代数的环同态。以下五个
性质彼此等价：
\begin{enumerate}
  \item[$a)$] $B$ 是一个忠实平坦 $A$-模（换言之，
    $M_{(B)}$ 关于 $M$ 是一个正合且忠实的函子）。
  \item[$b)$] 同态 $\varphi$ 是单射，而 $A$-模
    $B/\varphi(A)$ 是平坦的。

\oldpage[0\textsubscript{I}]{59}

  \item[$c)$] $A$-模 $B$ 是平坦的（换言之，函子 $M_{(B)}$
    正合），并且对每个 $A$-模 $M$，从 $M$ 到 $M_{(B)}$ 的同态
    $x\mapsto1\otimes x$ 是单射。
  \item[$d)$] $A$-模 $B$ 是平坦的，并且对 $A$ 的每个理想
    $\mathfrak a$，有 $\varphi^{-1}(\mathfrak aB)=\mathfrak a$。
  \item[$e)$] $A$-模 $B$ 是平坦的，并且对 $A$ 的每个极大理想
    $\mathfrak m$，存在 $B$ 的一个极大理想 $\mathfrak n$，使
    $\varphi^{-1}(\mathfrak n)=\mathfrak m$。
\end{enumerate}

当这些条件成立时，我们把 $A$ 等同于 $B$ 的一个子环。
\end{env}

\begin{env}[6.6.2]
\label{0.6.6.2-fr}
设 $A$ 是一个局部环，$\mathfrak m$ 是它的极大理想，$B$ 是一个
$A$-代数，并且 $\mathfrak mB\neq B$（例如，当 $B$ 是一个局部环，
而 $A\to B$ 是一个局部同态时，这一条件成立）。若 $B$ 是一个平坦
$A$-模，则 $B$ 是一个忠实平坦 $A$-模。事实上，因为
$\mathfrak mB\neq B$，所以 $B$ 中存在一个包含 $\mathfrak mB$
的极大理想 $\mathfrak n$；由于
$\varphi^{-1}(\mathfrak n)\cap A$ 包含 $\mathfrak m$ 而不包含
$1$，故 $\varphi^{-1}(\mathfrak n)=\mathfrak m$，从而可应用
(6.6.1) 的判据 $e)$。因此，在上述条件下，若 $B$ 是 Noether 环，
则 $A$ 也是 Noether 环 (6.5.2)。
\end{env}

\begin{env}[6.6.3]
\label{0.6.6.3-fr}
设 $B$ 是一个忠实平坦 $A$-模的 $A$-代数。对每个 $A$-模 $M$ 以及
$M$ 的每个 $A$-子模 $M'$，都有
$M'=M\cap M'_{(B)}$（这里把 $M$ 等同于 $M_{(B)}$ 的一个
$A$-子模）。$M$ 是平坦（相应地，忠实平坦）$A$-模的充要条件是
$M_{(B)}$ 是平坦（相应地，忠实平坦）$B$-模。
\end{env}

\begin{env}[6.6.4]
\label{0.6.6.4-fr}
设 $B$ 是一个 $A$-代数，$N$ 是一个忠实平坦 $B$-模。$B$ 是平坦
（相应地，忠实平坦）$A$-模的充要条件是 $N$ 具有同一性质。

特别地，设 $C$ 是一个 $B$-代数；若环 $C$ 在 $B$ 上忠实平坦，而
$B$ 在 $A$ 上忠实平坦，则 $C$ 在 $A$ 上忠实平坦；若 $C$ 在 $B$
上以及在 $A$ 上都忠实平坦，则 $B$ 在 $A$ 上忠实平坦。
\end{env}

\subsection{赋环空间的平坦态射}
\label{subsection:0.6.7-fr}

\begin{env}[6.7.1]
\label{0.6.7.1-fr}
设 $f:X\to Y$ 是赋环空间的一个态射，$\mathcal F$ 是一个
$\mathcal O_X$-模层。若 $\mathcal F_x$ 是一个平坦
$\mathcal O_{f(x)}$-模，则称 $\mathcal F$ 在点 $x\in X$ 处
\emph{$f$-平坦}（在不致混淆 $f$ 时，也称\emph{$Y$-平坦}）；
若 $\mathcal F$ 在所有点 $x\in f^{-1}(y)$ 处都 $f$-平坦，则称
$\mathcal F$ 在 $y\in Y$ \emph{上方 $f$-平坦}；若 $\mathcal F$
在 $X$ 的所有点处都 $f$-平坦，则称 $\mathcal F$ \emph{$f$-平坦}。
若 $\mathcal O_X$ 在 $x\in X$ 处 $f$-平坦（相应地，在 $y\in Y$
上方 $f$-平坦、$f$-平坦），则称态射 $f$ 在 $x$ 处平坦
（相应地，在 $y$ 上方平坦、平坦）。
\end{env}

\begin{env}[6.7.2]
\label{0.6.7.2-fr}
采用 (6.7.1) 的记号，若 $\mathcal F$ 在 $x$ 处 $f$-平坦，则对
$y=f(x)$ 的每个开邻域 $U$，关于 $\mathcal G$ 的函子
$(f^*(\mathcal G)\otimes_{\mathcal O_X}\mathcal F)_x$ 在
$(\mathcal O_Y|U)$-模层范畴中是正合的；事实上，这个茎典范地等同于
$\mathcal G_y\otimes_{\mathcal O_y}\mathcal F_x$，所断言的结论来自
定义。特别地，若 $f$ 是一个平坦态射，则函子 $f^*$ 在
$\mathcal O_Y$-模层范畴中是正合的。
\end{env}

\begin{env}[6.7.3]
\label{0.6.7.3-fr}
反过来，假设环层 $\mathcal O_Y$ 是凝聚的，并且对 $y$ 的每个开邻域
$U$，关于 $\mathcal G$ 的函子
$(f^*(\mathcal G)\otimes_{\mathcal O_X}\mathcal F)_x$ 在凝聚
$(\mathcal O_Y|U)$-模层范畴中都是正合的。那么 $\mathcal F$
在 $x$ 处 $f$-平坦。事实上，只须证明：对 $\mathcal O_y$ 的每个
有限型理想 $\mathfrak J$，典范同态
$\mathfrak J\otimes_{\mathcal O_y}\mathcal F_x\to\mathcal F_x$
都是单射 (6.1.1)。可是，我们知道 (5.3.8)，此时存在 $y$ 的一个

\oldpage[0\textsubscript{I}]{60}

开邻域 $U$，以及 $\mathcal O_Y|U$ 的一个凝聚理想层
$\mathcal J$，使 $\mathcal J_y=\mathfrak J$；结论由此成立。
\end{env}

\begin{env}[6.7.4]
\label{0.6.7.4-fr}
(6.1) 中关于平坦模的各个结果立即转述为关于在一点处 $f$-平坦的层
的命题：

若 $0\to\mathcal F'\to\mathcal F\to\mathcal F''\to0$ 是一个
$\mathcal O_X$-模层正合列，而 $\mathcal F''$ 在点 $x\in X$
处 $f$-平坦，则对 $y=f(x)$ 的每个开邻域 $U$ 以及每个
$(\mathcal O_Y|U)$-模层 $\mathcal G$，序列
\[
  0\longrightarrow
  (f^*(\mathcal G)\otimes_{\mathcal O_X}\mathcal F')_x
  \longrightarrow
  (f^*(\mathcal G)\otimes_{\mathcal O_X}\mathcal F)_x
  \longrightarrow
  (f^*(\mathcal G)\otimes_{\mathcal O_X}\mathcal F'')_x
  \longrightarrow0
\]
都是正合的。在这种情形下，$\mathcal F$ 在 $x$ 处 $f$-平坦的充要
条件是 $\mathcal F'$ 具有同一性质。对于 $\mathcal O_X$-模层在
$y\in Y$ 上方 $f$-平坦，或 $\mathcal O_X$-模层处处 $f$-平坦的
相应概念，也有类似陈述。
\end{env}

\begin{env}[6.7.5]
\label{0.6.7.5-fr}
设 $f:X\to Y$、$g:Y\to Z$ 是赋环空间的两个态射；设 $x\in X$、
$y=f(x)$，并设 $\mathcal F$ 是一个 $\mathcal O_X$-模层。若
$\mathcal F$ 在点 $x$ 处 $f$-平坦，而态射 $g$ 在点 $y$ 处平坦，
则 $\mathcal F$ 在 $x$ 处 $(g\circ f)$-平坦 (6.2.1)。特别地，
若 $f$、$g$ 是平坦态射，则 $g\circ f$ 平坦。
\end{env}

\begin{env}[6.7.6]
\label{0.6.7.6-fr}
设 $X$、$Y$ 是两个赋环空间，$f:X\to Y$ 是一个平坦态射。那么双函子
的典范同态 (4.4.6)
\[
  f^*(\mathcal{H}om_{\mathcal O_Y}(\mathcal F,\mathcal G))
  \longrightarrow
  \mathcal{H}om_{\mathcal O_X}(f^*(\mathcal F),f^*(\mathcal G))
  \tag{6.7.6.1}
  \label{0.6.7.6.1-fr}
\]
在 $\mathcal F$ 有有限表示时是一个同构 (5.2.5)。

事实上，由于问题是局部的，可以假设存在一个正合列
$\mathcal O_Y^m\to\mathcal O_Y^n\to\mathcal F\to0$。可是，由
关于 $f$ 的假设，(6.7.6.1) 两端关于 $\mathcal F$ 都是左正合函子；
于是归结到 $\mathcal F=\mathcal O_Y$ 时证明命题，而在这种情形下
结论是平凡的。
\end{env}

\begin{env}[6.7.8]
\label{0.6.7.8-fr}
若赋环空间的一个态射 $f:X\to Y$ 是满射，并且对每个 $x\in X$，
$\mathcal O_x$ 都是一个忠实平坦 $\mathcal O_{f(x)}$-模，则称
$f$ \emph{忠实平坦}。当 $X$、$Y$ 是局部赋环空间 (5.5.1) 时，
这等价于说态射 $f$ 是满射且平坦 (6.6.2)。当 $f$ 忠实平坦时，
$f^*$ 在 $\mathcal O_Y$-模层范畴中是一个正合且忠实的函子
(6.6.1, $a)$)；而 $\mathcal O_Y$-模层 $\mathcal G$ 为
$Y$-平坦的充要条件是 $f^*(\mathcal G)$ 为平坦模层 (6.6.3)。
\end{env}

\section{进制环}
\label{section:0.7-fr}

\subsection{可容环}
\label{subsection:0.7.1-fr}

\begin{env}[7.1.1]
\label{0.7.1.1-fr}
回忆一下，在一个拓扑环 $A$（不必分离）中，称元素 $x$
\emph{拓扑幂零}，如果 $0$ 是序列 $(x^n)_{n\geq0}$ 的一个极限。
若拓扑环 $A$ 中存在一个由理想（必为开理想）组成的 $0$ 的邻域基，
则称 $A$ 是一个\emph{线性拓扑化环}。
\end{env}

\begin{definition}[7.1.2]
\label{0.7.1.2-fr}
在线性拓扑化环 $A$ 中，称理想 $\mathfrak J$ 为一个\emph{定义理想}，
如果 $\mathfrak J$ 是开理想，并且对每个邻域 $V$（这里指 $0$ 的邻域），都存在
一个整数 $n>0$ 使得

\oldpage[0\textsubscript{I}]{61}

$\mathfrak J^n\subset V$（滥用语言时，也把这个条件表述为序列
$(\mathfrak J^n)$ 趋于 $0$）。
线性拓扑化环 $A$ 若在 $A$ 中存在定义理想，则称为\emph{预可容环}；
预可容而且分离且完备的环 $A$ 称为\emph{可容环}。
\end{definition}

显然，若 $\mathfrak J$ 是一个定义理想，$\mathfrak L$ 是 $A$ 的一个
开理想，则 $\mathfrak J\cap\mathfrak L$ 仍是一个定义理想；因此，
预可容环 $A$ 的定义理想组成 $0$ 的一个邻域基。

\begin{lemma}[7.1.3]
\label{0.7.1.3-fr}
设 $A$ 是一个线性拓扑化环。
\begin{enumerate}
  \item[{\upshape(i)}] 元素 $x\in A$ 拓扑幂零的充要条件是：对 $A$
    的每个开理想 $\mathfrak J$，$x$ 在 $A/\mathfrak J$ 中的典范像
    都是幂零元。$A$ 的所有拓扑幂零元组成的集合 $\mathfrak T$
    是一个理想。
  \item[{\upshape(ii)}] 再假设 $A$ 预可容，并设 $\mathfrak J$ 是
    $A$ 的一个定义理想。元素 $x\in A$ 拓扑幂零的充要条件是：
    它在 $A/\mathfrak J$ 中的典范像是幂零元；理想 $\mathfrak T$
    是 $A/\mathfrak J$ 的幂零根的原像，因而是开理想。
\end{enumerate}
\end{lemma}

(i) 立即由定义得到。为证明 (ii)，只须注意，对 $A$ 中 $0$ 的每个
邻域 $V$，都存在 $n>0$ 使得 $\mathfrak J^n\subset V$；若
$x\in A$ 满足 $x^m\in\mathfrak J$，则对 $q\geq n$ 有
$x^{mq}\in V$，故 $x$ 拓扑幂零。

\begin{proposition}[7.1.4]
\label{0.7.1.4-fr}
设 $A$ 是一个预可容环，$\mathfrak J$ 是 $A$ 的一个定义理想。
\begin{enumerate}
  \item[{\upshape(i)}] $A$ 的理想 $\mathfrak J'$ 包含在某个定义理想
    中的充要条件是：存在一个整数 $n>0$ 使得
    ${\mathfrak J'}^n\subset\mathfrak J$。
  \item[{\upshape(ii)}] 元素 $x\in A$ 包含在某个定义理想中的充要
    条件是：它拓扑幂零。
\end{enumerate}
\end{proposition}

(i) 若 ${\mathfrak J'}^n\subset\mathfrak J$，则对 $A$ 中 $0$ 的每个
开邻域 $V$，都存在 $m$ 使得 $\mathfrak J^m\subset V$，因而
${\mathfrak J'}^{mn}\subset V$。

(ii) 该条件显然是必要的；它也是充分的，因为当它成立时，存在 $n$
使得 $x^n\in\mathfrak J$，于是
$\mathfrak J'=\mathfrak J+Ax$ 是一个定义理想：它是开理想，并且
${\mathfrak J'}^n\subset\mathfrak J$。

\begin{corollary}[7.1.5]
\label{0.7.1.5-fr}
在预可容环 $A$ 中，每个开素理想都包含所有定义理想。
\end{corollary}

\begin{corollary}[7.1.6]
\label{0.7.1.6-fr}
沿用 (7.1.4) 的记号和假设，$A$ 的理想 $\mathfrak J_0$ 的下列性质
等价：
\begin{enumerate}
  \item[$a)$] $\mathfrak J_0$ 是 $A$ 的最大定义理想；
  \item[$b)$] $\mathfrak J_0$ 是一个极大定义理想；
  \item[$c)$] $\mathfrak J_0$ 是一个定义理想，并且环
    $A/\mathfrak J_0$ 是约化环。
\end{enumerate}
存在具有这些性质的理想 $\mathfrak J_0$ 的充要条件是：
$A/\mathfrak J$ 的幂零根是幂零理想；此时，$\mathfrak J_0$ 等于
$A$ 的拓扑幂零元所成的理想 $\mathfrak T$。
\end{corollary}

显然，$a)$ 蕴涵 $b)$，而由 (7.1.4, (ii)) 和 (7.1.3, (ii))，
$b)$ 蕴涵 $c)$；基于同一理由，$c)$ 蕴涵 $a)$。最后一个断言由
(7.1.4, (i)) 和 (7.1.3, (ii)) 得到。

当 $A/\mathfrak J$ 的幂零根 $\mathfrak T/\mathfrak J$ 是幂零理想
时，以 $A_{\mathrm{red}}$ 表示约化商环 $A/\mathfrak T$。

\oldpage[0\textsubscript{I}]{62}

\begin{corollary}[7.1.7]
\label{0.7.1.7-fr}
每个 Noether 预可容环都有最大定义理想。
\end{corollary}

\begin{corollary}[7.1.8]
\label{0.7.1.8-fr}
设预可容环 $A$ 具有如下性质：对某个定义理想 $\mathfrak J$，
各次幂 $\mathfrak J^n$（$n>0$）组成 $0$ 的一个邻域基。那么，对
$A$ 的任意定义理想 $\mathfrak J'$，各次幂 ${\mathfrak J'}^n$
也构成这样的邻域基。
\end{corollary}

\begin{definition}[7.1.9]
\label{0.7.1.9-fr}
若预可容环 $A$ 有一个定义理想 $\mathfrak J$，使得各
$\mathfrak J^n$ 组成 $A$ 中 $0$ 的一个邻域基（等价地，各
$\mathfrak J^n$ 都是开理想），则称 $A$ 是一个\emph{预进制环}。
分离且完备的预进制环称为\emph{进制环}。
\end{definition}

若 $\mathfrak J$ 是预进制（相应地，进制）环 $A$ 的一个定义理想，
则也称 $A$ 是一个\emph{$\mathfrak J$-预进制环}（相应地，
\emph{$\mathfrak J$-进制环}），并称它的拓扑为
\emph{$\mathfrak J$-预进制拓扑}（相应地，
\emph{$\mathfrak J$-进制拓扑}）。更一般地，若 $M$ 是一个
$A$-模，则以各子模 $\mathfrak J^nM$ 作为 $0$ 的邻域基所得的
$M$ 上的拓扑称为\emph{$\mathfrak J$-预进制拓扑}（相应地，
\emph{$\mathfrak J$-进制拓扑}）。由 (7.1.8)，这些拓扑与定义理想
$\mathfrak J$ 的选择无关。

\begin{proposition}[7.1.10]
\label{0.7.1.10-fr}
设 $A$ 是一个可容环，$\mathfrak J$ 是 $A$ 的一个定义理想。
那么，$\mathfrak J$ 包含在 $A$ 的 Jacobson 根中。
\end{proposition}

这一陈述等价于以下任一推论：

\begin{corollary}[7.1.11]
\label{0.7.1.11-fr}
对每个 $x\in\mathfrak J$，$1+x$ 在 $A$ 中可逆。
\end{corollary}

\begin{corollary}[7.1.12]
\label{0.7.1.12-fr}
元素 $f\in A$ 在 $A$ 中可逆的充要条件是：它在
$A/\mathfrak J$ 中的典范像在 $A/\mathfrak J$ 中可逆。
\end{corollary}

\begin{corollary}[7.1.13]
\label{0.7.1.13-fr}
对每个有限型 $A$-模 $M$，关系 $M=\mathfrak JM$（等价于
$M\otimes_A(A/\mathfrak J)=0$）蕴涵 $M=0$。
\end{corollary}

\begin{corollary}[7.1.14]
\label{0.7.1.14-fr}
设 $u:M\to N$ 是一个 $A$-模同态，其中 $N$ 是有限型模；
$u$ 为满射的充要条件是
$u\otimes1:M\otimes_A(A/\mathfrak J)\to N\otimes_A(A/\mathfrak J)$
为满射。
\end{corollary}

事实上，(7.1.10) 与 (7.1.11) 的等价性由 Bourbaki，
\emph{Alg.}，第 VIII 章，\textsection6，第 3 号，定理 1 得到，
而 (7.1.10) 与 (7.1.13) 的等价性由\emph{同上}，定理 2 得到；
(7.1.10) 蕴涵 (7.1.14) 来自\emph{同上}，命题 6 的推论 4；另一方面，
把 (7.1.14) 应用于零同态，可见它蕴涵 (7.1.13)。最后，(7.1.10)
蕴涵：若 $f$ 在 $A/\mathfrak J$ 中可逆，则 $f$ 不包含在 $A$
的任何极大理想中，因而 $f$ 在 $A$ 中可逆；换言之，(7.1.10)
蕴涵 (7.1.12)。反过来，(7.1.12) 蕴涵 (7.1.11)。

因此，只须证明 (7.1.11)。由于 $A$ 分离且完备，而序列
$(\mathfrak J^n)$ 趋于 $0$，级数
$\sum_{n=0}^{\infty}(-1)^n x^n$ 显然在 $A$ 中收敛；若 $y$ 是它的
和，则 $y(1+x)=1$。

\subsection{进制环与逆极限}
\label{subsection:0.7.2-fr}

\begin{env}[7.2.1]
\label{0.7.2.1-fr}
离散环的任意逆极限显然都是线性拓扑化、分离且完备的环。反过来，
设 $A$ 是一个线性拓扑化环，并设 $(\mathfrak J_\lambda)$ 是由理想
组成的 $A$ 中 $0$ 的开邻域基。

\oldpage[0\textsubscript{I}]{63}

各典范映射 $\varphi_\lambda:A\to A/\mathfrak J_\lambda$ 组成一个
连续表示的逆系统，因而定义一个连续表示
$\varphi:A\to\varprojlim A/\mathfrak J_\lambda$；若 $A$ 分离，则
$\varphi$ 是从 $A$ 到 $\varprojlim A/\mathfrak J_\lambda$ 的一个
稠密子环上的拓扑同构；若 $A$ 还完备，则 $\varphi$ 是从 $A$ 到
$\varprojlim A/\mathfrak J_\lambda$ 的拓扑同构。
\end{env}

\begin{lemma}[7.2.2]
\label{0.7.2.2-fr}
一个线性拓扑化环可容的充要条件是：它同构于一个逆极限
$A=\varprojlim A_\lambda$，其中 $(A_\lambda,u_{\lambda\mu})$ 是一个
离散环的逆系统，其指标集为关于 $\leq$ 的滤过有序集 $L$；该集合
有一个记作 $0$ 的最小元，并满足以下条件：1\textsuperscript{o}
各 $u_\lambda:A\to A_\lambda$ 都是满射；2\textsuperscript{o}
$u_{0\lambda}:A_\lambda\to A_0$ 的核 $\mathfrak J_\lambda$ 是幂零
理想。此时，$u_0:A\to A_0$ 的核 $\mathfrak J$ 等于
$\varprojlim\mathfrak J_\lambda$。
\end{lemma}

条件的必要性由 (7.2.1) 得到：取由包含在某个定义理想
$\mathfrak J_0$ 中的定义理想所组成的 $0$ 的一个邻域基
$(\mathfrak J_\lambda)$，并应用 (7.1.4, (i))。充分性由逆极限的定义
和 (7.1.2) 得到，最后一个断言则是显然的。

\begin{env}[7.2.3]
\label{0.7.2.3-fr}
设 $A$ 是一个可容拓扑环，$\mathfrak J$ 是 $A$ 的一个包含在某个
定义理想中的理想（换言之，由 (7.1.4)，序列 $(\mathfrak J^n)$
趋于 $0$）；可以在 $A$ 上考虑如下环拓扑：各次幂
$\mathfrak J^n$（$n>0$）组成 $0$ 的一个邻域基；我们仍称它为
\emph{$\mathfrak J$-预进制拓扑}。$A$ 可容这一假设蕴涵
$\bigcap_n\mathfrak J^n=0$，所以 $A$ 上的 $\mathfrak J$-预进制
拓扑是分离的；令
$\widehat A=\varprojlim A/\mathfrak J^n$ 为 $A$ 关于这一拓扑的
完备化（其中各 $A/\mathfrak J^n$ 都赋予离散拓扑），并以 $u$ 表示
环同态 $A\to\widehat A$（不必连续），即同态序列
$u_n:A\to A/\mathfrak J^n$ 的逆极限。另一方面，$A$ 上的
$\mathfrak J$-预进制拓扑比给定拓扑 $\mathcal T$ 更细；由于
$A$ 关于 $\mathcal T$ 分离且完备，可以把 $A$（赋予
$\mathfrak J$-预进制拓扑）到赋予 $\mathcal T$ 的 $A$ 的恒等映射
连续延拓；由此得到一个连续表示 $v:\widehat A\to A$。
\end{env}

\begin{proposition}[7.2.4]
\label{0.7.2.4-fr}
若 $A$ 是一个可容环，而 $\mathfrak J$ 包含在 $A$ 的某个定义
理想中，则 $A$ 关于 $\mathfrak J$-预进制拓扑分离且完备。
\end{proposition}

事实上，采用 (7.2.3) 的记号，$v\circ u$ 显然是 $A$ 的恒等映射。
另一方面，
$u_n\circ v:\widehat A\to A/\mathfrak J^n$ 是典范映射 $u_n$ 的
连续延拓（这里 $A$ 赋予 $\mathfrak J$-预进制拓扑，而
$A/\mathfrak J^n$ 赋予离散拓扑）；换言之，它是从
$\widehat A=\varprojlim_k A/\mathfrak J^k$ 到
$A/\mathfrak J^n$ 的典范映射；因此，$u\circ v$ 是这一映射序列的
逆极限，也就是按定义为 $\widehat A$ 的恒等映射；命题得证。

\begin{corollary}[7.2.5]
\label{0.7.2.5-fr}
在 (7.2.3) 的假设下，下列条件等价：
\begin{enumerate}
  \item[$a)$] 同态 $u$ 连续；

\oldpage[0\textsubscript{I}]{64}

  \item[$b)$] 同态 $v$ 是双连续的；
  \item[$c)$] $A$ 是一个 $\mathfrak J$-进制环。
\end{enumerate}
\end{corollary}

\begin{corollary}[7.2.6]
\label{0.7.2.6-fr}
设 $A$ 是一个可容环，$\mathfrak J$ 是 $A$ 的一个定义理想。
$A$ 为 Noether 环的充要条件是：$A/\mathfrak J$ 是 Noether 环，
并且 $\mathfrak J/\mathfrak J^2$ 是一个有限型
$(A/\mathfrak J)$-模。
\end{corollary}

这些条件显然是必要的。反过来，假设它们成立；由 (7.2.4)，$A$
关于 $\mathfrak J$-预进制拓扑完备，所以 $A$ 为 Noether 环的
充要条件是，与滤过 $\mathfrak J^n$ 相伴的分次环
$\operatorname{grad}(A)$ 为 Noether 环（[I]，第 18-07 页，定理 4）。
任取元素 $a_1,\ldots,a_n$，它们属于 $\mathfrak J$，并使它们模
$\mathfrak J^2$ 的类作为 $A/\mathfrak J$-模生成
$\mathfrak J/\mathfrak J^2$。立即由归纳法可知，对 $a_i$
$(1\leq i\leq n)$ 所作的总次数为 $m$ 的单项式模
$\mathfrak J^{m+1}$ 的类，组成 $(A/\mathfrak J)$-模
$\mathfrak J^m/\mathfrak J^{m+1}$ 的一个生成元组。由此可见，
$\operatorname{grad}(A)$ 同构于 $(A/\mathfrak J)[T_1,\ldots,T_n]$
的一个商环（各 $T_i$ 是不定元），证明完毕。

\begin{proposition}[7.2.7]
\label{0.7.2.7-fr}
设 $(A_i,u_{ij})$ 是一个离散环的逆系统 $(i\in\mathbf N)$；对每个
整数 $i$，设 $\mathfrak J_i$ 是同态
$u_{0i}:A_i\to A_0$ 在 $A_i$ 中的核。假设：
\begin{enumerate}
  \item[$a)$] 当 $i\leq j$ 时，$u_{ij}$ 为满射，其核为
    $\mathfrak J_j^{i+1}$（所以 $A_i$ 同构于
    $A_j/\mathfrak J_j^{i+1}$）。
  \item[$b)$] $\mathfrak J_1/\mathfrak J_1^2$（$=\mathfrak J_1$）
    是 $A_0=A_1/\mathfrak J_1$ 上的有限型模。
\end{enumerate}
令 $A=\varprojlim_i A_i$；对每个整数 $n\geq0$，设 $u_n$ 为典范同态
$A\to A_n$，并设 $\mathfrak J^{(n+1)}\subset A$ 为其核。在这些条件下：
\begin{enumerate}
  \item[(i)] $A$ 是一个 进制环，其定义理想为
    $\mathfrak J=\mathfrak J^{(1)}$。
  \item[(ii)] 对每个 $n\geq1$，有
    $\mathfrak J^{(n)}=\mathfrak J^n$。
  \item[(iii)] $\mathfrak J/\mathfrak J^2$ 同构于
    $\mathfrak J_1=\mathfrak J_1/\mathfrak J_1^2$，因而是
    $A_0=A/\mathfrak J$ 上的有限型模。
\end{enumerate}
\end{proposition}

各 $u_{ij}$ 为满射这一假设蕴涵 $u_n$ 为满射；此外，假设 $a)$
蕴涵 $\mathfrak J_j^{j+1}=0$，所以 $A$ 是一个可容环 (7.2.2)；
按定义，各 $\mathfrak J^{(n)}$ 组成 $A$ 中 $0$ 的一个邻域基，
所以 (ii) 蕴涵 (i)。此外，有
$\mathfrak J=\varprojlim_i\mathfrak J_i$，而各映射
$\mathfrak J\to\mathfrak J_i$ 都是满射，故 (ii) 蕴涵 (iii)，
从而只须证明 (ii)。按定义，$\mathfrak J^{(n)}$ 由 $A$ 中满足
$x_k=0$（当 $k<n$ 时）的元素 $(x_k)_{k\geq0}$ 组成，所以
$\mathfrak J^{(n)}\mathfrak J^{(m)}\subset\mathfrak J^{(n+m)}$；
换言之，各 $\mathfrak J^{(n)}$ 构成 $A$ 的一个\emph{滤过}。
另一方面，$\mathfrak J^{(n)}/\mathfrak J^{(n+1)}$ 同构于
$\mathfrak J^{(n)}$ 在 $A_n$ 中的投影；由于
$\mathfrak J^{(n)}=\varprojlim_{i>n}\mathfrak J_i^n$，这个投影正是
$\mathfrak J_n^n$，它是 $A_0=A_n/\mathfrak J_n$ 上的模。现取
$\mathfrak J=\mathfrak J^{(1)}$ 的 $r$ 个元素
$a_j=(a_{jk})_{k\geq0}$，使得 $a_{11},\ldots,a_{r1}$ 组成
$\mathfrak J_1$ 作为 $A_0$-模的一个生成元组；我们来证明，由各
$a_j$ 所作的总次数为 $n$ 的单项式集合 $S_n$ 生成 $A$ 的理想
$\mathfrak J^{(n)}$。由于 $\mathfrak J_i^{i+1}=0$，首先显然有
$S_n\subset\mathfrak J^{(n)}$；又因 $A$ 关于滤过
$(\mathfrak J^{(m)})$ 完备，只须证明 $S_n$ 中各元素模
$\mathfrak J^{(n+1)}$ 的类所成的集合 $\overline S_n$，在与上述滤过
相伴的分次环 $\operatorname{grad}(A)$ 上生成分次模
$\operatorname{grad}(\mathfrak J^{(n)})$（[I]，第 18-06 页，引理）；
根据 $\operatorname{grad}(A)$ 中乘法的定义，

\oldpage[0\textsubscript{I}]{65}

只须证明，对每个 $m$，$\overline S_m$ 是 $A_0$-模
$\mathfrak J^{(m)}/\mathfrak J^{(m+1)}$ 的一个生成元组，或者等价地，
$\mathfrak J_m^m$ 由关于 $a_{jm}$（$1\leq j\leq r$）的次数为 $m$
的单项式生成。为此，还须证明 $\mathfrak J_m$ 作为 $A_m$-模由
关于 $a_{jm}$ 的次数 $\leq m$ 的单项式生成。命题按定义对 $m=1$
显然成立；对 $m$ 作归纳，并设 $\mathfrak J'_m$ 为由这些单项式生成
的 $\mathfrak J_m$ 的 $A_m$-子模。关系
$\mathfrak J_{m-1}=\mathfrak J_m/\mathfrak J_m^m$ 与归纳假设说明
$\mathfrak J_m=\mathfrak J'_m+\mathfrak J_m^m$，由此结合
$\mathfrak J_m^{m+1}=0$，得到
$\mathfrak J_m^m={\mathfrak J'_m}^{\,m}$，最后得到
$\mathfrak J_m=\mathfrak J'_m$。

\begin{corollary}[7.2.8]
\label{0.7.2.8-fr}
在 (7.2.7) 的条件下，$A$ 为 Noether 环的充要条件是 $A_0$ 为
Noether 环。
\end{corollary}

这立即由 (7.2.6) 得到。

\begin{proposition}[7.2.9]
\label{0.7.2.9-fr}
假设 (7.2.7) 的假设成立：对每个整数 $i$，设 $M_i$ 是一个
$A_i$-模；当 $i\leq j$ 时，设
$v_{ij}:M_j\to M_i$ 是一个 $u_{ij}$-同态，使得 $(M_i,v_{ij})$
是一个逆系统。再假设 $M_0$ 是一个有限型 $A_0$-模，$v_{ij}$
为满射，并且其核为 $\mathfrak J_j^{i+1}M_j$。那么
$M=\varprojlim M_i$ 是一个有限型 $A$-模，而满的 $u_n$-同态
$v_n:M\to M_n$ 的核为 $\mathfrak J^{n+1}M$（从而把 $M_n$ 等同于
$M/\mathfrak J^{n+1}M=M\otimes_A(A/\mathfrak J^{n+1})$）。
\end{proposition}

取 $M$ 的 $s$ 个元素 $z_h=(z_{hk})_{k\geq0}$，使
$z_{h0}$（$1\leq h\leq s$）组成 $M_0$ 的一个生成元组；我们来证明
各 $z_h$ 生成 $A$-模 $M$。$A$-模 $M$ 关于由各 $M^{(n)}$ 构成的
滤过分离且完备，其中 $M^{(n)}$ 是 $M$ 中满足
$y_k=0$（当 $k<n$ 时）的元素 $y=(y_k)_{k\geq0}$ 的集合；显然有
$\mathfrak J^{(n)}M\subset M^{(n)}$，并且
$M^{(n)}/M^{(n+1)}=\mathfrak J_n^nM_n$。因此，只须证明各 $z_h$
% EGA1-ERR-0014: French print says modulo M^(0); the degree-zero classes require M^(1). Canonical translation preserves M^(0).
模 $M^{(0)}$ 的类在分次环 $\operatorname{grad}(A)$ 上生成与上述滤过
相伴的分次模 $\operatorname{grad}(M)$（[I]，第 18-06 页，引理）；
为此，不难看出，只须再证明各 $z_{hn}$（$1\leq h\leq s$）生成
$A_n$-模 $M_n$。再次对 $n$ 作归纳；按定义，命题对 $n=0$ 显然。
关系 $M_{n-1}=M_n/\mathfrak J_n^nM_n$ 和归纳假设说明：若 $M'_n$
是 $M_n$ 中由各 $z_{hn}$ 生成的子模，则
$M_n=M'_n+\mathfrak J_n^nM_n$；由于 $\mathfrak J_n$ 幂零，
这蕴涵 $M_n=M'_n$。同样的相伴分次模论证表明，从
$\mathfrak J^{(n)}M$ 到 $M^{(n)}$ 的典范映射是满射（因而是双射）；
换言之，$\mathfrak J^{(n)}M=\mathfrak J^nM$ 是 $M\to M_n$ 的核。

\begin{corollary}[7.2.10]
\label{0.7.2.10-fr}
设 $(N_i,w_{ij})$ 是另一个满足 (7.2.9) 条件的 $A_i$-模逆系统，
并设 $N=\varprojlim N_i$。$A_i$-同态的逆系统 $(h_i)$
$h_i:M_i\to N_i$ 与 $A$-模同态 $h:M\to N$ 一一对应（后者关于
$\mathfrak J$-进制拓扑必为连续映射）。
\end{corollary}

显然，若 $h:M\to N$ 是一个 $A$-同态，则
$h(\mathfrak J^nM)\subset\mathfrak J^nN$，故 $h$ 连续；取商后，
$h$ 因而对应于 $A_i$-同态 $h_i:M_i\to N_i$ 的一个逆系统，而
$h$ 是它的逆极限，推论得证。

\begin{remark}[7.2.11]
\label{0.7.2.11-fr}
设 $A$ 是一个 进制环，具有一个定义理想 $\mathfrak J$，并且
$\mathfrak J/\mathfrak J^2$ 是一个有限型 $(A/\mathfrak J)$-模；
显然，各 $A_i=A/\mathfrak J^{i+1}$ 满足

\oldpage[0\textsubscript{I}]{66}

(7.2.7) 的条件；由于 $A$ 是各 $A_i$ 的逆极限，可见命题 (7.2.7)
描述了所考察类型的\emph{所有} 进制环（特别是所有
\emph{Noether 进制环}）。
\end{remark}

\begin{example}[7.2.12]
\label{0.7.2.12-fr}
设 $B$ 是一个环，$\mathfrak J$ 是 $B$ 的一个理想，并且
$\mathfrak J/\mathfrak J^2$ 是 $B/\mathfrak J$ 上（等价地，$B$ 上）
的一个有限型模；令 $A=\varprojlim_n B/\mathfrak J^{n+1}$；$A$
是赋予 $\mathfrak J$-预进制拓扑的 $B$ 的分离完备化。若
$A_n=B/\mathfrak J^{n+1}$，则各 $A_n$ 显然满足 (7.2.7) 的条件；
因此，$A$ 是一个 进制环；若 $\overline{\mathfrak J}$ 是
$\mathfrak J$ 的典范像在 $A$ 中的闭包，则
$\overline{\mathfrak J}$ 是 $A$ 的一个定义理想，
$\overline{\mathfrak J^n}$ 是 $\mathfrak J^n$ 的典范像的闭包，
$A/\overline{\mathfrak J^n}$ 等同于 $B/\mathfrak J^n$，而
$\overline{\mathfrak J}/\overline{\mathfrak J^2}$ 作为
$(A/\overline{\mathfrak J})$-模同构于 $\mathfrak J/\mathfrak J^2$。
同样，若 $N$ 满足 $N/\mathfrak JN$ 是一个有限型 $B$-模，并令
$M_i=N/\mathfrak J^{i+1}N$，则 $M=\varprojlim M_i$ 是一个有限型
$A$-模，并同构于 $N$ 关于 $\mathfrak J$-预进制拓扑的分离完备化；
$\overline{\mathfrak J^n}M$ 等同于 $\mathfrak J^nN$ 的典范像的闭包，
而 $M/\overline{\mathfrak J^n}M$ 等同于 $N/\mathfrak J^nN$。
\end{example}

\subsection{Noether 预进制环}
\label{subsection:0.7.3-fr}

\begin{env}[7.3.1]
\label{0.7.3.1-fr}
设 $A$ 是一个环，$\mathfrak J$ 是 $A$ 的一个理想，$M$ 是一个
$A$-模；我们以
$\widehat A=\varprojlim_n A/\mathfrak J^n$（相应地，
% EGA1-ERR-0015: French print and diplomatic TeX use varprojlim_u here; the inverse system is indexed by n. Canonical translation preserves u.
$\widehat M=\varprojlim_u M/\mathfrak J^nM$）表示 $A$（相应地，$M$）
关于 $\mathfrak J$-预进制拓扑的分离完备化。设
\[
  M'\xrightarrow{u}M\xrightarrow{v}M''\to0
\]
是一个 $A$-模正合列；由于
$M/\mathfrak J^nM=M\otimes_A(A/\mathfrak J^n)$，序列
\[
  M'/\mathfrak J^nM'\xrightarrow{u_n}
  M/\mathfrak J^nM\xrightarrow{v_n}
  M''/\mathfrak J^nM''\to0
\]
对每个 $n$ 都正合。此外，由于
$v(\mathfrak J^nM)=\mathfrak J^nv(M)=\mathfrak J^nM''$，
$\widehat v=\varprojlim v_n$ 为满射（Bourbaki，
\emph{Top. gén.}，第 IX 章，第 2 版，第 60 页，推论 2）。
另一方面，若 $z=(z_k)$ 是 $\widehat v$ 的核中的一个元素，则对每个
整数 $k$，存在 $z'_k\in M'/\mathfrak J^kM'$ 使得
$u_k(z'_k)=z_k$；由此可知，存在
$z'=(z'_n)\in\widehat{M'}$，使 $\widehat u(z')$ 的前 $k$ 个分量
与 $z$ 的相应分量相同；换言之，$\widehat{M'}$ 在 $\widehat u$
下的像在 $\widehat v$ 的核中\emph{稠密}。

若再假设 $A$ \emph{是 Noether 环}，则由 (7.2.12)，$\widehat A$
也是 Noether 环，因为此时 $\mathfrak J/\mathfrak J^2$ 是有限型
$A$-模。此外：
\end{env}

\begin{namedtheorem}[7.3.2]{Krull 定理}
\label{0.7.3.2-fr}
设 $A$ 是一个 Noether 环，$\mathfrak J$ 是 $A$ 的一个理想，$M$
是一个有限型 $A$-模，$M'$ 是 $M$ 的一个子模；那么，$M$ 的
$\mathfrak J$-预进制拓扑在 $M'$ 上诱导的拓扑，与 $M'$ 的
$\mathfrak J$-预进制拓扑相同。
\end{namedtheorem}

这立即由下述引理得到：

\begin{namedlemma}[7.3.2.1]{Artin--Rees 引理}
\label{0.7.3.2.1-fr}
在 (7.3.2) 的假设下，存在一个整数 $p$，使得当 $n\geq p$ 时，
\[
  M'\cap\mathfrak J^nM=\mathfrak J^{n-p}
  (M'\cap\mathfrak J^pM).
\]
\end{namedlemma}

证明见（[I]，第 2-04 页）。

\oldpage[0\textsubscript{I}]{67}

\begin{corollary}[7.3.3]
\label{0.7.3.3-fr}
在 (7.3.2) 的假设下，典范映射
$M\otimes_A\widehat A\to\widehat M$ 是双射；在有限型 $A$-模范畴中，
函子 $M\otimes_A\widehat A$ 关于 $M$ 是正合的；因此，
$\mathfrak J$-进制 分离完备化 $\widehat A$ 是一个平坦 $A$-模
(6.1.1)。
\end{corollary}

先注意，在有限型 $A$-模范畴中，$\widehat M$ 关于 $M$ 是一个
\emph{正合}函子。事实上，设
\[
  0\to M'\xrightarrow{u}M\xrightarrow{v}M''\to0
\]
是一个正合列；我们已经知道
$\widehat v:\widehat M\to\widehat{M''}$ 是满射 (7.3.1)。另一方面，
若 $i$ 是典范同态 $M\to\widehat M$，则由 Krull 定理，$i(u(M'))$
在 $\widehat M$ 中的闭包等同于 $M'$ 关于 $\mathfrak J$-预进制
拓扑的分离完备化；因此，$\widehat u$ 是单射，而且由 (7.3.1)，
$\widehat u$ 的像等于 $\widehat v$ 的核。

现在，典范映射 $M\otimes_A\widehat A\to\widehat M$ 是对各映射
\[
  M\otimes_A\widehat A\to
  M\otimes_A(A/\mathfrak J^n)=M/\mathfrak J^nM.
\]
取逆极限得到的。当 $M$ 形如 $A^p$ 时，这个映射显然是双射。
若 $M$ 是一个有限型 $A$-模，则存在正合列
$A^p\to A^q\to M\to0$。函子 $M\otimes_A\widehat A$ 和
$\widehat M$ 在有限型 $A$-模范畴中关于 $M$ 都\emph{右}正合，
因而得到交换图
\[
\begin{tikzcd}[column sep=large,row sep=large]
  A^p\otimes_A\widehat A \arrow[r] \arrow[d] &
  A^q\otimes_A\widehat A \arrow[r] \arrow[d] &
  M\otimes_A\widehat A \arrow[r] \arrow[d] &
  0 \\
  \widehat A^{\,p} \arrow[r] &
  \widehat A^{\,q} \arrow[r] &
  \widehat M \arrow[r] &
  0
\end{tikzcd}
\]
其中两行都正合，前两个竖直箭头都是同构；结论立即得到。

\begin{corollary}[7.3.4]
\label{0.7.3.4-fr}
设 $A$ 是一个 Noether 环，$\mathfrak J$ 是 $A$ 的一个理想，$M$、
$N$ 是两个有限型 $A$-模；存在函子性的典范同构
\[
  (M\otimes_A N)^\wedge
  \xrightarrow{\sim}
  \widehat M\otimes_{\widehat A}\widehat N,
  \qquad
  \bigl(\operatorname{Hom}_A(M,N)\bigr)^\wedge
  \xrightarrow{\sim}
  \operatorname{Hom}_{\widehat A}(\widehat M,\widehat N).
\]
\end{corollary}

这由 (7.3.3)、(6.2.1) 和 (6.2.2) 得到。

\begin{corollary}[7.3.5]
\label{0.7.3.5-fr}
设 $A$ 是一个 Noether 环，$\mathfrak J$ 是 $A$ 的一个理想。下列
条件等价：
\begin{enumerate}
  \item[$a)$] $\mathfrak J$ 包含在 $A$ 的 Jacobson 根中。
  \item[$b)$] $\widehat A$ 是一个忠实平坦 $A$-模 (6.4.1)。
  \item[$c)$] 每个有限型 $A$-模关于 $\mathfrak J$-预进制拓扑
  都是分离的。
  \item[$d)$] 有限型 $A$-模的每个子模关于 $\mathfrak J$-预进制
  拓扑都是闭的。
\end{enumerate}
\end{corollary}

由于 $\widehat A$ 是一个平坦 $A$-模，条件 $b)$ 和 $c)$ 等价；
事实上，$b)$ 等价于说：若 $M$ 是一个有限型 $A$-模，则典范映射
$M\to\widehat M=M\otimes_A\widehat A$ 是单射 (6.6.1, $c)$。
$c)$ 显然蕴涵 $d)$，因为若 $N$ 是有限型 $A$-模 $M$ 的一个子模，
则 $M/N$ 关于 $\mathfrak J$-预进制拓扑分离，故 $N$ 在 $M$ 中
是闭的。来证明 $d)$ 蕴涵 $a)$：若 $\mathfrak m$ 是 $A$ 的一个
极大理想，则 $\mathfrak m$ 关于 $A$ 的 $\mathfrak J$-预进制拓扑
是闭的，所以
\[
  \mathfrak m=\bigcap_{p\geq0}(\mathfrak m+\mathfrak J^p),
\]
而 $\mathfrak m+\mathfrak J^p$ 必等于 $A$ 或 $\mathfrak m$；因此，
当 $p$ 足够

\oldpage[0\textsubscript{I}]{68}

大时，$\mathfrak m+\mathfrak J^p=\mathfrak m$，从而
$\mathfrak J^p\subset\mathfrak m$；又因 $\mathfrak m$ 是素理想，
所以 $\mathfrak J\subset\mathfrak m$。最后，$a)$ 蕴涵 $b)$：
事实上，设 $P$ 是有限型 $A$-模 $M$ 中 $\{0\}$ 关于
$\mathfrak J$-预进制拓扑的闭包；由 Krull 定理 (7.3.2)，$M$ 的
$\mathfrak J$-预进制拓扑在 $P$ 上诱导的拓扑就是 $P$ 的
$\mathfrak J$-预进制拓扑，所以 $\mathfrak JP=P$；由于 $P$ 是
有限型模，Nakayama 引理给出 $P=0$（$\mathfrak J$ 包含在 $A$
的 Jacobson 根中）。

注意，当 $A$ 是一个\emph{Noether 局部环}，而
$\mathfrak J\neq A$ 是 $A$ 的任意理想时，(7.3.5) 的各条件都成立。

\begin{corollary}[7.3.6]
\label{0.7.3.6-fr}
若 $A$ 是一个 Noether $\mathfrak J$-进制环，则每个有限型 $A$-模
关于 $\mathfrak J$-预进制拓扑都是分离且完备的。
\end{corollary}

此时 $\widehat A=A$，所以这立即由 (7.3.3) 得到。

由此可见，命题 (7.2.9) 描述了 Noether 进制环上的\emph{所有}
有限型模。

\begin{corollary}[7.3.7]
\label{0.7.3.7-fr}
在 (7.3.2) 的假设下，典范映射
$M\to\widehat M=M\otimes_A\widehat A$ 的核，是被 $1+\mathfrak J$
中的某个元素消去的全体 $x\in M$。
\end{corollary}

事实上，$x\in M$ 属于该核的充要条件是：子模 $Ax$ 的分离完备化
等于 0（由 (7.3.2)）；换言之，$x\in\mathfrak Jx$。

\subsection{局部环上的拟有限模}
\label{subsection:0.7.4-fr}

\begin{definition}[7.4.1]
\label{0.7.4.1-fr}
给定局部环 $A$，其极大理想为 $\mathfrak m$；若 $A$-模 $M$ 的
$M/\mathfrak mM$ 是剩余域 $k=A/\mathfrak m$ 上的有限维向量空间，
则称该模（在 $A$ 上）是拟有限的。
\end{definition}

当 $A$ 为 Noether 环时，$M$ 关于 $\mathfrak m$-预进制拓扑的
分离完备化 $\widehat M$ 是一个有限型 $\widehat A$-模；事实上，
此时 $\mathfrak m/\mathfrak m^2$ 是有限型 $A$-模，故该结论由
(7.2.12) 以及关于 $M/\mathfrak mM$ 的假设得到。

特别地，若再假设 $A$ 是\emph{完备的}，并且 $M$ 关于
$\mathfrak m$-预进制拓扑是\emph{分离的}（换言之，
$\bigcap_n\mathfrak m^nM=0$），则 $M$ 本身是有限型 $A$-模：
事实上，$\widehat M$ 是有限型 $A$-模，而 $M$ 可等同于
$\widehat M$ 的一个子模，故 $M$ 也为有限型（并且由 (7.3.6)，
它其实等同于自身的完备化）。

\begin{proposition}[7.4.2]
\label{0.7.4.2-fr}
设 $A$、$B$ 是两个局部环，$\mathfrak m$、$\mathfrak n$ 分别是
它们的极大理想，并假设 $B$ 为 Noether 环。设
$\varphi:A\to B$ 是局部同态，$M$ 是有限型 $B$-模。若 $M$ 是
拟有限 $A$-模，则 $M$ 上的 $\mathfrak m$-预进制拓扑与
$\mathfrak n$-预进制拓扑相同，因而是分离的。
\end{proposition}

注意，按假设，$M/\mathfrak mM$ 作为 $A$-模具有\emph{有限长度}，
所以作为 $B$-模当然也具有有限长度。由此可知，$\mathfrak n$ 是
$B$ 中\emph{唯一包含} $M/\mathfrak mM$ 的零化子的\emph{素理想}：
事实上，由 (I.7.4) 和 (I.7.2)，立即可约化到
$M/\mathfrak mM$ 是\emph{单模}的情形；这时它必同构于
$B/\mathfrak n$，而所述断言显然成立。另一方面，因为 $M$ 是有限型
$B$-模，包含 $M/\mathfrak mM$ 的零化子的素理想，正是包含
$\mathfrak mB+\mathfrak b$ 的素理想，其中 $\mathfrak b$ 表示
$B$-模 $M$ 的零化子 (I.7.5)。由于 $B$ 为 Noether 环，因而
由 ([III]，第 127 页，推论 4) 可知，$\mathfrak mB+\mathfrak b$
是 $B$ 的一个

\oldpage[0\textsubscript{I}]{69}

定义理想；换言之，存在 $k>0$ 使得
$\mathfrak n^k\subset\mathfrak mB+\mathfrak b\subset\mathfrak n$。
于是，对每个 $h>0$，
\[
  \mathfrak n^{hk}M
  \subset(\mathfrak mB+\mathfrak b)^hM
  =\mathfrak m^hM
  \subset\mathfrak n^hM,
\]
这就证明 $M$ 上的 $\mathfrak m$-预进制拓扑与
$\mathfrak n$-预进制拓扑相同；而由 (7.3.5)，后一拓扑是分离的。

\begin{corollary}[7.4.3]
\label{0.7.4.3-fr}
在 (7.4.2) 的假设下，若 $A$ 还为 Noether 环并且关于
$\mathfrak m$-预进制拓扑完备，则 $M$ 是有限型 $A$-模。
\end{corollary}

事实上，此时 $M$ 关于 $\mathfrak m$-预进制拓扑分离，断言由
(7.4.1) 中的说明得到。

\begin{env}[7.4.4]
\label{0.7.4.4-fr}
(7.4.2) 最重要的应用情形，是 $B$ 本身为一个拟有限 $A$-模；
这等价于说，$B/\mathfrak mB$ 是 $k=A/\mathfrak m$ 上的
\emph{有限维代数}。由前述结果，这一条件还可分解为以下两个条件
同时成立：
\begin{enumerate}
  \item[(i)] $\mathfrak mB$ \emph{是 $B$ 的一个定义理想}；
  \item[(ii)] $B/\mathfrak n$ \emph{是域 $A/\mathfrak m$ 的
  有限维扩张}。
\end{enumerate}
在这种情形下，每个有限型 $B$-模显然都是拟有限 $A$-模。
\end{env}

\begin{corollary}[7.4.5]
\label{0.7.4.5-fr}
在 (7.4.2) 的假设下，若 $\mathfrak b$ 是 $B$-模 $M$ 的零化子，
则 $B/\mathfrak b$ 是拟有限 $A$-模。
\end{corollary}

假设 $M\neq0$（否则推论显然成立）。可以把 $M$ 看作局部 Noether 环
$B/\mathfrak b$ 上的模；这时它的零化子约化为 $0$，而 (7.4.2) 的
证明表明，$\mathfrak m(B/\mathfrak b)$ 是 $B/\mathfrak b$ 的一个
定义理想。另一方面，$M/\mathfrak nM$ 是 $A/\mathfrak m$ 上的
有限维向量空间，因为它是 $M/\mathfrak mM$ 的一个商，而后者按假设
在 $A/\mathfrak m$ 上有限维。又因为 $M\neq0$，Nakayama 引理给出
$M\neq\mathfrak nM$。由于 $M/\mathfrak nM$ 是 $B/\mathfrak n$
上的向量空间 $\neq0$，它在 $A/\mathfrak m$ 上有限维这一事实蕴涵
$B/\mathfrak n$ 在 $A/\mathfrak m$ 上也有限维；因此，对环
$B/\mathfrak b$ 应用 (7.4.4) 即得结论。

\subsection{设限形式幂级数环}
\label{subsection:0.7.5-fr}

\begin{env}[7.5.1]
\label{0.7.5.1-fr}
设 $A$ 是线性拓扑化、分离且完备的拓扑环；设
$(\mathfrak J_\lambda)$ 是由（开）理想组成的 $A$ 中 $0$ 的一个
基本邻域系，使得 $A$ 典范地等同于
$\varprojlim A/\mathfrak J_\lambda$ (7.2.1)。对每个 $\lambda$，令
$B_\lambda=(A/\mathfrak J_\lambda)[T_1,\ldots,T_r]$，其中各 $T_i$
是不定元；显然，各 $B_\lambda$ 构成离散环的一个逆系统。我们记
$\varprojlim B_\lambda=A\{T_1,\ldots,T_r\}$；下面将说明，这个拓扑环
与所选的基本理想系 $(\mathfrak J_\lambda)$ 无关。更确切地，令 $A'$
为形式幂级数环 $A[\![T_1,\ldots,T_r]\!]$ 的如下子环：它由满足
$\lim_\alpha c_\alpha=0$ 的形式幂级数
$\sum_\alpha c_\alpha T^\alpha$ 组成，其中
$\alpha=(\alpha_1,\ldots,\alpha_r)\in\mathbf{N}^r$，该极限沿
$\mathbf{N}^r$ 的有限子集之补集所成的滤子取得。我们称这些幂级数为
以 $A$ 中元素为系数、关于各 $T_i$ 的\emph{设限形式幂级数}。

\oldpage[0\textsubscript{I}]{70}

对 $A$ 中 $0$ 的每个邻域 $V$，令 $V'$ 为所有满足下述条件的
$x=\sum_\alpha c_\alpha T^\alpha\in A'$ 所成的集合：对每个 $\alpha$，都有
$c_\alpha\in V$。立即可验证，各 $V'$ 构成 $0$ 的一个基本邻域系，
并在 $A'$ 上定义一个\emph{分离的}环拓扑。下面将典范地定义从环
$A\{T_1,\ldots,T_r\}$ 到 $A'$ 的一个\emph{拓扑同构}。对每个
$\alpha\in\mathbf{N}^r$ 和每个 $\lambda$，令
$\varphi_{\lambda,\alpha}$ 为从
$(A/\mathfrak J_\lambda)[T_1,\ldots,T_r]$ 到
$A/\mathfrak J_\lambda$ 的映射，它把前一环中的每个多项式映到该
多项式中 $T^\alpha$ 的系数。显然，各 $\varphi_{\lambda,\alpha}$
构成 $(A/\mathfrak J_\lambda)$-模同态的一个逆系统，其逆极限是连续同态
$\varphi_\alpha:A\{T_1,\ldots,T_r\}\to A$。我们来证明，对每个
$y\in A\{T_1,\ldots,T_r\}$，形式幂级数
$\sum_\alpha\varphi_\alpha(y)T^\alpha$ 是\emph{设限的}。事实上，
若 $y_\lambda$ 是 $y$ 在 $B_\lambda$ 中的分量，并以 $H_\lambda$
表示使多项式 $y_\lambda$ 的系数非零的全体
$\alpha\in\mathbf{N}^r$ 所成的有限集，则当
$\mathfrak J_\mu\subset\mathfrak J_\lambda$ 且
$\alpha\notin H_\lambda$ 时，有
$\varphi_{\lambda,\alpha}(y_\mu)\in\mathfrak J_\lambda$；取极限得到，
当 $\alpha\notin H_\lambda$ 时，
$\varphi_\alpha(y)\in\mathfrak J_\lambda$。因此，令
$\varphi(y)=\sum_\alpha\varphi_\alpha(y)T^\alpha$，便定义了环同态
$\varphi:A\{T_1,\ldots,T_r\}\to A'$；$\varphi$ 显然连续。反过来，若
$\theta_\lambda$ 是典范同态 $A\to A/\mathfrak J_\lambda$，则对每个
$z=\sum_\alpha c_\alpha T^\alpha\in A'$ 以及每个 $\lambda$，只有
有限多个指标 $\alpha$ 满足 $\theta_\lambda(c_\alpha)\neq0$；因而
$\psi_\lambda(z)=\sum_\alpha\theta_\lambda(c_\alpha)T^\alpha$ 属于
$B_\lambda$。各 $\psi_\lambda$ 连续，并构成同态的一个逆系统；其逆极限
是连续同态 $\psi:A'\to A\{T_1,\ldots,T_r\}$。最后只须验证
$\varphi\circ\psi$ 和 $\psi\circ\varphi$ 都是恒等自同构，而这显然成立。
\end{env}

\begin{env}[7.5.2]
\label{0.7.5.2-fr}
我们借助 (7.5.1) 中定义的同构，把 $A\{T_1,\ldots,T_r\}$ 等同于
设限形式幂级数环 $A'$。各典范同构
\[
  ((A/\mathfrak J_\lambda)[T_1,\ldots,T_r])
  [T_{r+1},\ldots,T_s]
  \xrightarrow{\sim}
  (A/\mathfrak J_\lambda)[T_1,\ldots,T_s]
\]
在取逆极限后定义一个典范同构
% EGA1-ERR-0016: French print/source has an unmatched opening delimiter; preserved literally pending canon disposition.
\[
  \Bigl[(A\{T_1,\ldots,T_r\})\{T_{r+1},\ldots,T_s\}
  \xrightarrow{\sim}
  A\{T_1,\ldots,T_s\}
\]
\end{env}

\begin{env}[7.5.3]
\label{0.7.5.3-fr}
设 $u:A\to B$ 是从 $A$ 到线性拓扑化、分离且完备的环 $B$ 的任意
连续同态，并且 $(b_1,\ldots,b_r)$ 是 $B$ 中任意 $r$ 个元素所成的系；
则存在\emph{唯一的连续同态}
$\overline u:A\{T_1,\ldots,T_r\}\to B$，使得对每个 $a\in A$ 都有
$\overline u(a)=u(a)$，并且当 $1\leq j\leq r$ 时，
$\overline u(T_j)=b_j$。事实上，只须取
\[
  \overline u\Bigl(\sum_\alpha c_\alpha T^\alpha\Bigr)
  =\sum_\alpha u(c_\alpha)b_1^{\alpha_1}\cdots b_r^{\alpha_r}~;
\]
容易验证，族 $(u(c_\alpha)b_1^{\alpha_1}\cdots b_r^{\alpha_r})$
在 $B$ 中可求和，并且 $\overline u$ 连续；这些验证留给读者。应当注意，
这一性质（其中 $B$ 与各 $b_j$ 任意）在相差唯一同构的意义下
\emph{刻画}拓扑环 $A\{T_1,\ldots,T_r\}$。
\end{env}

\begin{proposition}[7.5.4]
\label{0.7.5.4-fr}
\begin{enumerate}
  \item[(i)] 若 $A$ 是可容环，则
  $A'=A\{T_1,\ldots,T_r\}$ 也是可容环。
  \item[(ii)] 设 $A$ 是进制环，$\mathfrak J$ 是 $A$ 的一个定义理想，
  并且 $\mathfrak J/\mathfrak J^2$ 是有限型

\oldpage[0\textsubscript{I}]{71}

  $A/\mathfrak J$-模。若令 $\mathfrak J'=\mathfrak JA'$，则 $A'$ 是
  $\mathfrak J'$-进制环，并且
  $\mathfrak J'/\mathfrak J'^{\,2}$ 是有限型
  $A'/\mathfrak J'$-模。若 $A$ 还是 Noether 环，则 $A'$ 也是
  Noether 环。
\end{enumerate}
\end{proposition}

\begin{enumerate}
  \item[(i)] 若 $\mathfrak J$ 是 $A$ 的一个理想，并且 $\mathfrak J'$
  是 $A'$ 中所有满足下述条件的 $\sum_\alpha c_\alpha T^\alpha$
  所成的理想：对每个 $\alpha$，都有 $c_\alpha\in\mathfrak J$；
  则 $(\mathfrak J')^n\subset(\mathfrak J^n)'$。所以，若 $\mathfrak J$
  是 $A$ 的定义理想，则 $\mathfrak J'$ 是 $A'$ 的定义理想。

  \item[(ii)] 令 $A_i=A/\mathfrak J^{i+1}$；当 $i\leq j$ 时，令
  $u_{ij}$ 为典范同态
  $A/\mathfrak J^{j+1}\to A/\mathfrak J^{i+1}$。再令
  $A'_i=A_i[T_1,\ldots,T_r]$；并令 $u'_{ij}$ 为把 $u_{ij}$ 施于
  $A'_j$ 各多项式系数所得到的同态 $A'_j\to A'_i$（$i\leq j$）。
  我们来证明逆系统 $(A'_i,u'_{ij})$ 满足 (7.2.7) 的条件。由于
  $\mathfrak J'$ 是 $A'\to A'_0$ 的核，这将证明 (ii) 的第一个断言。
  显然，各 $u'_{ij}$ 都是满射。$u'_{0i}$ 的核 $\mathfrak J'_i$
  是 $A_i[T_1,\ldots,T_r]$ 中所有系数均属于
  $\mathfrak J/\mathfrak J^{i+1}$ 的多项式所成的集合；特别地，
  $\mathfrak J'_1$ 是 $A_1[T_1,\ldots,T_r]$ 中所有系数均属于
  $\mathfrak J/\mathfrak J^2$ 的多项式所成的集合。因为
  $\mathfrak J/\mathfrak J^2$ 是有限型 $A_1=A/\mathfrak J^2$-模，
  所以 $\mathfrak J'_1/\mathfrak J_1'^{\,2}$ 是有限型 $A'_1$-模
  （等价地，是有限型 $A'_0=A'_1/\mathfrak J'_1$-模）。我们来证明
  $u'_{ij}$ 的核为 $\mathfrak J_j'^{\,i+1}$。显然，
  $\mathfrak J_j'^{\,i+1}$ 包含在该核中。另一方面，设
  $a_1,\ldots,a_m$ 是 $\mathfrak J$ 中的一些元素，其模
  $\mathfrak J^2$ 的类生成 $\mathfrak J/\mathfrak J^2$。立即可验证，
  各 $a_k$（$1\leq k\leq m$）中次数 $\leq j$ 的单项式模
  $\mathfrak J^{j+1}$ 的类生成 $\mathfrak J/\mathfrak J^{j+1}$；
  所以，次数 $>i$ 且 $\leq j$ 的单项式的类生成
  $\mathfrak J^{i+1}/\mathfrak J^{j+1}$。因此，以这样的元素为系数的
  各 $T_k$ 中的单项式，是 $\mathfrak J'_j$ 中 $i+1$ 个元素的乘积，
  从而证明了所述断言。最后，若 $A$ 为 Noether 环，则
  $A'/\mathfrak J'=(A/\mathfrak J)[T_1,\ldots,T_r]$ 也是 Noether 环，
  因而 $A'$ 为 Noether 环 (7.2.8)。
\end{enumerate}

\begin{proposition}[7.5.5]
\label{0.7.5.5-fr}
设 $A$ 是 Noether $\mathfrak J$-进制环，$B$ 是可容拓扑环，
$\varphi:A\to B$ 是使 $B$ 成为 $A$-代数的连续同态。下列条件等价：
\begin{enumerate}
  \item[a)] $B$ 是 Noether $\mathfrak JB$-进制环，并且
  $B/\mathfrak JB$ 是有限型 $A/\mathfrak J$-代数。

  \item[b)] $B$ 在拓扑意义下 $A$-同构于 $\varprojlim B_n$，其中当
  $m\geq n$ 时，$B_n=B_m/\mathfrak J^{n+1}B_m$，并且 $B_1$ 是
  有限型 $A_1=A/\mathfrak J^2$-代数。

  \item[c)] $B$ 在拓扑意义下 $A$-同构于形如
  \linebreak $A\{T_1,\ldots,T_r\}$ 的代数关于某个理想的商
  （由 (7.3.6) 和 (7.5.4 (ii))，该理想必为闭理想）。
\end{enumerate}
\end{proposition}

因为 $A$ 为 Noether 环，$A'=A\{T_1,\ldots,T_r\}$ 也为 Noether 环
(7.5.4)，所以 c) 蕴涵 $B$ 为 Noether 环。又因为
$\mathfrak J'=\mathfrak JA'$ 是 $A'$ 中 0 的一个开邻域，并且各
$\mathfrak J'^{n}$ 构成 0 的一个基本邻域系，所以各
$\mathfrak J'^{n}$ 的像 $\mathfrak J^nB$ 构成 $B$ 中 0 的一个基本邻域系；而 $B$ 分离且完备，
故 $B$ 是 $\mathfrak JB$-进制环。最后，$B/\mathfrak JB$ 是
$A/\mathfrak J$ 上的代数，并且是
$A'/\mathfrak JA'=(A/\mathfrak J)[T_1,\ldots,T_r]$ 的商，因而为
有限型代数。这就证明了 c) 蕴涵 a)。

若 $B$ 为 Noether $\mathfrak JB$-进制环，则 $B$ 同构于
$\varprojlim B_n$，其中 $B_n=B/\mathfrak J^{n+1}B$ (7.2.11)，并且
$\mathfrak JB/\mathfrak J^2B$ 是有限型 $B/\mathfrak JB$-模。设
$(a_j)_{1\leq j\leq s}$ 是 $(B/\mathfrak JB)$-模
$\mathfrak JB/\mathfrak J^2B$ 的一个生成元系；并设
$(c_i)_{1\leq i\leq r}$ 是 $B/\mathfrak J^2B$ 中的一族元素，其类

\oldpage[0\textsubscript{I}]{72}

模 $\mathfrak JB/\mathfrak J^2B$ 构成 $(A/\mathfrak J)$-代数
$B/\mathfrak JB$ 的一个生成元系。立即可见，
% EGA1-ERR-0017: French print/source has c_i a_j; retained literally pending canon disposition.
各 $c_ia_j$ 构成 $(A/\mathfrak J^2)$-代数 $B/\mathfrak J^2B$ 的
一个生成元系，故 a) 蕴涵 b)。

还须证明 b) 蕴涵 c)。该假设蕴涵 $B_1$ 是 Noether 环；又因为
$B_1=B_2/\mathfrak J^2B_2$，所以 $\mathfrak J^2B_1=0$，从而
$\mathfrak JB_1=\mathfrak JB_1/\mathfrak J^2B_1$ 是有限型
$B_0$-模。因此，逆系统 $(B_n)$ 满足 (7.2.7) 的条件，而 $B$ 是
$\mathfrak JB$-进制环。设 $(c_i)_{1\leq i\leq r}$ 是 $B$ 中的一个
有限元素系，其模 $\mathfrak JB$ 的类生成 $(A/\mathfrak J)$-代数
$B/\mathfrak JB$；并且，使这些元素的、系数属于 $\mathfrak J$ 的
线性组合满足：它们模 $\mathfrak J^2B$ 的类生成 $B_0$-模
$\mathfrak JB/\mathfrak J^2B$。存在从
$A'=A\{T_1,\ldots,T_r\}$ 到 $B$ 的连续 $A$-同态 $u$，它在 $A$
上约化为 $\varphi$，并且当 $1\leq i\leq r$ 时，$u(T_i)=c_i$
(7.5.3)。若证明 $u$ 为\emph{满射}，则 c) 成立；因为从
$u(A')=B$ 可推出 $u(\mathfrak J^nA')=\mathfrak J^nB$。换言之，
$u$ 是拓扑环的严格态射，所以 $B$ 同构于 $A'$ 关于某个闭理想的商。
由于 $B$ 关于 $\mathfrak JB$-进制拓扑完备，只须证明由 $u$ 对
$A'$ 与 $B$ 上的 $\mathfrak J$-进制滤过典范诱导的同态
$\operatorname{grad}(A')\to\operatorname{grad}(B)$ 为满射
([I]，第 18-07 页)。可是，按定义，由 $u$ 诱导的同态
$A'/\mathfrak JA'\to B/\mathfrak JB$ 和
$\mathfrak JA'/\mathfrak J^2A'\to\mathfrak JB/\mathfrak J^2B$
都是满射；对 $n$ 作归纳，立即可知
$\mathfrak JA'/\mathfrak J^nA'\to\mathfrak JB/\mathfrak J^nB$
也为满射，因而更有
$\mathfrak J^nA'/\mathfrak J^{n+1}A'\to
\mathfrak J^nB/\mathfrak J^{n+1}B$ 为满射。证明完毕。

\subsection{完备分式环}
\label{subsection:0.7.6-fr}

\begin{env}[7.6.1]
\label{0.7.6.1-fr}
设 $A$ 是一个线性拓扑化环，$(\mathfrak J_\lambda)$ 是 $A$ 中由理想
组成的一个 0 的基本邻域系，$S$ 是 $A$ 的一个乘法子集。设
$u_\lambda$ 为典范同态
$A\to A_\lambda=A/\mathfrak J_\lambda$；当
$\mathfrak J_\mu\subset\mathfrak J_\lambda$ 时，设
$u_{\lambda\mu}$ 为典范同态 $A_\mu\to A_\lambda$。令
$S_\lambda=u_\lambda(S)$，于是
$u_{\lambda\mu}(S_\mu)=S_\lambda$。各 $u_{\lambda\mu}$ 典范地给出
满同态
$S_\mu^{-1}A_\mu\to S_\lambda^{-1}A_\lambda$，这些环由此构成一个
逆系统；以 $A\{S^{-1}\}$ 表示该系统的逆极限。这个定义与所选的
基本邻域系 $(\mathfrak J_\lambda)$ 无关；事实上：
\end{env}

\begin{proposition}[7.6.2]
\label{0.7.6.2-fr}
环 $A\{S^{-1}\}$ 拓扑同构于环 $S^{-1}A$ 关于下述拓扑的分离
完备化：该拓扑的一个 0 的基本邻域系由各
$S^{-1}\mathfrak J_\lambda$ 组成。
\end{proposition}

事实上，若 $v_\lambda$ 是由 $u_\lambda$ 诱导的典范同态
$S^{-1}A\to S_\lambda^{-1}A_\lambda$，则 $v_\lambda$ 的核为
$S^{-1}\mathfrak J_\lambda$，且 $v_\lambda$ 为满射；命题由
(7.2.1) 得到。

\begin{corollary}[7.6.3]
\label{0.7.6.3-fr}
若 $S'$ 是 $S$ 在 $A$ 的分离完备化 $\widehat A$ 中的典范像，则
$A\{S^{-1}\}$ 典范地等同于
$\widehat A\{{S'}^{-1}\}$。
\end{corollary}

应当注意，即使 $A$ 分离且完备，$S^{-1}A$ 对于由各
$S^{-1}\mathfrak J_\lambda$ 定义的拓扑也未必如此。例如，取 $S$
为各 $f^n$（$n\geq0$）所成的集合，其中 $f$ 拓扑幂零但并不幂零：
事实上，$S^{-1}A$ 不为零环；另一方面，对每个 $\lambda$ 都存在
$n$ 使得 $f^n\in\mathfrak J_\lambda$，因而
$1=f^n/f^n\in S^{-1}\mathfrak J_\lambda$，且
$S^{-1}\mathfrak J_\lambda=S^{-1}A$。

\oldpage[0\textsubscript{I}]{73}

\begin{corollary}[7.6.4]
\label{0.7.6.4-fr}
若在 $A$ 中 $0$ 不在 $S$ 的闭包中，则环 $A\{S^{-1}\}$ 不等于 $0$。
\end{corollary}

事实上，在环 $S^{-1}A$ 中，$0$ 不在 $\{1\}$ 的闭包中；否则，对
$A$ 的每个开理想 $\mathfrak J_\lambda$ 都有
$1\in S^{-1}\mathfrak J_\lambda$，从而对每个 $\lambda$ 都有
$\mathfrak J_\lambda\cap S\neq\varnothing$，这与假设矛盾。

\begin{env}[7.6.5]
\label{0.7.6.5-fr}
我们称 $A\{S^{-1}\}$ 为 $A$ 的、分母属于 $S$ 的
\emph{完备分式环}。沿用上述记号，$S^{-1}\mathfrak J_\lambda$
在 $A$ 中的逆像显然包含 $\mathfrak J_\lambda$，故典范映射
$A\to S^{-1}A$ 连续；将它与典范映射
$S^{-1}A\to A\{S^{-1}\}$ 复合，便得到一个典范连续同态
$A\to A\{S^{-1}\}$，它是各同态
$A\to S_\lambda^{-1}A_\lambda$ 的逆极限。
\end{env}

\begin{env}[7.6.6]
\label{0.7.6.6-fr}
由 $A\{S^{-1}\}$ 与典范映射 $A\to A\{S^{-1}\}$ 组成的二元组由
下述\emph{泛性质}刻画：设 u 是从 $A$ 到一个分离且完备的
线性拓扑化环 $B$ 的连续同态，并且 $u(S)$ 由 $B$ 中的可逆元组成，
则 $u$ 以唯一方式分解为
$A\to A\{S^{-1}\}\xrightarrow{u'}B$，其中 $u'$ 连续。事实上，
$u$ 以唯一方式分解为
$A\to S^{-1}A\xrightarrow{v'}B$；对 $B$ 的每个开理想
$\mathfrak K$，$u^{-1}(\mathfrak K)$ 包含某个
$\mathfrak J_\lambda$，所以 ${v'}^{-1}(\mathfrak K)$ 必然包含
$S^{-1}\mathfrak J_\lambda$，因而 $v'$ 连续。由于 $B$ 分离且完备，
$v'$ 以唯一方式分解为
$S^{-1}A\to A\{S^{-1}\}\xrightarrow{u'}B$，其中 $u'$ 连续；
断言得证。
\end{env}

\begin{env}[7.6.7]
\label{0.7.6.7-fr}
设 $B$ 是另一个线性拓扑化环，$T$ 是 $B$ 的一个乘法子集，
$\varphi:A\to B$ 是满足 $\varphi(S)\subset T$ 的连续同态。由上可知，
连续同态
$A\xrightarrow{\varphi}B\to B\{T^{-1}\}$ 以唯一方式分解为
$A\to A\{S^{-1}\}\xrightarrow{\varphi'}B\{T^{-1}\}$，其中
$\varphi'$ 连续。特别地，若 $B=A$ 且 $\varphi$ 为恒等同态，则当
$S\subset T$ 时，有连续同态
$\rho^{T,S}:A\{S^{-1}\}\to A\{T^{-1}\}$；它由
$S^{-1}A\to T^{-1}A$ 取分离完备化而得。若 $U$ 是 $A$ 的第三个
乘法子集，并且 $S\subset T\subset U$，则
$\rho^{U,S}=\rho^{U,T}\circ\rho^{T,S}$。
\end{env}

\begin{env}[7.6.8]
\label{0.7.6.8-fr}
设 $S_1$、$S_2$ 是 $A$ 的两个乘法子集，$S_2'$ 是 $S_2$ 在
$A\{S_1^{-1}\}$ 中的典范像；则有典范拓扑同构
$A\{(S_1S_2)^{-1}\}\xrightarrow{\sim}A\{S_1^{-1}\}\{{S_2'}^{-1}\}$。
这可从双连续的典范同构
$(S_1S_2)^{-1}A\xrightarrow{\sim}{S_2''}^{-1}(S_1^{-1}A)$
看出，其中 $S_2''$ 是 $S_2$ 在 $S_1^{-1}A$ 中的典范像。
\end{env}

\begin{env}[7.6.9]
\label{0.7.6.9-fr}
设 $\mathfrak a$ 是 $A$ 的一个\emph{开}理想；可以假设对每个
$\lambda$ 都有 $\mathfrak J_\lambda\subset\mathfrak a$，从而在环
$S^{-1}A$ 中
$S^{-1}\mathfrak J_\lambda\subset S^{-1}\mathfrak a$。换言之，
$S^{-1}\mathfrak a$ 是 $S^{-1}A$ 的一个\emph{开}理想。以
$\mathfrak a\{S^{-1}\}$ 表示它的分离完备化；它等于
$\varprojlim(S^{-1}\mathfrak a/S^{-1}\mathfrak J_\lambda)$，是
$A\{S^{-1}\}$ 的一个\emph{开}理想，并同构于
$S^{-1}\mathfrak a$ 的典范像的闭包。此外，
\emph{离散环 $A\{S^{-1}\}/\mathfrak a\{S^{-1}\}$ 典范同构于
$S^{-1}A/S^{-1}\mathfrak a=S^{-1}(A/\mathfrak a)$}。反过来，若
$\mathfrak a'$ 是 $A\{S^{-1}\}$ 的一个开理想，则
$\mathfrak a'$ 包含一个形如 $\mathfrak J_\lambda\{S^{-1}\}$ 的
理想，因而是 $S^{-1}A/S^{-1}\mathfrak J_\lambda$ 的某个理想的
逆像；该理想必然形如 $S^{-1}\mathfrak a$ (1.2.6)，其中
$\mathfrak a\supset\mathfrak J_\lambda$。由此得到
$\mathfrak a'=\mathfrak a\{S^{-1}\}$。特别地，由 (1.2.6) 有：
\end{env}

\begin{proposition}[7.6.10]
\label{0.7.6.10-fr}
映射 $\mathfrak p\mapsto\mathfrak p\{S^{-1}\}$ 是从 $A$ 的满足
$\mathfrak p\cap S=\varnothing$ 的开素理想 $\mathfrak p$ 所成集合
到 $A\{S^{-1}\}$ 的开素理想所成集合的一个保序双射；此外，

\oldpage[0\textsubscript{I}]{74}

$A\{S^{-1}\}/\mathfrak p\{S^{-1}\}$ 的分式域典范同构于
$A/\mathfrak p$ 的分式域。
\end{proposition}

\begin{proposition}[7.6.11]
\label{0.7.6.11-fr}
\begin{enumerate}
  \item[\rm(i)] 若 $A$ 是可容环，则
    $A'=A\{S^{-1}\}$ 也是可容环，并且对 $A$ 的每个定义理想
    $\mathfrak J$，$\mathfrak J'=\mathfrak J\{S^{-1}\}$ 是 $A'$
    的一个定义理想。
  \item[\rm(ii)] 设 $A$ 是进制环，$\mathfrak J$ 是 $A$ 的一个
    定义理想，并且 $\mathfrak J/\mathfrak J^2$ 是
    $A/\mathfrak J$ 上的有限型模；则 $A'$ 是
    $\mathfrak J'$-进制环，而
    $\mathfrak J'/\mathfrak J'^2$ 是
    $A'/\mathfrak J'$ 上的有限型模。若 $A$ 还是 Noether 环，则
    $A'$ 也是 Noether 环。
\end{enumerate}
\end{proposition}

\begin{enumerate}
  \item[\rm(i)] 若 $\mathfrak J$ 是 $A$ 的一个定义理想，则
    $S^{-1}\mathfrak J$ 显然是拓扑环 $S^{-1}A$ 的一个定义理想，
    因为
    $(S^{-1}\mathfrak J)^n=S^{-1}\mathfrak J^n$。设 $A''$ 是与
    $S^{-1}A$ 相伴的分离环，$\mathfrak J''$ 是
    $S^{-1}\mathfrak J$ 在 $A''$ 中的像；$S^{-1}\mathfrak J^n$
    的像为 $\mathfrak J''^n$，所以 $\mathfrak J''^n$ 在 $A''$
    中趋于 0。由于 $\mathfrak J'$ 是 $\mathfrak J''$ 在 $A'$
    中的闭包，$\mathfrak J'^n$ 包含于 $\mathfrak J''^n$ 的闭包，
    因而在 $A'$ 中趋于 0。
  \item[\rm(ii)] 令 $A_i=A/\mathfrak J^{i+1}$；当 $i\leq j$ 时，
    设 $u_{ij}$ 为典范同态
    $A/\mathfrak J^{j+1}\to A/\mathfrak J^{i+1}$。设 $S_i$ 是
    $S$ 在 $A_i$ 中的典范像，并令 $A_i'=S_i^{-1}A_i$；最后，
    设 $u_{ij}':A_j'\to A_i'$ 为由 $u_{ij}$ 典范诱导的同态。
    我们来证明逆系统 $(A_i',u_{ij}')$ 满足命题 (7.2.7) 的条件：
    各 $u_{ij}'$ 显然为满射；另一方面，$u_{ij}'$ 的核为
    $S_j^{-1}(\mathfrak J^{i+1}/\mathfrak J^{j+1})$ (1.3.2)，
    它等于 $\mathfrak J_j'^{i+1}$，其中
    $\mathfrak J_j'=S_j^{-1}(\mathfrak J/\mathfrak J^{j+1})$；
    最后，
    $\mathfrak J_1'/\mathfrak J_1'^2=S_1^{-1}(\mathfrak J/\mathfrak J^2)$。
    由于 $\mathfrak J/\mathfrak J^2$ 是
    $A/\mathfrak J^2$ 上的有限型模，
    $\mathfrak J_1'/\mathfrak J_1'^2$ 是 $A_1'$ 上的有限型模。
    最后，若 $A$ 为 Noether 环，则
    $A_0'=S_0^{-1}(A/\mathfrak J)$ 也为 Noether 环；
    % EGA1-ERR-0018: French print/source cites (7.2.8); retained literally pending canon disposition.
    这就完成了命题 (7.2.8) 的证明。
\end{enumerate}

\begin{corollary}[7.6.12]
\label{0.7.6.12-fr}
在 (7.6.11, (ii)) 的假设下，有
$(\mathfrak J\{S^{-1}\})^n=\mathfrak J^n\{S^{-1}\}$。
\end{corollary}

这确由 (7.2.7) 以及 (7.6.11) 的证明得到。

\begin{proposition}[7.6.13]
\label{0.7.6.13-fr}
设 $A$ 是一个进制 Noether 环，$S$ 是 $A$ 的一个乘法子集；则
$A\{S^{-1}\}$ 是一个平坦 $A$-模。
\end{proposition}

事实上，若 $\mathfrak J$ 是 $A$ 的一个定义理想，则
$A\{S^{-1}\}$ 是 Noether 环 $S^{-1}A$ 关于
$S^{-1}\mathfrak J$-预进制拓扑的分离完备化；因而由 (7.3.3)，
$A\{S^{-1}\}$ 是一个平坦 $S^{-1}A$-模。又因为 $S^{-1}A$ 是
平坦 $A$-模 (6.3.1)，命题由平坦性的传递性 (6.2.1) 得到。

\begin{corollary}[7.6.14]
\label{0.7.6.14-fr}
在 (7.6.13) 的假设下，设 $S'\subset S$ 是 $A$ 的另一个乘法
子集；则 $A\{S^{-1}\}$ 是一个平坦
$A\{{S'}^{-1}\}$-模。
\end{corollary}

事实上，由 (7.6.8)，$A\{S^{-1}\}$ 典范地等同于
$A\{{S'}^{-1}\}\{S_0^{-1}\}$，其中 $S_0$ 是 $S$ 在
$A\{{S'}^{-1}\}$ 中的典范像，而 $A\{{S'}^{-1}\}$ 是 Noether 环
(7.6.11)。

\begin{env}[7.6.15]
\label{0.7.6.15-fr}
对线性拓扑化环 $A$ 的每个元素 $f$，以 $A_{\{f\}}$ 表示完备
分式环 $A\{S_f^{-1}\}$，其中 $S_f$ 是各 $f^n$（$n\geq 0$）
组成的乘法集；对 $A$ 的每个开理想 $\mathfrak a$，以
$\mathfrak a_{\{f\}}$ 代替 $\mathfrak a\{S_f^{-1}\}$。若 $g$
是 $A$ 的另一个元素，则有典范连续同态
$A_{\{f\}}\to A_{\{fg\}}$ (7.6.7)。当 $f$ 遍历 $A$ 的一个
乘法子集 $S$ 时，各 $A_{\{f\}}$ 关于上述同态构成一个滤过归纳
环系；令
$A_{\{S\}}=\varinjlim_{f\in S}A_{\{f\}}$。对每个 $f\in S$，都有
同态 $A_{\{f\}}\to A\{S^{-1}\}$ (7.6.7)，并且

\oldpage[0\textsubscript{I}]{75}

这些同态构成一个归纳系统；取归纳极限，便定义了典范同态
$A_{\{S\}}\to A\{S^{-1}\}$。
\end{env}

\begin{proposition}[7.6.16]
\label{0.7.6.16-fr}
若 $A$ 是 Noether 环，则 $A\{S^{-1}\}$ 是
$A_{\{S\}}$ 上的平坦模。
\end{proposition}

事实上，由 (7.6.14)，对每个 $f\in S$，$A\{S^{-1}\}$ 在环
$A_{\{f\}}$ 上平坦；结论由 (6.2.3) 得到。

\begin{proposition}[7.6.17]
\label{0.7.6.17-fr}
设 $\mathfrak p$ 是可容环 $A$ 的一个开素理想，并令
$S=A-\mathfrak p$。则环 $A\{S^{-1}\}$ 与 $A_{\{S\}}$ 都是
局部环，典范同态 $A_{\{S\}}\to A\{S^{-1}\}$ 是局部同态，并且
$A_{\{S\}}$ 与 $A\{S^{-1}\}$ 的剩余域都典范同构于
$A/\mathfrak p$ 的分式域。
\end{proposition}

事实上，设 $\mathfrak J\subset\mathfrak p$ 是 $A$ 的一个定义
理想；有
$S^{-1}\mathfrak J\subset S^{-1}\mathfrak p=\mathfrak p
A_{\mathfrak p}$，所以 $A_{\mathfrak p}/S^{-1}\mathfrak J$ 是
局部环；由 (7.1.12)、(7.6.9) 和 (7.6.11, (i)) 可知，
$A\{S^{-1}\}$ 是局部环。令
$\mathfrak m=\varinjlim_{f\in S}\mathfrak p_{\{f\}}$，这是
$A_{\{S\}}$ 的一个理想；我们来证明 $A_{\{S\}}$ 中每个不属于
$\mathfrak m$ 的元素都可逆。事实上，这样的元素是某个
$f\in S$ 所对应的一个元素
$z\in A_{\{f\}}$ 在 $A_{\{S\}}$ 中的像，其中它不属于
$\mathfrak p_{\{f\}}$；所以它在
$A_{\{f\}}/\mathfrak J_{\{f\}}=S_f^{-1}(A/\mathfrak J)$ 中的
典范像 $z_0$ 不属于
$S_f^{-1}(\mathfrak p/\mathfrak J)$ (7.6.9)。这意味着
$z_0=\overline{x}/\overline{f}^k$，其中
$x\notin\mathfrak p$，而 $\overline{x},\overline{f}$ 是 $x,f$
模 $\mathfrak J$ 的类。由于 $x\in S$，有 $g=xf\in S$；在
$S_g^{-1}A$ 中，$x/f^k\in S_f^{-1}A$ 的典范像
$y_0=x^{k+1}/g^k$ 以 $x^{k-1}f^{2k}/g^k$ 为逆。这更蕴涵
$y_0$ 在 $S_g^{-1}A/S_g^{-1}\mathfrak J$ 中的像可逆；所以由
(7.6.9) 和 (7.1.12)，$z$ 在 $A_{\{g\}}$ 中的典范像 $y$
可逆。因此，$z$ 在 $A_{\{S\}}$ 中的像（它等于 $y$ 的像）
可逆。由此可见，$A_{\{S\}}$ 是以 $\mathfrak m$ 为极大理想的
局部环；此外，$\mathfrak p_{\{f\}}$ 在 $A\{S^{-1}\}$ 中的像
包含于该环的极大理想 $\mathfrak p\{S^{-1}\}$，因而
$\mathfrak m$ 在 $A\{S^{-1}\}$ 中的像也包含于
$\mathfrak p\{S^{-1}\}$；所以典范同态
$A_{\{S\}}\to A\{S^{-1}\}$ 是局部同态。最后，由于
$A\{S^{-1}\}/\mathfrak p\{S^{-1}\}$ 的每个元素都是某个适当的
$f\in S$ 所对应的环 $S_f^{-1}A$ 的一个元素的像，同态
$A_{\{S\}}\to A\{S^{-1}\}/\mathfrak p\{S^{-1}\}$ 为满射；取商后，
它给出剩余域的一个同构。

\begin{corollary}[7.6.18]
\label{0.7.6.18-fr}
在 (7.6.17) 的假设下，若还假设 $A$ 是进制 Noether 环，则局部环
$A\{S^{-1}\}$ 与 $A_{\{S\}}$ 都是 Noether 环，并且
$A\{S^{-1}\}$ 是一个忠实平坦 $A_{\{S\}}$-模。
\end{corollary}

我们已经知道 $A\{S^{-1}\}$ 是 Noether 环 (7.6.11, (ii))，且是
平坦 $A_{\{S\}}$-模 (7.6.16)；由于同态
$A_{\{S\}}\to A\{S^{-1}\}$ 是局部同态，可知 $A\{S^{-1}\}$ 是
忠实平坦 $A_{\{S\}}$-模 (6.6.2)，因而 $A_{\{S\}}$ 是 Noether 环
(6.5.2)。

\subsection{完备张量积}
\label{subsection:0.7.7-fr}

\begin{env}[7.7.1]
\label{0.7.7.1-fr}
设 $A$ 是一个线性拓扑化环，$M$、$N$ 是两个线性拓扑化 $A$-模。
设 $\mathfrak J$、$V$、$W$ 分别是 $A$、$M$、$N$ 中的开
$0$-邻域，它们都是 $A$-模，并且满足
$\mathfrak JM\subset V$、$\mathfrak JN\subset W$，从而
$M/V$ 与 $N/W$ 可以视为 $(A/\mathfrak J)$-模。当
$\mathfrak J$、$V$、$W$ 遍历满足上述条件的开邻域系时，显然各模
$(M/V)\otimes_{A/\mathfrak J}(N/W)$ 在逆环系

\oldpage[0\textsubscript{I}]{76}

$A/\mathfrak J$ 上构成一个逆模系；取逆极限，便得到 $A$ 的
分离完备化 $\widehat A$ 上的一个模，称为 $M$ 与 $N$ 的
\emph{完备张量积}，记为 $(M\otimes_A N)^\wedge$。注意到
$M/V$ 典范同构于 $\widehat M/\overline V$，其中 $\widehat M$ 是
$M$ 的分离完备化，而 $\overline V$ 是 $V$ 的像在 $\widehat M$
中的闭包；于是可见，完备张量积
$(M\otimes_A N)^\wedge$ 典范地等同于
$(\widehat M\otimes_{\widehat A}\widehat N)^\wedge$，后者也记为
$\widehat M\widehat\otimes_{\widehat A}\widehat N$。
\end{env}

\begin{env}[7.7.2]
\label{0.7.7.2-fr}
\begin{sloppypar}
沿用上述记号，张量积
$(M/V)\otimes_A(N/W)$ 与
$(M/V)\otimes_{A/\mathfrak J}(N/W)$ 典范地等同；因此，它们也
典范地等同于
$(M\otimes_A N)/(\Im(V\otimes_A N)+\Im(M\otimes_A W))$。
由此可知，$(M\otimes_A N)^\wedge$ 是\emph{ $A$-模
$M\otimes_A N$ 关于下述拓扑的分离完备化：各子模}
\[
  \Im(V\otimes_A N)+\Im(M\otimes_A W)
\]
\emph{组成 $0$ 的一个基本邻域系}（其中 $V$ 与 $W$ 分别遍历
$M$ 与 $N$ 的开子模）；为简便起见，我们称该拓扑为 $M$ 与 $N$
上给定拓扑的\emph{张量积}。
\end{sloppypar}
\end{env}

\begin{env}[7.7.3]
\label{0.7.7.3-fr}
设 $M'$、$N'$ 是两个线性拓扑化 $A$-模，
$u:M\to M'$、$v:N\to N'$ 是两个连续同态；显然，$u\otimes v$
关于 $M\otimes_A N$ 与 $M'\otimes_A N'$ 上各自的张量积拓扑
连续。取分离完备化，便得到连续同态
$(M\otimes_A N)^\wedge\to(M'\otimes_A N')^\wedge$，我们以
$u\widehat\otimes v$ 表示它。因此，$(M\otimes_A N)^\wedge$
在 $M$ 与 $N$ 上是线性拓扑化 $A$-模范畴中的一个\emph{双函子}。
\end{env}

\begin{env}[7.7.4]
\label{0.7.7.4-fr}
同样可以定义任意有限多个线性拓扑化 $A$-模的完备张量积；该张量积
显然具有通常的结合性与交换性。
\end{env}

\begin{env}[7.7.5]
\label{0.7.7.5-fr}
若 $B$、$C$ 是两个线性拓扑化的拓扑 $A$-代数，则
$B\otimes_A C$ 上的张量积拓扑以代数 $B\otimes_A C$ 的各
\emph{理想}
$\Im(\mathfrak K\otimes_A C)+\Im(B\otimes_A\mathfrak L)$ 为
$0$ 的基本邻域系，其中 $\mathfrak K$（相应地，$\mathfrak L$）
遍历 $B$（相应地，$C$）的开理想。因此，
$(B\otimes_A C)^\wedge$ 具有一个
\emph{$\widehat A$-拓扑代数}结构，它是各
$(A/\mathfrak J)$-代数
$(B/\mathfrak K)\otimes_{A/\mathfrak J}(C/\mathfrak L)$ 所成
逆系统的逆极限（$\mathfrak J$ 是 $A$ 的一个满足
$\mathfrak J.B\subset\mathfrak K$、
$\mathfrak J.C\subset\mathfrak L$ 的开理想；这样的理想总是存在）。
称这个代数为代数 $B$ 与 $C$ 的\emph{完备张量积}。
\end{env}

\begin{env}[7.7.6]
\label{0.7.7.6-fr}
从 $B$ 与 $C$ 到 $B\otimes_A C$ 的 $A$-代数同态
$b\mapsto b\otimes 1$、$c\mapsto 1\otimes c$，在后一个代数上
赋予张量积拓扑时连续；将它们与从 $B\otimes_A C$ 到其分离
完备化的典范同态复合，便得到\emph{典范同态}
$\rho:B\to(B\otimes_A C)^\wedge$、
$\sigma:C\to(B\otimes_A C)^\wedge$。此外，代数
$(B\otimes_A C)^\wedge$ 与同态 $\rho$、$\sigma$ 具有下述
\emph{泛性质}：对每个分离且完备的线性拓扑化 $A$-代数 $D$，
以及每一对连续 $A$-同态 $u:B\to D$、$v:C\to D$，存在唯一的
% EGA1-ERR-0019: French print/source duplicates “et un seul”; retained literally pending canon disposition.
连续 $A$-同态 $w:(B\otimes_A C)^\wedge\to D$，并且唯一地满足
$u=w\circ\rho$ 与 $v=w\circ\sigma$。事实上，已经存在唯一的
$A$-同态 $w_0:B\otimes_A C\to D$，使得
$u(b)=w_0(b\otimes 1)$ 与 $v(c)=w_0(1\otimes c)$；只须证明
$w_0$ 连续，因为这样一来，取分离完备化后它便给出连续同态 $w$。

\oldpage[0\textsubscript{I}]{77}

若 $\mathfrak M$ 是 $D$ 的一个开理想，则由假设存在开理想
$\mathfrak K\subset B$、$\mathfrak L\subset C$，使得
$u(\mathfrak K)\subset\mathfrak M$、
$v(\mathfrak L)\subset\mathfrak M$；在 $w_0$ 下，
$\Im(\mathfrak K\otimes_A C)+\Im(B\otimes_A\mathfrak L)$ 的像
仍包含于 $\mathfrak M$；断言得证。
\end{env}

\begin{proposition}[7.7.7]
\label{0.7.7.7-fr}
若 $B$ 与 $C$ 是两个预可容 $A$-代数，则
$(B\otimes_A C)^\wedge$ 可容；若 $\mathfrak K$（相应地，
$\mathfrak L$）是 $B$（相应地，$C$）的一个定义理想，则
$\mathfrak H=\Im(\mathfrak K\otimes_A C)+
\Im(B\otimes_A\mathfrak L)$ 的典范像在
$(B\otimes_A C)^\wedge$ 中的闭包是一个定义理想。
\end{proposition}

只须证明 $\mathfrak H^n$ 关于张量积拓扑趋于 $0$；这立即由包含
关系
\[
  \mathfrak H^{2n}\subset
  \Im(\mathfrak K^n\otimes_A C)+\Im(B\otimes_A\mathfrak L^n).
\]
得到。

\begin{proposition}[7.7.8]
\label{0.7.7.8-fr}
设 $A$ 是预进制环，$\mathfrak J$ 是 $A$ 的一个定义理想，
$M$ 是一个\emph{有限型} $A$-模，并赋予 $\mathfrak J$-预进制
拓扑。对每个进制 Noether 拓扑 $A$-代数 $B$，$B\otimes_A M$
等同于完备张量积 $(B\otimes_A M)^\wedge$。
\end{proposition}

若 $\mathfrak K$ 是 $B$ 的一个定义理想，则由假设存在整数 $m$，
使得 $\mathfrak J^mB\subset\mathfrak K$，从而
$\Im(B\otimes_A\mathfrak J^{nm}M)=
\Im(\mathfrak J^{nm}B\otimes_A M)\subset
\Im(\mathfrak K^nB\otimes_A M)=\mathfrak K^n(B\otimes_A M)$；
由此可知，在 $B\otimes_A M$ 上，$B$ 与 $M$ 的拓扑的张量积是
$\mathfrak K$-预进制拓扑。由于 $B\otimes_A M$ 是有限型 $B$-模，
命题立即由 (7.3.6) 得到。

\subsection{同态模上的拓扑}
\label{subsection:0.7.8-fr}

\begin{env}[7.8.1]
\label{0.7.8.1-fr}
设 $A$ 是一个 $\mathfrak J$-进制 Noether 环，$M$ 与 $N$ 是两个
有限型 $A$-模，并赋予 $\mathfrak J$-预进制拓扑；我们知道
(7.3.6)，它们分离且完备。此外，每个 $A$-同态 $M\to N$ 都自动
连续，而 $A$-模 $\operatorname{Hom}_A(M,N)$ 是有限型的。对每个
整数 $i\geq 0$，令
$A_i=A/\mathfrak J^{i+1}$、$M_i=M/\mathfrak J^{i+1}M$、
$N_i=N/\mathfrak J^{i+1}N$；当 $i\leq j$ 时，每个同态
$u_j:M_j\to N_j$ 把 $\mathfrak J^{i+1}M_j$ 映到
$\mathfrak J^{i+1}N_j$ 中，所以取商后给出同态
$u_i:M_i\to N_i$，由此定义一个典范同态
$\operatorname{Hom}_{A_j}(M_j,N_j)\to
\operatorname{Hom}_{A_i}(M_i,N_i)$。此外，各
$\operatorname{Hom}_{A_i}(M_i,N_i)$ 关于这些同态构成一个
\emph{逆系统}；由 (7.2.10) 可知，存在典范同构
\[
  \varphi:\operatorname{Hom}_A(M,N)\longrightarrow
  \varprojlim_i\operatorname{Hom}_{A_i}(M_i,N_i).
\]
此外：
\end{env}

\begin{proposition}[7.8.2]
\label{0.7.8.2-fr}
若 $M$ 与 $N$ 是 $\mathfrak J$-进制 Noether 环 $A$ 上的有限型模，
则各子模 $\operatorname{Hom}_A(M,\mathfrak J^{i+1}N)$ 组成
$\operatorname{Hom}_A(M,N)$ 关于 $\mathfrak J$-进制拓扑的一个
$0$ 的基本邻域系，并且典范同态
\[
  \varphi:\operatorname{Hom}_A(M,N)\longrightarrow
  \varprojlim_i\operatorname{Hom}_{A_i}(M_i,N_i)
\]
是拓扑同构。
\end{proposition}

事实上，可以把 $M$ 视为有限型自由 $A$-模 $L$ 的一个商，从而把
$\operatorname{Hom}_A(M,N)$ 等同于
$\operatorname{Hom}_A(L,N)$ 的一个子模。在这个等同下，
$\operatorname{Hom}_A(M,\mathfrak J^{i+1}N)$ 是
$\operatorname{Hom}_A(M,N)$ 与
$\operatorname{Hom}_A(L,\mathfrak J^{i+1}N)$ 的交；由于
$\operatorname{Hom}_A(L,N)$ 上的 $\mathfrak J$-进制拓扑在
$\operatorname{Hom}_A(M,N)$ 上诱导的拓扑就是
$\mathfrak J$-进制拓扑 (7.3.2)，只须对 $M=L=A^m$ 证明第一个
断言。

\oldpage[0\textsubscript{I}]{78}

可是在这种情形下，
$\operatorname{Hom}_A(L,N)=N^m$，
$\operatorname{Hom}_A(L,\mathfrak J^{i+1}N)=
(\mathfrak J^{i+1}N)^m=\mathfrak J^{i+1}N^m$，结论显然。为证明
第二个断言，注意到当 $j\leq i$ 时，
$\operatorname{Hom}_A(M,\mathfrak J^{i+1}N)$ 在
$\operatorname{Hom}_{A_j}(M_j,N_j)$ 中的像为零，故 $\varphi$
连续；反过来，$\operatorname{Hom}_{A_i}(M_i,N_i)$ 中 $0$ 的
逆像在 $\operatorname{Hom}_A(M,N)$ 中就是
$\operatorname{Hom}_A(M,\mathfrak J^{i+1}N)$，所以 $\varphi$
双连续。

若仅假设 $A$ 是一个 $\mathfrak J$-预进制 Noether 环，$M$ 与
$N$ 是两个关于 $\mathfrak J$-预进制拓扑分离的有限型 $A$-模，
则上述证明表明 (7.8.2) 的第一个断言仍然成立，而 $\varphi$ 是从
$\operatorname{Hom}_A(M,N)$ 到
$\varprojlim_i\operatorname{Hom}_{A_i}(M_i,N_i)$ 的一个子模上的
拓扑同构。

\begin{proposition}[7.8.3]
\label{0.7.8.3-fr}
在 (7.8.2) 的假设下，从 $M$ 到 $N$ 的单同态（相应地，满同态、
双射同态）所成集合是 $\operatorname{Hom}_A(M,N)$ 的一个开子集。
\end{proposition}

事实上，由 (7.3.5) 和 (7.1.14)，$u$ 为满射的充要条件是相应的
同态 $u_0:M/\mathfrak JM\to N/\mathfrak JN$ 为满射；所以从
$M$ 到 $N$ 的满同态所成集合，是连续映射
$\operatorname{Hom}_A(M,N)\to
\operatorname{Hom}_{A_0}(M_0,N_0)$ 下某个离散空间子集的逆像。
现在证明单同态所成集合为开集。设 $v$ 是一个这样的同态，并令
$M'=v(M)$；由 Artin--Rees 引理 (7.3.2.1)，存在整数 $k\geq 0$，
使得对每个 $m>0$ 都有
$M'\cap\mathfrak J^{m+k}N\subset\mathfrak J^mM'$。我们将证明，
对每个 $w\in\mathfrak J^{k+1}\operatorname{Hom}_A(M,N)$，
$u=v+w$ 都是单射，这将完成证明。事实上，设 $x\in M$ 且
$u(x)=0$；我们证明，对每个 $i\geq 0$，关系
$x\in\mathfrak J^iM$ 都蕴涵 $x\in\mathfrak J^{i+1}M$；于是确有
\[
  x\in\bigcap_{i\geq 0}\mathfrak J^iM=(0).
\]
事实上，此时 $w(x)\in\mathfrak J^{i+k+1}N$；另一方面，
$w(x)=-v(x)\in M'$，所以
$v(x)\in M'\cap\mathfrak J^{i+k+1}N\subset\mathfrak J^{i+1}M'$。
由于 $v$ 是从 $M$ 到 $M'$ 的同构，故
$x\in\mathfrak J^{i+1}M$。证毕。

\begin{flushright}
\emph{（待续。）}
\end{flushright}
\clearpage
\oldpage[I]{79}
\phantomsection
\label{chapter:I-fr}

\begin{center}
{\large 第一章\par}
\vspace{1.5em}
{\Large\bfseries 概形语言\par}
\vspace{1.8em}
{\bfseries 目录\par}
\end{center}

\medskip
\begin{tabular}{@{}r@{\ }l@{}}
\S~1.  & 仿射概形。\\
\S~2.  & 预概形及预概形的态射。\\
\S~3.  & 预概形的积。\\
\S~4.  & 子预概形与浸入态射。\\
\S~5.  & 约化预概形；分离条件。\\
\S~6.  & 有限性条件。\\
\S~7.  & 有理映射。\\
\S~8.  & Chevalley 概形。\\
\S~9.  & 准凝聚层补充。\\
\S~10. & 形式概形。
\end{tabular}

\medskip
第 1 至第 8 节基本上只是在发展一种语言，即此后通篇使用的语言。
不过应当指出，依照本论著的一般精神，第 7、8 节的使用将少于其余各节，而且其作用也不那么本质；
我们之所以讨论 Chevalley 概形，只是为了同 Chevalley~\cite{I-1} 和 Nagata~\cite{I-9} 的语言衔接。
第 9 节给出关于准凝聚层的一些定义和结果，其中有些已不再只是把交换代数中的已知概念译成“几何”语言，
而是具有真正的整体性质；从后续各章开始，它们在态射的整体研究中不可缺少。
最后，第 10 节引入概形概念的一种推广。第三章中，我们将以它为中介，
方便地表述并证明固有态射上同调研究的基本结果；此外还应指出，
形式概形的概念对于表达“模空间理论”中的某些事实（代数簇的分类问题）似乎不可缺少。
第 10 节的结果在第三章第 3 节以前不会使用，因此建议读者在读到那里以前略过本节。

\clearpage
% French authority: source/ega1/ega1-1-fr.tex
% Sequential translation cursor: through 1.7.5 (source lines 1--1624; printed pages 80--97; complete source file).

\setcounter{section}{0}
\section{仿射概形}
\label{section:I.1-fr}

\setcounter{subsection}{0}
\subsection{环的素谱}
\label{subsection:I.1.1-fr}

\oldpage[I]{80}
\begin{env}[1.1.1]
\label{I.1.1.1-fr}
\emph{记号}：设 $A$ 是一个（交换）环，$M$ 是一个 $A$-模。在本章及以后各章中，我们将经常使用下列记号：

\smallskip
\noindent
$\operatorname{Spec}(A)={}$ $A$ 的\emph{素理想集合}，也称为 $A$ 的\emph{素谱}；对
$x\in X=\operatorname{Spec}(A)$，常把 $x$ 写成
$\mathfrak{j}_x$ 更为方便。$\operatorname{Spec}(A)$ 为\emph{空集}的充要条件是环 $A$ 退化为 $0$。

\smallskip
\noindent
$A_x=A_{\mathfrak{j}_x}={}$\emph{分式（局部）环 $S^{-1}A$}，其中
$S=A-\mathfrak{j}_x$。

\smallskip
\noindent
$\mathfrak{m}_x=\mathfrak{j}_xA_{\mathfrak{j}_x}={}$ $A_x$ 的\emph{极大理想}。

\smallskip
\noindent
$k(x)=A_x/\mathfrak{m}_x={}$ $A_x$ 的\emph{剩余域}；它典范同构于整环
$A/\mathfrak{j}_x$ 的分式域，并与后者视为同一。

\smallskip
\noindent
$f(x)={}$ $f$ 在 $A/\mathfrak{j}_x\subset k(x)$ 中的\emph{模
$\mathfrak{j}_x$ 剩余类}，其中 $f\in A$、$x\in X$。也称
$f(x)$ 为 $f$ 在点 $x\in\operatorname{Spec}(A)$ 处的\emph{值}；
关系 $f(x)=0$ 与 $f\in\mathfrak{j}_x$ \emph{等价}。

\smallskip
\noindent
$M_x=M\otimes_AA_x={}$\emph{分母取自
$A-\mathfrak{j}_x$ 的分式模}。

\smallskip
\noindent
$\mathfrak{r}(E)={}$ $A$ 中由子集 $E$ 生成的理想的\emph{根}，其中
所指子集是 $A$ 的一个子集。

\smallskip
\noindent
$V(E)={}$\emph{满足 $E\subset\mathfrak{j}_x$ 的 $x\in X$ 所成集合}
（也就是使每个 $f\in E$ 都满足 $f(x)=0$ 的 $x\in X$ 所成集合），其中
$E\subset A$。因此
\begin{equation}
  \mathfrak{r}(E)=\bigcap_{x\in V(E)}\mathfrak{j}_x.
  \tag{1.1.1.1}\label{I.1.1.1.1-fr}
\end{equation}

\smallskip
\noindent
$V(f)=V(\{f\})$，其中 $f\in A$。

\smallskip
\noindent
$D(f)=X-V(f)={}$\emph{使 $f(x)\neq0$ 的 $x\in X$ 所成集合}。
\end{env}

\begin{proposition}[1.1.2]
\label{I.1.1.2-fr}
有下列性质：
\begin{enumerate}
  \item[\rm(i)] $V(0)=X$，$V(1)=\varnothing$。
  \item[\rm(ii)] 关系 $E\subset E'$ 蕴含 $V(E)\supset V(E')$。
  \item[\rm(iii)] 对 $A$ 的任一子集族 $(E_\lambda)$，
  \[
    V\left(\bigcup_\lambda E_\lambda\right)
      =V\left(\sum_\lambda E_\lambda\right)
      =\bigcap_\lambda V(E_\lambda).
  \]
  \item[\rm(iv)] $V(EE')=V(E)\cup V(E')$。
  \item[\rm(v)] $V(E)=V(\mathfrak{r}(E))$。
\end{enumerate}
\end{proposition}

性质 (i)、(ii)、(iii) 是平凡的，而 (v) 由 (ii) 与公式
(\hyperref[I.1.1.1.1-fr]{1.1.1.1}) 得到。显然
$V(EE')\supset V(E)\cup V(E')$；反过来，若 $x\notin V(E)$ 且
$x\notin V(E')$，则存在 $f\in E$ 与 $f'\in E'$，使得在 $k(x)$ 中
$f(x)\neq0$ 且 $f'(x)\neq0$，从而 $f(x)f'(x)\neq0$，也就是说
$x\notin V(EE')$；这就证明了 (iv)。

命题 (\hyperref[I.1.1.2-fr]{1.1.2}) 特别表明，形如 $V(E)$ 的集合
（其中 $E$ 遍历 $A$ 的所有子集）构成 $X$ 上某个拓扑的\emph{闭集}；
我们称这个拓扑为\emph{谱拓扑}\footnote{Zariski 把这个拓扑引入代数几何；
因此，它也常称为 $X$ 上的“Zariski 拓扑”。}。除非明确说明相反情形，
总假定 $X=\operatorname{Spec}(A)$ 赋有谱拓扑。

\oldpage[I]{81}
\begin{env}[1.1.3]
\label{I.1.1.3-fr}
对 $X$ 的任一子集 $Y$，以 $\mathfrak{j}(Y)$ 表示所有满足
$f(y)=0$（对每个 $y\in Y$）的 $f\in A$ 所成集合；等价地说，
$\mathfrak{j}(Y)$ 是 $y\in Y$ 时各素理想 $\mathfrak{j}_y$ 的交。
显然，关系 $Y\subset Y'$ 蕴含
$\mathfrak{j}(Y)\supset\mathfrak{j}(Y')$，并且对 $X$ 的任一子集族
$(Y_\lambda)$ 都有
\begin{equation}
  \mathfrak{j}\left(\bigcup_\lambda Y_\lambda\right)
    =\bigcap_\lambda\mathfrak{j}(Y_\lambda)
  \tag{1.1.3.1}\label{I.1.1.3.1-fr}
\end{equation}
最后，
\begin{equation}
  \mathfrak{j}(\{x\})=\mathfrak{j}_x.
  \tag{1.1.3.2}\label{I.1.1.3.2-fr}
\end{equation}
\end{env}

\begin{proposition}[1.1.4]
\label{I.1.1.4-fr}
\begin{enumerate}
  \item[\rm(i)] 对 $A$ 的任一子集 $E$，有
  $\mathfrak{j}(V(E))=\mathfrak{r}(E)$。
  \item[\rm(ii)] 对 $X$ 的任一子集 $Y$，
  $V(\mathfrak{j}(Y))=\overline{Y}$，后者是 $Y$ 在 $X$ 中的闭包。
\end{enumerate}
\end{proposition}

(i) 是定义与 (\hyperref[I.1.1.1.1-fr]{1.1.1.1}) 的直接推论；
另一方面，$V(\mathfrak{j}(Y))$ 是包含 $Y$ 的闭集。反过来，若
$Y\subset V(E)$，则对每个 $f\in E$ 与每个 $y\in Y$ 都有
$f(y)=0$，故 $E\subset\mathfrak{j}(Y)$，
$V(E)\supset V(\mathfrak{j}(Y))$；这就证明了 (ii)。

\begin{corollary}[1.1.5]
\label{I.1.1.5-fr}
$X=\operatorname{Spec}(A)$ 的闭子集与 $A$ 的根式理想（也就是说，
素理想的交）通过反序映射
$Y\mapsto\mathfrak{j}(Y)$、$\mathfrak{a}\mapsto V(\mathfrak{a})$
一一对应；两个闭子集的并 $Y_1\cup Y_2$ 对应于
$\mathfrak{j}(Y_1)\cap\mathfrak{j}(Y_2)$，而任意闭子集族
$(Y_\lambda)$ 的交对应于各 $\mathfrak{j}(Y_\lambda)$ 之和的根。
\end{corollary}

\begin{corollary}[1.1.6]
\label{I.1.1.6-fr}
若 $A$ 是 Noether 环，则 $X=\operatorname{Spec}(A)$ 是 Noether 空间。
\end{corollary}

应注意，这个推论的逆命题不成立。例如，一个非 Noether 整环可以只有一个
素理想 $\neq\{0\}$；秩为 1 的非离散赋值环就是一例。

作为素谱不是 Noether 空间的环 $A$ 的例子，可以取无限紧空间 $Y$ 上的
实值连续函数环 $\mathcal{C}(Y)$。已知作为集合，$Y$ 可与 $A$ 的极大理想集合
等同；容易看出，$X=\operatorname{Spec}(A)$ 的拓扑在 $Y$ 上诱导的拓扑正是
$Y$ 原有的拓扑。由于 $Y$ 不是 Noether 空间，$X$ 也不是 Noether 空间。

\begin{corollary}[1.1.7]
\label{I.1.1.7-fr}
对任一 $x\in X$，$\{x\}$ 的闭包是所有满足
$\mathfrak{j}_x\subset\mathfrak{j}_y$ 的 $y\in X$ 所成集合。
$\{x\}$ 为闭集的充要条件是 $\mathfrak{j}_x$ 为极大理想。
\end{corollary}

\begin{corollary}[1.1.8]
\label{I.1.1.8-fr}
空间 $X=\operatorname{Spec}(A)$ 是 Kolmogorov 空间。
\end{corollary}

事实上，若 $x$、$y$ 是 $X$ 的两个不同点，则
$\mathfrak{j}_x\not\subset\mathfrak{j}_y$ 与
$\mathfrak{j}_y\not\subset\mathfrak{j}_x$ 至少有一个成立，故
$x$、$y$ 中至少有一点不属于另一点的闭包。

\begin{env}[1.1.9]
\label{I.1.1.9-fr}
由命题 (\hyperref[I.1.1.2-fr]{1.1.2}, (iv))，对 $A$ 的两个元素
$f$、$g$，有
\begin{equation}
  D(fg)=D(f)\cap D(g).
  \tag{1.1.9.1}\label{I.1.1.9.1-fr}
\end{equation}
还应注意，按照命题 (\hyperref[I.1.1.4-fr]{1.1.4}, (i)) 与命题
(\hyperref[I.1.1.2-fr]{1.1.2}, (v))，关系 $D(f)=D(g)$ 意味着
$\mathfrak{r}(f)=\mathfrak{r}(g)$；等价地，包含 $(f)$ 与 $(g)$ 的
极小素理想分别相同。特别地，当 $f=ug$ 且 $u$ 可逆时，确实如此。
\end{env}

\oldpage[I]{82}
\begin{proposition}[1.1.10]
\label{I.1.1.10-fr}
\begin{enumerate}
  \item[\rm(i)] 当 $f$ 遍历 $A$ 时，各集合 $D(f)$ 构成 $X$ 的拓扑的一个基。
  \item[\rm(ii)] 对任一 $f\in A$，$D(f)$ 都是拟紧的。特别地，
  $X=D(1)$ 是拟紧的。
\end{enumerate}
\end{proposition}

(i) 设 $U$ 是 $X$ 中的开集。由定义，$U=X-V(E)$，其中 $E$ 是 $A$ 的一个
子集；又有 $V(E)=\bigcap_{f\in E}V(f)$，故
$U=\bigcup_{f\in E}D(f)$。

(ii) 由 (i)，只须证明：若 $(f_\lambda)_{\lambda\in L}$ 是 $A$ 中的元素族，
满足 $D(f)\subset\bigcup_{\lambda\in L}D(f_\lambda)$，则存在 $L$ 的有限子集
$J$，使得 $D(f)\subset\bigcup_{\lambda\in J}D(f_\lambda)$。设
$\mathfrak{a}$ 为 $A$ 中由各 $f_\lambda$ 生成的理想；由假设，
$V(f)\supset V(\mathfrak{a})$，从而
$\mathfrak{r}(f)\subset\mathfrak{r}(\mathfrak{a})$。由于
$f\in\mathfrak{r}(f)$，存在整数 $n\geq0$ 使得
$f^n\in\mathfrak{a}$。但此时，$f^n$ 属于某个有限子族
$(f_\lambda)_{\lambda\in J}$ 所生成的理想 $\mathfrak{b}$，并且
$V(f)=V(f^n)\supset V(\mathfrak{b})=\bigcap_{\lambda\in J}V(f_\lambda)$，
也就是说 $D(f)\subset\bigcup_{\lambda\in J}D(f_\lambda)$。

\begin{proposition}[1.1.11]
\label{I.1.1.11-fr}
对 $A$ 的任一理想 $\mathfrak{a}$，
$\operatorname{Spec}(A/\mathfrak{a})$ 典范地等同于
$\operatorname{Spec}(A)$ 的闭子空间 $V(\mathfrak{a})$。
\end{proposition}

事实上，$A/\mathfrak{a}$ 的理想（相应地，素理想）与 $A$ 中包含
$\mathfrak{a}$ 的理想（相应地，素理想）之间存在保持包含序的典范一一对应。

回忆：$A$ 的所有幂零元所成集合 $\mathfrak{N}$（也就是 $A$ 的\emph{幂零根}）
是理想；它等于 $\mathfrak{r}(0)$，即 $A$ 的所有素理想的交
(\hyperref[0.1.1.1-fr]{0, 1.1.1})。

\begin{corollary}[1.1.12]
\label{I.1.1.12-fr}
拓扑空间 $\operatorname{Spec}(A)$ 与
$\operatorname{Spec}(A/\mathfrak{N})$ 典范同胚。
\end{corollary}

\begin{proposition}[1.1.13]
\label{I.1.1.13-fr}
$X=\operatorname{Spec}(A)$ 不可约
(\hyperref[0.2.1.1-fr]{0, 2.1.1}) 的充要条件是环
$A/\mathfrak{N}$ 为整环（等价地，理想 $\mathfrak{N}$ 为素理想）。
\end{proposition}

由推论 (\hyperref[I.1.1.12-fr]{1.1.12})，只须考虑
$\mathfrak{N}=0$ 的情形。若 $X$ 可约，则存在两个不同于 $X$ 的闭子集
$Y_1$、$Y_2$，使得 $X=Y_1\cup Y_2$；从而
$\mathfrak{j}(X)=\mathfrak{j}(Y_1)\cap\mathfrak{j}(Y_2)=0$，而理想
$\mathfrak{j}(Y_1)$ 与 $\mathfrak{j}(Y_2)$ 都不同于 $(0)$
(\hyperref[I.1.1.5-fr]{1.1.5})；故 $A$ 不是整环。反过来，若 $A$ 中存在
$f\neq0$、$g\neq0$ 使得 $fg=0$，则 $V(f)\neq X$、$V(g)\neq X$
（因为 $A$ 的所有素理想的交是 $(0)$），并且
$X=V(fg)=V(f)\cup V(g)$。

\begin{corollary}[1.1.14]
\label{I.1.1.14-fr}
\begin{enumerate}
  \item[\rm(i)] 在 $X=\operatorname{Spec}(A)$ 的闭子集与 $A$ 的根式理想之间的
  一一对应中，$X$ 的不可约闭子集对应于 $A$ 的素理想。特别地，
  $X$ 的不可约分支对应于 $A$ 的极小素理想。
  \item[\rm(ii)] 映射 $x\mapsto\overline{\{x\}}$ 在 $X$ 与 $X$ 的不可约闭子集
  所成集合之间建立一一对应（也就是说，$X$ 的每个不可约闭子集恰有一个泛点）。
\end{enumerate}
\end{corollary}

(i) 立即由 (\hyperref[I.1.1.13-fr]{1.1.13}) 与
(\hyperref[I.1.1.11-fr]{1.1.11}) 得到。为证明 (ii)，由
(\hyperref[I.1.1.11-fr]{1.1.11})，只须考虑 $X$ 不可约的情形；
此时由命题 (\hyperref[I.1.1.13-fr]{1.1.13})，$A$ 中存在最小素理想
$\mathfrak{N}$，它对应于 $X$ 的一个泛点
\oldpage[I]{83}
；此外，$X$ 只能有一个泛点，因为它是 Kolmogorov 空间
（(\hyperref[I.1.1.8-fr]{1.1.8}) 与
(\hyperref[0.2.1.3-fr]{0, 2.1.3})）。

\begin{proposition}[1.1.15]
\label{I.1.1.15-fr}
若 $\mathfrak{J}$ 是 $A$ 的一个包含在 Jacobson 根
$\mathfrak{R}(A)$ 中的理想，则 $V(\mathfrak{J})$ 在
$X=\operatorname{Spec}(A)$ 中的唯一邻域是整个空间 $X$。
\end{proposition}

% EGA1-ERR-0020: French authority omits the proper-ideal restriction in the following sentence; kept source-literal.
事实上，按定义，$A$ 的每个极大理想都属于 $V(\mathfrak{J})$。由于 $A$ 的每个
理想 $\mathfrak{a}$ 都包含在一个极大理想中，有
$V(\mathfrak{a})\cap V(\mathfrak{J})\neq\varnothing$；命题得证。

\subsection{环的素谱的函子性质}
\label{subsection:I.1.2-fr}

\begin{env}[1.2.1]
\label{I.1.2.1-fr}
设 $A$、$A'$ 是两个环，
\[
  \varphi:A'\longrightarrow A
\]
是一个环同态。对任一素理想
$x=\mathfrak{j}_x\in\operatorname{Spec}(A)=X$，环
$A'/\varphi^{-1}(\mathfrak{j}_x)$ 典范同构于
$A/\mathfrak{j}_x$ 的一个子环，因而是整环；换言之，
$\varphi^{-1}(\mathfrak{j}_x)$ 是 $A'$ 的一个素理想。把它记作
${}^a\varphi(x)$；这样就定义了一个映射
\[
  {}^a\varphi:X=\operatorname{Spec}(A)\longrightarrow
  X'=\operatorname{Spec}(A')
\]
（也记作 $\operatorname{Spec}(\varphi)$），称为同态
$\varphi$ 的\emph{伴随映射}。以 $\varphi^x$ 表示由 $\varphi$
取商后诱导的从 $A'/\varphi^{-1}(\mathfrak{j}_x)$ 到
$A/\mathfrak{j}_x$ 的单同态，也表示它向域单同态的
典范延拓
% EGA1-ERR-0023: EGA II Errata orders uppercase K replaced by lowercase k twice; kept source-literal.
\[
  \varphi^x:K({}^a\varphi(x))\longrightarrow K(x)~;
\]
因此，按定义，对任一 $f'\in A'$ 有
\begin{equation}
  \varphi^x\bigl(f'({}^a\varphi(x))\bigr)=\bigl(\varphi(f')\bigr)(x)
  \qquad(x\in X).
  \tag{1.2.1.1}\label{I.1.2.1.1-fr}
\end{equation}
\end{env}

\begin{proposition}[1.2.2]
\label{I.1.2.2-fr}
\begin{enumerate}
  \item[\rm(i)] 对 $A'$ 的任一子集 $E'$，有
  \begin{equation}
    {}^a\varphi^{-1}(V(E'))=V(\varphi(E'))
    \tag{1.2.2.1}\label{I.1.2.2.1-fr}
  \end{equation}
  特别地，对任一 $f'\in A'$，
  \begin{equation}
    {}^a\varphi^{-1}(D(f'))=D(\varphi(f')).
    \tag{1.2.2.2}\label{I.1.2.2.2-fr}
  \end{equation}
  \item[\rm(ii)] 对 $A$ 的任一理想 $\mathfrak{a}$，有
  \begin{equation}
    \overline{{}^a\varphi(V(\mathfrak{a}))}
    =V(\varphi^{-1}(\mathfrak{a})).
    \tag{1.2.2.3}\label{I.1.2.2.3-fr}
  \end{equation}
\end{enumerate}
\end{proposition}

事实上，按定义，关系 ${}^a\varphi(x)\in V(E')$ 等价于
$E'\subset\varphi^{-1}(\mathfrak{j}_x)$，从而等价于
$\varphi(E')\subset\mathfrak{j}_x$，最终等价于
$x\in V(\varphi(E'))$；这就证明了 (i)。为证明 (ii)，可以假设
$\mathfrak{a}$ 等于它的根，因为
$V(\mathfrak{r}(\mathfrak{a}))=V(\mathfrak{a})$
(\hyperref[I.1.1.2-fr]{1.1.2 (v)}) 且
$\varphi^{-1}(\mathfrak{r}(\mathfrak{a}))
=\mathfrak{r}(\varphi^{-1}(\mathfrak{a}))$。令
$Y=V(\mathfrak{a})$、$\mathfrak{a}'=\mathfrak{j}({}^a\varphi(Y))$，则
$\overline{{}^a\varphi(Y)}=V(\mathfrak{a}')$
（命题 (\hyperref[I.1.1.4-fr]{1.1.4 (ii)}））。按定义，关系
$f'\in\mathfrak{a}'$ 等价于对每个 $x'\in{}^a\varphi(Y)$ 都有
$f'(x')=0$；由公式 (\hyperref[I.1.2.1.1-fr]{1.2.1.1})，这又等价于
对每个 $x\in Y$ 都有 $\varphi(f')(x)=0$，也等价于
$\varphi(f')\in\mathfrak{j}(Y)=\mathfrak{a}$，因为
$\mathfrak{a}$ 等于它的根；故 (ii) 成立。

\begin{corollary}[1.2.3]
\label{I.1.2.3-fr}
映射 ${}^a\varphi$ 连续。
\end{corollary}

应注意，若 $A''$ 是第三个环，而 $\varphi'$ 是同态
$A''\to A'$，则
% EGA1-ERR-0021: the printed composition order is ill-typed under the stated arrows; kept source-literal.
${}^a(\varphi'\circ\varphi)={}^a\varphi\circ{}^a\varphi'$。这个结果与推论
(\hyperref[I.1.2.3-fr]{1.2.3}) 表明，$\operatorname{Spec}(A)$ 关于
$A$ 是从环范畴到拓扑空间范畴的一个反变函子。

\begin{corollary}[1.2.4]
\label{I.1.2.4-fr}
\oldpage[I]{84}
设 $\varphi$ 满足：每个 $f\in A$ 都可写成
$f=h\varphi(f')$，其中 $h$ 在 $A$ 中可逆（当 $\varphi$ 为满射时尤其如此）。
则 ${}^a\varphi$ 是从 $X$ 到 ${}^a\varphi(X)$ 的同胚。
\end{corollary}

我们证明：对任一子集 $E\subset A$，存在 $A'$ 的子集 $E'$，使得
$V(E)=V(\varphi(E'))$。由公理 $(T_0)$
(\hyperref[I.1.1.8-fr]{1.1.8}) 与公式
(\hyperref[I.1.2.2.1-fr]{1.2.2.1})，这首先蕴含
${}^a\varphi$ 为单射，再由同一公式蕴含 ${}^a\varphi$ 为同胚。
事实上，对每个 $f\in E$，只须取 $f'\in A'$，使得
$h\varphi(f')=f$，其中 $h$ 在 $A$ 中可逆；由这些 $f'$ 组成的集合
$E'$ 就符合要求。

\begin{env}[1.2.5]
\label{I.1.2.5-fr}
特别地，当 $\varphi$ 是从 $A$ 到商环 $A/\mathfrak{a}$ 的典范同态时，
就重新得到 (\hyperref[I.1.1.12-fr]{1.1.12})；而
${}^a\varphi$ 是从与 $\operatorname{Spec}(A/\mathfrak{a})$ 等同的
$V(\mathfrak{a})$ 到 $X=\operatorname{Spec}(A)$ 的典范嵌入。
\end{env}

(\hyperref[I.1.2.4-fr]{1.2.4}) 的另一个特殊情形给出：

\begin{corollary}[1.2.6]
\label{I.1.2.6-fr}
若 $S$ 是 $A$ 的一个乘法子集，则谱
$\operatorname{Spec}(S^{-1}A)$ 典范地（连同其拓扑）等同于
$X=\operatorname{Spec}(A)$ 中由满足
$\mathfrak{j}_x\cap S=\varnothing$ 的 $x$ 所成的子空间。
\end{corollary}

事实上，已知 (\hyperref[0.1.2.6-fr]{0, 1.2.6})，
$S^{-1}A$ 的素理想正是满足 $\mathfrak{j}_x\cap S=\varnothing$ 的各理想
$S^{-1}\mathfrak{j}_x$，并且
$\mathfrak{j}_x=({}^S i_A)^{-1}(S^{-1}\mathfrak{j}_x)$。因此，只须把推论
(\hyperref[I.1.2.4-fr]{1.2.4}) 应用于 ${}^S i_A$。

\begin{corollary}[1.2.7]
\label{I.1.2.7-fr}
${}^a\varphi(X)$ 在 $X'$ 中稠密的充要条件是核
$\operatorname{Ker}\varphi$ 的每个元素都幂零。
\end{corollary}

事实上，把公式 (\hyperref[I.1.2.2.3-fr]{1.2.2.3}) 应用于理想
$\mathfrak{a}=(0)$，得到
$\overline{{}^a\varphi(X)}=V(\operatorname{Ker}\varphi)$；而
% EGA1-ERR-0022: the print uses X and A although Ker(varphi) is an ideal of A'; kept source-literal.
$V(\operatorname{Ker}\varphi)=X$ 的充要条件是
$\operatorname{Ker}\varphi$ 包含在 $A$ 的每个素理想中，也就是说，
包含在 $A$ 的幂零根 $\mathfrak{N}$ 中。

\subsection{模所确定的层}
\label{subsection:I.1.3-fr}
\label{I.1.3-fr}

\begin{env}[1.3.1]
\label{I.1.3.1-fr}
设 $A$ 是交换环，$M$ 是 $A$-模，$f$ 是 $A$ 的一个元素，
$S_f$ 是由各 $f^n$（其中 $n\geq0$）组成的乘法集。回忆我们记
$A_f=S_f^{-1}A$、$M_f=S_f^{-1}M$。若 $S'_f$ 是 $A$ 中由那些整除
$S_f$ 某个元素的 $g\in A$ 组成的饱和乘法子集，则已知 $A_f$ 与 $M_f$
分别典范地等同于 ${S'_f}^{-1}A$ 与 ${S'_f}^{-1}M$
（\hyperref[0.1.4.3-fr]{0, 1.4.3}）。
\end{env}

\begin{lemma}[1.3.2]
\label{I.1.3.2-fr}
下列条件等价：
\begin{enumerate}
  \item[\rm(a)] $g\in S'_f$；
  \item[\rm(b)] $S'_g\subset S'_f$；
  \item[\rm(c)] $f\in\mathfrak{r}(g)$；
  \item[\rm(d)] $\mathfrak{r}(f)\subset\mathfrak{r}(g)$；
  \item[\rm(e)] $V(g)\subset V(f)$；
  \item[\rm(f)] $D(f)\subset D(g)$。
\end{enumerate}
\end{lemma}

这由定义及（\hyperref[I.1.1.5-fr]{1.1.5}）立即得到。

\begin{env}[1.3.3]
\label{I.1.3.3-fr}
若 $D(f)=D(g)$，则引理（\hyperref[I.1.3.2-fr]{1.3.2, b)}）表明
$M_f=M_g$。更一般地，若 $D(f)\supset D(g)$，从而
$S'_f\subset S'_g$，则已知（\hyperref[0.1.4.1-fr]{0, 1.4.1}）存在一个
典范且具有函子性的同态
\[
  \rho_{g,f}:M_f\longrightarrow M_g,
\]
并且当 $D(f)\supset D(g)\supset D(h)$ 时，有
（\hyperref[0.1.4.4-fr]{0, 1.4.4}）
\begin{equation}
  \rho_{h,g}\circ\rho_{g,f}=\rho_{h,f}.
  \tag{1.3.3.1}\label{I.1.3.3.1-fr}
\end{equation}
\end{env}

\oldpage[I]{85}
当 $f$ 遍历 $A-\mathfrak{j}_x$（其中 $x$ 是
$X=\operatorname{Spec}(A)$ 中给定的一点）时，各 $S'_f$ 构成
$A-\mathfrak{j}_x$ 的子集所成的一个向上滤过族；事实上，对
$A-\mathfrak{j}_x$ 中任意两个元素 $f$、$g$，$S'_f$ 与 $S'_g$ 都包含在
$S'_{fg}$ 中。又因为所有 $f\in A-\mathfrak{j}_x$ 所对应的 $S'_f$ 的并
等于 $A-\mathfrak{j}_x$，故由（\hyperref[0.1.4.5-fr]{0, 1.4.5}）可知，
$A_x$-模 $M_x$ 典范地等同于相对于同态族 $(\rho_{g,f})$ 的
\emph{归纳极限} $\varinjlim M_f$。对
$f\in A-\mathfrak{j}_x$（等价地，$x\in D(f)$），我们用
\[
  \rho_x^f:M_f\longrightarrow M_x
\]
表示这个典范同态。

\begin{definition}[1.3.4]
\label{I.1.3.4-fr}
在由所有 $f\in A$ 所对应的 $D(f)$ 构成的 $X$ 的基
$\mathfrak{B}$ 上，把预层 $D(f)\mapsto A_f$（相应地，
$D(f)\mapsto M_f$）层化而得的环层（相应地，$\widetilde{A}$-模），
称为素谱 $X=\operatorname{Spec}(A)$ 的\emph{结构层}（相应地，
\emph{由 $A$-模 $M$ 所确定的层}），并记为 $\widetilde{A}$ 或
$\mathscr{O}_X$（相应地，$\widetilde{M}$）
（（\hyperref[I.1.1.10-fr]{1.1.10}）、
（\hyperref[0.3.2.1-fr]{0, 3.2.1} 与
\hyperref[0.3.5.6-fr]{3.5.6}））。
\end{definition}

我们已经看到（\hyperref[0.3.2.4-fr]{0, 3.2.4}），茎
$\widetilde{A}_x$（相应地，$\widetilde{M}_x$）\emph{等同于环}
$A_x$（相应地，\emph{等同于 $A_x$-模} $M_x$）。我们用
\[
  \theta_f:A_f\longrightarrow\Gamma(D(f),\widetilde{A})
\]
\[
  \text{（相应地，}\theta_f:M_f\longrightarrow
  \Gamma(D(f),\widetilde{M})\text{）}
\]
表示典范映射；于是对任意 $x\in D(f)$ 与任意 $\xi\in M_f$，有
\begin{equation}
  (\theta_f(\xi))_x=\rho_x^f(\xi).
  \tag{1.3.4.1}\label{I.1.3.4.1-fr}
\end{equation}

\begin{proposition}[1.3.5]
\label{I.1.3.5-fr}
$\widetilde{M}$ 关于 $M$ 是从 $A$-模范畴到 $\widetilde{A}$-模范畴的
正合协变函子。
\end{proposition}

事实上，设 $M$、$N$ 是两个 $A$-模，$u$ 是同态 $M\to N$。对每个
$f\in A$，$u$ 典范地对应于从 $A_f$-模 $M_f$ 到 $A_f$-模 $N_f$ 的
同态 $u_f$，并且下图（其中 $D(g)\subset D(f)$）
\[
  \xymatrix{
    M_f\ar[r]^{u_f}\ar[d]_{\rho_{g,f}} & N_f\ar[d]^{\rho_{g,f}}\\
    M_g\ar[r]_{u_g} & N_g
  }
\]
交换（\hyperref[0.1.4.1-fr]{0, 1.4.1}）；因而这些同态定义了一个
$\widetilde{A}$-模同态 $\widetilde{u}:\widetilde{M}\to\widetilde{N}$
（\hyperref[0.3.2.3-fr]{0, 3.2.3}）。此外，对每个 $x\in X$，
$\widetilde{u}_x$ 是所有满足 $x\in D(f)$（$f\in A$）的 $u_f$ 的归纳极限；
因此由（\hyperref[0.1.4.5-fr]{0, 1.4.5}），若分别把
$\widetilde{M}_x$ 与 $\widetilde{N}_x$ 典范地等同于 $M_x$ 与 $N_x$，
则 $\widetilde{u}_x$ 等同于由 $u$ 典范导出的同态 $u_x$。若 $P$ 是第三个
$A$-模，$v$ 是同态 $N\to P$，且 $w=v\circ u$，则立即有
$w_x=v_x\circ u_x$，从而 $\widetilde{w}=\widetilde{v}\circ\widetilde{u}$。
因此，我们确实定义了一个关于 $M$ 的\emph{协变函子}
$\widetilde{M}$，它从 $A$-模范畴取值于 $\widetilde{A}$-模范畴。
\emph{这个函子是正合的}，因为对每个 $x\in X$，$M_x$ 关于 $M$ 是正合函子
（\hyperref[0.1.3.2-fr]{0, 1.3.2}）；此外，由两端的定义
（\hyperref[0.1.7.1-fr]{0, 1.7.1} 与
\hyperref[0.3.1.6-fr]{3.1.6}），有
$\operatorname{Supp}(M)=\operatorname{Supp}(\widetilde{M})$。

\oldpage[I]{86}
\begin{proposition}[1.3.6]
\label{I.1.3.6-fr}
对每个 $f\in A$，开集 $D(f)\subset X$ 典范地等同于素谱
$\operatorname{Spec}(A_f)$，而由 $A_f$-模 $M_f$ 所确定的层
$\widetilde{M_f}$ 典范地等同于限制 $\widetilde{M}|D(f)$。
\end{proposition}

第一项断言是（\hyperref[I.1.2.6-fr]{1.2.6}）的特殊情形。此外，若
$g\in A$ 满足 $D(g)\subset D(f)$，则 $M_g$ 典范地等同于 $M_f$ 的分式模，
其中分母是 $g$ 在 $A_f$ 中的典范像的各次幂
（\hyperref[0.1.4.6-fr]{0, 1.4.6}）。由定义即得
$\widetilde{M_f}$ 与 $\widetilde{M}|D(f)$ 的典范等同。

\begin{theorem}[1.3.7]
\label{I.1.3.7-fr}
对任意 $A$-模 $M$ 与任意 $f\in A$，同态
\[
  \theta_f:M_f\longrightarrow\Gamma(D(f),\widetilde{M})
\]
是双射（换言之，预层 $D(f)\mapsto M_f$ 是一个\emph{层}）。特别地，
$M$ 通过 $\theta_1$ 等同于 $\Gamma(X,\widetilde{M})$。
\end{theorem}

应当指出，若 $M=A$，则 $\theta_f$ 是环同态；因而定理
（\hyperref[I.1.3.7-fr]{1.3.7}）将表明，若借助 $\theta_f$ 把环
$A_f$ 与 $\Gamma(D(f),\widetilde{A})$ 等同，则同态
$\theta_f:M_f\to\Gamma(D(f),\widetilde{M})$ 是\emph{模}同构。

先证明 $\theta_f$ 是\emph{单射}。事实上，若 $\xi\in M_f$ 满足
$\theta_f(\xi)=0$，这意味着对 $A_f$ 的每个素理想 $\mathfrak{p}$，
存在 $h\notin\mathfrak{p}$ 使得 $h\xi=0$；由于 $\xi$ 的零化子不包含在
$A_f$ 的任何素理想中，它只能是整个 $A_f$，故 $\xi=0$。

还须证明 $\theta_f$ 是\emph{满射}。可归约到 $f=1$ 的情形；一般情形由
（\hyperref[I.1.3.6-fr]{1.3.6}）作“局部化”即得。设 $s$ 是
$\widetilde{M}$ 在 $X$ 上的一个截面。由
（\hyperref[I.1.3.4-fr]{1.3.4}）及
（\hyperref[I.1.1.10-fr]{1.1.10, (ii)}），存在 $X$ 的一个\emph{有限}
覆盖 $(D(f_i))_{i\in I}$（$f_i\in A$），使得对每个 $i\in I$，限制
$s_i=s|D(f_i)$ 都形如 $\theta_{f_i}(\xi_i)$，其中
$\xi_i\in M_{f_i}$。若 $i$、$j$ 是 $I$ 的两个指标，把 $s_i$ 与 $s_j$
在 $D(f_i)\cap D(f_j)=D(f_i f_j)$ 上的限制相等这一事实写出，按照
$\widetilde{M}$ 的定义便得到
\begin{equation}
  \rho_{f_i f_j,f_i}(\xi_i)=\rho_{f_i f_j,f_j}(\xi_j).
  \tag{1.3.7.1}\label{I.1.3.7.1-fr}
\end{equation}
按定义，对每个 $i\in I$ 可写
$\xi_i=z_i/f_i^{n_i}$，其中 $z_i\in M$；由于 $I$ 有限，给每个 $z_i$
乘以 $f_i$ 的一个幂以后，可设所有 $n_i$ 都等于同一个 $n$。此时按定义，
（\hyperref[I.1.3.7.1-fr]{1.3.7.1}）意味着存在整数
$m_{ij}\geq0$，使得
$(f_i f_j)^{m_{ij}}(f_j^n z_i-f_i^n z_j)=0$；还可以假设所有
$m_{ij}$ 都等于同一个整数 $m$。于是把 $z_i$ 替换为 $f_i^m z_i$，便归约到
$m=0$ 的情形，也就是
\begin{equation}
  f_j^n z_i=f_i^n z_j
  \tag{1.3.7.2}\label{I.1.3.7.2-fr}
\end{equation}
对任意 $i$、$j$ 都成立的情形。现在 $D(f_i^n)=D(f_i)$，并且因为各
$D(f_i)$ 覆盖 $X$，由各 $f_i^n$ 生成的理想就是 $A$；换言之，存在
$g_i\in A$ 使得 $\sum_i g_i f_i^n=1$。考虑 $M$ 的元素
$z=\sum_i g_i z_i$；由（\hyperref[I.1.3.7.2-fr]{1.3.7.2}），有
$f_i^n z=\sum_j g_j f_i^n z_j=(\sum_j g_j f_j^n)z_i=z_i$，故按定义，
$M_{f_i}$ 中 $\xi_i=z/1$。因此
\oldpage[I]{87}
$s_i$ 是 $\theta_1(z)$ 在 $D(f_i)$ 上的限制，这就证明了
$s=\theta_1(z)$，并完成证明。

\begin{corollary}[1.3.8]
\label{I.1.3.8-fr}
设 $M$、$N$ 是两个 $A$-模；从 $\operatorname{Hom}_A(M,N)$ 到
$\operatorname{Hom}_{\widetilde{A}}(\widetilde{M},\widetilde{N})$ 的典范同态
$u\mapsto\widetilde{u}$ 是双射。特别地，关系 $M=0$ 与
$\widetilde{M}=0$ 等价。
\end{corollary}

事实上，考虑从
$\operatorname{Hom}_{\widetilde{A}}(\widetilde{M},\widetilde{N})$ 到
$\operatorname{Hom}_{\Gamma(\widetilde{A})}
(\Gamma(\widetilde{M}),\Gamma(\widetilde{N}))$ 的典范同态
$v\mapsto\Gamma(v)$；由定理（\hyperref[I.1.3.7-fr]{1.3.7}），后一个模
典范地等同于 $\operatorname{Hom}_A(M,N)$。只须验证
$u\mapsto\widetilde{u}$ 与 $v\mapsto\Gamma(v)$ 互为逆映射。由
$\widetilde{u}$ 的定义，显然 $\Gamma(\widetilde{u})=u$。另一方面，若对
$v\in\operatorname{Hom}_{\widetilde{A}}(\widetilde{M},\widetilde{N})$ 记
$u=\Gamma(v)$，则由 $v$ 典范导出的映射
$w:\Gamma(D(f),\widetilde{M})\to\Gamma(D(f),\widetilde{N})$ 使下图交换：
\[
  \xymatrix{
    M\ar[r]^u\ar[d]_{\rho_{f,1}} & N\ar[d]^{\rho_{f,1}}\\
    M_f\ar[r]_w & N_f
  }
\]
因此，对每个 $f\in A$，必有 $w=u_f$
（\hyperref[0.1.2.4-fr]{0, 1.2.4}），这就表明
$\widetilde{\Gamma(v)}=v$。

\begin{corollary}[1.3.9]
\label{I.1.3.9-fr}
\begin{enumerate}
  \item[\rm(i)] 设 $u$ 是从 $A$-模 $M$ 到 $A$-模 $N$ 的一个同态；则由
  $\operatorname{Ker}u$、$\operatorname{Im}u$、$\operatorname{Coker}u$
  所确定的层分别为 $\operatorname{Ker}\widetilde{u}$、
  $\operatorname{Im}\widetilde{u}$、$\operatorname{Coker}\widetilde{u}$。
  特别地，$\widetilde{u}$ 为单射（相应地，满射、双射）的充要条件是
  $u$ 具有同一性质。
  \item[\rm(ii)] 若 $M$ 是一族 $A$-模 $(M_\lambda)$ 的归纳极限
  （相应地，直和），则在典范同构的意义下，$\widetilde{M}$ 是
  $(\widetilde{M_\lambda})$ 的归纳极限（相应地，直和）。
\end{enumerate}
\end{corollary}

(i) 只须把 $\widetilde{M}$ 关于 $M$ 是正合函子这一事实
（\hyperref[I.1.3.5-fr]{1.3.5}）应用于两条 $A$-模正合列
\[
  0\to\operatorname{Ker}u\to M\to\operatorname{Im}u\to0
\]
\[
  0\to\operatorname{Im}u\to N\to\operatorname{Coker}u\to0
\]
第二项断言于是由定理（\hyperref[I.1.3.7-fr]{1.3.7}）得到。

(ii) 设 $(M_\lambda,g_{\mu\lambda})$ 是归纳极限为 $M$ 的 $A$-模归纳系，
$g_\lambda$ 是典范同态 $M_\lambda\to M$。由于当
$\lambda\leq\mu\leq\nu$ 时有
$\widetilde{g_{\nu\mu}}\circ\widetilde{g_{\mu\lambda}}
=\widetilde{g_{\nu\lambda}}$ 及
$\widetilde{g_\lambda}=\widetilde{g_\mu}\circ
\widetilde{g_{\mu\lambda}}$，
$(\widetilde{M_\lambda},\widetilde{g_{\mu\lambda}})$ 是 $X$ 上的层归纳系。
若以 $h_\lambda$ 表示典范同态
$\widetilde{M_\lambda}\to\varinjlim\widetilde{M_\lambda}$，则存在唯一同态
$v:\varinjlim\widetilde{M_\lambda}\to\widetilde{M}$，使得
$v\circ h_\lambda=\widetilde{g_\lambda}$。为证明 $v$ 是双射，只须验证对每个
$x\in X$，$v_x$ 是从 $(\varinjlim\widetilde{M_\lambda})_x$ 到
$\widetilde{M}_x$ 的双射；但 $\widetilde{M}_x=M_x$，并且
\[
  (\varinjlim\widetilde{M_\lambda})_x
  =\varinjlim(\widetilde{M_\lambda})_x
  =\varinjlim(M_\lambda)_x=M_x
  \quad(\hyperref[0.1.3.3-fr]{0, 1.3.3}).
\]
另一方面，由定义可知，$(\widetilde{g_\lambda})_x$ 与 $(h_\lambda)_x$ 都等于
从 $(M_\lambda)_x$ 到 $M_x$ 的典范映射；由于
$(\widetilde{g_\lambda})_x=v_x\circ(h_\lambda)_x$，$v_x$ 是恒等映射。
\oldpage[I]{88}
最后，若 $M$ 是两个 $A$-模 $N$、$P$ 的直和，则显然
$\widetilde{M}=\widetilde{N}\oplus\widetilde{P}$；由于任意直和都是有限直和的
归纳极限，(ii) 的断言得证。

应当指出，（\hyperref[I.1.3.8-fr]{1.3.8}）证明了：与由 $A$-模所确定的层
同构的各层构成一个\emph{阿贝尔范畴}（T, I, 1.4）。

还应指出，由（\hyperref[I.1.3.9-fr]{1.3.9}）可知，若 $M$ 是
\emph{有限型} $A$-模，也就是说存在满同态 $A^n\to M$，则存在满同态
$\widetilde{A^n}\to\widetilde{M}$；换言之，$\widetilde{A}$-模
$\widetilde{M}$ \emph{由 $X$ 上有限多个截面生成}
（\hyperref[0.5.1.1-fr]{0, 5.1.1}），反之亦然。

\begin{env}[1.3.10]
\label{I.1.3.10-fr}
若 $N$ 是 $A$-模 $M$ 的一个子模，则典范单射 $j:N\to M$ 由
（\hyperref[I.1.3.9-fr]{1.3.9}）给出一个单同态
$\widetilde{N}\to\widetilde{M}$，从而可把 $\widetilde{N}$ 典范地等同于
$\widetilde{M}$ 的一个\emph{$\widetilde{A}$-子模}；以后总假定已经作了这一
等同。若 $N$ 与 $P$ 是 $M$ 的两个子模，则
\begin{equation}
  \widetilde{(N+P)}=\widetilde{N}+\widetilde{P}
  \tag{1.3.10.1}\label{I.1.3.10.1-fr}
\end{equation}
\begin{equation}
  \widetilde{(N\cap P)}=\widetilde{N}\cap\widetilde{P}
  \tag{1.3.10.2}\label{I.1.3.10.2-fr}
\end{equation}
这是因为 $N+P$ 与 $N\cap P$ 分别是典范同态 $N\oplus P\to M$ 的像，
以及典范同态 $M\to(M/N)\oplus(M/P)$ 的核；只须应用
（\hyperref[I.1.3.9-fr]{1.3.9}）。

由（\hyperref[I.1.3.10.1-fr]{1.3.10.1}）与
（\hyperref[I.1.3.10.2-fr]{1.3.10.2}）可知，若
$\widetilde{N}=\widetilde{P}$，则 $N=P$。
\end{env}

\begin{corollary}[1.3.11]
\label{I.1.3.11-fr}
在与由 $A$-模所确定的层同构的各层所成的范畴中，函子 $\Gamma$ 是正合的。
\end{corollary}

事实上，设
$\widetilde{M}\xrightarrow{\widetilde{u}}\widetilde{N}
\xrightarrow{\widetilde{v}}\widetilde{P}$ 是一条正合列，它对应于两个
$A$-模同态 $u:M\to N$、$v:N\to P$。若
$Q=\operatorname{Im}u$ 且 $R=\operatorname{Ker}v$，则
$\widetilde{Q}=\operatorname{Im}\widetilde{u}
=\operatorname{Ker}\widetilde{v}=\widetilde{R}$
（推论（\hyperref[I.1.3.9-fr]{1.3.9}）），故 $Q=R$。

\begin{corollary}[1.3.12]
\label{I.1.3.12-fr}
设 $M$、$N$ 是两个 $A$-模。
\begin{enumerate}
  \item[\rm(i)] 由 $M\otimes_A N$ 所确定的层典范地等同于
  $\widetilde{M}\otimes_{\widetilde{A}}\widetilde{N}$。
  \item[\rm(ii)] 若 $M$ 还有一个有限表示，则由
  $\operatorname{Hom}_A(M,N)$ 所确定的层典范地等同于
  $\mathscr{H}\!om_{\widetilde{A}}(\widetilde{M},\widetilde{N})$。
\end{enumerate}
\end{corollary}

(i) 层 $\mathscr{F}=\widetilde{M}\otimes_{\widetilde{A}}\widetilde{N}$
由下列预层层化而得：
\[
  U\mapsto\mathscr{F}(U)=\Gamma(U,\widetilde{M})
  \otimes_{\Gamma(U,\widetilde{A})}\Gamma(U,\widetilde{N})
\]
这里 $U$ 遍历由所有 $f\in A$ 所对应的 $D(f)$ 构成的 $X$ 的基
（\hyperref[I.1.1.10-fr]{1.1.10, (i)}）。由
（\hyperref[I.1.3.7-fr]{1.3.7}）与
（\hyperref[I.1.3.6-fr]{1.3.6}），$\mathscr{F}(D(f))$ 典范地等同于
$M_f\otimes_{A_f}N_f$。另一方面，已知 $A_f$-模
$M_f\otimes_{A_f}N_f$ 典范同构于 $(M\otimes_A N)_f$
（\hyperref[0.1.3.4-fr]{0, 1.3.4}），而后者又典范同构于
$\Gamma(D(f),\widetilde{(M\otimes_A N)})$
（\hyperref[I.1.3.7-fr]{1.3.7} 与
\hyperref[I.1.3.6-fr]{1.3.6}）。此外，立即可以验证，这样得到的典范同构
\[
  \mathscr{F}(D(f))\simeq\Gamma(D(f),\widetilde{(M\otimes_A N)})
\]
\oldpage[I]{89}
满足与限制映射的相容条件（\hyperref[0.1.4.2-fr]{0, 1.4.2}），因而定义了
一个典范且具有函子性的同构
\[
  \widetilde{M}\otimes_{\widetilde{A}}\widetilde{N}
  \simeq\widetilde{(M\otimes_A N)}
\]

\begingroup
\setlength{\emergencystretch}{4em}
(ii) 层
$\mathscr{G}=\mathscr{H}\!om_{\widetilde{A}}(\widetilde{M},\widetilde{N})$
由下列预层层化而得：
\[
  U\longmapsto\mathscr{G}(U)
  =\operatorname{Hom}_{\widetilde{A}|U}
  (\widetilde{M}|U,\widetilde{N}|U)
\]
这里 $U$ 遍历由各 $D(f)$ 构成的 $X$ 的基。由
（\hyperref[I.1.3.6-fr]{1.3.6} 与
\hyperref[I.1.3.8-fr]{1.3.8}），$\mathscr{G}(D(f))$ 典范地等同于
$\operatorname{Hom}_{A_f}(M_f,N_f)$；由于对 $M$ 的假设，后者又典范地
等同于 $(\operatorname{Hom}_A(M,N))_f$
（\hyperref[0.1.3.5-fr]{0, 1.3.5}）。最后，
$(\operatorname{Hom}_A(M,N))_f$ 典范地等同于
$\Gamma(D(f),\widetilde{(\operatorname{Hom}_A(M,N))})$
（\hyperref[I.1.3.6-fr]{1.3.6} 与
\hyperref[I.1.3.7-fr]{1.3.7}）；这样得到的典范同构
$\mathscr{G}(D(f))\simeq
\Gamma(D(f),\widetilde{(\operatorname{Hom}_A(M,N))})$ 与限制映射相容
（\hyperref[0.1.4.2-fr]{0, 1.4.2}），故它们定义了典范同构
\[
  \mathscr{H}\!om_{\widetilde{A}}(\widetilde{M},\widetilde{N})
  \simeq\widetilde{(\operatorname{Hom}_A(M,N))}.
\]
\endgroup

\begin{env}[1.3.13]
\label{I.1.3.13-fr}
现在设 $B$ 是一个（交换）$A$-代数。也可以把这理解为：$B$ 是一个
$A$-模，并且给定了一个元素 $e\in B$ 与一个 $A$-同态
$\varphi:B\otimes_A B\to B$，使得：$1^\circ$ 下列图交换：
\[
  \xymatrix{
    B\otimes_A B\otimes_A B\ar[r]^{\varphi\otimes 1}
      \ar[d]_{1\otimes\varphi} &
    B\otimes_A B\ar[d]^\varphi & &
    B\otimes_A B\ar[rr]^\sigma\ar[rd]_\varphi & &
    B\otimes_A B\ar[dl]^\varphi\\
    B\otimes_A B\ar[r]_\varphi &
    B & & & B
  }
\]
（$\sigma$ 是典范对称映射）；$2^\circ$
$\varphi(e\otimes x)=\varphi(x\otimes e)=x$。由
（\hyperref[I.1.3.12-fr]{1.3.12}），$\widetilde{A}$-模同态
$\widetilde{\varphi}:\widetilde{B}\otimes_{\widetilde{A}}\widetilde{B}
\to\widetilde{B}$ 满足类似的条件，因而在 $\widetilde{B}$ 上定义了
$\widetilde{A}$-代数结构。同样，给定一个 $B$-模 $N$，等价于给定一个
$A$-模 $N$ 与一个 $A$-同态 $\psi:B\otimes_A N\to N$，使得下图交换：
\[
  \xymatrix{
    B\otimes_A B\otimes_A N\ar[r]^{\varphi\otimes 1}
      \ar[d]_{1\otimes\psi} &
    B\otimes_A N\ar[d]^\psi\\
    B\otimes_A N\ar[r]_\psi & N
  }
\]
并且 $\psi(e\otimes n)=n$。同态
$\widetilde{\psi}:\widetilde{B}\otimes_{\widetilde{A}}\widetilde{N}
\to\widetilde{N}$ 满足类似条件，故在 $\widetilde{N}$ 上定义了
$\widetilde{B}$-模结构。

同样可见，若 $u:B\to B'$（相应地，$v:N\to N'$）是 $A$-代数同态
（相应地，$B$-模同态），则 $\widetilde{u}$（相应地，
$\widetilde{v}$）是 $\widetilde{A}$-代数同态（相应地，
$\widetilde{B}$-模同态）；$\operatorname{Ker}\widetilde{u}$ 是
$\widetilde{B}$-理想（相应地，$\operatorname{Ker}\widetilde{v}$、
$\operatorname{Coker}\widetilde{v}$ 与 $\operatorname{Im}\widetilde{v}$ 是
$\widetilde{B}$-模）。若 $N$ 是 $B$-模，则 $\widetilde{N}$ 为有限型
$\widetilde{B}$-模的充要条件是 $N$ 为有限型 $B$-模
（\hyperref[0.5.2.3-fr]{0, 5.2.3}）。
\oldpage[I]{90}
若 $M$、$N$ 是两个 $B$-模，则 $\widetilde{B}$-模
$\widetilde{M}\otimes_{\widetilde{B}}\widetilde{N}$ 典范地等同于
$\widetilde{(M\otimes_B N)}$；同样，当 $M$ 有有限表示时，
$\mathscr{H}\!om_{\widetilde{B}}(\widetilde{M},\widetilde{N})$ 典范地等同于
$\widetilde{(\operatorname{Hom}_B(M,N))}$；证明与
（\hyperref[I.1.3.12-fr]{1.3.12}）相同。

若 $\mathfrak{J}$ 是 $B$ 的一个理想，$N$ 是 $B$-模，则
$\widetilde{(\mathfrak{J}N)}=\widetilde{\mathfrak{J}}\cdot\widetilde{N}$。

最后，若 $B$ 是由各 $A$-子模 $B_n$（$n\in\mathbb{Z}$）分次的
$A$-代数，则 $\widetilde{A}$-代数 $\widetilde{B}$ 是各
$\widetilde{A}$-模 $\widetilde{B_n}$ 的直和
（\hyperref[I.1.3.9-fr]{1.3.9}），并由这些 $\widetilde{A}$-子模分次；
分次公理表示同态 $B_m\otimes_A B_n\to B$ 的像包含在 $B_{m+n}$ 中。
同样，若 $M$ 是由各子模 $M_n$ 分次的 $B$-模，则 $\widetilde{M}$ 是由各
$\widetilde{M_n}$ 分次的 $\widetilde{B}$-模。
\end{env}

\begin{env}[1.3.14]
\label{I.1.3.14-fr}
若 $B$ 是 $A$-代数，$M$ 是 $B$ 的一个子模，则由 $\widetilde{M}$ 生成的
$\widetilde{B}$ 的 $\widetilde{A}$-子代数
（\hyperref[0.4.1.3-fr]{0, 4.1.3}）就是 $\widetilde{A}$-子代数
$\widetilde{C}$；这里 $C$ 表示由
$M$ 生成的 $B$ 的子代数。事实上，$C$ 是 $B$ 中各子模的和，而这些子模是
同态 $\bigotimes_A^n M\to B$（$n\geq0$）的像；只须应用
（\hyperref[I.1.3.9-fr]{1.3.9}）与
（\hyperref[I.1.3.12-fr]{1.3.12}）。
\end{env}

\subsection{素谱上的拟凝聚层}
\label{subsection:I.1.4-fr}

\begin{theorem}[1.4.1]
\label{I.1.4.1-fr}
设 $X$ 是环 $A$ 的素谱，$V$ 是 $X$ 的一个拟紧开子集，
$\mathscr{F}$ 是一个 $(\mathscr{O}_X|V)$-模。下列四个条件等价：
\begin{enumerate}
  \item[\textit{a})] 存在一个 $A$-模 $M$，使得 $\mathscr{F}$ 同构于
  $\widetilde{M}|V$。
  \item[\textit{b})] 存在 $V$ 的一个有限开覆盖 $(V_i)$，其中每一项
  都是包含在 $V$ 中的形如 $D(f_i)$（$f_i\in A$）的集合，并且对每个
  $i$，$\mathscr{F}|V_i$ 都同构于形如 $\widetilde{M_i}$ 的层，这里 $M_i$
  是一个 $A_{f_i}$-模。
  \item[\textit{c})] 层 $\mathscr{F}$ 是拟凝聚的
  （\hyperref[0.5.1.3-fr]{0, 5.1.3}）。
  \item[\textit{d})] 下列两个性质成立：
  \begin{enumerate}
    \item[\textit{d}\,1)] 对每个满足 $D(f)\subset V$ 的 $f\in A$，以及每个
    截面 $s\in\Gamma(D(f),\mathscr{F})$，存在整数 $n\geq0$，使得
    $f^n s$ 可延拓为 $\mathscr{F}$ 在 $V$ 上的一个截面。
    \item[\textit{d}\,2)] 对每个满足 $D(f)\subset V$ 的 $f\in A$，以及每个
    截面 $t\in\Gamma(V,\mathscr{F})$，若 $t$ 在 $D(f)$ 上的限制为 $0$，
    则存在整数 $n\geq0$，使得 $f^n t=0$。
  \end{enumerate}
\end{enumerate}
\end{theorem}

（在条件 \textit{d}\,1) 与 \textit{d}\,2) 的写法中，我们依
\hyperref[I.1.3.7-fr]{1.3.7} 暗中把 $A$ 与 $\Gamma(\widetilde{A})$
等同起来。）

\textit{a}) 推出 \textit{b})，这是
（\hyperref[I.1.3.6-fr]{1.3.6}）以及各 $D(f_i)$ 构成 $X$ 的拓扑的一组基
（\hyperref[I.1.1.10-fr]{1.1.10}）的直接推论。由于每个 $A$-模都同构于
某个同态 $A^{(I)}\to A^{(J)}$ 的余核，
（\hyperref[I.1.3.9-fr]{1.3.9}）表明，由任一 $A$-模所确定的层都是
拟凝聚的；故 \textit{b}) 推出 \textit{c})。反过来，若 $\mathscr{F}$
是拟凝聚的，则每个 $x\in V$ 都有一个形如 $D(f)\subset V$ 的邻域，使得
$\mathscr{F}|D(f)$ 同构于某个同态
$\widetilde{A_f}^{(I)}\to\widetilde{A_f}^{(J)}$ 的余核，因而同构于由模
% EGA1-ERR-0024: printed/source tilde retained; probable local type error.
$\widetilde{N}$ 所确定的层；这个模是相应同态
$A_f^{(I)}\to A_f^{(J)}$ 的余核
（\hyperref[I.1.3.8-fr]{1.3.8} 与
\hyperref[I.1.3.9-fr]{1.3.9}）。由于 $V$ 是拟紧的，显然
\textit{c}) 推出 \textit{b})。
\oldpage[I]{91}
为了证明 \textit{b}) 推出 \textit{d}\,1) 与 \textit{d}\,2)，先设
$V=D(g)$，其中 $g\in A$，并设 $\mathscr{F}$ 同构于由某个 $A_g$-模
$N$ 所确定的层 $\widetilde{N}$。把 $X$ 换成 $V$，把 $A$ 换成 $A_g$
（\hyperref[I.1.3.6-fr]{1.3.6}），即可归约到 $g=1$ 的情形。此时
$\Gamma(D(f),\widetilde{N})$ 与 $N_f$ 典范地等同
（\hyperref[I.1.3.6-fr]{1.3.6} 与
\hyperref[I.1.3.7-fr]{1.3.7}），故截面
$s\in\Gamma(D(f),\widetilde{N})$ 可等同于形如 $z/f^n$ 的元素，其中
$z\in N$；截面 $f^n s$ 等同于 $N_f$ 中的元素 $z/1$，从而是
$\widetilde{N}$ 在 $X$ 上与元素 $z\in N$ 等同的截面在 $D(f)$ 上的限制；
这就证明了此情形下的 \textit{d}\,1)。同样，
$t\in\Gamma(X,\widetilde{N})$ 等同于某个元素 $z'\in N$，$t$ 在
$D(f)$ 上的限制等同于 $z'$ 在 $N_f$ 中的像 $z'/1$；说这个像为零，
就是存在 $n\geq0$，使得 $N$ 中 $f^n z'=0$，等价地，$f^n t=0$。

为了完成 \textit{b}) 推出 \textit{d}\,1) 与 \textit{d}\,2) 的证明，
只须建立下列引理：

\begin{lemma}[1.4.1.1]
\label{I.1.4.1.1-fr}
设 $V$ 是有限多个形如 $D(g_i)$ 的集合之并，并且每个层
$\mathscr{F}|D(g_i)$ 以及
$\mathscr{F}|(D(g_i)\cap D(g_j))=\mathscr{F}|D(g_i g_j)$ 都满足
\textit{d}\,1) 与 \textit{d}\,2)。则 $\mathscr{F}$ 具有下列两个性质：
\begin{enumerate}
  \item[\textit{d}'\,1)] 对每个 $f\in A$ 及每个截面
  $s\in\Gamma(D(f)\cap V,\mathscr{F})$，存在整数 $n\geq0$，使得
  $f^n s$ 可延拓为 $\mathscr{F}$ 在 $V$ 上的一个截面。
  \item[\textit{d}'\,2)] 对每个 $f\in A$ 及每个截面
  $t\in\Gamma(V,\mathscr{F})$，若 $t$ 在 $D(f)\cap V$ 上的限制为 $0$，
  则存在整数 $n\geq0$，使得 $f^n t=0$。
\end{enumerate}
\end{lemma}

先证明 \textit{d}'\,2)。由于 $D(f)\cap D(g_i)=D(fg_i)$，对每个 $i$
都存在整数 $n_i$，使得 $(fg_i)^{n_i}t$ 在 $D(g_i)$ 上的限制为零。
由于 $g_i$ 在 $A_{g_i}$ 中的像可逆，$f^{n_i}t$ 在 $D(g_i)$ 上的限制
也为零；取 $n$ 为各 $n_i$ 中最大的一个，即得 \textit{d}'\,2)。

为证明 \textit{d}'\,1)，把 \textit{d}\,1) 应用于层
$\mathscr{F}|D(g_i)$：存在整数 $n_i\geq0$ 以及 $\mathscr{F}$ 在
$D(g_i)$ 上的一个截面 $s_i'$，它延拓了 $(fg_i)^{n_i}s$ 在
$D(fg_i)$ 上的限制。由于 $g_i$ 在 $A_{g_i}$ 中的像可逆，存在
$\mathscr{F}$ 在 $D(g_i)$ 上的截面 $s_i$，使得
$s_i'=g_i^{n_i}s_i$，并且 $s_i$ 延拓了 $f^{n_i}s$ 在 $D(fg_i)$ 上的
限制；还可以假设所有 $n_i$ 都等于同一个整数 $n$。按照构造，
$s_i-s_j$ 在
$D(f)\cap D(g_i)\cap D(g_j)=D(fg_i g_j)$ 上的限制为零。把
\textit{d}\,2) 应用于层 $\mathscr{F}|D(g_i g_j)$，可得整数
$m_{ij}\geq0$，使得 $(fg_i g_j)^{m_{ij}}(s_i-s_j)$ 在
$D(g_i g_j)$ 上的限制为零。由于 $g_i g_j$ 在 $A_{g_i g_j}$ 中的像
可逆，$f^{m_{ij}}(s_i-s_j)$ 在 $D(g_i g_j)$ 上的限制为零。于是可以
假设所有 $m_{ij}$ 都等于同一个整数 $m$，从而存在一个截面
$s'\in\Gamma(V,\mathscr{F})$，它延拓了各 $f^m s_i$；因此这个截面
延拓了 $f^{n+m}s$，故 \textit{d}'\,1) 成立。

还须证明 \textit{d}\,1) 与 \textit{d}\,2) 推出 \textit{a})。先说明
\textit{d}\,1) 与 \textit{d}\,2) 蕴含：对每个满足
$D(g)\subset V$ 的 $g\in A$，层 $\mathscr{F}|D(g)$ 也满足这两个条件。
对 \textit{d}\,1) 而言这是显然的。另一方面，若
$t\in\Gamma(D(g),\mathscr{F})$ 在 $D(f)\subset D(g)$ 上的限制为零，
则依 \textit{d}\,1)，存在整数 $m\geq0$，使得 $g^m t$ 可
\oldpage[I]{92}
延拓为 $\mathscr{F}$ 在 $V$ 上的一个截面 $s$。应用
\textit{d}\,2)，可知存在整数 $n\geq0$，使得 $f^n g^m t=0$；又因
$g$ 在 $A_g$ 中的像可逆，故 $f^n t=0$。

现在，由于 $V$ 拟紧，引理（\hyperref[I.1.4.1.1-fr]{1.4.1.1}）表明
条件 \textit{d}'\,1) 与 \textit{d}'\,2) 成立。考虑 $A$-模
$M=\Gamma(V,\mathscr{F})$，并定义一个 $\widetilde{A}$-模同态
$u:\widetilde{M}\to j_*(\mathscr{F})$，其中 $j$ 是典范嵌入
$V\to X$。由于各 $D(f)$ 构成 $X$ 的拓扑的一组基，只须对每个
$f\in A$ 定义同态
$u_f:M_f\to\Gamma(D(f),j_*(\mathscr{F}))
=\Gamma(D(f)\cap V,\mathscr{F})$，并使之满足通常的相容条件
（\hyperref[0.3.2.5-fr]{0, 3.2.5}）。由于 $f$ 在 $A_f$ 中的典范像可逆，
限制同态
$M=\Gamma(V,\mathscr{F})\to\Gamma(D(f)\cap V,\mathscr{F})$ 分解为
$M\longrightarrow M_f\xrightarrow{u_f}
\Gamma(D(f)\cap V,\mathscr{F})$
（\hyperref[0.1.2.4-fr]{0, 1.2.4}），而对 $D(g)\subset D(f)$ 验证
相容条件是立即的。于是，条件 \textit{d}'\,1)（相应地，
\textit{d}'\,2)）蕴含每个 $u_f$ 都是满射（相应地，单射），因而 $u$
是\emph{双射}，从而 $\mathscr{F}$ 是某个同构于 $\widetilde{M}$ 的
$\widetilde{A}$-模在 $V$ 上的限制。事实上，若
$s\in\Gamma(D(f)\cap V,\mathscr{F})$，则依 \textit{d}'\,1)，存在整数
$n\geq0$，使得 $f^n s$ 可延拓为一个截面 $z\in M$；于是
$u_f(z/f^n)=s$，故 $u_f$ 为满射。同样，若 $z\in M$ 满足
$u_f(z/1)=0$，这表示截面 $z$ 在 $D(f)\cap V$ 上的限制为零；依
\textit{d}'\,2)，存在整数 $n\geq0$，使得 $f^n z=0$，故 $M_f$ 中
$z/1=0$，从而 $u_f$ 为单射。

\par\hfill 证毕。\par

\begin{corollary}[1.4.2]
\label{I.1.4.2-fr}
$X$ 的任一拟紧开子集上的每个拟凝聚层，都是由 $X$ 上的某个拟凝聚层
诱导而来的。
\end{corollary}

\begin{corollary}[1.4.3]
\label{I.1.4.3-fr}
$X=\operatorname{Spec}(A)$ 上的每个拟凝聚 $\mathscr{O}_X$-代数，都同构于
形如 $\widetilde{B}$ 的 $\mathscr{O}_X$-代数，其中 $B$ 是一个 $A$-代数；
每个拟凝聚 $\widetilde{B}$-模，都同构于形如 $\widetilde{N}$ 的
$\widetilde{B}$-模，其中 $N$ 是一个 $B$-模。
\end{corollary}

事实上，一个拟凝聚 $\mathscr{O}_X$-代数是一个拟凝聚
$\mathscr{O}_X$-模，因而形如 $\widetilde{B}$，其中 $B$ 是一个 $A$-模；
$B$ 为 $A$-代数这一事实，由借助 $\widetilde{A}$-模同态
$\widetilde{B}\otimes_{\widetilde{A}}\widetilde{B}\to\widetilde{B}$ 对
$\mathscr{O}_X$-代数结构所作的刻画，以及
（\hyperref[I.1.3.12-fr]{1.3.12}）得出。若 $\mathscr{G}$ 是一个拟凝聚
$\widetilde{B}$-模，只须证明 $\mathscr{G}$ 也是一个拟凝聚
$\widetilde{A}$-模，便可再以同样方式得出结论。由于这个问题是局部的，
可以把 $X$ 限制到一个形如 $D(f)$ 的开集，并假设 $\mathscr{G}$ 是某个
$\widetilde{B}$-模同态
$\widetilde{B}^{(I)}\to\widetilde{B}^{(J)}$ 的余核（从而当然也是
$\widetilde{A}$-模同态的余核）；此时命题由
（\hyperref[I.1.3.8-fr]{1.3.8}）与
（\hyperref[I.1.3.9-fr]{1.3.9}）得出。

\subsection{素谱上的凝聚层}
\label{subsection:I.1.5-fr}

\begin{theorem}[1.5.1]
\label{I.1.5.1-fr}
设 $A$ 是一个 \emph{Noether} 环，$X=\operatorname{Spec}(A)$ 是它的素谱，
$V$ 是 $X$ 的一个开子集，$\mathscr{F}$ 是一个
$(\mathscr{O}_X|V)$-模。下列条件等价：
\begin{enumerate}
  \item[\textit{a})] $\mathscr{F}$ 是凝聚的。
  \item[\textit{b})] $\mathscr{F}$ 是有限型且拟凝聚的。
  \item[\textit{c})] 存在一个有限型 $A$-模 $M$，使得 $\mathscr{F}$
  同构于层 $\widetilde{M}|V$。
\end{enumerate}
\end{theorem}

\oldpage[I]{93}
\textit{a}) 显然推出 \textit{b})。为说明 \textit{b}) 推出
\textit{c})，先注意，由于 $V$ 拟紧
（\hyperref[0.2.2.3-fr]{0, 2.2.3}），$\mathscr{F}$ 同构于层
$\widetilde{N}|V$，其中 $N$ 是一个 $A$-模
（\hyperref[I.1.4.1-fr]{1.4.1}）。又有
$N=\varinjlim M_\lambda$，其中 $M_\lambda$ 遍历 $N$ 的所有有限型
$A$-子模，故由（\hyperref[I.1.3.9-fr]{1.3.9}），
$\mathscr{F}=\widetilde{N}|V=\varinjlim\widetilde{M_\lambda}|V$。但由于
$\mathscr{F}$ 是有限型的且 $V$ 拟紧，存在一个指标 $\lambda$，使得
$\mathscr{F}=\widetilde{M_\lambda}|V$
（\hyperref[0.5.2.3-fr]{0, 5.2.3}）。

最后证明 \textit{c}) 推出 \textit{a})。此时 $\mathscr{F}$ 显然是有限型的
（\hyperref[I.1.3.6-fr]{1.3.6} 与
\hyperref[I.1.3.9-fr]{1.3.9}）。此外，由于这个问题是局部的，可以只考虑
$V=D(f)$、$f\in A$ 的情形。由于 $A_f$ 是 Noether 环，最终只须证明：
若 $M$ 是一个 $A$-模，则同态
$\widetilde{A^n}\to\widetilde{M}$ 的核是有限型的。这样的同态必形如
$\widetilde{u}$，其中 $u$ 是一个同态 $A^n\to M$
（\hyperref[I.1.3.8-fr]{1.3.8}）；若 $P=\operatorname{Ker}u$，则
$\widetilde{P}=\operatorname{Ker}\widetilde{u}$
（\hyperref[I.1.3.9-fr]{1.3.9}）。由于 $A$ 是 Noether 环，$P$ 是有限型的，
证明至此完成。

\begin{corollary}[1.5.2]
\label{I.1.5.2-fr}
在（\hyperref[I.1.5.1-fr]{1.5.1}）的假设下，层 $\mathscr{O}_X$ 是一个
凝聚环层。
\end{corollary}

\begin{corollary}[1.5.3]
\label{I.1.5.3-fr}
在（\hyperref[I.1.5.1-fr]{1.5.1}）的假设下，$X$ 的一个开集上的任一凝聚层，
都是由 $X$ 上的某个凝聚层诱导而来的。
\end{corollary}

\begin{corollary}[1.5.4]
\label{I.1.5.4-fr}
在（\hyperref[I.1.5.1-fr]{1.5.1}）的假设下，每个拟凝聚
$\mathscr{O}_X$-模 $\mathscr{F}$，都是 $\mathscr{F}$ 的各凝聚
$\mathscr{O}_X$-子模的归纳极限。
\end{corollary}

事实上，$\mathscr{F}=\widetilde{M}$，其中 $M$ 是一个 $A$-模，而 $M$ 是其
各有限型子模的归纳极限；由（\hyperref[I.1.3.9-fr]{1.3.9}）与
（\hyperref[I.1.5.1-fr]{1.5.1}）即得结论。

\subsection{素谱上拟凝聚层的函子性质}
\label{subsection:I.1.6-fr}

\begin{env}[1.6.1]
\label{I.1.6.1-fr}
设 $A$、$A'$ 是两个环，
\[
  \varphi:A'\to A
\]
是一个同态，而
\[
  {}^a\varphi:X=\operatorname{Spec}(A)\to
  X'=\operatorname{Spec}(A')
\]
是与 $\varphi$ 相伴随的连续映射
（\hyperref[I.1.2.1-fr]{1.2.1}）。我们将定义一个环层的
\emph{典范同态}
\[
  \widetilde{\varphi}:\mathscr{O}_{X'}\to
  {}^a\varphi_*(\mathscr{O}_X)
\]
对每个 $f'\in A'$，令 $f=\varphi(f')$；则
${}^a\varphi^{-1}(D(f'))=D(f)$
（\hyperref[I.1.2.2.2-fr]{1.2.2.2}）。环
$\Gamma(D(f'),\widetilde{A'})$ 与 $\Gamma(D(f),\widetilde{A})$ 分别等同于
$A'_{f'}$ 与 $A_f$
（\hyperref[I.1.3.6-fr]{1.3.6} 与
\hyperref[I.1.3.7-fr]{1.3.7}）。而同态 $\varphi$ 典范地定义了同态
$\varphi_{f'}:A'_{f'}\to A_f$
（\hyperref[0.1.5.1-fr]{0, 1.5.1}）；换言之，有环同态
\[
  \Gamma(D(f'),\widetilde{A'})\to
  \Gamma({}^a\varphi^{-1}(D(f')),\widetilde{A})
  =\Gamma(D(f'),{}^a\varphi_*(\widetilde{A})).
\]

\oldpage[I]{94}
此外，这些同态满足通常的相容条件：当 $D(f')\supset D(g')$ 时，图
\[
  \xymatrix{
    \Gamma(D(f'),\widetilde{A'})\ar[r]\ar[d] &
    \Gamma(D(f'),{}^a\varphi_*(\widetilde{A}))\ar[d]\\
    \Gamma(D(g'),\widetilde{A'})\ar[r] &
    \Gamma(D(g'),{}^a\varphi_*(\widetilde{A}))
  }
\]
交换（\hyperref[0.1.5.1-fr]{0, 1.5.1}）。由于各 $D(f')$ 构成 $X'$ 的
拓扑的一组基（\hyperref[0.3.2.3-fr]{0, 3.2.3}），由此确实定义了一个
$\mathscr{O}_{X'}$-代数同态。因此，偶
$\Phi=({}^a\varphi,\widetilde{\varphi})$ 是一个赋环空间的\emph{态射}
\[
  \Phi:(X,\mathscr{O}_X)\to(X',\mathscr{O}_{X'})
\]
（\hyperref[0.4.1.1-fr]{0, 4.1.1}）。

还应注意，若令 $x'={}^a\varphi(x)$，则同态
$\widetilde{\varphi}_x^\#$（\hyperref[0.3.7.1-fr]{0, 3.7.1}）正是由
$\varphi:A'\to A$ 典范地导出的同态
\[
  \varphi_x:A'_{x'}\to A_x
\]
（\hyperref[0.1.5.1-fr]{0, 1.5.1}）。事实上，每个 $z'\in A'_{x'}$
都可写成 $g'/f'$，其中 $f'$、$g'$ 属于 $A'$ 且
$f'\notin\mathfrak{j}_{x'}$；故 $D(f')$ 是 $X'$ 中 $x'$ 的一个邻域，
而由 $\widetilde{\varphi}$ 导出的同态
$\Gamma(D(f'),\widetilde{A'})\to
\Gamma({}^a\varphi^{-1}(D(f')),\widetilde{A})$ 正是
$\varphi_{f'}$。考虑与 $g'/f'\in A'_{f'}$ 相对应的截面
$s'\in\Gamma(D(f'),\widetilde{A'})$，便得到 $A_x$ 中的等式
$\widetilde{\varphi}_x^\#(z')=\varphi(g')/\varphi(f')$。
\end{env}

\begin{example}[1.6.2]
\label{I.1.6.2-fr}
设 $S$ 是 $A$ 的一个乘法子集，$\varphi$ 是典范同态
$A\to S^{-1}A$。我们已经看到
（\hyperref[I.1.2.6-fr]{1.2.6}），${}^a\varphi$ 是从
$Y=\operatorname{Spec}(S^{-1}A)$ 到 $X=\operatorname{Spec}(A)$ 中满足
$\mathfrak{j}_x\cap S=\varnothing$ 的各点 $x$ 所成子空间的一个
\emph{同胚}。此外，对属于这个子空间的每个 $x$，它形如
${}^a\varphi(y)$，其中 $y\in Y$，同态
$\widetilde{\varphi}_y^\#:\mathscr{O}_x\to\mathscr{O}_y$ 是\emph{双射}
（\hyperref[0.1.2.6-fr]{0, 1.2.6}）；换言之，$\mathscr{O}_Y$ 等同于
$\mathscr{O}_X$ 在 $Y$ 上诱导的层。
\end{example}

\begin{proposition}[1.6.3]
\label{I.1.6.3-fr}
对每个 $A$-模 $M$，存在一个从 $\mathscr{O}_{X'}$-模
$(M_{[\varphi]})^{\sim}$ 到直像 $\Phi_*(\widetilde{M})$ 的典范函子性同构。
\end{proposition}

简记 $M'=M_{[\varphi]}$，并且对每个 $f'\in A'$ 仍令
$f=\varphi(f')$。截面模 $\Gamma(D(f'),\widetilde{M'})$ 与
$\Gamma(D(f),\widetilde{M})$ 分别等同于模 $M'_{f'}$ 与 $M_f$（分别作为
$A'_{f'}$-模与 $A_f$-模）。此外，$A'_{f'}$-模
$(M_f)_{[\varphi_{f'}]}$ 典范同构于 $M'_{f'}$
（\hyperref[0.1.5.2-fr]{0, 1.5.2}）。因此，有
$\Gamma(D(f'),\widetilde{A'})$-模的函子性同构
$\Gamma(D(f'),\widetilde{M'})\xrightarrow{\sim}
\Gamma({}^a\varphi^{-1}(D(f')),\widetilde{M})_{[\varphi_{f'}]}$，而这些
同构满足与限制映射之间通常的相容条件
（\hyperref[0.1.5.6-fr]{0, 1.5.6}），故定义了所断言的函子性同构。
更准确地说，若 $u:M_1\to M_2$ 是 $A$-模同态，则它也可视为
$A'$-模同态 $(M_1)_{[\varphi]}\to(M_2)_{[\varphi]}$；若将这个同态记为
$u_{[\varphi]}$，则 $\Phi_*(\widetilde{u})$ 等同于
$(u_{[\varphi]})^{\sim}$。

这个证明还表明，对每个 $A$-代数 $B$，典范函子性同构
\oldpage[I]{95}
$(B_{[\varphi]})^{\sim}\xrightarrow{\sim}\Phi_*(\widetilde{B})$ 是一个
$\mathscr{O}_{X'}$-代数同构；若 $M$ 是一个 $B$-模，则典范函子性同构
$(M_{[\varphi]})^{\sim}\xrightarrow{\sim}\Phi_*(\widetilde{M})$ 是一个
$\Phi_*(\widetilde{B})$-模同构。

\begin{corollary}[1.6.4]
\label{I.1.6.4-fr}
直像函子 $\Phi_*$ 在拟凝聚 $\mathscr{O}_X$-模范畴上是正合的。
\end{corollary}

事实上，$M_{[\varphi]}$ 关于 $M$ 是正合函子，而 $\widetilde{M'}$ 关于
$M'$ 是正合函子（\hyperref[I.1.3.5-fr]{1.3.5}）。

\begin{proposition}[1.6.5]
\label{I.1.6.5-fr}
设 $N'$ 是一个 $A'$-模，$N$ 是 $A$-模
$N'\otimes_{A'}A_{[\varphi]}$；则存在一个从 $\mathscr{O}_X$-模
$\Phi^*(\widetilde{N'})$ 到 $\widetilde{N}$ 的典范函子性同构。
\end{proposition}

先注意，$j:z'\mapsto z'\otimes 1$ 是从 $N'$ 到 $N_{[\varphi]}$ 的一个
$A'$-同态。事实上，按照定义，对 $f'\in A'$，有
$(f'z')\otimes 1=z'\otimes\varphi(f')=\varphi(f')(z'\otimes 1)$。由此根据
（\hyperref[I.1.3.8-fr]{1.3.8}）得到一个 $\mathscr{O}_{X'}$-模同态
$\widetilde{j}:\widetilde{N'}\to(N_{[\varphi]})^{\sim}$；依
（\hyperref[I.1.6.3-fr]{1.6.3}），可以把 $\widetilde{j}$ 看作从
$\widetilde{N'}$ 到 $\Phi_*(\widetilde{N})$ 的映射。与这个同态
$\widetilde{j}$ 典范对应的是从 $\Phi^*(\widetilde{N'})$ 到
$\widetilde{N}$ 的同态 $h=\widetilde{j}^{\#}$
（\hyperref[0.4.4.3-fr]{0, 4.4.3}）；我们将看到，对每个茎，$h_x$ 都是
\emph{双射}。令 $x'={}^a\varphi(x)$，并取满足 $x'\in D(f')$ 的
$f'\in A'$；令 $f=\varphi(f')$。环 $\Gamma(D(f),\widetilde{A})$
等同于 $A_f$，而模 $\Gamma(D(f),\widetilde{N})$ 与
$\Gamma(D(f'),\widetilde{N'})$ 分别等同于 $N_f$ 与 $N'_{f'}$。设
$s'\in\Gamma(D(f'),\widetilde{N'})$ 等同于 $n'/{f'}^p$（$n'\in N'$），
$s$ 是它在 $\Gamma(D(f),\widetilde{N})$ 中经 $\widetilde{j}$ 所得的像；
$s$ 等同于 $(n'\otimes 1)/f^p$。另一方面，设
$t\in\Gamma(D(f),\widetilde{A})$ 等同于 $g/f^q$（$g\in A$）；则按照
定义，有 $h_x(s'_{x'}\otimes t_x)=t_x\cdot s_x$
（\hyperref[0.4.4.3-fr]{0, 4.4.3}）。但可以把 $N_f$ 典范地等同于
$N'_{f'}\otimes_{A'_{f'}}(A_f)_{[\varphi_{f'}]}$
（\hyperref[0.1.5.4-fr]{0, 1.5.4}）；此时 $s$ 对应于元素
$(n'/{f'}^p)\otimes 1$，而截面 $y\mapsto t_y\cdot s_y$ 对应于
$(n'/{f'}^p)\otimes(g/f^q)$。（\hyperref[0.1.5.6-fr]{0, 1.5.6}）中的
相容图表明，$h_x$ 正是典范同构
\[
  \label{I.1.6.5.1-fr}
  N'_{x'}\otimes_{A'_{x'}}(A_x)_{[\varphi_x]}
  \xrightarrow{\sim}N_x=(N'\otimes_{A'}A_{[\varphi]})_x.
  \tag{1.6.5.1}
\]

此外，设 $v:N'_1\to N'_2$ 是一个 $A'$-模同态。由于对每个
$x'\in X'$ 都有 $\widetilde{v}_{x'}=v_{x'}$，由上文立即可知，
$\Phi^*(\widetilde{v})$ 典范地等同于 $(v\otimes 1)^{\sim}$；这就完成了
（\hyperref[I.1.6.5-fr]{1.6.5}）的证明。

若 $B'$ 是一个 $A'$-代数，则从 $\Phi^*(\widetilde{B'})$ 到
$(B'\otimes_{A'}A_{[\varphi]})^{\sim}$ 的典范同构是一个
$\mathscr{O}_X$-代数同构；若此外 $N'$ 是一个 $B'$-模，则从
$\Phi^*(\widetilde{N'})$ 到 $(N'\otimes_{A'}A_{[\varphi]})^{\sim}$ 的
典范同构是一个 $\Phi^*(\widetilde{B'})$-模同构。

\begin{corollary}[1.6.6]
\label{I.1.6.6-fr}
由各截面 $s'$ 的典范像所给出的 $\Phi^*(\widetilde{N'})$ 的截面，其中
$s'$ 遍历 $A'$-模 $\Gamma(\widetilde{N'})$，生成 $A$-模
$\Gamma(\Phi^*(\widetilde{N'}))$。
\end{corollary}

事实上，当分别把 $N'$ 与 $N$ 等同于 $\Gamma(\widetilde{N'})$ 与
$\Gamma(\widetilde{N})$（\hyperref[I.1.3.7-fr]{1.3.7}），并让 $z'$ 遍历
$N'$ 时，这些像就等同于 $N$ 中的各元素 $z'\otimes 1$。

\begin{env}[1.6.7]
\label{I.1.6.7-fr}
在（\hyperref[I.1.6.5-fr]{1.6.5}）的证明中，我们顺便证明了典范映射
（\hyperref[0.4.4.3.2-fr]{0, 4.4.3.2}）
$\rho:\widetilde{N'}\to\Phi_*(\Phi^*(\widetilde{N'}))$ 正是同态
$\widetilde{j}$，
\oldpage[I]{96}
其中 $j:N'\to N'\otimes_{A'}A_{[\varphi]}$ 是同态
$z'\mapsto z'\otimes 1$。同样，典范映射
（\hyperref[0.4.4.3.3-fr]{0, 4.4.3.3}）
$\sigma:\Phi^*(\Phi_*(\widetilde{M}))\to\widetilde{M}$ 正是
$\widetilde{p}$，其中
$p:M_{[\varphi]}\otimes_{A'}A_{[\varphi]}\to M$ 是典范同态，它把每个
张量积 $z\otimes a$（$z\in M$，$a\in A$）映到 $a\cdot z$；这由定义
（\hyperref[0.3.7.1-fr]{0, 3.7.1}、
\hyperref[0.4.4.3-fr]{0, 4.4.3} 与
\hyperref[I.1.3.7-fr]{I, 1.3.7}）立即得出。

由此可知（\hyperref[0.4.4.3-fr]{0, 4.4.3} 与
\hyperref[0.3.5.4.4-fr]{3.5.4.4}），若
$v:N'\to M_{[\varphi]}$ 是一个 $A'$-同态，则
$\widetilde{v}^{\#}=(v\otimes 1)^{\sim}$。
\end{env}

\begin{env}[1.6.8]
\label{I.1.6.8-fr}
设 $N'_1$、$N'_2$ 是两个 $A'$-模，并假设 $N'_1$ 有一个\emph{有限表示}。
由（\hyperref[I.1.6.7-fr]{1.6.7}）与
（\hyperref[I.1.3.12-fr]{1.3.12, (ii)}）可知，典范同态
（\hyperref[0.4.4.6-fr]{0, 4.4.6}）
\[
  \Phi^*\bigl(\mathscr{H}\!om_{\widetilde{A'}}
    (\widetilde{N'_1},\widetilde{N'_2})\bigr)
  \longrightarrow
  \mathscr{H}\!om_{\widetilde A}
    \bigl(\Phi^*(\widetilde{N'_1}),\Phi^*(\widetilde{N'_2})\bigr)
\]
正是 $\widetilde{\gamma}$；这里 $\gamma$ 表示 $A$-模的典范同态
\[
  \operatorname{Hom}_{A'}(N'_1,N'_2)\otimes_{A'}A
  \longrightarrow
  \operatorname{Hom}_A(N'_1\otimes_{A'}A,N'_2\otimes_{A'}A).
\]
\end{env}

\begin{env}[1.6.9]
\label{I.1.6.9-fr}
设 $\mathfrak{J}'$ 是 $A'$ 的一个理想，$M$ 是一个 $A$-模。按照定义，
$\widetilde{\mathfrak{J}'}\widetilde M$ 是典范同态
$\Phi^*(\widetilde{\mathfrak{J}'})\otimes_{\widetilde A}\widetilde M
\to\widetilde M$ 的像，故由（\hyperref[I.1.6.5-fr]{1.6.5}）与
（\hyperref[I.1.3.12-fr]{1.3.12, (i)}），
$\widetilde{\mathfrak{J}'}\widetilde M$ 典范地等同于
$(\mathfrak{J}'M)^{\sim}$。特别地，
$\Phi^*(\widetilde{\mathfrak{J}'})\widetilde A$ 等同于
$(\mathfrak{J}'A)^{\sim}$；再考虑到函子 $\Phi^*$ 的右正合性，
$\widetilde A$-代数 $\Phi^*((A'/\mathfrak{J}')^{\sim})$ 等同于
$(A/\mathfrak{J}'A)^{\sim}$。
\end{env}

\begin{env}[1.6.10]
\label{I.1.6.10-fr}
设 $A''$ 是第三个环，$\varphi'$ 是同态 $A''\to A'$，并令
$\varphi''=\varphi\circ\varphi'$。由定义立即可得
${}^a\varphi''=({}^a\varphi')\circ({}^a\varphi)$ 以及
$\widetilde{\varphi''}=\widetilde\varphi\circ\widetilde{\varphi'}$
（\hyperref[0.1.5.7-fr]{0, 1.5.7}）。因而有
$\Phi''=\Phi'\circ\Phi$；换言之，$(\operatorname{Spec} A,\widetilde A)$
是从环范畴到赋环空间范畴的一个关于 $A$ 的\emph{函子}。
\end{env}

\subsection{仿射概形态射的刻画}
\label{subsection:I.1.7-fr}

\begin{definition}[1.7.1]
\label{I.1.7.1-fr}
若一个赋环空间 $(X,\mathscr{O}_X)$ 同构于形如
$(\operatorname{Spec}(A),\widetilde A)$ 的赋环空间，其中 $A$ 是一个环，
就称它为一个\emph{仿射概形}。此时，$\Gamma(X,\mathscr{O}_X)$ 典范地
等同于环 $A$（\hyperref[I.1.3.7-fr]{1.3.7}），称为仿射概形
$(X,\mathscr{O}_X)$ 的环；在不会引起混淆时，记为 $A(X)$。
\end{definition}

按照语言上的滥用，当我们谈到\emph{仿射概形} $\operatorname{Spec}(A)$ 时，
总是指赋环空间 $(\operatorname{Spec}(A),\widetilde A)$。

\begin{env}[1.7.2]
\label{I.1.7.2-fr}
设 $A$、$B$ 是两个环，$(X,\mathscr{O}_X)$、$(Y,\mathscr{O}_Y)$ 是分别与
素谱 $X=\operatorname{Spec}(A)$、$Y=\operatorname{Spec}(B)$ 相应的仿射概形。
我们已经看到（\hyperref[I.1.6.1-fr]{1.6.1}），每个环同态
$\varphi:B\to A$ 都对应于一个态射
$\Phi=({}^a\varphi,\widetilde\varphi)=\operatorname{Spec}(\varphi):
(X,\mathscr{O}_X)\to(Y,\mathscr{O}_Y)$。应注意，$\varphi$ 完全由 $\Phi$
决定，因为按照定义有
\[
  \varphi=\Gamma(\widetilde\varphi):
  \Gamma(\widetilde B)\longrightarrow
  \Gamma({}^a\varphi_*(\widetilde A))=\Gamma(\widetilde A).
\]
\end{env}

\begin{theorem}[1.7.3]
\label{I.1.7.3-fr}
设 $(X,\mathscr{O}_X)$、$(Y,\mathscr{O}_Y)$ 是两个仿射概形。赋环空间的态射
$(\psi,\theta):(X,\mathscr{O}_X)\to(Y,\mathscr{O}_Y)$ 形如
$({}^a\varphi,\widetilde\varphi)$，其中 $\varphi$ 是环同态
$A(Y)\to A(X)$，当且仅当对每个 $x\in X$，$\theta_x^{\#}$ 都是局部同态
$\mathscr{O}_{\psi(x)}\to\mathscr{O}_x$。
\end{theorem}

\oldpage[I]{97}
令 $A=A(X)$、$B=A(Y)$。这一条件是必要的，因为我们已经看到
（\hyperref[I.1.6.1-fr]{1.6.1}），$\widetilde\varphi_x^{\#}$ 是由
$\varphi$ 典范地导出的从 $B_{{}^a\varphi(x)}$ 到 $A_x$ 的同态；而按照
${}^a\varphi(x)=\varphi^{-1}(\mathfrak{j}_x)$ 的定义，这个同态是局部的。

证明这一条件是充分的。按照定义，$\theta$ 是一个同态
$\mathscr{O}_Y\to\psi_*(\mathscr{O}_X)$，由此典范地得到环同态
\[
  \varphi=\Gamma(\theta):
  B=\Gamma(Y,\mathscr{O}_Y)
  \longrightarrow
  \Gamma(Y,\psi_*(\mathscr{O}_X))=\Gamma(X,\mathscr{O}_X)=A.
\]

关于 $\theta_x^{\#}$ 的假设表明，取商以后，这个同态导出从剩余域
$k(\psi(x))$ 到剩余域 $k(x)$ 的单同态 $\theta^x$，并且对每个截面
$f\in\Gamma(Y,\mathscr{O}_Y)=B$，都有
$\theta^x(f(\psi(x)))=\varphi(f)(x)$。因此，关系 $f(\psi(x))=0$ 等价于
$\varphi(f)(x)=0$，这就是说
$\mathfrak{j}_{\psi(x)}=\mathfrak{j}_{{}^a\varphi(x)}$；也可写成：对每个
$x\in X$，都有 $\psi(x)={}^a\varphi(x)$，即 $\psi={}^a\varphi$。我们还知道图
\phantomsection\label{I.1.7.3.diagram-fr}
\[
  \xymatrix{
    B=\Gamma(Y,\mathscr{O}_Y)\ar[r]^\varphi\ar[d] &
    \Gamma(X,\mathscr{O}_X)=A\ar[d]\\
    B_{\psi(x)}\ar[r]_{\theta_x^{\#}} & A_x
  }
\]
交换（\hyperref[0.3.7.2-fr]{0, 3.7.2}）；这就是说，$\theta_x^{\#}$ 等于由
$\varphi$ 典范地导出的同态 $\varphi_x:B_{\psi(x)}\to A_x$
（\hyperref[0.1.5.1-fr]{0, 1.5.1}）。由于各 $\theta_x^{\#}$ 完全刻画
$\theta^{\#}$，从而也完全刻画 $\theta$
（\hyperref[0.3.7.1-fr]{0, 3.7.1}），由 $\widetilde\varphi$ 的定义
（\hyperref[I.1.6.1-fr]{1.6.1}），可知 $\theta=\widetilde\varphi$。

我们称满足（\hyperref[I.1.7.3-fr]{1.7.3}）中条件的赋环空间态射
$(\psi,\theta)$ 为一个\emph{仿射概形的态射}。

\begin{corollary}[1.7.4]
\label{I.1.7.4-fr}
若 $(X,\mathscr{O}_X)$、$(Y,\mathscr{O}_Y)$ 是两个仿射概形，则仿射概形的
态射集 $\operatorname{Hom}((X,\mathscr{O}_X),(Y,\mathscr{O}_Y))$ 与从 $B$
到 $A$ 的环同态集之间存在一个典范同构，其中
$A=\Gamma(\mathscr{O}_X)$、$B=\Gamma(\mathscr{O}_Y)$。
\end{corollary}

还可以说，关于 $A$ 的函子 $(\operatorname{Spec}(A),\widetilde A)$ 与关于
$(X,\mathscr{O}_X)$ 的函子 $\Gamma(X,\mathscr{O}_X)$，给出了交换环范畴与
仿射概形范畴的对偶范畴之间的一个\emph{等价}（T, I, 1.2）。

\begin{corollary}[1.7.5]
\label{I.1.7.5-fr}
若 $\varphi:B\to A$ 是满射，则相应的态射
$({}^a\varphi,\widetilde\varphi)$ 是赋环空间的一个单态射
\emph{（参见（\hyperref[I.4.1.7-fr]{4.1.7}））}。
\end{corollary}

事实上，已知 ${}^a\varphi$ 是单射（\hyperref[I.1.2.5-fr]{1.2.5}）；又因
$\varphi$ 是满射，对每个 $x\in X$，从 $\varphi$ 通过转到分式环而导出的
$\varphi_x^{\#}:B_{{}^a\varphi(x)}\to A_x$ 也是满射
（\hyperref[0.1.5.1-fr]{0, 1.5.1}）；结论由
（\hyperref[0.4.1.1-fr]{0, 4.1.1}）得出。
% EGA I mainland Simplified-Chinese translation (producer output; uncertified).
% Sequential translation cursor: complete through 2.5.5 (source lines 1--592; printed pages 98--104).

\section{预概形与预概形的态射}
\label{section:I.2-fr}

\subsection{预概形的定义}
\label{subsection:I.2.1-fr}

\begin{env}[2.1.1]
\label{I.2.1.1-fr}
给定一个赋环空间 $(X,\mathscr{O}_X)$，若 $X$ 的一个开子集 $V$ 所给出的
赋环空间 $(V,\mathscr{O}_X|V)$ 是仿射概形
（\hyperref[I.1.7.1-fr]{1.7.1}），就称这个开子集为\emph{仿射开集}。
\end{env}

\begin{definition}[2.1.2]
\label{I.2.1.2-fr}
若赋环空间 $(X,\mathscr{O}_X)$ 的每个点都有一个仿射开邻域，就称它为一个
\emph{预概形}。
\end{definition}

\oldpage[I]{98}

\begin{proposition}[2.1.3]
\label{I.2.1.3-fr}
若 $(X,\mathscr{O}_X)$ 是一个预概形，则各仿射开集构成 $X$ 的拓扑的一组基。
\end{proposition}

事实上，若 $V$ 是 $x\in X$ 的任一开邻域，则按照假设，存在 $x$ 的一个
开邻域 $W$，使得 $(W,\mathscr{O}_X|W)$ 是仿射概形；以 $A$ 表示它的环。
在空间 $W$ 中，$V\cap W$ 是 $x$ 的一个开邻域；因此存在 $f\in A$，使得
$D(f)$ 是包含于 $V\cap W$ 的 $x$ 的一个开邻域
（\hyperref[I.1.1.10-fr]{1.1.10 (i)}）。于是，赋环空间
$(D(f),\mathscr{O}_X|D(f))$ 是一个环同构于 $A_f$ 的仿射概形
（\hyperref[I.1.3.6-fr]{1.3.6}），命题得证。

\begin{proposition}[2.1.4]
\label{I.2.1.4-fr}
预概形 $X$ 的底空间是一个 Kolmogorov 空间。
\end{proposition}

事实上，若 $x,y$ 是预概形 $X$ 的两个不同点，并且 $x,y$ 不属于同一个仿射
开集，则显然存在其中一点的一个开邻域不包含另一点；若它们属于同一个
仿射开集，则结论由（\hyperref[I.1.1.8-fr]{1.1.8}）得出。

\begin{proposition}[2.1.5]
\label{I.2.1.5-fr}
若 $(X,\mathscr{O}_X)$ 是一个预概形，则 $X$ 的每个不可约闭子集都恰有一个
一般点；因此，映射 $x\longmapsto\overline{\{x\}}$ 是从 $X$ 到其不可约闭
子集全体的一个双射。
\end{proposition}

事实上，若 $Y$ 是 $X$ 的一个不可约闭子集，$y\in Y$，而 $U$ 是 $X$ 中
$y$ 的一个仿射开邻域，则 $U\cap Y$ 在 $Y$ 中稠密且不可约
（\hyperref[0.2.1.1-fr]{0, 2.1.1} 与
\hyperref[0.2.1.4-fr]{2.1.4}）。因此，由
（\hyperref[I.1.1.14-fr]{1.1.14}），$U\cap Y$ 是某点 $x$ 在 $U$ 中的
闭包，进而 $Y\subset\overline{U}$ 是 $x$ 在 $X$ 中的闭包。
% EGA1-ERR-0025: source/print says X here; retained literally pending canon disposition.
$X$ 的一般点的唯一性由（\hyperref[I.2.1.4-fr]{2.1.4}）与
（\hyperref[0.2.1.3-fr]{0, 2.1.3}）得出。

\begin{env}[2.1.6]
\label{I.2.1.6-fr}
若 $Y$ 是 $X$ 的一个不可约闭子集，$y$ 是它的一般点，则局部环
$\mathscr{O}_y$ 也记作 $\mathscr{O}_{X/Y}$，称为\emph{$X$ 沿 $Y$ 的局部环}，
或\emph{$Y$ 在 $X$ 中的局部环}。

若 $X$ 本身不可约，而 $x$ 是它的一般点，则也称 $\mathscr{O}_x$ 为
\emph{$X$ 上的有理函数环}（参见 \hyperref[section:I.7-fr]{§~7}）。
\end{env}

\begin{proposition}[2.1.7]
\label{I.2.1.7-fr}
若 $(X,\mathscr{O}_X)$ 是一个预概形，则对 $X$ 的每个开子集 $U$，赋环空间
$(U,\mathscr{O}_X|U)$ 都是一个预概形。
\end{proposition}

这由定义（\hyperref[I.2.1.2-fr]{2.1.2}）与命题
（\hyperref[I.2.1.3-fr]{2.1.3}）立即得出。

称 $(U,\mathscr{O}_X|U)$ 为 $(X,\mathscr{O}_X)$ 在 $U$ 上\emph{诱导}的
预概形，或 $(X,\mathscr{O}_X)$ 在 $U$ 上的\emph{限制}。

\begin{env}[2.1.8]
\label{I.2.1.8-fr}
若预概形 $(X,\mathscr{O}_X)$ 的底空间 $X$ 不可约（相应地，连通），就称
它\emph{不可约}（相应地，\emph{连通}）。若预概形\emph{不可约且既约}
（参见（\hyperref[I.5.1.4-fr]{5.1.4}）），就称它为\emph{整预概形}。
若预概形 $(X,\mathscr{O}_X)$ 的每个点 $x\in X$ 都有一个开邻域 $U$，使得
$(X,\mathscr{O}_X)$ 在 $U$ 上诱导的预概形是整预概形，就称
该预概形\emph{局部为整的}。
\end{env}

\subsection{预概形的态射}
\label{subsection:I.2.2-fr}
\label{I.2.2-fr}

\begin{definition}[2.2.1]
\label{I.2.2.1-fr}
给定两个预概形 $(X,\mathscr{O}_X)$、$(Y,\mathscr{O}_Y)$，凡是满足下述条件的
赋环空间态射 $(\psi,\theta)$，都称为从 $(X,\mathscr{O}_X)$ 到
$(Y,\mathscr{O}_Y)$ 的一个\emph{态射（预概形的态射）}：对每个 $x\in X$，
$\theta_x^\sharp$ 都是一个局部同态
$\mathscr{O}_{\psi(x)}\longrightarrow\mathscr{O}_x$。
\end{definition}

\oldpage[I]{99}

取商以后，映射 $\theta_x^\sharp:\mathscr{O}_{\psi(x)}\to\mathscr{O}_x$
因而给出一个单同态 $\theta^x:k(\psi(x))\to k(x)$；由此可以把 $k(x)$ 看作
域 $k(\psi(x))$ 的一个\emph{扩张}。

\begin{env}[2.2.2]
\label{I.2.2.2-fr}
两个预概形态射 $(\psi,\theta)$ 与 $(\psi',\theta')$ 的复合
$(\psi'',\theta'')$ 仍是预概形态射，这是公式
${\theta''}^\sharp=\theta^\sharp\circ\psi^*({\theta'}^\sharp)$
（\hyperref[0.3.5.5-fr]{0, 3.5.5}）的直接结果。由此可知，预概形组成一个
\emph{范畴}；按照一般记号，用 $\operatorname{Hom}(X,Y)$ 表示从预概形 $X$
到预概形 $Y$ 的所有态射所成的集合。
\end{env}

\begin{example}[2.2.3]
\label{I.2.2.3-fr}
若 $U$ 是 $X$ 的一个开子集，则诱导预概形 $(U,\mathscr{O}_X|U)$ 到
$(X,\mathscr{O}_X)$ 的典范嵌入（\hyperref[0.4.1.2-fr]{0, 4.1.2}）是一个
预概形态射；而且，由（\hyperref[0.4.1.1-fr]{0, 4.1.1}）立即可知，它是
赋环空间的一个\emph{单态射}（从而更是预概形的一个单态射）。
\end{example}

\begin{proposition}[2.2.4]
\label{I.2.2.4-fr}
设 $(X,\mathscr{O}_X)$ 是一个预概形，$(S,\mathscr{O}_S)$ 是环 $A$ 所对应的
一个仿射概形。则从预概形 $(X,\mathscr{O}_X)$ 到预概形
$(S,\mathscr{O}_S)$ 的态射，与从 $A$ 到环 $\Gamma(X,\mathscr{O}_X)$ 的
同态之间存在一个典范双射对应。
\end{proposition}

先注意，若 $(X,\mathscr{O}_X)$ 与 $(Y,\mathscr{O}_Y)$ 是任意两个赋环空间，
则从 $(X,\mathscr{O}_X)$ 到 $(Y,\mathscr{O}_Y)$ 的一个态射
$(\psi,\theta)$ 典范地定义一个环同态
$\Gamma(\theta):\Gamma(Y,\mathscr{O}_Y)\to
\Gamma(Y,\psi_*(\mathscr{O}_X))=\Gamma(X,\mathscr{O}_X)$。在这里所考虑的
情形，一切归结为证明：任意同态
$\varphi:A\to\Gamma(X,\mathscr{O}_X)$ 都恰好以一种方式写成
$\Gamma(\theta)$，也就是说，恰好由一个 $\theta$ 给出。事实上，依假设，
$X$ 有一个由仿射开集组成的覆盖
$(V_\alpha)$；将 $\varphi$ 与限制同态
$\Gamma(X,\mathscr{O}_X)\to
\Gamma(V_\alpha,\mathscr{O}_X|V_\alpha)$ 复合，得到同态
$\varphi_\alpha:A\to\Gamma(V_\alpha,\mathscr{O}_X|V_\alpha)$；根据
（\hyperref[I.1.7.3-fr]{1.7.3}），它对应于从预概形
$(V_\alpha,\mathscr{O}_X|V_\alpha)$ 到 $(S,\mathscr{O}_S)$ 的唯一态射
$(\psi_\alpha,\theta_\alpha)$。此外，对每一对指标 $(\alpha,\beta)$，
$V_\alpha\cap V_\beta$ 的每个点都有一个包含于 $V_\alpha\cap V_\beta$ 的
仿射开邻域 $W$（\hyperref[I.2.1.3-fr]{2.1.3}）。显然，将
$\varphi_\alpha$ 与 $\varphi_\beta$ 分别同到 $W$ 的限制同态复合，所得的是
同一个同态
$\Gamma(S,\mathscr{O}_S)\to\Gamma(W,\mathscr{O}_X|W)$。因此，由关系
$(\theta_\alpha^\sharp)_x=(\varphi_\alpha)_x$（对所有 $x\in V_\alpha$
及所有 $\alpha$）（\hyperref[I.1.6.1-fr]{1.6.1}），态射
$(\psi_\alpha,\theta_\alpha)$ 与 $(\psi_\beta,\theta_\beta)$ 在 $W$ 上的
限制相同。由此可知，存在唯一的赋环空间态射
$(\psi,\theta):(X,\mathscr{O}_X)\to(S,\mathscr{O}_S)$，使其在每个
$V_\alpha$ 上的限制都是 $(\psi_\alpha,\theta_\alpha)$；显然，它是一个
预概形态射，并满足 $\Gamma(\theta)=\varphi$。

设 $u:A\to\Gamma(X,\mathscr{O}_X)$ 是一个环同态，而
$v=(\psi,\theta)$ 是相应的态射
$(X,\mathscr{O}_X)\to(S,\mathscr{O}_S)$。对每个 $f\in A$，有
\begin{equation}
  \psi^{-1}(D(f))=X_{u(f)}
  \tag{2.2.4.1}\label{I.2.2.4.1-fr}
\end{equation}
这里采用（\hyperref[0.5.5.2-fr]{0, 5.5.2}）中关于局部自由层
$\mathscr{O}_X$ 的记号。事实上，只需在 $X$ 本身是仿射的情形验证这一公式，
而此时它正是（\hyperref[I.1.2.2.2-fr]{1.2.2.2}）。

\begin{proposition}[2.2.5]
\label{I.2.2.5-fr}
在（\hyperref[I.2.2.4-fr]{2.2.4}）的假设下，设
$\varphi:A\to\Gamma(X,\mathscr{O}_X)$ 是一个环同态，
$f:(X,\mathscr{O}_X)\to(S,\mathscr{O}_S)$ 是相应的预概形态射，
$\mathscr{G}$（相应地，$\mathscr{F}$）是一个拟凝聚
$\mathscr{O}_X$-模（相应地，拟凝聚 $\mathscr{O}_S$-模），并令
$M=\Gamma(S,\mathscr{F})$。则 $f$-态射
\oldpage[I]{100}
$\mathscr{F}\to\mathscr{G}$（\hyperref[0.4.4.1-fr]{0, 4.4.1}）与
$A$-同态 $M\to(\Gamma(X,\mathscr{G}))_{[\varphi]}$ 之间存在一个典范双射对应。
\end{proposition}

事实上，像在（\hyperref[I.2.2.4-fr]{2.2.4}）中那样论证，立即归结到
$X$ 为仿射的情形；于是命题由（\hyperref[I.1.6.3-fr]{1.6.3}）与
（\hyperref[I.1.3.8-fr]{1.3.8}）得出。

\begin{env}[2.2.6]
\label{I.2.2.6-fr}
若对 $X$ 的每个开子集 $U$（相应地，$X$ 的每个闭子集 $F$），$\psi(U)$ 在
$Y$ 中是开集（相应地，$\psi(F)$ 在 $Y$ 中是闭集），就称预概形态射
$(\psi,\theta):(X,\mathscr{O}_X)\to(Y,\mathscr{O}_Y)$ 是\emph{开态射}
（相应地，\emph{闭态射}）。若 $\psi(X)$ 在 $Y$ 中稠密，就称
$(\psi,\theta)$ 是\emph{支配态射}；若 $\psi$ 是满射，就称它是
\emph{满态射}。应注意，这些条件只涉及连续映射 $\psi$。
\end{env}

\begin{proposition}[2.2.7]
\label{I.2.2.7-fr}
设
\[
  f=(\psi,\theta):(X,\mathscr{O}_X)\to(Y,\mathscr{O}_Y),
  \qquad
  g=(\psi',\theta'):(Y,\mathscr{O}_Y)\to(Z,\mathscr{O}_Z)
\]
是两个预概形态射。
\begin{enumerate}
  \item[\rm(i)] 若 $f$ 与 $g$ 都是开态射（相应地，闭态射、支配态射、
    满态射），则 $g\circ f$ 也是如此。
  \item[\rm(ii)] 若 $f$ 是满态射且 $g\circ f$ 是闭态射，则 $g$ 是闭态射。
  \item[\rm(iii)] 若 $g\circ f$ 是满态射，则 $g$ 是满态射。
\end{enumerate}
\end{proposition}

断言 (i) 与 (iii) 是显然的。记 $g\circ f=(\psi'',\theta'')$。若 $F$ 是
$Y$ 中的闭集，则 $\psi^{-1}(F)$ 是 $X$ 中的闭集，故
$\psi''(\psi^{-1}(F))$ 是 $Z$ 中的闭集；但因 $\psi$ 是满射，
$\psi(\psi^{-1}(F))=F$，所以
$\psi''(\psi^{-1}(F))=\psi'(F)$，这就证明了 (ii)。

\begin{proposition}[2.2.8]
\label{I.2.2.8-fr}
设 $f=(\psi,\theta)$ 是态射
$(X,\mathscr{O}_X)\to(Y,\mathscr{O}_Y)$，$(U_\alpha)$ 是 $Y$ 的一个
开覆盖。则 $f$ 是开态射（相应地，闭态射、满态射、支配态射）的充要条件是：
对每个诱导预概形
$(\psi^{-1}(U_\alpha),\mathscr{O}_X|\psi^{-1}(U_\alpha))$，把该态射在其上的
限制看作从该诱导预概形到诱导预概形
$(U_\alpha,\mathscr{O}_Y|U_\alpha)$ 的态射时，这一限制是开态射（相应地，
闭态射、满态射、支配态射）。
\end{proposition}

该命题由定义立即得出，只需注意：$Y$ 的一个子集 $F$ 在 $Y$ 中闭
（相应地，开、稠密）的充要条件是，每个 $F\cap U_\alpha$ 在
$U_\alpha$ 中闭（相应地，开、稠密）。

\begin{env}[2.2.9]
\label{I.2.2.9-fr}
设 $(X,\mathscr{O}_X)$、$(Y,\mathscr{O}_Y)$ 是两个预概形；假设 $X$ 与
$Y$ 各自恰有相同有限数目的不可约分支 $X_i$（相应地，$Y_i$）
$(1\leq i\leq n)$；令 $\xi_i$（相应地，$\eta_i$）为 $X_i$
（相应地，$Y_i$）的一般点（\hyperref[I.2.1.5-fr]{2.1.5}）。若态射
\[
  f=(\psi,\theta):(X,\mathscr{O}_X)\to(Y,\mathscr{O}_Y)
\]
对每个 $i$ 都满足 $\psi^{-1}(\eta_i)=\{\xi_i\}$，并且
$\theta_{\xi_i}^{\sharp}:\mathscr{O}_{\eta_i}\to\mathscr{O}_{\xi_i}$ 是
一个\emph{同构}，就称它是\emph{双有理态射}。显然，双有理态射是
\emph{支配态射}（\hyperref[0.2.1.8-fr]{0, 2.1.8}），因而若它还是闭态射，
它就是满态射。
\end{env}

\begin{notation}[2.2.10]
\label{I.2.2.10-fr}
在本书以下各处，只要不会引起混淆，我们将在预概形（相应地，态射）的记号中
\emph{省略}结构层（相应地，结构层的态射）。若 $U$ 是预概形 $X$ 的底空间的
一个开子集，则每当把 $U$ 作为预概形谈及时，始终指 $U$ 上的诱导预概形。
\end{notation}

\subsection{预概形的黏合}
\label{subsection:I.2.3-fr}

\begin{env}[2.3.1]
\label{I.2.3.1-fr}
\oldpage[I]{101}
由定义（\hyperref[I.2.1.2-fr]{2.1.2}）可知，把预概形\emph{黏合}起来
（\hyperref[0.4.1.6-fr]{0, 4.1.6}）所得到的任何赋环空间仍是一个预概形。
特别地，因为每个预概形按定义都有一个由仿射开集组成的覆盖，所以每个预概形
都可以通过\emph{黏合仿射概形}得到。
\end{env}

\begin{example}[2.3.2]
\label{I.2.3.2-fr}
设 $K$ 是一个域，$B=K[s]$、$C=K[t]$ 是 $K$ 上两个单变元多项式环，
并设 $X_1=\operatorname{Spec}(B)$、$X_2=\operatorname{Spec}(C)$；它们是两个
同构的仿射概形。在 $X_1$（相应地，$X_2$）中，令 $U_{12}$（相应地，
$U_{21}$）为仿射开集 $D(s)$（相应地，$D(t)$），其环 $B_s$（相应地，
$C_t$）由形如 $f(s)/s^m$（相应地，$g(t)/t^n$）的有理分式组成，其中
$f\in B$（相应地，$g\in C$）。令 $u_{12}$ 为预概形同构
$U_{21}\to U_{12}$；根据（\hyperref[I.2.2.4-fr]{2.2.4}），它对应于从
$B_s$ 到 $C_t$ 的同构，而该同构把 $f(s)/s^m$ 映为有理分式
$f(1/t)/(1/t^m)$。可以借助 $u_{12}$，沿 $U_{12}$ 与 $U_{21}$ 黏合
$X_1$ 与 $X_2$，因为显然没有黏合条件。稍后我们将再次遇到这样得到的
预概形 $X$，它是一般构造方法的一个特例
（\hyperref[II.2.4.3-fr]{II, 2.4.3}）。这里只证明 $X$\emph{不是仿射概形}；
这是因为环 $\Gamma(X,\mathscr{O}_X)$\emph{同构于} $K$，故其谱缩为一个点。
事实上，$\mathscr{O}_X$ 在 $X$ 上的一个截面限制到 $X_1$（相应地，
$X_2$）上时——这里把它们视为 $X$ 的仿射开集——是多项式 $f(s)$
（相应地，$g(t)$）；由定义必须有 $g(t)=f(1/t)$，而这只有在
$f=g\in K$ 时才可能。
\end{example}

\subsection{局部概形}
\label{subsection:I.2.4-fr}

\begin{env}[2.4.1]
\label{I.2.4.1-fr}
环 $A$ 为局部环的仿射概形称为\emph{局部概形}；此时
$X=\operatorname{Spec}(A)$ 中只有一个\emph{闭点} $a$，而对每个其他点
$b\in X$，都有 $a\in\overline{\{b\}}$
（\hyperref[I.1.1.7-fr]{1.1.7}）。
\end{env}

对每个预概形 $Y$ 及每个点 $y\in Y$，局部概形
$\operatorname{Spec}(\mathscr{O}_y)$ 称为\emph{$Y$ 在点 $y$ 处的局部概形}。
设 $V$ 是 $Y$ 的一个包含 $y$ 的仿射开集，$B$ 是仿射概形 $V$ 的环；
$\mathscr{O}_y$ 典范地等同于 $B_y$（\hyperref[I.1.3.4-fr]{1.3.4}），
因此典范同态 $B\to B_y$ 对应于一个预概形态射
$\operatorname{Spec}(\mathscr{O}_y)\to V$
（\hyperref[I.1.6.1-fr]{1.6.1}）。将此态射与典范嵌入 $V\to Y$ 复合，
就得到一个态射 $\operatorname{Spec}(\mathscr{O}_y)\to Y$；它\emph{不依赖于}
所选择的包含 $y$ 的仿射开集 $V$。事实上，若 $V'$ 是另一个包含 $y$ 的
仿射开集，则存在第三个包含 $y$ 的仿射开集 $W$，使得
$W\subset V\cap V'$（\hyperref[I.2.1.3-fr]{2.1.3}）；因而只需考察
$V\subset V'$ 的情形。若 $B'$ 是 $V'$ 的环，只需注意下图
\phantomsection\label{I.2.4.1.diagram-fr}
\[
  \xymatrix{
    B'\ar[rr]\ar[dr] & & B\ar[dl]\\
    & \mathscr{O}_y
  }
\]
交换（\hyperref[0.1.5.1-fr]{0, 1.5.1}）。这样定义的态射
\[
  \operatorname{Spec}(\mathscr{O}_y)\to Y
\]
称为\emph{典范态射}。

\begin{proposition}[2.4.2]
\label{I.2.4.2-fr}
\oldpage[I]{102}
设 $(Y,\mathscr{O}_Y)$ 是一个预概形；对每个 $y\in Y$，令
$(\psi,\theta)$ 为典范态射
$(\operatorname{Spec}(\mathscr{O}_y),\widetilde{\mathscr{O}}_y)
\to(Y,\mathscr{O}_Y)$。则 $\psi$ 是从
$\operatorname{Spec}(\mathscr{O}_y)$ 到 $Y$ 的子空间 $S_y$ 的同胚，其中
该子空间由满足 $y\in\overline{\{z\}}$ 的点 $z$ 组成（换言之，就是 $y$ 的
\emph{一般化点}（\hyperref[0.2.1.2-fr]{0, 2.1.2}））；此外，若
$z=\psi(\mathfrak{p})$，则
$\theta_z^{\sharp}:\mathscr{O}_z\to(\mathscr{O}_y)_{\mathfrak{p}}$ 是同构；
所以 $(\psi,\theta)$ 是赋环空间的一个单态射。
\end{proposition}

因为 $\operatorname{Spec}(\mathscr{O}_y)$ 的唯一闭点 $a$ 属于该空间中
每一点的闭包，并且 $\psi(a)=y$，所以连续映射 $\psi$ 所给出的
$\operatorname{Spec}(\mathscr{O}_y)$ 的像包含于 $S_y$。又因为 $S_y$ 包含于
每个包含 $y$ 的仿射开集，所以可以归结到 $Y$ 是仿射概形的情形；在这种情形，
命题由（\hyperref[I.1.6.2-fr]{1.6.2}）得出。

\emph{因此，由（\hyperref[I.2.1.5-fr]{2.1.5}）可见，
$\operatorname{Spec}(\mathscr{O}_y)$ 与 $Y$ 中包含 $y$ 的不可约闭子集全体之间
存在一个双射对应。}

\begin{corollary}[2.4.3]
\label{I.2.4.3-fr}
$y\in Y$ 是 $Y$ 的某个不可约分支的一般点的充要条件是：局部环
$\mathscr{O}_y$ 的唯一素理想是其极大理想（换言之，$\mathscr{O}_y$ 的维数为零）。
\end{corollary}

\begin{proposition}[2.4.4]
\label{I.2.4.4-fr}
设 $(X,\mathscr{O}_X)$ 是环 $A$ 所对应的局部概形，$a$ 是其唯一闭点，
$(Y,\mathscr{O}_Y)$ 是一个预概形。每个态射
$u=(\psi,\theta):(X,\mathscr{O}_X)\to(Y,\mathscr{O}_Y)$ 都唯一地分解为
$X\to\operatorname{Spec}(\mathscr{O}_{\psi(a)})\to Y$，其中第二个箭头表示
典范态射，第一个箭头对应于一个局部同态
$\mathscr{O}_{\psi(a)}\to A$。这建立了态射全体
$(X,\mathscr{O}_X)\to(Y,\mathscr{O}_Y)$ 与局部同态全体
$\mathscr{O}_y\to A$、$(y\in Y)$ 之间的一个典范双射对应。
\end{proposition}

事实上，对每个 $x\in X$，有 $a\in\overline{\{x\}}$，从而
$\psi(a)\in\overline{\{\psi(x)\}}$；这证明 $\psi(X)$ 包含于每个包含
$\psi(a)$ 的仿射开集。因此可以归结到 $(Y,\mathscr{O}_Y)$ 是环 $B$ 所对应的
仿射概形的情形，此时
$u=({}^a\varphi,\widetilde{\varphi})$，其中
$\varphi\in\operatorname{Hom}(B,A)$
（\hyperref[I.1.7.3-fr]{1.7.3}）。此外，
$\varphi^{-1}(\mathfrak{j}_a)=\mathfrak{j}_{\psi(a)}$，因而
$B-\mathfrak{j}_{\psi(a)}$ 的每个元素在 $\varphi$ 下的像都在局部环 $A$ 中
可逆；所以命题中的分解由分式环的泛性质
（\hyperref[0.1.2.4-fr]{0, 1.2.4}）得出。反过来，每个局部同态
$\mathscr{O}_y\to A$ 都对应于唯一的态射
$(\psi,\theta):X\to\operatorname{Spec}(\mathscr{O}_y)$，使得
$\psi(a)=y$（\hyperref[I.1.7.3-fr]{1.7.3}）；将其与典范态射
$\operatorname{Spec}(\mathscr{O}_y)\to Y$ 复合，得到态射 $X\to Y$，
从而完成命题的证明。

\begin{env}[2.4.5]
\label{I.2.4.5-fr}
环是一个\emph{域} $K$ 的仿射概形，其底空间缩为一个点。若 $A$ 是一个局部环，
其极大理想为 $\mathfrak{m}$，则每个局部同态 $A\to K$ 的核都等于
$\mathfrak{m}$，因而分解为 $A\to A/\mathfrak{m}\to K$，其中第二个箭头是
单同态。所以态射 $\operatorname{Spec}(K)\to\operatorname{Spec}(A)$ 与域的
单同态 $A/\mathfrak{m}\to K$ 之间存在双射对应。
\end{env}

设 $(Y,\mathscr{O}_Y)$ 是一个预概形；对每个 $y\in Y$ 及
$\mathscr{O}_y$ 的每个理想 $\mathfrak{a}_y$，典范同态
$\mathscr{O}_y\to\mathscr{O}_y/\mathfrak{a}_y$ 定义一个态射
$\operatorname{Spec}(\mathscr{O}_y/\mathfrak{a}_y)
\to\operatorname{Spec}(\mathscr{O}_y)$；将其与典范态射
$\operatorname{Spec}(\mathscr{O}_y)\to Y$ 复合，就得到仍称为\emph{典范的}
态射 $\operatorname{Spec}(\mathscr{O}_y/\mathfrak{a}_y)\to Y$。当
$\mathfrak{a}_y=\mathfrak{m}_y$，即
$\mathscr{O}_y/\mathfrak{a}_y=k(y)$ 时，命题
（\hyperref[I.2.4.4-fr]{2.4.4}）因而推出：

\begin{corollary}[2.4.6]
\label{I.2.4.6-fr}
\oldpage[I]{103}
设 $(X,\mathscr{O}_X)$ 是环为域 $K$ 的局部概形，$\xi$ 是 $X$ 的唯一点，
$(Y,\mathscr{O}_Y)$ 是一个预概形。每个态射
$u:(X,\mathscr{O}_X)\to(Y,\mathscr{O}_Y)$ 都唯一地分解为
$X\to\operatorname{Spec}(k(\psi(\xi)))\to Y$，其中第二个箭头表示典范态射，
第一个箭头对应于一个单同态 $k(\psi(\xi))\to K$。这建立了态射全体
$(X,\mathscr{O}_X)\to(Y,\mathscr{O}_Y)$ 与单同态全体
$k(y)\to K$、$(y\in Y)$ 之间的一个典范双射对应。
\end{corollary}

\begin{corollary}[2.4.7]
\label{I.2.4.7-fr}
对每个 $y\in Y$，每个典范态射
$\operatorname{Spec}(\mathscr{O}_y/\mathfrak{a}_y)\to Y$ 都是赋环空间的一个
单态射。
\end{corollary}

当 $\mathfrak{a}_y=0$ 时已经证明了这一点
（\hyperref[I.2.4.2-fr]{2.4.2}），只需应用
（\hyperref[I.1.7.5-fr]{1.7.5}）。

\begin{remark}[2.4.8]
\label{I.2.4.8-fr}
设 $X$ 是一个局部概形，$a$ 是其唯一闭点。因为每个包含 $a$ 的仿射开集必然是
整个 $X$，所以每个\emph{可逆} $\mathscr{O}_X$-模
（\hyperref[0.5.4.1-fr]{0, 5.4.1}）都必然\emph{同构于
$\mathscr{O}_X$}（也就是说，它是\emph{平凡的}）。对于任意仿射概形
$\operatorname{Spec}(A)$，这一性质一般并不成立；我们将在第 V 章看到，若
$A$ 是正规环，则当 $A$ 是\emph{唯一分解整环}时，这一性质成立。
\end{remark}

\subsection{预概形之上的预概形}
\label{subsection:I.2.5-fr}

\begin{definition}[2.5.1]
\label{I.2.5.1-fr}
给定一个预概形 $S$，若给定一个预概形 $X$ 及一个预概形态射
$\varphi:X\to S$，就称它们定义了一个\emph{预概形 $S$ 之上的预概形} $X$，
或一个\emph{$S$-预概形}；称 $S$ 为 $S$-预概形 $X$ 的\emph{基预概形}。
态射 $\varphi$ 称为 $S$-预概形 $X$ 的\emph{结构态射}。当 $S$ 是环 $A$
所对应的仿射概形时，也称带有 $\varphi$ 的 $X$ 是\emph{环 $A$ 之上的预概形}
（或一个\emph{$A$-预概形}）。
\end{definition}

由（\hyperref[I.2.2.4-fr]{2.2.4}）可知，给定环 $A$ 之上的一个预概形，等价于
给定一个预概形 $(X,\mathscr{O}_X)$，其结构层 $\mathscr{O}_X$ 是
$A$-\emph{代数层}。\emph{因此，任意预概形总能以唯一的方式视为一个
$\mathbf{Z}$-预概形。}

若 $\varphi:X\to S$ 是 $S$-预概形 $X$ 的结构态射，并且
$\varphi(x)=s$，就称点 $x\in X$\emph{位于点 $s\in S$ 之上}。若
$\varphi$ 是支配态射，就称 $X$\emph{支配} $S$
（\hyperref[I.2.2.6-fr]{2.2.6}）。

\begin{env}[2.5.2]
\label{I.2.5.2-fr}
设 $X$、$Y$ 是两个 $S$-预概形；若预概形态射 $u:X\to Y$ 使图表
\phantomsection\label{I.2.5.2.diagram-fr}
\[
  \xymatrix{
    X\ar[rr]^u\ar[dr] & & Y\ar[dl]\\
    & S &
  }
\]
交换（其中斜箭头是结构态射），就称该态射是一个\emph{$S$ 之上的预概形态射}
（或一个 $S$-\emph{态射}）。这意味着，对每个 $s\in S$ 以及每个
位于 $s$ 之上的 $x\in X$，$u(x)$ 也必须位于 $s$ 之上。
\end{env}

由此定义立即可知，两个 $S$-态射的复合是一个 $S$-态射；因此
$S$-预概形组成一个\emph{范畴}。

用 $\operatorname{Hom}_S(X,Y)$ 表示从 $S$-预概形 $X$ 到 $S$-预概形 $Y$
的所有 $S$-态射所成的集合；$S$-预概形 $X$ 的恒等态射记作 $1_X$。

当 $S$ 是环 $A$ 所对应的仿射概形时，也用 $A$-\emph{态射}代替
$S$-态射这一说法。

\begin{env}[2.5.3]
\label{I.2.5.3-fr}
\oldpage[I]{104}
若 $X$ 是一个 $S$-预概形，而 $v:X'\to X$ 是一个预概形态射，则复合态射
$X'\xrightarrow{v}X\to S$ 把 $X'$ 定义为一个 $S$-预概形；特别地，
$X$ 的开集 $U$ 上的每个诱导预概形，都可以借助典范嵌入视为一个
$S$-预概形。

若 $u:X\to Y$ 是 $S$-预概形之间的一个 $S$-态射，则 $u$ 在 $X$ 的开集
$U$ 上所诱导的预概形上的限制，是一个 $S$-态射 $U\to Y$。反过来，设
$(U_\alpha)$ 是 $X$ 的一个开覆盖，并且对每个 $\alpha$，
$u_\alpha:U_\alpha\to Y$ 是一个 $S$-态射；若对每一对指标
$(\alpha,\beta)$，$u_\alpha$ 与 $u_\beta$ 在 $U_\alpha\cap U_\beta$ 上的
限制相同，则存在唯一的 $S$-态射 $X\to Y$，它在每个 $U_\alpha$ 上的
限制都是 $u_\alpha$。

若 $u:X\to Y$ 是一个满足 $u(X)\subset V$ 的 $S$-态射，其中 $V$ 是 $Y$
的一个开子集，则把 $u$ 看作从 $X$ 到 $V$ 的态射时，它仍是一个
$S$-态射。
\end{env}

\begin{env}[2.5.4]
\label{I.2.5.4-fr}
设 $S'\to S$ 是一个预概形态射；对每个 $S'$-预概形，复合态射
$X\to S'\to S$ 把 $X$ 定义为一个 $S$-预概形。反过来，假设 $S'$ 是 $S$
的一个开集上所诱导的预概形；设 $X$ 是一个 $S$-预概形，并假设其结构态射
$f:X\to S$ 满足 $f(X)\subset S'$；则可以把 $X$ 看作一个
$S'$-预概形。在后一种情形，若 $Y$ 是一个 $S$-预概形，并且其结构态射也把
底空间映入 $S'$，则从 $X$ 到 $Y$ 的每个 $S$-态射也都是一个
$S'$-态射。
\end{env}

\begin{env}[2.5.5]
\label{I.2.5.5-fr}
% EGA1-ERR-0026: print/source says S-morphism here; retained literally pending canon disposition.
若 $X$ 是一个 $S$-态射，$\varphi:X\to S$ 是结构态射，则从 $S$ 到 $X$
的一个 $S$-态射称为 $X$ 的一个 $S$-\emph{截面}；也就是说，它是一个
满足 $\varphi\circ\psi$ 为 $S$ 的恒等态射的预概形态射 $\psi:S\to X$。
用 $\Gamma(X/S)$ 表示 $X$ 的所有 $S$-截面所成的集合。
\end{env}
% EGA I mainland Simplified-Chinese translation (producer output; uncertified).
% Sequential translation cursor: through 3.7.3 (source lines 1--1282; printed pages 104--119).

\section{预概形的积}
\label{section:I.3-fr}

\subsection{预概形的不交并}
\label{subsection:I.3.1-fr}

\phantomsection\label{I.3.1.1-fr}
设 $(X_\alpha)$ 是任意一族预概形；令 $X$ 为各底空间 $X_\alpha$ 的拓扑
\emph{不交并}。于是，$X$ 是两两不交的开子空间 $X'_\alpha$ 的并，并且对
每个 $\alpha$，都有一个从 $X_\alpha$ 到 $X'_\alpha$ 的同胚
$\varphi_\alpha$。在每个 $X'_\alpha$ 上赋予层
$(\varphi_\alpha)_*(\mathscr{O}_{X_\alpha})$ 以后，显然 $X$ 成为一个预概形；
它称为预概形族 $(X_\alpha)$ 的\emph{不交并}，记作
$\bigsqcup_\alpha X_\alpha$。若 $Y$ 是一个预概形，则映射
$f\mapsto(f\circ\varphi_\alpha)$ 给出从集合 $\operatorname{Hom}(X,Y)$
到积集合 $\prod_\alpha\operatorname{Hom}(X_\alpha,Y)$ 的一个\emph{双射}。
特别地，若各 $X_\alpha$ 都是结构态射为 $\psi_\alpha$ 的 $S$-预概形，
则 $X$ 通过唯一满足 $\psi\circ\varphi_\alpha=\psi_\alpha$（对每个
$\alpha$）的态射 $\psi:X\to S$ 成为一个 $S$-预概形。两个预概形 $X$、
$Y$ 的不交并记作 $X\sqcup Y$。显然，若 $X=\operatorname{Spec}(A)$、
$Y=\operatorname{Spec}(B)$，则 $X\sqcup Y$ 典范地等同于
$\operatorname{Spec}(A\times B)$。

\subsection{预概形的积}
\label{subsection:I.3.2-fr}

\begin{definition}[3.2.1]
\label{I.3.2.1-fr}
给定两个 $S$-预概形 $X$、$Y$，设三元组 $(Z,p_1,p_2)$ 由一个
$S$-预概形 $Z$ 以及两个 $S$-态射 $p_1:Z\to X$、$p_2:Z\to Y$ 组成。
若对每个 $S$-预概形 $T$，映射
\oldpage[I]{105}
$f\mapsto(p_1\circ f,p_2\circ f)$ 都是从 $T$ 到 $Z$ 的 $S$-态射全体到
由一个 $S$-态射 $T\to X$ 与一个 $S$-态射 $T\to Y$ 组成的有序对全体的
一个双射（换言之，是一个双射
\[
  \operatorname{Hom}_S(T,Z)\xrightarrow{\sim}
  \operatorname{Hom}_S(T,X)\times\operatorname{Hom}_S(T,Y)).
\]
就称该三元组是 $S$-预概形 $X$ 与 $Y$ 的一个积。
\end{definition}

这正是范畴中两个对象的\emph{积}这一一般概念应用于 $S$-预概形范畴的结果
（T, I, 1.1）；特别地，两个 $S$-预概形的积在唯一的 $S$-同构意义下
\emph{唯一}。由于这种唯一性，通常用 $X\times_S Y$ 表示 $S$-预概形 $X$、
$Y$ 的一个积（若不会引起混淆，则简记为 $X\times Y$），并在记号中省略态射
$p_1$、$p_2$（它们分别称为从 $X\times_S Y$ 到 $X$、$Y$ 的
\emph{典范投影}）。若 $g:T\to X$、$h:T\to Y$ 是两个 $S$-态射，则用
$(g,h)_S$，或简记为 $(g,h)$，表示满足 $p_1\circ f=g$、$p_2\circ f=h$ 的
$S$-态射 $f:T\to X\times_S Y$。若 $X'$、$Y'$ 是两个 $S$-预概形，
$p'_1$、$p'_2$ 是 $X'\times_S Y'$（假设它存在）的典范投影，而
$u:X'\to X$、$v:Y'\to Y$ 是两个 $S$-态射，则用 $u\times_S v$
（或简记为 $u\times v$）表示从 $X'\times_S Y'$ 到 $X\times_S Y$ 的
$S$-态射 $(u\circ p'_1,v\circ p'_2)_S$。

当 $S$ 是环 $A$ 所对应的仿射概形时，在上述记号中常以 $A$ 代替 $S$。

\begin{proposition}[3.2.2]
\label{I.3.2.2-fr}
设 $X$、$Y$、$S$ 是三个仿射概形，其环分别为 $B$、$C$、$A$。令
$Z=\operatorname{Spec}(B\otimes_A C)$，并令 $p_1$、$p_2$ 是分别对应于
典范 $A$-同态 $u:b\mapsto b\otimes1$ 与 $v:c\mapsto1\otimes c$ 的
$S$-态射（\hyperref[I.2.2.4-fr]{2.2.4}）；这两个同态从 $B$、$C$ 映入
$B\otimes_A C$。则 $(Z,p_1,p_2)$ 是 $X$ 与 $Y$ 的一个积。
\end{proposition}

根据（\hyperref[I.2.2.4-fr]{2.2.4}），只需验证：把每个 $A$-同态
$f:B\otimes_A C\to L$（其中 $L$ 是一个 $A$-代数）映为有序对
$(f\circ u,f\circ v)$，就定义了一个双射
\[
  \operatorname{Hom}_A(B\otimes_A C,L)\xrightarrow{\sim}
  \operatorname{Hom}_A(B,L)\times\operatorname{Hom}_A(C,L)\footnotemark,
\]
\footnotetext{这里，记号 $\operatorname{Hom}_A$ 表示 $A$-代数同态全体。}
这由定义以及关系 $b\otimes c=(b\otimes1)(1\otimes c)$ 立即得出。

\begin{corollary}[3.2.3]
\label{I.3.2.3-fr}
设 $T$ 是环 $D$ 所对应的仿射概形，
$\alpha=({}^a\rho,\widetilde{\rho})$
（相应地，$\beta=({}^a\sigma,\widetilde{\sigma})$）是一个 $S$-态射
$T\to X$（相应地，$T\to Y$），其中 $\rho$（相应地，$\sigma$）是从
$B$（相应地，$C$）到 $D$ 的一个 $A$-同态。则
$(\alpha,\beta)_S=({}^a\tau,\widetilde{\tau})$，其中 $\tau$ 是满足
$\tau(b\otimes c)=\rho(b)\sigma(c)$ 的同态 $B\otimes_A C\to D$。
\end{corollary}

\begin{proposition}[3.2.4]
\label{I.3.2.4-fr}
设 $f:S'\to S$ 是预概形的一个单态射（T, I, 1.1），$X$、$Y$ 是两个
$S'$-预概形，并借助 $f$ 也把它们看作 $S$-预概形。于是，$S$-预概形
$X$、$Y$ 的每个积也是 $S'$-预概形 $X$、$Y$ 的一个积，反之亦然。
\end{proposition}

设 $\varphi:X\to S'$、$\psi:Y\to S'$ 是结构态射。若 $T$ 是一个
$S$-预概形，而 $u:T\to X$、$v:T\to Y$ 是两个 $S$-态射，则按定义有
$f\circ\varphi\circ u=f\circ\psi\circ v=\theta$，其中右端是 $T$ 的
结构态射；由关于 $f$ 的假设可得
$\varphi\circ u=\psi\circ v=\theta'$。可见，可以把 $T$ 看作结构态射为
$\theta'$ 的 $S'$-预概形，并把 $u$ 与 $v$ 看作 $S'$-态射。结合
（\hyperref[I.3.2.1-fr]{3.2.1}），命题的结论立即得出。

\begin{corollary}[3.2.5]
\label{I.3.2.5-fr}
设 $X$、$Y$ 是两个 $S$-预概形，$\varphi:X\to S$、$\psi:Y\to S$ 是它们的
结构态射，$S'$ 是 $S$ 的一个满足 $\varphi(X)\subset S'$、
$\psi(Y)\subset S'$ 的开子集。则 $S$-预概形 $X$、$Y$ 的每个积也都是
$S'$-预概形 $X$、$Y$ 的一个积，反之亦然。
\end{corollary}

\oldpage[I]{106}
只需把（\hyperref[I.3.2.4-fr]{3.2.4}）应用于典范嵌入 $S'\to S$。

\begin{theorem}[3.2.6]
\label{I.3.2.6-fr}
给定两个 $S$-预概形 $X$、$Y$，存在一个积 $X\times_S Y$。
\end{theorem}

我们将分几个步骤来证明。

\begin{lemma}[3.2.6.1]
\label{I.3.2.6.1-fr}
设 $(Z,p,q)$ 是 $X$ 与 $Y$ 的一个积，$U$、$V$ 分别是 $X$、$Y$ 的
开子集。令 $W=p^{-1}(U)\cap q^{-1}(V)$，则由 $W$ 以及 $p$、$q$ 在 $W$
上的限制所组成的三元组\textup{（分别把这些限制视为态射 $W\to U$、
$W\to V$）}是 $U$ 与 $V$ 的一个积。
\end{lemma}

事实上，若 $T$ 是一个 $S$-预概形，就可以把 $S$-态射 $T\to W$ 与那些
把 $T$ 映入 $W$ 的 $S$-态射 $T\to Z$ 等同起来。若
$g:T\to U$、$h:T\to V$ 是任意两个 $S$-态射，则可分别把它们看作从
$T$ 到 $X$、$Y$ 的 $S$-态射；因此，依假设，存在唯一的 $S$-态射
$f:T\to Z$，使得 $g=p\circ f$、$h=q\circ f$。由于
$p(f(T))\subset U$、$q(f(T))\subset V$，有
\[
  f(T)\subset p^{-1}(U)\cap q^{-1}(V)=W,
\]
断言得证。

\begin{lemma}[3.2.6.2]
\label{I.3.2.6.2-fr}
设 $Z$ 是一个 $S$-预概形，$p:Z\to X$、$q:Z\to Y$ 是两个 $S$-态射，
$(U_\alpha)$ 是 $X$ 的一个开覆盖，$(V_\lambda)$ 是 $Y$ 的一个开覆盖。
假设对每一对 $(\alpha,\lambda)$，$S$-预概形
$W_{\alpha\lambda}=p^{-1}(U_\alpha)\cap q^{-1}(V_\lambda)$ 与 $p$、$q$ 在
$W_{\alpha\lambda}$ 上的限制构成 $U_\alpha$ 与 $V_\lambda$ 的一个积。
则 $(Z,p,q)$ 是 $X$ 与 $Y$ 的一个积。
\end{lemma}

先证明：若 $f_1$、$f_2$ 是两个 $S$-态射 $T\to Z$，则关系
$p\circ f_1=p\circ f_2$ 与 $q\circ f_1=q\circ f_2$ 蕴含 $f_1=f_2$。
事实上，$Z$ 是各 $W_{\alpha\lambda}$ 的并，所以各
$f_1^{-1}(W_{\alpha\lambda})$ 构成 $T$ 的一个开覆盖，各
$f_2^{-1}(W_{\alpha\lambda})$ 也是如此。此外，依假设有
\[
\begin{aligned}
  f_1^{-1}(W_{\alpha\lambda})
  &=f_1^{-1}(p^{-1}(U_\alpha))\cap
    f_1^{-1}(q^{-1}(V_\lambda)) \\
  &=f_2^{-1}(p^{-1}(U_\alpha))\cap
    f_2^{-1}(q^{-1}(V_\lambda))
   =f_2^{-1}(W_{\alpha\lambda})
\end{aligned}
\]
因而只需证明，对每一对指标，$f_1$ 与 $f_2$ 在
$f_1^{-1}(W_{\alpha\lambda})=f_2^{-1}(W_{\alpha\lambda})$ 上的限制相同。
但这些限制可以看作从 $f_1^{-1}(W_{\alpha\lambda})$ 到
$W_{\alpha\lambda}$ 的 $S$-态射，故断言由假设与定义
（\hyperref[I.3.2.1-fr]{3.2.1}）得出。

现在给定两个 $S$-态射 $g:T\to X$、$h:T\to Y$。令
$T_{\alpha\lambda}=g^{-1}(U_\alpha)\cap h^{-1}(V_\lambda)$；各
$T_{\alpha\lambda}$ 构成 $T$ 的一个开覆盖。依假设，存在一个 $S$-态射
$f_{\alpha\lambda}$，使得 $p\circ f_{\alpha\lambda}$ 与
$q\circ f_{\alpha\lambda}$ 分别是 $g$、$h$ 在 $T_{\alpha\lambda}$ 上的
限制。此外，只需证明 $f_{\alpha\lambda}$ 与 $f_{\beta\mu}$ 在
$T_{\alpha\lambda}\cap T_{\beta\mu}$ 上的限制相同，即可完成
（\hyperref[I.3.2.6.2-fr]{3.2.6.2}）的证明。依定义，
$T_{\alpha\lambda}\cap T_{\beta\mu}$ 在 $f_{\alpha\lambda}$ 与
$f_{\beta\mu}$ 下的像都包含于
$W_{\alpha\lambda}\cap W_{\beta\mu}$。由于
\[
  W_{\alpha\lambda}\cap W_{\beta\mu}
  =p^{-1}(U_\alpha\cap U_\beta)\cap
   q^{-1}(V_\lambda\cap V_\mu),
\]
由（\hyperref[I.3.2.6.1-fr]{3.2.6.1}）可知，
$W_{\alpha\lambda}\cap W_{\beta\mu}$ 连同 $p$、$q$ 在该预概形上的限制，
构成 $U_\alpha\cap U_\beta$ 与 $V_\lambda\cap V_\mu$ 的一个积。又因为
$p\circ f_{\alpha\lambda}$ 与 $p\circ f_{\beta\mu}$ 在
$T_{\alpha\lambda}\cap T_{\beta\mu}$ 上相同，而
$q\circ f_{\alpha\lambda}$ 与 $q\circ f_{\beta\mu}$ 也是如此，所以
$f_{\alpha\lambda}$ 与 $f_{\beta\mu}$ 在
$T_{\alpha\lambda}\cap T_{\beta\mu}$ 上相同，证毕。

\oldpage[I]{107}
\begin{lemma}[3.2.6.3]
\label{I.3.2.6.3-fr}
设 $(U_\alpha)$ 是 $X$ 的一个开覆盖，$(V_\lambda)$ 是 $Y$ 的一个开覆盖，
并假设对每一对 $(\alpha,\lambda)$，$U_\alpha$ 与 $V_\lambda$ 的一个积存在；
则 $X$ 与 $Y$ 的一个积存在。
\end{lemma}

把引理（\hyperref[I.3.2.6.1-fr]{3.2.6.1}）应用于开集
$U_\alpha\cap U_\beta$ 与 $V_\lambda\cap V_\mu$，可知由 $X$、$Y$ 分别在
这些开集上诱导的 $S$-预概形的一个积存在。此外，积的唯一性表明，若令
$i=(\alpha,\lambda)$、$j=(\beta,\mu)$，则有一个典范同构 $h_{ij}$
（相应地，$h_{ji}$），把这个积映到由 $U_\alpha\times_S V_\lambda$
（相应地，$U_\beta\times_S V_\mu$）在某个开集上诱导的一个 $S$-预概形
$W_{ij}$（相应地，$W_{ji}$）上；因此
$f_{ij}=h_{ij}\circ h_{ji}^{-1}$ 是从 $W_{ji}$ 到 $W_{ij}$ 的一个同构。
此外，对于第三个有序对 $k=(\gamma,\nu)$，在
$W_{ki}\cap W_{kj}$ 上有 $f_{ik}=f_{ij}\circ f_{jk}$；这是把
（\hyperref[I.3.2.6.1-fr]{3.2.6.1}）分别应用于 $U_\beta$、$V_\mu$ 中的
开集 $U_\alpha\cap U_\beta\cap U_\gamma$ 与
$V_\lambda\cap V_\mu\cap V_\nu$ 所得。因而存在一个预概形 $Z$、$Z$ 的底空间的
一个开覆盖 $(Z_i)$，以及对每个 $i$ 的一个同构 $g_i$，把 $Z_i$ 上的诱导
预概形映到预概形 $U_\alpha\times_S V_\lambda$ 上，使得对每一对 $(i,j)$，
都有 $f_{ij}=g_i\circ g_j^{-1}$（\hyperref[I.2.3.1-fr]{2.3.1}）；此外，
$g_i(Z_i\cap Z_j)=W_{ij}$。设 $p_i$、$q_i$、$\theta_i$ 是 $S$-预概形
$U_\alpha\times_S V_\lambda$ 的两个投影及结构态射。立即可见，在
$Z_i\cap Z_j$ 上有 $p_i\circ g_i=p_j\circ g_j$，另外两个态射也是如此。
因此，可以定义预概形态射 $p:Z\to X$（相应地，$q:Z\to Y$、
$\theta:Z\to S$），使 $p$（相应地，$q$、$\theta$）在每个 $Z_i$ 上都与
$p_i\circ g_i$（相应地，$q_i\circ g_i$、$\theta_i\circ g_i$）相同；赋予
$Z$ 以 $\theta$ 后，它便是一个 $S$-预概形。现在证明
$Z'_i=p^{-1}(U_\alpha)\cap q^{-1}(V_\lambda)$ 等于 $Z_i$。对于每个指标
$j=(\beta,\mu)$，有
$Z_j\cap Z'_i=g_j^{-1}(p_j^{-1}(U_\alpha)\cap q_j^{-1}(V_\lambda))$。另一方面，
\[
  p_j^{-1}(U_\alpha)\cap q_j^{-1}(V_\lambda)
  =p_j^{-1}(U_\alpha\cap U_\beta)\cap
   q_j^{-1}(V_\lambda\cap V_\mu)~;
\]
根据（\hyperref[I.3.2.6.1-fr]{3.2.6.1}），$p_j$、$q_j$ 在
$p_j^{-1}(U_\alpha)\cap q_j^{-1}(V_\lambda)$ 上的限制，使这个
$S$-预概形成为 $U_\alpha\cap U_\beta$ 与 $V_\lambda\cap V_\mu$ 的一个积；
而积的唯一性遂给出
$p_j^{-1}(U_\alpha)\cap q_j^{-1}(V_\lambda)=W_{ji}$。故对每个 $j$，均有
$Z_j\cap Z'_i=Z_j\cap Z_i$，从而 $Z'_i=Z_i$。最后由
（\hyperref[I.3.2.6.2-fr]{3.2.6.2}）可知，$(Z,p,q)$ 是 $X$ 与 $Y$ 的一个积。

\begin{lemma}[3.2.6.4]
\label{I.3.2.6.4-fr}
设 $\varphi:X\to S$、$\psi:Y\to S$ 是 $X$、$Y$ 的结构态射，$(S_i)$ 是
$S$ 的一个开覆盖，并令 $X_i=\varphi^{-1}(S_i)$、$Y_i=\psi^{-1}(S_i)$。
若每个积 $X_i\times_S Y_i$ 都存在，则 $X\times_S Y$ 存在。
\end{lemma}

根据（\hyperref[I.3.2.6.3-fr]{3.2.6.3}），只需证明积
$X_i\times_S Y_j$ 对任意 $i$、$j$ 都存在。令
$X_{ij}=X_i\cap X_j=\varphi^{-1}(S_i\cap S_j)$、
$Y_{ij}=Y_i\cap Y_j=\psi^{-1}(S_i\cap S_j)$；由
（\hyperref[I.3.2.6.1-fr]{3.2.6.1}），积
$Z_{ij}=X_{ij}\times_S Y_{ij}$ 存在。现在注意，若 $T$ 是一个
$S$-预概形，而 $g:T\to X_i$、$h:T\to Y_j$ 是 $S$-态射，则按 $S$-态射的
定义必有 $\varphi(g(T))=\psi(h(T))\subset S_i\cap S_j$，所以
$g(T)\subset X_{ij}$ 且 $h(T)\subset Y_{ij}$；于是立即可见，$Z_{ij}$ 是
$X_i$ 与 $Y_j$ 的一个积。

\begin{env}[3.2.6.5]
\label{I.3.2.6.5-fr}
现在可以完成定理（\hyperref[I.3.2.6-fr]{3.2.6}）的证明。若 $S$ 是一个
\emph{仿射概形}，则 $X$、$Y$ 分别有由仿射开集组成的覆盖 $(U_\alpha)$、
$(V_\lambda)$；由（\hyperref[I.3.2.2-fr]{3.2.2}），
$U_\alpha\times_S V_\lambda$ 存在，故由
（\hyperref[I.3.2.6.3-fr]{3.2.6.3}），$X\times_S Y$ 也存在。若 $S$ 是任意
预概形，则 $S$ 有一个由仿射开集组成的覆盖 $(S_i)$。若
$\varphi:X\to S$、$\psi:Y\to S$ 是结构态射，并令
$X_i=\varphi^{-1}(S_i)$、$Y_i=\psi^{-1}(S_i)$，则由前述结果，积
$X_i\times_{S_i}Y_i$ 存在；但是积 $X_i\times_S Y_i$ 也存在
（\hyperref[I.3.2.5-fr]{3.2.5}），因此由
\oldpage[I]{108}
（\hyperref[I.3.2.6.4-fr]{3.2.6.4}），$X\times_S Y$ 也存在。
\end{env}

\begin{corollary}[3.2.7]
\label{I.3.2.7-fr}
设 $Z=X\times_S Y$ 是两个 $S$-预概形的积，$p$、$q$ 是从 $Z$ 到 $X$、
$Y$ 的投影，$\varphi$（相应地，$\psi$）是 $X$（相应地，$Y$）的结构态射。
设 $S'$ 是 $S$ 的一个开子集，$U$（相应地，$V$）是 $X$（相应地，$Y$）
的一个包含于 $\varphi^{-1}(S')$（相应地，$\psi^{-1}(S')$）的开子集。
则积 $U\times_{S'}V$ 典范地等同于由 $Z$ 在
$p^{-1}(U)\cap q^{-1}(V)$ 上诱导的预概形（把它看作 $S'$-预概形）。此外，
若 $f:T\to X$、$g:T\to Y$ 是满足 $f(T)\subset U$、$g(T)\subset V$ 的
$S$-态射，则 $S'$-态射 $(f,g)_{S'}$ 等同于 $(f,g)_S$ 在
$p^{-1}(U)\cap q^{-1}(V)$ 上的限制。
\end{corollary}

这由（\hyperref[I.3.2.5-fr]{3.2.5}）和
（\hyperref[I.3.2.6.1-fr]{3.2.6.1}）得出。

\begin{env}[3.2.8]
\label{I.3.2.8-fr}
设 $(X_\alpha)$、$(Y_\lambda)$ 是两个 $S$-预概形族，$X$（相应地，$Y$）
是族 $(X_\alpha)$（相应地，$(Y_\lambda)$）的不交并
（\hyperref[subsection:I.3.1-fr]{3.1}）。则 $X\times_S Y$ 等同于族
$(X_\alpha\times_S Y_\lambda)$ 的\emph{不交并}；这立即由
（\hyperref[I.3.2.6.3-fr]{3.2.6.3}）得出。
\end{env}

\subsection{积的形式性质；预概形的换基}
\label{subsection:I.3.3-fr}

\begin{env}[3.3.1]
\label{I.3.3.1-fr}
读者会注意到，本节所述的全部性质，除
（\hyperref[I.3.3.13-fr]{3.3.13}）和
（\hyperref[I.3.3.15-fr]{3.3.15}）外，在\emph{任意范畴}中均可原样成立，
只要陈述中所涉及的积\emph{存在}（因为显然，对于范畴的任意对象 $S$，
都可以完全仿照（\hyperref[subsection:I.2.5-fr]{2.5}）定义 $S$-对象与
$S$-态射的概念）。
\end{env}

\begin{env}[3.3.2]
\label{I.3.3.2-fr}
首先，在 $S$-预概形范畴中，$X\times_S Y$ 关于 $X$ 和 $Y$ 是一个
\emph{协变双函子}：事实上，只需注意图表
\[
  \xymatrix{
    X\times Y\ar[r]^{f\times 1}\ar[d] &
    X'\times Y\ar[r]^{f'\times 1}\ar[d] &
    X''\times Y\ar[d]\\
    X\ar[r]_f & X'\ar[r]_{f'} & X''
  }
\]
是交换的。
\end{env}

\begin{proposition}[3.3.3]
\label{I.3.3.3-fr}
对于任意 $S$-预概形 $X$，$X\times_S S$（相应地，$S\times_S X$）的第一
（相应地，第二）投影，是从 $X\times_S S$（相应地，$S\times_S X$）
到 $X$ 的函子性同构；其逆同构为 $(1_X,\varphi)_S$（相应地，
$(\varphi,1_X)_S$），其中 $\varphi$ 表示结构态射 $X\to S$。因此，在相差
一个典范同构的意义下，可以写成
\[
  X\times_S S=S\times_S X=X.
\]
\end{proposition}

只需证明三元组 $(X,1_X,\varphi)$ 是 $X$ 与 $S$ 的一个积。事实上，若
$T$ 是一个 $S$-预概形，则从 $T$ 到 $S$ 的唯一 $S$-态射必然是结构态射
$\psi:T\to S$。若 $f$ 是从 $T$ 到 $X$ 的一个 $S$-态射，则必有
$\psi=\varphi\circ f$，由此即得所断言者。

\begin{corollary}[3.3.4]
\label{I.3.3.4-fr}
设 $X$、$Y$ 是两个 $S$-预概形，$\varphi:X\to S$、$\psi:Y\to S$ 是它们的
结构态射。若典范地把 $X$ 与 $X\times_S S$ 等同，并把 $Y$ 与
$S\times_S Y$ 等同，则投影 $X\times_S Y\to X$ 与
$X\times_S Y\to Y$ 分别等同于 $1_X\times\psi$ 与
$\varphi\times 1_Y$。
\end{corollary}

验证是立即的，留给读者。

\begin{env}[3.3.5]
\label{I.3.3.5-fr}
可以仿照（\hyperref[subsection:I.3.2-fr]{3.2}）定义任意有限个
\oldpage[I]{109}
$S$-预概形（设其数目为 $n$）的积；由（\hyperref[I.3.2.6-fr]{3.2.6}），
对 $n$ 作归纳，
并注意到
$(X_1\times_S X_2\times_S\cdots\times_S X_{n-1})\times_S X_n$ 满足积的
定义，即得这些积的存在性。如同在任意范畴中一样，积的唯一性蕴含它的
\emph{交换性}和\emph{结合性}。例如，若 $p_1$、$p_2$、$p_3$ 表示
$X_1\times_S X_2\times_S X_3$ 的各投影，并且把这个预概形等同于
$(X_1\times_S X_2)\times_S X_3$，则到 $X_1\times_S X_2$ 的投影等同于
$(p_1,p_2)_S$。
\end{env}

\begin{env}[3.3.6]
\label{I.3.3.6-fr}
设 $S$、$S'$ 是两个预概形，$\varphi:S'\to S$ 是一个态射，它使 $S'$ 成为
一个 $S$-预概形。对于任意 $S$-预概形 $X$，考虑积 $X\times_S S'$，并令
$p$、$\pi'$ 分别为它到 $X$、$S'$ 的投影。赋予 $\pi'$ 后，这个积是一个
$S'$-预概形；把它看作这样的预概形时，记作 $X_{(S')}$ 或
$X_{(\varphi)}$，并称它为经由态射 $\varphi$ 从 $S$ 到 $S'$ \emph{换基}
所得的预概形，或称为 $X$ 在 $\varphi$ 下的\emph{逆像}。注意，若 $\pi$
是 $X$ 的结构态射，而 $\theta$ 是把 $X\times_S S'$ 看作 $S$-预概形时的
结构态射，则图表
\phantomsection\label{I.3.3.6.base-change-diagram-fr}
\[
  \xymatrix{
    X\ar[d]_{\pi} &
    X_{(S')}\ar[l]_{p}\ar[ld]_{\theta}\ar[d]^{\pi'}\\
    S & S'\ar[l]^{\varphi}
  }
\]
是交换的。
\end{env}

\begin{env}[3.3.7]
\label{I.3.3.7-fr}
沿用（\hyperref[I.3.3.6-fr]{3.3.6}）的记号，对于任意 $S$-态射
$f:X\to Y$，仍以 $f_{(S')}$ 表示 $S'$-态射
$f\times_S 1:X_{(S')}\to Y_{(S')}$，并称 $f_{(S')}$ 为态射 $f$ 在
$\varphi$ 下的\emph{逆像}。因此，$X_{(S')}$ 关于 $X$ 是一个从
$S$-预概形范畴到 $S'$-预概形范畴的\emph{协变函子}。
\end{env}

\begin{env}[3.3.8]
\label{I.3.3.8-fr}
预概形 $X_{(S')}$ 还可以看作一个\emph{泛映射问题}的解：任意
$S'$-预概形 $T$ 经由 $\varphi$ 也成为一个 $S$-预概形；任意 $S$-态射
$g:T\to X$ 都能唯一地写成 $g=p\circ f$，其中 $f$ 是一个 $S'$-态射
$T\to X_{(S')}$；这是由把积的定义应用于
% EGA1-ERR-0027: French print/source has f; the defining pair is (g,\psi); canonical Chinese retains f.
$S$-态射 $f$ 和 $\psi:T\to S'$（$T$ 的结构态射）所得。
\end{env}

\begin{proposition}[3.3.9]
\label{I.3.3.9-fr}
\textup{（“换基的传递性”。）}设 $S''$ 是一个预概形，
$\varphi':S''\to S'$ 是一个态射。对于任意 $S$-预概形 $X$，从
$S''$-预概形 $(X_{(\varphi)})_{(\varphi')}$ 到 $S''$-预概形
$X_{(\varphi\circ\varphi')}$，存在一个函子性的典范同构。
\end{proposition}

事实上，设 $T$ 是一个 $S''$-预概形，$\psi$ 是它的结构态射，$g$ 是从
$T$ 到 $X$ 的一个 $S$-态射（这里把 $T$ 看作结构态射为
$\varphi\circ\varphi'\circ\psi$ 的 $S$-预概形）。由于 $T$ 也是结构态射
为 $\varphi'\circ\psi$ 的一个 $S'$-预概形，可以写成 $g=p\circ g'$，
其中 $g'$ 是一个 $S'$-态射 $T\to X_{(\varphi)}$；进而写成
$g'=p'\circ g''$，其中 $g''$ 是一个 $S''$-态射
$T\to(X_{(\varphi)})_{(\varphi')}$：
\phantomsection\label{I.3.3.9.transitivity-diagram-fr}
\[
  \xymatrix{
    X\ar[d]_{\pi} &
    X_{(\varphi)}\ar[l]_{p}\ar[d]_{\pi'} &
    (X_{(\varphi)})_{(\varphi')}\ar[l]_{p'}\ar[d]_{\pi''}\\
    S & S'\ar[l]^{\varphi} & S''\ar[l]^{\varphi'}
  }
\]
\oldpage[I]{110}
由于泛映射问题的解具有唯一性，由此即得命题。

若不存在混淆，这个结果还可以写成（在相差一个典范同构的意义下）的等式
$(X_{(S')})_{(S'')}=X_{(S'')}$，或者写成
\begin{equation}
  (X\times_S S')\times_{S'}S''=X\times_S S''~;
  \tag{3.3.9.1}\label{I.3.3.9.1-fr}
\end{equation}
（\hyperref[I.3.3.9-fr]{3.3.9}）中所定义同构的函子性，也可用下列
\emph{态射逆像的传递性}公式表示：
\begin{equation}
  (f_{(S')})_{(S'')}=f_{(S'')}
  \tag{3.3.9.2}\label{I.3.3.9.2-fr}
\end{equation}
其中 $f:X\to Y$ 是任意 $S$-态射。

\begin{corollary}[3.3.10]
\label{I.3.3.10-fr}
若 $X$、$Y$ 是两个 $S$-预概形，则存在一个函子性的典范同构，把
$S'$-预概形 $X_{(S')}\times_{S'}Y_{(S')}$ 映到 $S'$-预概形
$(X\times_S Y)_{(S')}$ 上。
\end{corollary}

事实上，在相差典范同构的意义下，有
\[
  (X\times_S S')\times_{S'}(Y\times_S S')
  =X\times_S(Y\times_S S')=(X\times_S Y)\times_S S'
\]
这是由（\hyperref[I.3.3.9.1-fr]{3.3.9.1}）以及 $S$-预概形的积的结合性得出。

（\hyperref[I.3.3.10-fr]{3.3.10}）中所定义同构的函子性由公式
\begin{equation}
  (u_{(S')},v_{(S')})_{S'}=((u,v)_S)_{(S')}
  \tag{3.3.10.1}\label{I.3.3.10.1-fr}
\end{equation}
表示，其中 $u:T\to X$、$v:T\to Y$ 是任意一对 $S$-态射。

换言之，逆像函子 $X_{(S')}$ \emph{与取积可交换}；它也与取不交并可交换
（\hyperref[I.3.2.8-fr]{3.2.8}）。

\begin{corollary}[3.3.11]
\label{I.3.3.11-fr}
设 $Y$ 是一个 $S$-预概形，$f:X\to Y$ 是一个使 $X$ 成为 $Y$-预概形
（从而也成为 $S$-预概形）的态射。于是，预概形 $X_{(S')}$ 等同于积
$X\times_Y Y_{(S')}$，而投影 $X\times_Y Y_{(S')}\to Y_{(S')}$ 等同于
$f_{(S')}$。
\end{corollary}

设 $\psi:Y\to S$ 是 $Y$ 的结构态射；有交换图表
\phantomsection\label{I.3.3.11.diagram-fr}
\[
  \xymatrix{
    S'\ar[d] &
    Y_{(S')}\ar[l]_{\psi_{(S')}}\ar[d] &
    X_{(S')}\ar[l]_{f_{(S')}}\ar[d]\\
    S & Y\ar[l]^{\psi} & X\ar[l]^{f}
  }
\]
而 $Y_{(S')}$ 等同于 $S'_{(\psi)}$，$X_{(S')}$ 等同于
$S'_{(\psi\circ f)}$；考虑到（\hyperref[I.3.3.9-fr]{3.3.9}）和
（\hyperref[I.3.3.4-fr]{3.3.4}），即得推论。

\begin{env}[3.3.12]
\label{I.3.3.12-fr}
设 $f:X\to X'$、$g:Y\to Y'$ 是两个 $S$-态射，并且都是预概形的
\emph{单态射}（T, I, 1.1）；则 $f\times_S g$ 是一个\emph{单态射}。
事实上，若 $p$、$q$ 是 $X\times_S Y$ 的投影，$p'$、$q'$ 是
$X'\times_S Y'$ 的投影，而 $u$、$v$ 是两个 $S$-态射
$T\to X\times_S Y$，则关系
$(f\times_S g)\circ u=(f\times_S g)\circ v$ 蕴含
$p'\circ(f\times_S g)\circ u=p'\circ(f\times_S g)\circ v$，换言之，
$f\circ p\circ u=f\circ p\circ v$；由于 $f$ 是单态射，便有
$p\circ u=p\circ v$。利用 $g$ 是单态射这一事实，同理得到
$q\circ u=q\circ v$，因而 $u=v$。
\oldpage[I]{111}
由此可知，对于底预概形的任意换基 $S'\to S$，
\phantomsection\label{I.3.3.12.base-change-monomorphism-fr}
\[
% EGA1-ERR-0028: French print/source has Y_(S'); the base-changed codomain of f:X->X' is X'_(S'). Canonical Chinese retains Y_(S').
  f_{(S')}:X_{(S')}\to Y_{(S')}
\]
是一个单态射。
\end{env}

\begin{env}[3.3.13]
\label{I.3.3.13-fr}
设 $S$、$S'$ 是环分别为 $A$、$A'$ 的两个仿射概形；于是，一个态射
$S'\to S$ 对应于一个环同态 $A\to A'$。若 $X$ 是一个 $S$-预概形，
则也把 $S'$-预概形 $X_{(S')}$ 记作 $X_{(A')}$ 或 $X\otimes_A A'$；
当 $X$ 本身是环为 $B$ 的仿射概形时，$X_{(A')}$ 是环为
$B_{(A')}=B\otimes_A A'$ 的仿射概形，这个环是把 $A$-代数 $B$ 的
标量环扩张到 $A'$ 所得。
\end{env}

\begin{env}[3.3.14]
\label{I.3.3.14-fr}
沿用（\hyperref[I.3.3.6-fr]{3.3.6}）的记号，对于任意
\emph{$S$-态射} $f:S'\to X$，$f'=(f,1_{S'})_S$ 是一个 $S'$-态射
$S'\to X'=X_{(S')}$，满足 $p\circ f'=f$、$\pi'\circ f'=1_{S'}$；换言之，
它是 $X'$ 的一个\emph{$S'$-截面}。反过来，若 $f'$ 是这样的
$S'$-截面，则 $f=p\circ f'$ 是一个 $S$-态射 $S'\to X$。这样就定义了
一个典范的\emph{双射对应}
\phantomsection\label{I.3.3.14.correspondence-fr}
\[
  \operatorname{Hom}_S(S',X)\overset{\sim}{\longrightarrow}
  \operatorname{Hom}_{S'}(S',X').
\]
称 $f'$ 为 $f$ 的\emph{图态射}，并把它记作 $\Gamma_f$。
\end{env}

\begin{env}[3.3.15]
\label{I.3.3.15-fr}
给定一个预概形 $X$，总可以把它看作一个 $\mathbf{Z}$-预概形；特别地，
由（\hyperref[I.3.3.14-fr]{3.3.14}）可知，
$X\otimes_{\mathbf{Z}}\mathbf{Z}[T]$ 的各 $X$-截面（其中 $T$ 是一个不定元）
与各\emph{态射}
% EGA1-ERR-0029: French print/source reverses the required scheme-morphism direction; canonical Chinese retains Z[T]->X.
$\mathbf{Z}[T]\to X$ 一一对应。下面证明，这些 $X$-截面也与
\emph{结构层 $\mathcal{O}_X$ 在 $X$ 上的截面}一一对应。事实上，设
$(U_\alpha)$ 是 $X$ 的一个仿射开覆盖；设
$u:X\to X\otimes_{\mathbf{Z}}\mathbf{Z}[T]$ 是一个 $X$-态射，并令
$u_\alpha$ 为它在 $U_\alpha$ 上的限制。若 $A_\alpha$ 是仿射概形
$U_\alpha$ 的环，则 $U_\alpha\otimes_{\mathbf{Z}}\mathbf{Z}[T]$ 是环为
$A_\alpha[T]$ 的一个仿射概形（\hyperref[I.3.2.2-fr]{3.2.2}），而
$u_\alpha$ 典范地对应于一个 $A_\alpha$-同态
$A_\alpha[T]\to A_\alpha$（1.7.3）。这样的同态完全由 $T$ 在
$A_\alpha$ 中的像决定；把这个像记作
$s_\alpha\in A_\alpha=\Gamma(U_\alpha,\mathcal{O}_X)$。若写出
$u_\alpha$、$u_\beta$ 在仿射开集 $V\subset U_\alpha\cap U_\beta$ 上的
限制相同，就立即看出 $s_\alpha$ 与 $s_\beta$ 在 $V$ 上相同。因此，
族 $(s_\alpha)$ 由 $\mathcal{O}_X$ 在 $X$ 上的一个截面 $s$ 在各
$U_\alpha$ 上的限制组成。反过来，显然这样的截面定义一个态射族
$(u_\alpha)$，它们是某个 $X$-态射
$X\to X\otimes_{\mathbf{Z}}\mathbf{Z}[T]$ 在各 $U_\alpha$ 上的限制。
这个结果将在 II, 1.7.12 中推广。
\end{env}

\subsection{预概形的取值在预概形中的点；几何点}
\label{subsection:I.3.4-fr}

\begin{env}[3.4.1]
\label{I.3.4.1-fr}
设 $X$ 是一个预概形；对于任意预概形 $T$，仍以 $X(T)$ 表示态射
$T\to X$ 的集合 $\operatorname{Hom}(T,X)$，这个集合的元素也称为
\emph{$X$ 的取值于 $T$ 的点}。若对每个态射 $f:T\to T'$，令它对应于从
$X(T')$ 到 $X(T)$ 的映射 $u'\mapsto u'\circ f$，则可见，在 $X$ 固定时，
$X(T)$ 是从预概形范畴到集合范畴的一个\emph{关于 $T$ 的逆变函子}。
此外，每个预概形态射 $g:X\to Y$ 都定义一个函子性同态
$X(T)\to Y(T)$，它把 $v\in X(T)$ 映到 $g\circ v$。
\end{env}

\begin{env}[3.4.2]
\label{I.3.4.2-fr}
给定三个集合 $P$、$Q$、$R$ 以及两个映射
$\varphi:P\to R$、$\psi:Q\to R$，称（相对于 $\varphi$ 和 $\psi$）
\emph{$P$ 与 $Q$ 在 $R$ 上的纤维积}为
\oldpage[I]{112}
积集合 $P\times Q$ 中由所有满足 $\varphi(p)=\psi(q)$ 的偶 $(p,q)$
所组成的子集，并把它记作 $P\times_R Q$。采用
（\hyperref[I.3.4.1-fr]{3.4.1}）的记号，
$S$-预概形之积的定义（\hyperref[I.3.2.1-fr]{3.2.1}）还可以写成公式
\phantomsection\label{I.3.4.2.1-fr}
\[
  (X\times_S Y)(T)=X(T)\times_{S(T)}Y(T)
  \tag{3.4.2.1}
\]
其中映射 $X(T)\to S(T)$、$Y(T)\to S(T)$ 分别对应于结构态射
$X\to S$、$Y\to S$。
\end{env}

\begin{env}[3.4.3]
\label{I.3.4.3-fr}
若给定一个预概形 $S$，并且只考虑 $S$-预概形和 $S$-态射，则仍以
$X(T)_S$ 表示 $S$-态射 $T\to X$ 的集合 $\operatorname{Hom}_S(T,X)$；
在不会混淆时，容许省略下标 $S$。仍称 $X(T)_S$ 的元素为
\emph{$S$-预概形 $X$ 的取值于 $S$-预概形 $T$ 的点}；在可能混淆时也称为
\emph{$S$-点}。特别地，$X$ 的一个\emph{$S$-截面}无非就是 $X$ 的一个
\emph{取值于 $S$ 的点}。公式（\hyperref[I.3.4.2.1-fr]{3.4.2.1}）也可写成
\phantomsection\label{I.3.4.3.1-fr}
\[
  (X\times_S Y)(T)_S=X(T)_S\times Y(T)_S~;
  \tag{3.4.3.1}
\]
更一般地，若 $Z$ 是一个 $S$-预概形，而 $X$、$Y$、$T$ 是
$Z$-预概形（因而当然也是 $S$-预概形），则有
\phantomsection\label{I.3.4.3.2-fr}
\[
  (X\times_Z Y)(T)_S=X(T)_S\times_{Z(T)_S}Y(T)_S.
  \tag{3.4.3.2}
\]

注意，为了证明一个由 $S$-预概形 $W$ 和两个 $S$-态射
$r:W\to X$、$s:W\to Y$ 组成的三元组 $(W,r,s)$ 是 $X$ 与 $Y$
（在 $Z$ 上）的积，按照定义，只需验证：对于\emph{任意 $S$-预概形 $T$}，
图表
\phantomsection\label{I.3.4.3.product-diagram-fr}
\[
  \xymatrix{
    W(T)_S\ar[r]^{r'}\ar[d]_{s'} & X(T)_S\ar[d]^{\varphi'}\\
    Y(T)_S\ar[r]_{\psi'} & Z(T)_S
  }
\]
使 $W(T)_S$ 成为 $X(T)_S$ 与 $Y(T)_S$ 在 $Z(T)_S$ 上的纤维积；这里
$r'$、$s'$ 分别对应于 $r$、$s$，$\varphi'$、$\psi'$ 分别对应于结构态射
$\varphi:X\to Z$、$\psi:Y\to Z$。
\end{env}

\begin{env}[3.4.4]
\label{I.3.4.4-fr}
在上述情形中，若 $T$（相应地，$S$）是环为 $B$（相应地，$A$）的仿射概形，
就在前面的记号中以 $B$（相应地，$A$）代替 $T$（相应地，$S$）；此时，
分别称 $X(B)$ 和 $X(B)_A$ 的元素为\emph{$X$ 的取值于环 $B$ 的点}以及
\emph{$A$-预概形 $X$ 的取值于 $A$-代数 $B$ 的点}。注意，$X(B)$ 和
$X(B)_A$ 现在是关于 $B$ 的\emph{协变函子}。类似地，以 $X(T)_A$ 表示
$A$-预概形 $X$ 的取值于 $A$-预概形 $T$ 的点的集合。
\end{env}

\begin{env}[3.4.5]
\label{I.3.4.5-fr}
特别考虑 $T$ 形如 $\operatorname{Spec}(A)$ 的情形，其中 $A$ 是一个
\emph{局部}环；
此时，$X(A)$ 的元素与局部同态 $\mathcal{O}_x\to A$（$x\in X$）一一对应
（\hyperref[I.2.4.4-fr]{2.4.4}）。称 $X$ 的底空间中的点 $x$ 为它所对应的
$X$ 的 $A$-值点的\emph{所在点}。

更特别地，称一个预概形 $X$ 的\emph{几何点}为\emph{$X$ 的取值于某个域
$K$ 的点}：因此，给定这样的点等价于给定它在 $X$ 的底空间中的
\oldpage[I]{113}
所在点 $x$，以及 $k(x)$ 的一个\emph{扩张} $K$；称 $K$ 为相应几何点的
\emph{取值域}，并且也说这个几何点\emph{位于 $x$ 上}。这样就定义了一个映射
$X(K)\to X$，它把取值于 $K$ 的几何点映到其所在点。

若 $S'=\operatorname{Spec}(K)$ 是一个 $S$-预概形（换言之，若把 $K$
看作某个剩余域 $k(s)$ 的扩张，其中 $s\in S$），而 $X$ 是一个
$S$-预概形，则 $X(K)_S$ 的一个元素，或者说一个\emph{$X$ 的位于 $s$
上方且取值于 $K$ 的几何点}，就是给定一个从剩余域 $k(x)$ 到 $K$ 的
$k(s)$-单同态，其中 $x$ 是 $X$ 的一个\emph{位于 $s$ 上方的点}
（因而 $k(x)$ 是 $k(s)$ 的一个扩张）。

更特别地，若 $S=\operatorname{Spec}(K)=\{\xi\}$，则
\emph{$X$ 的取值于 $K$ 的几何点可以与满足 $k(x)=K$ 的点 $x\in X$
等同}；也称后者为\emph{$K$-预概形 $X$ 的 $K$-有理点}。若 $K'$ 是
$K$ 的一个扩张，则 $X$ 的取值于 $K'$ 的几何点与
$X'=X_{(K')}$ 的 $K'$-有理点一一对应
（\hyperref[I.3.3.14-fr]{3.3.14}）。
\end{env}

\begin{lemma}[3.4.6]
\label{I.3.4.6-fr}
设 $X_i$（$1\leq i\leq n$）是 $S$-预概形，$s$ 是 $S$ 的一个点，
$x_i$（$1\leq i\leq n$）是 $X_i$ 的位于 $s$ 上方的点。于是，存在
$k(s)$ 的一个扩张 $K$，以及积
$Y=X_1\times_S X_2\times_S\cdots\times_S X_n$ 的一个取值于 $K$
的几何点，使它的各个投影分别位于 $x_i$ 上。
\end{lemma}

事实上，在 $k(s)$ 的某个共同扩张 $K$ 中，存在 $k(s)$-单同态
$k(x_i)\to K$（Bourbaki，\emph{Alg.}，第 V 章，\S~4，命题 2）。
各复合 $k(s)\to k(x_i)\to K$ 都相同，所以与 $k(x_i)\to K$ 对应的态射
$\operatorname{Spec}(K)\to X_i$ 都是 $S$-态射，由此可知，它们定义一个唯一的
态射 $\operatorname{Spec}(K)\to Y$。若 $y$ 是 $Y$ 的相应点，则显然它在每个
$X_i$ 上的投影都是 $x_i$。

\begin{proposition}[3.4.7]
\label{I.3.4.7-fr}
设 $X_i$（$1\leq i\leq n$）是 $S$-预概形，并且对每个指标 $i$，给定
$X_i$ 的一个点 $x_i$。存在积
$Y=X_1\times_S X_2\times_S\cdots\times_S X_n$ 的一个点 $y$，使得
其第 $i$ 个投影是 $x_i$（$1\leq i\leq n$）的充要条件是：各 $x_i$
都位于 $S$ 的同一个点 $s$ 上方。
\end{proposition}

这个条件显然是必要的；引理（\hyperref[I.3.4.6-fr]{3.4.6}）证明了它是充分的。

\phantomsection\label{I.3.4.7.underlying-map-fr}
换言之，若以 $(X)$ 表示 $X$ 的底集合，则有一个典范的\emph{满射}
$(X\times_S Y)\to (X)\times_{(S)}(Y)$；应当注意，这个映射一般
\emph{不是单射}。也就是说，\emph{$X\times_S Y$ 中可以存在多个互异的点
$z$，它们具有相同的投影 $x\in X$、$y\in Y$}。当 $S$、$X$、$Y$
分别是域 $k$、$K$、$K'$ 的素谱时，已经可以看出这一点，因为张量积
$K\otimes_k K'$ 一般具有多个互异的素理想
（参见 \hyperref[I.3.4.9-fr]{3.4.9}）。

\begin{corollary}[3.4.8]
\label{I.3.4.8-fr}
设 $f:X\to Y$ 是一个 $S$-态射，而
$f_{(S')}:X_{(S')}\to Y_{(S')}$ 是由底预概形的扩张 $S'\to S$ 从 $f$
导出的 $S'$-态射。设 $p$（相应地，$q$）是投影
$X_{(S')}\to X$（相应地，$Y_{(S')}\to Y$）；则对于 $X$ 的任意子集
$M$，都有
\phantomsection\label{I.3.4.8.base-change-image-fr}
\[
  q^{-1}(f(M))=f_{(S')}(p^{-1}(M)).
\]
\end{corollary}

\oldpage[I]{114}
事实上，根据（\hyperref[I.3.3.11-fr]{3.3.11}），通过交换图表
\phantomsection\label{I.3.4.8.proof-diagram-fr}
\[
  \xymatrix{
    X\ar[d]_f & X_{(S')}\ar[l]_p\ar[d]^{f_{(S')}}\\
    Y & Y_{(S')}\ar[l]^q
  }
\]
$X_{(S')}$ 等同于积 $X\times_Y Y_{(S')}$。由
（\hyperref[I.3.4.7-fr]{3.4.7}），对于 $x\in M$、$y'\in Y_{(S')}$，
关系 $q(y')=f(x)$ 等价于存在 $x'\in X_{(S')}$，使
$p(x')=x$ 且 $f_{(S')}(x')=y'$，推论即得。

引理（\hyperref[I.3.4.6-fr]{3.4.6}）可以进一步明确如下：

\begin{proposition}[3.4.9]
\label{I.3.4.9-fr}
设 $X$、$Y$ 是两个 $S$-预概形，$x$ 是 $X$ 的一个点，$y$ 是 $Y$
的一个点，并且二者都位于同一点 $s\in S$ 上方。$X\times_S Y$ 中以
$x$、$y$ 为投影的点的集合，与把 $k(x)$、$k(y)$ 看作 $k(s)$ 的扩张时
它们的复合扩张类型的集合存在典范的一一对应
（Bourbaki，\emph{Alg.}，第 VIII 章，\S~8，命题 2）。
\end{proposition}

设 $p$（相应地，$q$）是从 $X\times_S Y$ 到 $X$（相应地，$Y$）的投影，
而 $E$ 是 $X\times_S Y$ 的底空间的子空间
$p^{-1}(x)\cap q^{-1}(y)$。首先注意，态射
$\operatorname{Spec}(k(x))\to S$ 和 $\operatorname{Spec}(k(y))\to S$
分别分解为
$\operatorname{Spec}(k(x))\to\operatorname{Spec}(k(s))\to S$ 和
$\operatorname{Spec}(k(y))\to\operatorname{Spec}(k(s))\to S$；由于
$\operatorname{Spec}(k(s))\to S$ 是一个单态射
（\hyperref[I.2.4.7-fr]{2.4.7}），由（\hyperref[I.3.2.4-fr]{3.2.4}）得到
\phantomsection\label{I.3.4.9.tensor-product-fibre-fr}
\[
  P=\operatorname{Spec}(k(x))\times_S\operatorname{Spec}(k(y))
  =\operatorname{Spec}(k(x))\times_{\operatorname{Spec}(k(s))}
   \operatorname{Spec}(k(y))
  =\operatorname{Spec}\bigl(k(x)\otimes_{k(s)}k(y)\bigr).
\]

我们将定义两个互为逆映射的映射 $\alpha:P_0\to E$、$\beta:E\to P_0$
（其中 $P_0$ 表示预概形 $P$ 的底集合）。若
$i:\operatorname{Spec}(k(x))\to X$、$j:\operatorname{Spec}(k(y))\to Y$
是典范态射（\hyperref[I.2.4.5-fr]{2.4.5}），则取 $\alpha$ 为态射
$i\times_S j$ 所对应的底空间映射。另一方面，每个 $z\in E$ 按照假设都定义
两个 $k(s)$-单同态 $k(x)\to k(z)$、$k(y)\to k(z)$，因而定义一个
% EGA1-ERR-0030: French print/source calls the induced tensor-product map a monomorphism; it need only be a homomorphism. Canonical Chinese retains 单同态.
$k(s)$-单同态 $k(x)\otimes_{k(s)}k(y)\to k(z)$，从而定义一个态射
$\operatorname{Spec}(k(z))\to P$；令 $\beta(z)$ 为 $z$ 经这个态射在
$P_0$ 中的像。由（\hyperref[I.2.4.5-fr]{2.4.5}）和积的定义
（\hyperref[I.3.2.1-fr]{3.2.1}），立即验证 $\alpha\circ\beta$ 和
$\beta\circ\alpha$ 都是恒等映射。最后，已知 $P_0$ 与 $k(x)$、$k(y)$
的复合扩张类型的集合一一对应
（Bourbaki，\emph{Alg.}，第 VIII 章，\S~8，命题 1）。

\subsection{映满和含容}
\label{subsection:I.3.5-fr}

\begin{env}[3.5.1]
\label{I.3.5.1-fr}
一般地，考虑预概形态射的一个性质 $\mathbf{P}$，以及以下两个命题：
\begin{enumerate}
  \item[(i)] 若 $f:X\to X'$、$g:Y\to Y'$ 是两个具有性质
  $\mathbf{P}$ 的 $S$-态射，则 $f\times_S g$ 具有性质 $\mathbf{P}$。
  \item[(ii)] 若 $f:X\to Y$ 是一个具有性质 $\mathbf{P}$ 的 $S$-态射，
  则由底预概形的扩张从 $f$ 导出的任意 $S'$-态射
  $f_{(S')}:X_{(S')}\to Y_{(S')}$ 都具有性质 $\mathbf{P}$。
\end{enumerate}

\oldpage[I]{115}
由于 $f_{(S')}=f\times_S 1_{S'}$，可见，若对每个预概形 $X$，
\emph{恒等态射} $1_X$ 都具有性质 $\mathbf{P}$，则 (i) 蕴含 (ii)。
另一方面，由于 $f\times_S g$ 是复合态射
\phantomsection\label{I.3.5.1.product-composition-fr}
\[
  X\times_S Y\xrightarrow{f\times 1_Y}X'\times_S Y
  \xrightarrow{1_{X'}\times g}X'\times_S Y'
\]
可见，若两个具有性质 $\mathbf{P}$ 的态射之\emph{复合}也具有这个性质，
则 (ii) 蕴含 (i)。
\end{env}

这个说明的第一个应用是

\begin{proposition}[3.5.2]
\label{I.3.5.2-fr}
\begin{enumerate}
  \item[(i)] 若 $f:X\to X'$、$g:Y\to Y'$ 是满射的 $S$-态射，则
  $f\times_S g$ 是满射。
  \item[(ii)] 若 $f:X\to Y$ 是一个满射的 $S$-态射，则对于底预概形的
  任意扩张 $S'$，$f_{(S')}$ 都是满射。
\end{enumerate}
\end{proposition}

两个满射的复合是满射，因此只需证明 (ii)；而这个命题立即由
（\hyperref[I.3.4.8-fr]{3.4.8}）对 $M=X$ 的应用得出。

\begin{proposition}[3.5.3]
\label{I.3.5.3-fr}
态射 $f:X\to Y$ 是满射的充要条件是：对于任意域 $K$ 和任意态射
$\operatorname{Spec}(K)\to Y$，存在 $K$ 的一个扩张 $K'$ 和一个态射
$\operatorname{Spec}(K')\to X$，使图表
\phantomsection\label{I.3.5.3.surjectivity-diagram-fr}
\[
  \xymatrix{
    X\ar[d]_f & \operatorname{Spec}(K')\ar[l]\ar[d]\\
    Y & \operatorname{Spec}(K)\ar[l]
  }
\]
交换。
\end{proposition}

这个条件是充分的，因为对于任意 $y\in Y$，只需把它应用于与某个单同态
$k(y)\to K$ 对应的态射 $\operatorname{Spec}(K)\to Y$，其中 $K$ 是
$k(y)$ 的一个扩张（\hyperref[I.2.4.6-fr]{2.4.6}）。反过来，设 $f$
是满射，并令 $y\in Y$ 为 $\operatorname{Spec}(K)$ 的唯一一点的像；
存在 $x\in X$ 使 $f(x)=y$。考虑相应的单同态
$k(y)\to k(x)$（\hyperref[I.2.2.1-fr]{2.2.1}）；此时只需取
$k(y)$ 的一个扩张 $K'$，使得存在从 $k(x)$ 和 $K$ 到 $K'$ 的
$k(y)$-单同态（Bourbaki，\emph{Alg.}，第 V 章，\S~4，命题 2）。
与 $k(x)\to K'$ 对应的态射 $\operatorname{Spec}(K')\to X$ 即满足要求。

采用（\hyperref[I.3.4.5-fr]{3.4.5}）引入的语言，也可以说：
\emph{$Y$ 的每个取值于 $K$ 的几何点，都来自 $X$ 的一个取值于 $K$
的某个扩张的几何点}。

\begin{definition}[3.5.4]
\label{I.3.5.4-fr}
若对于任意域 $K$，相应的映射 $X(K)\to Y(K)$ 都是单射，则称预概形态射
$f:X\to Y$ 是\emph{普遍单射的}，或称它为一个\emph{根态射}。
\end{definition}

由定义立即可知，预概形的每个\emph{单态射}（T, 1.1）都是根态射。

\begin{env}[3.5.5]
\label{I.3.5.5-fr}
为了使态射 $f:X\to Y$ 成为根态射，只需定义
（\hyperref[I.3.5.4-fr]{3.5.4}）中的条件对每个\emph{代数闭}域成立。
事实上，设 $K$ 是任意域，$K'$ 是 $K$ 的一个代数闭扩张；图表
\phantomsection\label{I.3.5.5.algebraic-closure-diagram-fr}
\[
  \xymatrix{
    X(K)\ar[r]^\alpha\ar[d]_\varphi & Y(K)\ar[d]^{\varphi'}\\
    X(K')\ar[r]_{\alpha'} & Y(K')
  }
\]

\oldpage[I]{116}
交换；其中 $\varphi$、$\varphi'$ 来自态射
$\operatorname{Spec}(K')\to\operatorname{Spec}(K)$，而 $\alpha$、
$\alpha'$ 对应于 $f$。映射 $\varphi$ 是单射，而根据假设
$\alpha'$ 也是单射；因此 $\alpha$ 必为单射。
\end{env}

\begin{proposition}[3.5.6]
\label{I.3.5.6-fr}
设 $f:X\to Y$、$g:Y\to Z$ 是两个预概形态射。
\begin{enumerate}
  \item[(i)] 若 $f$ 和 $g$ 都是根态射，则 $g\circ f$ 也是根态射。
  \item[(ii)] 反过来，若 $g\circ f$ 是根态射，则 $f$ 也是根态射。
\end{enumerate}
\end{proposition}

考虑到定义（\hyperref[I.3.5.4-fr]{3.5.4}），这个命题归结为映射
$X(K)\to Y(K)\to Z(K)$ 的相应断言，而这些断言是显然的。

\begin{proposition}[3.5.7]
\label{I.3.5.7-fr}
\begin{enumerate}
  \item[(i)] 若 $S$-态射 $f:X\to X'$、$g:Y\to Y'$ 都是根态射，则
  $f\times_S g$ 也是根态射。
  \item[(ii)] 若 $S$-态射 $f:X\to Y$ 是根态射，则对于底预概形的任意
  扩张 $S'\to S$，$f_{(S')}:X_{(S')}\to Y_{(S')}$ 也是根态射。
\end{enumerate}
\end{proposition}

根据（\hyperref[I.3.5.1-fr]{3.5.1}），只需证明 (i)。我们已经看到
（\hyperref[I.3.4.2.1-fr]{3.4.2.1}）
\phantomsection\label{I.3.5.7.field-valued-product-fr}
\[
  (X\times_S Y)(K)=X(K)\times_{S(K)}Y(K),\qquad
  (X'\times_S Y')(K)=X'(K)\times_{S(K)}Y'(K)
\]
而与 $f\times_S g$ 对应的映射
$(X\times_S Y)(K)\to(X'\times_S Y')(K)$ 因而等同于
$(u,v)\mapsto(f\circ u,g\circ v)$，命题立即得证。

\begin{proposition}[3.5.8]
\label{I.3.5.8-fr}
态射 $f=(\psi,\theta):X\to Y$ 是根态射的充要条件是：$\psi$ 是单射，
并且对于每个 $x\in X$，单同态
$\theta^x:k(\psi(x))\to k(x)$ 使 $k(x)$ 成为 $k(\psi(x))$ 的
纯不可分扩张。
\end{proposition}

设 $f$ 是根态射。首先证明关系 $\psi(x_1)=\psi(x_2)=y$ 必然蕴含
$x_1=x_2$。事实上，存在 $k(y)$ 的一个扩张域 $K$，以及
$k(y)$-单同态 $k(x_1)\to K$、$k(x_2)\to K$
（Bourbaki，\emph{Alg.}，第 V 章，\S~4，命题 2）；相应的态射
$u_1:\operatorname{Spec}(K)\to X$、$u_2:\operatorname{Spec}(K)\to X$
满足 $f\circ u_1=f\circ u_2$，所以由假设有 $u_1=u_2$，这特别蕴含
$x_1=x_2$。现在通过 $\theta^x$ 把 $k(x)$ 看作 $k(\psi(x))$ 的扩张：
若 $k(x)$ 不是 $k(\psi(x))$ 的纯不可分扩张，则存在从 $k(x)$ 到
$k(\psi(x))$ 的某个代数闭扩张 $K$ 的两个互异的
$k(\psi(x))$-单同态，而相应的两个态射
$\operatorname{Spec}(K)\to X$ 将违反假设。反过来，考虑到
（\hyperref[I.2.4.6-fr]{2.4.6}），命题中的条件足以使 $f$ 成为根态射，
这是立即可见的。

\begin{corollary}[3.5.9]
\label{I.3.5.9-fr}
设 $A$ 是一个环，$S$ 是 $A$ 的一个乘性子集；则典范态射
$\operatorname{Spec}(S^{-1}A)\to\operatorname{Spec}(A)$ 是根态射。
\end{corollary}

事实上，这个态射是单态射（\hyperref[I.1.6.2-fr]{1.6.2}）。

\begin{corollary}[3.5.10]
\label{I.3.5.10-fr}
设 $f:X\to Y$ 是一个根态射，$g:Y'\to Y$ 是一个态射，并令
$X'=X_{(Y')}=X\times_Y Y'$。于是，根态射 $f_{(Y')}$
（\hyperref[I.3.5.7-fr]{3.5.7 (ii)}）是从 $X'$ 的底空间到
$g^{-1}(f(X))$ 的一个双射；此外，对于任意域 $K$，集合 $X'(K)$
在 $Y'(K)$ 中等同于 $Y(K)$ 的子集 $X(K)$ 在与 $g$ 对应的映射
$Y'(K)\to Y(K)$ 下的逆像。
\end{corollary}

第一个断言立即由（\hyperref[I.3.5.8-fr]{3.5.8}）和
（\hyperref[I.3.4.8-fr]{3.4.8}）得出；第二个断言由图表
\oldpage[I]{117}
\phantomsection\label{I.3.5.10.base-change-points-diagram-fr}
\[
  \xymatrix{
    X'(K)\ar[r]\ar[d] & Y'(K)\ar[d]\\
    X(K)\ar[r] & Y(K)
  }
\]
的交换性得出。

\begin{remark}[3.5.11]
\label{I.3.5.11-fr}
若预概形态射 $f=(\psi,\theta)$ 中的映射 $\psi$ 是单射，则称这个态射
是\emph{单射的}。态射 $f=(\psi,\theta):X\to Y$ 是根态射的充要条件是：
对于任意态射 $Y'\to Y$，态射 $f_{(Y')}:X_{(Y')}\to Y'$ 都是单射
（这说明了“普遍单射态射”这一术语）。事实上，根据
（\hyperref[I.3.5.7-fr]{3.5.7, (ii)}）和
（\hyperref[I.3.5.8-fr]{3.5.8}），这个条件是必要的。反过来，这个条件
首先蕴含 $\psi$ 是单射；若对某个 $x\in X$，单同态
$\theta^x:k(\psi(x))\to k(x)$ 不是纯不可分的，则存在
$k(\psi(x))$ 的一个扩张 $K$ 以及两个互异的态射
$\operatorname{Spec}(K)\to X$，它们对应于同一个态射
$\operatorname{Spec}(K)\to Y$（\hyperref[I.3.5.8-fr]{3.5.8}）。
但是，令 $Y'=\operatorname{Spec}(K)$，就会得到 $X_{(Y')}$ 的两个互异的
$Y'$-截面（\hyperref[I.3.3.14-fr]{3.3.14}），这与
$f_{(Y')}$ 是单射的假设矛盾。
\end{remark}

\subsection{纤维}
\label{subsection:I.3.6-fr}

\begin{proposition}[3.6.1]
\label{I.3.6.1-fr}
设 $f:X\to Y$ 是一个态射，$y$ 是 $Y$ 的一个点，而
$\mathfrak{a}_y$ 是 $\mathcal{O}_y$ 对于
$\mathfrak{m}_y$-预进制拓扑的一个定义理想。投影
$p:X\times_Y\operatorname{Spec}(\mathcal{O}_y/\mathfrak{a}_y)\to X$
是从预概形
$X\times_Y\operatorname{Spec}(\mathcal{O}_y/\mathfrak{a}_y)$ 的底空间
到配备了由 $X$ 的底空间拓扑所诱导拓扑的纤维 $f^{-1}(y)$ 的一个同胚。
\end{proposition}

由于 $\operatorname{Spec}(\mathcal{O}_y/\mathfrak{a}_y)\to Y$ 是根态射
（\hyperref[I.3.5.4-fr]{3.5.4} 和 \hyperref[I.2.4.7-fr]{2.4.7}），
并且 $\operatorname{Spec}(\mathcal{O}_y/\mathfrak{a}_y)$ 只含一个点
（按假设，理想 $\mathfrak{m}_y/\mathfrak{a}_y$ 是诣零的，
见 \hyperref[I.1.1.12-fr]{1.1.12}），我们已经知道
（\hyperref[I.3.5.10-fr]{3.5.10} 和 \hyperref[I.3.3.4-fr]{3.3.4}），
$p$ 作为\emph{集合}的映射把
$X\times_Y\operatorname{Spec}(\mathcal{O}_y/\mathfrak{a}_y)$ 的底空间
与 $f^{-1}(y)$ 等同；问题归结为证明 $p$ 是同胚。根据
（\hyperref[I.3.2.7-fr]{3.2.7}），这个问题关于 $X$ 和 $Y$ 是局部的，
因而可以假设 $X=\operatorname{Spec}(B)$、$Y=\operatorname{Spec}(A)$，
其中 $B$ 是一个 $A$-代数。此时，态射 $p$ 对应于同态
$1\otimes\varphi:B\to B\otimes_A A'$，其中
$A'=A_y/\mathfrak{a}_y$，而 $\varphi$ 是从 $A$ 到 $A'$ 的典范映射。
$B\otimes_A A'$ 的每个元素都可写成
\phantomsection\label{I.3.6.1.localization-fraction-fr}
\[
  \sum_i b_i\otimes\varphi(a_i)/\varphi(s)
  =\left(\sum_i(a_i b_i\otimes 1)\right)(1\otimes\varphi(s))^{-1},
\]
其中 $s\notin\mathfrak{j}_y$，因而可以应用命题
（\hyperref[I.1.2.4-fr]{1.2.4}）。

\begin{env}[3.6.2]
\label{I.3.6.2-fr}
在本论著的下文中，每当我们考虑态射的一个纤维 $f^{-1}(y)$，并给它配备
$k(y)$-预概形结构时，\emph{所指总是借助到 $X$ 的投影，把
$X\times_Y\operatorname{Spec}(k(y))$ 的结构转移过去所得到的预概形}。
也把后一个积写成 $X\otimes_Y k(y)$ 或
$X\otimes_{\mathcal{O}_Y}k(y)$；更一般地，若 $B$ 是一个
$\mathcal{O}_y$-代数，则以 $X\otimes_Y B$ 或
$X\otimes_{\mathcal{O}_Y}B$ 表示积
$X\times_Y\operatorname{Spec}(B)$。
\end{env}

根据上述约定，由（\hyperref[I.3.5.10-fr]{3.5.10}）可知，$X$ 的取值于
$k(y)$ 的一个扩张 $K$ 的点，与\emph{$f^{-1}(y)$ 的取值于 $K$ 的点}
等同。

\begin{env}[3.6.3]
\label{I.3.6.3-fr}
设 $f:X\to Y$、$g:Y\to Z$ 是两个态射，$h=g\circ f$ 是它们的复合；
对于任意 $z\in Z$，纤维 $h^{-1}(z)$ 是一个同构于
\phantomsection\label{I.3.6.3.fibre-composition-fr}
\[
  X\times_Z\operatorname{Spec}(k(z))
  =(X\times_Y Y)\times_Z\operatorname{Spec}(k(z))
  =X\times_Y g^{-1}(z)
\]
的预概形。
\oldpage[I]{118}
特别地，若 $U$ 是 $X$ 的一个开子集，则预概形 $f^{-1}(y)$ 在
$U\cap f^{-1}(y)$ 上诱导的预概形同构于 $f_U^{-1}(y)$
（$f_U$ 表示 $f$ 在 $U$ 上的限制）。
\end{env}

\begin{proposition}[3.6.4)（“纤维的传递性”）]
\label{I.3.6.4-fr}
设 $f:X\to Y$、$g:Y'\to Y$ 是两个态射；令
$X'=X\times_Y Y'=X_{(Y')}$，并令
$f'=f_{(Y')}:X'\to Y'$。对于任意 $y'\in Y'$，若令
$y=g(y')$，则预概形 $f'^{-1}(y')$ 同构于
$f^{-1}(y)\otimes_{k(y)}k(y')$。
\end{proposition}

事实上，只需注意，根据（\hyperref[I.3.3.9.1-fr]{3.3.9.1}），两个预概形
$(X\otimes_Y k(y))\otimes_{k(y)}k(y')$ 和
$(X\times_Y Y')\otimes_{Y'}k(y')$ 都典范地同构于
$X\times_Y\operatorname{Spec}(k(y'))$。

特别地，若 $V$ 是 $Y$ 中 $y$ 的一个开邻域，并以 $f_V$ 表示 $f$
在 $f^{-1}(V)$ 上所诱导的预概形上的限制，则预概形
$f^{-1}(y)$ 和 $f_V^{-1}(y)$ 典范地等同。

\begin{proposition}[3.6.5]
\label{I.3.6.5-fr}
设 $f:X\to Y$ 是一个态射，$y$ 是 $Y$ 的一个点，
$Z$ 是局部预概形 $\operatorname{Spec}(\mathcal{O}_y)$，
$p=(\psi,\theta)$ 是投影 $X\times_Y Z\to X$。当把 $Z$ 的底空间
等同于 $Y$ 的一个子空间时（参见
（\hyperref[I.2.4.2-fr]{2.4.2}）），$p$ 是从 $X\times_Y Z$ 的底空间
到 $X$ 的子空间 $f^{-1}(Z)$ 的一个同胚；而且，对于任意
$t\in X\times_Y Z$，若令 $z=\psi(t)$，则 $\theta_t^\sharp$ 是从
$\mathcal{O}_z$ 到 $\mathcal{O}_t$ 的一个同构。
\end{proposition}

由于 $Z$（等同于 $Y$ 的一个子空间）包含在每个含 $y$ 的仿射开集中
（\hyperref[I.2.4.2-fr]{2.4.2}），可以像
（\hyperref[I.3.6.1-fr]{3.6.1}）那样，把问题归结到
$X=\operatorname{Spec}(A)$、$Y=\operatorname{Spec}(B)$ 都是仿射概形的
情形，其中 $A$ 是一个 $B$-代数。此时，$X\times_Y Z$ 是
$A\otimes_B B_y$ 的素谱，而这个环典范地等同于 $S^{-1}A$，其中
$S$ 是 $B-\mathfrak{j}_y$ 在 $A$ 中的像（0, 1.5.2）。由于 $p$
此时对应于典范同态 $A\to S^{-1}A$，命题由
（\hyperref[I.1.6.2-fr]{1.6.2}）得出。

\subsection{应用：预概形的模
\texorpdfstring{$\mathfrak{J}$\protect\footnotemark}{J} 约化}
\label{subsection:I.3.7-fr}
\label{I.3.7-fr}
\footnotetext{本节使用第一章后文及第二章的概念和结果；本书后文不再使用
本节，本节只供熟悉经典代数几何的读者阅读。}

\begin{env}[3.7.1]
\label{I.3.7.1-fr}
设 $A$ 是一个环，$X$ 是一个 $A$-预概形，$\mathfrak{J}$ 是 $A$ 的一个
理想；则 $X_0=X\otimes_A(A/\mathfrak{J})$ 是一个
$(A/\mathfrak{J})$-预概形，有时称它由 $X$ 通过模 $\mathfrak{J}$
\emph{约化}而得。
\end{env}

\begin{env}[3.7.2]
\label{I.3.7.2-fr}
这个术语主要用于 $A$ 是一个\emph{局部环}而 $\mathfrak{J}$ 是其极大
理想的情形；于是 $X_0$ 是 $A$ 的剩余域 $k=A/\mathfrak{J}$ 上的一个
预概形。

若此外 $A$ 是整环，其分式域为 $K$，还可以考虑 $K$-预概形
$X'=X\otimes_A K$。按照一种我们并不采用的语言滥用，以往通常说
$X_0$ 是由 $X'$ 模 $\mathfrak{J}$ 约化而得。在使用这种说法的情形中，
$A$ 是一个维数为 1 的局部环（通常是离散赋值环），并且（或多或少
明确地）默认给定的 $K$-预概形 $X'$ 是某个 $K$-预概形 $P'$ 的闭子预概形
（事实上，是形如 $\mathbf{P}_K^r$ 的射影空间，参见 II, 4.1.1），而
后者本身形如 $P'=P\otimes_A K$，其中 $P$ 是一个给定的 $A$-预概形
（事实上，按 II, 4.1.1 的记号，它是 $A$-概形 $\mathbf{P}_A^r$）。
用我们的语言，由 $X'$ 定义 $X_0$ 可表述如下：

考虑仿射概形 $Y=\operatorname{Spec}(A)$，它由两个点组成：唯一的闭点
\oldpage[I]{119}
$y=\mathfrak{J}$ 和一般点 $(0)$；只含一般点的集合 $U$ 因而是 $Y$ 中的
开集 $U=\operatorname{Spec}(K)$。若 $X$ 是一个 $A$-预概形（换言之，
一个 $Y$-预概形），并以 $\psi$ 表示结构态射 $X\to Y$，则
$X\otimes_A K=X'$ 不过是 $X$ 在 $\psi^{-1}(U)$ 上所诱导的预概形。
特别地，若 $\varphi$ 是结构态射 $P\to Y$，则
$P'=\varphi^{-1}(U)$ 的闭子预概形 $X'$ 是 $P$ 的一个（局部闭）
子预概形。若 $P$ 是 Noether 的（例如，若 $A$ 是 Noether 的，且 $P$
在 $A$ 上为有限型），则存在 $P$ 中包含 $X$ 的最小闭子预概形
% EGA1-ERR-0031: NUMDAM printed p.119 omits the prime; canonical Chinese remains print-literal.
$X=\overline{X'}$（9.5.10）；而 $X'$ 是 $X$ 在开集
$\varphi^{-1}(U)\cap X$ 上所诱导的预概形，因而同构于
$X\otimes_A K$（9.5.10）。
\emph{$X'$ 到 $P'=P\otimes_A K$ 中的浸入，因而使我们可以典范地把
$X'$ 看成形如 $X'=X\otimes_A K$，其中 $X$ 是一个 $A$-预概形。}
这时可以考虑模 $\mathfrak{J}$ 约化所得的预概形
$X_0=X\otimes_A k$；它正是闭点 $y$ 的纤维 $\psi^{-1}(y)$。以往由于
缺少适当的术语，人们避免显式引入 $A$-预概形 $X$。然而应当注意，
通常关于“模 $\mathfrak{J}$ 约化所得的”预概形 $X_0$ 所作的一切断言，
都应看作关于 $X$ 本身的更完整断言的推论；只有这样解释，才能令人满意地
表述和理解这些断言。此外，所作的假设似乎总可归结为对 $X$ 本身的假设
（与预先给定的 $X'$ 到 $\mathbf{P}_K^r$ 中的浸入无关），从而可以给出
更为内蕴的陈述。
\end{env}

\begin{env}[3.7.3]
\label{I.3.7.3-fr}
最后指出一个十分特殊的事实；这一事实无疑曾使这里所讨论情形的概念澄清
受到延误：若 $A$ 是一个离散赋值环，且 $X$ \emph{在 $A$ 上固有}
（若 $X$ 是 $\mathbf{P}_A^r$ 的一个闭子预概形，确实如此，参见
II, 5.5.4），则 $X$ 的 $A$-值点与 $X'$ 的 $K$-值点一一对应
（II, 7.3.8）。正因如此，人们常以为证明了关于 $X'$ 的结果，实际上
证明的却是关于 $X$ 的命题；即使不再假设底局部环的维数为 1，这些命题
（以这种形式）仍然成立。
\end{env}
% EGA I mainland Simplified-Chinese translation (producer output; uncertified).
% Sequential translation cursor: through 4.1.10 (source lines 1--217; printed pages 119--122).

\section{子预概形与浸入态射}
\label{section:I.4-fr}

\subsection{子预概形}
\label{subsection:I.4.1-fr}

\begin{env}[4.1.1]
\label{I.4.1.1-fr}
由于拟凝聚层的概念（\hyperref[0.5.1.3-fr]{0, 5.1.3}）是局部的，预概形
$X$ 上的拟凝聚 $\mathscr{O}_X$-模层 $\mathscr{F}$ 可以由下述条件定义：
对 $X$ 的每个仿射开集 $V$，$\mathscr{F}|V$ 都同构于某个
$\Gamma(V,\mathscr{O}_X)$-模所关联的层
（\hyperref[I.1.4.1-fr]{1.4.1}）。显然，在预概形 $X$ 上，结构层
$\mathscr{O}_X$ 是拟凝聚的；拟凝聚 $\mathscr{O}_X$-模层同态的核、
余核和像，以及拟凝聚 $\mathscr{O}_X$-模层的归纳极限与直和，也都是
拟凝聚的（\hyperref[I.1.3.7-fr]{1.3.7} 与
\hyperref[I.1.3.9-fr]{1.3.9}）。
\end{env}

\begin{proposition}[4.1.2]
\label{I.4.1.2-fr}
设 $X$ 是预概形，$\mathscr{I}$ 是 $\mathscr{O}_X$ 中的一个拟凝聚理想层。
则层 $\mathscr{O}_X/\mathscr{I}$ 的支集 $Y$ 是闭的；若以
$\mathscr{O}_Y$ 表示 $\mathscr{O}_X/\mathscr{I}$ 在 $Y$ 上的限制，
则 $(Y,\mathscr{O}_Y)$ 是一个预概形。
\end{proposition}

\oldpage[I]{120}
显然，只需（\hyperref[I.2.1.3-fr]{2.1.3}）考虑 $X$ 是仿射概形的情形，
并证明此时 $Y$ 在 $X$ 中是闭的，而且是一个\emph{仿射概形}。事实上，
若 $X=\operatorname{Spec}(A)$，则 $\mathscr{O}_X=\widetilde{A}$ 且
$\mathscr{I}=\widetilde{\mathfrak{I}}$，其中 $\mathfrak{I}$ 是 $A$ 的一个
理想（\hyperref[I.1.4.1-fr]{1.4.1}）；这时 $Y$ 等于 $X$ 的闭子集
$V(\mathfrak{I})$，并等同于环 $B=A/\mathfrak{I}$ 的素谱
（\hyperref[I.1.1.11-fr]{1.1.11}）。此外，若 $\varphi$ 是典范同态
$A\to B=A/\mathfrak{I}$，则直像 ${}^a\varphi_*(\widetilde{B})$ 典范地
等同于层
$\widetilde{A}/\widetilde{\mathfrak{I}}=\mathscr{O}_X/\mathscr{I}$
（\hyperref[I.1.6.3-fr]{1.6.3} 与
\hyperref[I.1.3.9-fr]{1.3.9}），从而证明完毕。

我们称 $(Y,\mathscr{O}_Y)$ 为 $(X,\mathscr{O}_X)$ 的\emph{子预概形}，
它\emph{由理想层 $\mathscr{I}$ 定义}；这是\emph{子预概形}的一般概念的
一个特例：

\begin{definition}[4.1.3]
\label{I.4.1.3-fr}
若赋环空间 $(Y,\mathscr{O}_Y)$ 满足下列条件，就称它是预概形
$(X,\mathscr{O}_X)$ 的一个\emph{子预概形}：
\begin{enumerate}
  \item[$1^\circ$] $Y$ 是 $X$ 的一个局部闭子空间；
  \item[$2^\circ$] 若 $U$ 表示 $X$ 中包含 $Y$ 且使 $Y$ 在 $U$ 中为闭集的
  最大开集（换言之，它是 $Y$ 关于 $\overline{Y}$ 的边界在 $X$ 中的
  补集），则 $(Y,\mathscr{O}_Y)$ 是 $(U,\mathscr{O}_X|U)$ 的一个由
  $\mathscr{O}_X|U$ 的拟凝聚理想层定义的子预概形。
\end{enumerate}
若 $Y$ 在 $X$ 中是闭的（此时 $U=X$），就称 $(X,\mathscr{O}_X)$ 的
子预概形 $(Y,\mathscr{O}_Y)$ 是\emph{闭的}。
\end{definition}

由此定义和（\hyperref[I.4.1.2-fr]{4.1.2}）立即可知，$X$ 的闭子预概形
与 $\mathscr{O}_X$ 的拟凝聚理想层 $\mathscr{J}$ 之间有一个\emph{典范的
一一对应}。事实上，若两个这样的层 $\mathscr{J}$、$\mathscr{J}'$ 具有
相同的（闭）支集 $Y$，并且 $\mathscr{O}_X/\mathscr{J}$ 与
$\mathscr{O}_X/\mathscr{J}'$ 在 $Y$ 上的限制相同，就立即有
$\mathscr{J}'=\mathscr{J}$。

\begin{env}[4.1.4]
\label{I.4.1.4-fr}
设 $(Y,\mathscr{O}_Y)$ 是 $X$ 的一个子预概形，$U$ 是 $X$ 中包含 $Y$
且使 $Y$ 在其中为闭集的最大开集，而 $V$ 是 $X$ 的一个包含于 $U$ 的
开集；则 $V\cap Y$ 在 $V$ 中是闭的。此外，若 $Y$ 由
$\mathscr{O}_X|U$ 的拟凝聚理想层 $\mathscr{J}$ 定义，则
$\mathscr{J}|V$ 是 $\mathscr{O}_X|V$ 的一个拟凝聚理想层；而且显然，
$Y$ 在 $Y\cap V$ 上诱导的预概形，就是 $V$ 的由理想层
$\mathscr{J}|V$ 定义的闭子预概形。反过来：
\end{env}

\begin{proposition}[4.1.5]
\label{I.4.1.5-fr}
设 $(Y,\mathscr{O}_Y)$ 是一个赋环空间，其中 $Y$ 是 $X$ 的子空间；又设
$Y$ 有一个由 $X$ 的开集构成的覆盖 $(V_\alpha)$，使得对每个 $\alpha$，
$Y\cap V_\alpha$ 在 $V_\alpha$ 中是闭的，并且赋环空间
$(Y\cap V_\alpha,\mathscr{O}_Y|(Y\cap V_\alpha))$ 是 $X$ 在
$V_\alpha$ 上诱导的预概形的一个闭子预概形。则
$(Y,\mathscr{O}_Y)$ 是 $X$ 的一个子预概形。
\end{proposition}

上述假设蕴涵 $Y$ 在 $X$ 中局部闭，而且包含 $Y$ 并使 $Y$ 在其中为闭集
的最大开集 $U$ 包含所有 $V_\alpha$；因此可归结为 $U=X$ 且 $Y$ 在
$X$ 中为闭集的情形。这时定义 $\mathscr{O}_X$ 的一个拟凝聚理想层
$\mathscr{J}$：令 $\mathscr{J}|V_\alpha$ 为 $\mathscr{O}_X|V_\alpha$
中定义闭子预概形
$(Y\cap V_\alpha,\mathscr{O}_Y|(Y\cap V_\alpha))$ 的理想层；而对 $X$
中每个不与 $Y$ 相交的开集 $W$，令
$\mathscr{J}|W=\mathscr{O}_X|W$。由（\hyperref[I.4.1.3-fr]{4.1.3}）
与（\hyperref[I.4.1.4-fr]{4.1.4}）立即验证：满足这些条件的理想层
$\mathscr{J}$ 存在且唯一，并且它定义闭子预概形
$(Y,\mathscr{O}_Y)$。

特别地，$X$ 在 $X$ 的一个\emph{开集}上所\emph{诱导}的预概形，是 $X$
的一个\emph{子预概形}。

\begin{proposition}[4.1.6]
\label{I.4.1.6-fr}
$X$ 的一个子预概形（相应地，闭子预概形）的子\oldpage[I]{121}预概形
（相应地，闭子预概形），典范地等同于 $X$ 的一个子预概形
（相应地，闭子预概形）。
\end{proposition}

$X$ 的局部闭子空间的一个局部闭子集仍是 $X$ 的局部闭子空间；因而显然，
由（\hyperref[I.4.1.5-fr]{4.1.5}）可知问题是局部的，所以可以假设 $X$
为仿射的。于是命题归结为如下典范等同：若 $\mathfrak{J}$、
$\mathfrak{J}'$ 是环 $A$ 的两个理想且
$\mathfrak{J}\subset\mathfrak{J}'$，则 $A/\mathfrak{J}'$ 典范地等同于
$(A/\mathfrak{J})/(\mathfrak{J}'/\mathfrak{J})$。

下文总采用上述等同。

\begin{env}[4.1.7]
\label{I.4.1.7-fr}
设 $Y$ 是预概形 $X$ 的一个子预概形，并以 $\psi$ 表示\emph{底空间}的
典范嵌入 $Y\to X$；已知逆像 $\psi^*(\mathscr{O}_X)$ 就是限制
$\mathscr{O}_X|Y$（\hyperref[0.3.7.1-fr]{0, 3.7.1}）。对每个
$y\in Y$，以 $\omega_y$ 表示典范同态
$(\mathscr{O}_X)_y\to(\mathscr{O}_Y)_y$；这些同态是一个\emph{满的}
环层同态 $\omega$ 在各茎上的限制：
$\mathscr{O}_X|Y\to\mathscr{O}_Y$。事实上，只需
在 $Y$ 上局部验证这一点，也就是说，可以假设 $X$ 为仿射的且子预概形
$Y$ 为闭的；若此时 $\mathscr{I}$ 是 $\mathscr{O}_X$ 中定义 $Y$ 的理想层，
则各 $\omega_y$ 无非是同态
$\mathscr{O}_X|Y\to(\mathscr{O}_X/\mathscr{I})|Y$ 在各茎上的限制。

由此定义了赋环空间的一个\emph{单态射}
（\hyperref[0.4.1.1-fr]{0, 4.1.1}）$j=(\psi,\omega^b)$；它显然是预概形
的一个态射 $Y\to X$（\hyperref[I.2.2.1-fr]{2.2.1}），称为
\emph{典范嵌入态射}。

若 $f:X\to Z$ 是一个态射，我们也称复合态射
$Y\xrightarrow{j}X\xrightarrow{f}Z$ 为 $f$ 在子预概形 $Y$ 上的
\emph{限制}。
\end{env}

\begin{env}[4.1.8]
\label{I.4.1.8-fr}
依照一般定义（T, I, 1.1），设 $Y$ 是 $X$ 的一个子预概形，而
$j:Y\to X$ 是其嵌入态射；若预概形态射 $f:Z\to X$ 可分解为
$Z\xrightarrow{g}Y\xrightarrow{j}X$，其中 $g$ 是预概形态射，就称
$f$ \emph{可经由该嵌入态射分解}。由于 $j$ 是单态射，$g$ 必然\emph{唯一}。
\end{env}

\begin{proposition}[4.1.9]
\label{I.4.1.9-fr}
态射 $f:Z\to X$ \emph{可经嵌入态射 $j:Y\to X$ 分解}的充要条件是：
$f(Z)\subset Y$，并且对每个 $z\in Z$，若令 $y=f(z)$，则与 $f$ 对应的
同态 $(\mathscr{O}_X)_y\to\mathscr{O}_z$ 可分解为
$(\mathscr{O}_X)_y\to(\mathscr{O}_Y)_y\to\mathscr{O}_z$（或者等价地，
$(\mathscr{O}_X)_y\to\mathscr{O}_z$ 的核包含
$(\mathscr{O}_X)_y\to(\mathscr{O}_Y)_y$ 的核）。
\end{proposition}

这些条件显然是必要的。为证明它们充分，必要时把 $X$ 换成一个使 $Y$
在 $U$ 中为闭集的开集 $U$（\hyperref[I.4.1.3-fr]{4.1.3}），即可归结为
$Y$ 是 $X$ 的\emph{闭}子预概形的情形；此时 $Y$ 由
$\mathscr{O}_X$ 的一个拟凝聚理想层 $\mathscr{I}$ 定义。令
$f=(\psi,\theta)$，并以 $\mathscr{J}$ 表示
$\psi^*(\mathscr{O}_X)$ 的理想层，即同态
$\theta^\sharp:\psi^*(\mathscr{O}_X)\to\mathscr{O}_Z$ 的核。考虑到函子
$\psi^*$ 的性质（\hyperref[0.3.7.2-fr]{0, 3.7.2}），上述假设蕴涵对
每个 $z\in Z$ 都有
$(\psi^*(\mathscr{I}))_z\subset\mathscr{J}_z$，从而
$\psi^*(\mathscr{I})\subset\mathscr{J}$。因此，$\theta^\sharp$ 分解为
\[
  \psi^*(\mathscr{O}_X)\longrightarrow
  \psi^*(\mathscr{O}_X)/\psi^*(\mathscr{I})
  =\psi^*(\mathscr{O}_X/\mathscr{I})
  \xrightarrow{\omega}\mathscr{O}_Z,
\]
其中第一个箭头是典范同态。设 $\psi'$ 是与 $\psi$ 相同的连续映射
$Z\to Y$；显然
${\psi'}^*(\mathscr{O}_Y)=\psi^*(\mathscr{O}_X/\mathscr{I})$。另一方面，
$\omega$ 显然是局部同态，故 $g=(\psi',\omega^b)$ 是预概形的一个态射
$Z\to Y$\oldpage[I]{122}（\hyperref[I.2.2.1-fr]{2.2.1}）；由前述构造，
它满足 $f=j\circ g$，命题得证。

\begin{corollary}[4.1.10]
\label{I.4.1.10-fr}
嵌入态射 $Z\to X$ \emph{可经嵌入态射 $Y\to X$ 分解}的充要条件是：$Z$ 是
$Y$ 的一个子预概形。
\end{corollary}

此时写作 $Z\leq Y$；这个关系显然是 $X$ 的子预概形全体上的一个
\emph{序关系}。

\subsection{浸入态射}
\label{subsection:I.4.2-fr}

\begin{definition}[4.2.1]
\label{I.4.2.1-fr}
若态射 $f:Y\to X$ 可分解为
$Y\xrightarrow{g}Z\xrightarrow{j}X$，其中 $g$ 是同构，$Z$ 是 $X$
的一个子预概形（相应地，一个闭子预概形、由某个开集诱导的子预概形），
而 $j$ 是嵌入态射，则称该态射为一个\emph{浸入}（相应地，一个
\emph{闭浸入}、一个\emph{开浸入}）。
\end{definition}

此时，子预概形 $Z$ 和同构 $g$ 以\emph{唯一的}方式确定。事实上，若
$Z'$ 是 $X$ 的另一个子预概形，$j'$ 是嵌入态射 $Z'\to X$，而
$g'$ 是满足 $j\circ g=j'\circ g'$ 的同构 $Y\to Z'$，则可推出
$j'=j\circ g\circ {g'}^{-1}$，因而
$Z'\leq Z$（\hyperref[I.4.1.10-fr]{4.1.10}）；同理可证
$Z\leq Z'$，故 $Z'=Z$；又因 $j$ 是预概形的单态射，所以 $g'=g$。

称 $f=j\circ g$ 为浸入 $f$ 的\emph{典范分解}，并称子预概形 $Z$
和同构 $g$ \emph{与 $f$ 相伴}。

显然，浸入是预概形的\emph{单态射}
（\hyperref[I.4.1.7-fr]{4.1.7}），\emph{从而更是}一个\emph{根态射}
（\hyperref[I.3.5.4-fr]{3.5.4}）。

\begin{proposition}[4.2.2]
\label{I.4.2.2-fr}
\begin{enumerate}
  \item[a)] 态射 $f=(\psi,\theta):Y\to X$ 是\emph{开浸入}的充要条件
  是：$\psi$ 是从 $Y$ 到 $X$ 的一个开子集的同胚，并且对于每个
  $y\in Y$，同态
  $\theta_y^\sharp:\mathscr{O}_{\psi(y)}\to\mathscr{O}_y$ 是双射。
  \item[b)] 态射 $f=(\psi,\theta):Y\to X$ 是\emph{浸入}
  （相应地，\emph{闭浸入}）的充要条件是：$\psi$ 是从 $Y$ 到 $X$
  的一个局部闭子集（相应地，闭子集）的同胚，并且对于每个
  $y\in Y$，同态
  $\theta_y^\sharp:\mathscr{O}_{\psi(y)}\to\mathscr{O}_y$ 是满射。
\end{enumerate}
\end{proposition}

\begin{enumerate}
  \item[a)] 这些条件显然是必要的。反过来，若它们成立，则显然
  % EGA1-ERR-0032: source-literal reversed sheaf order retained; see ERRATA_CANDIDATES.jsonl.
  $\theta^\sharp$ 是从 $\mathscr{O}_Y$ 到
  $\psi^*(\mathscr{O}_X)$ 的同构，而
  $\psi^*(\mathscr{O}_X)$ 是从 $\mathscr{O}_X|\psi(Y)$ 出发、借助
  $\psi^{-1}$ 作结构传递所得的层；结论由此得出。
  \item[b)] 这些条件显然是必要的；现证明它们是充分的。先考虑如下
  特殊情形：假设 $X$ 是仿射概形，并且 $Z=\psi(Y)$ 在 $X$ 中是
  \emph{闭的}。已知（\hyperref[0.3.4.6-fr]{0, 3.4.6}），此时
  $\psi_*(\mathscr{O}_Y)$ 的支集为 $Z$；若以
  $\mathscr{O}'_Z$ 表示它在 $Z$ 上的限制，则赋环空间
  $(Z,\mathscr{O}'_Z)$ 是由 $(Y,\mathscr{O}_Y)$ 借助同胚 $\psi$
  作结构传递得到的，其中把该同胚看作从 $Y$ 到 $Z$ 的映射。
  下面证明 $f_*(\mathscr{O}_Y)=\psi_*(\mathscr{O}_Y)$ 是一个\emph{拟凝聚}
  $\mathscr{O}_X$-模。事实上，对于每个
  $x\notin Z$，$\psi_*(\mathscr{O}_Y)$ 在 $x$ 的某个适当邻域上的
  限制为零。反之，若 $x\in Z$，则存在唯一确定的 $y\in Y$ 使
  $x=\psi(y)$；设 $V$ 为 $Y$ 中 $y$ 的一个仿射开邻域；
  $\psi(V)$ 是 $Z$ 中的开集，因而是 $X$ 的某个开集 $U$ 在 $Z$
  上的迹；$\psi_*(\mathscr{O}_Y)$ 在 $U$ 上的限制，等同于下述直
\oldpage[I]{123}
  像 $(\psi_V)_*(\mathscr{O}_Y|V)$ 在 $U$ 上的限制，其中
  $\psi_V$ 是 $\psi$ 在 $V$ 上的限制。态射 $(\psi,\theta)$ 在
  $(V,\mathscr{O}_Y|V)$ 上的限制，是从这个预概形到
  $(X,\mathscr{O}_X)$ 的态射，因而具有形式
  $({}^a\varphi,\widetilde{\varphi})$，其中 $\varphi$ 是从环
  $A=\Gamma(X,\mathscr{O}_X)$ 到环
  $\Gamma(V,\mathscr{O}_Y)$ 的同态
  （\hyperref[I.1.7.3-fr]{1.7.3}）；由此可知，
  $(\psi_V)_*(\mathscr{O}_Y|V)$ 是一个拟凝聚
  $\mathscr{O}_X$-模（\hyperref[I.1.6.3-fr]{1.6.3}）；由于拟凝聚层
  的性质是局部的，这就证明了我们的断言。此外，$\psi$ 是同胚这一
  假设蕴含（\hyperref[0.3.4.5-fr]{0, 3.4.5}）：对于每个 $y\in Y$，
  $\psi_y$ 是同构
  $(\psi_*(\mathscr{O}_Y))_{\psi(y)}\to\mathscr{O}_y$；由于图表
  \phantomsection\label{I.4.2.2.theta-diagram-fr}
  \[
    \xymatrix{
      \mathscr{O}_{\psi(y)}\ar[r]^{\theta_{\psi(y)}}
        \ar[d]_{\psi_y\circ\rho_{\psi(y)}} &
      (\psi_*(\mathscr{O}_Y))_{\psi(y)}\ar[d]^{\psi_y}\\
      (\psi^*(\mathscr{O}_X))_y\ar[r]^{\theta_y^\sharp} &
      \mathscr{O}_y
    }
  \]
  是交换的，并且两个竖直箭头都是同构
  （\hyperref[0.3.7.2-fr]{0, 3.7.2}），$\theta_y^\sharp$ 是满射这一
  假设蕴含 $\theta_{\psi(y)}$ 也是满射。由于
  $\psi_*(\mathscr{O}_Y)$ 的支集是 $Z=\psi(Y)$，$\theta$ 是从
  $\mathscr{O}_X=\widetilde{A}$ 到拟凝聚 $\mathscr{O}_X$-模
  $f_*(\mathscr{O}_Y)$ 的一个\emph{满同态}。因此，存在从某个商层
  $\widetilde{A}/\widetilde{\mathfrak{J}}$（其中 $\mathfrak{J}$
  是 $A$ 的理想）到 $f_*(\mathscr{O}_Y)$ 的唯一同构 $\omega$；
  它与典范同态
  $\widetilde{A}\to\widetilde{A}/\widetilde{\mathfrak{J}}$ 的复合
  等于 $\theta$（\hyperref[I.1.3.8-fr]{1.3.8}）。若
  $\mathscr{O}_Z$ 表示
  $\widetilde{A}/\widetilde{\mathfrak{J}}$ 在 $Z$ 上的限制，则
  $(Z,\mathscr{O}_Z)$ 是 $(X,\mathscr{O}_X)$ 的一个子预概形，而
  $f$ 分解为该子预概形到 $X$ 的典范嵌入态射与同构
  $(\psi_0,\omega_0)$ 的复合；其中 $\psi_0$ 是把 $\psi$ 看作从
  $Y$ 到 $Z$ 的映射，$\omega_0$ 是 $\omega$ 在
  $\mathscr{O}_Z$ 上的限制。

  现在转入一般情形。设 $U$ 为 $X$ 的一个仿射开集，使得
  $U\cap\psi(Y)$ 在 $U$ 中闭且非空。把 $f$ 限制到由 $Y$ 在开集
  $\psi^{-1}(U)$ 上诱导的预概形，并把它视为从这个预概形到由 $X$
  在 $U$ 上诱导的预概形的态射，就归结到第一种情形；所以 $f$ 在
  $\psi^{-1}(U)$ 上的限制是闭浸入
  $\psi^{-1}(U)\to U$，其典范分解为 $j_U\circ g_U$，其中 $g_U$
  是从预概形 $\psi^{-1}(U)$ 到 $U$ 的一个子预概形 $Z_U$ 的同构，
  $j_U$ 是典范嵌入态射 $Z_U\to U$。设 $V$ 是 $X$ 的另一个仿射
  开集，且 $V\subset U$；由于 $Z_U$ 在 $V$ 上的限制 $Z'_V$ 是
  预概形 $V$ 的子预概形，$f$ 在 $\psi^{-1}(V)$ 上的限制分解为
  $j'_V\circ g'_V$，其中 $j'_V$ 是典范嵌入态射 $Z'_V\to V$，
  $g'_V$ 是从 $\psi^{-1}(V)$ 到 $Z'_V$ 的同构。由浸入的典范分解
  的唯一性（\hyperref[I.4.2.1-fr]{4.2.1}），必有
  $Z'_V=Z_V$ 且 $g'_V=g_V$。由此可知
  （\hyperref[I.4.1.5-fr]{4.1.5}），存在 $X$ 的一个子预概形 $Z$，
  其底空间为 $\psi(Y)$，并且它在每个 $U\cap\psi(Y)$ 上的限制都是
  $Z_U$；此时，各 $g_U$ 是某个同构 $g:Y\to Z$ 在
  $\psi^{-1}(U)$ 上的限制，而且 $f=j\circ g$，其中 $j$ 是典范
  嵌入态射 $Z\to X$。
\end{enumerate}

\begin{corollary}[4.2.3]
\label{I.4.2.3-fr}
设 $X$ 为仿射概形。态射
$f=(\psi,\theta):Y\to X$ 是闭浸入的充要条件是：$Y$ 是仿射概形，
并且同态
% EGA1-ERR-0033: source-literal Gamma(psi) retained; see ERRATA_CANDIDATES.jsonl.
$\Gamma(\psi):\Gamma(\mathscr{O}_X)\to\Gamma(\mathscr{O}_Y)$ 是满射。
\end{corollary}

\begin{corollary}[4.2.4]
\label{I.4.2.4-fr}
\begin{enumerate}
  \item[a)] 设 $f$ 为态射 $Y\to X$，$(V_\lambda)$ 为由 $X$ 的开集
  组成的 $f(Y)$ 的一个覆盖。$f$ 是浸入（相应地，开浸入）的充要条件
\oldpage[I]{124}
  是：它在每个诱导预概形 $f^{-1}(V_\lambda)$ 上的限制，是到
  $V_\lambda$ 中的浸入（相应地，开浸入）。
  \item[b)] 设 $f$ 为态射 $Y\to X$，$(V_\lambda)$ 为 $X$ 的一个
  开覆盖。$f$ 是闭浸入的充要条件是：它在每个诱导预概形
  $f^{-1}(V_\lambda)$ 上的限制，是到 $V_\lambda$ 中的闭浸入。
\end{enumerate}
\end{corollary}

设 $f=(\psi,\theta)$；在情形 a) 中，对于每个 $y\in Y$，
$\theta_y^\sharp$ 是满射（相应地，双射）；在情形 b) 中，对于每个
$y\in Y$，$\theta_y^\sharp$ 是满射。因此，只须验证：在情形 a)
中，$\psi$ 是从 $Y$ 到 $X$ 的一个局部闭子集（相应地，开子集）的
同胚；而在情形 b) 中，该映射是从 $Y$ 到 $X$ 的一个闭子集的同胚。
根据假设，$\psi$ 显然是单射，并且对于每个 $y\in Y$，它把 $Y$ 中
$y$ 的每个邻域映为 $\psi(Y)$ 中 $\psi(y)$ 的一个邻域。在情形 a)
中，$\psi(Y)\cap V_\lambda$ 在 $V_\lambda$ 中是局部闭的
（相应地，开的），因而 $\psi(Y)$ 在所有 $V_\lambda$ 的并中是局部
闭的（相应地，开的），从而在 $X$ 中\emph{更是}如此；在情形 b) 中，
$\psi(Y)\cap V_\lambda$ 在 $V_\lambda$ 中是闭的，因而由于
$X=\bigcup_\lambda V_\lambda$，$\psi(Y)$ 在 $X$ 中是闭的。

\begin{proposition}[4.2.5]
\label{I.4.2.5-fr}
两个浸入（相应地，两个开浸入、两个闭浸入）的复合仍是浸入
（相应地，开浸入、闭浸入）。
\end{proposition}

这由（\hyperref[I.4.1.6-fr]{4.1.6}）立即得出。

\subsection{浸入的积}
\label{subsection:I.4.3-fr}

\begin{proposition}[4.3.1]
\label{I.4.3.1-fr}
设 $\alpha:X'\to X$、$\beta:Y'\to Y$ 为两个 $S$-态射；若 $\alpha$
和 $\beta$ 是浸入（相应地，开浸入、闭浸入），则
$\alpha\times_S\beta$ 是浸入（相应地，开浸入、闭浸入）。此外，若
$\alpha$（相应地，$\beta$）把 $X'$（相应地，$Y'$）等同于一个
子预概形 $X''$（相应地，$Y''$），它是 $X$（相应地，$Y$）的
子预概形，则
$\alpha\times_S\beta$ 把 $X'\times_S Y'$ 的底空间等同于子空间
$p^{-1}(X'')\cap q^{-1}(Y'')$，后者是 $X\times_S Y$ 的底空间的
子空间；这里 $p$ 和 $q$ 分别表示从
$X\times_S Y$ 到 $X$ 和 $Y$ 的投影。
\end{proposition}

根据定义（\hyperref[I.4.2.1-fr]{4.2.1}），可以只考虑 $X'$ 和 $Y'$
是子预概形、$\alpha$ 和 $\beta$ 是嵌入态射的情形。对于在开集上诱导
的子预概形，命题已经得到证明（\hyperref[I.3.2.7-fr]{3.2.7}）；由于
每个子预概形都是由某个开集诱导的预概形的闭子预概形
（\hyperref[I.4.1.3-fr]{4.1.3}），所以归结到 $X'$ 和 $Y'$ 是
\emph{闭}子预概形的情形。

首先证明可以假设 $S$ 是\emph{仿射的}。事实上，设 $(S_\lambda)$ 为
由仿射开集组成的 $S$ 的一个覆盖；若 $\varphi$ 和 $\psi$ 分别是
$X$ 和 $Y$ 的结构态射，令 $X_\lambda=\varphi^{-1}(S_\lambda)$、
$Y_\lambda=\psi^{-1}(S_\lambda)$。$X'_\lambda$（相应地，
$Y'_\lambda$）是 $X'$（相应地，$Y'$）在
$X_\lambda\cap X'$（相应地，$Y_\lambda\cap Y'$）上的限制，并且是
$X_\lambda$
（相应地，$Y_\lambda$）的闭子预概形；预概形
$X_\lambda$、$Y_\lambda$、$X'_\lambda$、$Y'_\lambda$ 都可视为
$S_\lambda$-预概形，并且积 $X_\lambda\times_S Y_\lambda$ 与
$X_\lambda\times_{S_\lambda}Y_\lambda$（相应地，
$X'_\lambda\times_S Y'_\lambda$ 与
$X'_\lambda\times_{S_\lambda}Y'_\lambda$）相同
（\hyperref[I.3.2.5-fr]{3.2.5}）。若命题在 $S$ 仿射时成立，则
$\alpha\times_S\beta$ 在每个 $X'_\lambda\times_S Y'_\lambda$
上的限制都是浸入（\hyperref[I.3.2.7-fr]{3.2.7}）。由于积
$X'_\lambda\times_S Y'_\mu$（相应地，
$X_\lambda\times_S Y_\mu$）等同于
$(X'_\lambda\cap X'_\mu)\times_S(Y'_\lambda\cap Y'_\mu)$
（相应地，$(X_\lambda\cap X_\mu)\times_S(Y_\lambda\cap Y_\mu)$）
（\hyperref[I.3.2.6.4-fr]{3.2.6.4}），$\alpha\times_S\beta$
\oldpage[I]{125}
在每个 $X'_\lambda\times_S Y'_\mu$ 上的限制仍为浸入；因此，根据
（\hyperref[I.4.2.4-fr]{4.2.4}），$\alpha\times_S\beta$ 本身也是
浸入。

其次证明还可以假设 $X$ 和 $Y$ 是\emph{仿射的}。事实上，设
$(U_i)$（相应地，$(V_j)$）为由仿射开集组成的 $X$（相应地，$Y$）
的一个覆盖，并设 $X'_i$（相应地，$Y'_j$）为 $X'$（相应地，$Y'$）
在 $X'\cap U_i$（相应地，$Y'\cap V_j$）上的限制；它是 $U_i$
（相应地，$V_j$）的闭子预概形。$U_i\times_S V_j$ 等同于
$X\times_S Y$ 在 $p^{-1}(U_i)\cap q^{-1}(V_j)$ 上的限制
（\hyperref[I.3.2.7-fr]{3.2.7}）；同样，若 $p'$、$q'$ 是
$X'\times_S Y'$ 的投影，则 $X'_i\times_S Y'_j$ 等同于
$X'\times_S Y'$ 在
${p'}^{-1}(X'_i)\cap{q'}^{-1}(Y'_j)$ 上的限制。令
$\gamma=\alpha\times_S\beta$；按照定义有
$p\circ\gamma=\alpha\circ p'$、$q\circ\gamma=\beta\circ q'$；由于
$X'_i=\alpha^{-1}(U_i)$、$Y'_j=\beta^{-1}(V_j)$，还得到
${p'}^{-1}(X'_i)=\gamma^{-1}(p^{-1}(U_i))$、
${q'}^{-1}(Y'_j)=\gamma^{-1}(q^{-1}(V_j))$，从而
\[
  {p'}^{-1}(X'_i)\cap{q'}^{-1}(Y'_j)
  =\gamma^{-1}\bigl(p^{-1}(U_i)\cap q^{-1}(V_j)\bigr)
  =\gamma^{-1}(U_i\times_S V_j)
\]
于是可按论证的第一部分得出结论。

现在假设 $X$、$Y$、$S$ 都是仿射的，并设 $B$、$C$、$A$ 分别为
它们的环。于是 $B$ 和 $C$ 是 $A$-代数，而 $X'$ 和 $Y'$ 是仿射
概形，其环分别为商代数 $B'$、$C'$；二者分别是 $B$ 和 $C$ 的商。
此外，有
$\alpha=({}^a\rho,\widetilde{\rho})$、
$\beta=({}^a\sigma,\widetilde{\sigma})$，其中 $\rho$ 和 $\sigma$
分别是典范同态 $B\to B'$、$C\to C'$
（\hyperref[I.1.7.3-fr]{1.7.3}）。由此已知，$X\times_S Y$
（相应地，$X'\times_S Y'$）是环为 $B\otimes_A C$
（相应地，$B'\otimes_A C'$）的仿射概形，并且
$\alpha\times_S\beta=({}^a\tau,\widetilde{\tau})$，其中 $\tau$ 是
同态 $\rho\otimes\sigma$，它从 $B\otimes_A C$ 映到
$B'\otimes_A C'$
（\hyperref[I.3.2.2-fr]{3.2.2} 和
\hyperref[I.3.2.3-fr]{3.2.3}）；由于这个同态是满射，
$\alpha\times_S\beta$ 确为浸入。此外，若 $\mathfrak{b}$
（相应地，$\mathfrak{c}$）为 $\rho$（相应地，$\sigma$）的核，则
% EGA1-ERR-0034: source-literal generated-ideal ambiguity retained; see ERRATA_CANDIDATES.jsonl.
$\tau$ 的核是 $u(\mathfrak{b})+v(\mathfrak{c})$，其中 $u$
（相应地，$v$）是同态 $b\mapsto b\otimes1$
（相应地，$c\mapsto1\otimes c$）。由于
$p=({}^a u,\widetilde{u})$ 且 $q=({}^a v,\widetilde{v})$，根据
（\hyperref[I.1.2.2.1-fr]{1.2.2.1} 和
\hyperref[I.1.1.2-fr]{1.1.2 (iii)}），这个核在
$B\otimes_A C$ 的素谱中对应于闭集
$p^{-1}(X')\cap q^{-1}(Y')$；证明至此完成。

\begin{corollary}[4.3.2]
\label{I.4.3.2-fr}
若 $f:X\to Y$ 是浸入（相应地，开浸入、闭浸入）并且是
$S$-态射，则 $f_{(S')}$ 都是浸入（相应地，开浸入、闭浸入），
其中 $S'\to S$ 是基预概形的任意扩张。
\end{corollary}

\subsection{子预概形的逆像}
\label{subsection:I.4.4-fr}

\begin{proposition}[4.4.1]
\label{I.4.4.1-fr}
设 $f:X\to Y$ 为态射，$Y'$ 为 $Y$ 的一个子预概形（相应地，一个
闭子预概形、在某个开集上诱导的预概形），$j:Y'\to Y$ 为嵌入态射。
则投影 $p:X\times_Y Y'\to X$ 是浸入（相应地，闭浸入、开浸入）；
与 $p$ 相联系的 $X$ 的子预概形以 $f^{-1}(Y')$ 为底空间；此外，
若 $j'$ 是这个子预概形到 $X$ 中的嵌入态射，则态射 $h:Z\to X$
使 $f\circ h:Z\to Y$ 可经由 $j$ 分解的充要条件是：$h$ 可经由
$j'$ 分解。
\end{proposition}

由于 $p=1_X\times_Y j$（\hyperref[I.3.3.4-fr]{3.3.4}），第一个断言
由（\hyperref[I.4.3.1-fr]{4.3.1}）得出；第二个断言是
（\hyperref[I.3.5.10-fr]{3.5.10}）的一个特殊情形（其中需要交换
$X$ 与 $Y'$ 的角色）。最后，若有 $f\circ h=j\circ h'$，其中 $h'$
是态射 $Z\to Y'$，则由纤维积的定义可知 $h=p\circ u$，其中 $u$
是态射 $Z\to X\times_Y Y'$；最后一个断言由此得出。

我们称这样定义的 $X$ 的子预概形为 $Y$ 的子预概形 $Y'$ 在态射
$f$ 下的\emph{逆像}；这一术语与
\oldpage[I]{126}
在（\hyperref[I.3.3.6-fr]{3.3.6}）中以更一般方式引入的术语一致。
当我们把 $f^{-1}(Y')$ 作为 $X$ 的一个子预概形来谈论时，所指的
总是这个子预概形。

当预概形 $f^{-1}(Y')$ 与 $X$ 相同时，$j'$ 是恒等态射，因而每个
态射 $h:Z\to X$ 都可经由 $j'$ 分解；所以
\emph{态射 $f:X\to Y$ 此时可分解为}
$X\xrightarrow{g}Y'\xrightarrow{j}Y$。

当 $y$ 是 $Y$ 的一个\emph{闭}点，并且
$Y'=\operatorname{Spec}(k(y))$ 是以 $\{y\}$ 为底空间的 $Y$ 的
最小闭子预概形（\hyperref[I.4.1.9-fr]{4.1.9}）时，闭子预概形
$f^{-1}(Y')$ 典范同构于在（\hyperref[I.3.6.2-fr]{3.6.2}）中定义的
\emph{纤维} $f^{-1}(y)$；今后将二者等同。

\begin{corollary}[4.4.2]
\label{I.4.4.2-fr}
设 $f:X\to Y$、$g:Y\to Z$ 为两个态射，$h=g\circ f$ 为它们的复合。
对于 $Z$ 的任意子预概形 $Z'$，$X$ 的子预概形
$f^{-1}(g^{-1}(Z'))$ 与 $h^{-1}(Z')$ 相同。
\end{corollary}

这是因为存在典范同构
$X\times_Y(Y\times_Z Z')\xrightarrow{\sim}X\times_Z Z'$
（\hyperref[I.3.3.9.1-fr]{3.3.9.1}）。

\begin{corollary}[4.4.3]
\label{I.4.4.3-fr}
设 $X'$、$X''$ 为 $X$ 的两个子预概形，$j':X'\to X$、
$j'':X''\to X$ 为嵌入态射；则 ${j'}^{-1}(X'')$ 与
${j''}^{-1}(X')$ 都等于 $X'$ 与 $X''$ 在子预概形的序关系下的
下确界 $\inf(X',X'')$，并且该下确界典范同构于 $X'\times_X X''$。
\end{corollary}

这立即由（\hyperref[I.4.4.1-fr]{4.4.1}）和
（\hyperref[I.4.1.10-fr]{4.1.10}）得出。

\begin{corollary}[4.4.4]
\label{I.4.4.4-fr}
设 $f:X\to Y$ 为态射，$Y'$、$Y''$ 为 $Y$ 的两个子预概形；则
\[
  f^{-1}(\inf(Y',Y''))=\inf(f^{-1}(Y'),f^{-1}(Y'')).
\]
\end{corollary}

这是因为 $(X\times_Y Y')\times_X(X\times_Y Y'')$ 与
$X\times_Y(Y'\times_Y Y'')$ 之间存在典范同构
（\hyperref[I.3.3.9.1-fr]{3.3.9.1}）。

\begin{proposition}[4.4.5]
\label{I.4.4.5-fr}
设 $f:X\to Y$ 为态射，$Y'$ 为 $Y$ 的一个闭子预概形，它由
$\mathscr{O}_Y$ 的一个拟凝聚理想层 $\mathscr{K}$ 定义
（\hyperref[I.4.1.3-fr]{4.1.3}）；则 $X$ 的闭子预概形
$f^{-1}(Y')$ 由 $\mathscr{O}_X$ 的拟凝聚理想层
$f^*(\mathscr{K})\mathscr{O}_X$ 定义。
\end{proposition}

这个问题在 $X$ 和 $Y$ 上显然都是局部的；因而只须注意，
% EGA1-ERR-0035: source-literal A/B-algebra reversal retained; see ERRATA_CANDIDATES.jsonl.
若 $B$ 是一个 $A$-代数且 $\mathfrak{K}$ 是 $B$ 的一个理想，则
$A\otimes_B(B/\mathfrak{K})=A/\mathfrak{K}A$，再应用
（\hyperref[I.1.6.9-fr]{1.6.9}）即可。

\begin{corollary}[4.4.6]
\label{I.4.4.6-fr}
设 $X'$ 为 $X$ 的一个闭子预概形，它由 $\mathscr{O}_X$ 的拟凝聚
理想层 $\mathscr{J}$ 定义，$i$ 为嵌入态射 $X'\to X$；$f$ 在 $X'$
上的限制 $f\circ i$ 可经由嵌入态射 $j:Y'\to Y$ 分解（换言之，
可分解为 $j\circ g$，其中 $g$ 是态射 $X'\to Y'$）的充要条件是
$f^*(\mathscr{K})\mathscr{O}_X\subset\mathscr{J}$。
\end{corollary}

只须对 $i$ 应用命题（\hyperref[I.4.4.1-fr]{4.4.1}），并计及
（\hyperref[I.4.4.5-fr]{4.4.5}）即可。

\subsection{局部浸入与局部同构}
\label{subsection:I.4.5-fr}

\begin{definition}[4.5.1]
\label{I.4.5.1-fr}
设 $f:X\to Y$ 为预概形的一个态射。若存在 $X$ 中点 $x$ 的一个开邻域
$U$ 和 $Y$ 中 $f(x)$ 的一个开邻域 $V$，使得 $f$ 在诱导预概形 $U$
上的限制是从 $U$ 到诱导预概形 $V$ 的一个闭浸入，则称 $f$
\emph{在点 $x\in X$ 是局部浸入}。若 $f$ 在 $X$ 的每一点都是
局部浸入，则称 $f$ 是一个局部浸入。
\end{definition}

\begin{definition}[4.5.2]
\label{I.4.5.2-fr}
若存在 $X$ 中点 $x$ 的一个开邻域 $U$，使得态射 $f:X\to Y$ 在
\oldpage[I]{127}
诱导预概形 $U$ 上的限制是从 $U$ 到 $Y$ 的一个开浸入，则称 $f$
在点 $x\in X$ 是局部同构。若 $f$ 在 $X$ 的每一点都是
局部同构，则称 $f$ 是一个局部同构。
\end{definition}

\begin{env}[4.5.3]
\label{I.4.5.3-fr}
一个浸入（相应地，闭浸入）$f:X\to Y$ 因而可以刻画为满足下述条件
的局部浸入：$f$ 把 $X$ 的底空间同胚地映到 $Y$ 的一个子集
（相应地，闭子集）上。一个开浸入 $f$ 可以刻画为一个\emph{单射的}
局部同构。
\end{env}

\begin{proposition}[4.5.4]
\label{I.4.5.4-fr}
设 $X$ 是不可约预概形，$f:X\to Y$ 是一个单射的支配态射。若 $f$
是局部浸入，则 $f$ 是浸入，并且 $f(X)$ 在 $Y$ 中是开的。
\end{proposition}

事实上，设 $x\in X$，并设 $U$ 是 $x$ 的一个开邻域，$V$ 是 $Y$ 中
$f(x)$ 的一个开邻域，使得 $f$ 在 $U$ 上的限制是到 $V$ 中的闭浸入；
由于 $U$ 在 $X$ 中稠密，根据假设 $f(U)$ 在 $Y$ 中稠密，故
$f(U)=V$，并且 $f$ 是从 $U$ 到 $V$ 的同胚；$f$ 为单射这一假设蕴涵
$f^{-1}(V)=U$，命题立即得证。

\begin{proposition}[4.5.5]
\label{I.4.5.5-fr}
\begin{enumerate}
  \item[\rm(i)] 两个局部浸入（相应地，两个局部同构）的复合是局部
    浸入（相应地，局部同构）。
  \item[\rm(ii)] 设 $f:X\to X'$、$g:Y\to Y'$ 为两个 $S$-态射。
    若 $f$ 和 $g$ 是局部浸入（相应地，局部同构），则
    $f\times_S g$ 也是如此。
  \item[\rm(iii)] 若一个 $S$-态射 $f$ 是局部浸入（相应地，局部
    同构），则对于基预概形的任意扩张 $S'\to S$，$f_{(S')}$
    也是如此。
\end{enumerate}
\end{proposition}

根据（\hyperref[I.3.5.1-fr]{3.5.1}），只须证明 (i) 和 (ii)。

% EGA1-ERR-0036: source-literal cross-reference 4.2.4 retained; see ERRATA_CANDIDATES.jsonl.
(i) 立即由闭浸入（相应地，开浸入）的传递性
（\hyperref[I.4.2.4-fr]{4.2.4}）以及下述事实得出：若 $f$ 是从 $X$
到 $Y$ 的一个闭子集的同胚，则对于任意开集 $U\subset X$，$f(U)$
在 $f(X)$ 中是开的，故存在 $Y$ 的一个开子集 $V$ 使
$f(U)=V\cap f(X)$，从而 $f(U)$ 在 $V$ 中是闭的。

% EGA1-ERR-0037: source-literal unintroduced z,z' retained; see ERRATA_CANDIDATES.jsonl.
为证明 (ii)，设 $p$、$q$ 为 $X\times_S Y$ 的投影，$p'$、$q'$ 为
$X'\times_S Y'$ 的投影。根据假设，分别存在
$x=p(z)$、$x'=p'(z')$、$y=q(z)$、$y'=q'(z')$ 的开邻域
$U$、$U'$、$V$、$V'$，使得 $f$ 和 $g$ 分别在 $U$ 和 $V$ 上的限制
是到 $U'$ 和 $V'$ 中的闭浸入（相应地，开浸入）。由于
$U\times_S V$ 和 $U'\times_S V'$ 的底空间分别等同于 $z$ 和 $z'$
的开邻域 $p^{-1}(U)\cap q^{-1}(V)$ 和
${p'}^{-1}(U')\cap{q'}^{-1}(V')$
（\hyperref[I.3.2.7-fr]{3.2.7}），命题由
（\hyperref[I.4.3.1-fr]{4.3.1}）得出。
% EGA I mainland Simplified-Chinese translation (producer output; uncertified).
% Sequential translation cursor: through 5.1.2 (source lines 1--63; printed pages 127--128).

\section{约化预概形；分离条件}
\label{section:I.5-fr}

\subsection{约化预概形}
\label{subsection:I.5.1-fr}

\begin{proposition}[5.1.1]
\label{I.5.1.1-fr}
设 $(X,\mathscr{O}_X)$ 是一个预概形，$\mathscr{B}$ 是一个拟凝聚
$\mathscr{O}_X$-代数。存在唯一的拟凝聚 $\mathscr{O}_X$-模
$\mathscr{N}$，使得对每个 $x\in X$，茎 $\mathscr{N}_x$ 都是环
$\mathscr{B}_x$ 的幂零根。当 $X$ 是仿射的，因而
$\mathscr{B}=\widetilde{B}$，其中 $B$ 是 $A(X)$ 上的一个代数时，有
$\mathscr{N}=\widetilde{\mathfrak{N}}$，其中 $\mathfrak{N}$ 是 $B$ 的
幂零根。
\end{proposition}

\oldpage[I]{128}
由于问题是局部的，归结为证明最后一个断言。已知
$\widetilde{\mathfrak{N}}$ 是一个拟凝聚 $\mathscr{O}_X$-模
（\hyperref[I.1.4.1-fr]{1.4.1}），并且它在点 $x\in X$ 处的茎是
分式环 $B_x$ 的理想 $\mathfrak{N}_x$；因而只须证明 $B_x$ 的幂零根
包含在 $\mathfrak{N}_x$ 中，反向包含关系是显然的。事实上，设
$z/s$ 是 $B_x$ 的幂零根的一个元素，其中 $z\in B$、
$s\notin\mathfrak{j}_x$；依假设，存在整数 $k$ 使得
$(z/s)^k=0$，这意味着存在 $t\notin\mathfrak{j}_x$ 使得
$tz^k=0$。由此得到 $(tz)^k=0$，因而
$z/s=(tz)/(ts)$ 属于 $\mathfrak{N}_x$。

我们称这样定义的拟凝聚 $\mathscr{O}_X$-模 $\mathscr{N}$ 为
$\mathscr{O}_X$-代数 $\mathscr{B}$ 的\emph{幂零根}；特别地，以
$\mathscr{N}_X$ 表示 $\mathscr{O}_X$ 的幂零根。

\begin{corollary}[5.1.2]
\label{I.5.1.2-fr}
设 $X$ 是一个预概形；由理想层 $\mathscr{N}_X$ 定义的 $X$ 的闭子预概形，
是以 $X$ 为底空间的 $X$ 的唯一约化子预概形
（\hyperref[0.4.1.4-fr]{0, 4.1.4}）；它也是以 $X$ 为底空间的 $X$ 的最小子预概形。
\end{corollary}

由 $\mathscr{N}_X$ 定义的闭子预概形 $Y$ 的结构层是
$\mathscr{O}_X/\mathscr{N}_X$；由于对每个 $x\in X$ 都有
$\mathscr{N}_x\ne\mathscr{O}_x$，立即可知 $Y$ 是约化的，并且以
$X$ 为底空间。为证明其余断言，注意，以 $X$ 为底空间的 $X$ 的
一个子预概形 $Z$，由一个理想层 $\mathscr{J}$
（\hyperref[I.4.1.3-fr]{4.1.3}）定义，并且对每个 $x\in X$ 都有
$\mathscr{J}_x\ne\mathscr{O}_x$。只须考虑 $X$ 是仿射的情形；设
$X=\operatorname{Spec}(A)$，且
$\mathscr{J}=\widetilde{\mathfrak{J}}$，其中 $\mathfrak{J}$ 是 $A$ 的
一个理想。于是对每个 $x\in X$，有
$\mathfrak{J}\subset\mathfrak{j}_x$，所以 $\mathfrak{J}$ 包含在
$A$ 的所有素理想中，也就是说，包含在它们的交 $\mathfrak{N}$（即 $A$
的幂零根）中。这证明了 $Y$ 是以 $X$ 为底空间的 $X$ 的最小子预概形
（\hyperref[I.4.1.9-fr]{4.1.9}）；此外，若 $Z$ 不同于
$Y$，则必至少存在一个 $x\in X$，使得
$\mathscr{J}_x\ne\mathscr{N}_x$，因而由
（\hyperref[I.5.1.1-fr]{5.1.1}）可知 $Z$ 不是约化的。

\begin{definition}[5.1.3]
\label{I.5.1.3-fr}
称以 $X$ 为底空间的 $X$ 的唯一约化子预概形为与预概形 $X$ 相伴的
约化预概形，并记为 $X_{\mathrm{red}}$。
\end{definition}

因此，说预概形 $X$ 是约化的，就是说
$X=X_{\mathrm{red}}$。

% EGA1-ERR-0038: retain printed 2.1.7; the integral-prescheme definition is 2.1.8.
\begin{proposition}[5.1.4]
\label{I.5.1.4-fr}
环 $A$ 的素谱为约化（相应地，整）预概形
（\hyperref[I.2.1.7-fr]{2.1.7}）的充要条件是，$A$ 为约化环
（相应地，整环）。
\end{proposition}

事实上，由（\hyperref[I.5.1.1-fr]{5.1.1}）立即可知，条件
$\mathscr{N}=(0)$ 是 $X=\operatorname{Spec}(A)$ 为约化预概形的充要条件；
关于整环的断言于是由（\hyperref[I.1.1.13-fr]{1.1.13}）得出。

由于整环的每个分式环 $\ne\{0\}$ 都是整环，由
（\hyperref[I.5.1.4-fr]{5.1.4}）可知，对每个\emph{局部整}预概形
$X$，$\mathscr{O}_x$ 对每个 $x\in X$ 都是\emph{整环}。当 $X$ 的
底空间\emph{局部 Noether} 时，逆命题成立：事实上，此时 $X$ 是约化的；
若 $U$ 是 $X$ 的一个同时为 Noether 空间的仿射开集，则 $U$ 只有有限个
不可约分支，因而它的环 $A$ 只有有限个极小素理想
（\hyperref[I.1.1.14-fr]{1.1.14}）。若这些分支中的两个 $U_i$ 有一个
公共点 $x$，则 $\mathscr{O}_x$ 至少有两个不同的极小素理想，因而不是
整环；所以各 $U_i$ 是两两不交的开集，并且每一个都是整的。

\begin{env}[5.1.5]
\label{I.5.1.5-fr}
设 $f=(\psi,\theta):X\to Y$ 是预概形的一个态射；同态
\oldpage[I]{129}
$\theta_x^\sharp:\mathscr{O}_{\psi(x)}\to\mathscr{O}_x$ 把
$\mathscr{O}_{\psi(x)}$ 的每个幂零元映到 $\mathscr{O}_x$ 的一个
幂零元；转到商以后，因而由 $\theta^\sharp$ 导出一个同态
\[
  \omega:\psi^*(\mathscr{O}_Y/\mathscr{N}_Y)
  \longrightarrow \mathscr{O}_X/\mathscr{N}_X
\]
显然，对每个 $x\in X$，
$\omega_x:\mathscr{O}_{\psi(x)}/\mathscr{N}_{\psi(x)}
\to\mathscr{O}_x/\mathscr{N}_x$ 都是一个局部同态，因而
$(\psi,\omega^b)$ 是预概形的一个态射
$X_{\mathrm{red}}\to Y_{\mathrm{red}}$；我们把它记为 $f_{\mathrm{red}}$，
并称为与 $f$ 相伴的\emph{约化}态射。立即可知，对两个态射
$f:X\to Y$、$g:Y\to Z$，有
$(g\circ f)_{\mathrm{red}}=g_{\mathrm{red}}\circ f_{\mathrm{red}}$，
所以我们把 $X_{\mathrm{red}}$ 定义成关于 $X$ 的一个\emph{协变函子}。

上述定义表明，图表
\[
  \xymatrix{
    X_{\mathrm{red}}\ar[r]^{f_{\mathrm{red}}}\ar[d] &
    Y_{\mathrm{red}}\ar[d]\\
    X\ar[r]_f & Y
  }
\]
是交换的，其中竖直箭头是嵌入态射；换言之，
$X_{\mathrm{red}}\to X$ 是一个\emph{函子性的}态射。特别应注意，
若 $X$ 是\emph{约化的}，则每个态射 $f:X\to Y$ 都分解为
$X\xrightarrow{f_{\mathrm{red}}}Y_{\mathrm{red}}\to Y$；换言之，
$f$ \emph{可经由嵌入态射 $Y_{\mathrm{red}}\to Y$ 分解}。
\end{env}

\begin{proposition}[5.1.6]
\label{I.5.1.6-fr}
设 $f:X\to Y$ 是一个态射；若 $f$ 是满射（相应地，根态射、浸入、
闭浸入、开浸入、局部浸入、局部同构），则 $f_{\mathrm{red}}$ 也是如此。
反过来，若 $f_{\mathrm{red}}$ 是满射（相应地，根态射），则 $f$ 也是如此。
\end{proposition}

当 $f$ 是满射时，命题是平凡的；当 $f$ 是根态射时，它由下述事实得出：
对每个 $x\in X$，域 $k(x)$ 对预概形 $X$ 和
$X_{\mathrm{red}}$ 是相同的（\hyperref[I.3.5.8-fr]{3.5.8}）。最后，
若 $f=(\psi,\theta)$ 是浸入、闭浸入或局部浸入（相应地，开浸入或
局部同构），则命题由下述事实得出：若 $\theta_x^\sharp$ 是满射
（相应地，双射），则把 $\mathscr{O}_{\psi(x)}$ 和 $\mathscr{O}_x$
分别对其幂零根取商后所得的同态也是如此
（\hyperref[I.5.1.2-fr]{5.1.2} 与 \hyperref[I.4.2.2-fr]{4.2.2}）
（参见（\hyperref[I.5.5.12-fr]{5.5.12}））。

\begin{proposition}[5.1.7]
\label{I.5.1.7-fr}
若 $X$、$Y$ 是两个 $S$-预概形，则预概形
$X_{\mathrm{red}}\times_{S_{\mathrm{red}}}Y_{\mathrm{red}}$ 与
$X_{\mathrm{red}}\times_S Y_{\mathrm{red}}$ 相同，并且典范等同于
$X\times_S Y$ 的一个与该积具有相同底空间的子预概形。
\end{proposition}

$X_{\mathrm{red}}\times_S Y_{\mathrm{red}}$ 与
$X\times_S Y$ 的一个具有相同底空间的子预概形之间的典范等同，
由（\hyperref[I.4.3.1-fr]{4.3.1}）得出。另一方面，若 $\varphi$ 和
$\psi$ 分别是结构态射 $X_{\mathrm{red}}\to S$、
$Y_{\mathrm{red}}\to S$，则它们都经由 $S_{\mathrm{red}}$ 分解
（\hyperref[I.5.1.5-fr]{5.1.5}）；由于 $S_{\mathrm{red}}\to S$ 是
单态射，第一个断言遂由（\hyperref[I.3.2.4-fr]{3.2.4}）得出。

\begin{corollary}[5.1.8]
\label{I.5.1.8-fr}
预概形 $(X\times_S Y)_{\mathrm{red}}$ 与
$(X_{\mathrm{red}}\times_{S_{\mathrm{red}}}Y_{\mathrm{red}})_{\mathrm{red}}$
典范等同。
\end{corollary}

这由（\hyperref[I.5.1.2-fr]{5.1.2} 与
\hyperref[I.5.1.7-fr]{5.1.7}）得出。

应注意，即使 $X$ 和 $Y$ 是约化预概形，$X\times_S Y$ 也未必约化，
因为两个约化代数的张量积可以有幂零元。

\begin{proposition}[5.1.9]
\label{I.5.1.9-fr}
\oldpage[I]{130}
设 $X$ 是一个预概形，$\mathscr{J}$ 是 $\mathscr{O}_X$ 的一个拟凝聚
理想层，并且对某个整数 $n>0$ 有 $\mathscr{J}^n=0$。设 $X_0$ 是
$X$ 的闭子预概形 $(X,\mathscr{O}_X/\mathscr{J})$；$X$ 是仿射概形的
充要条件是 $X_0$ 是仿射概形。
\end{proposition}

该条件显然是必要的；我们来证明它是充分的。令
$X_k=(X,\mathscr{O}_X/\mathscr{J}^{k+1})$；只须对 $k$ 归纳证明各
$X_k$ 都是仿射的，因而归结为 $\mathscr{J}^2=0$ 的情形。令
\[
  A=\Gamma(X,\mathscr{O}_X),\qquad
  A_0=\Gamma(X_0,\mathscr{O}_{X_0})
      =\Gamma(X,\mathscr{O}_X/\mathscr{J}).
\]
由典范同态 $\mathscr{O}_X\to\mathscr{O}_X/\mathscr{J}$ 导出环同态
$\varphi:A\to A_0$。下面将看到，$\varphi$ 是\emph{满射}，因而序列
\[
  0\to\Gamma(X,\mathscr{J})\to\Gamma(X,\mathscr{O}_X)
  \to\Gamma(X,\mathscr{O}_X/\mathscr{J})\to 0
  \tag{5.1.9.1}\label{I.5.1.9.1-fr}
\]
是\emph{正合的}。先假定这一点已经成立，并证明它蕴含本命题。注意，
$\mathfrak{K}=\Gamma(X,\mathscr{J})$ 是 $A$ 中一个平方为零的理想，
因而是 $A_0=A/\mathfrak{K}$ 上的一个模。依假设，
$X_0=\operatorname{Spec}(A_0)$；由于底拓扑空间 $X_0$ 和 $X$ 相同，
$\mathfrak{K}=\Gamma(X_0,\mathscr{J})$。此外，由于 $\mathscr{J}^2=0$，
$\mathscr{J}$ 是一个拟凝聚 $(\mathscr{O}_X/\mathscr{J})$-模，所以有
$\mathscr{J}\simeq\widetilde{\mathfrak{K}}$，并且对每个 $x\in X_0$ 有
$\mathfrak{K}_x=\mathscr{J}_x$
（\hyperref[I.1.4.1-fr]{1.4.1}）。

现在令 $X'=\operatorname{Spec}(A)$，并考虑与恒等映射
$A\to\Gamma(X,\mathscr{O}_X)$ 相对应的预概形态射
$f=(\psi,\theta):X\to X'$（\hyperref[I.2.2.4-fr]{2.2.4}）。对 $X$
中的每个仿射开集 $V$，图表
\[
  \xymatrix{
    A\ar[r]\ar[d] & \Gamma(V,\mathscr{O}_X|V)\ar[d]\\
    A_0=A/\mathfrak{K}\ar[r] & \Gamma(V,\mathscr{O}_{X_0}|V)
  }
\]
是交换的，由此可知图表
\[
  \xymatrix{
    X' & X\ar[l]_f\\
    X'_0\ar[u]^{j'} & X_0\ar[l]^{f_0}\ar[u]_j
  }
\]
是交换的；这里，$X'_0$ 是由拟凝聚理想层
$\widetilde{\mathfrak{K}}$ 定义的 $X'$ 的闭子预概形，$j$、$j'$ 是
典范嵌入态射。但由于 $X_0$ 是仿射的，$f_0$ 是同构；而 $j$ 和 $j'$
所对应的底连续映射都是恒等映射，所以首先可见
$\psi:X\to X'$ 是同胚。此外，关系 $\mathfrak{K}_x=\mathscr{J}_x$
表明，$\theta^\sharp:\psi^*(\mathscr{O}_{X'})\to\mathscr{O}_X$ 的限制，
是从 $\psi^*(\widetilde{\mathfrak{K}})$ 到 $\mathscr{J}$ 的
\emph{同构}；另一方面，转到商以后，$\theta^\sharp$ 给出\emph{同构}
$\psi^*(\mathscr{O}_{X'}/\widetilde{\mathfrak{K}})
\to\mathscr{O}_X/\mathscr{J}$，因为 $f_0$ 是同构。由五引理
（M, I, 1.1）立即可知，$\theta^\sharp$ 本身也是同构，因而 $f$ 是
\emph{同构}，所以 $X$ 是仿射的。

因此，只须证明（\hyperref[I.5.1.9.1-fr]{5.1.9.1}）的正合性；这将由

\oldpage[I]{131}
$H^1(X,\mathscr{J})=0$ 得出。事实上，
$H^1(X,\mathscr{J})=H^1(X_0,\mathscr{J})$，并且已经看到，
$\mathscr{J}$ 是一个拟凝聚 $\mathscr{O}_{X_0}$-模。因而，我们的断言
由下述引理得出。

\begin{lemma}[5.1.9.2]
\label{I.5.1.9.2-fr}
若 $Y$ 是仿射概形，$\mathscr{F}$ 是一个拟凝聚 $\mathscr{O}_Y$-模，
则 $H^1(Y,\mathscr{F})=0$。
\end{lemma}

这个引理将在第三章 § 1 中证明，作为更一般的定理
$H^i(Y,\mathscr{F})=0$ 对每个 $i>0$ 都成立的推论。为给出一个独立证明，
注意，$H^1(Y,\mathscr{F})$ 可与下述模等同：
$\operatorname{Ext}^1_{\mathscr{O}_Y}
(Y;\mathscr{O}_Y,\mathscr{F})$，即以 $\mathscr{O}_Y$-模
$\mathscr{F}$ 扩张 $\mathscr{O}_Y$-模 $\mathscr{O}_Y$ 所得扩张类的模
（T, 4.2.3）；因而只须证明这样的扩张 $\mathscr{G}$ 是平凡的。事实上，
对每个 $y\in Y$，在 $Y$ 中存在 $y$ 的一个邻域 $V$，使得
% EGA1-ERR-0039: retain printed \mathscr{F}|Y; the local restriction must be \mathscr{F}|V.
$\mathscr{G}|V$ 同构于 $\mathscr{F}|Y\oplus\mathscr{O}_Y|V$
（\hyperref[0.5.4.9-fr]{0, 5.4.9}）；由此可知，$\mathscr{G}$ 是一个
拟凝聚 $\mathscr{O}_Y$-模。若 $A$ 是 $Y$ 的环，则
$\mathscr{F}=\widetilde{M}$、$\mathscr{G}=\widetilde{N}$，其中 $M$ 和
$N$ 是 $A$-模；依假设，$N$ 是以 $A$-模 $M$ 扩张 $A$-模 $A$ 所得的
扩张（\hyperref[I.1.3.11-fr]{1.3.11}）。由于这个扩张必然平凡，引理得证，
因而（\hyperref[I.5.1.9-fr]{5.1.9}）也得证。

\begin{corollary}[5.1.10]
\label{I.5.1.10-fr}
设 $X$ 是一个使 $\mathscr{N}_X$ 为幂零的预概形。$X$ 是仿射概形的
充要条件是 $X_{\mathrm{red}}$ 是仿射概形。
\end{corollary}

\subsection{具有给定底空间的子预概形的存在性。}
\label{subsection:I.5.2-fr}

\begin{proposition}[5.2.1]
\label{I.5.2.1-fr}
对预概形 $X$ 的底空间的每个局部闭子空间 $Y$，存在唯一的
$X$ 的约化子预概形，其底空间为 $Y$。
\end{proposition}

唯一性由（\hyperref[I.5.1.2-fr]{5.1.2}）得出，故只需证明所述子预概形
的存在性。

若 $X$ 是环 $A$ 的仿射概形，且 $Y$ 在 $X$ 中为闭的，则命题立即得出：
$\mathfrak{j}(Y)$ 是满足 $V(\mathfrak{a})=Y$ 的最大理想
$\mathfrak{a}\subset A$，并且它等于其根
（\hyperref[I.1.1.4-fr]{1.1.4 (i)}），因而
$A/\mathfrak{j}(Y)$ 是约化环。

在一般情形下，对每个使 $U\cap Y$ 在 $U$ 中为闭的仿射开集
$U\subset X$，考虑 $U$ 的闭子预概形 $Y_U$，它由与 $A(U)$ 的
理想 $\mathfrak{j}(U\cap Y)$ 相关联的理想层定义，并且是约化的。
我们证明，若 $V$ 是包含于 $U$ 的 $X$ 的一个仿射开集，则
$Y_V$ 是 $Y_U$ 在 $V\cap Y$ 上所\emph{诱导}的预概形；事实上，
这个诱导预概形是 $V$ 的一个约化闭子预概形，并且以
$V\cap Y$ 为底空间；$Y_V$ 的唯一性因而得出我们的断言。

\begin{proposition}[5.2.2]
\label{I.5.2.2-fr}
设 $X$ 是约化预概形，$f:X\to Y$ 是一个态射，$Z$ 是 $Y$ 的
一个闭子预概形，使得 $f(X)\subset Z$；则 $f$ 可分解为
$X\xrightarrow{g}Z\xrightarrow{j}Y$，其中 $j$ 是嵌入态射。
\end{proposition}

由假设可知，$X$ 的闭子预概形 $f^{-1}(Z)$ 的底空间是整个
$X$（\hyperref[I.4.4.1-fr]{4.4.1}）；由于 $X$ 是约化的，这个闭子预概形
与 $X$ 重合（\hyperref[I.5.1.2-fr]{5.1.2}），故命题由
（\hyperref[I.4.4.1-fr]{4.4.1}）得出。

\begin{corollary}[5.2.3]
\label{I.5.2.3-fr}
设 $X$ 是预概形 $Y$ 的一个约化子预概形；若 $Z$ 是 $Y$ 的
以 $\overline{X}$ 为底空间的约化闭子预概形，则 $X$ 是 $Z$ 在其一个
开子空间上所诱导的子预概形。
\end{corollary}

\oldpage[I]{132}
事实上，存在 $Y$ 的一个开集 $U$，使得
$X=U\cap\overline{X}$；根据（\hyperref[I.5.2.2-fr]{5.2.2}），
$X$ 是 $Z$ 的一个约化子预概形，故由唯一性
（\hyperref[I.5.2.1-fr]{5.2.1}），子预概形 $X$ 就是 $Z$ 在开子空间
$X$ 上所诱导的预概形。

\begin{corollary}[5.2.4]
\label{I.5.2.4-fr}
设 $f:X\to Y$ 是一个态射，$X'$ （相应地，$Y'$）是 $X$（相应地，
$Y$）的一个闭子预概形，它由 $\mathscr{O}_X$（相应地，
$\mathscr{O}_Y$）的拟凝聚理想层 $\mathscr{J}$（相应地，$\mathscr{K}$）
定义。假设 $X'$ 是约化的，且 $f(X')\subset Y'$。则
$f^*(\mathscr{K})\mathscr{O}_X\subset\mathscr{J}$。
\end{corollary}

根据（\hyperref[I.5.2.2-fr]{5.2.2}），$f$ 在 $X'$ 上的限制可分解为
$X'\to Y'\to Y$，因而只须应用
（\hyperref[I.4.4.6-fr]{4.4.6}）。

\subsection{对角线；态射的图。}
\label{subsection:I.5.3-fr}

\begin{env}[5.3.1]
\label{I.5.3.1-fr}
设 $X$ 是一个 $S$-预概形；称 $S$-态射 $(1_X,1_X)_S$ 为 $X$ 到
$X\times_S X$ 中的\emph{对角态射}，并记为 $\Delta_{X|S}$，或
$\Delta_X$，在不会混淆时甚至记为 $\Delta$；换言之，它是唯一使得
\[
  p_1\circ\Delta_X=p_2\circ\Delta_X=1_X
  \tag{5.3.1.1}\label{I.5.3.1.1-fr}
\]
的 $S$-态射 $\Delta_X$，其中 $p_1$、$p_2$ 表示 $X\times_S X$ 的投影
（定义 \hyperref[I.3.2.1-fr]{3.2.1}）。若 $f:T\to X$、$g:T\to Y$ 是两个
$S$-态射，则立即验证得
\[
  (f,g)_S=(f\times_S g)\circ\Delta_{T|S}
  \tag{5.3.1.2}\label{I.5.3.1.2-fr}
\]

读者应注意，上述定义以及编号
（\hyperref[I.5.3.1-fr]{5.3.1}）至（5.3.8）所述的结果在任意范畴中都成立，
\emph{只要其中出现的积在该范畴中存在}。
\end{env}

\begin{proposition}[5.3.2]
\label{I.5.3.2-fr}
设 $X$、$Y$ 是两个 $S$-预概形；若把积
$(X\times Y)\times(X\times Y)$ 典范地等同于
$(X\times X)\times(Y\times Y)$，则态射 $\Delta_{X\times Y}$ 等同于
$\Delta_X\times\Delta_Y$。
\end{proposition}

事实上，若 $p_1$、$q_1$ 分别是第一投影
$X\times X\to X$、$Y\times Y\to Y$，则第一投影
$(X\times Y)\times(X\times Y)\to X\times Y$ 等同于 $p_1\times q_1$，并且
\[
  (p_1\times q_1)\circ(\Delta_X\times\Delta_Y)
  =(p_1\circ\Delta_X)\times(q_1\circ\Delta_Y)=1_{X\times Y}
\]
对第二投影作同样的论证即可。

\begin{corollary}[5.3.4]
\label{I.5.3.4-fr}
对基预概形的任意扩张 $S'\to S$，
$\Delta_{X_{(S')}}$ 典范地等同于 $(\Delta_X)_{(S')}$。
\end{corollary}

只须注意，$(X\times_S X)_{(S')}$ 典范地等同于
$X_{(S')}\times_{S'}X_{(S')}$
（\hyperref[I.3.3.10-fr]{3.3.10}）。

\begin{proposition}[5.3.5]
\label{I.5.3.5-fr}
设 $X$、$Y$ 是两个 $S$-预概形，$\varphi:S\to T$ 是一个态射，从而使
每个 $S$-预概形成为一个 $T$-预概形。
% EGA1-ERR-0040: retain the printed omission of the name g before Y\to S.
设 $f:X\to S$、$Y\to S$ 是结构态射，$p$、$q$ 是 $X\times_S Y$ 的投影，
$\pi=f\circ p=g\circ q$ 是结构态射 $X\times_S Y\to S$。则图表
\[
  \xymatrix{
    X\times_S Y\ar[r]^{(p,q)_T}\ar[d]_\pi &
    X\times_T Y\ar[d]^{f\times_T g}\\
    S\ar[r]_{\Delta_{S|T}} &
    S\times_T S
  }
  \tag{5.3.5.1}\label{I.5.3.5.1-fr}
\]
\oldpage[I]{133}
是交换的，并且把 $X\times_S Y$ 等同于两个
$(S\times_T S)$-预概形 $S$ 与 $X\times_T Y$ 的积，其投影分别等同于
$\pi$ 与 $(p,q)_T$。
\end{proposition}

根据（\hyperref[I.3.4.3-fr]{3.4.3}），只须证明\emph{集合范畴中}相应的命题，
其中把 $X$、$Y$、$S$ 分别换成 $X(Z)_T$、$Y(Z)_T$、$S(Z)_T$，
$Z$ 是任意 $T$-预概形。但对于集合范畴，验证是立即的，留给读者。

\begin{corollary}[5.3.6]
\label{I.5.3.6-fr}
态射 $(p,q)_T$ 等同于（令 $P=S\times_T S$）
$1_{X\times_T Y}\times_P\Delta_S$。
\end{corollary}

这由（\hyperref[I.5.3.5-fr]{5.3.5}）和
（\hyperref[I.3.3.4-fr]{3.3.4}）得出。

\begin{corollary}[5.3.7]
\label{I.5.3.7-fr}
若 $f:X\to Y$ 是一个 $S$-态射，则图表
\[
  \xymatrix{
    X\ar[r]^{(1_X,f)_S}\ar[d]_f &
    X\times_S Y\ar[d]^{f\times_S 1_Y}\\
    Y\ar[r]_{\Delta_Y} &
    Y\times_S Y
  }
\]
是交换的，并且把 $X$ 等同于两个
$(Y\times_S Y)$-预概形 $Y$ 与 $X\times_S Y$ 的积。
\end{corollary}

只须应用（\hyperref[I.5.3.5-fr]{5.3.5}），其中把 $S$ 换成 $Y$、
$T$ 换成 $S$，并注意 $X\times_Y Y=X$
（\hyperref[I.3.3.3-fr]{3.3.3}）。

\begin{proposition}[5.3.8]
\label{I.5.3.8-fr}
$f:X\to Y$ 是预概形的单态射的充要条件是：
$\Delta_{X|Y}$ 是 $X$ 到 $X\times_Y X$ 上的一个同构。
\end{proposition}

事实上，$f$ 是单态射，就是说对每个 $Y$-预概形 $Z$，相应的映射
$f':X(Z)_Y\to Y(Z)_Y$ 是单射；又因 $Y(Z)_Y$ 缩为一个元素，
% EGA1-ERR-0041: retain the printed claim that X(Z)_Y also has one element; injectivity only gives at most one.
这意味着 $X(Z)_Y$ 也是如此。但这也可表述为：
$X(Z)_Y\times X(Z)_Y$ 典范同构于 $X(Z)_Y$；而这两个集合中的
第一个是 $(X\times_Y X)(Z)_Y$
（\hyperref[I.3.4.3.1-fr]{3.4.3.1}），这就意味着 $\Delta_{X|Y}$ 是一个同构。

\begin{proposition}[5.3.9]
\label{I.5.3.9-fr}
对角态射 $\Delta_X$ 是 $X$ 到 $X\times_S X$ 中的一个浸入。
\end{proposition}

% EGA1-ERR-0042: retain the printed proof; it does not establish that Delta_X(X) is locally closed.
事实上，由于底空间的连续映射 $p_1$ 与 $\Delta_X$ 满足
$p_1\circ\Delta_X$ 是恒等映射，$\Delta_X$ 是 $X$ 到 $\Delta_X(X)$ 上的
一个同胚。同样，与 $p_1$ 和 $\Delta_X$ 对应的同态的复合同态
$\mathscr{O}_x\to\mathscr{O}_{\Delta_X(x)}\to\mathscr{O}_x$ 是恒等的，
因而与 $\Delta_X$ 对应的同态是满射；所以命题由
（\hyperref[I.4.2.2-fr]{4.2.2}）得出。

称与浸入 $\Delta_X$（\hyperref[I.4.2.1-fr]{4.2.1}）相关联的
$X\times_S X$ 的子预概形为 $X\times_S X$ 的\emph{对角线}。

\begin{corollary}[5.3.10]
\label{I.5.3.10-fr}
在（\hyperref[I.5.3.5-fr]{5.3.5}）的假设下，$(p,q)_T$ 是一个浸入。
\end{corollary}

这由（\hyperref[I.5.3.6-fr]{5.3.6}）和
（\hyperref[I.4.3.1-fr]{4.3.1}）得出。

称（在（\hyperref[I.5.3.5-fr]{5.3.5}）的假设下）$(p,q)_T$ 为
$X\times_S Y$ 到 $X\times_T Y$ 中的\emph{典范浸入}。

\begin{corollary}[5.3.11]
\label{I.5.3.11-fr}
设 $X$、$Y$ 是两个 $S$-预概形，$f:X\to Y$ 是一个 $S$-态射；则 $f$ 的
图态射 $\Gamma_f=(1_X,f)_S$（\hyperref[I.3.3.14-fr]{3.3.14}）是 $X$ 到
$X\times_S Y$ 中的一个浸入。
\end{corollary}

这是推论（\hyperref[I.5.3.10-fr]{5.3.10}）的特殊情形，其中把
$S$ 换成 $Y$、$T$ 换成 $S$
（参见（\hyperref[I.5.3.7-fr]{5.3.7}））。

\oldpage[I]{134}
与浸入 $\Gamma_f$（\hyperref[I.4.2.1-fr]{4.2.1}）相关联的
$X\times_S Y$ 的子预概形称为态射 $f$ 的\emph{图}；
$X\times_S Y$ 的那些是态射 $X\to Y$ 的图的子预概形，可由下述性质刻画：
投影 $p_1:X\times_S Y\to X$ 在这样的子预概形 $G$ 上的限制，是 $G$ 到
$X$ 上的一个\emph{同构} $g$；此时 $G$ 是态射 $p_2\circ g^{-1}$ 的图，
其中 $p_2$ 是投影 $X\times_S Y\to Y$。

特别地，当取 $X=S$ 时，$S$-态射 $S\to Y$，也就是 $Y$ 的
$S$-截面（\hyperref[I.2.5.5-fr]{2.5.5}），等于它们的图态射；$Y$ 的
那些是 $S$-截面的图的子预概形（换言之，那些经由结构态射
$Y\to S$ 的限制而与 $S$ 同构的子预概形），也称为这些截面的
\emph{像}，或滥用语言称为 $Y$ 的\emph{$S$-截面}。

\begin{corollary}[5.3.12]
\label{I.5.3.12-fr}
假设和记号如（\hyperref[I.5.3.11-fr]{5.3.11}）所述。对每个态射
$g:S'\to S$，令 $f'$ 是 $f$ 在 $g$ 下的逆像
（\hyperref[I.3.3.7-fr]{3.3.7}）；则 $\Gamma_{f'}$ 是 $\Gamma_f$ 在 $g$ 下的逆像。
\end{corollary}

这是公式（\hyperref[I.3.3.10.1-fr]{3.3.10.1}）的一个特殊情形。

\begin{corollary}[5.3.13]
\label{I.5.3.13-fr}
设 $f:X\to Y$、$g:Y\to Z$ 是两个态射；若 $g\circ f$ 是一个浸入
（相应地，一个局部浸入），则 $f$ 也是如此。
\end{corollary}

事实上，$f$ 可分解为
\[
  X\xrightarrow{\Gamma_f}X\times_Z Y\xrightarrow{p_2}Y.
\]
另一方面，$p_2$ 等同于 $(g\circ f)\times_Z 1_Y$
（\hyperref[I.3.3.4-fr]{3.3.4}）；若 $g\circ f$ 是一个浸入（相应地，
一个局部浸入），则 $p_2$ 也是如此
（\hyperref[I.4.3.1-fr]{4.3.1} 和 \hyperref[I.4.5.5-fr]{4.5.5}）；又因
$\Gamma_f$ 是一个浸入（\hyperref[I.5.3.11-fr]{5.3.11}），故由
% EGA1-ERR-0043: retain printed 4.2.4; composition/transitivity requires 4.2.5.
（\hyperref[I.4.2.4-fr]{4.2.4}）（相应地，
（\hyperref[I.4.5.5-fr]{4.5.5}））得出结论。

\begin{corollary}[5.3.14]
\label{I.5.3.14-fr}
设 $j:X\to Y$、$g:X\to Z$ 是两个 $S$-态射。若 $j$ 是一个浸入
（相应地，一个局部浸入），则 $(j,g)_S$ 也是如此。
\end{corollary}

事实上，若 $p:Y\times_S Z\to Y$ 是第一投影，则
$j=p\circ(j,g)_S$，因而只须应用
（\hyperref[I.5.3.13-fr]{5.3.13}）。

\begin{proposition}[5.3.15]
\label{I.5.3.15-fr}
若 $f:X\to Y$ 是一个 $S$-态射，则图表
\[
  \xymatrix{
    X\ar[r]^{\Delta_X}\ar[d]_f &
    X\times_S X\ar[d]^{f\times_S f}\\
    Y\ar[r]_{\Delta_Y} &
    Y\times_S Y
  }
  \tag{5.3.15.1}\label{I.5.3.15.1-fr}
\]
是交换的（换言之，$\Delta_X$ 是预概形范畴中的一个函子性态射）。
\end{proposition}

验证是立即的，留给读者。

\begin{corollary}[5.3.16]
\label{I.5.3.16-fr}
若 $X$ 是 $Y$ 的一个子预概形，则对角线 $\Delta_X(X)$ 等同于
$\Delta_Y(Y)$ 的一个子预概形，其底空间等同于
\[
  \Delta_Y(Y)\cap p_1^{-1}(X)=\Delta_Y(Y)\cap p_2^{-1}(X)
\]
（$p_1$、$p_2$ 是 $Y\times_S Y$ 的投影）。
\end{corollary}

对嵌入态射 $f:X\to Y$ 应用（\hyperref[I.5.3.15-fr]{5.3.15}）；于是知道
$f\times_S f$ 是一个浸入，把 $X\times_S X$ 的底空间等同于
$Y\times_S Y$ 的子空间 $p_1^{-1}(X)\cap p_2^{-1}(X)$
（\hyperref[I.4.3.1-fr]{4.3.1}）；此外，若
$z\in\Delta_Y(Y)\cap p_1^{-1}(X)$，则 $z=\Delta_Y(y)$

\oldpage[I]{135}
并且 $y=p_1(z)\in X$，因而 $y=f(y)$，且 $z=\Delta_Y(f(y))$ 根据图表
（\hyperref[I.5.3.15.1-fr]{5.3.15.1}）的交换性属于 $\Delta_X(X)$。

\begin{corollary}[5.3.17]
\label{I.5.3.17-fr}
设 $f_1:Y\to X$、$f_2:Y\to X$ 是两个 $S$-态射，$y$ 是 $Y$ 的一个点，
使得 $f_1(y)=f_2(y)=x$，并且与 $f_1$ 和 $f_2$ 对应的同态
$k(x)\to k(y)$ 相同。则，若 $f=(f_1,f_2)_S$，点 $f(y)$ 属于对角线
$\Delta_{X|S}(X)$。
\end{corollary}

与 $f_i$（$i=1,2$）对应的两个同态 $k(x)\to k(y)$ 定义两个
$S$-态射 $g_i:\operatorname{Spec}(k(y))\to\operatorname{Spec}(k(x))$，使图表
\[
  \xymatrix{
    \operatorname{Spec}(k(y))\ar[r]^{g_i}\ar[d] &
    \operatorname{Spec}(k(x))\ar[d]\\
    Y\ar[r]_{f_i} & X
  }
\]
交换。因而图表
\[
  \xymatrix{
    \operatorname{Spec}(k(y))\ar[rr]^{(g_1,g_2)_S}\ar[d] & &
    \operatorname{Spec}(k(x))\times_S\operatorname{Spec}(k(x))\ar[d]\\
    Y\ar[rr]_{(f_1,f_2)_S} & & X\times_S X
  }
\]
也是交换的。而由等式 $g_1=g_2$ 可知，$\operatorname{Spec}(k(y))$ 的
唯一一点在 $(g_1,g_2)_S$ 下的像属于
$\operatorname{Spec}(k(x))\times_S\operatorname{Spec}(k(x))$ 的对角线；因而结论由
（\hyperref[I.5.3.15-fr]{5.3.15}）得出。

\subsection{分离态射与分离预概形。}
\label{subsection:I.5.4-fr}

\begin{definition}[5.4.1]
\label{I.5.4.1-fr}
称预概形的态射 $f:X\to Y$ 是\emph{分离的}，若对角态射
$X\to X\times_Y X$ 是一个\emph{闭浸入}；这时也称 $X$ 是一个
\emph{$Y$ 上的分离预概形}，或一个\emph{$Y$-概形}。称预概形 $X$
是分离的，若它在 $\operatorname{Spec}(\mathbf{Z})$ 上分离；这时也称
$X$ 是一个\emph{概形}（参见（\hyperref[I.5.5.7-fr]{5.5.7}））。
\end{definition}

根据（\hyperref[I.5.3.9-fr]{5.3.9}），$X$ 在 $Y$ 上分离的充要条件是
$\Delta_X(X)$ 为 $X\times_Y X$ 的底空间中的一个\emph{闭子空间}。

\begin{proposition}[5.4.2]
\label{I.5.4.2-fr}
设 $S\to T$ 是一个分离态射。若 $X$ 和 $Y$ 是两个 $S$-预概形，则
典范浸入 $X\times_S Y\to X\times_T Y$
（\hyperref[I.5.3.10-fr]{5.3.10}）是闭浸入。
\end{proposition}

事实上，参照图表（\hyperref[I.5.3.5.1-fr]{5.3.5.1}），可见
$(p,q)_T$ 可视为由 $\Delta_{S|T}$ 经基变换
$f\times_T g:X\times_T Y\to S\times_T S$ 得到（这里的基预概形为
$S\times_T S$）；于是命题由
（\hyperref[I.4.3.2-fr]{4.3.2}）得出。

\begin{corollary}[5.4.3]
\label{I.5.4.3-fr}
设 $Y$ 是一个 $S$-概形，$f:X\to Y$ 是一个 $S$-态射。则图态射
$\Gamma_f:X\to X\times_S Y$（\hyperref[I.5.3.11-fr]{5.3.11}）
是一个闭浸入。
\end{corollary}

这是（\hyperref[I.5.4.2-fr]{5.4.2}）中以 $Y$ 代替 $S$、以 $S$
代替 $T$ 的特殊情形。

\begin{corollary}[5.4.4]
\label{I.5.4.4-fr}
设 $f:X\to Y$、$g:Y\to Z$ 是两个态射，并且 $g$ 分离。若 $g\circ f$
是一个闭浸入，则 $f$ 也是闭浸入。
\end{corollary}

从（\hyperref[I.5.4.3-fr]{5.4.3}）出发的证明，与
（\hyperref[I.5.3.13-fr]{5.3.13}）的、从
（\hyperref[I.5.3.11-fr]{5.3.11}）出发的证明相同。

\oldpage[I]{136}
\begin{corollary}[5.4.5]
\label{I.5.4.5-fr}
设 $Z$ 是一个 $S$-概形，$j:X\to Y$、$g:X\to Z$ 是两个 $S$-态射。
若 $j$ 是一个闭浸入，则
$(j,g)_S:X\to Y\times_S Z$ 也是闭浸入。
\end{corollary}

从（\hyperref[I.5.4.4-fr]{5.4.4}）出发的证明，与
（\hyperref[I.5.3.14-fr]{5.3.14}）的、从
（\hyperref[I.5.3.13-fr]{5.3.13}）出发的证明相同。

\begin{corollary}[5.4.6]
\label{I.5.4.6-fr}
若 $X$ 是一个 $S$-概形，则 $X$ 的每个 $S$-截面
（\hyperref[I.2.5.5-fr]{2.5.5}）都是闭浸入。
\end{corollary}

若 $\varphi:X\to S$ 是结构态射，$\psi:S\to X$ 是 $X$ 的一个
$S$-截面，只需把（\hyperref[I.5.4.5-fr]{5.4.5}）应用于
$\varphi\circ\psi=1_S$。

\begin{corollary}[5.4.7]
\label{I.5.4.7-fr}
设 $S$ 是一个整预概形，$s$ 是它的泛点，$X$ 是一个 $S$-概形。
若 $X$ 的两个 $S$-截面 $f$、$g$ 满足 $f(s)=g(s)$，则 $f=g$。
\end{corollary}

事实上，若 $x=f(s)=g(s)$，则与 $f$ 和 $g$ 对应的同态
$k(x)\to k(s)$ 必然相同。若 $h=(f,g)_S$，则由
（\hyperref[I.5.3.17-fr]{5.3.17}）可知 $h(s)$ 属于对角线
$Z=\Delta_X(X)$；但由于 $S=\overline{\{s\}}$，且 $Z$ 依假设为闭集，
所以 $h(S)\subset Z$。于是由（\hyperref[I.5.2.2-fr]{5.2.2}）可知，
$h$ 可分解为 $S\to Z\to X\times_S X$，再由对角线的定义得出
$f=g$。

\begin{remark}[5.4.8]
\label{I.5.4.8-fr}
反过来，若假设在 $f=1_Y$ 时（\hyperref[I.5.4.3-fr]{5.4.3}）的结论
成立，则可知 $Y$ 在 $S$ 上分离；同样，若假设
% EGA1-ERR-0047: printed 5.4.5 retained; the displayed composable pair invokes the cancellation conclusion of 5.4.4.
（\hyperref[I.5.4.5-fr]{5.4.5}）的结论适用于两个态射
\[
  Y\xrightarrow{\Delta_Y}Y\times_Z Y\xrightarrow{p_1}Y,
\]
则可知 $\Delta_Y$ 是一个闭浸入，因而 $Y$ 在 $Z$ 上分离；最后，若
（\hyperref[I.5.4.6-fr]{5.4.6}）的结论对 $Y$-预概形
$Y\times_S Y\to Y$ 的 $Y$-截面 $\Delta_Y$ 成立，则蕴含 $Y$ 在
$S$ 上分离。
\end{remark}

\subsection{分离判据。}
\label{subsection:I.5.5-fr}

\begin{proposition}[5.5.1]
\label{I.5.5.1-fr}
\begin{enumerate}
  \item[\rm(i)] 预概形的每个单态射（特别是每个浸入）都是分离态射。
  \item[\rm(ii)] 两个分离态射的复合是分离的。
  \item[\rm(iii)] 若 $f:X\to X'$、$g:Y\to Y'$ 是两个分离的
    $S$-态射，则 $f\times_S g$ 分离。
  \item[\rm(iv)] 若 $f:X\to Y$ 是一个分离的 $S$-态射，则对于基预概形
    的每个扩张 $S'\to S$，$S'$-态射 $f_{(S')}$ 都分离。
  \item[\rm(v)] 若两个态射的复合 $g\circ f$ 分离，则 $f$ 分离。
  \item[\rm(vi)] 态射 $f$ 分离的充要条件是
    $f_{\mathrm{red}}$（\hyperref[I.5.1.5-fr]{5.1.5}）分离。
\end{enumerate}
\end{proposition}

(i) 立即由（\hyperref[I.5.3.8-fr]{5.3.8}）得出。若 $f:X\to Y$、
$g:Y\to Z$ 是两个态射，则图表
\[
  \xymatrix{
    X\ar[rr]^{\Delta_{X|Z}}\ar[dr]_{\Delta_{X|Y}} & &
    X\times_Z X\\
    & X\times_Y X\ar[ur]_j
  }
  \tag{5.5.1.1}\label{I.5.5.1.1-fr}
\]
是交换的，这一点可以立即验证；这里 $j$ 表示典范浸入
（\hyperref[I.5.3.10-fr]{5.3.10}）。若 $f$、$g$ 分离，则
$\Delta_{X|Y}$ 按定义是闭浸入，而 $j$ 由
（\hyperref[I.5.4.2-fr]{5.4.2}）也是闭浸入，因而
$\Delta_{X|Z}$ 由（\hyperref[I.4.2.4-fr]{4.2.4}）是闭浸入，从而

\oldpage[I]{137}
证明了 (ii)。由 (i) 和 (ii)，(iii) 与 (iv) 等价
（\hyperref[I.3.5.1-fr]{3.5.1}），所以只需证明 (iv)。而根据
（\hyperref[I.3.3.11-fr]{3.3.11}）和
（\hyperref[I.3.3.9.1-fr]{3.3.9.1}），
$X_{(S')}\times_{Y_{(S')}}X_{(S')}$ 典范地等同于
$(X\times_Y X)\times_Y Y_{(S')}$；并且立即验证，对角态射
$\Delta_{X_{(S')}}$ 这时等同于
$\Delta_X\times_Y 1_{Y_{(S')}}$；故命题由
（\hyperref[I.4.3.1-fr]{4.3.1}）得出。

为证明 (v)，像在（\hyperref[I.5.3.13-fr]{5.3.13}）中那样，考虑
$f$ 的分解
$X\xrightarrow{\Gamma_f}X\times_Z Y\xrightarrow{p_2}Y$，并注意到
$p_2=(g\circ f)\times_Z 1_Y$；$g\circ f$ 分离这一假设，由 (iii) 和
(i) 推出 $p_2$ 分离；又因为 $\Gamma_f$ 是一个浸入，由 (i) 可知
$\Gamma_f$ 分离，故由 (ii) 得 $f$ 分离。最后，为证明 (vi)，回忆
预概形 $X_{\mathrm{red}}\times_{Y_{\mathrm{red}}}X_{\mathrm{red}}$ 与
$X_{\mathrm{red}}\times_Y X_{\mathrm{red}}$ 典范等同
（\hyperref[I.5.1.7-fr]{5.1.7}）；若以 $j$ 表示嵌入
$X_{\mathrm{red}}\to X$，则图表
\[
  \xymatrix{
    X_{\mathrm{red}}\ar[r]^{\Delta_{X_{\mathrm{red}}}}\ar[d]_j &
    X_{\mathrm{red}}\times_Y X_{\mathrm{red}}\ar[d]^{j\times_Y j}\\
    X\ar[r]_{\Delta_X} & X\times_Y X
  }
\]
交换（\hyperref[I.5.3.15-fr]{5.3.15}）；命题由竖直箭头都是底空间的
同胚（\hyperref[I.4.3.1-fr]{4.3.1}）这一事实得出。

\begin{corollary}[5.5.2]
\label{I.5.5.2-fr}
若 $f:X\to Y$ 分离，则 $f$ 在 $X$ 的任意子预概形上的限制都分离。
\end{corollary}

这由（\hyperref[I.5.5.1-fr]{5.5.1, (i) 和 (ii)}）得出。

\begin{corollary}[5.5.3]
\label{I.5.5.3-fr}
若 $X$、$Y$ 是两个 $S$-预概形，并且 $Y$ 在 $S$ 上分离，则
$X\times_S Y$ 在 $X$ 上分离。
\end{corollary}

这是（\hyperref[I.5.5.1-fr]{5.5.1, (iv)}）的一个特殊情形。

\begin{proposition}[5.5.4]
\label{I.5.5.4-fr}
设 $X$ 是一个预概形，并假设其底空间是有限个闭子集 $X_k$
$(1\leq k\leq n)$ 的并；对每个 $k$，考虑 $X$ 的以 $X_k$ 为
底空间的约化子预概形（\hyperref[I.5.2.1-fr]{5.2.1}），并仍把它记为
$X_k$。设 $f:X\to Y$ 是一个态射；对每个 $k$，设 $Y_k$ 是 $Y$ 的
一个闭子集，使得 $f(X_k)\subset Y_k$；仍以 $Y_k$ 表示 $Y$ 的以
$Y_k$ 为底空间的约化子预概形，于是 $f$ 在 $X_k$ 上的限制
$X_k\to Y$ 可分解为
$X_k\xrightarrow{f_k}Y_k\to Y$（\hyperref[I.5.2.2-fr]{5.2.2}）。
$f$ 分离的充要条件是所有 $f_k$ 都分离。
\end{proposition}

必要性由（\hyperref[I.5.5.1-fr]{5.5.1, (i)、(ii) 和 (v)}）得出。
反过来，若命题中的条件成立，则 $f$ 的每个限制 $X_k\to Y$ 都分离
（\hyperref[I.5.5.1-fr]{5.5.1, (i) 和 (ii)}）；若 $p_1$、$p_2$ 是
$X\times_Y X$ 的投影，则子空间 $\Delta_{X_k}(X_k)$ 等同于
$X\times_Y X$ 的底空间的子空间
$\Delta_X(X)\cap p_1^{-1}(X_k)$
（\hyperref[I.5.3.16-fr]{5.3.16}）；这些子空间在 $X\times_Y X$ 中
都是闭的，故它们的并 $\Delta_X(X)$ 也是闭的。

特别地，假设 $X_k$ 是 $X$ 的各\emph{不可约分支}；这时可以假设
每个 $Y_k$ 都是 $Y$ 的一个不可约分支（\hyperref[0.2.1.5-fr]{0, 2.1.5}）；
因此命题（\hyperref[I.5.5.4-fr]{5.5.4}）在此情形下把分离概念归结为
\emph{整}预概形的情形（\hyperref[I.2.1.7-fr]{2.1.7}）。

\oldpage[I]{138}

\begin{proposition}[5.5.5]
\label{I.5.5.5-fr}
设 $(Y_\lambda)$ 是预概形 $Y$ 的一个开覆盖；态射 $f:X\to Y$ 分离的
充要条件是它的每个限制 $f^{-1}(Y_\lambda)\to Y_\lambda$ 都分离。
\end{proposition}

若置 $X_\lambda=f^{-1}(Y_\lambda)$，则考虑到
（\hyperref[I.4.2.4-fr]{4.2.4, b)}）以及积
$X_\lambda\times_Y X_\lambda$ 与
$X_\lambda\times_{Y_\lambda}X_\lambda$ 的等同
（\hyperref[I.3.2.5-fr]{3.2.5}），问题归结为证明
$X_\lambda\times_Y X_\lambda$ 构成 $X\times_Y X$ 的一个覆盖。若置
$Y_{\lambda\mu}=Y_\lambda\cap Y_\mu$ 和
$X_{\lambda\mu}=X_\lambda\cap X_\mu=f^{-1}(Y_{\lambda\mu})$，则
$X_\lambda\times_Y X_\mu$ 等同于积
$X_{\lambda\mu}\times_{Y_{\lambda\mu}}X_{\lambda\mu}$
（\hyperref[I.3.2.6.4-fr]{3.2.6.4}），因而也等同于
$X_{\lambda\mu}\times_Y X_{\lambda\mu}$
（\hyperref[I.3.2.5-fr]{3.2.5}），最后等同于
$X_\lambda\times_Y X_\lambda$ 的一个开集；这就证明了断言
（\hyperref[I.3.2.7-fr]{3.2.7}）。

% EGA1-ERR-0044: the three printed 5.5.4 references in this reduction and the proof of 5.5.9 are retained; sidecar evidence shows 5.5.5 is intended.
命题（\hyperref[I.5.5.4-fr]{5.5.4}）允许取 $Y$ 的一个仿射开覆盖，
从而把分离态射的研究归结为取值于仿射概形的分离态射的研究。

\begin{proposition}[5.5.6]
\label{I.5.5.6-fr}
设 $Y$ 是一个仿射概形，$X$ 是一个预概形，$(U_\alpha)$ 是 $X$ 的
一个仿射开覆盖。态射 $f:X\to Y$ 分离的充要条件是：对每一对指标
$(\alpha,\beta)$，$U_\alpha\cap U_\beta$ 都是仿射开集，并且环
$\Gamma(U_\alpha\cap U_\beta,\mathscr{O}_X)$ 由环
$\Gamma(U_\alpha,\mathscr{O}_X)$ 与
$\Gamma(U_\beta,\mathscr{O}_X)$ 的典范像之并生成。
\end{proposition}

$U_\alpha\times_Y U_\beta$ 构成 $X\times_Y X$ 的一个开覆盖
（\hyperref[I.3.2.7-fr]{3.2.7}）；以 $p$、$q$ 表示
$X\times_Y X$ 的投影，则
\[
  \begin{aligned}
    \Delta_X^{-1}(U_\alpha\times_Y U_\beta)
      &=\Delta_X^{-1}\bigl(p^{-1}(U_\alpha)\cap q^{-1}(U_\beta)\bigr)\\
      &=\Delta_X^{-1}\bigl(p^{-1}(U_\alpha)\bigr)
        \cap\Delta_X^{-1}\bigl(q^{-1}(U_\beta)\bigr)
        =U_\alpha\cap U_\beta~;
  \end{aligned}
\]
因此问题归结为表明 $\Delta_X$ 在 $U_\alpha\cap U_\beta$ 上的限制，
是到 $U_\alpha\times_Y U_\beta$ 中的闭浸入。由定义可知，这个限制
恰是 $(j_\alpha,j_\beta)_Y$；这里 $j_\alpha$（相应地 $j_\beta$）表示
$U_\alpha\cap U_\beta$ 到 $U_\alpha$（相应地 $U_\beta$）的嵌入态射。
由于 $U_\alpha\times_Y U_\beta$ 是一个仿射概形，其环典范同构于
$\Gamma(U_\alpha,\mathscr{O}_X)
 \otimes_{\Gamma(Y,\mathscr{O}_Y)}
 \Gamma(U_\beta,\mathscr{O}_X)$
（\hyperref[I.3.2.2-fr]{3.2.2}），可知 $U_\alpha\cap U_\beta$ 必须是
仿射概形，并且从环 $A(U_\alpha\times_Y U_\beta)$ 到
$\Gamma(U_\alpha\cap U_\beta,\mathscr{O}_X)$ 的映射
$h_\alpha\otimes h_\beta\to h_\alpha h_\beta$ 必须是满射
（\hyperref[I.4.2.3-fr]{4.2.3}）；证明完毕。

\begin{corollary}[5.5.7]
\label{I.5.5.7-fr}
仿射概形是分离的（因而是概形；这就说明了
（\hyperref[I.5.4.1-fr]{5.4.1}）中术语的理由）。
\end{corollary}

\begin{corollary}[5.5.8]
\label{I.5.5.8-fr}
设 $Y$ 是一个仿射概形；态射 $f:X\to Y$ 分离的充要条件是 $X$
分离（换言之，$X$ 是一个概形）。
\end{corollary}

事实上应注意，（\hyperref[I.5.5.6-fr]{5.5.6}）的判据与 $f$ 无关。

\begin{corollary}[5.5.9]
\label{I.5.5.9-fr}
态射 $f:X\to Y$ 分离的必要条件是：对于 $Y$ 在其上诱导出分离
预概形的每个开集 $U$，诱导预概形 $f^{-1}(U)$ 都分离；充分条件是：
这对每个仿射开集 $U\subset Y$ 都成立。
\end{corollary}

% EGA1-ERR-0044: source-literal 5.5.4 retained; the locality proposition required here is 5.5.5.
条件的必要性由（\hyperref[I.5.5.4-fr]{5.5.4}）和
（\hyperref[I.5.5.1-fr]{5.5.1, (ii)}）得出；充分性由
% EGA1-ERR-0044: source-literal 5.5.4 retained; the locality proposition required here is 5.5.5.
（\hyperref[I.5.5.4-fr]{5.5.4}）和
（\hyperref[I.5.5.8-fr]{5.5.8}）得出，同时要用到 $Y$ 存在仿射开覆盖。

特别地，若 $X$ 和 $Y$ 都是仿射概形，则\emph{每个}态射
$X\to Y$ 都分离。

\begin{proposition}[5.5.10]
\label{I.5.5.10-fr}
设 $Y$ 是一个概形，$f:X\to Y$ 是一个态射。对于 $X$ 的每个仿射
开集 $U$ 和 $Y$ 的每个仿射开集 $V$，$U\cap f^{-1}(V)$ 都是仿射的。
\end{proposition}

设 $p_1$、$p_2$ 是 $X\times_{\mathbf{Z}}Y$ 的投影；子空间
$U\cap f^{-1}(V)$ 是
$\Gamma_f(X)\cap p_1^{-1}(U)\cap p_2^{-1}(V)$ 在 $p_1$ 下的像。
而 $p_1^{-1}(U)\cap p_2^{-1}(V)$ 等同于

\oldpage[I]{139}
预概形 $U\times_{\mathbf{Z}}V$
（\hyperref[I.3.2.7-fr]{3.2.7}）的底空间，因而是一个仿射概形
（\hyperref[I.3.2.2-fr]{3.2.2}）；由于 $\Gamma_f(X)$ 在
$X\times_{\mathbf{Z}}Y$ 中为闭集（\hyperref[I.5.4.3-fr]{5.4.3}），
$\Gamma_f(X)\cap p_1^{-1}(U)\cap p_2^{-1}(V)$ 在
$U\times_{\mathbf{Z}}V$ 中为闭集；因而，与 $\Gamma_f$ 相伴的
$X\times_{\mathbf{Z}}Y$ 的子预概形（\hyperref[I.4.2.1-fr]{4.2.1}）
在其底空间的开部分
$\Gamma_f(X)\cap p_1^{-1}(U)\cap p_2^{-1}(V)$ 上诱导的预概形，
是一个仿射概形的闭子预概形，故为仿射概形
（\hyperref[I.4.2.3-fr]{4.2.3}）。于是命题由 $\Gamma_f$ 是一个浸入
这一事实得出。

\begin{examples}[5.5.11]
\label{I.5.5.11-fr}
例（\hyperref[I.2.3.2-fr]{2.3.2}）中的预概形
（“域 $K$ 上的射影直线”）是\emph{分离的}，因为对于 $X$ 的仿射开覆盖
$(X_1,X_2)$，$X_1\cap X_2=U_{12}$ 是仿射的，并且
$\Gamma(U_{12},\mathscr{O}_X)$，即形如 $f(s)/s^m$、其中
$f\in K[s]$ 的有理分式所成的环，由 $K[s]$ 和 $1/s$ 生成；故
（\hyperref[I.5.5.6-fr]{5.5.6}）的条件成立。

仍像例（\hyperref[I.2.3.2-fr]{2.3.2}）中那样选择 $X_1$、$X_2$、
$U_{12}$ 和 $U_{21}$，但这次取 $u_{12}$ 为把 $f(s)$ 对应到 $f(t)$ 的
同构；由粘合得到一个\emph{整的非分离}预概形 $X$，因为
（\hyperref[I.5.5.6-fr]{5.5.6}）的第一个条件成立，而第二个条件不成立。
这里立即可见，
$\Gamma(X,\mathscr{O}_X)\to\Gamma(X_1,\mathscr{O}_X)=K[s]$ 是一个
同构；其逆同构定义一个满射态射
$f:X\to\operatorname{Spec}(K[s])$，并且对于每个
% EGA1-ERR-0045: the printed ideal (0) is retained; the doubled-origin fibre is over (s), while (0) is generic.
$y\in\operatorname{Spec}(K[s])$，若 $\mathfrak{j}_y\neq(0)$，则
$f^{-1}(y)$ 缩为一点；但当 $\mathfrak{j}_y=(0)$ 时，$f^{-1}(y)$ 由
\emph{两个}不同的点组成（称 $X$ 为“把点 $0$ 加倍的域 $K$ 上的仿射直线”）。

% EGA1-ERR-0046: the printed claim that neither condition holds is retained; direct algebra shows only affineness of the overlap fails.
也可以给出（\hyperref[I.5.5.6-fr]{5.5.6}）的两个条件\emph{都不成立}的
例子。先注意，在域 $K$ 上的二元多项式环 $A=K[s,t]$ 的素谱 $Y$ 中，
开集 $U$，即 $D(s)$ 与 $D(t)$ 的并，\emph{不是仿射开集}。事实上，若 $z$ 是
$\mathscr{O}_Y$ 在 $U$ 上的一个截面，则存在两个整数 $m\geq0$、
$n\geq0$，使 $s^m z$ 和 $t^n z$ 是 $s$、$t$ 的多项式在 $U$ 上的限制
（\hyperref[I.1.4.1-fr]{1.4.1}）；这显然只有在截面 $z$ 能延拓为整个
$Y$ 上的一个截面，并等同于 $s$、$t$ 的一个多项式时才可能。若 $U$
是仿射开集，则嵌入态射 $U\to Y$ 将是一个同构
（\hyperref[I.1.7.3-fr]{1.7.3}），这与 $U\neq Y$ 矛盾。

在此基础上，取两个仿射概形 $Y_1$、$Y_2$，它们分别是环
$A_1=K[s_1,t_1]$、$A_2=K[s_2,t_2]$ 的素谱；取
$U_{12}=D(s_1)\cup D(t_1)$、$U_{21}=D(s_2)\cup D(t_2)$，并取
$u_{12}$ 为同构 $Y_2\to Y_1$ 在 $U_{21}$ 上的限制；该同构对应于把
$f(s_1,t_1)$ 对应到 $f(s_2,t_2)$ 的环同构。这样就得到一个
% EGA1-ERR-0046: source-literal conclusion retained; the second generation condition is in fact satisfied.
（\hyperref[I.5.5.6-fr]{5.5.6}）的两个条件都不成立的例子（所得整预概形
称为“把点 $0$ 加倍的域 $K$ 上的仿射平面”）。
\end{examples}

\begin{remark}[5.5.12]
\label{I.5.5.12-fr}
\normalfont 给定预概形态射的一个性质 \textbf{P}，考虑下列命题：
\begin{enumerate}
  \item[\rm(i)] \emph{每个闭浸入都具有性质 \textbf{P}。}
  \item[\rm(ii)] \emph{两个具有性质 \textbf{P} 的态射之复合具有性质
    \textbf{P}。}
  \item[\rm(iii)] \emph{若 $f:X\to X'$、$g:Y\to Y'$ 是两个具有性质
    \textbf{P} 的 $S$-态射，则 $f\times_S g$ 具有性质 \textbf{P}。}

\oldpage[I]{140}
  \item[\rm(iv)] \emph{若 $f:X\to Y$ 是一个具有性质 \textbf{P} 的
    $S$-态射，则由基预概形的扩张 $S'\to S$ 得到的每个 $S'$-态射
    $f_{(S')}$ 都具有性质 \textbf{P}。}
  \item[\rm(v)] \emph{若两个态射 $f:X\to Y$、$g:Y\to Z$ 的复合
    $g\circ f$ 具有性质 \textbf{P}，并且 $g$ 分离，则 $f$ 具有性质
    \textbf{P}。}
  \item[\rm(vi)] \emph{若态射 $f:X\to Y$ 具有性质 \textbf{P}，则
    $f_{\mathrm{red}}$（\hyperref[I.5.1.5-fr]{5.1.5}）也具有性质
    \textbf{P}。}
\end{enumerate}

在这些条件下，若假设 (i) 和 (ii) 成立，则 (iii) 与 (iv)
\emph{等价}，而 (v) 与 (vi) 都是 (i)、(ii) 和 (iii) 的\emph{推论}。

第一个断言已经证明（\hyperref[I.3.5.1-fr]{3.5.1}）。考虑
（\hyperref[I.5.3.13-fr]{5.3.13}）中 $f$ 的分解
$X\xrightarrow{\Gamma_f}X\times_Z Y\xrightarrow{p_2}Y$；关系
$p_2=(g\circ f)\times_Z 1_Y$ 表明，若 $g\circ f$ 具有性质 \textbf{P}，
则由 (iii)，$p_2$ 也具有性质 \textbf{P}；若 $g$ 分离，则
$\Gamma_f$ 是一个闭浸入（\hyperref[I.5.4.3-fr]{5.4.3}），故由 (i)
它也具有性质 \textbf{P}；最后，由 (ii)，$f$ 具有性质 \textbf{P}。

最后，考虑交换图表
\[
  \xymatrix{
    X_{\mathrm{red}}\ar[r]^{f_{\mathrm{red}}}\ar[d] &
    Y_{\mathrm{red}}\ar[d]\\
    X\ar[r]_f &
    Y
  }
\]
其中竖直箭头都是闭浸入（\hyperref[I.5.1.5-fr]{5.1.5}），故由 (i)
具有性质 \textbf{P}。$f$ 具有性质 \textbf{P} 这一假设，由 (ii) 推出
$X_{\mathrm{red}}\xrightarrow{f_{\mathrm{red}}}Y_{\mathrm{red}}\to Y$
具有性质 \textbf{P}；最后，由于闭浸入是分离的
（\hyperref[I.5.5.1-fr]{5.5.1, (i)}），由 (v) 可知
$f_{\mathrm{red}}$ 具有性质 \textbf{P}。

应注意，若考虑命题：
\begin{enumerate}
  \item[\rm(i')] \emph{每个浸入都具有性质 \textbf{P}。}
  \item[\rm(v')] \emph{若 $g\circ f$ 具有性质 \textbf{P}，则 $f$
    也具有性质 \textbf{P}；}
\end{enumerate}
则上述论证表明，(v') 是 (i')、(ii) 和 (iii) 的推论。
\end{remark}

\begin{env}[5.5.13]
\label{I.5.5.13-fr}
应注意，(v) 与 (vi) 仍然是 (i)、(iii) 和
\begin{enumerate}
  \item[\rm(ii')] \emph{若 $j:X\to Y$ 是一个闭浸入，$g:Y\to Z$
    是一个具有性质 \textbf{P} 的态射，则 $g\circ j$ 具有性质
    \textbf{P}。}
\end{enumerate}
的推论。同样，(v') 是 (i')、(iii) 和
\begin{enumerate}
  \item[\rm(ii'')] \emph{若 $j:X\to Y$ 是一个浸入，$g:Y\to Z$
    是一个具有性质 \textbf{P} 的态射，则 $g\circ j$ 具有性质
    \textbf{P}。}
\end{enumerate}
的推论。事实上，这立即由（\hyperref[I.5.5.12-fr]{5.5.12}）中的
论证得出。
\end{env}
\section{有限性条件}
\label{section:I.6-fr}

\subsection{Noether预概形与局部Noether预概形。}
\label{subsection:I.6.1-fr}

\begin{definition}[6.1.1]
\label{I.6.1.1-fr}
称预概形 $X$ 是\emph{Noether的}（相应地，\emph{局部Noether的}），如果它是
仿射开集 $V_\alpha$ 的有限并（相应地，并），并且在每个 $V_\alpha$
上诱导的概形之环都是Noether环。
\end{definition}

由（\hyperref[I.1.5.2-fr]{1.5.2}）立即可知，若 $X$ 局部Noether，则结构层
$\mathscr{O}_X$ 是一个\emph{凝聚环层}，因为这是一个局部问题。一个
\emph{凝聚的} $\mathscr{O}_X$-模 $\mathscr{F}$ 的\emph{任意拟凝聚
$\mathscr{O}_X$-子}

\oldpage[I]{141}
\emph{模（相应地，任意拟凝聚的 $\mathscr{O}_X$-商模）}因而都是
\emph{凝聚的}；因为问题仍是局部的，只需应用
（\hyperref[I.1.5.1-fr]{1.5.1}）、
（\hyperref[I.1.4.1-fr]{1.4.1}）和
（\hyperref[I.1.3.10-fr]{1.3.10}），再结合Noether环上有限型模的子模
（相应地，商模）仍是有限型模这一事实。特别地，$\mathscr{O}_X$ 的
\emph{任意拟凝聚理想层}都是\emph{凝聚的}。

若预概形 $X$ 是开集 $W_\lambda$ 的有限并（相应地，并），并且在这些
$W_\lambda$ 上诱导的预概形都是Noether的（相应地，局部Noether的），则显然
$X$ 是Noether的（相应地，局部Noether的）。

\begin{proposition}[6.1.2]
\label{I.6.1.2-fr}
预概形 $X$ 为Noether的充要条件是它局部Noether且其底空间拟紧。这时 $X$ 的
底空间是Noether空间。
\end{proposition}

\begin{proof}
第一个断言立即由定义和（\hyperref[I.1.1.10-fr]{1.1.10 (ii)}）得出。
第二个断言由（\hyperref[I.1.1.6-fr]{1.1.6}）以及有限个Noether子空间的并
仍是Noether空间这一事实（\hyperref[0.2.2.3-fr]{0, 2.2.3}）得出。
\end{proof}

\begin{proposition}[6.1.3]
\label{I.6.1.3-fr}
设 $X$ 是环 $A$ 的仿射概形。下列条件等价：a) $X$ 是Noether的；
b) $X$ 是局部Noether的；c) $A$ 是Noether环。
\end{proposition}

\begin{proof}
a) 与 b) 的等价性由（\hyperref[I.6.1.2-fr]{6.1.2}）以及每个仿射概形
都有拟紧底空间这一事实（\hyperref[I.1.1.10-fr]{1.1.10}）得出；此外，
c) 蕴含 a) 是显然的。为证明 a) 蕴含 c)，注意到 $X$ 有一个由仿射开集
组成的有限覆盖 $(V_i)$，使得在 $V_i$ 上诱导的预概形之环 $A_i$ 是Noether环。
现设 $(\mathfrak{a}_n)$ 是 $A$ 的一个理想升链；由典范的一一对应
（\hyperref[I.1.3.7-fr]{1.3.7}），它对应于
$\widetilde{A}\equiv\mathscr{O}_X$ 中的一个理想层升链
$(\widetilde{\mathfrak{a}}_n)$。要证明序列 $(\mathfrak{a}_n)$ 稳定，只需
证明序列 $(\widetilde{\mathfrak{a}}_n)$ 稳定。现在，限制
$\widetilde{\mathfrak{a}}_n|V_i$ 是 $\mathscr{O}_X|V_i$ 中的一个拟凝聚
理想层，因为它是 $\widetilde{\mathfrak{a}}_n$ 沿典范浸入 $V_i\to X$
的逆像（\hyperref[0.5.1.4-fr]{0, 5.1.4}）；因此
$\widetilde{\mathfrak{a}}_n|V_i$ 形如 $\widetilde{\mathfrak{a}}_{ni}$，
其中 $\mathfrak{a}_{ni}$ 是 $A_i$ 的一个理想
（\hyperref[I.1.3.7-fr]{1.3.7}）。由于 $A_i$ 是Noether环，序列
$(\mathfrak{a}_{ni})$ 对每个 $i$ 都稳定，命题得证。
\end{proof}

应注意，上述论证还证明了：\emph{若 $X$ 是Noether预概形，则
$\mathscr{O}_X$ 的凝聚理想层的每个升链都稳定}。

\begin{proposition}[6.1.4]
\label{I.6.1.4-fr}
Noether（相应地，局部Noether）预概形的每个子预概形都是Noether的（相应地，
局部Noether的）。
\end{proposition}

\begin{proof}
只需对Noether预概形 $X$ 给出证明；此外，由定义
（\hyperref[I.6.1.1-fr]{6.1.1}），立即归结到 $X$ 为仿射概形的情形。
由于 $X$ 的每个子预概形都是某个开集上诱导的预概形的闭子预概形
（\hyperref[I.4.1.3-fr]{4.1.3}），故只需考虑子预概形 $Y$ 为闭子预概形
或由 $X$ 的一个开集诱导的情形。当 $Y$ 为闭子预概形时，结论是立即的：
若 $A$ 是 $X$ 的环，则已知 $Y$ 是环 $A/\mathfrak{J}$ 的仿射概形，其中
$\mathfrak{J}$ 是 $A$ 的一个理想（\hyperref[I.4.2.3-fr]{4.2.3}）；因为
$A$ 是Noether环（\hyperref[I.6.1.3-fr]{6.1.3}），$A/\mathfrak{J}$ 也是
Noether环。

现设 $Y$ 在 $X$ 中为开集；$Y$ 的底空间是Noether空间
（\hyperref[I.6.1.2-fr]{6.1.2}），因而拟紧，所以是有限个开集 $D(f_i)$
（$f_i\in A$）的并；
\oldpage[I]{142}
问题归结到 $Y=D(f)$ 且 $f\in A$ 的情形。但此时 $Y$ 是一个仿射概形，
其环同构于 $A_f$（\hyperref[I.1.3.6-fr]{1.3.6}）；由于 $A$ 是Noether环
（\hyperref[I.6.1.3-fr]{6.1.3}），$A_f$ 也是Noether环。
\end{proof}

\begin{env}[6.1.5]
\label{I.6.1.5-fr}
应注意，两个 Noether $S$-预概形的\emph{积}未必是 Noether 的，即使这些预概形
都是仿射的；这是因为两个 Noether 代数的张量积未必是 Noether 环
（参见（\hyperref[I.6.3.8-fr]{6.3.8}））。
\end{env}

\begin{proposition}[6.1.6]
\label{I.6.1.6-fr}
若 $X$ 是 Noether 预概形，则 $\mathscr{O}_X$ 的幂零根
$\mathscr{N}_X$ 是幂零的。
\end{proposition}

事实上，可以用有限个仿射开集 $U_i$ 覆盖 $X$，并且只须证明存在整数
$n_i$ 使得 $(\mathscr{N}_X|U_i)^{n_i}=0$；若 $n$ 是各 $n_i$ 中最大的，
则有 $\mathscr{N}_X^n=0$。因此归结到
$X=\operatorname{Spec}(A)$ 为仿射的情形，其中 $A$ 是 Noether 环；由
（\hyperref[I.5.1.1-fr]{5.1.1}）和
（\hyperref[I.1.3.13-fr]{1.3.13}），只须注意到 $A$ 的幂零根是幂零的
（\cite{I-11}, p.~127, cor.~4）。

\begin{corollary}[6.1.7]
\label{I.6.1.7-fr}
设 $X$ 是 Noether 预概形；$X$ 为仿射概形的充要条件是
$X_{\mathrm{red}}$ 为仿射概形。
\end{corollary}

这由（\hyperref[I.6.1.6-fr]{6.1.6}）和
（\hyperref[I.5.1.10-fr]{5.1.10}）得出。

\begin{lemma}[6.1.8]
\label{I.6.1.8-fr}
设 $X$ 是拓扑空间，$x$ 是 $X$ 的一点，$U$ 是 $x$ 的一个只有有限多个
不可约分支的开邻域。则存在 $x$ 的一个邻域 $V$，使得包含于 $V$ 中的
$x$ 的每个开邻域都是连通的。
\end{lemma}

设 $U_i$（$1\leq i\leq m$）是 $U$ 中不含 $x$ 的不可约分支；这些
$U_i$ 之并在 $U$ 中的补集是 $x$ 在 $U$ 中的一个开邻域 $V$，因而后者在
$X$ 中也为开邻域；
% EGA1-ERR-0048: the next printed/source-literal equality is retained diplomatically; see ERRATA_CANDIDATES.jsonl.
此外，它也是 $X$ 中不含 $x$ 的各不可约分支之并在 $X$ 中的补集
（\hyperref[0.2.1.6-fr]{0, 2.1.6}）。现设 $W$ 是包含于 $V$ 中的 $x$ 的一个
开邻域。$W$ 的不可约分支是 $U$ 中与 $W$ 相交的各不可约分支在 $W$
上的迹（\hyperref[0.2.1.6-fr]{0, 2.1.6}），所以这些分支都含有 $x$；由于
它们是连通的，$W$ 也连通。

\begin{corollary}[6.1.9]
\label{I.6.1.9-fr}
局部 Noether 拓扑空间是局部连通的（这尤其蕴含它的连通分支都是开集）。
\end{corollary}

\begin{proposition}[6.1.10]
\label{I.6.1.10-fr}
设 $X$ 是局部 Noether 拓扑空间。下列条件等价：
\begin{enumerate}
  \item[a)] $X$ 的不可约分支都是开集。
  \item[b)] $X$ 的不可约分支与其连通分支相同。
  \item[c)] $X$ 的连通分支都是不可约的。
  \item[d)] $X$ 的两个不同不可约分支互不相交。
\end{enumerate}

最后，若 $X$ 是预概形，则这些条件还等价于：
\begin{enumerate}
  \item[e)] 对每个 $x\in X$，$\operatorname{Spec}(\mathscr{O}_x)$
    都是不可约的（换言之，$\mathscr{O}_x$ 的幂零根是素理想）。
\end{enumerate}
\end{proposition}

a) 蕴含 b) 是立即的，因为不可约空间是连通的，并且 a) 蕴含 $X$ 的
不可约分支都是既开又闭的集合。b) 蕴含 c) 是显然的；反过来，若闭集
$F$ 包含
\oldpage[I]{143}
$X$ 的一个连通分支 $C$ 且不等于 $C$，则该闭集不可能不可约，
因为这个集合不连通，故它是两个互不相交的非空集合之并，这两个集合在
$F$ 中既开又闭，因而在 $X$ 中闭；所以 c) 蕴含 b)。由此立即得出 c)
蕴含 d)，因为两个不同连通分支没有公共点。

到此为止，我们还没有用到 $X$ 局部 Noether 这一事实。现假设这一条件成立，
并证明 d) 蕴含 a)：由（\hyperref[0.2.1.6-fr]{0, 2.1.6}），可以归结到
空间 $X$ 为 Noether 的情形，因而它只有有限多个不可约分支。由于这些分支
都是闭集并且两两不交，它们都是开集。

最后，即使不假设预概形 $X$ 的底空间局部 Noether，d) 与 e) 仍然等价。
事实上，由（\hyperref[0.2.1.6-fr]{0, 2.1.6}），可以归结到
$X=\operatorname{Spec}(A)$ 为仿射的情形；说 $x$ 只包含在 $X$ 的一个
不可约分支中，这时意味着 $\mathfrak{j}_x$ 只包含 $A$ 的一个极小素理想
（\hyperref[I.1.1.14-fr]{1.1.14}），这等价于说
$\mathfrak{j}_x\mathscr{O}_x$ 只包含 $\mathscr{O}_x$ 的一个极小素理想，
结论由此得出。

\begin{corollary}[6.1.11]
\label{I.6.1.11-fr}
设 $X$ 是局部 Noether 空间。$X$ 不可约的充要条件是：$X$ 连通且非空，并且
$X$ 的两个不同不可约分支互不相交。若 $X$ 是预概形，则最后一个条件
等价于对每个 $x\in X$，$\operatorname{Spec}(\mathscr{O}_x)$ 都不可约。
\end{corollary}

最后一个断言已在（\hyperref[I.6.1.10-fr]{6.1.10}）中证明；因此只须证明
第一个断言中诸条件的充分性。但由
（\hyperref[I.6.1.10-fr]{6.1.10}），这些条件蕴含 $X$ 的不可约分支就是
它的连通分支；又因为 $X$ 连通且非空，所以它是不可约的。

% EGA1-ERR-0049: the next printed/source-literal criterion omits nonemptiness; see ERRATA_CANDIDATES.jsonl.
\begin{corollary}[6.1.12]
\label{I.6.1.12-fr}
设 $X$ 是局部 Noether 预概形。$X$ 为整的充要条件是：$X$ 连通，并且对每个
$x\in X$，$\mathscr{O}_x$ 都是整环。
\end{corollary}

\begin{proposition}[6.1.13]
\label{I.6.1.13-fr}
设 $X$ 是局部 Noether 预概形，$x\in X$ 是一点，并且 $\mathscr{O}_x$ 的
幂零根 $\mathscr{N}_x$ 是素理想（相应地，$\mathscr{O}_x$ 是约化环，相应地，
是整环）；则存在 $x$ 的一个开邻域 $U$，它是不可约的
（相应地，约化的，相应地，整的）。
\end{proposition}

只须考虑 $\mathscr{N}_x$ 是素理想以及 $\mathscr{N}_x=0$ 这两种情形，
因为第三个假设是前两个假设的合取。若 $\mathscr{N}_x$ 是素理想，则
$x$ 只属于 $X$ 的一个不可约分支 $Y$
（\hyperref[I.6.1.10-fr]{6.1.10}）；$X$ 中不含 $x$ 的各不可约分支之并
是闭集（因为这些分支所成的族局部有限），故该并的补集 $U$ 是开集且
包含于 $Y$ 中，因而不可约
（\hyperref[0.2.1.6-fr]{0, 2.1.6}）。若 $\mathscr{N}_x=0$，则在 $x$
的某个邻域中，对每个 $y$ 也有 $\mathscr{N}_y=0$，因为
$\mathscr{N}$ 是拟凝聚的（\hyperref[I.5.1.1-fr]{5.1.1}），并且由于
$X$ 局部 Noether，它因而是凝聚的；结论于是由
（\hyperref[0.5.2.2-fr]{0, 5.2.2}）得出。

\subsection{Artin预概形。}
\label{subsection:I.6.2-fr}

\begin{definition}[6.2.1]
\label{I.6.2.1-fr}
称一个预概形是\emph{Artin的}，如果它是仿射的，并且它的环是Artin环。
\end{definition}

\oldpage[I]{144}
\begin{proposition}[6.2.2]
\label{I.6.2.2-fr}
给定一个预概形 $X$，下列条件等价：
\begin{enumerate}
  \item[a)] $X$ 是一个Artin概形；
  \item[b)] $X$ 是Noether的，并且其底空间是离散的；
  \item[c)] $X$ 是Noether的，并且其底空间的点都是闭点
    （条件 $\mathrm{T}_1$）。
\end{enumerate}

当这些条件成立时，$X$ 的底空间是有限的，并且 $X$ 的环 $A$ 是 $X$
各点的局部（Artin）环的直积。
\end{proposition}

已知 a) 蕴含最后一个断言（\cite{I-13}, p.~205, th.~3）；此时 $A$ 的每个
素理想都是极大理想，并且是 $A$ 的某个局部分量之极大理想的逆像，因而空间
$X$ 是有限离散的；所以 a) 蕴含 b)，而 b) 显然蕴含 c)。为了证明 c) 蕴含
a)，先证明此时 $X$ 是有限的；事实上，可以归结到 $X$ 为仿射的情形，而我们
知道，所有素理想都是极大理想的Noether环是Artin环（\cite{I-13}, p.~203），
这就证明了该断言。于是底空间 $X$ 是离散的，是有限多个点 $x_i$ 的拓扑和，
并且局部环 $\mathscr{O}_{x_i}=A_i$ 都是Artin环；显然，$X$ 同构于环 $A$
的素谱这一仿射概形，其中该环是各 $A_i$ 的直积
（\hyperref[I.1.7.3-fr]{1.7.3}）。

\subsection{有限型态射。}
\label{subsection:I.6.3-fr}

\begin{definition}[6.3.1]
\label{I.6.3.1-fr}
称态射 $f:X\to Y$ 是\emph{有限型的}，如果 $Y$ 是一族仿射开集
$(V_\alpha)$ 的并，并且这些开集具有下列性质：
\begin{enumerate}
  \item[(P)] $f^{-1}(V_\alpha)$ 是有限多个仿射开集 $U_{\alpha i}$ 的并，
    并且每个环 $A(U_{\alpha i})$ 都是 $A(V_\alpha)$ 上的有限型代数。
\end{enumerate}
这时也称 $X$ 是 $Y$ 上的有限型预概形，或有限型 $Y$-预概形。
\end{definition}

\begin{proposition}[6.3.2]
\label{I.6.3.2-fr}
若 $f:X\to Y$ 是有限型态射，则 $Y$ 的每个仿射开集 $W$ 都具有定义
（\hyperref[I.6.3.1-fr]{6.3.1}）中的性质 (P)。
\end{proposition}

我们先证明下述引理。

\begin{lemma}[6.3.2.1]
\label{I.6.3.2.1-fr}
若 $T\subset Y$ 是具有性质 (P) 的仿射开集，则对每个 $g\in A(T)$，
$D(g)$ 也具有性质 (P)。
\end{lemma}

事实上，依假设，$f^{-1}(T)$ 是有限多个仿射开集 $Z_j$ 的并，并且
$A(Z_j)$ 是 $A(T)$ 上的有限型代数；设
$\varphi_j:A(T)\to A(Z_j)$ 是与 $f$ 在 $Z_j$ 上的限制相对应的环同态
（\hyperref[I.2.2.4-fr]{2.2.4}），并令 $g_j=\varphi_j(g)$；于是
$f^{-1}(D(g))\cap Z_j=D(g_j)$
（\hyperref[I.1.2.2.2-fr]{1.2.2.2}）。而
$A(D(g_j))=A(Z_j)_{g_j}=A(Z_j)[1/g_j]$ 是 $A(Z_j)$ 上的有限型代数，
由假设，\emph{尤其}是 $A(T)$ 上的有限型代数，因而也是
$A(D(g))=A(T)[1/g]$ 上的有限型代数；引理得证。

% EGA1-ERR-0050: published EGA II erratum inserts D(g_i)\subset W here; the diplomatic Chinese remains source-literal.
既已证明该引理，又因为 $W$ 拟紧
（\hyperref[I.1.1.10-fr]{1.1.10}），所以 $W$ 存在一个由形如
$D(g_i)$ 的集合组成的有限覆盖，其中每个 $g_i$ 都属于某个环
$A(V_{\alpha(i)})$。每个 $D(g_i)$ 都拟紧，因而是有限多个集合
$D(h_{ik})$ 的并，其中 $h_{ik}\in A(W)$；若
$\varphi_i:A(W)\to A(D(g_i))$ 是典范映射，则由
（\hyperref[I.1.2.2.2-fr]{1.2.2.2}）有
$D(h_{ik})=D(\varphi_i(h_{ik}))$。由
（\hyperref[I.6.3.2.1-fr]{6.3.2.1}），每个
$f^{-1}(D(h_{ik}))$ 都有一个由有限多个仿射开集 $U_{ijk}$ 组成的覆盖，
并且 $A(U_{ijk})$ 是
$A(D(h_{ik}))=A(W)[1/h_{ik}]$ 上的有限型代数；命题得证。

\oldpage[I]{145}
因此可以说，$Y$ 上有限型预概形这一概念\emph{对 $Y$ 是局部的}。

\begin{proposition}[6.3.3]
\label{I.6.3.3-fr}
设 $X$、$Y$ 是两个仿射概形；$X$ 在 $Y$ 上为有限型的充要条件是
$A(X)$ 为 $A(Y)$ 上的有限型代数。
\end{proposition}

该条件显然是充分的；现证明其必要性。令 $A=A(Y)$、$B=A(X)$；由
（\hyperref[I.6.3.2-fr]{6.3.2}），$X$ 有一个有限仿射开覆盖 $(V_i)$，
使得每个环 $A(V_i)$ 都是 $A$ 上的有限型代数。此外，由于 $V_i$ 拟紧，
可以用有限多个形如 $D(g_{ij})\subset V_i$ 的开集覆盖其中每一个，其中
$g_{ij}\in B$；若 $\varphi_i$ 是与典范浸入 $V_i\to X$ 相对应的同态
$B\to A(V_i)$，则
\[
  B_{g_{ij}}=(A(V_i))_{\varphi_i(g_{ij})}
  =A(V_i)[1/\varphi_i(g_{ij})],
\]
所以 $B_{g_{ij}}$ 是 $A$ 上的有限型代数。因此可以归结到
$V_i=D(g_i)$ 且 $g_i\in B$ 的情形。依假设，存在 $B$ 的有限子集 $F_i$
和整数 $n_i\geq 0$，使得 $B_{g_i}$ 是由 $F_i$ 中各 $b_i$ 所对应的元素
$b_i/g_i^{n_i}$ 在 $A$ 上生成的代数。由于 $g_i$ 只有有限多个，不妨还可
假设所有 $n_i$ 都等于同一个整数 $n$。此外，由于 $D(g_i)$ 构成 $X$
的覆盖，$g_i$ 在 $B$ 中生成的理想等于 $B$；换言之，存在 $h_i\in B$，
使得 $\sum_i h_i g_i=1$。现设 $F$ 是 $B$ 的有限子集，它是各 $F_i$、
各 $g_i$ 组成的集合和各 $h_i$ 组成的集合之并；我们证明 $B$ 的子环
$B'=A[F]$
等于 $B$。依假设，对每个 $b\in B$ 和每个 $i$，$b$ 在 $B_{g_i}$ 中的
典范像都形如 $b'_i/g_i^{m_i}$，其中 $b'_i\in B'$；将 $b'_i$ 乘以
$g_i$ 的适当幂以后，不妨仍可假设所有 $m_i$ 都等于同一个整数 $m$。
由分式环的定义，因而存在一个依赖于 $b$ 的整数 $N$，使得 $N\geq m$，
并且对每个 $i$ 都有 $g_i^N b\in B'$；但\emph{在环 $B'$ 中}，各
$g_i^N$ 生成的理想就是 $B'$，因为各 $g_i$ 已具有这一性质（各 $h_i$
都属于 $B'$）；所以存在 $c_i\in B'$，使得
$\sum_i c_i g_i^N=1$，从而 $b=\sum_i c_i g_i^N b\in B'$，证毕。

\begin{proposition}[6.3.4]
\label{I.6.3.4-fr}
\begin{enumerate}
  \item[(i)] 每个闭浸入都是有限型的。
  \item[(ii)] 两个有限型态射的复合是有限型的。
  \item[(iii)] 若 $f:X\to X'$、$g:Y\to Y'$ 是两个有限型 $S$-态射，
    则 $f\times_S g$ 是有限型的。
  \item[(iv)] 若 $f:X\to Y$ 是有限型 $S$-态射，则对于基预概形的任意扩张
    $g:S'\to S$，$f_{(S')}$ 都是有限型的。
  \item[(v)] 若两个态射的复合 $g\circ f$ 是有限型的，并且 $g$ 是分离的，
    则 $f$ 是有限型的。
  \item[(vi)] 若态射 $f$ 是有限型的，则 $f_{\mathrm{red}}$ 也是有限型的。
\end{enumerate}
\end{proposition}

由（\hyperref[I.5.5.12-fr]{5.5.12}），只需证明 (i)、(ii) 和 (iv)。

为证明 (i)，可以只考虑典范浸入 $X\to Y$ 的情形，其中 $X$ 是 $Y$ 的
闭子预概形；此外，由（\hyperref[I.6.3.2-fr]{6.3.2}），可设 $Y$ 为
仿射的。这时 $X$ 也为仿射的（\hyperref[I.4.2.3-fr]{4.2.3}），并且其环
同构于商环 $A/\mathfrak{J}$，其中 $A$ 是 $Y$ 的环，而
$\mathfrak{J}$ 是 $A$ 的一个理想；由于 $A/\mathfrak{J}$ 是 $A$ 上的
有限型代数，结论随即成立。

现证明 (ii)。设 $f:X\to Y$、$g:Y\to Z$ 是两个有限型态射，并设 $U$
是 $Z$ 的一个仿射开集；$g^{-1}(U)$ 有一个由有限多个仿射开集 $V_i$
组成的覆盖，使得 $A(V_i)$ 是 $A(U)$ 上的有限型代数
（\hyperref[I.6.3.2-fr]{6.3.2}）；同样，

\oldpage[I]{146}
每个 $f^{-1}(V_i)$ 都有一个由有限多个仿射开集 $W_{ij}$ 组成的覆盖，
使得 $A(W_{ij})$ 是 $A(V_i)$ 上的有限型代数，因而也是 $A(U)$ 上的
有限型代数；结论由此成立。

最后，为证明 (iv)，可以只考虑 $S=Y$ 的情形；事实上，把 $f$ 看作
$Y$-态射，并取基扩张 $Y_{(S')}\to Y$，则 $f_{(S')}$ 也等于
$f_{(Y_{(S')})}$（\hyperref[I.3.3.9-fr]{3.3.9}）。令 $p$、$q$ 分别为投影
$X_{(S')}\to X$ 和 $X_{(S')}\to S'$。设 $V$ 是 $S$ 中的一个仿射开集；
$f^{-1}(V)$ 是有限多个仿射开集 $W_i$ 的并，并且每个 $A(W_i)$ 都是
$A(V)$ 上的有限型代数（\hyperref[I.6.3.2-fr]{6.3.2}）。设 $V'$ 是
$S'$ 中包含于 $g^{-1}(V)$ 的一个仿射开集；由于 $f\circ p=g\circ q$，
$q^{-1}(V')$ 包含于各 $p^{-1}(W_i)$ 的并；另一方面，交集
$p^{-1}(W_i)\cap q^{-1}(V')$ 可与纤维积 $W_i\times_V V'$ 认同
（\hyperref[I.3.2.7-fr]{3.2.7}），后者是一个仿射概形，其环同构于
$A(W_i)\otimes_{A(V)}A(V')$（\hyperref[I.3.2.2-fr]{3.2.2}）；依假设，
这个环是 $A(V')$ 上的有限型代数，命题得证。

\begin{corollary}[6.3.5]
\label{I.6.3.5-fr}
设 $f:X\to Y$ 是一个浸入态射。若 $Y$（相应地，$X$）的底空间是
局部 Noether 的（相应地，Noether 的），则 $f$ 是有限型的。
\end{corollary}

总可以假设 $Y$ 是仿射的（\hyperref[I.6.3.2-fr]{6.3.2}）；若 $Y$ 的
底空间是局部 Noether 的，还可以假设它是 Noether 的，于是其子空间
$X$ 的底空间也是 Noether 的。换言之，可设 $Y$ 为仿射的，并且 $X$
的底空间为 Noether 的；这时 $X$ 有一个由有限多个仿射开集
$D(g_i)\subset Y$ 组成的覆盖，其中 $g_i\in A(Y)$，并且
$X\cap D(g_i)$ 在 $D(g_i)$ 中闭（因而是仿射概形
（\hyperref[I.4.2.3-fr]{4.2.3}）），这是因为 $X$ 在 $Y$ 中局部闭
（\hyperref[I.4.1.3-fr]{4.1.3}）。于是由
（\hyperref[I.6.3.4-fr]{6.3.4, (i)}）和
（\hyperref[I.6.3.3-fr]{6.3.3}），$A(X\cap D(g_i))$ 是
$A(D(g_i))$ 上的有限型代数，而
$A(D(g_i))=A(Y)_{g_i}=A(Y)[1/g_i]$ 是 $A(Y)$ 上的有限型代数，证明完毕。

\begin{corollary}[6.3.6]
\label{I.6.3.6-fr}
设 $f:X\to Y$、$g:Y\to Z$ 是两个态射。若 $g\circ f$ 是有限型的，
并且 $X$ 是 Noether 的，或 $X\times_Z Y$ 是局部 Noether 的，则 $f$
是有限型的。
\end{corollary}

这立即由（\hyperref[I.5.5.12-fr]{5.5.12}）的证明以及将
（\hyperref[I.6.3.5-fr]{6.3.5}）应用于浸入态射 $\Gamma_f$ 得出。

\begin{proposition}[6.3.7]
\label{I.6.3.7-fr}
设 $f:X\to Y$ 是有限型态射；若 $Y$ 是 Noether 的（相应地，局部
Noether 的），则 $X$ 是 Noether 的（相应地，局部 Noether 的）。
\end{proposition}

只需在 $Y$ 为 Noether 的情形下证明。此时 $Y$ 是有限多个仿射开集
$V_i$ 的并，并且各 $A(V_i)$ 都是 Noether 环。由
（\hyperref[I.6.3.2-fr]{6.3.2}），每个 $f^{-1}(V_i)$ 都是有限多个
仿射开集 $W_{ij}$ 的并，并且各 $A(W_{ij})$ 都是 $A(V_i)$ 上的有限型
代数，因而是 Noether 环；这就证明了 $X$ 是 Noether 的。

\begin{corollary}[6.3.8]
\label{I.6.3.8-fr}
设 $X$ 是 $S$ 上的有限型预概形。对于基的任意扩张 $S'\to S$，若
$S'$ 是 Noether 的（相应地，局部 Noether 的），则 $X_{(S')}$ 是
Noether 的（相应地，局部 Noether 的）。
\end{corollary}

这由（\hyperref[I.6.3.7-fr]{6.3.7}）得出，因为由
（\hyperref[I.6.3.4-fr]{6.3.4, (iv)}），$X_{(S')}$ 在 $S'$ 上为有限型。

还可以说，在 $S$-预概形的积 $X\times_S Y$ 中，若因子 $X$、$Y$ 中
\emph{一个}
\oldpage[I]{147}
在 $S$ 上\emph{有限型}，而\emph{另一个为 Noether 的}（相应地，
\emph{局部 Noether 的}），则 $X\times_S Y$ 是\emph{Noether 的}
（相应地，\emph{局部 Noether 的}）。

\begin{corollary}[6.3.9]
\label{I.6.3.9-fr}
设 $X$ 是局部 Noether 预概形 $S$ 上的有限型预概形。则每个
$S$-态射 $f:X\to Y$ 都是有限型的。
\end{corollary}

事实上，可设 $S$ 是 Noether 的；若 $\varphi:X\to S$、
$\psi:Y\to S$ 是结构态射，则 $\varphi=\psi\circ f$，并且由
（\hyperref[I.6.3.7-fr]{6.3.7}），$X$ 是 Noether 的；所以由
（\hyperref[I.6.3.6-fr]{6.3.6}），$f$ 是有限型的。

\begin{proposition}[6.3.10]
\label{I.6.3.10-fr}
设 $f:X\to Y$ 是有限型态射。$f$ 为满射的充要条件是：对于每个代数
闭域 $\Omega$，由 $f$ 对应的映射 $X(\Omega)\to Y(\Omega)$
（\hyperref[I.3.4.1-fr]{3.4.1}）都是满射。
\end{proposition}

该条件是充分的；只要对每个 $y\in Y$，取 $k(y)$ 的一个代数闭扩域
$\Omega$，并考虑交换图
\[
  \xymatrix{
    X\ar[dd]_f\\
    & \operatorname{Spec}(\Omega)\ar[ul]\ar[dl]\\
    Y
  }
\]
即可看出（\emph{参见}（\hyperref[I.3.5.3-fr]{3.5.3}））。反过来，
设 $f$ 为满射，并设 $g:\{\xi\}=\operatorname{Spec}(\Omega)\to Y$ 为
一个态射，其中 $\Omega$ 是代数闭域。考虑交换图
\[
  \xymatrix{
    X\ar[d]_f &
    X_{(\Omega)}\ar[l]\ar[d]^{f_{(\Omega)}}\\
    Y &
    \operatorname{Spec}(\Omega)\ar[l]
  }
\]
于是只需证明 $X_{(\Omega)}$ 中存在一个\emph{$\Omega$ 上的有理点}
（\hyperref[I.3.3.14-fr]{3.3.14}、\hyperref[I.3.4.3-fr]{3.4.3} 和
\hyperref[I.3.4.4-fr]{3.4.4}）。由于 $f$ 是满射，$X_{(\Omega)}$ 非空
（\hyperref[I.3.5.10-fr]{3.5.10}）；又由于 $f$ 是有限型的，由
（\hyperref[I.6.3.4-fr]{6.3.4, (iv)}），$f_{(\Omega)}$ 也是有限型的；
因此 $X_{(\Omega)}$ 含有一个非空仿射开集 $Z$，使得 $A(Z)$ 是
$\Omega$ 上的非零有限型代数。由 Hilbert 零点定理 \cite{I-21}，存在一个
$\Omega$-同态 $A(Z)\to\Omega$，也就是 $X_{(\Omega)}$ 在
$\operatorname{Spec}(\Omega)$ 上的一个截面；命题得证。

\subsection{代数预概形。}
\label{subsection:I.6.4-fr}

\begin{definition}[6.4.1]
\label{I.6.4.1-fr}
给定一个域 $K$，称 $K$ 上的有限型预概形 $X$ 为\emph{$K$-代数预概形}；
$K$ 称为 $X$ 的\emph{基域}。若 $X$ 还是一个概形（或者等价地
（\hyperref[I.5.5.8-fr]{5.5.8}），若 $X$ 是一个 $K$-概形），也称 $X$
为\emph{$K$-代数概形}。
\end{definition}

每个 $K$-代数预概形都是\emph{Noether 的}
（\hyperref[I.6.3.7-fr]{6.3.7}）。

\begin{proposition}[6.4.2]
\label{I.6.4.2-fr}
设 $X$ 是一个 $K$-代数预概形。点 $x\in X$ 为闭点的充要条件是
$k(x)$ 为 $K$ 的有限次代数扩张。
\end{proposition}

可设 $X$ 为仿射的，并且 $X$ 的环 $A$ 是有限型 $K$-代数。事实上，
$X$ 中满足 $A(U)$ 为有限型 $K$-代数的仿射开集 $U$ 构成 $X$ 的一个
覆盖（\hyperref[I.6.3.1-fr]{6.3.1}）。于是 $X$ 的闭点就是使
$\mathfrak{j}_x$ 为
\oldpage[I]{148}
$A$ 的极大理想的点，换言之，就是使 $A/\mathfrak{j}_x$ 为域的点
（该域必然等于 $k(x)$）。由于 $A/\mathfrak{j}_x$ 是有限型
$K$-代数，可见若 $x$ 为闭点，则 $k(x)$ 是一个作为 $K$-代数有限型的
域，因而必然是\emph{有限秩} $K$-代数 \cite{I-21}。反过来，若 $k(x)$
在 $K$ 上秩有限，则其子环 $A/\mathfrak{j}_x\subset k(x)$ 也如此；
又因为每个作为 $K$-代数秩有限的整环都是域，所以
$A/\mathfrak{j}_x=k(x)$，从而 $x$ 为闭点。

\begin{corollary}[6.4.3]
\label{I.6.4.3-fr}
设 $K$ 是代数闭域，$X$ 是一个 $K$-代数预概形；则 $X$ 的闭点就是
$K$-有理点（\hyperref[I.3.4.4-fr]{3.4.4}），并可典范地与 $X$ 的
$K$-值点等同。
\end{corollary}

\begin{proposition}[6.4.4]
\label{I.6.4.4-fr}
设 $X$ 是域 $K$ 上的一个代数预概形。下列性质等价：
\begin{enumerate}
  \item[\rm(a)] $X$ 是 Artin 的。
  \item[\rm(b)] $X$ 的底空间是离散的。
  \item[\rm(c)] $X$ 的底空间只有有限多个闭点。
  \item[\rm(c')] $X$ 的底空间是有限的。
  \item[\rm(d)] $X$ 的点都是闭点。
  \item[\rm(e)] $X$ 同构于 $\operatorname{Spec}(A)$，其中 $A$ 是一个
    有限秩 $K$-代数。
\end{enumerate}
\end{proposition}

由于 $X$ 是 Noether 的，由（\hyperref[I.6.2.2-fr]{6.2.2}）可知条件
a)、b)、d) 等价，并且它们蕴含 c) 和 c')；另一方面，e) 显然蕴含 a)。
只需再证明 c) 蕴含 d) 和 e)；可以只考虑 $X$ 为仿射的情形。这时
$A(X)$ 是有限型 $K$-代数（\hyperref[I.6.3.3-fr]{6.3.3}），因而是
Jacobson 环（\cite[p.~3-11 et 3-12]{I-1}），而依假设，其中只有有限多个
极大理想。由于素理想的有限交只有在等于其中之一时才可能仍是素理想，
$A(X)$ 的每个素理想都是极大理想，从而得到 d)。此外，由
（\hyperref[I.6.2.2-fr]{6.2.2}）知道，$A(X)$ 是有限型 Artin
$K$-代数，因而必然\emph{秩有限} \cite{I-21}。

\begin{env}[6.4.5]
\label{I.6.4.5-fr}
当（\hyperref[I.6.4.4-fr]{6.4.4}）中的条件成立时，称 $X$ 是一个
\emph{$K$ 上的有限概形}（\emph{参见}（\hyperref[II.6.1.1-fr]{II, 6.1.1}）），
或一个\emph{有限 $K$-概形}；其\emph{秩}为 $[A:K]$，也记作
$\operatorname{rg}_K(X)$。若 $X$、$Y$ 是两个 $K$ 上的有限概形，则
\[
  \operatorname{rg}_K(X\amalg Y)
  =\operatorname{rg}_K(X)+\operatorname{rg}_K(Y)
  \tag{6.4.5.1}\label{I.6.4.5.1-fr}
\]
\[
  \operatorname{rg}_K(X\times_K Y)
  =\operatorname{rg}_K(X)\operatorname{rg}_K(Y)
  \tag{6.4.5.2}\label{I.6.4.5.2-fr}
\]
这是（\hyperref[I.3.2.2-fr]{3.2.2}）的直接推论。
\end{env}

\begin{corollary}[6.4.6]
\label{I.6.4.6-fr}
设 $X$ 是域 $K$ 上的有限概形。对于 $K$ 的每个扩张 $K'$，
$X\otimes_K K'$ 是 $K'$ 上的有限概形，并且它在 $K'$ 上的秩等于 $X$
在 $K$ 上的秩。
\end{corollary}

事实上，若 $A=A(X)$，则
$[A\otimes_K K':K']=[A:K]$。

\begin{corollary}[6.4.7]
\label{I.6.4.7-fr}
设 $X$ 是域 $K$ 上的有限概形；置
$n=\sum_{x\in X}[k(x):K]_s$
\emph{（回顾：若 $K'$ 是 $K$ 的一个扩域，则 $[K':K]_s$ 是 $K'$ 在
$K$ 上的\emph{可分秩}，即 $K'$ 中所含 $K$ 的最大可分代数扩张的秩）}；
\oldpage[I]{149}
则对于 $K$ 的任意代数闭扩域 $\Omega$，$X\otimes_K\Omega$ 的底空间
恰有 $n$ 个点，并且这些点与 $X$ 的 $\Omega$-值点等同。
\end{corollary}

显然只需考虑环 $A=A(X)$ 为局部环的情形
（\hyperref[I.6.2.2-fr]{6.2.2}）；设 $\mathfrak m$ 是它的极大理想，
$L=A/\mathfrak m$ 是它的剩余域，这是 $K$ 的一个代数扩张。于是 $X$
的 $\Omega$-值点一一对应于 $X\otimes_K\Omega$ 的 $\Omega$-截面
（\hyperref[I.3.4.1-fr]{3.4.1} 和 \hyperref[I.3.3.14-fr]{3.3.14}），
也一一对应于从 $L$ 到 $\Omega$ 的 $K$-同态
（\hyperref[I.1.7.3-fr]{1.7.3}）；结合
（\hyperref[I.6.4.3-fr]{6.4.3}），结论由此得出（Bourbaki，
\emph{Alg.}，第~V 章，\textsection7，n\textsuperscript{o}~5，命题~8）。

\begin{env}[6.4.8]
\label{I.6.4.8-fr}
（\hyperref[I.6.4.7-fr]{6.4.7}）中定义的数 $n$ 称为 $A$（或 $X$）
在 $K$ 上的\emph{可分秩}，也称为 $X$ 的\emph{几何点数}；因此它等于
$X(\Omega)_K$ 的元素个数。由该定义立即可知，对于 $K$ 的任意扩域
$K'$，$X\otimes_K K'$ 与 $X$ 有相同的几何点数。将这个数记为
$n(X)$，则显然对于 $K$ 上的任意两个有限概形 $X$、$Y$，有
\[
  n(X\amalg Y)=n(X)+n(Y).
  \tag{6.4.8.1}\label{I.6.4.8.1-fr}
\]

在同样的假设下，还有
\[
  n(X\times_K Y)=n(X)n(Y)
  \tag{6.4.8.2}\label{I.6.4.8.2-fr}
\]
这是把 $n(X)$ 解释为 $X(\Omega)_K$ 的元素个数，并应用公式
（\hyperref[I.3.4.3.1-fr]{3.4.3.1}）后立即得到的。
\end{env}

\begin{proposition}[6.4.9]
\label{I.6.4.9-fr}
设 $K$ 是一个域，$X$、$Y$ 是两个 $K$-代数预概形，$f:X\to Y$ 是一个
$K$-态射，$\Omega$ 是 $K$ 的一个代数闭扩域，并且在 $K$ 上的超越次数
无限。则 $f$ 为满射的充要条件是：由 $f$ 对应的映射
$X(\Omega)_K\to Y(\Omega)_K$（\hyperref[I.3.4.1-fr]{3.4.1}）为满射。
\end{proposition}

必要性由（\hyperref[I.6.3.10-fr]{6.3.10}）得出，因为 $f$ 必然是有限型的
（\hyperref[I.6.3.9-fr]{6.3.9}）。为证明该条件充分，按
（\hyperref[I.6.3.10-fr]{6.3.10}）中的方式论证，并注意到对于每个
$y\in Y$，$k(y)$ 都是 $K$ 的一个有限型扩张，因而 $K$-同构于
$\Omega$ 的一个子域。

\begin{remark}[6.4.10]
\label{I.6.4.10-fr}
我们将在第~IV 章看到，即使不对 $\Omega$ 在 $K$ 上的超越次数作任何
假设，（\hyperref[I.6.4.9-fr]{6.4.9}）的结论仍然成立。
\end{remark}

\begin{proposition}[6.4.11]
\label{I.6.4.11-fr}
若 $f:X\to Y$ 是一个有限型态射，则对于每个 $y\in Y$，纤维
$f^{-1}(y)$ 是剩余域 $k(y)$ 上的一个代数预概形，并且对于每个
$x\in f^{-1}(y)$，$k(x)$ 是 $k(y)$ 的一个有限型扩张。
\end{proposition}

由于 $f^{-1}(y)=X\otimes_Y k(y)$
（\hyperref[I.3.6.3-fr]{3.6.3}），该命题由
（\hyperref[I.6.3.4-fr]{6.3.4, (iv)}）和
（\hyperref[I.6.3.3-fr]{6.3.3}）得出。

\begin{proposition}[6.4.12]
\label{I.6.4.12-fr}
设 $f:X\to Y$、$g:Y'\to Y$ 是两个态射；置 $X'=X\times_Y Y'$，并设
$f'=f_{(Y')}:X'\to Y'$。设 $y'\in Y'$、$y=g(y')$；若纤维
$f^{-1}(y)$ 是 $k(y)$ 上的一个有限代数概形，则纤维 $f'^{-1}(y')$
是 $k(y')$ 上的一个有限代数概形，并且它与 $f^{-1}(y)$ 有相同的秩和
几何点数。
\end{proposition}

结合纤维的传递性（\hyperref[I.3.6.5-fr]{3.6.5}），这立即由
（\hyperref[I.6.4.6-fr]{6.4.6}）和
（\hyperref[I.6.4.8-fr]{6.4.8}）得出。

\begin{env}[6.4.13]
\label{I.6.4.13-fr}
命题（\hyperref[I.6.4.11-fr]{6.4.11}）表明，有限型态射直观上对应于
“代数簇的代数族”，$Y$ 的点起
\oldpage[I]{150}
“参数”的作用，这赋予这些态射一种“几何”意义。非有限型态射在后文
主要出现于“改变基预概形”的问题中，例如局部化或完备化。
\end{env}

\subsection{态射的局部确定。}
\label{subsection:I.6.5-fr}

\begin{proposition}[6.5.1]
\label{I.6.5.1-fr}
设 $X$、$Y$ 是两个 $S$-预概形，其中 $Y$ 在 $S$ 上为有限型；设
$x\in X$、$y\in Y$ 位于同一点 $s\in S$ 上方。
\begin{enumerate}
  \item[(i)] 若从 $X$ 到 $Y$ 的两个 $S$-态射 $f=(\psi,\theta)$、
    $f'=(\psi',\theta')$ 满足 $\psi(x)=\psi'(x)=y$，并且从
    $\mathscr{O}_y$ 到 $\mathscr{O}_x$ 的（局部）$\mathscr{O}_s$-同态
    $\theta_x^\sharp$ 与 ${\theta'}_x^\sharp$ 相同，则 $f$ 与 $f'$ 在
    $x$ 的某个开邻域上一致。
  \item[(ii)] 再设 $S$ 是局部 Noether 的。对于每个局部
    $\mathscr{O}_s$-同态
    $\varphi:\mathscr{O}_y\to\mathscr{O}_x$，存在 $X$ 中 $x$ 的一个
    开邻域 $U$ 和一个从 $U$ 到 $Y$ 的 $S$-态射 $f=(\psi,\theta)$，
    使得 $\psi(x)=y$ 且 $\theta_x^\sharp=\varphi$。
\end{enumerate}
\end{proposition}

\begin{enumerate}
  \item[(i)] 由于问题对 $S$、$X$ 和 $Y$ 都是局部的，可以设 $S$、
    $X$、$Y$ 均为仿射的，其环依次为 $A$、$B$、$C$，而 $f$ 与 $f'$
    分别具有形式 $({}^a\varphi,\widetilde{\varphi})$ 和
    $({}^a\varphi',\widetilde{\varphi}')$；这里 $\varphi$、$\varphi'$ 是
    从 $C$ 到 $B$ 的两个 $A$-同态，满足
    $\varphi^{-1}(\mathfrak{j}_x)={\varphi'}^{-1}(\mathfrak{j}_x)
    =\mathfrak{j}_y$，并且由 $\varphi$、$\varphi'$ 导出的从 $C_y$ 到
    $B_x$ 的同态 $\varphi_x$、$\varphi'_x$ 相同；还可以设 $C$ 是有限型
    $A$-代数。设 $c_i$（$1\leq i\leq n$）是 $A$-代数 $C$ 的一组
    生成元，并置 $b_i=\varphi(c_i)$、$b'_i=\varphi'(c_i)$；依假设，在
    分式环 $B_x$ 中有 $b_i/1=b'_i/1$（$1\leq i\leq n$）。这意味着
    存在元素 $s_i\in B-\mathfrak{j}_x$，使得对于
    $1\leq i\leq n$ 有 $s_i(b_i-b'_i)=0$；显然可以设所有 $s_i$ 都等于
    同一个元素 $g\in B-\mathfrak{j}_x$。由此可知，在分式环 $B_g$ 中，
    对于 $1\leq i\leq n$ 有 $b_i/1=b'_i/1$；若 $i_g$ 是典范同态
    $B\to B_g$，则有 $i_g\circ\varphi=i_g\circ\varphi'$；所以 $f$ 与
    $f'$ 在 $D(g)$ 上的限制相同。
  \item[(ii)] 可以化归到 (i) 中的同一情形，并且还可设环 $A$ 是
    Noether 的。设 $c_i$（$1\leq i\leq n$）是 $A$-代数 $C$ 的一组
    生成元，并设
    $\alpha:A[X_1,\ldots,X_n]\to C$ 是从多项式代数
    $A[X_1,\ldots,X_n]$ 到 $C$ 的同态，它对于
    $1\leq i\leq n$ 将 $X_i$ 送到 $c_i$。另一方面，设 $i_y$ 是典范
    同态 $C\to C_y$，并考虑复合同态
    \[
      \beta:A[X_1,\ldots,X_n]
      \xrightarrow{\alpha} C
      \xrightarrow{i_y} C_y
      \xrightarrow{\varphi} B_x.
    \]
    以 $\mathfrak{a}$ 表示 $\beta$ 的核；由于 $A$ 是 Noether 的，
    $A[X_1,\ldots,X_n]$ 也是 Noether 的，因而 $\mathfrak{a}$ 有一组
    有限生成元 $Q_j(X_1,\ldots,X_n)$（$1\leq j\leq m$）。另一方面，
    每个元素 $\varphi(i_y(c_i))$ 都可写为 $b_i/s_i$，其中 $b_i\in B$
    且 $s_i\notin\mathfrak{j}_x$；还可以设所有 $s_i$ 都等于同一个元素
    $g\in B-\mathfrak{j}_x$。于是依假设，在 $B_x$ 中有
    $Q_j(b_1/g,\ldots,b_n/g)=0$；置
    \[
      Q_j(X_1/T,\ldots,X_n/T)
      =P_j(X_1,\ldots,X_n,T)/T^{k_j}
    \]
    其中 $P_j$ 是 $k_j$ 次齐次多项式。再置
    $d_j=P_j(b_1,\ldots,b_n,g)\in B$。依假设，对于某个
    $t_j\in B-\mathfrak{j}_x$，有 $t_jd_j=0$（$1\leq j\leq m$）；
    显然可以设所有 $t_j$
\oldpage[I]{151}
    都等于同一个元素 $h\in B-\mathfrak{j}_x$。由此可知，对于
    $1\leq j\leq m$，有 $P_j(hb_1,\ldots,hb_n,hg)=0$。现在考虑从
    $A[X_1,\ldots,X_n]$ 到分式环 $B_{hg}$ 的同态 $\rho$，它对于
    $1\leq i\leq n$ 将 $X_i$ 送到 $hb_i/hg$；$\mathfrak{a}$ 在这个
    同态下的像为 0，因而 $\alpha^{-1}(0)$ 在 $\rho$ 下的像\emph{更是如此}。
    所以 $\rho$ 分解为
    $A[X_1,\ldots,X_n]\xrightarrow{\alpha}C\xrightarrow{\gamma}B_{hg}$，
    其中 $\gamma(c_i)=hb_i/hg$；而且显然，若 $i_x$ 是典范同态
    $B_{hg}\to B_x$，则图
    \[
      \xymatrix{
        C\ar[r]^\gamma\ar[d]_{i_y} & B_{hg}\ar[d]^{i_x}\\
        C_y\ar[r]_\varphi & B_x
      }
      \tag{6.5.1.1}\label{I.6.5.1.1-fr}
    \]
    交换；所以 $\varphi=\gamma_x$，并且由于 $\varphi$ 是局部同态，
    ${}^a\gamma(x)=y$；于是 $f=({}^a\gamma,\widetilde{\gamma})$ 是从
    $x$ 的邻域 $D(hg)$ 到 $Y$ 的一个满足要求的 $S$-态射。
\end{enumerate}

\begin{corollary}[6.5.2]
\label{I.6.5.2-fr}
在（\hyperref[I.6.5.1-fr]{6.5.1, (ii)}）的假设下，若 $X$ 还在 $S$
上为有限型，则可以设态射 $f$ 为有限型。
\end{corollary}

这由（\hyperref[I.6.3.6-fr]{6.3.6}）得出。

\begin{corollary}[6.5.3]
\label{I.6.5.3-fr}
设（\hyperref[I.6.5.1-fr]{6.5.1, (ii)}）的假设成立，并且还设 $Y$
是整的，$\varphi$ 是单射同态。则可以设
$f=({}^a\gamma,\widetilde{\gamma})$，其中 $\gamma$ 是单射的。
\end{corollary}

事实上，可以设 $C$ 是整的（\hyperref[I.5.1.4-fr]{5.1.4}），因而
$i_y$ 是单射的；这时从图（\hyperref[I.6.5.1.1-fr]{6.5.1.1}）可知
$\gamma$ 是单射的。

\begin{proposition}[6.5.4]
\label{I.6.5.4-fr}
设 $f=(\psi,\theta):X\to Y$ 是一个有限型态射，$x$ 是 $X$ 的一个点，
$y=\psi(x)$。
\begin{enumerate}
  \item[(i)] $f$ 在点 $x$ 为局部浸入
    （\hyperref[I.4.5.1-fr]{4.5.1}）的充要条件是
    $\theta_x^\sharp:\mathscr{O}_y\to\mathscr{O}_x$ 为满射。
  \item[(ii)] 再设 $Y$ 是局部 Noether 的。$f$ 在点 $x$ 为局部同构
    （\hyperref[I.4.5.2-fr]{4.5.2}）的充要条件是
    $\theta_x^\sharp$ 为同构。
\end{enumerate}
\end{proposition}

\begin{enumerate}
  \item[(ii)] 由（\hyperref[I.6.5.1-fr]{6.5.1}），存在 $y$ 的一个开邻域
    $V$ 和一个态射 $g:V\to X$，使得 $g\circ f$（相应地，$f\circ g$）
    有定义，并在 $x$（相应地，$y$）的一个邻域上与恒等态射一致；由此
    不难推出 $f$ 是局部同构。
  \item[(i)] 由于问题对 $X$ 和 $Y$ 都是局部的，可以设 $X$ 和 $Y$
    均为仿射的，其环依次为 $A$、$B$；有
    $f=({}^a\varphi,\widetilde{\varphi})$，其中 $\varphi$ 是一个环同态
    $B\to A$，使 $A$ 成为有限型 $B$-代数；有
    $\varphi^{-1}(\mathfrak{j}_x)=\mathfrak{j}_y$，并且由 $\varphi$
    导出的同态 $\varphi_x:B_y\to A_x$ 是满射。设 $(t_i)$
    （$1\leq i\leq n$）是 $B$-代数 $A$ 的一组生成元；关于
    $\varphi_x$ 的假设蕴含存在 $b_i\in B$ 和
    $c\in B-\mathfrak{j}_y$，使得在分式环 $A_x$ 中，对于
    $1\leq i\leq n$ 有 $t_i/1=\varphi(b_i)/\varphi(c)$。于是由
    （\hyperref[I.1.3.3-fr]{1.3.3}），存在
    $a\in A-\mathfrak{j}_x$，使得若置 $g=a\varphi(c)$，则还有
    $t_i/1=a\varphi(b_i)/g$，\emph{这是在分式环 $A_g$ 中的等式}。
    另一方面，依假设存在一个系数属于环 $\varphi(B)$ 的多项式
    $Q(X_1,\ldots,X_n)$，
\oldpage[I]{152}
    使得 $a=Q(t_1,\ldots,t_n)$；置
    \[
      Q(X_1/T,\ldots,X_n/T)
      =P(X_1,\ldots,X_n,T)/T^m,
    \]
    其中 $P$ 是 $m$ 次齐次多项式。在环 $A_g$ 中有
    \[
      a/1
      =a^mP(\varphi(b_1),\ldots,\varphi(b_n),\varphi(c))/g^m
      =a^m\varphi(d)/g^m
    \]
    其中 $d\in B$。由于在 $A_g$ 中，
    $g/1=(a/1)(\varphi(c)/1)$ 按定义是可逆的，所以 $a/1$ 和
    $\varphi(c)/1$ 也都是可逆的，从而可以写成
    \[
      a/1=(\varphi(d)/1)(\varphi(c)/1)^{-m}.
    \]
    由此可知 $\varphi(d)/1$ 在 $A_g$ 中也可逆。现在置 $h=cd$；由于
    $\varphi(h)/1$ 在 $A_g$ 中可逆，复合同态
    \[
      B\xrightarrow{\varphi}A\longrightarrow A_g
    \]
    分解为
    \[
      B\longrightarrow B_h\xrightarrow{\gamma}A_g
    \]
    （\hyperref[0.1.2.4-fr]{0, 1.2.4}）。我们证明 $\gamma$ 是
    \emph{满射}；只需验证 $B_h$ 在 $A_g$ 中的像含有 $t_i/1$ 和
    $(g/1)^{-1}$。但有
    \[
      (g/1)^{-1}
      =(\varphi(c)/1)^{m-1}(\varphi(d)/1)^{-1}
      =\gamma(c^m/h),
    \]
    以及
    \[
      a/1=\gamma(d^{m+1}/h^m),
    \]
    因而
    \[
      (a\varphi(b_i))/1=\gamma(b_i d^{m+1}/h^m),
    \]
    又因为 $t_i/1=(a\varphi(b_i)/1)(g/1)^{-1}$，断言得证。$h$ 的选择
    蕴含 $\psi(D(g))\subset D(h)$，而 $f$ 在 $D(g)$ 上的限制等于
    $({}^a\gamma,\widetilde{\gamma})$；由于 $\gamma$ 是满射，这一限制
    是从 $D(g)$ 到 $D(h)$ 的闭浸入
    （\hyperref[I.4.2.3-fr]{4.2.3}）。
\end{enumerate}

\begin{corollary}[6.5.5]
\label{I.6.5.5-fr}
设 $f=(\psi,\theta):X\to Y$ 是一个有限型态射。设 $X$ 不可约，以
$x$ 表示它的泛点，并置 $y=\psi(x)$。
\begin{enumerate}
  \item[(i)] $f$ 在 $X$ 的某一点为局部浸入的充要条件是
    $\theta_x^\sharp:\mathscr{O}_y\to\mathscr{O}_x$ 为满射。
  \item[(ii)] 再设 $Y$ 不可约且局部 Noether。$f$ 在 $X$ 的某一点为
    局部同构的充要条件是 $y$ 为 $Y$ 的泛点（或者等价地，由
    （\hyperref[0.2.1.4-fr]{0, 2.1.4}），$f$ 是\emph{支配态射}），并且
    $\theta_x^\sharp$ 是同构（换言之，$f$ 是\emph{双有理的}
    （\hyperref[I.2.2.9-fr]{2.2.9}））。
\end{enumerate}
\end{corollary}

显然，考虑到 $X$ 的每个非空开集都含有 $x$，(i) 由
（\hyperref[I.6.5.4-fr]{6.5.4, (i)}）得出；同理，(ii) 由
（\hyperref[I.6.5.4-fr]{6.5.4, (ii)}）得出。

\subsection{拟紧态射和局部有限型态射。}
\label{subsection:I.6.6-fr}

\begin{definition}[6.6.1]
\label{I.6.6.1-fr}
若 $Y$ 的每个拟紧开集在态射 $f:X\to Y$ 下的逆像都是拟紧的，则称
$f$ 是\emph{拟紧的}。
\end{definition}

设 $\mathfrak{B}$ 是由拟紧开集（例如仿射开集）构成的 $Y$ 的一个拓扑基；
$f$ 为拟紧态射的充要条件是：$\mathfrak{B}$ 中每个集合在 $f$ 下的逆像
都是拟紧的（或者等价地，是仿射开集的\emph{有限}并），因为 $Y$ 的每个
拟紧开集都是 $\mathfrak{B}$ 中有限多个集合的并。例如，若 $X$ 是
\emph{拟紧的}而 $Y$ 是\emph{仿射的}，则\emph{每个}态射 $f:X\to Y$
都是拟紧的：事实上，$X$ 是有限多个仿射开集 $U_i$ 的并，并且对于 $Y$
的每个仿射开集 $V$，$U_i\cap f^{-1}(V)$ 都是仿射的
（\hyperref[I.5.5.10-fr]{5.5.10}），因而是拟紧的。

若 $f:X\to Y$ 是拟紧态射，则显然对于 $Y$ 的每个开集 $V$，$f$ 在
$f^{-1}(V)$ 上的限制都是一个拟紧态射 $f^{-1}(V)\to V$。反过来，若
$(U_\alpha)$ 是 $Y$ 的一个开覆盖，而态射 $f:X\to Y$ 满足各限制
$f^{-1}(U_\alpha)\to U_\alpha$ 都是拟紧的，则 $f$ 是拟紧的。

\begin{definition}[6.6.2]
\label{I.6.6.2-fr}
% EGA1-ERR-0051: source-literal undefined y retained; see 06_errata ledger.
若对于每个 $x\in X$，都存在 $x$ 的一个开邻域 $U$ 和 $y$ 的一个
开邻域 $V\supset f(U)$，使得 $f$ 在 $U$ 上的
\oldpage[I]{153}
限制是一个从 $U$ 到 $V$ 的有限型态射，则称态射 $f:X\to Y$ 是
\emph{局部有限型的}。这时也称 $X$ 是 $Y$ 上的局部有限型预概形，或
局部有限型 $Y$-预概形。
\end{definition}

由（\hyperref[I.6.3.2-fr]{6.3.2}）立即可知，若 $f$ 是局部有限型的，
则对于 $Y$ 的每个开集 $W$，$f$ 在 $f^{-1}(W)$ 上的限制都是一个
局部有限型态射 $f^{-1}(W)\to W$。

若 $Y$ 是局部 Noether 的，而 $X$ 在 $Y$ 上为局部有限型，则由
（\hyperref[I.6.3.7-fr]{6.3.7}），$X$ 是局部 Noether 的。

\begin{proposition}[6.6.3]
\label{I.6.6.3-fr}
态射 $f:X\to Y$ 为有限型的充要条件是：它是拟紧且局部有限型的。
\end{proposition}

考虑到（\hyperref[I.6.3.1-fr]{6.3.1}）以及
（\hyperref[I.6.6.1-fr]{6.6.1}）之后的说明，条件的必要性是立即的。
反过来，设这些条件成立，并设 $U$ 是 $Y$ 的一个仿射开集，其环为 $A$；
对于每个 $x\in f^{-1}(U)$，依假设存在 $x$ 的一个邻域
$V(x)\subset f^{-1}(U)$，以及 $y=f(x)$ 的一个邻域 $W(x)\subset U$，
它含有 $f(V(x))$，并且 $f$ 在 $V(x)$ 上的限制是一个有限型态射
$V(x)\to W(x)$。
% EGA1-ERR-0052: source-literal x retained although W_1(x) is a subset of Y.
将 $W(x)$ 换成形如 $D(g)$（其中 $g\in A$）的 $x$ 的一个邻域
$W_1(x)\subset W(x)$，并将 $V(x)$ 换成
$V(x)\cap f^{-1}(W_1(x))$，就可以设 $W(x)$ 形如 $D(g)$，因而在
$U$ 上为有限型（因为它的环可写为 $A[1/g]$）；所以 $V(x)$ 在 $U$
上为有限型。此外，依假设 $f^{-1}(U)$ 是拟紧的，因而是有限多个开集
$V(x_i)$ 的并；证明完毕。

\begin{proposition}[6.6.4]
\label{I.6.6.4-fr}
\begin{enumerate}
  \item[(i)] 浸入 $X\to Y$ 在下列任一情形都是拟紧的：它是闭浸入；
    $Y$ 的底空间是局部 Noether 的；或 $X$ 的底空间是 Noether 的。
  \item[(ii)] 两个拟紧态射的复合是拟紧态射。
  \item[(iii)] 若 $f:X\to Y$ 是拟紧 $S$-态射，则对于基预概形的任意
    扩张 $g:S'\to S$，$f_{(S')}:X_{(S')}\to Y_{(S')}$ 也是拟紧的。
  \item[(iv)] 若 $f:X\to X'$ 和 $g:Y\to Y'$ 是两个拟紧 $S$-态射，
    则 $f\times_S g$ 是拟紧的。
  \item[(v)] 若两个态射 $f:X\to Y$、$g:Y\to Z$ 的复合是拟紧的，
    并且 $g$ 是分离的，或者 $X$ 的底空间是局部 Noether 的，则 $f$
    是拟紧的。
  \item[(vi)] 态射 $f$ 为拟紧的充要条件是 $f_{\mathrm{red}}$ 为拟紧的。
\end{enumerate}
\end{proposition}

应注意，(vi) 是显然的，因为态射是否拟紧只依赖于相应的底空间连续映射。
我们同样证明 (v) 中 $X$ 的底空间被假设为局部 Noether 的情形。置
$h=g\circ f$，并设 $U$ 是 $Y$ 中的一个拟紧开集；$g(U)$ 在 $Z$ 中
是拟紧的（但未必是开的），因而包含在有限多个拟紧开集 $V_j$ 的并中
（\hyperref[I.2.1.3-fr]{2.1.3}），于是 $f^{-1}(U)$ 包含在各
$h^{-1}(V_j)$ 的并中；这些集合是 $X$ 的拟紧子空间，因而是 Noether
子空间。由（\hyperref[0.2.2.3-fr]{0, 2.2.3}）可知 $f^{-1}(U)$ 是
Noether 空间，\emph{从而更是}拟紧的。

为证明其他断言，只需证明 (i)、(ii) 和 (iii)
（\hyperref[I.5.5.12-fr]{5.5.12}）。但 (ii) 是显然的；当 $Y$ 的空间
是局部 Noether 的或 $X$ 的空间是 Noether 的时，(i) 由
（\hyperref[I.6.3.5-fr]{6.3.5}）得出，而对于闭浸入，(i) 是显然的。
为证明 (iii)，
\oldpage[I]{154}
可以只考虑 $S=Y$ 的情形（\hyperref[I.3.3.11-fr]{3.3.11}）；置
$f'=f_{(S')}$，并设 $U'$ 是 $S'$ 中的一个拟紧开集。对于每个
$s'\in U'$，设 $T$ 是 $S$ 中 $g(s')$ 的一个仿射开邻域，并设 $W$
是 $s'$ 的一个仿射开邻域，且包含在 $U'\cap g^{-1}(T)$ 中；只需证明
$f'^{-1}(W)$ 是拟紧的；换言之，可以化归为证明：当 $S$ 和 $S'$ 均
为仿射的时，$X\times_S S'$ 的底空间是拟紧的。但这时依假设，$X$
是有限多个仿射开集 $V_j$ 的并，所以 $X\times_S S'$ 是各仿射概形
$V_j\times_S S'$ 的底空间之并
（\hyperref[I.3.2.2-fr]{3.2.2} 和 \hyperref[I.3.2.7-fr]{3.2.7}）；
命题得证。

还应注意，若 $X=X'\amalg X''$ 是两个预概形之和，则态射 $f:X\to Y$
为拟紧的充要条件是它在 $X'$ 和 $X''$ 上的限制都为拟紧的。

\begin{proposition}[6.6.5]
\label{I.6.6.5-fr}
设 $f:X\to Y$ 是拟紧态射。则 $f$ 为支配态射的充要条件是：对于 $Y$
的每个不可约分支的泛点 $y$，$f^{-1}(y)$ 都含有 $X$ 的某个不可约分支
的泛点。
\end{proposition}

显然，该条件是充分的（无须假设 $f$ 拟紧）。为证明它是必要的，考虑
$y$ 的一个仿射开邻域 $U$；$f^{-1}(U)$ 是拟紧的，因而是\emph{有限}
多个仿射开集 $V_i$ 的并；$f$ 为支配态射这一假设蕴含 $y$ 属于某个
$f(V_i)$ 在 $U$ 中的闭包。显然可以设 $X$ 和 $Y$ 是约化的；由于
$V_i$ 的一个不可约分支在 $X$ 中的闭包是 $X$ 的一个不可约分支
（\hyperref[0.2.1.6-fr]{0, 2.1.6}），可以将 $X$ 换成 $V_i$，并将 $Y$
换成 $U$ 的约化闭子预概形，其底空间为
$\overline{f(V_i)}\cap U$（\hyperref[I.5.2.1-fr]{5.2.1}）；这样就化归为
$X=\operatorname{Spec}(A)$ 和 $Y=\operatorname{Spec}(B)$ 均为仿射且
约化的情形。由于 $f$ 是支配态射，此时 $B$ 是 $A$ 的一个子环
（\hyperref[I.1.2.7-fr]{1.2.7}），而命题由下述事实得出：$B$ 的每个
极小素理想都是 $A$ 的某个极小素理想与 $B$ 的交
（\hyperref[0.1.5.8-fr]{0, 1.5.8}）。

\begin{proposition}[6.6.6]
\label{I.6.6.6-fr}
\begin{enumerate}
  \item[(i)] 每个局部浸入都是局部有限型的。
  \item[(ii)] 若两个态射 $f:X\to Y$、$g:Y\to Z$ 都是局部有限型的，
    则 $g\circ f$ 也是局部有限型的。
  \item[(iii)] 若 $f:X\to Y$ 是局部有限型 $S$-态射，则对于基预概形的
    任意扩张 $S'\to S$，$f_{(S')}:X_{(S')}\to Y_{(S')}$ 都是局部
    有限型的。
  \item[(iv)] 若 $f:X\to X'$ 和 $g:Y\to Y'$ 是两个局部有限型
    $S$-态射，则 $f\times_S g$ 是局部有限型的。
  \item[(v)] 若两个态射的复合 $g\circ f$ 是局部有限型的，则 $f$ 是
    局部有限型的。
  \item[(vi)] 若态射 $f$ 是局部有限型的，则 $f_{\mathrm{red}}$ 也是。
\end{enumerate}
\end{proposition}

由（\hyperref[I.5.5.12-fr]{5.5.12}），只需证明 (i)、(ii) 和 (iii)。
若 $j:X\to Y$ 是一个局部浸入，则对于每个 $x\in X$，存在 $Y$ 中
$j(x)$ 的一个开邻域 $V$ 和 $X$ 中 $x$ 的一个开邻域 $U$，使得 $j$
在 $U$ 上的限制是闭浸入 $U\to V$
（\hyperref[I.4.5.1-fr]{4.5.1}），所以这一限制为有限型。为证明 (ii)，
考虑一点 $x\in X$；依假设，存在 $g(f(x))$ 的一个开邻域 $W$ 和
$f(x)$ 的一个开邻域 $V$，使得 $g(V)\subset W$ 且 $V$ 在 $W$ 上为
有限型；此外，$f^{-1}(V)$ 在 $V$ 上为局部有限型
（\hyperref[I.6.6.2-fr]{6.6.2}），所以存在 $x$ 的一个开邻域 $U$，它
\oldpage[I]{155}
包含在 $f^{-1}(V)$ 中且在 $V$ 上为有限型；于是
$g(f(U))\subset W$，并且 $U$ 在 $W$ 上为有限型
（\hyperref[I.6.3.4-fr]{6.3.4, (ii)}）。最后，为证明 (iii)，可以只
考虑 $Y=S$ 的情形（\hyperref[I.3.3.11-fr]{3.3.11}）；对于每个
$x'\in X'=X_{(S')}$，设 $x$ 为 $x'$ 在 $X$ 中的像，$s$ 为 $x$ 在
$S$ 中的像，$T$ 为 $s$ 的一个开邻域，$T'$ 为它在 $S'$ 中的逆像，
$U$ 为 $x$ 的一个开邻域，其像包含在 $T$ 中且在 $T$ 上为有限型；则
$U\times_S T'=U\times_T T'$ 是 $x'$ 的一个开邻域
（\hyperref[I.3.2.7-fr]{3.2.7}），并且在 $T'$ 上为有限型
（\hyperref[I.6.3.4-fr]{6.3.4, (iv)}）。

\begin{corollary}[6.6.7]
\label{I.6.6.7-fr}
设 $X$、$Y$ 是两个在 $S$ 上局部有限型的 $S$-预概形。若 $S$ 是局部
Noether 的，则 $X\times_S Y$ 是局部 Noether 的。
\end{corollary}

事实上，由于 $X$ 在 $S$ 上为局部有限型，所以它是局部 Noether
的；而 $X\times_S Y$ 在 $X$ 上为局部有限型，因而也是局部 Noether
的。

\begin{remark}[6.6.8]
\label{I.6.6.8-fr}
命题（\hyperref[I.6.3.10-fr]{6.3.10}）及其证明立即推广到仅假设态射
$f$ 为局部有限型的情形。同样，当只假设命题
（\hyperref[I.6.4.2-fr]{6.4.2}）和（\hyperref[I.6.4.9-fr]{6.4.9}）
的陈述中出现的预概形 $X$、$Y$ 在域 $K$ 上为局部有限型时，这两个
命题仍然成立。
\end{remark}
\section{有理映射}
\label{section:I.7-fr}

\subsection{有理映射和有理函数。}
\label{subsection:I.7.1-fr}

\begin{env}[7.1.1]
\label{I.7.1.1-fr}
设 $X$、$Y$ 是两个预概形，$U$、$V$ 是 $X$ 中的两个稠密开集，$f$
（相应地 $g$）是从 $U$（相应地 $V$）到 $Y$ 的态射；若 $f$ 和 $g$
在 $U\cap V$ 的某个稠密开集中重合，则称它们\emph{等价}。由于 $X$ 中
有限多个稠密开集的交仍是 $X$ 中的稠密开集，显然这是一种
\emph{等价关系}。
\end{env}

\begin{definition}[7.1.2]
\label{I.7.1.2-fr}
给定两个预概形 $X$、$Y$，把从 $X$ 的稠密开集到 $Y$ 的态射关于
（\hyperref[I.7.1.1-fr]{7.1.1}）中所定义关系的一个等价类称为
\emph{从 $X$ 到 $Y$ 的有理映射}。若 $X$ 和 $Y$ 是 $S$-预概形，且
这样的等价类有一个代表元是 $S$-态射，则称它为
\emph{$S$-有理映射}。把从 $S$ 到 $X$ 的任意 $S$-有理映射称为
\emph{$X$ 的有理 $S$-截面}。把 $X$-预概形
$X\otimes_{\mathbf{Z}}\mathbf{Z}[T]$（其中 $T$ 是一个不定元）的任意
有理 $X$-截面称为\emph{预概形 $X$ 上的有理函数}。

作为语言上的简便，当只讨论 $S$-预概形时，只要不致混淆，就用
“有理映射”代替“$S$-有理映射”。

设 $f$ 是从 $X$ 到 $Y$ 的有理映射，$U$ 是 $X$ 的开集；若 $f_1$、
$f_2$ 是属于 $f$ 这一等价类的态射，分别定义在 $X$ 的稠密开集 $V$、
$W$ 上，则限制 $f_1|(U\cap V)$、$f_2|(U\cap W)$ 在 $U\cap V\cap W$
上重合，而后者在 $U$ 中稠密；因此态射类 $f$ 定义了从 $U$ 到 $Y$
的有理映射，称为 $f$ \emph{在 $U$ 上的限制}，记作 $f|U$。

把每个 $S$-态射 $f:X\to Y$ 对应到含有 $f$ 的那个 $S$-有理映射，
就定义了从
\oldpage[I]{156}
$\operatorname{Hom}_S(X,Y)$ 到从 $X$ 到 $Y$ 的 $S$-有理映射集合的一个
典范映射。用 $\Gamma_{\mathrm{rat}}(X/Y)$ 表示 $X$ 的有理 $Y$-截面
集合，于是有典范映射
$\Gamma(X/Y)\to\Gamma_{\mathrm{rat}}(X/Y)$。此外显然，若 $X$、$Y$
是两个 $S$-预概形，则从 $X$ 到 $Y$ 的 $S$-有理映射集合典范地等同于
$\Gamma_{\mathrm{rat}}((X\times_S Y)/X)$
（\hyperref[I.3.3.14-fr]{3.3.14}）。
\end{definition}

\begin{env}[7.1.3]
\label{I.7.1.3-fr}
由（\hyperref[I.7.1.2-fr]{7.1.2}）和
（\hyperref[I.3.3.14-fr]{3.3.14}）立即可知，$X$ 上的有理函数典范地
等同于 $X$ 的处处稠密开集上的\emph{结构层 $\mathscr{O}_X$ 的截面的
等价类}；两个这样的截面等价，是指它们在一个处处稠密开集中重合，
该开集包含在二者定义域的交中。特别地，$X$ 上的有理函数构成一个
\emph{环} $R(X)$。
\end{env}

\begin{env}[7.1.4]
\label{I.7.1.4-fr}
当 $X$ 是\emph{不可约}预概形时，每个非空开集都在 $X$ 中稠密；也可以
说，$X$ 的非空开集就是 $X$ 的\emph{泛点 $x$ 的开邻域}。因此在这种
情形，说从 $X$ 的两个非空开部分到 $Y$ 的态射等价，就是说它们在点
$x$ 处具有\emph{相同的芽}。换言之，有理映射（相应地 $S$-有理映射）
$X\to Y$ 等同于从 $X$ 的非空开部分到 $Y$ 的态射（相应地
$S$-态射）在 $X$ 的泛点 $x$ 处的\emph{芽}。特别地：
\end{env}

\begin{proposition}[7.1.5]
\label{I.7.1.5-fr}
若 $X$ 是不可约预概形，则 $X$ 上的有理函数环 $R(X)$ 典范地等同于
$X$ 的泛点 $x$ 的局部环 $\mathscr{O}_x$。这是一个维数为 0 的局部
环，因而当 $X$ 为 Noether 的时是 Artin 局部环；当 $X$ 为整的时它是
一个域；若 $X$ 还是仿射概形，则它等同于 $A(X)$ 的分式域。
\end{proposition}

考虑到以上所述以及有理函数与 $\mathscr{O}_X$ 在处处稠密开集上的截面
的等同，第一项断言无非是层在一点的茎的定义。对于其他断言，可以只
考虑 $X$ 为仿射且其环为 $A$ 的情形；这时 $\mathfrak{j}_x$ 是 $A$ 的
幂零根，所以 $\mathscr{O}_x$ 的维数为 0；若 $A$ 是整环，则
$\mathfrak{j}_x=(0)$，故 $\mathscr{O}_x$ 是 $A$ 的分式域。最后，若
$A$ 是 Noether 的，由（\cite{I-11}, p.~127, cor.~4）可知
$\mathfrak{j}_x$ 是幂零的，并且
$\mathscr{O}_x=A_{\mathfrak{j}_x}$ 是 Artin 的。

若 $X$ 是\emph{整的}，则对于每个 $z\in X$，环 $\mathscr{O}_z$ 都是
整环；每个含有 $z$ 的仿射开集 $U$ 也含有 $x$，而等于 $A(U)$ 分式域
的 $R(U)$ 因而等同于 $R(X)$；由此可知 $R(X)$ 也等同于
\emph{$\mathscr{O}_z$ 的分式域}：$\mathscr{O}_z$ 与 $R(X)$ 的一个
子环之间的典范等同，是把截面芽 $s\in\mathscr{O}_z$ 对应到 $X$ 上
唯一的有理函数；这个有理函数是 $\mathscr{O}_X$ 的一个截面的等价类，
该截面必定义在处处稠密开集上，并在点 $z$ 处具有芽 $s$。

\begin{env}[7.1.6]
\label{I.7.1.6-fr}
现在设 $X$ 有\emph{有限}多个不可约分支 $X_i$（$1\leq i\leq n$）
（当 $X$ 的底空间为 \emph{Noether} 空间时即属此情形）；设 $X'_i$ 是
$X$ 的开集，它在 $X_i$ 中是各 $X_j\cap X_i$（$j\neq i$）之并的
补集；$X'_i$ 是不可约的，它的泛点 $x_i$ 就是 $X_i$ 的泛点，而且各
$X'_i$ 两两不交，其并在 $X$ 中稠密
（\hyperref[0.2.1.6-fr]{0, 2.1.6}）。对于
\oldpage[I]{157}
$X$ 的每个处处稠密开集 $U$，$U_i=U\cap X'_i$ 是 $X'_i$ 中的非空
稠密开集，各 $U_i$ 两两没有公共点，所以 $U'=\bigcup_i U_i$ 在 $X$
中稠密。给定一个从 $U'$ 到 $Y$ 的态射，等价于任意给定从每个 $U_i$
到 $Y$ 的态射。因此：
\end{env}

\begin{proposition}[7.1.7]
\label{I.7.1.7-fr}
设 $X$、$Y$ 是两个预概形（相应地 $S$-预概形），且 $X$ 有有限多个
不可约分支 $X_i$，其泛点为 $x_i$（$1\leq i\leq n$）。若 $R_i$ 是
从 $X$ 的开部分到 $Y$ 的态射（相应地 $S$-态射）在点 $x_i$ 处的芽
的集合，则从 $X$ 到 $Y$ 的有理映射（相应地 $S$-有理映射）集合等同于
各 $R_i$ 的积（$1\leq i\leq n$）。
\end{proposition}

\begin{corollary}[7.1.8]
\label{I.7.1.8-fr}
设 $X$ 是 Noether 预概形。$X$ 上的有理函数环是 Artin 环；它的各局部
分量是 $X$ 的不可约分支的泛点 $x_i$ 的环 $\mathscr{O}_{x_i}$。
\end{corollary}

\begin{corollary}[7.1.9]
\label{I.7.1.9-fr}
设 $A$ 是 Noether 环，$X=\operatorname{Spec}(A)$。若 $Q$ 是 $A$ 的
所有极小素理想之并的补集，则 $X$ 上的有理函数环典范地等同于分式环
$Q^{-1}A$。
\end{corollary}

这由下述引理得出：

\begin{lemma}[7.1.9.1]
\label{I.7.1.9.1-fr}
对于元素 $f\in A$，$D(f)$ 在 $X$ 中处处稠密的充要条件是 $f\in Q$；
$X$ 的每个稠密开集都含有形如 $D(f)$ 的开集，其中 $f\in Q$。
\end{lemma}

事实上，设该引理已经证明。由于截面环
$\Gamma(D(f),\mathscr{O}_X)$ 等同于 $A_f$
（\hyperref[I.1.3.6-fr]{1.3.6} 和
\hyperref[I.1.3.7-fr]{1.3.7}），且 $f\in Q$ 时的 $D(f)$ 在 $X$ 的
稠密开集按 $\supset$ 排成的有序集中构成共尾族，所以由定义
（\hyperref[I.7.1.1-fr]{7.1.1}）可知，$X$ 上的有理函数环等同于
$f\in Q$ 时各 $A_f$ 的归纳极限（预序关系为“$g$ 是 $f$ 的倍数”），
即 $Q^{-1}A$（\hyperref[0.1.4.5-fr]{0, 1.4.5}）。

为证明（\hyperref[I.7.1.9.1-fr]{7.1.9.1}），仍用 $X_i$ 表示 $X$ 的
不可约分支（$1\leq i\leq n$）；$D(f)$ 在 $X$ 中稠密，当且仅当
$D(f)\cap X_i\neq\varnothing$（$1\leq i\leq n$）；置
$\mathfrak{p}_i=\mathfrak{j}(X_i)$，这又等价于
$f\notin\mathfrak{p}_i$（$1\leq i\leq n$）。由于各
$\mathfrak{p}_i$ 是 $A$ 的极小素理想
（\hyperref[I.1.1.14-fr]{1.1.14}），关系
$f\notin\mathfrak{p}_i$（$1\leq i\leq n$）等价于 $f\in Q$，引理的
第一项断言得证。另一方面，若 $U$ 是 $X$ 的稠密开集，则 $U$ 的补集
形如 $V(\mathfrak{a})$，其中理想 $\mathfrak{a}$ 不包含在任何
$\mathfrak{p}_i$ 中；所以它不包含在它们的并中
（\cite{I-10}, p.~13），从而存在 $f\in\mathfrak{a}$ 属于 $Q$；于是
$D(f)\subset U$，证明完毕。

\begin{env}[7.1.10]
\label{I.7.1.10-fr}
再次设 $X$ 不可约，泛点为 $x$。由于 $X$ 的每个非空开集 $U$ 都含有
$x$，因而也含有每个满足 $x\in\overline{\{z\}}$ 的 $z\in X$，所以
每个态射 $U\to Y$ 都可以与典范态射
$\operatorname{Spec}(\mathscr{O}_x)\to X$
（\hyperref[I.2.4.1-fr]{2.4.1}）复合；从 $X$ 的两个非空开部分到 $Y$
的两个态射，如果在 $X$ 的某个非空开部分上重合，则复合后给出同一个
态射 $\operatorname{Spec}(\mathscr{O}_x)\to Y$。换言之，每个从 $X$
到 $Y$ 的有理映射都由此对应于一个确定的态射
$\operatorname{Spec}(\mathscr{O}_x)\to Y$。
\end{env}

\oldpage[I]{158}
\begin{proposition}[7.1.11]
\label{I.7.1.11-fr}
设 $X$、$Y$ 是两个 $S$-预概形；假设 $X$ 不可约，泛点为 $x$，且 $Y$
在 $S$ 上为有限型。从 $X$ 到 $Y$ 的两个 $S$-有理映射，若对应于同一个
$S$-态射 $\operatorname{Spec}(\mathscr{O}_x)\to Y$，则二者相同。若还
假设 $S$ 是局部 Noether 的，则每个从
$\operatorname{Spec}(\mathscr{O}_x)$ 到 $Y$ 的 $S$-态射都对应于从 $X$
到 $Y$ 的唯一一个 $S$-有理映射。
\end{proposition}

考虑到 $X$ 的每个非空开集都处处稠密，这立即由
（\hyperref[I.6.5.1-fr]{6.5.1}）得出。

\begin{corollary}[7.1.12]
\label{I.7.1.12-fr}
假设 $S$ 是局部 Noether 的，并且
（\hyperref[I.7.1.11-fr]{7.1.11}）的其他假设均成立。这时从 $X$ 到
$Y$ 的 $S$-有理映射等同于 $S$-预概形 $Y$ 的取值于 $S$-预概形
$\operatorname{Spec}(\mathscr{O}_x)$ 的点。
\end{corollary}

这正是（\hyperref[I.7.1.11-fr]{7.1.11}），采用
（\hyperref[I.3.4.1-fr]{3.4.1}）中引入的术语即可。

\begin{corollary}[7.1.13]
\label{I.7.1.13-fr}
假设（\hyperref[I.7.1.12-fr]{7.1.12}）的条件成立。设 $s$ 是 $x$ 在
$S$ 中的像。给定从 $X$ 到 $Y$ 的一个 $S$-有理映射，等价于给定 $Y$
中位于 $s$ 上的一点 $y$，以及一个局部 $\mathscr{O}_s$-同态
$\mathscr{O}_y\to\mathscr{O}_x=R(X)$。
\end{corollary}

这由（\hyperref[I.7.1.11-fr]{7.1.11}）和
（\hyperref[I.2.4.4-fr]{2.4.4}）得出。

特别地：

\begin{corollary}[7.1.14]
\label{I.7.1.14-fr}
在（\hyperref[I.7.1.12-fr]{7.1.12}）的条件下，从 $X$ 到 $Y$ 的
$S$-有理映射（$Y$ 给定）只依赖于 $S$-预概形
$\operatorname{Spec}(\mathscr{O}_x)$；特别地，对于每个 $z\in X$，
把 $X$ 换成 $\operatorname{Spec}(\mathscr{O}_z)$ 时，这些有理映射
保持不变。
\end{corollary}

事实上，由于 $z\in\overline{\{x\}}$，$x$ 是
$Z=\operatorname{Spec}(\mathscr{O}_z)$ 的泛点，并且
$\mathscr{O}_{X,x}=\mathscr{O}_{Z,x}$。

当 $X$ 为整的时，$R(X)=\mathscr{O}_x=k(x)$ 是一个域
（\hyperref[I.7.1.5-fr]{7.1.5}）；这时以上推论具体化为：

\begin{corollary}[7.1.15]
\label{I.7.1.15-fr}
假设（\hyperref[I.7.1.12-fr]{7.1.12}）的条件成立，并且 $X$ 是整的。
设 $s$ 是 $x$ 在 $S$ 中的像。这时，从 $X$ 到 $Y$ 的 $S$-有理映射
等同于 $Y\otimes_S k(s)$ 的取值于 $k(s)$ 的扩张 $R(X)$ 的几何点；
换言之，每个这样的映射等价于给定其中位于 $s$ 上的一点
$y\in Y$，以及一个从 $k(y)$ 到 $k(x)=R(X)$ 的 $k(s)$-单同态。
\end{corollary}

事实上，$Y$ 中位于 $s$ 上的点等同于 $Y\otimes_S k(s)$ 的点
（\hyperref[I.3.6.3-fr]{3.6.3}），而局部 $\mathscr{O}_s$-同态
$\mathscr{O}_y\to R(X)$ 等同于 $k(s)$-单同态 $k(y)\to R(X)$。

更特别地：

\begin{corollary}[7.1.16]
\label{I.7.1.16-fr}
设 $k$ 是一个域，$X$、$Y$ 是 $k$ 上的两个代数预概形
（\hyperref[I.6.4.1-fr]{6.4.1}）；还假设 $X$ 是整的。则从 $X$ 到
$Y$ 的 $k$-有理映射等同于 $Y$ 的取值于 $k$ 的扩张 $R(X)$ 的几何点
（\hyperref[I.3.4.4-fr]{3.4.4}）。
\end{corollary}

\subsection{有理映射的定义域。}
\label{subsection:I.7.2-fr}

\begin{env}[7.2.1]
\label{I.7.2.1-fr}
设 $X$、$Y$ 是两个预概形，$f$ 是从 $X$ 到 $Y$ 的有理映射。若存在
一个含有 $x$ 的处处稠密开集 $U$，以及属于等价类 $f$ 的态射
$U\to Y$，则称 $f$ \emph{在点 $x\in X$ 处有定义}。使 $f$
有定义的点 $x\in X$ 的集合称为 \emph{$f$ 的定义域}；显然它是 $X$
中的处处稠密开集。
\end{env}

\oldpage[I]{159}
\begin{proposition}[7.2.2]
\label{I.7.2.2-fr}
设 $X$、$Y$ 是两个 $S$-预概形，$X$ 是约化的，且 $Y$ 在 $S$ 上分离。
设 $f$ 是从 $X$ 到 $Y$ 的 $S$-有理映射，$U_0$ 是它的定义域。则存在
唯一一个属于类 $f$ 的 $S$-态射 $U_0\to Y$。
\end{proposition}

由于对于属于类 $f$ 的每个态射 $U\to Y$，必有 $U\subset U_0$，显然
该命题是下述引理的推论。

\begin{lemma}[7.2.2.1]
\label{I.7.2.2.1-fr}
在（\hyperref[I.7.2.2-fr]{7.2.2}）的假设下，设 $U_1$、$U_2$ 是 $X$
的两个处处稠密开集，$f_i:U_i\to Y$（$i=1,2$）是两个 $S$-态射；
假设存在一个在 $X$ 中稠密的开集 $V\subset U_1\cap U_2$，使 $f_1$
和 $f_2$ 在其中重合。则 $f_1$ 和 $f_2$ 在 $U_1\cap U_2$ 中重合。
\end{lemma}

显然可以只考虑 $X=U_1=U_2$ 的情形。由于 $X$（从而 $V$）是约化的，
$X$ 是 $X$ 的使 $V$ 的典范态射经由其分解的最小闭子预概形
（\hyperref[I.5.2.2-fr]{5.2.2}）。设
$g=(f_1,f_2)_S:X\to Y\times_S Y$；依假设，对角线
$T=\Delta_Y(Y)$ 是 $Y\times_S Y$ 的闭子预概形，所以
$Z=g^{-1}(T)$ 是 $X$ 的闭子预概形
（\hyperref[I.4.4.1-fr]{4.4.1}）。若 $h:V\to Y$ 是 $f_1$ 和 $f_2$
在 $V$ 上的共同限制，则 $g$ 在 $V$ 上的限制是 $g'=(h,h)_S$，它分解为
$g'=\Delta_Y\circ h$；由于 $\Delta_Y^{-1}(T)=Y$，有
$g'^{-1}(T)=V$，所以 $Z$ 是 $X$ 的一个在该开集上诱导 $V$ 的闭子预概形，
因而 $V$ 的典范态射经由 $Z$ 分解；这蕴含 $Z=X$。由关系
$g^{-1}(T)=X$ 及（\hyperref[I.4.4.1-fr]{4.4.1}）可知 $g$ 分解为
$\Delta_Y\circ f$，其中 $f$ 是态射 $X\to Y$；按照对角态射的定义，
这蕴含 $f_1=f_2=f$。

显然，（\hyperref[I.7.2.2-fr]{7.2.2}）中定义的态射 $U_0\to Y$ 是类
$f$ 中唯一一个\emph{不能延拓}成定义在 $X$ 的严格包含 $U_0$ 的开部分
上的态射。\emph{在（\hyperref[I.7.2.2-fr]{7.2.2}）的假设下}，因此
可以把从 $X$ 到 $Y$ 的有理映射与从 $X$ 的处处稠密开集到 $Y$ 的
\emph{不可延拓}态射（即不能延拓到严格更大的开集）\emph{等同}起来。
在这种等同下，命题（\hyperref[I.7.2.2-fr]{7.2.2}）蕴含：

\begin{corollary}[7.2.3]
\label{I.7.2.3-fr}
对 $X$ 和 $Y$ 的假设如（\hyperref[I.7.2.2-fr]{7.2.2}）。设 $U$ 是
$X$ 的处处稠密开集。则从 $U$ 到 $Y$ 的 $S$-态射与从 $X$ 到 $Y$ 且在
$U$ 的每一点都有定义的 $S$-有理映射之间存在典范双射。
\end{corollary}

事实上，由（\hyperref[I.7.2.2-fr]{7.2.2}），对于每个从 $U$ 到 $Y$
的 $S$-态射 $f$，存在唯一一个从 $X$ 到 $Y$ 且延拓 $f$ 的 $S$-有理
映射 $\overline f$。

\begin{corollary}[7.2.4]
\label{I.7.2.4-fr}
设 $S$ 是概形，$X$ 是约化 $S$-预概形，$Y$ 是 $S$-概形，
$f:U\to Y$ 是从 $X$ 的稠密开集 $U$ 到 $Y$ 的 $S$-态射。若
$\overline f$ 是从 $X$ 到 $Y$ 且延拓 $f$ 的 $\mathbf{Z}$-有理映射，
则 $\overline f$ 是 $S$-态射（从而就是从 $X$ 到 $Y$ 且延拓 $f$ 的
$S$-有理映射）。
\end{corollary}

事实上，若 $\varphi:X\to S$、$\psi:Y\to S$ 是结构态射，$U_0$ 是
$\overline f$ 的定义域，$j$ 是嵌入 $U_0\to X$，则只需证明
$\psi\circ\overline f=\varphi\circ j$；由于 $f$ 是 $S$-态射，这立即由
（\hyperref[I.7.2.2.1-fr]{7.2.2.1}）得出。

\begin{corollary}[7.2.5]
\label{I.7.2.5-fr}
设 $X$、$Y$ 是两个 $S$-预概形；假设 $X$ 是约化的，且 $X$、$Y$ 都在
$S$ 上分离。设 $p:Y\to X$ 是 $S$-态射（使 $Y$ 成为 $X$-预概形），
$U$ 是 $X$ 的处处稠密开集，$f$ 是 $Y$ 的一个 $U$-截面；则从 $X$
到 $Y$ 且延拓 $f$ 的有理映射 $\overline f$ 是 $Y$ 的有理 $X$-截面。
\end{corollary}

\oldpage[I]{160}
必须证明 $p\circ\overline f$ 在 $\overline f$ 的定义域中是恒等映射；
由于 $X$ 在 $S$ 上分离，这仍由
（\hyperref[I.7.2.2.1-fr]{7.2.2.1}）得出。

\begin{corollary}[7.2.6]
\label{I.7.2.6-fr}
设 $X$ 是约化预概形，$U$ 是 $X$ 的处处稠密开集。$\mathscr{O}_X$ 在
$U$ 上的截面与在 $U$ 的每一点都有定义的 $X$ 上有理函数 $f$ 之间存在
典范双射。
\end{corollary}

考虑到（\hyperref[I.7.2.3-fr]{7.2.3}）、
（\hyperref[I.7.1.2-fr]{7.1.2}）和（\hyperref[I.7.1.3-fr]{7.1.3}），
只需注意到 $X$-预概形 $X\otimes_{\mathbf Z}\mathbf Z[T]$ 在 $X$ 上分离
（\hyperref[I.5.5.1-fr]{5.5.1, (iv)}）。

\begin{corollary}[7.2.7]
\label{I.7.2.7-fr}
设 $Y$ 是约化预概形，$f:X\to Y$ 是分离态射，$U$ 是 $Y$ 的处处稠密
开集，$g:U\to f^{-1}(U)$ 是 $f^{-1}(U)$ 的一个 $U$-截面，$Z$ 是 $X$
的以 $\overline{g(U)}$ 为底空间的约化子预概形
（\hyperref[I.5.2.1-fr]{5.2.1}）。则 $g$ 是 $X$ 的一个 $Y$-截面的限制
（换言之，由（\hyperref[I.7.2.5-fr]{7.2.5}），延拓 $g$ 的从 $Y$ 到
$X$ 的有理映射处处有定义）的充要条件是：$f$ 在 $Z$ 上的限制是从
$Z$ 到 $Y$ 的同构。
\end{corollary}

$f$ 在 $f^{-1}(U)$ 上的限制是分离态射
（\hyperref[I.5.5.1-fr]{5.5.1, (i)}），所以 $g$ 是闭浸入
（\hyperref[I.5.4.6-fr]{5.4.6}），从而
$g(U)=Z\cap f^{-1}(U)$，并且 $Z$ 在其开集 $g(U)$ 上诱导的子预概形
与 $f^{-1}(U)$ 中由 $g$ 确定的闭子预概形相同
（\hyperref[I.5.2.1-fr]{5.2.1}）。于是陈述中的条件显然充分：若它成立，
$f_Z:Z\to Y$ 是 $f$ 在 $Z$ 上的限制，$\overline g:Y\to Z$ 是其逆同构，
则 $\overline g$ 延拓 $g$。反过来，若 $g$ 是 $X$ 的一个 $Y$-截面 $h$
在 $U$ 上的限制，则 $h$ 是闭浸入
（\hyperref[I.5.4.6-fr]{5.4.6}），所以 $h(Y)$ 是闭的；由于它包含在
$Z$ 中，故等于 $Z$，而由（\hyperref[I.5.2.1-fr]{5.2.1}）可知 $h$
必为从 $Y$ 到 $X$ 的闭子预概形 $Z$ 的同构。

\begin{env}[7.2.8]
\label{I.7.2.8-fr}
设 $X$、$Y$ 是两个 $S$-预概形，假设 $X$ 是约化的且 $Y$ 在 $S$ 上
分离。设 $f$ 是从 $X$ 到 $Y$ 的 $S$-有理映射，$x$ 是 $X$ 的一点；
只要 $f$ 的定义域在 $\operatorname{Spec}(\mathscr{O}_x)$ 上的迹在
$\operatorname{Spec}(\mathscr{O}_x)$ 中稠密，就可以把 $f$ 与典范
$S$-态射 $\operatorname{Spec}(\mathscr{O}_x)\to X$
（\hyperref[I.2.4.1-fr]{2.4.1}）复合。这里后者等同于满足
$x\in\overline{\{z\}}$ 的 $z\in X$ 的集合
（\hyperref[I.2.4.2-fr]{2.4.2}）。在下列情形中该条件成立：
\begin{enumerate}
  \item[1\textsuperscript{o}] $X$ 是不可约的（从而是整的），因为这时
    $X$ 的泛点 $\xi$ 是 $\operatorname{Spec}(\mathscr{O}_x)$ 的泛点；
    由于 $f$ 的定义域 $U$ 含有 $\xi$，
    $U\cap\operatorname{Spec}(\mathscr{O}_x)$ 含有 $\xi$，因而在
    $\operatorname{Spec}(\mathscr{O}_x)$ 中稠密。
  \item[2\textsuperscript{o}] $X$ 是局部 Noether 的；事实上这时我们的
    断言由下述引理得出。
\end{enumerate}
\end{env}

\begin{lemma}[7.2.8.1]
\label{I.7.2.8.1-fr}
设 $X$ 是底空间局部 Noether 的预概形，$x$ 是 $X$ 的一点。
$\operatorname{Spec}(\mathscr{O}_x)$ 的不可约分支，是含有 $x$ 的 $X$
的不可约分支在 $\operatorname{Spec}(\mathscr{O}_x)$ 上的迹。开集
$U\subset X$ 满足 $U\cap\operatorname{Spec}(\mathscr{O}_x)$ 在
$\operatorname{Spec}(\mathscr{O}_x)$ 中稠密的充要条件是：它与含有
$x$ 的 $X$ 的各不可约分支相交（特别地，当 $U$ 在 $X$ 中稠密时，
此条件成立）。
\end{lemma}

第二项断言显然由第一项得出，因此只需证明第一项。由于
$\operatorname{Spec}(\mathscr{O}_x)$ 包含在每个含有 $x$ 的仿射开集
$U$ 中，而且含有 $x$ 的 $U$ 的不可约分支，是含有 $x$ 的 $X$ 的
不可约分支在 $U$ 上的迹
（\hyperref[0.2.1.6-fr]{0, 2.1.6}），可以假设 $X$ 为仿射且其环为
$A$。由于 $A_x$ 的素理想与包含在 $\mathfrak{j}_x$ 中的 $A$ 的素理想
一一对应
\oldpage[I]{161}
（\hyperref[0.1.2.6-fr]{0, 1.2.6}），$A_x$ 的极小素理想与包含在
$\mathfrak{j}_x$ 中的 $A$ 的极小素理想对应；结合
（\hyperref[I.1.1.14-fr]{1.1.14}），引理得证。

现在假设处于上述两种情形之一。若 $U$ 是 $S$-有理映射 $f$ 的定义域，
用 $f'$ 表示从 $\operatorname{Spec}(\mathscr{O}_x)$ 到 $Y$ 的有理映射；
考虑到（\hyperref[I.2.4.2-fr]{2.4.2}），它在
$U\cap\operatorname{Spec}(\mathscr{O}_x)$ 上与 $f$ 重合。称这个有理
映射\emph{由 $f$ 诱导}。

\begin{proposition}[7.2.9]
\label{I.7.2.9-fr}
设 $S$ 是局部 Noether 预概形，$X$ 是约化 $S$-预概形，$Y$ 是有限型
$S$-概形。还假设 $X$ 不可约或局部 Noether。设 $f$ 是从 $X$ 到 $Y$
的 $S$-有理映射，$x$ 是 $X$ 的一点。则 $f$ 在点 $x$ 处有定义的充要
条件是：由 $f$ 诱导的从 $\operatorname{Spec}(\mathscr{O}_x)$ 到 $Y$
的有理映射 $f'$（\hyperref[I.7.2.8-fr]{7.2.8}）是一个态射。
\end{proposition}

该条件显然必要（因为 $\operatorname{Spec}(\mathscr{O}_x)$ 包含在每个
含有 $x$ 的开集中），下面证明它充分。由
（\hyperref[I.6.5.1-fr]{6.5.1}），存在 $X$ 中 $x$ 的一个开邻域 $U$，
以及从 $U$ 到 $Y$ 的 $S$-态射 $g$，它在
$\operatorname{Spec}(\mathscr{O}_x)$ 上诱导 $f'$。若 $X$ 不可约，则
$U$ 在 $X$ 中稠密；由（\hyperref[I.7.2.3-fr]{7.2.3}），可以假设 $g$
是 $S$-有理映射。此外，$X$ 的泛点既属于
$\operatorname{Spec}(\mathscr{O}_x)$，也属于 $f$ 的定义域，所以 $f$
和 $g$ 在该点重合，因而在 $X$ 的某个非空开集中重合
（\hyperref[I.6.5.1-fr]{6.5.1}）。但由于 $f$ 和 $g$ 都是 $S$-有理
映射，它们相同（\hyperref[I.7.2.3-fr]{7.2.3}），故 $f$ 在 $x$ 处
有定义。

现在设 $X$ 局部 Noether；可以假设 $U$ 是 Noether 的。这时 $X$ 中
含有 $x$ 的不可约分支 $X_i$ 只有有限多个
（\hyperref[I.7.2.8.1-fr]{7.2.8.1}）；必要时把 $U$ 换成更小的开集，
可以假设只有这些分支与 $U$ 相交（因为 $U$ 为 Noether 的，与 $U$
相交的 $X$ 的不可约分支只有有限多个）。于是如上可见，$f$ 和 $g$ 在
每个 $X_i$ 的某个非空开集中重合。考虑到每个 $X_i$ 都包含在
$\overline U$ 中，考虑态射 $f_1$：它定义在
$U\cup(X-\overline U)$ 的一个稠密开集中，在 $U$ 中等于 $g$，在
$X-\overline U$ 与 $f$ 的定义域的交中等于 $f$。由于
$U\cup(X-\overline U)$ 在 $X$ 中稠密，$f_1$ 与 $f$ 在 $X$ 的某个稠密
开集中重合；又因为 $f$ 是有理映射，所以 $f$ 是 $f_1$ 的延拓
（\hyperref[I.7.2.3-fr]{7.2.3}），故它在点 $x$ 处有定义。

\subsection{有理函数层。}
\label{subsection:I.7.3-fr}

\begin{env}[7.3.1]
\label{I.7.3.1-fr}
设 $X$ 是预概形。对于每个开集 $U\subset X$，用 $R(U)$ 表示 $U$ 上的
有理函数环（\hyperref[I.7.1.3-fr]{7.1.3}）；它是一个
$\Gamma(U,\mathscr{O}_X)$-代数。此外，若 $V\subset U$ 是 $X$ 的另一个
开集，则 $\mathscr{O}_X$ 在 $U$ 的一个处处稠密开部分上的任意截面，
限制到 $V$ 后给出 $V$ 的一个处处稠密开部分上的截面；若两个截面在
$U$ 的一个处处稠密开部分上重合，则它们在 $V$ 上的限制也在 $V$ 的
一个处处稠密开部分上重合。由此定义一个代数双同态
$\rho_{V,U}:R(U)\to R(V)$；显然，若 $U\supset V\supset W$ 是 $X$ 的
三个开集，则
$\rho_{W,U}=\rho_{W,V}\circ\rho_{V,U}$；所以各 $R(U)$ 定义了 $X$ 上的
一个代数\emph{预层}。
\end{env}

\oldpage[I]{162}
\begin{definition}[7.3.2]
\label{I.7.3.2-fr}
把预概形 $X$ 上由各 $R(U)$ 组成的预层所伴随的
$\mathscr{O}_X$-代数称为它的\emph{有理函数层}，记作
$\mathscr{R}(X)$。
\end{definition}

对于任意预概形 $X$ 和任意开集 $U\subset X$，显然诱导层
$\mathscr{R}(X)|U$ 就是 $\mathscr{R}(U)$。

\begin{proposition}[7.3.3]
\label{I.7.3.3-fr}
设 $X$ 是预概形，它的不可约分支族 $(X_\lambda)$ 是局部有限的（特别地，
当 $X$ 的底空间局部 Noether 时即属此情形）。则
$\mathscr{O}_X$-模 $\mathscr{R}(X)$ 是拟凝聚的；对于 $X$ 的任意只与
有限多个分支 $X_\lambda$ 相交的开集 $U$，$R(U)$ 等于
$\Gamma(U,\mathscr{R}(X))$，并且典范地等同于所有满足
$U\cap X_\lambda\neq\varnothing$ 的 $X_\lambda$ 的泛点的局部环的直积。
\end{proposition}

显然可以只考虑 $X$ 只有有限多个不可约分支 $X_i$，其泛点为 $x_i$
（$1\leq i\leq n$）的情形。由（\hyperref[I.7.1.7-fr]{7.1.7}）可知，
$R(U)$ 典范地等同于所有满足 $U\cap X_i\neq\varnothing$ 的
$\mathscr{O}_{x_i}=R(X_i)$ 的直积。再来证明预层 $U\to R(U)$ 满足层公理，
这将证明 $R(U)=\Gamma(U,\mathscr{R}(X))$。事实上，由前面的讨论，它满足
(F 1)。为证明它满足 (F 2)，考虑用各开集 $V_\alpha\subset U$ 覆盖 $X$
的一个开集 $U$；若 $s_\alpha\in R(V_\alpha)$，并且对于任意一对指标，
$s_\alpha$ 和 $s_\beta$ 在 $V_\alpha\cap V_\beta$ 上的限制重合，则对于
每个满足 $U\cap X_i\neq\varnothing$ 的指标 $i$，所有满足
$V_\alpha\cap X_i\neq\varnothing$ 的 $s_\alpha$ 在 $R(X_i)$ 中的分量
都相同。用 $t_i$ 表示这个分量，则显然，以各 $t_i$ 为分量的 $R(U)$ 中
元素在每个 $V_\alpha$ 上的限制都是 $s_\alpha$。最后，为证明
$\mathscr{R}(X)$ 拟凝聚，可以只考虑 $X=\operatorname{Spec}(A)$ 为仿射的
情形；取 $U$ 为形如 $D(f)$ 的仿射开集，其中 $f\in A$，由上述结论和
定义（\hyperref[I.1.3.4-fr]{1.3.4}）可知
$\mathscr{R}(X)=\widetilde M$，其中 $M$ 是各 $A$-模 $A_{x_i}$ 的直和。

\begin{corollary}[7.3.4]
\label{I.7.3.4-fr}
设 $X$ 是只有有限多个不可约分支的约化预概形，并设 $X_i$
（$1\leq i\leq n$）是 $X$ 的约化闭子预概形，其底空间为 $X$ 的各不可约
分支（\hyperref[I.5.2.1-fr]{5.2.1}）。若 $h_i$ 是典范嵌入
$X_i\to X$，则 $\mathscr{R}(X)$ 是各 $\mathscr{O}_X$-代数
$(h_i)_*(\mathscr{R}(X_i))$ 的直积。
\end{corollary}

\begin{corollary}[7.3.5]
\label{I.7.3.5-fr}
若 $X$ 不可约，则任意拟凝聚 $\mathscr{R}(X)$-模 $\mathscr{F}$ 都是常值层。
\end{corollary}

只需证明每个 $x\in X$ 都有一个邻域 $U$，使 $\mathscr{F}|U$ 为常值层
（\hyperref[0.3.6.2-fr]{0, 3.6.2}）；换言之，可以归结到 $X$ 为仿射的
情形。还可以假设 $\mathscr{F}$ 是某个同态
$(\mathscr{R}(X))^{(I)}\to(\mathscr{R}(X))^{(J)}$
的余核（\hyperref[0.5.1.3-fr]{0, 5.1.3}），于是只需证明
$\mathscr{R}(X)$ 是常值层；但这是显然的，因为对于每个非空开集 $U$，
都有 $\Gamma(U,\mathscr{R}(X))=R(X)$，而 $U$ 含有 $X$ 的泛点。

\begin{corollary}[7.3.6]
\label{I.7.3.6-fr}
若 $X$ 不可约，则对于任意拟凝聚 $\mathscr{O}_X$-模 $\mathscr{F}$，
$\mathscr{F}\otimes_{\mathscr{O}_X}\mathscr{R}(X)$ 是常值层；若 $X$ 还约化
（从而是整的），则 $\mathscr{F}\otimes_{\mathscr{O}_X}\mathscr{R}(X)$
同构于形如 $(\mathscr{R}(X))^{(I)}$ 的层。
\end{corollary}

第二项断言来自 $R(X)$ 这时是一个域。

\begin{proposition}[7.3.7]
\label{I.7.3.7-fr}
假设预概形 $X$ 是局部整的或局部
\oldpage[I]{163}
Noether 的。则 $\mathscr{R}(X)$ 是拟凝聚 $\mathscr{O}_X$-代数；若 $X$
还约化（当 $X$ 局部整时即属此情形），则典范同态
$\mathscr{O}_X\to\mathscr{R}(X)$ 是单射。
\end{proposition}

由于问题是局部的，第一项断言由（\hyperref[I.7.3.3-fr]{7.3.3}）得出；
第二项立即由（\hyperref[I.7.2.3-fr]{7.2.3}）得出。

% EGA1-ERR-0053: preserve the printed I.7.3.8 reading; EGA II Errata et Addenda replaces the whole paragraph.
\begin{env}[7.3.8]
\label{I.7.3.8-fr}
设 $X$、$Y$ 是两个各自只有有限多个不可约分支的预概形，
$f:X\to Y$ 是一个态射，并且它在 $X$ 的不可约分支泛点集合上的限制，
满射到 $Y$ 的不可约分支泛点集合上。则有
\[
  f^*(\mathscr{R}(Y))=\mathscr{R}(X).
  \tag{7.3.8.1}\label{I.7.3.8.1-fr}
\]
\end{env}

事实上，由（\hyperref[I.7.3.3-fr]{7.3.3}），可以归结到 $X$ 和 $Y$
均不可约的情形；设它们的泛点分别为 $x$、$y$，且 $f(x)=y$。于是
\[
  (f^*(\mathscr{R}(Y)))_x
  =\mathscr{O}_y\otimes_{\mathscr{O}_y}\mathscr{O}_x
  =\mathscr{O}_x
\]
（\hyperref[0.4.3.1-fr]{0, 4.3.1}）；由
（\hyperref[I.7.3.5-fr]{7.3.5}），这就证明了 (7.3.8.1)。

\subsection{挠层和无挠层。}
\label{subsection:I.7.4-fr}

\begin{env}[7.4.1]
\label{I.7.4.1-fr}
设 $X$ 是\emph{整的}预概形。对于任意 $\mathscr{O}_X$-模
$\mathscr{F}$，典范同态 $\mathscr{O}_X\to\mathscr{R}(X)$ 通过张量化定义
一个同态（仍称为\emph{典范}同态）
\[
  \mathscr{F}\to
  \mathscr{F}\otimes_{\mathscr{O}_X}\mathscr{R}(X)
\]
它在每个茎上就是从 $\mathscr{F}_x$ 到
$\mathscr{F}_x\otimes_{\mathscr{O}_x}R(X)$ 的同态 $z\to z\otimes 1$。
这个同态的\emph{核 $\mathscr{T}$} 是 $\mathscr{F}$ 的一个
$\mathscr{O}_X$-子模，称为 $\mathscr{F}$ 的\emph{挠层}；当
$\mathscr{F}$ 拟凝聚时，它也是拟凝聚的
（\hyperref[I.4.1.1-fr]{4.1.1} 和 \hyperref[I.7.3.6-fr]{7.3.6}）。若
$\mathscr{T}=0$，则称 $\mathscr{F}$ \emph{无挠}；若
$\mathscr{T}=\mathscr{F}$，则称 $\mathscr{F}$ 为\emph{挠层}。对于任意
$\mathscr{O}_X$-模 $\mathscr{F}$，$\mathscr{F}/\mathscr{T}$ 都无挠。
由（\hyperref[I.7.3.5-fr]{7.3.5}）可得：
\end{env}

\begin{proposition}[7.4.2]
\label{I.7.4.2-fr}
若 $X$ 是整预概形，则任意无挠拟凝聚 $\mathscr{O}_X$-模
$\mathscr{F}$ 都同构于形如 $(\mathscr{R}(X))^{(I)}$ 的常值层的一个子层
$\mathscr{G}$，并且该常值层作为 $\mathscr{R}(X)$-模由 $\mathscr{G}$ 生成。
\end{proposition}

$I$ 的基数称为 $\mathscr{F}$ 的\emph{秩}；对于 $X$ 的任意非空仿射开集
$U$，$\mathscr{F}$ 的秩等于 $\Gamma(U,\mathscr{F})$ 作为
$\Gamma(U,\mathscr{O}_X)$-模的秩；考虑到 $U$ 含有 $X$ 的泛点，立即可见
这一点。特别地：

\begin{corollary}[7.4.3]
\label{I.7.4.3-fr}
在整预概形 $X$ 上，任意秩为 $1$ 的无挠拟凝聚 $\mathscr{O}_X$-模
（特别地，任意可逆 $\mathscr{O}_X$-模）都同构于 $\mathscr{R}(X)$ 的一个
$\mathscr{O}_X$-子模，反之亦然。
\end{corollary}

\begin{corollary}[7.4.4]
\label{I.7.4.4-fr}
设 $X$ 是整预概形，$\mathscr{L}$、$\mathscr{L}'$ 是两个无挠
$\mathscr{O}_X$-模，$f$（相应地 $f'$）是 $\mathscr{L}$（相应地
$\mathscr{L}'$）在 $X$ 上的一个截面。$f\otimes f'=0$ 的充要条件是截面
$f$、$f'$ 中至少有一个为零。
\end{corollary}

设 $x$ 是 $X$ 的泛点；依假设
$(f\otimes f')_x=f_x\otimes f'_x=0$。由于 $\mathscr{L}_x$ 和
$\mathscr{L}'_x$ 等同于域 $\mathscr{O}_x$ 的某些
$\mathscr{O}_x$-子模，上式蕴含 $f_x=0$ 或 $f'_x=0$，从而蕴含
$f=0$ 或 $f'=0$，因为 $\mathscr{L}$ 和 $\mathscr{L}'$ 无挠
（\hyperref[I.7.3.5-fr]{7.3.5}）。

\begin{proposition}[7.4.5]
\label{I.7.4.5-fr}
设 $X$、$Y$ 是两个整预概形，$f:X\to Y$ 是支配态射。对于任意无挠
拟凝聚 $\mathscr{O}_X$-模 $\mathscr{F}$，$f_*(\mathscr{F})$ 是无挠
$\mathscr{O}_Y$-模。
\end{proposition}

\oldpage[I]{164}
\begin{proof}
由于 $f_*$ 左正合（\hyperref[0.4.2.1-fr]{0, 4.2.1}），由
（\hyperref[I.7.4.2-fr]{7.4.2}），只需在
$\mathscr{F}=(\mathscr{R}(X))^{(I)}$ 时证明该命题。$Y$ 的任意非空开集
$U$ 都含有 $Y$ 的泛点，所以 $f^{-1}(U)$ 含有 $X$ 的泛点
（\hyperref[0.2.1.5-fr]{0, 2.1.5}），从而
\[
  \Gamma(U,f_*(\mathscr{F}))
  =\Gamma(f^{-1}(U),\mathscr{F})=(R(X))^{(I)}~;
\]
换言之，$f_*(\mathscr{F})$ 是以 $(R(X))^{(I)}$ 为茎的常值层，视为
$\mathscr{R}(Y)$-模时显然无挠。
\end{proof}

\begin{proposition}[7.4.6]
\label{I.7.4.6-fr}
设 $X$ 是整预概形，$x$ 是它的泛点。对于任意有限型拟凝聚
$\mathscr{O}_X$-模 $\mathscr{F}$，下列条件等价：a) $\mathscr{F}$ 是挠层；
b) $\mathscr{F}_x=0$；c) $\operatorname{Supp}(\mathscr{F})\neq X$。
\end{proposition}

由（\hyperref[I.7.3.5-fr]{7.3.5}）和
（\hyperref[I.7.4.1-fr]{7.4.1}），关系 $\mathscr{F}_x=0$ 与
$\mathscr{F}\otimes_{\mathscr{O}_X}\mathscr{R}(X)=0$ 等价，所以 a) 和 b)
等价；另一方面，$\operatorname{Supp}(\mathscr{F})$ 是 $X$ 中的闭集
（\hyperref[0.5.2.2-fr]{0, 5.2.2}），而 $X$ 的每个非空开集都含有 $x$，
所以 b) 和 c) 等价。

\begin{env}[7.4.7]
\label{I.7.4.7-fr}
把（\hyperref[I.7.4.1-fr]{7.4.1}）的定义（滥用语言地）推广到 $X$ 是
只有\emph{有限}多个不可约分支的\emph{约化}预概形的情形；这时由
（\hyperref[I.7.3.4-fr]{7.3.4}）可知，
（\hyperref[I.7.4.6-fr]{7.4.6}）中 a) 与 c) 的等价性对这样的预概形
仍然成立。
\end{env}
\section{Chevalley概形}
\label{section:I.8-fr}

\subsection{同源的局部环。}
\label{subsection:I.8.1-fr}

对于任意局部环 $A$，我们用 $\mathfrak m(A)$ 表示 $A$ 的极大理想。

\begin{lemma}[8.1.1]
\label{I.8.1.1-fr}
设 $A$、$B$ 是两个局部环，且 $A\subset B$；则下列条件等价：
\begin{enumerate}
  \item[(i)] $\mathfrak m(B)\cap A=\mathfrak m(A)$；
  \item[(ii)] $\mathfrak m(A)\subset\mathfrak m(B)$；
  \item[(iii)] $1$ 不属于 $B$ 中由 $\mathfrak m(A)$ 生成的理想。
\end{enumerate}
\end{lemma}

显然，(i) 蕴含 (ii)，而 (ii) 蕴含 (iii)；最后，若 (iii) 成立，则
$\mathfrak m(B)\cap A$ 包含 $\mathfrak m(A)$ 且不包含 $1$，所以它等于
$\mathfrak m(A)$。

当（\hyperref[I.8.1.1-fr]{8.1.1}）的等价条件成立时，称 $B$ \emph{支配}
$A$；这等价于说嵌入 $A\to B$ 是一个\emph{局部}同态。显然，在环 $R$ 的
局部子环全体中，支配关系是一种序关系。

\begin{env}[8.1.2]
\label{I.8.1.2-fr}
现在考虑一个\emph{域} $R$。对于 $R$ 的任意子环 $A$，用 $L(A)$ 表示局部环
$A_\mathfrak p$ 的集合，其中 $\mathfrak p$ 遍历 $A$ 的素谱；把这些局部环等同于
$R$ 中包含 $A$ 的子环。由于
$\mathfrak p=(\mathfrak p A_\mathfrak p)\cap A$，映射
$\mathfrak p\mapsto A_\mathfrak p$ 给出从 $\operatorname{Spec}(A)$ 到 $L(A)$
的双射。
\end{env}

\begin{lemma}[8.1.3]
\label{I.8.1.3-fr}
设 $R$ 是一个域，$A$ 是 $R$ 的一个子环。$R$ 的局部子环 $M$ 支配某个环
$A_\mathfrak p\in L(A)$ 的充分必要条件是 $A\subset M$；此时，被 $M$ 支配的
局部环 $A_\mathfrak p$ 是唯一的，并且对应于
$\mathfrak p=\mathfrak m(M)\cap A$。
\end{lemma}

事实上，若 $M$ 支配 $A_\mathfrak p$，则由（\hyperref[I.8.1.1-fr]{8.1.1}），
有 $\mathfrak m(M)\cap A_\mathfrak p=\mathfrak p A_\mathfrak p$，从而得到
$\mathfrak p$ 的唯一性；另一方面，若 $A\subset M$，则
$\mathfrak m(M)\cap A=\mathfrak p$ 是 $A$ 中的素理想，并且由于
$A-\mathfrak p\subset M$，有 $A_\mathfrak p\subset M$ 以及
$\mathfrak p A_\mathfrak p\subset\mathfrak m(M)$，所以 $M$ 支配
$A_\mathfrak p$。

\oldpage[I]{165}
\begin{lemma}[8.1.4]
\label{I.8.1.4-fr}
设 $R$ 是一个域，$M$、$N$ 是 $R$ 的两个局部子环，$P$ 是 $R$ 中由
$M\cup N$ 生成的子环。下列条件等价：
\begin{enumerate}
  \item[(i)] 存在 $P$ 的一个素理想 $\mathfrak p$，使得
    $\mathfrak m(M)=\mathfrak p\cap M$，
    $\mathfrak m(N)=\mathfrak p\cap N$。
  \item[(ii)] $P$ 中由 $\mathfrak m(M)\cup\mathfrak m(N)$ 生成的理想
    $\mathfrak a$ 不等于 $P$。
  \item[(iii)] 存在 $R$ 的一个局部子环 $Q$，它同时支配 $M$ 和 $N$。
\end{enumerate}
\end{lemma}

显然 (i) 蕴含 (ii)；反之，若 $\mathfrak a\ne P$，则 $\mathfrak a$ 包含于
$P$ 的某个极大理想 $\mathfrak n$ 中；由于 $1\notin\mathfrak n$，
$\mathfrak n\cap M$ 包含 $\mathfrak m(M)$ 且不等于 $M$，所以
$\mathfrak n\cap M=\mathfrak m(M)$；同理，
$\mathfrak n\cap N=\mathfrak m(N)$。显然，若 $Q$ 支配 $M$ 和 $N$，则
$P\subset Q$，并且
\[
  \mathfrak m(M)=\mathfrak m(Q)\cap M
  =(\mathfrak m(Q)\cap P)\cap M,\qquad
  \mathfrak m(N)=(\mathfrak m(Q)\cap P)\cap N,
\]
所以 (iii) 蕴含 (i)；反过来取 $Q=P_\mathfrak p$ 即得 (i) 蕴含 (iii)。

当（\hyperref[I.8.1.4-fr]{8.1.4}）的条件成立时，沿用 C.~Chevalley 的
术语，称局部环 $M$ 与 $N$ \emph{同源}。

\begin{proposition}[8.1.5]
\label{I.8.1.5-fr}
设 $A$、$B$ 是域 $R$ 的两个子环，$C$ 是 $R$ 中由 $A\cup B$ 生成的子环。
下列条件等价：
\begin{enumerate}
  \item[(i)] 对于包含 $A$ 和 $B$ 的任意局部环 $Q$，置
    $\mathfrak p=\mathfrak m(Q)\cap A$、
    $\mathfrak q=\mathfrak m(Q)\cap B$，则
    $A_\mathfrak p=B_\mathfrak q$。
  \item[(ii)] 对于 $C$ 的任意素理想 $\mathfrak r$，置
    $\mathfrak p=\mathfrak r\cap A$、$\mathfrak q=\mathfrak r\cap B$，则
    $A_\mathfrak p=B_\mathfrak q$。
  \item[(iii)] 若 $M\in L(A)$ 与 $N\in L(B)$ 同源，则它们相同。
  \item[(iv)] 有 $L(A)\cap L(B)=L(C)$。
\end{enumerate}
\end{proposition}

引理（\hyperref[I.8.1.3-fr]{8.1.3}）和
（\hyperref[I.8.1.4-fr]{8.1.4}）证明了 (i) 与 (iii) 等价；把 (i) 应用于
$Q=C_\mathfrak r$，显然可知 (i) 蕴含 (ii)；反过来，(ii) 蕴含 (i)，因为若
$Q$ 包含 $A\cup B$，它也包含 $C$，而置
$\mathfrak r=\mathfrak m(Q)\cap C$，由
（\hyperref[I.8.1.3-fr]{8.1.3}）有
$\mathfrak p=\mathfrak r\cap A$ 和 $\mathfrak q=\mathfrak r\cap B$。
立即可知 (iv) 蕴含 (i)，因为若 $Q$ 包含 $A\cup B$，则由
（\hyperref[I.8.1.3-fr]{8.1.3}），它支配某个局部环
$C_\mathfrak r\in L(C)$；依假设有
$C_\mathfrak r\in L(A)\cap L(B)$，而
（\hyperref[I.8.1.1-fr]{8.1.1}）与
（\hyperref[I.8.1.3-fr]{8.1.3}）证明
$C_\mathfrak r=A_\mathfrak p=B_\mathfrak q$。最后证明 (iii) 蕴含 (iv)。
设 $Q\in L(C)$；由（\hyperref[I.8.1.3-fr]{8.1.3}），$Q$ 支配某个
$M\in L(A)$ 和某个 $N\in L(B)$，所以 $M$ 与 $N$ 同源，依假设它们相同。
于是 $C\subset M$，由（\hyperref[I.8.1.3-fr]{8.1.3}），$M$ 支配某个
$Q'\in L(C)$，因而 $Q$ 支配 $Q'$；再由
（\hyperref[I.8.1.3-fr]{8.1.3}）必有 $Q=Q'=M$，所以
$Q\in L(A)\cap L(B)$。反过来，若 $Q\in L(A)\cap L(B)$，则
$C\subset Q$，所以由（\hyperref[I.8.1.3-fr]{8.1.3}），$Q$ 支配某个
$Q''\in L(C)\subset L(A)\cap L(B)$；$Q$ 与 $Q''$ 同源，故二者相同，
于是 $Q''=Q\in L(C)$，证明完毕。

\subsection{整概形的局部环。}
\label{subsection:I.8.2-fr}

\begin{env}[8.2.1]
\label{I.8.2.1-fr}
设 $X$ 是一个\emph{整}预概形，$R$ 是它的有理函数域，并且等同于
$X$ 的一般点 $a$ 处的局部环；对于任意 $x\in X$，我们知道
$\mathscr{O}_x$ 可典范地识别为 $R$ 的一个子环
（\hyperref[I.7.1.5-fr]{7.1.5}）；对于任意有理函数 $f\in R$，$f$ 的定义域
$\delta(f)$ 是由满足 $f\in\mathscr{O}_x$ 的 $x\in X$ 构成的开集。因此，由
（\hyperref[I.7.2.6-fr]{7.2.6}），对于任意开集 $U\subset X$，有
\begin{equation}
  \Gamma(U,\mathscr{O}_X)=\bigcap_{x\in U}\mathscr{O}_x.
  \tag{8.2.1.1}\label{I.8.2.1.1-fr}
\end{equation}
\end{env}

\oldpage[I]{166}
\begin{proposition}[8.2.2]
\label{I.8.2.2-fr}
设 $X$ 是整预概形，$R$ 是它的有理函数域。$X$ 为概形的充分必要条件是：
对 $X$ 中的点 $x$、$y$，若局部环 $\mathscr{O}_x$ 与 $\mathscr{O}_y$ 同源
（\hyperref[I.8.1.4-fr]{8.1.4}），则 $x=y$。
\end{proposition}

假设这一条件成立，证明 $X$ 是分离的。设 $U$、$V$ 是 $X$ 中两个不同的仿射开集，
其环分别为 $A$、$B$，并将它们识别为 $R$ 的子环；于是由
（\hyperref[I.8.1.2-fr]{8.1.2}），$U$（相应地 $V$）可识别为 $L(A)$
（相应地 $L(B)$），并且由假设及（\hyperref[I.8.1.5-fr]{8.1.5}），若 $C$
是 $R$ 中由 $A\cup B$ 生成的子环，则
$W=U\cap V$ 可识别为 $L(A)\cap L(B)=L(C)$。此外，我们知道
（\cite{I-1}，p.~5-03，命题 4 \emph{bis}），$R$ 的每个子环 $E$ 都等于
属于 $L(E)$ 的局部环的交；因而 $C$ 可识别为 $z\in W$ 的局部环
$\mathscr{O}_z$ 的交，换言之，由（\hyperref[I.8.2.1.1-fr]{8.2.1.1}）可识别为
$\Gamma(W,\mathscr{O}_X)$。考虑 $X$ 在 $W$ 上诱导的子预概形；恒等同态
$\varphi:C\to\Gamma(W,\mathscr{O}_X)$ 对应于
（\hyperref[I.2.2.4-fr]{2.2.4}）中的态射
$\Phi=(\psi,\theta):W\to\operatorname{Spec}(C)$；我们将证明 $\Phi$ 是预概形的
\emph{同构}，从而 $W$ 是\emph{仿射}开集。将 $W$ 识别为
$L(C)=\operatorname{Spec}(C)$，可知 $\psi$ 是\emph{双射}。另一方面，对于每个
$x\in W$，若 $\mathfrak r=\mathfrak m_x\cap C$，则
$\theta_x^\sharp$ 是嵌入 $C_\mathfrak r\to\mathscr{O}_x$；按定义，
$C_\mathfrak r$ 已识别为 $\mathscr{O}_x$，所以 $\theta_x^\sharp$ 是双射。于是只需
证明 $\psi$ 是一个\emph{同胚}，也就是证明对每个闭集
$F\subset W$，$\psi(F)$ 是 $\operatorname{Spec}(C)$ 中的闭集。事实上，$F$ 是
$U$ 中形如 $V(\mathfrak a)$ 的闭集在 $W$ 上的迹，其中 $\mathfrak a$ 是 $A$ 的
一个理想；我们证明 $\psi(F)=V(\mathfrak a C)$，这就证明了断言。包含
$\mathfrak a C$ 的 $C$ 的素理想正是包含 $\mathfrak a$ 的 $C$ 的素理想，因而
它们是形如 $\psi(x)=\mathfrak m_x\cap C$ 的理想，其中
$\mathfrak a\subset\mathfrak m_x$ 且 $x\in W$；对于 $x\in U$，
$\mathfrak a\subset\mathfrak m_x$ 等价于
$x\in V(\mathfrak a)=W\cap F$，所以确有
$\psi(F)=V(\mathfrak a C)$。

由此得到 $X$ 是分离的，因为 $U\cap V$ 是仿射的，并且其环 $C$ 由 $U$、$V$
的环的并 $A\cup B$ 生成（\hyperref[I.5.5.6-fr]{5.5.6}）。

反过来，假设 $X$ 是分离的，并设 $x$、$y$ 是 $X$ 中满足
$\mathscr{O}_x$ 与 $\mathscr{O}_y$ 同源的两个点。令 $U$（相应地 $V$）为包含
$x$（相应地 $y$）的仿射开集，其环为 $A$（相应地 $B$）；我们知道 $U\cap V$
是仿射的，且其环 $C$ 由 $A\cup B$ 生成（\hyperref[I.5.5.6-fr]{5.5.6}）。
若 $\mathfrak p=\mathfrak m_x\cap A$、$\mathfrak q=\mathfrak m_y\cap B$，则
$A_\mathfrak p=\mathscr{O}_x$、$B_\mathfrak q=\mathscr{O}_y$；由于
$A_\mathfrak p$ 与 $B_\mathfrak q$ 同源，存在 $C$ 的素理想 $\mathfrak r$，使得
$\mathfrak p=\mathfrak r\cap A$、$\mathfrak q=\mathfrak r\cap B$
（\hyperref[I.8.1.4-fr]{8.1.4}）。但由于 $U\cap V$ 是仿射的，存在
$z\in U\cap V$ 使得 $\mathfrak r=\mathfrak m_z\cap C$；显然 $x=z$ 且
$y=z$，所以 $x=y$。

\begin{corollary}[8.2.3]
\label{I.8.2.3-fr}
设 $X$ 是整概形，$x$、$y$ 是 $X$ 的两个点。$x\in\overline{\{y\}}$ 的充分必要条件是
$\mathscr{O}_x\subset\mathscr{O}_y$，换言之，在 $x$ 处有定义的每个有理函数也在
$y$ 处有定义。
\end{corollary}

必要性显然，因为有理函数 $f\in R$ 的定义域 $\delta(f)$ 是开集；现在证明充分性。
若 $\mathscr{O}_x\subset\mathscr{O}_y$，则由（\hyperref[I.8.1.3-fr]{8.1.3}），
$\mathscr{O}_x$ 有一个素理想 $\mathfrak p$，使得 $\mathscr{O}_y$ 支配
$(\mathscr{O}_x)_\mathfrak p$；另一方面，由（\hyperref[I.2.4.2-fr]{2.4.2}），
存在 $z\in X$，使得 $x\in\overline{\{z\}}$ 且
$\mathscr{O}_z=(\mathscr{O}_x)_\mathfrak p$；由于 $\mathscr{O}_z$ 与
$\mathscr{O}_y$ 同源，由（\hyperref[I.8.2.2-fr]{8.2.2}）有 $z=y$，证毕。

\begin{corollary}[8.2.4]
\label{I.8.2.4-fr}
若 $X$ 是整概形，则映射 $x\mapsto\mathscr{O}_x$ 是单射；换言之，若 $x$、$y$
是 $X$ 中的两个不同点，则存在一个有理函数在其中一点处有定义而在另一点处
没有定义。
\end{corollary}

\oldpage[I]{167}
这由（\hyperref[I.8.2.3-fr]{8.2.3}）和公理 $(T_0)$
（\hyperref[I.2.1.4-fr]{2.1.4}）立即得到。

\begin{corollary}[8.2.5]
\label{I.8.2.5-fr}
设 $X$ 是整概形，且其底空间是 Noether 的；当 $f$ 遍历 $X$ 上有理函数的
域 $R$ 时，集合 $\delta(f)$ 生成 $X$ 的拓扑。
\end{corollary}

事实上，$X$ 的每个闭集都是有限个不可约闭集的并，也就是形如
$\overline{\{y\}}$ 的集合（\hyperref[I.2.1.5-fr]{2.1.5}）。若
$x\notin\overline{\{y\}}$，则由（\hyperref[I.8.2.3-fr]{8.2.3}），存在有理函数
$f$ 在 $x$ 处有定义而在 $y$ 处没有定义；换言之，$x\in\delta(f)$，且
$\delta(f)$ 与 $\overline{\{y\}}$ 不相交。因此
$\overline{\{y\}}$ 的补集是形如 $\delta(f)$ 的集合的并；根据前面的观察，
$X$ 的每个开集都是形如 $\delta(f)$ 的开集的有限交的并。

\begin{env}[8.2.6]
\label{I.8.2.6-fr}
推论（\hyperref[I.8.2.5-fr]{8.2.5}）表明，$X$ 的拓扑完全由局部环族
$(\mathscr{O}_x)_{x\in X}$ 决定，其中 $R$ 是这些局部环的分式域。换一种说法，
$X$ 的闭集按如下方式定义：给定 $X$ 的一个有限子集
$\{x_1,\ldots,x_n\}$，考虑由满足
$\mathscr{O}_y\subset\mathscr{O}_{x_i}$（至少对一个指标 $i$）的
$y\in X$ 组成的集合；对 $\{x_1,\ldots,x_n\}$ 的所有选择所得到的这些集合，
正好是 $X$ 的闭集。此外，一旦 $X$ 的拓扑已知，结构层
$\mathscr{O}_X$ 也由 $\mathscr{O}_x$ 族唯一决定，因为
$\Gamma(U,\mathscr{O}_X)=\bigcap_{x\in U}\mathscr{O}_x$
（\hyperref[I.8.2.1.1-fr]{8.2.1.1}）。所以，当 $X$ 是底空间为 Noether 空间的整
概形时，族 $(\mathscr{O}_x)_{x\in X}$ 完全决定预概形 $X$。
\end{env}

\begin{proposition}[8.2.7]
\label{I.8.2.7-fr}
设 $X$、$Y$ 是两个整概形，$f:X\to Y$ 是支配态射
（\hyperref[I.2.2.6-fr]{2.2.6}），$K$（相应地 $L$）是 $X$（相应地 $Y$）上的
有理函数域。则 $L$ 可识别为 $K$ 的一个子域，并且对每个 $x\in X$，
$\mathscr{O}_{f(x)}$ 是被 $\mathscr{O}_x$ 支配的 $Y$ 的唯一局部环。
\end{proposition}

事实上，若 $f=(\psi,\theta)$，且 $a$ 是 $X$ 的一般点，则
$\psi(a)$ 是 $Y$ 的一般点
（\hyperref[0.2.1.5-fr]{0, 2.1.5}）；因此
$\theta_a^\sharp$ 是从域 $L=\mathscr{O}_{\psi(a)}$ 到域 $K=\mathscr{O}_a$ 的单同态。
由于 $Y$ 的每个非空仿射开集 $U$ 都包含 $\psi(a)$，由
（\hyperref[I.2.2.4-fr]{2.2.4}）可知，对应于 $f$ 的同态
$\Gamma(U,\mathscr{O}_Y)\to\Gamma(\psi^{-1}(U),\mathscr{O}_X)$
是 $\theta_a^\sharp$ 在 $\Gamma(U,\mathscr{O}_Y)$ 上的限制。因此，对每个
$x\in X$，$\theta_x^\sharp$ 是 $\theta_a^\sharp$ 在
$\mathscr{O}_{\psi(x)}$ 上的限制，因而是单同态。此外，我们知道
$\theta_x^\sharp$ 是局部同态；所以，若通过 $\theta_a^\sharp$ 将 $L$ 识别为
$K$ 的子域，则由（\hyperref[I.8.1.1-fr]{8.1.1}），
$\mathscr{O}_{\psi(x)}$ 被 $\mathscr{O}_x$ 支配；而且它是被
$\mathscr{O}_x$ 支配的 $Y$ 的唯一局部环，因为 $Y$ 中两个同源的局部环必相同
（\hyperref[I.8.2.2-fr]{8.2.2}）。

\begin{proposition}[8.2.8]
\label{I.8.2.8-fr}
设 $X$ 是不可约预概形，$f:X\to Y$ 是局部浸入（相应地，局部同构）；此外假设
态射 $f$ 是分离的。则 $f$ 是浸入（相应地，开浸入）。
\end{proposition}

设 $f=(\psi,\theta)$；在两种情形中，只需证明 $\psi$ 是从 $X$ 到
$\psi(X)$ 的\emph{同胚}
（\hyperref[I.4.5.3-fr]{4.5.3}）。以 $f_{\mathrm{red}}$ 替代 $f$
（\hyperref[I.5.1.6-fr]{5.1.6} 及
\hyperref[I.5.5.1-fr]{5.5.1, (vi)}），可假设 $X$、$Y$ 都是约化的。若 $Y'$
是 $Y$ 的既约闭子预概形，其底空间为 $\overline{\psi(X)}$，则 $f$ 分解为
$X\xrightarrow{f'}Y'\xrightarrow{j}Y$，其中 $j$ 是典范嵌入
（\hyperref[I.5.2.2-fr]{5.2.2}）。由（\hyperref[I.5.5.1-fr]{5.5.1, (v)})），
$f'$ 仍是分离态射；此外，$f'$ 仍是局部浸入（相应地，局部同构），因为这一问题
在 $X$ 和 $Y$ 上都是局部的，可以将证明限制到 $f$ 是闭浸入（相应地，开浸入）
的情形，而此时断言由（\hyperref[I.4.2.2-fr]{4.2.2}）立即得到。

\oldpage[I]{168}
于是可以假设 $f$ 是\emph{支配}态射，这导致 $Y$ 也是不可约的
（\hyperref[0.2.1.5-fr]{0, 2.1.5}），从而 $X$、$Y$ 都是\emph{整}的。此外，由于
问题在 $Y$ 上是局部的，可以假设 $Y$ 是仿射概形；由于 $f$ 是分离的，$X$ 是概形
（\hyperref[I.5.5.1-fr]{5.5.1, (ii)}），于是我们最终处于
（\hyperref[I.8.2.7-fr]{8.2.7}）的假设下。对每个 $x\in X$，
$\theta_x^\sharp$ 是单射；但 $f$ 是局部浸入的假设意味着
$\theta_x^\sharp$ 是满射（\hyperref[I.4.2.2-fr]{4.2.2}），所以
$\theta_x^\sharp$ 是双射；换言之（使用
（\hyperref[I.8.2.7-fr]{8.2.7}）中的识别），有
$\mathscr{O}_{\psi(x)}=\mathscr{O}_x$。由
（\hyperref[I.8.2.4-fr]{8.2.4}）这意味着 $\psi$ 是\emph{单射}，当 $f$ 是局部同构时，
这已经证明了命题（\hyperref[I.4.5.3-fr]{4.5.3}）。若只假设 $f$ 是局部浸入，则对
每个 $x\in X$，存在 $X$ 中 $x$ 的开邻域 $U$ 和 $Y$ 中 $\psi(x)$ 的开邻域 $V$，
使得 $\psi$ 在 $U$ 上的限制是从 $U$ 到 $V$ 的一个\emph{闭集}上的同胚。但是
$U$ 在 $X$ 中稠密，所以 $\psi(U)$ 在 $Y$ 中稠密，从而\emph{尤其}在 $V$ 中稠密，
这证明 $\psi(U)=V$；由于 $\psi$ 是单射，$\psi^{-1}(V)=U$，于是 $\psi$ 是从
$X$ 到 $\psi(X)$ 的同胚。

\subsection{Chevalley概形。}
\label{subsection:I.8.3-fr}

\begin{env}[8.3.1]
\label{I.8.3.1-fr}
设 $X$ 是一个\emph{Noether整概形}，$R$ 是它的有理函数域；令 $X'$ 为局部子环
$\mathscr{O}_x\subset R$ 的集合，其中 $x$ 遍历 $X$。集合 $X'$ 满足下列三个条件：
\begin{enumerate}
  \item[(Sch. 1)] \emph{对于每个 $M\in X'$，$R$ 是 $M$ 的分式域。}
  \item[(Sch. 2)] \emph{存在 $R$ 的有限个 Noether 子环 $A_i$，使得
    $X'=\bigcup_i L(A_i)$，并且对于任意指标对 $i$、$j$，由
    $A_i\cup A_j$ 在 $R$ 中生成的子环 $A_{ij}$ 是 $A_i$ 上的有限型代数。}
  \item[(Sch. 3)] \emph{$X'$ 中任意两个同源的元素 $M$、$N$ 都相同。}
\end{enumerate}
\end{env}

事实上，我们在（\hyperref[I.8.2.1-fr]{8.2.1}）中已经看到 (Sch. 1) 成立，而
(Sch. 3) 由（\hyperref[I.8.2.2-fr]{8.2.2}）得出。为证明 (Sch. 2)，只需用有限个
仿射开集 $U_i$ 覆盖 $X$，这些开集的环是 Noether 环，并取
$A_i=\Gamma(U_i,\mathscr{O}_X)$；$X$ 是概形这一假设意味着 $U_i\cap U_j$ 是
仿射的，并且
$\Gamma(U_i\cap U_j,\mathscr{O}_X)=A_{ij}$
（\hyperref[I.5.5.6-fr]{5.5.6}）；此外，由于空间 $U_i$ 是 Noether 的，
浸入 $U_i\cap U_j\to U_i$ 是有限型的
（\hyperref[I.6.3.5-fr]{6.3.5}），所以 $A_{ij}$ 是有限型
$A_i$-代数（\hyperref[I.6.3.3-fr]{6.3.3}）。

\begin{env}[8.3.2]
\label{I.8.3.2-fr}
其公理为 (Sch. 1)、(Sch. 2) 和 (Sch. 3) 的结构，推广了 C.~Chevalley 意义下的
“概形”；Chevalley 还假设 $R$ 是域 $K$ 的有限型扩张，且 $A_i$ 是有限型
$K$-代数（这使 (Sch. 2) 的一部分变得多余）
\cite{I-1}。反过来，若集合 $X'$ 上有这样一种结构，则可以利用
（\hyperref[I.8.2.6-fr]{8.2.6}）中的备注赋予它一个整概形 $X$：$X$ 的底空间等于
带有（\hyperref[I.8.2.6-fr]{8.2.6}）所定义拓扑的 $X'$，并带有结构层
$\mathscr{O}_X$

\oldpage[I]{169}
使得对于每个开集 $U\subset X$，
$\Gamma(U,\mathscr{O}_X)=\bigcap_{x\in U}\mathscr{O}_x$，限制同态按显然方式定义。
我们把验证由此确实得到一个整概形的工作留给读者；该概形的局部环正是 $X'$ 的
元素；在下文中我们不会再使用这一结果。
\end{env}
\section{准凝聚层补充}
\label{section:I.9-fr}

\subsection{准凝聚层的张量积}
\label{subsection:I.9.1-fr}

\begin{proposition}[9.1.1]
\label{I.9.1.1-fr}
设 $X$ 是一个预概形（相应地，是一个局部诺特预概形）。设
$\mathscr{F}$ 与 $\mathscr{G}$ 是两个准凝聚的
$\mathscr{O}_X$-模层（相应地，凝聚的）；则
$\mathscr{F}\otimes_{\mathscr{O}_X}\mathscr{G}$ 是准凝聚的（相应地，
凝聚的）；若 $\mathscr{F}$ 与 $\mathscr{G}$ 都是有限型的，它还是有限型的。
若 $\mathscr{F}$ 有有限表示，且 $\mathscr{G}$ 是准凝聚的（相应地，凝聚的），则
$\mathscr{H}\!om_{\mathscr{O}_X}(\mathscr{F},\mathscr{G})$ 是准凝聚的
（相应地，凝聚的）。
\end{proposition}

问题是局部的，因此可以假设 $X$ 是仿射的（相应地，是诺特仿射的）；此外，
若 $\mathscr{F}$ 是凝聚的，可以假设它是同态
$\mathscr{O}_X^m\to\mathscr{O}_X^n$ 的余核。关于准凝聚层的断言于是由
（\hyperref[I.1.3.12-fr]{1.3.12}）和
（\hyperref[I.1.3.9-fr]{1.3.9}）推出；关于凝聚层的断言由（\hyperref[I.1.5.1-fr]{1.5.1}）
以及以下事实推出：若 $M$、$N$ 是诺特环 $A$ 上的有限型模，则
$M\otimes_A N$ 和 $\operatorname{Hom}_A(M,N)$ 都是有限型 $A$-模。

\begin{definition}[9.1.2]
\label{I.9.1.2-fr}
设 $X$、$Y$ 是两个 $S$-预概形，$p$、$q$ 是
$X\times_S Y$ 的投影，$\mathscr{F}$（相应地，$\mathscr{G}$）是
$\mathscr{O}_X$-模层（相应地，$\mathscr{O}_Y$-模层），且为准凝聚的。称
$\mathscr{F}$ 与 $\mathscr{G}$ 在 $\mathscr{O}_S$（或在 $S$）上的张量积，并记为
$\mathscr{F}\otimes_{\mathscr{O}_S}\mathscr{G}$（或
$\mathscr{F}\otimes_S\mathscr{G}$），是预概形 $X\times_S Y$ 上的张量积
$p^*(\mathscr{F})\otimes_{\mathscr{O}_{X\times_S Y}}q^*(\mathscr{G})$。
\end{definition}

若 $X_i$（$1\leq i\leq n$）是 $S$-预概形，$\mathscr{F}_i$ 是
$\mathscr{O}_{X_i}$-模层且准凝聚（$1\leq i\leq n$），则同样定义张量积
$\mathscr{F}_1\otimes_S\mathscr{F}_2\otimes_S\cdots\otimes_S\mathscr{F}_n$
在预概形
$Z=X_1\times_S X_2\times_S\cdots\times_S X_n$ 上；根据
（\hyperref[I.9.1.1-fr]{9.1.1}）和（\hyperref[0.5.1.4-fr]{0, 5.1.4}），它是
$\mathscr{O}_Z$-模层\emph{准凝聚}；若各 $\mathscr{F}_i$ 都凝聚且 $Z$ 是
\emph{局部诺特的}，则根据（\hyperref[I.9.1.1-fr]{9.1.1}）、
（\hyperref[0.5.3.11-fr]{0, 5.3.11}）和（\hyperref[I.6.1.1-fr]{6.1.1}），它是\emph{凝聚的}。

注意，若取 $X=Y=S$，定义（\hyperref[I.9.1.2-fr]{9.1.2}）就给出
$\mathscr{O}_S$-模层的张量积。此外，由于
$q^*(\mathscr{O}_Y)=\mathscr{O}_{X\times_S Y}$
（\hyperref[0.4.3.4-fr]{0, 4.3.4}），张量积
$\mathscr{F}\otimes_S\mathscr{O}_Y$ 典范地等同于
$p^*(\mathscr{F})$；同样，
$\mathscr{O}_X\otimes_S\mathscr{G}$ 典范地等同于
$q^*(\mathscr{G})$。特别地，若取 $Y=S$，并以 $f$ 表示结构态射 $X\to Y$，则有
$\mathscr{O}_X\otimes_Y\mathscr{G}=f^*(\mathscr{G})$：普通张量积与逆像
因此都是一般张量积的特殊情形。

定义（\hyperref[I.9.1.2-fr]{9.1.2}）立即推出：对固定的 $X$ 与 $Y$，
$\mathscr{F}\otimes_S\mathscr{G}$ 关于 $\mathscr{F}$ 和 $\mathscr{G}$ 是一个
\emph{加性的、协变的且右正合的双函子}。

\begin{proposition}[9.1.3]
\label{I.9.1.3-fr}
设 $S$、$X$、$Y$ 是三个仿射概形，其环分别为 $A$、$B$、$C$，其中
$B$、$C$ 是 $A$-代数。设 $M$（相应地，$N$）是 $B$-模（相应地，$C$-模），
$\mathscr{F}=\widetilde{M}$（相应地，$\mathscr{G}=\widetilde{N}$）是与之
关联的准凝聚层；则 $\mathscr{F}\otimes_S\mathscr{G}$ 典范同构于与
$(B\otimes_A C)$-模 $M\otimes_A N$ 关联的层。
\end{proposition}

\begin{proof}
\oldpage[I]{170}
事实上，根据（\hyperref[I.1.6.5-fr]{1.6.5}），
$\mathscr{F}\otimes_S\mathscr{G}$ 典范同构于与如下
$(B\otimes_A C)$-模关联的层：
\[
  \bigl(M\otimes_B(B\otimes_A C)\bigr)
  \otimes_{B\otimes_A C}
  \bigl((B\otimes_A C)\otimes_C N\bigr)
\]
由于张量积之间的典范同构，上式又同构于
\[
  M\otimes_B(B\otimes_A C)\otimes_C N
  =(M\otimes_B B)\otimes_A(C\otimes_C N)=M\otimes_A N.
\]
\end{proof}

\begin{proposition}[9.1.4]
\label{I.9.1.4-fr}
设 $f:T\to X$、$g:T\to Y$ 是两个 $S$-态射，$\mathscr{F}$（相应地，
$\mathscr{G}$）是准凝聚的 $\mathscr{O}_X$-模层（相应地，$\mathscr{O}_Y$-模层）。则
\[
  (f,g)_S^*(\mathscr{F}\otimes_S\mathscr{G})
  =f^*(\mathscr{F})\otimes_{\mathscr{O}_T}g^*(\mathscr{G}).
\]
\end{proposition}

\begin{proof}
若 $p$、$q$ 是 $X\times_S Y$ 的投影，则该公式事实上由关系
$(f,g)_S^*\circ p^*=f^*$ 和
$(f,g)_S^*\circ q^*=g^*$（\hyperref[0.3.5.5-fr]{0, 3.5.5}），以及
一个代数层张量积的逆像是其逆像层的张量积这一事实推出
（\hyperref[0.4.3.3-fr]{0, 4.3.3}）。
\end{proof}

\begin{corollary}[9.1.5]
\label{I.9.1.5-fr}
设 $f:X\to X'$、$g:Y\to Y'$ 是两个 $S$-态射，$\mathscr{F}'$（相应地，
$\mathscr{G}'$）是准凝聚的 $\mathscr{O}_{X'}$-模层（相应地，$\mathscr{O}_{Y'}$-模层）。则
\[
  (f\times_S g)^*(\mathscr{F}'\otimes_S\mathscr{G}')
  =f^*(\mathscr{F}')\otimes_S g^*(\mathscr{G}').
\]
\end{corollary}

\begin{proof}
这是由（\hyperref[I.9.1.4-fr]{9.1.4}）以及
$f\times_S g=(f\circ p,g\circ q)_S$（其中 $p$、$q$ 是
$X\times_S Y$ 的投影）推出的。
\end{proof}

\begin{corollary}[9.1.6]
\label{I.9.1.6-fr}
设 $X$、$Y$、$Z$ 是三个 $S$-预概形，$\mathscr{F}$（相应地，
$\mathscr{G}$、$\mathscr{H}$）是准凝聚的 $\mathscr{O}_X$-模层（相应地，
$\mathscr{O}_Y$-模层、$\mathscr{O}_Z$-模层）；层
$\mathscr{F}\otimes_S\mathscr{G}\otimes_S\mathscr{H}$ 是通过
$X\times_S Y\times_S Z$ 到 $(X\times_S Y)\times_S Z$ 的典范同构，
从 $(\mathscr{F}\otimes_S\mathscr{G})\otimes_S\mathscr{H}$ 取逆像所得。
\end{corollary}

\begin{proof}
事实上，若以 $p_1$、$p_2$、$p_3$ 表示 $X\times_S Y\times_S Z$ 的投影，
该同构写成 $(p_1,p_2)_S\times_S p_3$。

同样，$\mathscr{G}\otimes_S\mathscr{F}$ 通过
$X\times_S Y$ 到 $Y\times_S X$ 的典范同构取逆像，得到
$\mathscr{F}\otimes_S\mathscr{G}$。
\end{proof}

\begin{corollary}[9.1.7]
\label{I.9.1.7-fr}
若 $X$ 是一个 $S$-预概形，则每个准凝聚的 $\mathscr{O}_X$-模层
$\mathscr{F}$ 都是通过 $X$ 到 $X\times_S S$ 的典范同构
（\hyperref[I.3.3.3-fr]{3.3.3}），从 $\mathscr{F}\otimes_S\mathscr{O}_S$ 取逆像所得。
\end{corollary}

\begin{proof}
事实上，该同构为 $(1_X,\varphi)_S$，其中 $\varphi$ 是结构态射 $X\to S$；
推论由（\hyperref[I.9.1.4-fr]{9.1.4}）以及
$\varphi^*(\mathscr{O}_S)=\mathscr{O}_X$ 得出。
\end{proof}

\begin{env}[9.1.8]
\label{I.9.1.8-fr}
设 $X$ 是一个 $S$-预概形，$\mathscr{F}$ 是准凝聚的
$\mathscr{O}_X$-模层，$\varphi:S'\to S$ 是一个态射；把
$\mathscr{F}_{(\varphi)}$ 或 $\mathscr{F}_{(S')}$ 记为在
$X\times_S S'=X_{(\varphi)}=X_{(S')}$ 上的准凝聚层
$\mathscr{F}\otimes_S\mathscr{O}_{S'}$；因此
$\mathscr{F}_{(S')}=p^*(\mathscr{F})$，其中 $p$ 是投影
$X_{(S')}\to X$。
\end{env}

\begin{proposition}[9.1.9]
\label{I.9.1.9-fr}
设 $\varphi':S''\to S'$ 是一个态射。对 $S$-预概形 $X$ 上的任意准凝聚
$\mathscr{O}_X$-模层 $\mathscr{F}$，$(\mathscr{F}_{(\varphi)})_{(\varphi')}$ 是通过典范同构
$(X_{(\varphi)})_{(\varphi')}\xrightarrow{\sim}
X_{(\varphi\circ\varphi')}$（\hyperref[I.3.3.9-fr]{3.3.9}），
从 $\mathscr{F}_{(\varphi\circ\varphi')}$ 取逆像所得。
\end{proposition}

\begin{proof}
这立即由定义和（\hyperref[I.3.3.9-fr]{3.3.9}）推出，也可写成
\[
  (\mathscr{F}\otimes_S\mathscr{O}_{S'})
  \otimes_{S'}\mathscr{O}_{S''}
  =\mathscr{F}\otimes_S\mathscr{O}_{S''}.
  \tag{9.1.9.1}\label{I.9.1.9.1-fr}
\]
\end{proof}

\begin{proposition}[9.1.10]
\label{I.9.1.10-fr}
设 $Y$ 是一个 $S$-预概形，$f:X\to Y$ 是一个 $S$-态射。对任意准凝聚
$\mathscr{O}_Y$-模层 $\mathscr{G}$ 以及任意态射 $S'\to S$，有
\[
  (f_{(S')})^*(\mathscr{G}_{(S')})
  =(f^*(\mathscr{G}))_{(S')}.
\]
\end{proposition}

\begin{proof}
这立即由下图的交换性得到：
\oldpage[I]{171}
\[
  \xymatrix{
    X_{(S')}\ar[r]^{f_{(S')}}\ar[d] &
    Y_{(S')}\ar[d]\\
    X\ar[r]_f &
    Y
  }
\]
\end{proof}

\begin{corollary}[9.1.11]
\label{I.9.1.11-fr}
设 $X$、$Y$ 是两个 $S$-预概形，$\mathscr{F}$（相应地，$\mathscr{G}$）是
$\mathscr{O}_X$-模层（相应地，$\mathscr{O}_Y$-模层），且为准凝聚的。
通过典范同构
$(X\times_S Y)_{(S')}\xrightarrow{\sim}
(X_{(S')})\times_{S'}(Y_{(S')})$
（\hyperref[I.3.3.10-fr]{3.3.10}），对层
$(\mathscr{F}_{(S')}\otimes_{S'}\mathscr{G}_{(S')})$ 取逆像，所得层等于
$(\mathscr{F}\otimes_S\mathscr{G})_{(S')}$。
\end{corollary}

\begin{proof}
若 $p$、$q$ 是 $X\times_S Y$ 的投影，则所涉及的同构正是
$(p_{(S')},q_{(S')})_{S'}$；推论由命题
（\hyperref[I.9.1.4-fr]{9.1.4}）和（\hyperref[I.9.1.10-fr]{9.1.10}）推出。
\end{proof}

\begin{proposition}[9.1.12]
\label{I.9.1.12-fr}
沿用（\hyperref[I.9.1.2-fr]{9.1.2}）的记号，设 $z$ 是 $X\times_S Y$ 的一点，
$x=p(z)$、$y=q(z)$；则纤维
$(\mathscr{F}\otimes_S\mathscr{G})_z$ 同构于
$(\mathscr{F}_x\otimes_{\mathscr{O}_x}\mathscr{O}_z)
\otimes_{\mathscr{O}_z}
(\mathscr{G}_y\otimes_{\mathscr{O}_y}\mathscr{O}_z)
=\mathscr{F}_x\otimes_{\mathscr{O}_x}\mathscr{O}_z
\otimes_{\mathscr{O}_y}\mathscr{G}_y$。
\end{proposition}

\begin{proof}
由于可以归约到仿射情形，命题由公式
（\hyperref[I.1.6.5.1-fr]{1.6.5.1}）推出。
\end{proof}

\begin{corollary}[9.1.13]
\label{I.9.1.13-fr}
若 $\mathscr{F}$ 与 $\mathscr{G}$ 是有限型的，则
\[
  \operatorname{Supp}(\mathscr{F}\otimes_S\mathscr{G})
  =p^{-1}(\operatorname{Supp}(\mathscr{F}))
  \cap q^{-1}(\operatorname{Supp}(\mathscr{G})).
\]
\end{corollary}

\begin{proof}
由于 $p^*(\mathscr{F})$ 与 $q^*(\mathscr{G})$ 是
$\mathscr{O}_{X\times_S Y}$ 上的有限型模层，根据
（\hyperref[I.9.1.12-fr]{9.1.12}）和（\hyperref[0.1.7.5-fr]{0, 1.7.5}），
只需考虑 $\mathscr{G}=\mathscr{O}_Y$ 的特殊情形，即证明公式
\[
  \operatorname{Supp}(p^{-1}(\mathscr{F}))
  =p^{-1}(\operatorname{Supp}(\mathscr{F})).
  \tag{9.1.13.1}\label{I.9.1.13.1-fr}
\]

与（\hyperref[0.1.7.5-fr]{0, 1.7.5}）中相同的论证归结为验证：对每个
$z\in X\times_S Y$，有
$\mathscr{O}_z/\mathfrak{m}_x\mathscr{O}_z\neq0$（其中 $x=p(z)$）。这由以下事实推出：
按假设，同态 $\mathscr{O}_x\to\mathscr{O}_z$ 是\emph{局部}的。
\end{proof}

我们把本节中对两个因子证明的结果推广到任意多个因子的工作留给读者。

\subsection{准凝聚层的直像}
\label{subsection:I.9.2-fr}

\begin{proposition}[9.2.1]
\label{I.9.2.1-fr}
设 $f:X\to Y$ 是预概形的一个态射。假设 $Y$ 有一个仿射开集覆盖
$(Y_\alpha)$，满足以下性质：每个 $f^{-1}(Y_\alpha)$ 都有一个有限开仿射覆盖
$(X_{\alpha i})$，这些开集包含于 $f^{-1}(Y_\alpha)$ 中，并且每个交集
$X_{\alpha i}\cap X_{\alpha j}$ 本身都是有限个仿射开集的并。在这些条件下，对任意
准凝聚的 $\mathscr{O}_X$-模层 $\mathscr{F}$，$f_*(\mathscr{F})$ 是
准凝聚的 $\mathscr{O}_Y$-模层。
\end{proposition}

\begin{proof}
问题在 $Y$ 上是局部的，因此可以假设 $Y$ 等于某个 $Y_\alpha$，从而略去指标 $\alpha$。

\begin{enumerate}
\item[a)] 首先假设 $X_i\cap X_j$ 本身是\emph{仿射}开集。置
$\mathscr{F}_i=\mathscr{F}|X_i$，
$\mathscr{F}_{ij}=\mathscr{F}|(X_i\cap X_j)$，并令
$\mathscr{F}'_i$ 与 $\mathscr{F}'_{ij}$ 分别为 $\mathscr{F}_i$ 与
$\mathscr{F}_{ij}$ 通过 $f$ 在 $X_i$ 与 $X_i\cap X_j$ 上的限制所得的直像；已知
$\mathscr{F}'_i$ 与 $\mathscr{F}'_{ij}$ 是准凝聚的
（\hyperref[I.1.6.3-fr]{1.6.3}）。置
$\mathscr{G}=\bigoplus_i\mathscr{F}'_i$，
$\mathscr{H}=\bigoplus_{i,j}\mathscr{F}'_{ij}$；$\mathscr{G}$ 与
$\mathscr{H}$ 是准凝聚的 $\mathscr{O}_Y$-模层；我们将定义一个同态
$u:\mathscr{G}\to\mathscr{H}$，使得 $f_*(\mathscr{F})$ 是 $u$ 的\emph{核}；
于是根据（\hyperref[I.1.3.9-fr]{1.3.9}），$f_*(\mathscr{F})$ 是准凝聚的。只需把 $u$
\oldpage[I]{172}
定义为预层同态；考虑到 $\mathscr{G}$ 与 $\mathscr{H}$ 的定义，对每个开集
$W\subset Y$，只需定义同态
\[
  u_W:\bigoplus_i\Gamma(f^{-1}(W)\cap X_i,\mathscr{F})
  \longrightarrow
  \bigoplus_{i,j}\Gamma(f^{-1}(W)\cap X_i\cap X_j,\mathscr{F})
\]
并使之满足 $W$ 变化时通常的相容条件。若对每个
$s_i\in\Gamma(f^{-1}(W)\cap X_i,\mathscr{F})$，以 $s_{i|j}$ 表示其在
$f^{-1}(W)\cap X_i\cap X_j$ 上的限制，则令
\[
  u_W((s_i))=(s_{i|j}-s_{j|i})
\]
显然满足相容条件。为证明 $u$ 的核 $\mathscr{R}$ 就是 $f_*(\mathscr{F})$，
定义从 $f_*(\mathscr{F})$ 到 $\mathscr{R}$ 的同态：把每个
$s\in\Gamma(f^{-1}(W),\mathscr{F})$ 对应到族 $(s_i)$，其中 $s_i$ 是 $s$ 在
$f^{-1}(W)\cap X_i$ 上的限制；层（G, II, 1.1）的公理 (F 1) 和 (F 2) 表明，
这个同态是\emph{双射}，于是本情形的证明完成。

\item[b)] 一般情形中，只要证明 $\mathscr{F}'_{ij}$ 是准凝聚的，便可应用相同论证。
然而按假设，$X_i\cap X_j$ 是有限个仿射开集 $X_{ijk}$ 的并；又因为
$X_{ijk}$ 是概形中的\emph{仿射}开集，它们中任意两个的交仍是仿射开集
（\hyperref[I.5.5.6-fr]{5.5.6}）。因此可归约到第一种情形，命题（9.2.1）得证。
\end{enumerate}
\end{proof}

\begin{corollary}[9.2.2]
\label{I.9.2.2-fr}
（9.2.1）的结论在以下每一种情形中都成立：
\begin{enumerate}
\item[a)] $f$ 是分离且拟紧的。
\item[b)] $f$ 是分离且有限型的。
\item[c)] $f$ 是拟紧的，且 $X$ 的底空间是局部诺特的。
\end{enumerate}
\end{corollary}

\begin{proof}
在 a) 中，$X_{\alpha i}\cap X_{\alpha j}$ 是仿射的
（\hyperref[I.5.5.6-fr]{5.5.6}）。b) 是 a) 的特殊情形
（\hyperref[I.6.6.3-fr]{6.6.3}）。最后，在 c) 中，可以归约到 $Y$ 仿射且
$X$ 的底空间诺特的情形；此时 $X$ 有一个有限仿射开覆盖 $(X_i)$，而
$X_i\cap X_j$ 是拟紧的，因而是有限个仿射开集的并
（\hyperref[I.2.1.3-fr]{2.1.3}）。
\end{proof}

\subsection{准凝聚层截面的延拓}
\label{subsection:I.9.3-fr}

\begin{theorem}[9.3.1]
\label{I.9.3.1-fr}
设 $X$ 是底空间为诺特空间的预概形，或底空间为拟紧空间的概形。设
$\mathscr{L}$ 是一个可逆 $\mathscr{O}_X$-模层
（\hyperref[0.5.4.1-fr]{0, 5.4.1}），$f$ 是 $X$ 上 $\mathscr{L}$ 的一个截面，
$X_f$ 是满足 $f(x)\neq0$ 的 $x\in X$ 所成的开集
（\hyperref[0.5.5.1-fr]{0, 5.5.1}），$\mathscr{F}$ 是准凝聚的
$\mathscr{O}_X$-模层。
\begin{enumerate}
\item[(i)] 若 $s\in\Gamma(X,\mathscr{F})$ 且 $s|X_f=0$，则存在整数 $n>0$ 使得
$s\otimes f^{\otimes n}=0$。
\item[(ii)] 对任意 $s\in\Gamma(X_f,\mathscr{F})$，存在整数 $n>0$ 使得
$s\otimes f^{\otimes n}$ 可以延拓为 $X$ 上
$\mathscr{F}\otimes\mathscr{L}^{\otimes n}$ 的一个截面。
\end{enumerate}
\end{theorem}

\begin{proof}
\begin{enumerate}
\item[(i)] 由于 $X$ 的底空间是拟紧的，即有限个仿射开集 $U_i$ 的并，且
$\mathscr{L}|U_i$ 同构于 $\mathscr{O}_X|U_i$，可归约到 $X$ 仿射且
$\mathscr{L}=\mathscr{O}_X$ 的情形。此时 $f$ 可视为 $A(X)$ 中的元素，且
$X_f=D(f)$；$s$ 可视为 $A(X)$-模 $M$ 中的元素，而 $s|X_f$ 是 $M_f$ 中的相应元素；
根据分式模的定义，结论显然。
\oldpage[I]{173}

\item[(ii)] 再次，$X$ 是有限个仿射开集 $U_i$（$1\leq i\leq r$）的并，且
$\mathscr{L}|U_i\simeq\mathscr{O}_X|U_i$。对每个 $i$，
$(s\otimes f^{\otimes n})|(U_i\cap X_f)$ 在前述同构下对应于
$(f|(U_i\cap X_f))^n(s|(U_i\cap X_f))$。由（\hyperref[I.1.4.1-fr]{1.4.1}）可知，
存在整数 $n>0$，使得对每个 $i$，
$(s\otimes f^{\otimes n})|(U_i\cap X_f)$ 可延拓为 $U_i$ 上
$\mathscr{F}\otimes\mathscr{L}^{\otimes n}$ 的截面 $s_i$。令 $s_{i|j}$ 是 $s_i$ 在
$U_i\cap U_j$ 上的限制；按定义，在 $X_f\cap U_i\cap U_j$ 中
$s_{i|j}-s_{j|i}=0$。若 $X$ 是诺特空间，则 $U_i\cap U_j$ 拟紧；若 $X$ 是概形，
$U_i\cap U_j$ 是仿射开集（\hyperref[I.5.5.6-fr]{5.5.6}），因而仍拟紧。
由 (i)，存在与 $i$、$j$ 无关的整数 $m$，使得
$(s_{i|j}-s_{j|i})\otimes f^{\otimes m}=0$。于是立即得到一个
$X$ 上 $\mathscr{F}\otimes\mathscr{L}^{\otimes(n+m)}$ 的截面 $s'$，它在每个
$U_i$ 上诱导 $s_i\otimes f^{\otimes m}$，从而在 $X_f$ 上诱导
$s\otimes f^{\otimes(n+m)}$。
\end{enumerate}
\end{proof}

这些推论给定理（9.3.1）一个更代数的解释：

\begin{corollary}[9.3.2]
\label{I.9.3.2-fr}
在（9.3.1）的假设下，考虑分次环
$A_* = \Gamma_*(\mathscr{L})$ 以及分次 $A_*$-模
$M_* = \Gamma_*(\mathscr{L},\mathscr{F})$
（\hyperref[0.5.4.6-fr]{0, 5.4.6}）。若 $f\in A_n$，其中 $n\in\mathbb{Z}$，则有典范同构
\[
  \Gamma(X_f,\mathscr{F})\xrightarrow{\sim}((M_*)_f)_0
\]
（即分式模 $(M_*)_f$ 中由次数为 $0$ 的元素组成的子群）。
\end{corollary}

\begin{corollary}[9.3.3]
\label{I.9.3.3-fr}
假设（9.3.1）的条件成立，并进一步假设
$\mathscr{L}=\mathscr{O}_X$。置 $A=\Gamma(X,\mathscr{O}_X)$、
$M=\Gamma(X,\mathscr{F})$，则 $A_f$-模
$\Gamma(X_f,\mathscr{F})$ 典范同构于 $M_f$。
\end{corollary}

\begin{proposition}[9.3.4]
\label{I.9.3.4-fr}
设 $X$ 是诺特预概形，$\mathscr{F}$ 是一个凝聚的
$\mathscr{O}_X$-模层，$\mathscr{J}$ 是 $\mathscr{O}_X$ 中的凝聚理想层，且
 $\mathscr{F}$ 的支集包含于 $\mathscr{O}_X/\mathscr{J}$ 的支集。则存在整数 $n>0$，使得
$\mathscr{J}^n\mathscr{F}=0$。
\end{proposition}

\begin{proof}
由于 $X$ 是有限个仿射开集的并，且其环是诺特的，可以假设 $X$ 是环为诺特环
$A$ 的仿射概形；于是 $\mathscr{F}=\widetilde M$，其中
$M=\Gamma(X,\mathscr{F})$ 是有限型 $A$-模，且 $\mathscr{J}=\widetilde{\mathfrak{J}}$，其中
$\mathfrak{J}=\Gamma(X,\mathscr{J})$ 是 $A$ 的理想
（\hyperref[I.1.4.1-fr]{1.4.1} 和 \hyperref[I.1.5.1-fr]{1.5.1}）。由于 $A$ 是诺特的，
$\mathfrak{J}$ 有有限生成元组 $f_i$（$1\leq i\leq m$）。按假设，
$\mathscr{F}$ 在 $X$ 上的每个截面在每个 $D(f_i)$ 中都为零；若 $s_j$
（$1\leq j\leq q$）是生成 $M$ 的 $\mathscr{F}$ 截面，则存在与 $i$、$j$ 无关的整数 $h$，使得
$f_i^h s_j=0$（\hyperref[I.1.4.1-fr]{1.4.1}），因而对每个 $s\in M$ 都有 $f_i^h s=0$。
所以若 $n=mh$，则 $\mathfrak{J}^nM=0$，相应的
$\mathscr{O}_X$-模层
$\mathscr{J}^n\mathscr{F}=\widetilde{\mathfrak{J}^nM}$
（\hyperref[I.1.3.13-fr]{1.3.13}）为零。
\end{proof}

\begin{corollary}[9.3.5]
\label{I.9.3.5-fr}
在（9.3.4）的假设下，存在 $X$ 的一个闭子预概形 $Y$，其底空间是
$\mathscr{O}_X/\mathscr{J}$ 的支集，并且若 $j:Y\to X$ 是典范嵌入，则
$\mathscr{F}=j_*(j^*(\mathscr{F}))$。
\end{corollary}

\begin{proof}
先注意到 $\mathscr{O}_X/\mathscr{J}$ 与 $\mathscr{O}_X/\mathscr{J}^n$ 的支集相同：
若 $\mathscr{J}_x=\mathscr{O}_x$，则 $\mathscr{J}_x^n=\mathscr{O}_x$；另一方面，对每个
$x\in X$ 有 $\mathscr{J}_x^n\subset\mathscr{J}_x$。因此由（9.3.4）可假设
$\mathscr{J}\mathscr{F}=0$；取 $Y$ 为由 $\mathscr{J}$ 定义的 $X$ 的闭子预概形，
此时 $\mathscr{F}$ 是一个 $(\mathscr{O}_X/\mathscr{J})$-模层，结论立即得到。
\end{proof}

\oldpage[I]{174}

\subsection{准凝聚层的延拓}
\label{subsection:I.9.4-fr}

\begin{env}[9.4.1]
\label{I.9.4.1-fr}
设 $X$ 是拓扑空间，$\mathscr{F}$ 是 $X$ 上的集合层
（相应地，群层、环层），$U$ 是 $X$ 的开子集，
$\psi:U\to X$ 是典范嵌入，$\mathscr{G}$ 是
$\mathscr{F}|U=\psi^*(\mathscr{F})$ 的子层。由于 $\psi_*$ 左正合，
$\psi_*(\mathscr{G})$ 是 $\psi_*(\psi^*(\mathscr{F}))$ 的子层；考虑典范同态
$\rho:\mathscr{F}\to\psi_*(\psi^*(\mathscr{F}))$
（\hyperref[0.3.5.3-fr]{0, 3.5.3}），把
$\overline{\mathscr{G}}$ 定义为 $\mathscr{F}$ 的子层
$\rho^{-1}(\psi_*(\mathscr{G}))$。由定义立即得到，对 $X$ 的每个开集 $V$，
$\Gamma(V,\overline{\mathscr{G}})$ 由以下截面组成：
$s\in\Gamma(V,\mathscr{F})$，其在 $V\cap U$ 上的限制是
$V\cap U$ 上 $\mathscr{G}$ 的截面。因此
$\overline{\mathscr{G}}|U=\psi^*(\overline{\mathscr{G}})=\mathscr{G}$，且
$\overline{\mathscr{G}}$ 是在 $U$ 上诱导出 $\mathscr{G}$ 的 $\mathscr{F}$ 的\emph{最大}子层；
称 $\overline{\mathscr{G}}$ 为 $\mathscr{F}|U$ 的子层
$\mathscr{G}$ 到 $\mathscr{F}$ 中的\emph{典范延拓}。
\end{env}

\begin{proposition}[9.4.2]
\label{I.9.4.2-fr}
设 $X$ 是预概形，$U$ 是 $X$ 的开子集，且典范嵌入
$j:U\to X$ 是拟紧态射
\emph{（若 $X$ 的底空间\emph{局部诺特}，这对\emph{每个} $U$ 都成立
（\hyperref[I.6.6.4-fr]{6.6.4, (i)}））}。则：
\begin{enumerate}
\item[(i)] 对每个准凝聚的 $(\mathscr{O}_X|U)$-模层
  $\mathscr{G}$，$j_*(\mathscr{G})$ 是准凝聚的 $\mathscr{O}_X$-模层，且
  $j_*(\mathscr{G})|U=j^*(j_*(\mathscr{G}))=\mathscr{G}$。
\item[(ii)] 对每个准凝聚的 $\mathscr{O}_X$-模层 $\mathscr{F}$ 以及
  准凝聚的 $(\mathscr{O}_X|U)$-模层 $\mathscr{G}$，
  $\mathscr{G}$ 的典范延拓 $\overline{\mathscr{G}}$（9.4.1）是
  $\mathscr{F}$ 的准凝聚子 $\mathscr{O}_X$-模层。
\end{enumerate}
\end{proposition}

\begin{proof}
若 $j=(\psi,\theta)$（其中 $\psi$ 是底空间的嵌入 $U\to X$），按定义，对每个
$(\mathscr{O}_X|U)$-模层 $\mathscr{G}$ 有 $j_*(\mathscr{G})=\psi_*(\mathscr{G})$；
此外，由开集上诱导预概形的定义，对每个 $\mathscr{O}_X$-模层 $\mathscr{H}$ 有
$j^*(\mathscr{H})=\psi^*(\mathscr{H})=\mathscr{H}|U$。因此 (i) 是
（\hyperref[I.9.2.2-fr]{9.2.2, a)}) 的特例；同理，
$j_*(j^*(\mathscr{F}))$ 是准凝聚的，而由于
$\overline{\mathscr{G}}$ 是通过同态
$\rho:\mathscr{F}\to j_*(j^*(\mathscr{F}))$ 从 $j_*(\mathscr{G})$ 取逆像所得，
(ii) 由（\hyperref[I.4.1.1-fr]{4.1.1}）推出。
\end{proof}

注意，当开集 $U$ 是\emph{拟紧的}且 $X$ 是一个\emph{概形}时，
$j:U\to X$ 为拟紧态射的假设也成立：事实上，此时 $U$ 是有限个仿射开集
$U_i$ 的并，而对 $X$ 的每个仿射开集 $V$，$V\cap U_i$ 是仿射开集
（\hyperref[I.5.5.6-fr]{5.5.6}），因而拟紧。

\begin{corollary}[9.4.3]
\label{I.9.4.3-fr}
设 $X$ 是预概形，$U$ 是 $X$ 的拟紧开集，且嵌入态射 $j:U\to X$ 是拟紧的。
进一步假设每个准凝聚 $\mathscr{O}_X$-模层都是其有限型准凝聚
子 $\mathscr{O}_X$-模层的归纳极限
\emph{（当 $X$ 是\emph{仿射概形}时满足此条件）}。设 $\mathscr{F}$ 是
准凝聚的 $\mathscr{O}_X$-模层，$\mathscr{G}$ 是 $\mathscr{F}|U$ 的有限型
准凝聚子 $(\mathscr{O}_X|U)$-模层。则存在 $\mathscr{F}$ 的有限型
准凝聚子 $\mathscr{O}_X$-模层 $\mathscr{G}'$，使得
$\mathscr{G}'|U=\mathscr{G}$。
\end{corollary}

\begin{proof}
事实上，$\mathscr{G}=\overline{\mathscr{G}}|U$，而由（9.4.2），
$\overline{\mathscr{G}}$ 是准凝聚的，因而是其有限型准凝聚
子 $\mathscr{O}_X$-模层 $\mathscr{H}_\lambda$ 的归纳极限。因此
$\mathscr{G}$ 是 $\mathscr{H}_\lambda|U$ 的归纳极限；由于它是有限型的，
它等于某个 $\mathscr{H}_\lambda|U$
（\hyperref[0.5.2.3-fr]{0, 5.2.3}）。
\end{proof}

\begin{remark}[9.4.4]
\label{I.9.4.4-fr}
假设对每个仿射开集 $U\subset X$，嵌入态射 $U\to X$ 都是拟紧的。若对每个仿射开集
$U$ 以及 $\mathscr{F}|U$ 的每个有限型准凝聚子 $(\mathscr{O}_X|U)$-模层 $\mathscr{G}$，
（9.4.3）的结论都成立，
\oldpage[I]{175}
则 $\mathscr{F}$ 是其有限型准凝聚子 $\mathscr{O}_X$-模层的归纳极限。事实上，对每个仿射开集
$U\subset X$，有 $\mathscr{F}|U=\widetilde{M}$，其中 $M$ 是一个 $A(U)$-模；
由于后者是其有限型子模的归纳极限，$\mathscr{F}|U$ 是其有限型准凝聚
子 $(\mathscr{O}_X|U)$-模层的归纳极限
（\hyperref[I.1.3.9-fr]{1.3.9}）。按假设，这些子模层中的每一个都由
$\mathscr{F}$ 的一个有限型准凝聚子 $\mathscr{O}_X$-模层
$\mathscr{G}_{\lambda,U}$ 在 $U$ 上诱导。$\mathscr{G}_{\lambda,U}$ 的有限和
仍是有限型准凝聚的 $\mathscr{O}_X$-模层，因为问题是局部的，且 $X$ 仿射的情形已在
（\hyperref[I.1.3.10-fr]{1.3.10}）中处理；显然，$\mathscr{F}$ 是这些有限和的归纳极限，
断言得证。
\end{remark}

\begin{corollary}[9.4.5]
\label{I.9.4.5-fr}
在（9.4.3）的假设下，对每个有限型准凝聚的
$(\mathscr{O}_X|U)$-模层 $\mathscr{G}$，存在有限型准凝聚的
$\mathscr{O}_X$-模层 $\mathscr{G}'$，使得
$\mathscr{G}'|U=\mathscr{G}$。
\end{corollary}

\begin{proof}
由于 $\mathscr{F}=j_*(\mathscr{G})$ 是准凝聚的（9.4.2），且
$\mathscr{F}|U=\mathscr{G}$，只需把（9.4.3）应用于 $\mathscr{F}$。
\end{proof}

\begin{lemma}[9.4.6]
\label{I.9.4.6-fr}
设 $X$ 是预概形，$L$ 是良序集，$(V_\lambda)_{\lambda\in L}$ 是由仿射开集构成的
$X$ 的覆盖，$U$ 是 $X$ 的一个开集；对每个 $\lambda\in L$，置
$W_\lambda=\bigcup_{\mu<\lambda}V_\mu$。假设：
\begin{enumerate}
  \item[1\textsuperscript{o}] 对每个 $\lambda\in L$，$V_\lambda\cap W_\lambda$ 是拟紧的；
  \item[2\textsuperscript{o}] 浸入态射 $U\to X$ 是拟紧的。
\end{enumerate}
则，对每个准凝聚的 $\mathscr{O}_X$-模层 $\mathscr{F}$ 以及
$\mathscr{F}|U$ 的有限型准凝聚子 $(\mathscr{O}_X|U)$-模层 $\mathscr{G}$，存在
$\mathscr{F}$ 的有限型准凝聚子 $\mathscr{O}_X$-模层 $\mathscr{G}'$，使得
$\mathscr{G}'|U=\mathscr{G}$。
\end{lemma}

\begin{proof}
置 $U_\lambda=U\cup W_\lambda$；将用超限归纳定义一族
$(\mathscr{G}'_\lambda)$，其中 $\mathscr{G}'_\lambda$ 是
$\mathscr{F}|U_\lambda$ 的有限型准凝聚子 $(\mathscr{O}_X|U_\lambda)$-模层，满足
$\mathscr{G}'_\lambda|U_\mu=\mathscr{G}'_\mu$（$\mu<\lambda$）以及
$\mathscr{G}'_\lambda|U=\mathscr{G}$。唯一的 $\mathscr{F}$ 的子
$\mathscr{O}_X$-模层 $\mathscr{G}'$ 满足
$\mathscr{G}'|U_\lambda=\mathscr{G}'_\lambda$（对每个 $\lambda\in L$）
（\hyperref[0.3.3.1-fr]{0, 3.3.1}）即为所求。因此假设对 $\mu<\lambda$，
$\mathscr{G}'_\mu$ 已定义并具有上述性质；若 $\lambda$ 没有前驱，则取
$\mathscr{G}'_\lambda$ 为 $\mathscr{F}|U_\lambda$ 的唯一
$(\mathscr{O}_X|U_\lambda)$-子模层，使得对所有 $\mu<\lambda$ 有
$\mathscr{G}'_\lambda|U_\mu=\mathscr{G}'_\mu$；这是合法的，因为此时
$\mu<\lambda$ 的 $U_\mu$ 构成 $U_\lambda$ 的覆盖。若 $\lambda=\mu+1$，则
$U_\lambda=U_\mu\cup V_\mu$，只需定义 $\mathscr{F}|V_\mu$ 的一个有限型
准凝聚子 $(\mathscr{O}_X|V_\mu)$-模层 $\mathscr{G}''_\mu$，满足
\[
  \mathscr{G}''_\mu|(U_\mu\cap V_\mu)
  =\mathscr{G}'_\mu|(U_\mu\cap V_\mu)~;
\]
随后取 $\mathscr{F}|U_\lambda$ 的子 $(\mathscr{O}_X|U_\lambda)$-模层
$\mathscr{G}'_\lambda$，使得
$\mathscr{G}'_\lambda|U_\mu=\mathscr{G}'_\mu$ 且
$\mathscr{G}'_\lambda|V_\mu=\mathscr{G}''_\mu$
（\hyperref[0.3.3.1-fr]{0, 3.3.1}）。由于 $V_\mu$ 是仿射的，
一旦证明 $U_\mu\cap V_\mu$ 拟紧，（9.4.3）便保证 $\mathscr{G}''_\mu$ 存在；
而 $U_\mu\cap V_\mu$ 是 $U\cap V_\mu$ 与 $W_\mu\cap V_\mu$ 的并，
按假设这两者都是拟紧的。
\end{proof}

\begin{theorem}[9.4.7]
\label{I.9.4.7-fr}
设 $X$ 是预概形，$U$ 是 $X$ 的开集。假设以下条件之一成立：
\begin{enumerate}
  \item[(a)] $X$ 的底空间是局部诺特的；
  \item[(b)] $X$ 是拟紧概形且 $U$ 是拟紧开集。
\end{enumerate}
则对每个准凝聚的 $\mathscr{O}_X$-模层 $\mathscr{F}$ 以及 $\mathscr{F}|U$ 的
有限型准凝聚子 $(\mathscr{O}_X|U)$-模层 $\mathscr{G}$，存在 $\mathscr{F}$ 的
有限型准凝聚子 $\mathscr{O}_X$-模层 $\mathscr{G}'$，使得
$\mathscr{G}'|U=\mathscr{G}$。
\end{theorem}

\oldpage[I]{176}

\begin{proof}
设 $(V_\lambda)_{\lambda\in L}$ 是 $X$ 的仿射开覆盖；在情形 b) 中假设 $L$ 有限。
给 $L$ 赋予良序结构，只需验证引理（9.4.6）的条件。情形 a) 中显然成立，因为
$V_\lambda$ 是诺特的。在情形 b) 中，$V_\lambda\cap V_\mu$ 是仿射的
（\hyperref[I.5.5.6-fr]{5.5.6}），因而拟紧；又因为 $L$ 有限，
$V_\lambda\cap W_\lambda$ 是拟紧的。定理得证。
\end{proof}

\begin{corollary}[9.4.8]
\label{I.9.4.8-fr}
在（9.4.7）的假设下，对每个有限型准凝聚的
$(\mathscr{O}_X|U)$-模层 $\mathscr{G}$，存在有限型准凝聚的
$\mathscr{O}_X$-模层 $\mathscr{G}'$，使得
$\mathscr{G}'|U=\mathscr{G}$。
\end{corollary}

\begin{proof}
只需把（9.4.7）应用于 $\mathscr{F}=j_*(\mathscr{G})$；由（9.4.2），它是准凝聚的，
且 $\mathscr{F}|U=\mathscr{G}$。
\end{proof}

\begin{corollary}[9.4.9]
\label{I.9.4.9-fr}
设 $X$ 是底空间局部诺特的预概形，或是拟紧概形。则每个准凝聚的
$\mathscr{O}_X$-模层都是其有限型准凝聚子
$\mathscr{O}_X$-模层的归纳极限。
\end{corollary}

\begin{proof}
这由（9.4.7）和注记（9.4.4）推出。
\end{proof}

\begin{corollary}[9.4.10]
\label{I.9.4.10-fr}
在（9.4.9）的假设下，若一个准凝聚的 $\mathscr{O}_X$-模层
$\mathscr{F}$ 满足：它的每个有限型准凝聚子 $\mathscr{O}_X$-模层都由其在
$X$ 上的截面生成，则 $\mathscr{F}$ 也由其在 $X$ 上的截面生成。
\end{corollary}

\begin{proof}
事实上，设 $U$ 是 $x\in X$ 的一个仿射开邻域，设 $s$ 是 $U$ 上
$\mathscr{F}$ 的截面。由 $s$ 生成的 $\mathscr{F}|U$ 的子 $\mathscr{O}_X$-模层
$\mathscr{G}$ 是准凝聚且有限型，因此存在 $\mathscr{F}$ 的有限型准凝聚子
$\mathscr{O}_X$-模层 $\mathscr{G}'$，使得 $\mathscr{G}'|U=\mathscr{G}$（9.4.7）。
按假设，存在 $\mathscr{G}'$ 在 $X$ 上的有限个截面 $t_i$，以及 $x$ 的某个
$V\subset U$ 邻域上 $\mathscr{O}_X$ 的截面 $a_i$，使得
$s|V=\sum_i a_i\mathbin{\cdot}(t_i|V)$，推论得证。
\end{proof}

\subsection{预概形的闭像与子预概形的闭包}
\label{subsection:I.9.5-fr}

\begin{proposition}[9.5.1]
\label{I.9.5.1-fr}
设 $f:X\to Y$ 是预概形的一个态射，且 $f_*(\mathscr{O}_X)$ 是准凝聚的
$\mathscr{O}_Y$-模层（当 $f$ 拟紧且 $f$ 还分离或 $X$ 局部诺特时成立
（\hyperref[I.9.2.2-fr]{9.2.2}））。则存在 $Y$ 的最小子预概形 $Y'$，使得典范嵌入
$j:Y'\to Y$ 支配 $f$（或者，等价地（\hyperref[I.4.4.1-fr]{4.4.1}），
$X$ 中的子预概形 $f^{-1}(Y')$ 与 $X$\emph{相同}）。
\end{proposition}

更精确地说：

\begin{corollary}[9.5.2]
\label{I.9.5.2-fr}
在（9.5.1）的条件下，设 $f=(\psi,\theta)$，且 $\mathscr{J}$ 是同态
$\theta:\mathscr{O}_Y\to f_*(\mathscr{O}_X)$ 的（准凝聚）核。则由
$\mathscr{J}$ 定义的 $Y$ 的闭子预概形 $Y'$ 满足（9.5.1）的条件。
\end{corollary}

\begin{proof}
由于函子 $\psi^*$ 是正合的，典范分解
$\theta:\mathscr{O}_Y\to\mathscr{O}_Y/\mathscr{J}
\xrightarrow{\theta'}\psi_*(\mathscr{O}_X)$ 给出
（\hyperref[0.3.5.4.3-fr]{0, 3.5.4.3}）分解
\[
  \theta^\sharp:\psi^*(\mathscr{O}_Y)
  \to\psi^*(\mathscr{O}_Y)/\psi^*(\mathscr{J})
  \xrightarrow{{\theta'}^\sharp}\mathscr{O}_X~;
\]
由于对每个 $x\in X$，$\theta_x^\sharp$ 是局部同态，
${\theta'_x}^\sharp$ 也具有同样性质；若以 $\psi_0$ 表示将连续映射 $\psi$ 看作
从 $X$ 到 $X'$ 的映射，以 $\theta_0$ 表示限制
$\theta'|X':(\mathscr{O}_Y/\mathscr{J})|X'
\to\psi_*(\mathscr{O}_X)|X'=(\psi_0)_*(\mathscr{O}_X)$，
则可见 $f_0=(\psi_0,\theta_0)$ 是预概形态射 $X\to X'$
（\hyperref[I.2.2.1-fr]{2.2.1}），且 $f=j\circ f_0$。若现在 $X''$
\oldpage[I]{177}
是 $Y$ 的另一个闭子预概形，由 $\mathscr{O}_Y$ 的准凝聚理想层
$\mathscr{J}'$ 定义，并且嵌入 $j':X''\to Y$ 支配 $f$，则首先必须有
$X''\supset\psi(X)$，所以由于 $X''$ 是闭的，有 $X'\subset X''$。此外，对每个
$y\in X''$，$\theta$ 必须分解为
$\mathscr{O}_y\to\mathscr{O}_y/\mathscr{J}'_y
\to(\psi_*(\mathscr{O}_X))_y$；按定义，这意味着
$\mathscr{J}'_y\subset\mathscr{J}_y$，从而 $X'$ 是 $X''$ 的闭子预概形
（\hyperref[I.4.1.10-fr]{4.1.10}）。
\end{proof}

\begin{definition}[9.5.3]
\label{I.9.5.3-fr}
当存在 $Y$ 的最小闭子预概形 $Y'$，使得典范嵌入 $j:Y'\to Y$ 支配 $f$ 时，称
$Y'$ 是态射 $f$ 下 $X$ 的\emph{闭像预概形}。
\end{definition}

\begin{proposition}[9.5.4]
\label{I.9.5.4-fr}
若 $f_*(\mathscr{O}_X)$ 是准凝聚的 $\mathscr{O}_Y$-模层，则 $f$ 下 $X$ 的闭像的底空间
是 $Y$ 中 $\overline{f(X)}$。
\end{proposition}

\begin{proof}
由于 $f_*(\mathscr{O}_X)$ 的支集包含于 $\overline{f(X)}$，用
（\hyperref[I.9.5.2-fr]{9.5.2}）的记号，对 $y\notin\overline{f(X)}$ 有
$\mathscr{J}_y=\mathscr{O}_y$，所以 $\mathscr{O}_Y/\mathscr{J}$ 的支集包含于
$\overline{f(X)}$。另一方面，该支集是闭的并包含 $f(X)$：事实上，若 $y\in f(X)$，
环 $(\psi_*(\mathscr{O}_X))_y$ 的单位元不为零，因为它是截面
\[
  1\in\Gamma(X,\mathscr{O}_X)=\Gamma(Y,\psi_*(\mathscr{O}_X))~;
\]
在 $y$ 处的芽；由于它是 $\mathscr{O}_y$ 的单位元经 $\theta$ 的像，该单位元不属于
$\mathscr{J}_y$，故 $\mathscr{O}_y/\mathscr{J}_y\neq0$；命题得证。
\end{proof}

\begin{proposition}[9.5.5]
\label{I.9.5.5-fr}
\emph{（闭像的传递性。）} 设 $f:X\to Y$ 和 $g:Y\to Z$ 是两个预概形态射；假设
$f$ 下 $X$ 的闭像 $Y'$ 存在，并且若 $g'$ 是 $g$ 在 $Y'$ 上的限制，$g'$ 下 $Y'$ 的
闭像 $Z'$ 也存在。则 $g\circ f$ 下 $X$ 的闭像存在，且等于 $Z'$。
\end{proposition}

\begin{proof}
根据（\hyperref[I.9.5.1-fr]{9.5.1}），只需证明 $Z'$ 是 $Z$ 的最小闭子预概形 $Z_1$，
使得 $X$ 中的闭子预概形 $(g\circ f)^{-1}(Z_1)$（由
（\hyperref[I.4.4.2-fr]{4.4.2}），等于 $f^{-1}(g^{-1}(Z_1))$）等于 $X$；
这等价于说，$Z'$ 是 $Z$ 的最小闭子预概形，使得 $f$ 被嵌入
$g^{-1}(Z_1)\to Y$ 所支配（\hyperref[I.4.4.1-fr]{4.4.1}）。然而，由于 $Y'$ 的闭像存在，
具有此性质的每个 $Z_1$ 都使 $g^{-1}(Z_1)$ 支配 $Y'$，即
$j^{-1}(g^{-1}(Z_1))=g'^{-1}(Z_1)=Y'$，其中 $j$ 表示嵌入 $Y'\to Y$。
根据 $Z'$ 的定义，$Z'$ 正是满足上述条件的 $Z$ 的最小闭子预概形。
\end{proof}

\begin{corollary}[9.5.6]
\label{I.9.5.6-fr}
设 $f:X\to Y$ 是一个 $S$-态射，且 $Y$ 是 $f$ 下 $X$ 的闭像。设 $Z$ 是一个
$S$-概形；若 $Y$ 到 $Z$ 的两个 $S$-态射 $g_1$、$g_2$ 满足
$g_1\circ f=g_2\circ f$，则 $g_1=g_2$。
\end{corollary}

\begin{proof}
设 $h=(g_1,g_2)_S:Y\to Z\times_S Z$；由于对角线
$T=\Delta_Z(Z)$ 是 $Z\times_S Z$ 的闭子预概形，
$Y'=h^{-1}(T)$ 是 $Y$ 的闭子预概形
（\hyperref[I.4.4.1-fr]{4.4.1}）。置 $u=g_1\circ f=g_2\circ f$；按积的定义，
$h'=h\circ f=(u,u)_S$，故 $h\circ f=\Delta_Z\circ u$。由于
$\Delta_Z^{-1}(T)=Z$，有 $h'^{-1}(T)=u^{-1}(Z)=X$，从而
$f^{-1}(Y')=X$。由（\hyperref[I.4.4.1-fr]{4.4.1}），典范嵌入 $Y'\to Y$ 支配 $f$，
按假设 $Y'=Y$；于是根据（\hyperref[I.4.4.1-fr]{4.4.1}），$h$ 分解为
$\Delta_Z\circ v$，其中 $v$ 是 $Y\to Z$ 的态射，这导致 $g_1=g_2=v$。
\end{proof}

\begin{remark}[9.5.7]
\label{I.9.5.7-fr}
若 $X$ 与 $Y$ 是 $S$-概形，则命题
（\hyperref[I.9.5.6-fr]{9.5.6}）意味着：
\oldpage[I]{178}
当 $Y$ 是 $f$ 下 $X$ 的闭像时，$f$ 是 $S$-概形范畴中的
\emph{满态射}（T, 1.1）。我们将在第五章证明：反之，若 $f$ 下 $X$ 的闭像 $Y'$ 存在且
$f$ 是 $S$-概形的满态射，则必有 $Y'=Y$。
\end{remark}

\begin{proposition}[9.5.8]
\label{I.9.5.8-fr}
假设（\hyperref[I.9.5.1-fr]{9.5.1}）的条件成立，且设 $Y'$ 是 $f$ 下 $X$ 的闭像。
对 $Y$ 的每个开集 $V$，设 $f_V:f^{-1}(V)\to V$ 是 $f$ 的限制；则 $f_V$ 下
$f^{-1}(V)$ 在 $V$ 中的闭像存在，且等于 $Y'$ 在 $Y'$ 的开集 $V\cap Y'$ 上诱导的
预概形（换言之，等于 $Y$ 的子预概形 $\inf(V,Y')$
（\hyperref[I.4.4.3-fr]{4.4.3}））。
\end{proposition}

\begin{proof}
置 $X'=f^{-1}(V)$；由于 $\mathscr{O}_{X'}$ 经 $f_V$ 的直像正是
$f_*(\mathscr{O}_X)$ 在 $V$ 上的限制，故同态
$\mathscr{O}_V\to(f_V)_*(\mathscr{O}_{X'})$ 的核 $\mathscr{J}'$ 正是
$\mathscr{J}$ 在 $V$ 上的限制，命题立即得到。
\end{proof}

注意，这一结果可以解释为：当基预概形作一个\emph{开浸入}
$Y_1\to Y$ 的扩张时，闭像的构造与之交换。我们将在第四章看到，对一个为\emph{平坦}
态射的扩张 $Y_1\to Y$ 也有同样结论，只要 $f$ 是分离且拟紧的。

\begin{proposition}[9.5.9]
\label{I.9.5.9-fr}
设 $f:X\to Y$ 是一个态射，且 $f$ 下 $X$ 的闭像 $Y'$ 存在。
\begin{enumerate}
  \item[(i)] 若 $X$ 是约化的，则 $Y'$ 也是约化的。
  \item[(ii)] 若假设（\hyperref[I.9.5.1-fr]{9.5.1}）的条件成立，且 $X$ 是不可约的
  （相应地，整的），则 $Y'$ 也是不可约的（相应地，整的）。
\end{enumerate}
\end{proposition}

\begin{proof}
按假设，态射 $f$ 分解为
$X\xrightarrow{g}Y'\xrightarrow{j}Y$，其中 $j$ 是典范嵌入。由于 $X$ 是约化的，
$g$ 分解为
$X\xrightarrow{h}Y'_{\mathrm{red}}\xrightarrow{j'}Y'$，其中 $j'$ 是典范嵌入
（\hyperref[I.5.2.2-fr]{5.2.2}）；由 $Y'$ 的定义得到 $Y'_{\mathrm{red}}=Y'$。
若再满足（\hyperref[I.9.5.1-fr]{9.5.1}）的条件，则由
（\hyperref[I.9.5.4-fr]{9.5.4}），$f(X)$ 在 $Y'$ 中稠密；若 $X$ 不可约，
则 $Y'$ 也不可约（\hyperref[0.2.1.5-fr]{0, 2.1.5}）。关于整预概形的断言由前两者合取得到。
\end{proof}

\begin{proposition}[9.5.10]
\label{I.9.5.10-fr}
设 $Y$ 是预概形 $X$ 的一个子预概形，且典范嵌入 $i:Y\to X$ 是拟紧态射。则存在
$X$ 的最小闭子预概形 $\overline{Y}$ 支配 $Y$；其底空间是 $Y$ 的底空间的闭包；后者在
其闭包中是开集，而预概形 $Y$ 正是由 $\overline{Y}$ 在这个开集上诱导的。
\end{proposition}

\begin{proof}
只需把（\hyperref[I.9.5.1-fr]{9.5.1}）应用于嵌入 $j$；按假设它是分离的
（\hyperref[I.5.5.1-fr]{5.5.1}）且拟紧。因此（\hyperref[I.9.5.1-fr]{9.5.1}）证明
$\overline{Y}$ 存在，而（\hyperref[I.9.5.4-fr]{9.5.4}）说明其底空间是 $Y$ 在 $X$ 中的闭包；
由于 $Y$ 在 $X$ 中局部闭，故它在 $\overline{Y}$ 中是开集。最后一个断言由
（\hyperref[I.9.5.8-fr]{9.5.8}）应用于 $X$ 的一个开集 $V$ 得到，使得 $Y$ 在 $V$ 中闭。
\end{proof}

沿用这些记号，若嵌入 $V\to X$ 是拟紧的，且 $\mathscr{J}$ 是
$\mathscr{O}_X|V$ 的准凝聚理想层，定义 $V$ 中闭子预概形 $Y$，则由
（\hyperref[I.9.5.1-fr]{9.5.1}），定义 $\overline{Y}$ 的 $\mathscr{O}_X$ 的准凝聚理想层
是 $\mathscr{J}$ 的典范延拓
（\hyperref[I.9.4.1-fr]{9.4.1}）$\overline{\mathscr{J}}$，因为它显然是
在 $V$ 上诱导 $\mathscr{J}$ 的 $\mathscr{O}_X$ 的最大准凝聚理想子层。
\oldpage[I]{179}
\begin{corollary}[9.5.11]
\label{I.9.5.11-fr}
在（\hyperref[I.9.5.10-fr]{9.5.10}）的假设下，$\mathscr{O}_{\overline{Y}}$ 在
$\overline{Y}$ 的开集 $V$ 上、且在 $V\cap Y$ 中为零的每个截面，都必为零。
\end{corollary}

\begin{proof}
由（\hyperref[I.9.5.8-fr]{9.5.8}），可归约到 $V=\overline{Y}$ 的情形。注意到
$\mathscr{O}_{\overline{Y}}$ 在 $\overline{Y}$ 上的截面典范地对应于
$\overline{Y}\otimes_{\mathbb{Z}}\mathbb{Z}[T]$ 的 $\overline{Y}$-截面
（\hyperref[I.3.3.15-fr]{3.3.15}），且后者在 $\overline{Y}$ 上是分离的，
推论即为（\hyperref[I.9.5.6-fr]{9.5.6}）的特例。
\end{proof}

当存在支配 $X$ 的子预概形 $Y$ 的最小闭子预概形 $\overline{Y}$，其中该子预概形是 $X$ 的子预概形时，在不致混淆的情况下，
称 $\overline{Y}$ 为 $Y$ 在 $X$ 中的\emph{闭包}。

\subsection{准凝聚代数层；结构层的改变}
\label{subsection:I.9.6-fr}

\begin{proposition}[9.6.1]
\label{I.9.6.1-fr}
设 $X$ 是预概形，$\mathscr{B}$ 是准凝聚的 $\mathscr{O}_X$-代数
（\hyperref[0.5.1.3-fr]{0, 5.1.3}）。对于 $\mathscr{B}$-模层 $\mathscr{F}$，要使其在赋环空间
$(X,\mathscr{B})$ 上准凝聚，充要条件是 $\mathscr{F}$ 作为
$\mathscr{O}_X$-模层准凝聚。
\end{proposition}

\begin{proof}
由于问题是局部的，可以假设 $X$ 是环为 $A$ 的仿射概形；于是
$\mathscr{B}=\widetilde{B}$，其中 $B$ 是一个 $A$-代数
（\hyperref[I.1.4.3-fr]{1.4.3}）。若 $\mathscr{F}$ 在赋环空间
$(X,\mathscr{B})$ 上准凝聚，也可以假设 $\mathscr{F}$ 是某个
$\mathscr{B}$-同态 $\mathscr{B}^{(I)}\to\mathscr{B}^{(J)}$ 的余核；由于该同态
同时是 $\mathscr{O}_X$-模层之间的 $\mathscr{O}_X$-同态，且
$\mathscr{B}^{(I)}$ 与 $\mathscr{B}^{(J)}$ 是准凝聚的
$\mathscr{O}_X$-模层（\hyperref[I.1.3.9-fr]{1.3.9, (ii)}），所以
$\mathscr{F}$ 也是准凝聚的 $\mathscr{O}_X$-模层
（\hyperref[I.1.3.9-fr]{1.3.9, (i)}）。

反之，若 $\mathscr{F}$ 是准凝聚的 $\mathscr{O}_X$-模层，则
$\mathscr{F}=\widetilde{M}$，其中 $M$ 是一个 $B$-模
（\hyperref[I.1.4.3-fr]{1.4.3}）；$M$ 同构于某个
$B$-同态 $B^{(I)}\to B^{(J)}$ 的余核，因此 $\mathscr{F}$ 是一个
$\mathscr{B}$-模层，同构于相应同态
$\mathscr{B}^{(I)}\to\mathscr{B}^{(J)}$ 的余核
（\hyperref[I.1.3.13-fr]{1.3.13}），证明完成。
\end{proof}

特别地，若 $\mathscr{F}$ 与 $\mathscr{G}$ 是两个准凝聚的
$\mathscr{B}$-模层，则 $\mathscr{F}\otimes_{\mathscr{B}}\mathscr{G}$ 是
准凝聚的 $\mathscr{B}$-模层；若进一步假设 $\mathscr{F}$ 有有限表示，
$\mathscr{H}\!om_{\mathscr{B}}(\mathscr{F},\mathscr{G})$ 也同样如此
（\hyperref[I.1.3.13-fr]{1.3.13}）。

\begin{env}[9.6.2]
\label{I.9.6.2-fr}
给定预概形 $X$，称准凝聚的 $\mathscr{O}_X$-代数 $\mathscr{B}$ 是
\emph{有限型的}，如果对每个 $x\in X$，存在 $x$ 的一个\emph{仿射}开邻域
$U$，使得 $\Gamma(U,\mathscr{B})=B$ 是 $\Gamma(U,\mathscr{O}_X)=A$ 上的有限型代数。
此时 $\mathscr{B}|U=\widetilde{B}$，且对每个 $f\in A$，诱导的
$(\mathscr{O}_X|D(f))$-代数 $\mathscr{B}|D(f)$ 是有限型的，因为它同构于
$\widetilde{B_f}$，而 $B_f=B\otimes_A A_f$ 显然是 $A_f$ 上的有限型代数。
由于 $D(f)$ 构成 $U$ 的拓扑基，可得：若 $\mathscr{B}$ 是有限型的准凝聚
$\mathscr{O}_X$-代数，则对 $X$ 的每个开集 $V$，$\mathscr{B}|V$ 是
有限型的准凝聚 $(\mathscr{O}_X|V)$-代数。
\end{env}

\begin{proposition}[9.6.3]
\label{I.9.6.3-fr}
设 $X$ 是局部诺特预概形。则每个有限型准凝聚的
$\mathscr{O}_X$-代数 $\mathscr{B}$ 都是凝聚环层
（\hyperref[0.5.3.7-fr]{0, 5.3.7}）。
\end{proposition}

\begin{proof}
再次可限制到 $X$ 是环为诺特环 $A$ 的仿射概形，且
$\mathscr{B}=\widetilde{B}$，其中 $B$ 是 $A$ 上的有限型代数；因此 $B$ 是诺特环。
需要证明的是，某个 $\mathscr{B}$-同态 $\mathscr{B}^m\to\mathscr{B}$ 的核
$\mathscr{N}$ 是有限型的 $\mathscr{B}$-
\oldpage[I]{180}
模层；然而作为 $\mathscr{B}$-模层，它同构于
$\widetilde{N}$，其中 $N$ 是相应的 $B$-模同态 $B^m\to B$ 的核
（\hyperref[I.1.3.13-fr]{1.3.13}）。由于 $B$ 是诺特的，$B^m$ 的子模 $N$ 是有限型 $B$-模，
因此存在一个 $B^p\to B^m$ 的同态，其像为 $N$；序列
$B^p\to B^m\to B$ 正合，所以相应的序列
$\mathscr{B}^p\to\mathscr{B}^m\to\mathscr{B}$ 也正合
（\hyperref[I.1.3.5-fr]{1.3.5}），且由于 $\mathscr{N}$ 是
$\mathscr{B}^p\to\mathscr{B}^m$ 的像
（\hyperref[I.1.3.9-fr]{1.3.9, (i)}），命题得证。
\end{proof}

\begin{corollary}[9.6.4]
\label{I.9.6.4-fr}
在（\hyperref[I.9.6.3-fr]{9.6.3}）的假设下，$\mathscr{B}$-模层
$\mathscr{F}$ 是凝聚的充要条件是：它是准凝聚的 $\mathscr{O}_X$-模层，且是有限型
$\mathscr{B}$-模层。若满足此条件，且 $\mathscr{G}$ 是 $\mathscr{F}$ 的一个
$\mathscr{B}$-子模层或商 $\mathscr{B}$-模层，则 $\mathscr{G}$ 是凝聚的
$\mathscr{B}$-模层的充要条件是：$\mathscr{G}$ 是准凝聚的 $\mathscr{O}_X$-模层。
\end{corollary}

\begin{proof}
根据（\hyperref[I.9.6.1-fr]{9.6.1}），关于 $\mathscr{F}$ 的条件显然是必要的；证明其充分性。
可限制到 $X$ 是环为诺特环 $A$ 的仿射概形，
$\mathscr{B}=\widetilde{B}$，其中 $B$ 是有限型 $A$-代数，
$\mathscr{F}=\widetilde{M}$，其中 $M$ 是 $B$-模，并存在满的
$\mathscr{B}$-同态 $\mathscr{B}^m\to\mathscr{F}\to 0$。相应地有正合序列
$B^m\to M\to 0$，故 $M$ 是有限型 $B$-模；此外，$B^m\to M$ 的核 $P$ 是有限型
$B$-模，因为 $B$ 是诺特的。由（\hyperref[I.1.3.13-fr]{1.3.13}），$\mathscr{F}$ 是某个
$\mathscr{B}^n\to\mathscr{B}^m$ 的 $\mathscr{B}$-同态的余核；由于
$\mathscr{B}$ 是凝聚环层（\hyperref[0.5.3.4-fr]{0, 5.3.4}），故 $\mathscr{F}$ 凝聚。
同样的论证表明，$\mathscr{F}$ 的准凝聚 $\mathscr{B}$-子模层（相应地，准凝聚
$\mathscr{B}$-商模层）是有限型的，第二个断言得证。
\end{proof}

\begin{proposition}[9.6.5]
\label{I.9.6.5-fr}
设 $X$ 是拟紧概形，或是底空间诺特的预概形。对每个有限型准凝聚的
$\mathscr{O}_X$-代数 $\mathscr{B}$，存在 $\mathscr{B}$ 的有限型准凝聚
$\mathscr{O}_X$-子模层 $\mathscr{E}$，使得 $\mathscr{E}$ 生成
（\hyperref[0.4.1.4-fr]{0, 4.1.4}）$\mathscr{O}_X$-代数 $\mathscr{B}$。
\end{proposition}

\begin{proof}
按假设，$X$ 有一个有限仿射开覆盖 $(U_i)$，使得
$\Gamma(U_i,\mathscr{B})=B_i$ 是 $\Gamma(U_i,\mathscr{O}_X)=A_i$ 上的有限型代数；
令 $E_i$ 是 $B_i$ 的有限型 $A_i$-子模，生成 $A_i$-代数 $B_i$；由
（\hyperref[I.9.4.7-fr]{9.4.7}），存在 $\mathscr{B}$ 的准凝聚有限型
$\mathscr{O}_X$-子模层 $\mathscr{E}_i$，使得
$\mathscr{E}_i|U_i=\widetilde{E_i}$。显然，取 $\mathscr{E}_i$ 的和
$\mathscr{E}$ 即得所求。
\end{proof}

\begin{proposition}[9.6.6]
\label{I.9.6.6-fr}
设 $X$ 是底空间局部诺特的预概形，或是拟紧概形。则每个准凝聚的
$\mathscr{O}_X$-代数 $\mathscr{B}$ 都是其有限型准凝聚
$\mathscr{O}_X$-子代数的归纳极限。
\end{proposition}

\begin{proof}
事实上，由（\hyperref[I.9.4.9-fr]{9.4.9}），$\mathscr{B}$ 作为
$\mathscr{O}_X$-模层，是其有限型准凝聚
$\mathscr{O}_X$-子模层的归纳极限；这些子模层生成
$\mathscr{B}$ 的有限型准凝聚 $\mathscr{O}_X$-子代数
（\hyperref[I.1.3.14-fr]{1.3.14}），故 $\mathscr{B}$\emph{更不用说}是归纳极限。
\end{proof}
\section{形式概形}
\label{section:I.10-fr}

\subsection{仿射形式概形}
\label{subsection:I.10.1-fr}

\begin{env}[10.1.1]
\label{I.10.1.1-fr}
设 $A$ 是一个\emph{可容}拓扑环
（\hyperref[0.7.1.2-fr]{0, 7.1.2}）；对于 $A$ 的任一定义理想
$\mathfrak{J}$，$\operatorname{Spec}(A/\mathfrak{J})$ 与
$\operatorname{Spec}(A)$ 的闭子空间 $V(\mathfrak{J})$
等同（\hyperref[I.1.1.11-fr]{1.1.11}），后者是 $A$ 的
\emph{开}素理想的集合；这一拓扑空间不依赖于所取的
\oldpage[I]{181}
定义理想 $\mathfrak{J}$；将其记为 $\mathfrak{X}$。设
$(\mathfrak{J}_\lambda)$ 是 $A$ 中由定义理想组成的一个 0 的基本邻域系；
对每个 $\lambda$，设 $\mathscr{O}_\lambda$ 是
$\operatorname{Spec}(A/\mathfrak{J}_\lambda)$ 的结构层；该层由
$\widetilde{A}/\widetilde{\mathfrak{J}}_\lambda$ 在 $\mathfrak{X}$ 上诱导
（后者在 $\mathfrak{X}$ 外为零）。

若 $\mathfrak{J}_\mu\subset\mathfrak{J}_\lambda$，则典范同态
$A/\mathfrak{J}_\mu\to A/\mathfrak{J}_\lambda$ 定义了一个环层同态
$u_{\lambda\mu}:\mathscr{O}_\mu\to\mathscr{O}_\lambda$
（\hyperref[I.1.6.1-fr]{1.6.1}），而 $(\mathscr{O}_\lambda)$
关于这些同态构成一个\emph{环层的逆系统}。由于 $\mathfrak{X}$ 的拓扑
具有一个由拟紧开集组成的拓扑基，可以给每个 $\mathscr{O}_\lambda$
配上一个\emph{伪离散拓扑环层}（\hyperref[0.3.8.1-fr]{0, 3.8.1}），
其底层的不带拓扑的环层就是 $\mathscr{O}_\lambda$，并仍将其记作
$\mathscr{O}_\lambda$；这些 $\mathscr{O}_\lambda$ 仍构成一个
\emph{拓扑环层的逆系统}（\hyperref[0.3.8.2-fr]{0, 3.8.2}）。记
$\mathscr{O}_{\mathfrak{X}}$ 为 $\mathfrak{X}$ 上的\emph{拓扑环层}，
即系统 $(\mathscr{O}_\lambda)$ 的逆极限；对 $\mathfrak{X}$ 的任一
\emph{拟紧}开集 $U$，$\Gamma(U,\mathscr{O}_{\mathfrak{X}})$
因而是\emph{离散}环系统 $\Gamma(U,\mathscr{O}_\lambda)$ 的逆极限拓扑环
（\hyperref[0.3.2.6-fr]{0, 3.2.6}）。
\end{env}



\begin{definition}[10.1.2]
\label{I.10.1.2-fr}
给定一个可容拓扑环 $A$，称 $\operatorname{Spec}(A)$ 中由 $A$ 的开素理想
组成的闭子空间 $\mathfrak{X}$ 为 $A$ 的形式谱，并记作
$\operatorname{Spf}(A)$。如果一个拓扑赋环空间同构于形式谱
$\operatorname{Spf}(A)=\mathfrak{X}$，并且后者配有拓扑环层
$\mathscr{O}_{\mathfrak{X}}$，此层是伪离散环层
$(\widetilde{A}/\widetilde{\mathfrak{J}}_\lambda)|\mathfrak{X}$ 的逆极限，
其中 $\mathfrak{J}_\lambda$ 遍历 $A$ 的定义理想所成的滤过集，
则称该拓扑赋环空间为仿射形式概形。
\end{definition}

当我们把一个\emph{形式谱} $\mathfrak{X}=\operatorname{Spf}(A)$
称作仿射形式概形时，始终是指拓扑赋环空间
$(\mathfrak{X},\mathscr{O}_{\mathfrak{X}})$，其中
$\mathscr{O}_{\mathfrak{X}}$ 如上定义。

注意，每个\emph{仿射概形} $X=\operatorname{Spec}(A)$ 都可以唯一地看成
仿射形式概形：只需把 $A$ 看成离散拓扑环。此时，当 $U$ 拟紧时，拓扑环
$\Gamma(U,\mathscr{O}_X)$ 是离散的（但当 $U$ 是 $X$ 的任意开集时，
一般并非如此）。

\begin{proposition}[10.1.3]
\label{I.10.1.3-fr}
若 $\mathfrak{X}=\operatorname{Spf}(A)$，其中 $A$ 是可容环，则
$\Gamma(\mathfrak{X},\mathscr{O}_{\mathfrak{X}})$ 与 $A$ 拓扑同构。
\end{proposition}

\begin{proof}
事实上，由于 $\mathfrak{X}$ 在 $\operatorname{Spec}(A)$ 中是闭的，
它是拟紧的，因而 $\Gamma(\mathfrak{X},\mathscr{O}_{\mathfrak{X}})$
与离散环 $\Gamma(\mathfrak{X},\mathscr{O}_\lambda)$ 的逆极限拓扑同构；
而 $\Gamma(\mathfrak{X},\mathscr{O}_\lambda)$ 与
$A/\mathfrak{J}_\lambda$ 同构（\hyperref[I.1.3.7-fr]{1.3.7}）。
由于 $A$ 是分离且完备的，它与
$\varprojlim A/\mathfrak{J}_\lambda$ 拓扑同构
（\hyperref[0.7.2.1-fr]{0, 7.2.1}），命题得证。
\end{proof}

\begin{proposition}[10.1.4]
\label{I.10.1.4-fr}
设 $A$ 是可容环，$\mathfrak{X}=\operatorname{Spf}(A)$，且对每个
$f\in A$，令 $\mathfrak{D}(f)=D(f)\cap\mathfrak{X}$；则拓扑赋环空间
$(\mathfrak{D}(f),\mathscr{O}_{\mathfrak{X}}|\mathfrak{D}(f))$
同构于仿射形式谱 $\operatorname{Spf}(A_{\{f\}})$
（\hyperref[0.7.6.15-fr]{0, 7.6.15}）。
\end{proposition}

\begin{proof}
对 $A$ 的任一定义理想 $\mathfrak{J}$，离散环
$S_f^{-1}A/S_f^{-1}\mathfrak{J}$ 典范地等同于
$A_{\{f\}}/\mathfrak{J}_{\{f\}}$
（\hyperref[0.7.6.9-fr]{0, 7.6.9}），所以由
\hyperref[I.1.2.5-fr]{1.2.5} 和 \hyperref[I.1.2.6-fr]{1.2.6}，
拓扑空间 $\operatorname{Spf}(A_{\{f\}})$ 典范地等同于
$\mathfrak{D}(f)$。此外，对 $\mathfrak{X}$ 中包含于
$\mathfrak{D}(f)$ 的任一拟紧开集 $U$，
$\Gamma(U,\mathscr{O}_\lambda)$ 与
$\operatorname{Spec}(S_f^{-1}A/S_f^{-1}\mathfrak{J}_\lambda)$
的结构层在 $U$ 上的截面模等同（\hyperref[I.1.3.6-fr]{1.3.6}）。
因此，若令 $\mathfrak{Y}=\operatorname{Spf}(A_{\{f\}})$，则
$\Gamma(U,\mathscr{O}_{\mathfrak{X}})$ 与截面模
$\Gamma(U,\mathscr{O}_{\mathfrak{Y}})$ 等同，命题得证。
\end{proof}

\oldpage[I]{182}
\begin{env}[10.1.5]
\label{I.10.1.5-fr}
把 $\operatorname{Spf}(A)$ 的结构层 $\mathscr{O}_{\mathfrak{X}}$
看作\emph{不带拓扑}的环层时，对每个 $x\in\mathfrak{X}$，其茎由
\hyperref[I.10.1.4-fr]{10.1.4} 可等同于归纳极限
$\varinjlim_{f\notin\mathfrak{j}_x}A_{\{f\}}$。于是由
\hyperref[0.7.6.17-fr]{0, 7.6.17} 和
\hyperref[0.7.6.18-fr]{7.6.18} 可得：
\end{env}

\begin{proposition}[10.1.6]
\label{I.10.1.6-fr}
对每个 $x\in\mathfrak{X}=\operatorname{Spf}(A)$，茎
$\mathscr{O}_x$ 是局部环，其剩余域同构于
$k(x)=A_x/\mathfrak{j}_xA_x$。若进一步假设 $A$ 是 Noether 进制环，
则 $\mathscr{O}_x$ 是 Noether 环。
\end{proposition}

由于 $k(x)$ 不是零环，由此可知环层 $\mathscr{O}_{\mathfrak{X}}$
的\emph{支集}正好\emph{等于 $\mathfrak{X}$}。

\subsection{仿射形式概形的态射}
\label{subsection:I.10.2-fr}

\begin{env}[10.2.1]
\label{I.10.2.1-fr}
设 $A$、$B$ 是两个可容环，且
$\varphi:B\to A$ 是一个\emph{连续}同态。于是连续映射
${}^a\varphi:\operatorname{Spec}(A)\to\operatorname{Spec}(B)$
（\hyperref[I.1.2.1-fr]{1.2.1}）把
$\mathfrak{X}=\operatorname{Spf}(A)$ 映入
$\mathfrak{Y}=\operatorname{Spf}(B)$，因为 $A$ 的开素理想在
$\varphi$ 下的逆像是 $B$ 的开素理想。另一方面，对每个 $g\in B$，
$\varphi$ 由 \hyperref[I.10.1.4-fr]{10.1.4}、
\hyperref[I.10.1.3-fr]{10.1.3} 和
\hyperref[0.7.6.7-fr]{0, 7.6.7} 定义一个连续同态
\[
  \Gamma(\mathfrak{D}(g),\mathscr{O}_{\mathfrak{Y}})
  \longrightarrow
  \Gamma(\mathfrak{D}(\varphi(g)),\mathscr{O}_{\mathfrak{X}})
\]
这些同态满足把 $g$ 换成 $g$ 的倍数时相应限制映射的相容条件；又因为
$\mathfrak{D}(\varphi(g))={}^a\varphi^{-1}(\mathfrak{D}(g))$，
它们定义了拓扑环层的一个\emph{连续}同态
$\mathscr{O}_{\mathfrak{Y}}\to
{}^a\varphi_*(\mathscr{O}_{\mathfrak{X}})$
（\hyperref[0.3.2.5-fr]{0, 3.2.5}），仍将其记作
$\widetilde{\varphi}$；这样便得到拓扑赋环空间的态射
$\Phi=({}^a\varphi,\widetilde{\varphi})$，即
$\mathfrak{X}\to\mathfrak{Y}$。注意，把 $\widetilde{\varphi}$ 看成
不带拓扑的环层同态时，对每个 $x\in\mathfrak{X}$，它在茎上定义同态
$\widetilde{\varphi}_x^\sharp:
\mathscr{O}_{{}^a\varphi(x)}\to\mathscr{O}_x$。
\end{env}

\begin{proposition}[10.2.2]
\label{I.10.2.2-fr}
设 $A$、$B$ 是两个可容拓扑环，且
$\mathfrak{X}=\operatorname{Spf}(A)$、
$\mathfrak{Y}=\operatorname{Spf}(B)$。拓扑赋环空间的态射
$u=(\psi,\theta):\mathfrak{X}\to\mathfrak{Y}$ 具有形式
$({}^a\varphi,\widetilde{\varphi})$，其中 $\varphi$ 是连续环同态
$B\to A$，当且仅当对每个 $x\in\mathfrak{X}$，
$\theta_x^\sharp$ 都是局部同态
$\mathscr{O}_{\psi(x)}\to\mathscr{O}_x$。
\end{proposition}

\begin{proof}
条件是必要的。事实上，令
$\mathfrak{p}=\mathfrak{j}_x\in\operatorname{Spf}(A)$，并令
$\mathfrak{q}=\varphi^{-1}(\mathfrak{j}_x)$；若
$g\notin\mathfrak{q}$，则 $\varphi(g)\notin\mathfrak{p}$，而且显然，
由 $\varphi$ 诱导的同态
$B_{\{g\}}\to A_{\{\varphi(g)\}}$
（\hyperref[0.7.6.7-fr]{0, 7.6.7}）把
$\mathfrak{q}_{\{g\}}$ 映入
$\mathfrak{p}_{\{\varphi(g)\}}$。取归纳极限，并考虑
\hyperref[I.10.1.5-fr]{10.1.5} 和
\hyperref[0.7.6.17-fr]{0, 7.6.17}，便知
$\widetilde{\varphi}_x^\sharp$ 是局部同态。

反过来，设 $(\psi,\theta)$ 是满足命题条件的态射。由
\hyperref[I.10.1.3-fr]{10.1.3}，$\theta$ 定义连续环同态
\[
  \varphi=\Gamma(\theta):
  B=\Gamma(\mathfrak{Y},\mathscr{O}_{\mathfrak{Y}})
  \longrightarrow
  \Gamma(\mathfrak{X},\mathscr{O}_{\mathfrak{X}})=A.
\]
由关于 $\theta$ 的假设，$\mathscr{O}_{\mathfrak{X}}$ 在
$\mathfrak{X}$ 上的截面 $\varphi(g)$ 在点 $x$ 处的芽可逆，当且仅当
$g$ 在点 $\psi(x)$ 处的芽可逆。但由
\hyperref[0.7.6.17-fr]{0, 7.6.17}，
$\mathscr{O}_{\mathfrak{X}}$（相应地，$\mathscr{O}_{\mathfrak{Y}}$）
在 $\mathfrak{X}$（相应地，$\mathfrak{Y}$）上的截面中，在点 $x$
（相应地，$\psi(x)$）处的芽不可逆者，恰好是
$\mathfrak{j}_x$
\oldpage[I]{183}
（相应地，$\mathfrak{j}_{\psi(x)}$）中的元素。因此上述说明给出
${}^a\varphi=\psi$。最后，对每个 $g\in B$，图表
\[
  \xymatrix{
    B=\Gamma(\mathfrak{Y},\mathscr{O}_{\mathfrak{Y}})
      \ar[r]^\varphi\ar[d] &
    \Gamma(\mathfrak{X},\mathscr{O}_{\mathfrak{X}})=A\ar[d]\\
    B_{\{g\}}=\Gamma(\mathfrak{D}(g),\mathscr{O}_{\mathfrak{Y}})
      \ar[r]_{\Gamma(\theta_{\mathfrak{D}(g)})} &
    \Gamma(\mathfrak{D}(\varphi(g)),\mathscr{O}_{\mathfrak{X}})
      =A_{\{\varphi(g)\}}
  }
\]
交换。由完备分式环的泛性质（\hyperref[0.7.6.6-fr]{0, 7.6.6}），
对每个 $g\in B$，$\theta_{\mathfrak{D}(g)}$ 等于
$\widetilde{\varphi}_{\mathfrak{D}(g)}$；于是由
\hyperref[0.3.2.5-fr]{0, 3.2.5}，有
$\theta=\widetilde{\varphi}$。
\end{proof}

称满足 \hyperref[I.10.2.2-fr]{10.2.2} 条件的拓扑赋环空间态射
$(\psi,\theta)$ 为\emph{仿射形式概形的态射}。可以说，分别把
$A$ 送到 $\operatorname{Spf}(A)$、把 $\mathfrak{X}$ 送到
$\Gamma(\mathfrak{X},\mathscr{O}_{\mathfrak{X}})$ 的函子，定义了
可容环范畴与仿射形式概形范畴的对偶范畴之间的一个\emph{等价}
（T, I, 1.2）。

\begin{env}[10.2.3]
\label{I.10.2.3-fr}
作为 \hyperref[I.10.2.2-fr]{10.2.2} 的一个特例，注意，对
$f\in A$，由 $\mathfrak{X}$ 在 $\mathfrak{D}(f)$ 上诱导的仿射形式概形
到母空间的典范嵌入，对应于典范连续同态
$A\to A_{\{f\}}$。在 \hyperref[I.10.2.2-fr]{10.2.2} 的假设下，
设 $h$ 是 $B$ 的元素，$g$ 是 $A$ 中 $\varphi(h)$ 的倍数；则
$\psi(\mathfrak{D}(g))\subset\mathfrak{D}(h)$。把 $u$ 在
$\mathfrak{D}(g)$ 上的限制看成从 $\mathfrak{D}(g)$ 到
$\mathfrak{D}(h)$ 的态射，它是使图表
\[
  \xymatrix{
    \mathfrak{D}(g)\ar[r]^v\ar[d] &
    \mathfrak{D}(h)\ar[d]\\
    \mathfrak{X}\ar[r]_u &
    \mathfrak{Y}
  }
\]
交换的唯一态射 $v$。该态射对应于使图表
\[
  \xymatrix{
    A\ar[d] &
    B\ar[l]_\varphi\ar[d]\\
    A_{\{g\}} &
    B_{\{h\}}\ar[l]_{\varphi'}
  }
\]
交换的唯一连续同态
$\varphi':B_{\{h\}}\to A_{\{g\}}$
（\hyperref[0.7.6.7-fr]{0, 7.6.7}）。
\end{env}

\subsection{仿射形式概形的定义理想层}
\label{subsection:I.10.3-fr}

\begin{env}[10.3.1]
\label{I.10.3.1-fr}
设 $A$ 是可容环，$\mathfrak{J}$ 是 $A$ 的开理想，$\mathfrak{X}$ 是
仿射形式概形 $\operatorname{Spf}(A)$。令 $(\mathfrak{J}_\lambda)$
为 $A$ 中包含于 $\mathfrak{J}$ 的所有定义理想所成的集合；于是
$\widetilde{\mathfrak{J}}/\widetilde{\mathfrak{J}}_\lambda$ 是
$\widetilde{A}/\widetilde{\mathfrak{J}}_\lambda$ 的理想层。记
$\mathfrak{J}^\Delta$ 为这些
$\widetilde{\mathfrak{J}}/\widetilde{\mathfrak{J}}_\lambda$ 在
$\mathfrak{X}$ 上诱导的层的逆极限；它可等同于
$\mathscr{O}_{\mathfrak{X}}$ 的一个\emph{理想层}
（\hyperref[0.3.2.6-fr]{0, 3.2.6}）。对每个 $f\in A$，
$\Gamma(\mathfrak{D}(f),\mathfrak{J}^\Delta)$ 是
$S_f^{-1}\mathfrak{J}/S_f^{-1}\mathfrak{J}_\lambda$ 的逆极限；换言之，
它等同于环 $A_{\{f\}}$ 的开理想 $\mathfrak{J}_{\{f\}}$
（\hyperref[0.7.6.9-fr]{0, 7.6.9}），特别地，
$\Gamma(\mathfrak{X},\mathfrak{J}^\Delta)=\mathfrak{J}$。由于
$\mathfrak{D}(f)$ 构成 $\mathfrak{X}$ 的拓扑基，从而有
\[
  \mathfrak{J}^\Delta|\mathfrak{D}(f)
  =(\mathfrak{J}_{\{f\}})^\Delta.
  \tag{10.3.1.1}\label{I.10.3.1.1-fr}
\]
\end{env}

\oldpage[I]{184}
\begin{env}[10.3.2]
\label{I.10.3.2-fr}
沿用 \hyperref[I.10.3.1-fr]{10.3.1} 的记号。对每个 $f\in A$，
从 $A_{\{f\}}=\Gamma(\mathfrak{D}(f),\mathscr{O}_{\mathfrak{X}})$ 到
$\Gamma\bigl(\mathfrak{D}(f),
(\widetilde{A}/\widetilde{\mathfrak{J}})|\mathfrak{X}\bigr)
=S_f^{-1}A/S_f^{-1}\mathfrak{J}$ 的典范映射是\emph{满射}，其核为
$\Gamma(\mathfrak{D}(f),\mathfrak{J}^\Delta)=\mathfrak{J}_{\{f\}}$
（\hyperref[0.7.6.9-fr]{0, 7.6.9}）。因此这些映射定义一个称为
\emph{典范的}\emph{满}连续同态，从拓扑环层
$\mathscr{O}_{\mathfrak{X}}$ 到离散环层
$(\widetilde{A}/\widetilde{\mathfrak{J}})|\mathfrak{X}$，其核为
$\mathfrak{J}^\Delta$。此外，该同态正是
$\widetilde{\varphi}$（\hyperref[I.10.2.1-fr]{10.2.1}），其中
$\varphi$ 是连续同态 $A\to A/\mathfrak{J}$；仿射形式概形态射
$({}^a\varphi,\widetilde{\varphi}):
\operatorname{Spec}(A/\mathfrak{J})\to\mathfrak{X}$
（这里 ${}^a\varphi$ 正是 $\mathfrak{X}$ 到 $\mathfrak{X}$ 的恒等同胚）也称为
\emph{典范的}。由此得到\emph{典范同构}
\[
  \mathscr{O}_{\mathfrak{X}}/\mathfrak{J}^\Delta
  \xrightarrow{\sim}
  (\widetilde{A}/\widetilde{\mathfrak{J}})|\mathfrak{X}.
  \tag{10.3.2.1}\label{I.10.3.2.1-fr}
\]

显然（因为 $\Gamma(\mathfrak{X},\mathfrak{J}^\Delta)=\mathfrak{J}$），
映射 $\mathfrak{J}\to\mathfrak{J}^\Delta$ 是\emph{严格递增的}；
由上所述，当 $\mathfrak{J}\subset\mathfrak{J}'$ 时，层
${\mathfrak{J}'}^\Delta/\mathfrak{J}^\Delta$ 典范同构于
$\widetilde{\mathfrak{J}'}/\widetilde{\mathfrak{J}}
=\widetilde{\mathfrak{J}'/\mathfrak{J}}$。
\end{env}

\begin{env}[10.3.3]
\label{I.10.3.3-fr}
仍采用 \hyperref[I.10.3.1-fr]{10.3.1} 的假设和记号。称
$\mathscr{O}_{\mathfrak{X}}$ 的理想层 $\mathscr{J}$ 为
$\mathfrak{X}$ 的\emph{定义理想层}（或 $\mathfrak{X}$ 的一个
\emph{定义理想}），如果对每个 $x\in\mathfrak{X}$，都存在 $x$ 的一个
形如 $\mathfrak{D}(f)$ 的开邻域，其中 $f\in A$，使得
$\mathscr{J}|\mathfrak{D}(f)$ 形如 $\mathfrak{H}^\Delta$，这里
$\mathfrak{H}$ 是 $A_{\{f\}}$ 的一个定义理想。
\end{env}

\begin{proposition}[10.3.4]
\label{I.10.3.4-fr}
对每个 $f\in A$，$\mathfrak{X}$ 的任一定义理想在
$\mathfrak{D}(f)$ 上诱导一个定义理想。
\end{proposition}

\begin{proof}
这由 \hyperref[I.10.3.1.1-fr]{10.3.1.1} 立即得到。
\end{proof}

\begin{proposition}[10.3.5]
\label{I.10.3.5-fr}
若 $A$ 是可容环，则 $\mathfrak{X}=\operatorname{Spf}(A)$ 的每个定义理想
都具有形式 $\mathfrak{J}^\Delta$，其中 $\mathfrak{J}$ 是 $A$ 的一个
唯一确定的定义理想。
\end{proposition}

\begin{proof}
事实上，设 $\mathscr{J}$ 是 $\mathfrak{X}$ 的一个定义理想。由假设以及
$\mathfrak{X}$ 的拟紧性，存在有限多个元素 $f_i\in A$，使得
$\mathfrak{D}(f_i)$ 覆盖 $\mathfrak{X}$，且
$\mathscr{J}|\mathfrak{D}(f_i)=\mathfrak{H}_i^\Delta$，其中
$\mathfrak{H}_i$ 是 $A_{\{f_i\}}$ 的一个定义理想。对每个 $i$，存在
$A$ 的开理想 $\mathfrak{K}_i$，使
$(\mathfrak{K}_i)_{\{f_i\}}=\mathfrak{H}_i$
（\hyperref[0.7.6.9-fr]{0, 7.6.9}）；取 $A$ 的一个定义理想
$\mathfrak{K}$，使它包含于所有 $\mathfrak{K}_i$。由
\hyperref[I.10.3.2-fr]{10.3.2}，$\mathscr{J}/\mathfrak{K}^\Delta$
在 $\operatorname{Spec}(A/\mathfrak{K})$ 的结构层
$\widetilde{A/\mathfrak{K}}$ 中的典范像，在 $\mathfrak{D}(f_i)$ 上的
限制等于 $\widetilde{\mathfrak{K}_i/\mathfrak{K}}$ 的限制。因此该典范像是
$\operatorname{Spec}(A/\mathfrak{K})$ 上的准凝聚层，故具有形式
$\widetilde{\mathfrak{J}/\mathfrak{K}}$，其中 $\mathfrak{J}$ 是 $A$
中包含 $\mathfrak{K}$ 的理想（\hyperref[I.1.4.1-fr]{1.4.1}）。于是
$\mathscr{J}=\mathfrak{J}^\Delta$
（\hyperref[I.10.3.2-fr]{10.3.2}）。此外，对每个 $i$，存在整数
$n_i$，使 $\mathfrak{H}_i^{n_i}\subset\mathfrak{K}_{\{f_i\}}$；
若以 $n$ 表示这些 $n_i$ 中的最大者，则
$(\mathscr{J}/\mathfrak{K}^\Delta)^n=0$，从而由
\hyperref[I.10.3.2-fr]{10.3.2} 有
$(\widetilde{\mathfrak{J}/\mathfrak{K}})^n=0$，最终得到
$(\mathfrak{J}/\mathfrak{K})^n=0$
（\hyperref[I.1.3.13-fr]{1.3.13}）；这证明 $\mathfrak{J}$ 是 $A$ 的
定义理想（\hyperref[0.7.1.4-fr]{0, 7.1.4}）。
\end{proof}

\begin{proposition}[10.3.6]
\label{I.10.3.6-fr}
设 $A$ 是进制环，$\mathfrak{J}$ 是 $A$ 的定义理想，并且
$\mathfrak{J}/\mathfrak{J}^2$ 是有限型 $A/\mathfrak{J}$-模。
则对每个整数 $n>0$，有
$(\mathfrak{J}^\Delta)^n=(\mathfrak{J}^n)^\Delta$。
\end{proposition}

\begin{proof}
事实上，对每个 $f\in A$，由于 $\mathfrak{J}^n$ 是开理想，有
\[
  \bigl(\Gamma(\mathfrak{D}(f),\mathfrak{J}^\Delta)\bigr)^n
  =(\mathfrak{J}_{\{f\}})^n
  =(\mathfrak{J}^n)_{\{f\}}
  =\Gamma\bigl(\mathfrak{D}(f^n),(\mathfrak{J}^n)^\Delta\bigr)
\]
\oldpage[I]{185}
由 \hyperref[I.10.3.1.1-fr]{10.3.1.1} 和
\hyperref[0.7.6.12-fr]{0, 7.6.12} 即得。由于
$(\mathfrak{J}^\Delta)^n$ 伴随于预层
$U\mapsto(\Gamma(U,\mathfrak{J}^\Delta))^n$
（\hyperref[0.4.1.6-fr]{0, 4.1.6}），而 $\mathfrak{D}(f)$ 构成
$\mathfrak{X}$ 的拓扑基，结论得证。
\end{proof}

\begin{env}[10.3.7]
\label{I.10.3.7-fr}
称 $\mathfrak{X}$ 的定义理想族 $(\mathscr{J}_\lambda)$ 为一个
\emph{定义理想的基本系}，如果 $\mathfrak{X}$ 的每个定义理想都包含某个
$\mathscr{J}_\lambda$。由于
$\mathscr{J}_\lambda=\mathfrak{J}_\lambda^\Delta$，这等价于说
$\mathfrak{J}_\lambda$ 在 $A$ 中构成\emph{0 的基本邻域系}。设
$(f_\alpha)$ 是 $A$ 的一个元素族，使 $\mathfrak{D}(f_\alpha)$ 覆盖
$\mathfrak{X}$。若 $(\mathscr{J}_\lambda)$ 是
$\mathscr{O}_{\mathfrak{X}}$ 的一个滤过递减理想族，且对每个 $\alpha$，
族 $(\mathscr{J}_\lambda|\mathfrak{D}(f_\alpha))$ 是
$\mathfrak{D}(f_\alpha)$ 的定义理想的基本系，则
$(\mathscr{J}_\lambda)$ 是 $\mathfrak{X}$ 的定义理想的基本系。
事实上，对 $\mathfrak{X}$ 的任一定义理想 $\mathscr{J}$，存在一个由
$\mathfrak{D}(f_i)$ 组成的有限覆盖，使得对每个 $i$，
$\mathscr{J}_{\lambda_i}|\mathfrak{D}(f_i)$ 是
$\mathfrak{D}(f_i)$ 的定义理想，并包含于
$\mathscr{J}|\mathfrak{D}(f_i)$。若指标 $\mu$ 满足
$\mathscr{J}_\mu\subset\mathscr{J}_{\lambda_i}$ 对所有 $i$ 成立，
则由 \hyperref[I.10.3.3-fr]{10.3.3}，$\mathscr{J}_\mu$ 是
$\mathfrak{X}$ 的定义理想，且显然包含于 $\mathscr{J}$，断言得证。
\end{env}

\subsection{形式预概形及其态射}
\label{subsection:I.10.4-fr}

\begin{env}[10.4.1]
\label{I.10.4.1-fr}
给定拓扑赋环空间 $\mathfrak{X}$，若开集 $U\subset\mathfrak{X}$ 上由
$\mathfrak{X}$ 诱导的拓扑赋环空间是仿射形式概形（相应地，是环为进制环的
这种概形；相应地，是环为 Noether 进制环的这种概形），则称 $U$ 为
\emph{仿射形式开集}（相应地，\emph{进制仿射形式开集}；相应地，
\emph{Noether 仿射形式开集}）。
\end{env}

\begin{definition}[10.4.2]
\label{I.10.4.2-fr}
称拓扑赋环空间 $\mathfrak{X}$ 为形式预概形，如果它的每一点都有一个仿射
形式开邻域。如果 $\mathfrak{X}$ 的每一点都有一个进制仿射形式开邻域
（相应地，Noether 仿射形式开邻域），则称形式预概形 $\mathfrak{X}$ 是
进制的（相应地，局部 Noether 的）。如果 $\mathfrak{X}$ 局部 Noether，
且其底空间拟紧（因而是 Noether 空间），则称 $\mathfrak{X}$ 是 Noether 的。
\end{definition}

\begin{proposition}[10.4.3]
\label{I.10.4.3-fr}
若 $\mathfrak{X}$ 是形式预概形（相应地，局部 Noether 形式预概形），则仿射
形式开集（相应地，Noether 仿射开集）构成 $\mathfrak{X}$ 的拓扑基。
\end{proposition}

\begin{proof}
这由 \hyperref[I.10.4.2-fr]{10.4.2} 和
\hyperref[I.10.1.4-fr]{10.1.4} 得到；还要注意，若 $A$ 是 Noether
进制环，则对每个 $f\in A$，$A_{\{f\}}$ 也如此
（\hyperref[0.7.6.11-fr]{0, 7.6.11}）。
\end{proof}

\begin{corollary}[10.4.4]
\label{I.10.4.4-fr}
若 $\mathfrak{X}$ 是形式预概形（相应地，局部 Noether 形式预概形；相应地，
Noether 形式预概形），则在 $\mathfrak{X}$ 的任一开集上诱导的拓扑赋环空间
仍是形式预概形（相应地，局部 Noether 形式预概形；相应地，Noether
形式预概形）。
\end{corollary}

\begin{definition}[10.4.5]
\label{I.10.4.5-fr}
给定两个形式预概形 $\mathfrak{X}$、$\mathfrak{Y}$，称拓扑赋环空间的态射
$(\psi,\theta)$ 为从 $\mathfrak{X}$ 到 $\mathfrak{Y}$ 的（形式预概形）态射，
如果对每个 $x\in\mathfrak{X}$，$\theta_x^\sharp$ 都是局部同态
$\mathscr{O}_{\psi(x)}\to\mathscr{O}_x$。
\end{definition}

显然，两个形式预概形态射的复合仍是这样的态射；因此形式预概形构成一个
\emph{范畴}，并以 $\operatorname{Hom}(\mathfrak{X},\mathfrak{Y})$
表示从形式预概形 $\mathfrak{X}$ 到形式预概形 $\mathfrak{Y}$ 的态射集合。

\oldpage[I]{186}
若 $U$ 是 $\mathfrak{X}$ 的开子集，则由它在 $U$ 上诱导的
形式预概形到 $\mathfrak{X}$ 的典范嵌入是形式预概形态射（甚至是拓扑赋环
空间的\emph{单态射}（\hyperref[0.4.1.1-fr]{0, 4.1.1}））。

\begin{proposition}[10.4.6]
\label{I.10.4.6-fr}
设 $\mathfrak{X}$ 是形式预概形，
$\mathfrak{S}=\operatorname{Spf}(A)$ 是仿射形式概形。从形式预概形
$\mathfrak{X}$ 到形式预概形 $\mathfrak{S}$ 的态射，与从环 $A$ 到拓扑环
$\Gamma(\mathfrak{X},\mathscr{O}_{\mathfrak{X}})$ 的连续同态之间，存在典范
双射对应。
\end{proposition}

证明与 \hyperref[I.2.2.4-fr]{2.2.4} 的证明相同，只须把“同态”换成
“连续同态”，把“仿射开集”换成“仿射形式开集”，并用
\hyperref[I.10.2.2-fr]{10.2.2} 代替
\hyperref[I.1.7.3-fr]{1.7.3}；细节留给读者。

\begin{env}[10.4.7]
\label{I.10.4.7-fr}
给定形式预概形 $\mathfrak{S}$，若给出形式预概形 $\mathfrak{X}$ 和态射
$\varphi:\mathfrak{X}\to\mathfrak{S}$，则称这定义了一个
\emph{$\mathfrak{S}$ 上的}形式预概形 $\mathfrak{X}$，或一个
\emph{$\mathfrak{S}$-形式预概形}；$\varphi$ 称为该
$\mathfrak{S}$-形式预概形 $\mathfrak{X}$ 的\emph{结构态射}。若
$\mathfrak{S}=\operatorname{Spf}(A)$，其中 $A$ 是可容环，则也称
$\mathfrak{S}$-形式预概形 $\mathfrak{X}$ 为\emph{$A$-形式预概形}，
或称它为\emph{$A$ 上的}形式预概形。任一形式预概形总可以看成
$\mathbf{Z}$（赋予离散拓扑）上的形式预概形。

若 $\mathfrak{X}$、$\mathfrak{Y}$ 是两个 $\mathfrak{S}$-形式预概形，
且态射 $u:\mathfrak{X}\to\mathfrak{Y}$ 使图表
\[
  \xymatrix{
    \mathfrak{X}\ar[rr]^u\ar[rd] & &
    \mathfrak{Y}\ar[ld]\\
    & \mathfrak{S}
  }
\]
交换（其中斜箭头为结构态射），则称该态射为
\emph{$\mathfrak{S}$-态射}。在此定义下，对固定的 $\mathfrak{S}$，
$\mathfrak{S}$-形式预概形构成一个\emph{范畴}。以
$\operatorname{Hom}_{\mathfrak{S}}(\mathfrak{X},\mathfrak{Y})$
表示从 $\mathfrak{S}$-形式预概形 $\mathfrak{X}$ 到
$\mathfrak{S}$-形式预概形 $\mathfrak{Y}$ 的 $\mathfrak{S}$-态射集合。
当 $\mathfrak{S}=\operatorname{Spf}(A)$ 时，也用\emph{$A$-态射}
代替\emph{$\mathfrak{S}$-态射}。
\end{env}

\begin{env}[10.4.8]
\label{I.10.4.8-fr}
由于每个仿射概形都可以看成仿射形式概形
（\hyperref[I.10.1.2-fr]{10.1.2}），每个（通常的）预概形都可以看成形式
预概形。此外，由 \hyperref[I.10.4.5-fr]{10.4.5} 可知，对\emph{通常的}
预概形而言，\emph{形式}预概形的态射（相应地，$S$-态射）与
\S~\hyperref[section:I.2-fr]{2} 中定义的态射（相应地，$S$-态射）相同。
\end{env}

\subsection{形式预概形的定义理想}
\label{subsection:I.10.5-fr}

\begin{env}[10.5.1]
\label{I.10.5.1-fr}
设 $\mathfrak{X}$ 是形式预概形。若 $\mathscr{O}_{\mathfrak{X}}$-理想
$\mathscr{J}$ 满足：每个 $x\in\mathfrak{X}$ 都有一个仿射形式开邻域
$U$，使 $\mathscr{J}|U$ 是由 $\mathfrak{X}$ 在 $U$ 上诱导的仿射
形式概形的\emph{定义理想层}（或\emph{定义理想}）
（\hyperref[I.10.3.3-fr]{10.3.3}），则称其为
$\mathfrak{X}$ 的定义理想层（或定义理想）。由
\hyperref[I.10.3.1.1-fr]{10.3.1.1} 和
\hyperref[I.10.4.3-fr]{10.4.3}，对每个开集
$V\subset\mathfrak{X}$，$\mathscr{J}|V$ 都是由 $\mathfrak{X}$ 在
$V$ 上诱导的形式预概形的定义理想。

称 $\mathfrak{X}$ 的定义理想族 $(\mathscr{J}_\lambda)$ 为一个
\emph{基本系}
\oldpage[I]{187}
\emph{的定义理想}，如果存在由仿射形式开集组成的 $\mathfrak{X}$ 的覆盖
$(U_\alpha)$，使得对每个 $\alpha$，族
$\mathscr{J}_\lambda|U_\alpha$ 是由 $\mathfrak{X}$ 在 $U_\alpha$
上诱导的仿射形式概形的定义理想基本系
（\hyperref[I.10.3.6-fr]{10.3.6}）。由
\hyperref[I.10.3.7-fr]{10.3.7} 的最后说明，当 $\mathfrak{X}$ 是
仿射形式概形时，这一定义与 \hyperref[I.10.3.7-fr]{10.3.7} 中的定义
相同。对 $\mathfrak{X}$ 的每个开集 $V$，由
\hyperref[I.10.3.1.1-fr]{10.3.1.1}，这些限制
$\mathscr{J}_\lambda|V$ 构成在 $V$ 上诱导的形式预概形的定义理想基本系。
若 $\mathfrak{X}$ 是\emph{局部 Noether}形式预概形，且
$\mathscr{J}$ 是 $\mathfrak{X}$ 的定义理想，则由
\hyperref[I.10.3.6-fr]{10.3.6}，各次幂 $\mathscr{J}^n$ 构成
$\mathfrak{X}$ 的定义理想基本系。
\end{env}

\begin{env}[10.5.2]
\label{I.10.5.2-fr}
设 $\mathfrak{X}$ 是形式预概形，$\mathscr{J}$ 是 $\mathfrak{X}$ 的
定义理想。则赋环空间
$(\mathfrak{X},\mathscr{O}_{\mathfrak{X}}/\mathscr{J})$ 是一个
（通常的）\emph{预概形}；当 $\mathfrak{X}$ 是仿射形式概形
（相应地，局部 Noether 形式预概形；相应地，Noether 形式预概形）时，
它是仿射的（相应地，局部 Noether 的；相应地，Noether 的）。事实上，
问题立即归结到仿射情形，而该情形已在
\hyperref[I.10.3.2-fr]{10.3.2} 中证明。此外，若
$\theta:\mathscr{O}_{\mathfrak{X}}\to
\mathscr{O}_{\mathfrak{X}}/\mathscr{J}$ 是典范同态，则
$u=(1_{\mathfrak{X}},\theta)$ 是形式预概形的一个\emph{态射}
（称为\emph{典范态射}）
$(\mathfrak{X},\mathscr{O}_{\mathfrak{X}}/\mathscr{J})\to
(\mathfrak{X},\mathscr{O}_{\mathfrak{X}})$；因为问题仍可立即归结到已由
\hyperref[I.10.3.2-fr]{10.3.2} 处理的仿射情形。
\end{env}

\begin{proposition}[10.5.3]
\label{I.10.5.3-fr}
设 $\mathfrak{X}$ 是形式预概形，$(\mathscr{J}_\lambda)$ 是
$\mathfrak{X}$ 的定义理想基本系。则拓扑环层
$\mathscr{O}_{\mathfrak{X}}$ 是伪离散环层
（\hyperref[0.3.8.1-fr]{0, 3.8.1}）
$\mathscr{O}_{\mathfrak{X}}/\mathscr{J}_\lambda$ 的逆极限。
\end{proposition}

\begin{proof}
由于 $\mathfrak{X}$ 的拓扑有一个由拟紧仿射形式开集组成的拓扑基
（\hyperref[I.10.4.3-fr]{10.4.3}），问题归结到仿射情形；在该情形，
命题是 \hyperref[I.10.3.5-fr]{10.3.5}、
\hyperref[I.10.3.2-fr]{10.3.2} 以及定义
\hyperref[I.10.1.1-fr]{10.1.1} 的推论。
\end{proof}

并不能断定每个形式预概形都有定义理想。不过，有如下结果：

\begin{proposition}[10.5.4]
\label{I.10.5.4-fr}
设 $\mathfrak{X}$ 是局部 Noether 形式预概形。则 $\mathfrak{X}$ 有一个
最大的定义理想 $\mathscr{T}$；它是唯一使预概形
$(\mathfrak{X},\mathscr{O}_{\mathfrak{X}}/\mathscr{J})$ 约化的定义理想
$\mathscr{J}$。若 $\mathscr{J}$ 是 $\mathfrak{X}$ 的定义理想，则
$\mathscr{T}$ 是 $\mathscr{O}_{\mathfrak{X}}\to
\mathscr{O}_{\mathfrak{X}}/\mathscr{J}$ 下
$\mathscr{O}_{\mathfrak{X}}/\mathscr{J}$ 的幂零根的逆像。
\end{proposition}

\begin{proof}
先设 $\mathfrak{X}=\operatorname{Spf}(A)$，其中 $A$ 是 Noether 进制环。
$\mathscr{T}$ 的存在性及其性质由
% EGA1-ERR-0054: published EGA II erratum replaces 10.3.4 by 10.3.5; canonical Chinese remains source-literal.
\hyperref[I.10.3.4-fr]{10.3.4} 和
\hyperref[I.5.1.1-fr]{5.1.1} 立即得到；这里还使用了 $A$ 的最大定义理想
的存在性及其性质（\hyperref[0.7.1.6-fr]{0, 7.1.6} 和
\hyperref[0.7.1.7-fr]{7.1.7}）。

为证明一般情形下 $\mathscr{T}$ 的存在性及其性质，只须证明：若
$U\supset V$ 是 $\mathfrak{X}$ 的两个 Noether 仿射形式开集，则 $U$ 的
最大定义理想 $\mathscr{T}_U$ 在 $V$ 上诱导的正是其最大定义理想
$\mathscr{T}_V$。但由于
$(V,(\mathscr{O}_{\mathfrak{X}}|V)/(\mathscr{T}_U|V))$ 是约化的，
这由前述结果立即得到。
\end{proof}

以 $\mathfrak{X}_{\mathrm{red}}$ 表示约化的（通常）预概形
$(\mathfrak{X},\mathscr{O}_{\mathfrak{X}}/\mathscr{T})$。

\begin{corollary}[10.5.5]
\label{I.10.5.5-fr}
设 $\mathfrak{X}$ 是局部 Noether 形式预概形，$\mathscr{T}$ 是
$\mathfrak{X}$ 的最大定义理想；对 $\mathfrak{X}$ 的任一开集 $V$，
$\mathscr{T}|V$ 是由 $\mathfrak{X}$ 在 $V$ 上诱导的形式预概形的最大
定义理想。
\end{corollary}

\oldpage[I]{188}

\begin{proposition}[10.5.6]
\label{I.10.5.6-fr}
设 $\mathfrak{X}$、$\mathfrak{Y}$ 是两个形式预概形，$\mathscr{J}$
（相应地，$\mathscr{K}$）是 $\mathfrak{X}$（相应地，$\mathfrak{Y}$）
的定义理想，$f:\mathfrak{X}\to\mathfrak{Y}$ 是形式预概形态射。
\begin{enumerate}
  \item[\rm{(i)}] 若
    $f^*(\mathscr{K})\mathscr{O}_{\mathfrak{X}}\subset\mathscr{J}$，
    则存在唯一的通常预概形态射
    $f':(\mathfrak{X},\mathscr{O}_{\mathfrak{X}}/\mathscr{J})\to
    (\mathfrak{Y},\mathscr{O}_{\mathfrak{Y}}/\mathscr{K})$，使图表
    \[
      \label{I.10.5.6.1-fr}
      \xymatrix{
        (\mathfrak{X},\mathscr{O}_{\mathfrak{X}})\ar[r]^f &
        (\mathfrak{Y},\mathscr{O}_{\mathfrak{Y}})\\
        (\mathfrak{X},\mathscr{O}_{\mathfrak{X}}/\mathscr{J})
          \ar[r]_{f'}\ar[u] &
        (\mathfrak{Y},\mathscr{O}_{\mathfrak{Y}}/\mathscr{K})\ar[u]
      }
      \tag{10.5.6.1}
    \]
    交换，其中竖箭头是典范态射。
  \item[\rm{(ii)}] 设
    $\mathfrak{X}=\operatorname{Spf}(A)$、
    $\mathfrak{Y}=\operatorname{Spf}(B)$ 是仿射形式概形，
    $\mathscr{J}=\mathfrak{J}^{\Delta}$、
    $\mathscr{K}=\mathfrak{K}^{\Delta}$，其中 $\mathfrak{J}$
    （相应地，$\mathfrak{K}$）是 $A$（相应地，$B$）的定义理想，且
    $f=({}^a\varphi,\widetilde{\varphi})$，其中
    $\varphi:B\to A$ 是连续同态。条件
    $f^*(\mathscr{K})\mathscr{O}_{\mathfrak{X}}\subset\mathscr{J}$
    成立，当且仅当 $\varphi(\mathfrak{K})\subset\mathfrak{J}$；此时
    $f'$ 是态射 $({}^a\varphi',\widetilde{\varphi'})$，其中
    $\varphi':B/\mathfrak{K}\to A/\mathfrak{J}$ 是 $\varphi$ 取商后
    诱导的同态。
\end{enumerate}
\end{proposition}

\begin{proof}
\medskip\noindent
\begin{enumerate}
  \item[\rm{(i)}] 若 $f=(\psi,\theta)$，则假设意味着
    $\theta^\sharp:\psi^*(\mathscr{O}_{\mathfrak{Y}})\to
    \mathscr{O}_{\mathfrak{X}}$ 把
    $\psi^*(\mathscr{O}_{\mathfrak{Y}})$ 的理想层
    $\psi^*(\mathscr{K})$ 映入 $\mathscr{J}$
    （\hyperref[0.4.3.5-fr]{0, 4.3.5}）。取商后，$\theta^\sharp$
    因而诱导环层同态
    \[
      \omega:\psi^*(\mathscr{O}_{\mathfrak{Y}}/\mathscr{K})
      =\psi^*(\mathscr{O}_{\mathfrak{Y}})/\psi^*(\mathscr{K})
      \to\mathscr{O}_{\mathfrak{X}}/\mathscr{J}~;
    \]
    此外，由于对每个 $x\in\mathfrak{X}$，$\theta_x^\sharp$ 是
    \emph{局部}同态，$\omega_x$ 也如此。因此赋环空间态射
    $(\psi,\omega^b)$ 正是满足要求的唯一赋环空间态射 $f'$
    （\hyperref[I.2.2.1-fr]{2.2.1}）。
  \item[\rm{(ii)}] 仿射形式预概形态射与连续环同态之间的典范函子性对应
    （\hyperref[I.10.2.2-fr]{10.2.2}）表明，在当前情形，关系
    $f^*(\mathscr{K})\mathscr{O}_{\mathfrak{X}}\subset\mathscr{J}$
    蕴含 $f'=({}^a\varphi',\widetilde{\varphi'})$，其中
    $\varphi':B/\mathfrak{K}\to A/\mathfrak{J}$ 是使图表
    \[
      \label{I.10.5.6.2-fr}
      \xymatrix{
        B\ar[r]^\varphi\ar[d] & A\ar[d]\\
        B/\mathfrak{K}\ar[r]_{\varphi'} & A/\mathfrak{J}
      }
      \tag{10.5.6.2}
    \]
    交换的唯一同态。$\varphi'$ 的存在因而蕴含
    $\varphi(\mathfrak{K})\subset\mathfrak{J}$。反之，若此条件成立，
    以 $\varphi'$ 表示使图表
    \hyperref[I.10.5.6.2-fr]{10.5.6.2} 交换的唯一同态，并令
    $f'=({}^a\varphi',\widetilde{\varphi'})$，则图表
    \hyperref[I.10.5.6.1-fr]{10.5.6.1} 显然交换；再考察分别对应于
    $f$ 和 $f'$ 的同态
    ${}^a\varphi^*(\mathscr{O}_{\mathfrak{Y}})\to
    \mathscr{O}_{\mathfrak{X}}$ 与
    ${}^a{\varphi'}^*(\mathscr{O}_{\mathfrak{Y}}/\mathscr{K})\to
    \mathscr{O}_{\mathfrak{X}}/\mathscr{J}$，便知这又蕴含关系
    $f^*(\mathscr{K})\mathscr{O}_{\mathfrak{X}}\subset\mathscr{J}$。
\end{enumerate}
\end{proof}

显然，上述定义的对应 $f\mapsto f'$ 是\emph{函子性的}。

\subsection{作为预概形之归纳极限的形式预概形}
\label{subsection:I.10.6-fr}

\begin{env}[10.6.1]
\label{I.10.6.1-fr}
设 $\mathfrak{X}$ 是形式预概形，$(\mathscr{J}_\lambda)$ 是
$\mathfrak{X}$ 的定义理想基本系；对每个 $\lambda$，令 $f_\lambda$ 为
典范态射
$(\mathfrak{X},\mathscr{O}_{\mathfrak{X}}/\mathscr{J}_\lambda)
\to\mathfrak{X}$（\hyperref[I.10.5.2-fr]{10.5.2}）。当
$\mathscr{J}_\mu\subset\mathscr{J}_\lambda$ 时，典范同态
$\mathscr{O}_{\mathfrak{X}}/\mathscr{J}_\mu\to
\mathscr{O}_{\mathfrak{X}}/\mathscr{J}_\lambda$ 定义一个典范态射

\oldpage[I]{189}

$f_{\mu\lambda}:(\mathfrak{X},
\mathscr{O}_{\mathfrak{X}}/\mathscr{J}_\lambda)\to
(\mathfrak{X},\mathscr{O}_{\mathfrak{X}}/\mathscr{J}_\mu)$，它是通常预概形
之间的态射，并且 $f_\lambda=f_\mu\circ f_{\mu\lambda}$。因此，预概形
$X_\lambda=(\mathfrak{X},
\mathscr{O}_{\mathfrak{X}}/\mathscr{J}_\lambda)$ 和态射
$f_{\mu\lambda}$ 在形式预概形范畴中构成一个\emph{归纳系}
（由 \hyperref[I.10.4.8-fr]{10.4.8}）。
\end{env}

\begin{proposition}[10.6.2]
\label{I.10.6.2-fr}
采用 \hyperref[I.10.6.1-fr]{10.6.1} 的记号，形式预概形
$\mathfrak{X}$ 和态射 $f_\lambda$ 在形式预概形范畴中构成系统
$(X_\lambda,f_{\mu\lambda})$ 的一个归纳极限（T, I, 1.8）。
\end{proposition}

\begin{proof}
设 $\mathfrak{Y}$ 是形式预概形，并且对每个指标 $\lambda$，有一个态射
\[
  g_\lambda=(\psi_\lambda,\theta_\lambda):X_\lambda\to\mathfrak{Y}
\]
满足：当 $\mathscr{J}_\mu\subset\mathscr{J}_\lambda$ 时，
$g_\lambda=g_\mu\circ f_{\mu\lambda}$。这一条件以及 $X_\lambda$ 的定义
首先说明，所有 $\psi_\lambda$ 都等于同一个底空间上的连续映射
$\psi:\mathfrak{X}\to\mathfrak{Y}$；此外，各同态
$\theta_\lambda^\sharp:\psi^*(\mathscr{O}_{\mathfrak{Y}})\to
\mathscr{O}_{X_\lambda}=
\mathscr{O}_{\mathfrak{X}}/\mathscr{J}_\lambda$ 构成环层同态的一个
\emph{逆系统}。取逆极限便得到同态
\[
  \omega:\psi^*(\mathscr{O}_{\mathfrak{Y}})\to
  \varprojlim\mathscr{O}_{\mathfrak{X}}/\mathscr{J}_\lambda
  =\mathscr{O}_{\mathfrak{X}},
\]
而且显然，\emph{赋环空间}的态射 $g=(\psi,\omega^b)$ 是使各图表
\[
  \label{I.10.6.2.1-fr}
  \xymatrix{
    X_\lambda\ar[rr]^{g_\lambda}\ar[rd]_{f_\lambda} &&
    \mathfrak{Y}\\
    & \mathfrak{X}\ar[ru]_g
  }
  \tag{10.6.2.1}
\]
交换的\emph{唯一}态射。尚须证明 $g$ 是\emph{形式预概形}的态射。
由于问题对于 $\mathfrak{X}$ 和 $\mathfrak{Y}$ 是局部的，可以设
$\mathfrak{X}=\operatorname{Spf}(A)$，
$\mathfrak{Y}=\operatorname{Spf}(B)$，其中 $A$ 和 $B$ 是可容环，并且
$\mathscr{J}_\lambda=\mathfrak{J}_\lambda^{\Delta}$；这里
$(\mathfrak{J}_\lambda)$ 是 $A$ 的定义理想基本系
（\hyperref[I.10.3.5-fr]{10.3.5}）。由于
$A=\varprojlim A/\mathfrak{J}_\lambda$，使图表
\hyperref[I.10.6.2.1-fr]{10.6.2.1} 交换的仿射形式概形态射 $g$ 的存在性，
来自仿射形式概形态射与连续环同态之间的一一对应
（\hyperref[I.10.2.2-fr]{10.2.2}）以及逆极限的定义。不过，$g$ 作为
赋环空间态射的唯一性表明，它就是证明开头所记的同名态射。
\end{proof}

下一个命题给出附加条件，以保证给定的通常预概形归纳系在形式预概形
范畴中的归纳极限存在：

\begin{proposition}[10.6.3]
\label{I.10.6.3-fr}
设 $\mathfrak{X}$ 是拓扑空间，$(\mathscr{O}_i,u_{ji})$ 是
$\mathfrak{X}$ 上以 $\mathbf{N}$ 为指标集的环层逆系统。令
$\mathscr{J}_i$ 为 $u_{0i}:\mathscr{O}_i\to\mathscr{O}_0$ 的核。假设：
\begin{enumerate}
  \item[a)] 赋环空间 $(\mathfrak{X},\mathscr{O}_i)$ 是预概形 $X_i$。
  % EGA1-ERR-0055: published EGA II erratum replaces X by \mathfrak{X}; canonical Chinese remains source-literal.
  \item[b)] 对每个 $x\in X$ 和每个 $i$，存在 $x$ 在
    $\mathfrak{X}$ 中的开邻域 $U_i$，使得限制
    $\mathscr{J}_i|U_i$ 是幂零的。
  \item[c)] 各同态 $u_{ji}$ 都是满射。
\end{enumerate}

\oldpage[I]{190}

令 $\mathscr{O}_{\mathfrak{X}}$ 为各伪离散环层 $\mathscr{O}_i$ 的
逆极限所得的拓扑环层，并令
$u_i:\mathscr{O}_{\mathfrak{X}}\to\mathscr{O}_i$ 为典范同态。那么，
拓扑赋环空间 $(\mathfrak{X},\mathscr{O}_{\mathfrak{X}})$ 是形式预概形；
各同态 $u_i$ 都是满射；它们的核 $\mathscr{J}^{(i)}$ 构成
$\mathfrak{X}$ 的定义理想基本系，并且 $\mathscr{J}^{(0)}$ 是各理想层
$\mathscr{J}_i$ 的逆极限。
\end{proposition}

\begin{proof}
首先注意，在每个茎上，$u_{ji}$ 是满同态，\emph{更不用说}是局部同态；
因此，$v_{ij}=(1_{\mathfrak{X}},u_{ji})$ 是预概形态射
$X_j\to X_i$（$i\geq j$）（\hyperref[I.2.2.1-fr]{2.2.1}）。先假设
每个 $X_i$ 都是环 $A_i$ 的仿射概形。存在一个\emph{环}同态
$\varphi_{ji}:A_i\to A_j$，使得
$u_{ji}=\widetilde{\varphi_{ji}}$
（\hyperref[I.1.7.3-fr]{1.7.3}）。于是由
\hyperref[I.1.6.3-fr]{1.6.3}，层 $\mathscr{O}_j$ 是 $X_i$ 上的
准凝聚 $\mathscr{O}_i$-模层（其外运算由 $u_{ji}$ 定义），它关联于
借助 $\varphi_{ji}$ 看作 $A_i$-模的 $A_j$。对每个 $f\in A_i$，令
$f'=\varphi_{ji}(f)$。按假设，开集 $D(f)$ 和 $D(f')$ 在
$\mathfrak{X}$ 中相同，而由 $u_{ji}$ 对应的同态
$\Gamma(D(f),\mathscr{O}_i)=(A_i)_f$ 到
$\Gamma(D(f),\mathscr{O}_j)=(A_j)_{f'}$ 正是
$(\varphi_{ji})_f$（\hyperref[I.1.6.1-fr]{1.6.1}）。但把 $A_j$
看作 $A_i$-模时，$(A_j)_{f'}$ 就是 $(A_i)_f$-模 $(A_j)_f$，所以若
这次把 $\varphi_{ji}$ 看作 $A_i$-模同态，仍有
$u_{ji}=\widetilde{\varphi_{ji}}$。于是，由于 $u_{ji}$ 是满射，
$\varphi_{ji}$ 也是满射（\hyperref[I.1.3.9-fr]{1.3.9}）；若
$\mathfrak{J}_{ji}$ 是 $\varphi_{ji}$ 的核，则 $u_{ji}$ 的核是等于
$\widetilde{\mathfrak{J}_{ji}}$ 的准凝聚 $\mathscr{O}_i$-模层。特别地，
$\mathscr{J}_i=\widetilde{\mathfrak{J}_i}$，其中 $\mathfrak{J}_i$ 是
$\varphi_{0i}:A_i\to A_0$ 的核。假设 b) 蕴含 $\mathscr{J}_i$ 是
\emph{幂零的}：事实上，由于 $\mathfrak{X}$ 是拟紧的，可用有限多个开集
$U_k$ 覆盖 $\mathfrak{X}$，并使
$(\mathscr{J}_i|U_k)^{n_k}=0$；取各 $n_k$ 中的最大者为 $n$，便有
$\mathscr{J}_i^n=0$。于是 $\mathfrak{J}_i$ 是幂零的
（\hyperref[I.1.3.13-fr]{1.3.13}）。因此，环
$A=\varprojlim A_i$ 是可容的（\hyperref[0.7.2.2-fr]{0, 7.2.2}），
典范同态 $\varphi_i:A\to A_i$ 是满射，其核
$\mathfrak{J}^{(i)}$ 等于当 $k\geq i$ 时各 $\mathfrak{J}_{ik}$ 的逆极限；
各 $\mathfrak{J}^{(i)}$ 构成 $A$ 中 0 的邻域基本系。在这一情形，
\hyperref[I.10.6.3-fr]{10.6.3} 的断言由
\hyperref[I.10.1.1-fr]{10.1.1} 和
\hyperref[I.10.3.2-fr]{10.3.2} 得出，因为
$(\mathfrak{X},\mathscr{O}_{\mathfrak{X}})$ 正是
$\operatorname{Spf}(A)$。

仍在上述特殊情形中，注意若 $f=(f_i)$ 是逆极限
$A=\varprojlim A_i$ 的元素，则所有开集 $D(f_i)$（它是 $X_i$ 中的仿射
开集）都与 $\mathfrak{X}$ 的开集 $\mathfrak{D}(f)$ 等同；因而由
$X_i$ 在 $\mathfrak{D}(f)$ 上诱导的预概形，与仿射概形
$\operatorname{Spec}((A_i)_{f_i})$ 等同。

在一般情形中，先注意，对 $\mathfrak{X}$ 的每个拟紧开集 $U$，上述论证
表明每个 $\mathscr{J}_i|U$ 都是幂零的。我们将证明：对每个
$x\in\mathfrak{X}$，都存在 $x$ 在 $\mathfrak{X}$ 中的一个开邻域 $U$，
它对于所有 $X_i$ 都是一个\emph{仿射开集}。事实上，取 $U$ 为
$X_0$ 的仿射开集，并注意
$\mathscr{O}_{X_0}=\mathscr{O}_{X_i}/\mathscr{J}_i$。由于上述结果表明
$\mathscr{J}_i|U$ 是幂零的，故由
\hyperref[I.5.1.9-fr]{5.1.9}，$U$ 对每个 $X_i$ 也是仿射开集。因此，
对每个满足上述条件的 $U$，前面对仿射情形的研究表明
$(U,\mathscr{O}_{\mathfrak{X}}|U)$ 是形式预概形，其中各
$\mathscr{J}^{(i)}|U$ 构成定义理想基本系，而
$\mathscr{J}^{(0)}|U$ 是各 $\mathscr{J}_i|U$ 的逆极限；结论由此得证。
\end{proof}

\begin{corollary}[10.6.4]
\label{I.10.6.4-fr}
假设当 $i\geq j$ 时，$u_{ji}$ 的核为 $\mathscr{J}_i^{j+1}$，并且
$\mathscr{J}_1/\mathscr{J}_1^2$

\oldpage[I]{191}

在 $\mathscr{O}_0=\mathscr{O}_1/\mathscr{J}_1$ 上是有限型的。那么
$\mathfrak{X}$ 是进制形式预概形；若 $\mathscr{J}^{(n)}$ 是
$\mathscr{O}_{\mathfrak{X}}\to\mathscr{O}_n$ 的核，则
$\mathscr{J}^{(n)}=\mathscr{J}^{n+1}$，而
$\mathscr{J}/\mathscr{J}^2$ 同构于 $\mathscr{J}_1$。若此外 $X_0$ 是
局部 Noether 的（相应地，Noether 的），则 $\mathfrak{X}$ 是局部
Noether 的（相应地，Noether 的）。
\end{corollary}

\begin{proof}
由于 $\mathfrak{X}$ 与 $X_0$ 的底空间相同，问题是局部的，可以假设
所有 $X_i$ 都是仿射的。考虑关系
$\mathscr{J}_{ij}=\widetilde{\mathfrak{J}_{ji}}$（采用
\hyperref[I.10.6.3-fr]{10.6.3} 的记号），问题立即归结为
\hyperref[0.7.2.7-fr]{0, 7.2.7} 和
\hyperref[0.7.2.8-fr]{7.2.8} 的相应断言；还要注意，此时
$\mathfrak{J}_1/\mathfrak{J}_1^2$ 是有限型 $A_0$-模
（\hyperref[I.1.3.9-fr]{1.3.9}）。
\end{proof}

特别地，\emph{每个局部 Noether 形式预概形 $\mathfrak{X}$} 都是满足
\hyperref[I.10.6.3-fr]{10.6.3} 和
\hyperref[I.10.6.4-fr]{10.6.4} 条件的一列局部 Noether 通常预概形
$(X_n)$ 的归纳极限：只须取 $\mathfrak{X}$ 的一个定义理想
$\mathscr{J}$（\hyperref[I.10.5.4-fr]{10.5.4}），并令
$X_n=(\mathfrak{X},
\mathscr{O}_{\mathfrak{X}}/\mathscr{J}^{n+1})$
（\hyperref[I.10.5.1-fr]{10.5.1} 与
\hyperref[I.10.6.2-fr]{10.6.2}）。

\begin{corollary}[10.6.5]
\label{I.10.6.5-fr}
设 $A$ 是可容环。仿射形式概形
$\mathfrak{X}=\operatorname{Spf}(A)$ 为 Noether 的充分必要条件是 $A$
为 Noether 进制环。
\end{corollary}

\begin{proof}
此条件显然充分。反之，设 $\mathfrak{X}$ 是 Noether 的，
$\mathfrak{J}$ 是 $A$ 的定义理想，
$\mathscr{J}=\mathfrak{J}^\Delta$ 是 $\mathfrak{X}$ 对应的定义理想。
此时通常预概形
$X_n=(\mathfrak{X},
\mathscr{O}_{\mathfrak{X}}/\mathscr{J}^{n+1})$ 是仿射且 Noether 的，
故各环 $A_n=A/\mathfrak{J}^{n+1}$ 是 Noether 的
（\hyperref[I.6.1.3-fr]{6.1.3}），从而
$\mathfrak{J}/\mathfrak{J}^2$ 是有限型 $A/\mathfrak{J}$-模。由于各
$\mathscr{J}^n$ 构成 $\mathfrak{X}$ 的定义理想基本系
（\hyperref[I.10.5.1-fr]{10.5.1}），有
$\mathscr{O}_{\mathfrak{X}}
=\varprojlim(\mathscr{O}_{\mathfrak{X}}/\mathscr{J}^n)$
（\hyperref[I.10.5.3-fr]{10.5.3}）。由此以及
\hyperref[I.10.1.3-fr]{10.1.3} 可知，$A$ 拓扑同构于
$\varprojlim A/\mathfrak{J}^n$，所以是 Noether 进制环
（\hyperref[0.7.2.8-fr]{0, 7.2.8}）。
\end{proof}

\begin{remark}[10.6.6]
\label{I.10.6.6-fr}
采用 \hyperref[I.10.6.3-fr]{10.6.3} 的记号，设 $\mathscr{F}_i$ 是
$\mathscr{O}_i$-模层，并假设当 $i\geq j$ 时给定一个
$v_{ij}$-态射 $\theta_{ji}:\mathscr{F}_i\to\mathscr{F}_j$，使得当
$k\leq j\leq i$ 时，
$\theta_{kj}\circ\theta_{ji}=\theta_{ki}$。由于 $v_{ij}$ 的底连续映射
是恒等映射，$\theta_{ji}$ 是空间 $\mathfrak{X}$ 上 Abel 群层的同态。
此外，若 $\mathscr{F}$ 是 Abel 群层逆系统 $(\mathscr{F}_i)$ 的逆极限，
则各 $\theta_{ji}$ 为 $v_{ij}$-态射这一事实，使我们能通过取逆极限在
$\mathscr{F}$ 上定义 $\mathscr{O}_{\mathfrak{X}}$-模结构。赋予这一结构
后，称 $\mathscr{F}$ 为 $\mathscr{O}_i$-模系 $(\mathscr{F}_i)$（关于
$\theta_{ji}$）的\emph{逆极限}。在特殊情形
$v_{ij}^*(\mathscr{F}_i)=\mathscr{F}_j$ 且 $\theta_{ji}$ 为
\emph{恒等映射}时，我们简称 $\mathscr{F}$ 为一个满足
$v_{ij}^*(\mathscr{F}_i)=\mathscr{F}_j$（$j\leq i$）的系统
$(\mathscr{F}_i)$ 的逆极限，而不再提及 $\theta_{ji}$。
\end{remark}

\begin{env}[10.6.7]
\label{I.10.6.7-fr}
设 $\mathfrak{X}$、$\mathfrak{Y}$ 是两个形式预概形，$\mathscr{J}$
（相应地，$\mathscr{K}$）是 $\mathfrak{X}$（相应地，
$\mathfrak{Y}$）的定义理想，而
$f:\mathfrak{X}\to\mathfrak{Y}$ 是满足
$f^*(\mathscr{K})\mathscr{O}_{\mathfrak{X}}\subset\mathscr{J}$ 的态射。
于是，对每个整数 $n>0$，有
$f^*(\mathscr{K}^n)\mathscr{O}_{\mathfrak{X}}
=(f^*(\mathscr{K})\mathscr{O}_{\mathfrak{X}})^n
\subset\mathscr{J}^n$；因而由
\hyperref[I.10.5.6-fr]{10.5.6}，可从 $f$ 得到通常预概形的态射
$f_n:X_n\to Y_n$，其中
$X_n=(\mathfrak{X},
\mathscr{O}_{\mathfrak{X}}/\mathscr{J}^{n+1})$，
$Y_n=(\mathfrak{Y},
\mathscr{O}_{\mathfrak{Y}}/\mathscr{K}^{n+1})$。由定义立即可知，各图表
\[
  \label{I.10.6.7.1-fr}
  \xymatrix{
    X_m\ar[r]^{f_m}\ar[d] &
    Y_m\ar[d]\\
    X_n\ar[r]_{f_n} &
    Y_n
  }
  \tag{10.6.7.1}
\]

\oldpage[I]{192}

当 $m\leq n$ 时交换；换言之，$(f_n)$ 是态射的一个\emph{归纳系}。
\end{env}

\begin{env}[10.6.8]
\label{I.10.6.8-fr}
反之，设 $(X_n)$（相应地，$(Y_n)$）是满足
\hyperref[I.10.6.3-fr]{10.6.3} 的条件 b) 和 c) 的通常预概形归纳系，
而 $\mathfrak{X}$（相应地，$\mathfrak{Y}$）是它的归纳极限。由归纳极限
的定义，每个构成归纳系的态射列 $(f_n)$，$X_n\to Y_n$，都有一个
\emph{归纳极限} $f:\mathfrak{X}\to\mathfrak{Y}$；它是使各图表
\[
  \xymatrix{
    X_n\ar[r]^{f_n}\ar[d] &
    Y_n\ar[d]\\
    \mathfrak{X}\ar[r]_f &
    \mathfrak{Y}
  }
\]
交换的唯一形式预概形态射。
\end{env}

\begin{proposition}[10.6.9]
\label{I.10.6.9-fr}
设 $\mathfrak{X}$、$\mathfrak{Y}$ 是两个局部 Noether 形式预概形，
$\mathscr{J}$（相应地，$\mathscr{K}$）是 $\mathfrak{X}$
（相应地，$\mathfrak{Y}$）的定义理想。于是，在
\hyperref[I.10.6.7-fr]{10.6.7} 中定义的映射 $f\mapsto(f_n)$，给出从
满足 $f^*(\mathscr{K})\mathscr{O}_{\mathfrak{X}}\subset\mathscr{J}$ 的态射
$f:\mathfrak{X}\to\mathfrak{Y}$ 所成集合，到使各图表
\hyperref[I.10.6.7.1-fr]{10.6.7.1} 交换的态射列 $(f_n)$ 所成集合的
双射。
\end{proposition}

\begin{proof}
若 $f$ 是这样一个态射列的归纳极限，须证明
$f^*(\mathscr{K})\mathscr{O}_{\mathfrak{X}}\subset\mathscr{J}$。由于问题对
$\mathfrak{X}$ 和 $\mathfrak{Y}$ 是局部的，只须考虑仿射情形
$\mathfrak{X}=\operatorname{Spf}(A)$、
$\mathfrak{Y}=\operatorname{Spf}(B)$，其中 $A$、$B$ 是 Noether 进制环，
$\mathscr{J}=\mathfrak{J}^\Delta$，
$\mathscr{K}=\mathfrak{K}^\Delta$，而 $\mathfrak{J}$（相应地，
$\mathfrak{K}$）是 $A$（相应地，$B$）的定义理想。由
\hyperref[I.10.3.6-fr]{10.3.6} 和
\hyperref[I.10.3.2-fr]{10.3.2}，此时
$X_n=\operatorname{Spec}(A_n)$，
$Y_n=\operatorname{Spec}(B_n)$，其中
$A_n=A/\mathfrak{J}^{n+1}$ 和 $B_n=B/\mathfrak{K}^{n+1}$；并且
$f_n=({}^a\varphi_n,\widetilde{\varphi_n})$，其中各同态
$\varphi_n:B_n\to A_n$ 构成一个逆系统，因而
$f=({}^a\varphi,\widetilde{\varphi})$，其中
$\varphi=\varprojlim\varphi_n$。取 $m=0$，图表
\hyperref[I.10.6.7.1-fr]{10.6.7.1} 的交换性给出条件
$\varphi_n(\mathfrak{K}/\mathfrak{K}^{n+1})
\subset\mathfrak{J}/\mathfrak{J}^{n+1}$（对每个 $n$）。取逆极限便有
$\varphi(\mathfrak{K})\subset
\mathfrak{J}$，从而
$f^*(\mathscr{K})\mathscr{O}_{\mathfrak{X}}\subset\mathscr{J}$
（\hyperref[I.10.5.6-fr]{10.5.6, (ii)}）。
\end{proof}

\begin{corollary}[10.6.10]
\label{I.10.6.10-fr}
设 $\mathfrak{X}$、$\mathfrak{Y}$ 是两个局部 Noether 形式预概形，
$\mathscr{T}$ 是 $\mathfrak{X}$ 的最大定义理想
（\hyperref[I.10.5.4-fr]{10.5.4}）。
\begin{enumerate}
  \item[\rm{(i)}] 对 $\mathfrak{Y}$ 的每个定义理想 $\mathscr{K}$ 以及
  每个态射 $f:\mathfrak{X}\to\mathfrak{Y}$，都有
  $f^*(\mathscr{K})\mathscr{O}_{\mathfrak{X}}\subset\mathscr{T}$。
  \item[\rm{(ii)}] $\operatorname{Hom}(\mathfrak{X},\mathfrak{Y})$ 与
  使各图表 \hyperref[I.10.6.7.1-fr]{10.6.7.1} 交换的态射列
  $(f_n)$ 所成集合之间有典范一一对应，其中
  $X_n=(\mathfrak{X},
  \mathscr{O}_{\mathfrak{X}}/\mathscr{T}^{n+1})$，
  $Y_n=(\mathfrak{Y},
  \mathscr{O}_{\mathfrak{Y}}/\mathscr{K}^{n+1})$。
\end{enumerate}
\end{corollary}

\begin{proof}
(ii) 立即由 (i) 和 \hyperref[I.10.6.9-fr]{10.6.9} 得出。为证明 (i)，
只须考虑 $\mathfrak{X}=\operatorname{Spf}(A)$、
$\mathfrak{Y}=\operatorname{Spf}(B)$ 的情形，其中 $A$、$B$ 是 Noether
环，$\mathscr{T}=\mathfrak{T}^\Delta$，
$\mathscr{K}=\mathfrak{K}^\Delta$，$\mathfrak{T}$ 是 $A$ 的最大定义理想，
而 $\mathfrak{K}$ 是 $B$ 的定义理想。设
$f=({}^a\varphi,\widetilde{\varphi})$，其中
$\varphi:B\to A$ 是连续同态。由于 $\mathfrak{K}$ 的元素都是拓扑幂零的
（\hyperref[0.7.1.4-fr]{0, 7.1.4, (ii)}），
$\varphi(\mathfrak{K})$ 的元素也都是拓扑幂零的；又因为
$\mathfrak{T}$ 是 $A$ 中所有拓扑幂零元素所成的集合
（\hyperref[0.7.1.6-fr]{0, 7.1.6}），所以
$\varphi(\mathfrak{K})\subset\mathfrak{T}$。结论由
\hyperref[I.10.5.6-fr]{10.5.6, (ii)} 得出。
\end{proof}

\begin{corollary}[10.6.11]
\label{I.10.6.11-fr}
设 $\mathfrak{S}$、$\mathfrak{X}$、$\mathfrak{Y}$ 是三个局部 Noether
形式预概形，$f:\mathfrak{X}\to\mathfrak{S}$、
$g:\mathfrak{Y}\to\mathfrak{S}$ 是使 $\mathfrak{X}$ 和
$\mathfrak{Y}$ 成为形式 $\mathfrak{S}$-预概形的态射。设
$\mathscr{J}$（相应地，$\mathscr{K}$、$\mathscr{L}$）是
$\mathfrak{S}$（相应地，$\mathfrak{X}$、$\mathfrak{Y}$）的定义理想，
并假设
$f^*(\mathscr{J})\mathscr{O}_{\mathfrak{X}}\subset\mathscr{K}$，
$g^*(\mathscr{J})\mathscr{O}_{\mathfrak{Y}}=\mathscr{L}$。置
$S_n=(\mathfrak{S},
\mathscr{O}_{\mathfrak{S}}/\mathscr{J}^{n+1})$，
$X_n=(\mathfrak{X},
\mathscr{O}_{\mathfrak{X}}/\mathscr{K}^{n+1})$，
$Y_n=(\mathfrak{Y},
\mathscr{O}_{\mathfrak{Y}}/\mathscr{L}^{n+1})$。于是有典范一一对应

\oldpage[I]{193}

$\operatorname{Hom}_{\mathfrak{S}}(\mathfrak{X},\mathfrak{Y})$ 与使图表
\hyperref[I.10.6.7.1-fr]{10.6.7.1} 交换的 $S_n$-态射列
$(u_n)$，$u_n:X_n\to Y_n$，所成集合之间的典范一一对应。
\end{corollary}

\begin{proof}
对每个 $\mathfrak{S}$-态射 $u:\mathfrak{X}\to\mathfrak{Y}$，由定义有
$f=g\circ u$，所以
\[
  u^*(\mathscr{L})\mathscr{O}_{\mathfrak{X}}
  =u^*(g^*(\mathscr{J})\mathscr{O}_{\mathfrak{Y}})
    \mathscr{O}_{\mathfrak{X}}
  =f^*(\mathscr{J})\mathscr{O}_{\mathfrak{X}}
  \subset\mathscr{K}
\]
故推论由 \hyperref[I.10.6.9-fr]{10.6.9} 得出。
\end{proof}

注意，当 $m\leq n$ 时，给定一个态射 $f_n:X_n\to Y_n$，便唯一确定
一个使图表 \hyperref[I.10.6.7.1-fr]{10.6.7.1} 交换的态射
$f_m:X_m\to Y_m$；将问题归结到仿射情形即可立即看出这一点。这样定义了
映射 $\varphi_{mn}:\operatorname{Hom}_{S_n}(X_n,Y_n)
\to\operatorname{Hom}_{S_m}(X_m,Y_m)$，而各集合
$\operatorname{Hom}_{S_n}(X_n,Y_n)$ 关于 $\varphi_{mn}$ 构成一个
\emph{集合的逆系统}。因此，\hyperref[I.10.6.11-fr]{10.6.11} 也可表述为：
存在典范双射
\[
  \operatorname{Hom}_{\mathfrak{S}}(\mathfrak{X},\mathfrak{Y})
  \xrightarrow{\sim}
  \varprojlim_n\operatorname{Hom}_{S_n}(X_n,Y_n).
\]

\subsection{形式预概形的积}
\label{subsection:I.10.7-fr}

\begin{env}[10.7.1]
\label{I.10.7.1-fr}
设 $\mathfrak{S}$ 是一个形式预概形；形式 $\mathfrak{S}$-预概形构成一个
范畴，因而可以定义形式 $\mathfrak{S}$-预概形的\emph{积}概念。
\end{env}

\begin{proposition}[10.7.2]
\label{I.10.7.2-fr}
设 $\mathfrak{X}=\operatorname{Spf}(B)$、
$\mathfrak{Y}=\operatorname{Spf}(C)$ 是仿射形式概形
$\mathfrak{S}=\operatorname{Spf}(A)$ 之上的两个仿射形式概形。令
$\mathfrak{Z}=\operatorname{Spf}(B\widehat{\otimes}_A C)$，并令
$p_1$、$p_2$ 是对应于（\hyperref[I.10.2.2-fr]{10.2.2}）从 $B$ 和 $C$
到 $B\widehat{\otimes}_A C$ 的典范（连续）$A$-同态 $\rho$ 和 $\sigma$
的 $\mathfrak{S}$-态射；则 $(\mathfrak{Z},p_1,p_2)$ 是仿射形式
$\mathfrak{S}$-概形 $\mathfrak{X}$ 和 $\mathfrak{Y}$ 的一个积。
\end{proposition}

\begin{proof}
由 \hyperref[I.10.4.6-fr]{10.4.6}，只须验证：若对每个连续 $A$-同态
$\varphi:B\widehat{\otimes}_A C\to D$，其中 $D$ 是一个作为拓扑
$A$-代数的可容环，都配以有序对
$(\varphi\circ\rho,\varphi\circ\sigma)$，那么便定义了一个双射
\[
  \operatorname{Hom}_A(B\widehat{\otimes}_A C,D)
  \xrightarrow{\sim}
  \operatorname{Hom}_A(B,D)\times\operatorname{Hom}_A(C,D)
\]
而这正是完备张量积的泛性质
（\hyperref[0.7.7.6-fr]{0, 7.7.6}）。
\end{proof}

\begin{proposition}[10.7.3]
\label{I.10.7.3-fr}
给定两个形式 $\mathfrak{S}$-预概形 $\mathfrak{X}$、$\mathfrak{Y}$，
积 $\mathfrak{X}\times_{\mathfrak{S}}\mathfrak{Y}$ 存在。
\end{proposition}

\begin{proof}
证明与 \hyperref[I.3.2.6-fr]{3.2.6} 的证明相同，只须将其中的仿射概形
（相应地，仿射开集）换成仿射形式概形（相应地，仿射形式开集），
并将命题 \hyperref[I.3.2.2-fr]{3.2.2} 换成
\hyperref[I.10.7.2-fr]{10.7.2}。
\end{proof}

预概形的积所具有的一切形式性质（\hyperref[I.3.2.7-fr]{3.2.7} 和
\hyperref[I.3.2.8-fr]{3.2.8}，以及
\hyperref[I.3.3.1-fr]{3.3.1} 至
\hyperref[I.3.3.12-fr]{3.3.12}）对于形式预概形的积都原封不动地成立。

\begin{env}[10.7.4]
\label{I.10.7.4-fr}
设 $\mathfrak{S}$、$\mathfrak{X}$、$\mathfrak{Y}$ 是三个形式预概形，
而 $f:\mathfrak{X}\to\mathfrak{S}$、
$g:\mathfrak{Y}\to\mathfrak{S}$ 是两个态射。假设在
$\mathfrak{S}$、$\mathfrak{X}$、$\mathfrak{Y}$ 中分别存在三个具有同一
指标集 $I$ 的定义理想的基本系
$(\mathscr{J}_\lambda)$、$(\mathscr{K}_\lambda)$、
$(\mathscr{L}_\lambda)$，使得对于每个 $\lambda$，都有
$f^*(\mathscr{J}_\lambda)\mathscr{O}_{\mathfrak{X}}
\subset\mathscr{K}_\lambda$ 和
$g^*(\mathscr{J}_\lambda)\mathscr{O}_{\mathfrak{Y}}
\subset\mathscr{L}_\lambda$。置
$S_\lambda=(\mathfrak{S},
\mathscr{O}_{\mathfrak{S}}/\mathscr{J}_\lambda)$，
$X_\lambda=(\mathfrak{X},
\mathscr{O}_{\mathfrak{X}}/\mathscr{K}_\lambda)$，
$Y_\lambda=(\mathfrak{Y},
\mathscr{O}_{\mathfrak{Y}}/\mathscr{L}_\lambda)$；当
$\mathscr{J}_\mu\subset\mathscr{J}_\lambda$、
$\mathscr{K}_\mu\subset\mathscr{K}_\lambda$、
$\mathscr{L}_\mu\subset\mathscr{L}_\lambda$ 时，注意到
$S_\lambda$（相应地，$X_\lambda$、$Y_\lambda$）是
$S_\mu$（相应地，$X_\mu$、$Y_\mu$）的一个具有相同

\oldpage[I]{194}

底空间的闭子预概形（\hyperref[I.10.6.1-fr]{10.6.1}）。由于
$S_\lambda\to S_\mu$ 是预概形的单态射，先可看出积
$X_\lambda\times_{S_\lambda}Y_\lambda$ 与
$X_\lambda\times_{S_\mu}Y_\lambda$ 相同
（\hyperref[I.3.2.4-fr]{3.2.4}），继而可看出
$X_\lambda\times_{S_\mu}Y_\lambda$ 等同于
$X_\mu\times_{S_\mu}Y_\mu$ 的一个具有相同底空间的闭子预概形
（\hyperref[I.4.3.1-fr]{4.3.1}）。在此情形下，积
$\mathfrak{X}\times_{\mathfrak{S}}\mathfrak{Y}$ 是通常预概形
$X_\lambda\times_{S_\lambda}Y_\lambda$ 的\emph{归纳极限}：事实上，
与 \hyperref[I.10.6.2-fr]{10.6.2} 中一样，可将问题归结到
$\mathfrak{S}$、$\mathfrak{X}$ 和 $\mathfrak{Y}$ 都是仿射形式概形的
情形。考虑到 \hyperref[I.10.5.6-fr]{10.5.6, (ii)} 以及关于
$\mathfrak{S}$、$\mathfrak{X}$ 和 $\mathfrak{Y}$ 的定义理想基本系的
假设，便立即看出，我们的断言由两个代数的完备张量积的定义
（\hyperref[0.7.7.1-fr]{0, 7.7.1}）得出。

此外，设 $\mathfrak{Z}$ 是形式 $\mathfrak{S}$-预概形，
$(\mathscr{M}_\lambda)$ 是 $\mathfrak{Z}$ 的以 $I$ 为指标集的定义理想
基本系，$u:\mathfrak{Z}\to\mathfrak{X}$、
$v:\mathfrak{Z}\to\mathfrak{Y}$ 是两个 $\mathfrak{S}$-态射，满足
$u^*(\mathscr{K}_\lambda)\mathscr{O}_{\mathfrak{Z}}
\subset\mathscr{M}_\lambda$ 和
$v^*(\mathscr{L}_\lambda)\mathscr{O}_{\mathfrak{Z}}
\subset\mathscr{M}_\lambda$。若置
$Z_\lambda=(\mathfrak{Z},
\mathscr{O}_{\mathfrak{Z}}/\mathscr{M}_\lambda)$，并且
$u_\lambda:Z_\lambda\to X_\lambda$ 和
$v_\lambda:Z_\lambda\to Y_\lambda$ 是对应于 $u$ 和 $v$ 的
$S_\lambda$-态射（\hyperref[I.10.5.6-fr]{10.5.6}），则立即可验证，
$(u,v)_{\mathfrak{S}}$ 是各 $S_\lambda$-态射
$(u_\lambda,v_\lambda)_{S_\lambda}$ 的归纳极限。

本条的讨论尤其适用于 $\mathfrak{S}$、$\mathfrak{X}$ 和
$\mathfrak{Y}$ 都是局部 Noether 的情形；此时可取一个定义理想的各次幂
所成的系统（\hyperref[I.10.5.1-fr]{10.5.1}）作为定义理想的基本系。
但应当注意，$\mathfrak{X}\times_{\mathfrak{S}}\mathfrak{Y}$ 不一定是
局部 Noether 的（不过参见 \hyperref[I.10.13.5-fr]{10.13.5}）。
\end{env}

\subsection{概形沿着一个闭子集的形式完备化}
\label{subsection:I.10.8-fr}

\begin{env}[10.8.1]
\label{I.10.8.1-fr}
设 $X$ 是一个（通常的）局部 Noether 预概形，$X'$ 是 $X$ 的底空间的一个
闭子集；用 $\Phi$ 表示 $\mathscr{O}_X$ 中所有满足
$\mathscr{O}_X/\mathscr{J}$ 的支集为 $X'$ 的凝聚理想层
$\mathscr{J}$ 所成的集合。集合 $\Phi$ 非空
（\hyperref[I.5.2.1-fr]{5.2.1}、
\hyperref[I.4.1.4-fr]{4.1.4} 和
\hyperref[I.6.1.1-fr]{6.1.1}）；我们用关系 $\supset$ 对它排序。
\end{env}

\begin{lemma}[10.8.2]
\label{I.10.8.2-fr}
有序集 $\Phi$ 是滤过的；若 $X$ 是 Noether 的，则对于每个
$\mathscr{J}_0\in\Phi$，各幂 $\mathscr{J}_0^n$（$n>0$）所成的集合在
$\Phi$ 中共尾。
\end{lemma}

\begin{proof}
事实上，若 $\mathscr{J}_1$ 和 $\mathscr{J}_2$ 属于 $\Phi$，并置
$\mathscr{J}=\mathscr{J}_1\cap\mathscr{J}_2$，则因
$\mathscr{O}_X$ 是凝聚的，$\mathscr{J}$ 也是凝聚的
（\hyperref[I.6.1.1-fr]{6.1.1} 和
\hyperref[0.5.3.4-fr]{0, 5.3.4}）；并且对于每个 $x\in X$，都有
$\mathscr{J}_x=(\mathscr{J}_1)_x\cap(\mathscr{J}_2)_x$。因此，当
$x\notin X'$ 时 $\mathscr{J}_x=\mathscr{O}_x$，而当 $x\in X'$ 时
$\mathscr{J}_x\neq\mathscr{O}_x$，这就证明了 $\mathscr{J}\in\Phi$。
另一方面，若 $X$ 是 Noether 的，且 $\mathscr{J}_0$ 和
$\mathscr{J}$ 属于 $\Phi$，则存在整数 $n>0$，使得
$\mathscr{J}_0^n(\mathscr{O}_X/\mathscr{J})=0$
（\hyperref[I.9.3.4-fr]{9.3.4}），这就是说
$\mathscr{J}_0^n\subset\mathscr{J}$。
\end{proof}

\begin{env}[10.8.3]
\label{I.10.8.3-fr}
现在设 $\mathscr{F}$ 是一个凝聚 $\mathscr{O}_X$-模；对于每个
$\mathscr{J}\in\Phi$，
$\mathscr{F}\otimes_{\mathscr{O}_X}(\mathscr{O}_X/\mathscr{J})$
都是一个凝聚 $\mathscr{O}_X$-模
（\hyperref[I.9.1.1-fr]{9.1.1}），其支集包含在 $X'$ 中；我们通常把它
与其在 $X'$ 上的限制等同起来。当 $\mathscr{J}$ 遍历 $\Phi$ 时，
这些层构成一个阿贝尔群层的\emph{逆系统}。
\end{env}

\begin{definition}[10.8.4]
\label{I.10.8.4-fr}
给定局部 Noether 预概形 $X$ 的一个闭子集 $X'$ 以及一个凝聚
$\mathscr{O}_X$-模 $\mathscr{F}$，把层

\oldpage[I]{195}

$\varprojlim_{\Phi}
(\mathscr{F}\otimes_{\mathscr{O}_X}
(\mathscr{O}_X/\mathscr{J}))$ 在 $X'$ 上的限制称为
\emph{$\mathscr{F}$ 沿 $X'$ 的完备化}，并记作
$\mathscr{F}_{/X'}$，或者在不致混淆时记作
$\widehat{\mathscr{F}}$；称它在 $X'$ 上的截面为
\emph{$\mathscr{F}$ 沿 $X'$ 的形式截面}。
\end{definition}

立即可见，对于每个开集 $U\subset X$，都有
$(\mathscr{F}|U)_{/(U\cap X')}=(\mathscr{F}_{/X'})|(U\cap X')$。

取逆极限立即可见，$(\mathscr{O}_X)_{/X'}$ 是一个环层，而
$\mathscr{F}_{/X'}$ 可视为一个 $(\mathscr{O}_X)_{/X'}$-模。此外，
由于 $X'$ 的拓扑有一个由拟紧开集组成的基，可以把
$(\mathscr{O}_X)_{/X'}$（相应地，$\mathscr{F}_{/X'}$）视为一个
\emph{拓扑环层}（相应地，\emph{拓扑群层}）：它是各\emph{伪离散}环层
（相应地，群层）$\mathscr{O}_X/\mathscr{J}$（相应地，
$\mathscr{F}\otimes_{\mathscr{O}_X}
(\mathscr{O}_X/\mathscr{J})=\mathscr{F}/\mathscr{J}\mathscr{F}$）的
逆极限。再取逆极限，$\mathscr{F}_{/X'}$ 于是成为一个
\emph{拓扑 $(\mathscr{O}_X)_{/X'}$-模}
（\hyperref[0.3.8.1-fr]{0, 3.8.1} 和
\hyperref[0.3.8.2-fr]{3.8.2}）。回忆，对于每个\emph{拟紧}开集
$U\subset X$，$\Gamma(U\cap X',(\mathscr{O}_X)_{/X'})$
（相应地，$\Gamma(U\cap X',\mathscr{F}_{/X'})$）是离散环
（相应地，离散群）$\Gamma(U,\mathscr{O}_X/\mathscr{J})$
（相应地，$\Gamma(U,\mathscr{F}/\mathscr{J}\mathscr{F})$）的逆极限。

现在，若 $u:\mathscr{F}\to\mathscr{G}$ 是一个
$\mathscr{O}_X$-模同态，则对于每个 $\mathscr{J}\in\Phi$，都典范地得到
同态
$u_{\mathscr{J}}:
\mathscr{F}\otimes_{\mathscr{O}_X}(\mathscr{O}_X/\mathscr{J})
\to
\mathscr{G}\otimes_{\mathscr{O}_X}(\mathscr{O}_X/\mathscr{J})$，
并且这些同态构成一个逆系统。取逆极限并限制到 $X'$，它们给出一个连续
$(\mathscr{O}_X)_{/X'}$-同态
$\mathscr{F}_{/X'}\to\mathscr{G}_{/X'}$，记作 $u_{/X'}$ 或
$\widehat u$，并称作同态 $u$ 沿 $X'$ 的\emph{完备化}。显然，若
$v:\mathscr{G}\to\mathscr{H}$ 是另一个 $\mathscr{O}_X$-模同态，则有
$(v\circ u)_{/X'}=(v_{/X'})\circ(u_{/X'})$；因此，
$\mathscr{F}_{/X'}$ 是关于 $\mathscr{F}$ 的一个\emph{协变加性函子}，
它从凝聚 $\mathscr{O}_X$-模范畴取值于拓扑
$(\mathscr{O}_X)_{/X'}$-模范畴。

\begin{proposition}[10.8.5]
\label{I.10.8.5-fr}
$(\mathscr{O}_X)_{/X'}$ 的支集是 $X'$；拓扑赋环空间
$(X',(\mathscr{O}_X)_{/X'})$ 是一个局部 Noether 形式预概形，并且若
$\mathscr{J}\in\Phi$，则 $\mathscr{J}_{/X'}$ 是这个形式预概形的一个
定义理想。若 $X=\operatorname{Spec}(A)$ 是一个 Noether 环的仿射概形，
$\mathscr{J}=\widetilde{\mathfrak{J}}$，其中
$\mathfrak{J}$ 是 $A$ 的一个理想，并且
$X'=V(\mathfrak{J})$，则 $(X',(\mathscr{O}_X)_{/X'})$ 典范地等同于
$\operatorname{Spf}(\widehat A)$，其中 $\widehat A$ 是 $A$ 关于
$\mathfrak{J}$-预进制拓扑的分离完备化。
\end{proposition}

\begin{proof}
显然只须证明最后一个断言。我们知道
（\hyperref[0.7.3.3-fr]{0, 7.3.3}），$\mathfrak{J}$ 关于
$\mathfrak{J}$-预进制拓扑的分离完备化 $\widehat{\mathfrak{J}}$
等同于 $\widehat A$ 的理想 $\mathfrak{J}\widehat A$，而
$\widehat A$ 是一个 Noether 的 $\widehat{\mathfrak{J}}$-进制环，满足
$\widehat A/\widehat{\mathfrak{J}}^n=A/\mathfrak{J}^n$
（\hyperref[0.7.2.6-fr]{0, 7.2.6}）。后一关系表明，$\widehat A$ 的
开素理想就是各理想
$\widehat{\mathfrak p}=\mathfrak p\widehat A$，其中 $\mathfrak p$ 是
$A$ 中包含 $\mathfrak{J}$ 的素理想；并且
$\widehat{\mathfrak p}\cap A=\mathfrak p$，所以
$\operatorname{Spf}(\widehat A)=X'$。由于
$\mathscr{O}_X/\mathscr{J}^n=(A/\mathfrak{J}^n)^{\sim}$，命题立即由定义得出。
\end{proof}

称这样定义的形式预概形为\emph{$X$ 沿 $X'$ 的完备化}，并记作
$X_{/X'}$，或者在不致混淆时记作 $\widehat X$。当取 $X'=X$ 时，
可取 $\mathscr{J}=0$，因而 $X_{/X}=X$。

显然，若 $U$ 是 $X$ 的一个开集上所诱导的子预概形，则
$U_{/(U\cap X')}$ 典范地等同于 $X_{/X'}$ 在 $X'$ 的开集
$U\cap X'$ 上所诱导的形式子预概形。

\begin{corollary}[10.8.6]
\label{I.10.8.6-fr}
（通常的）预概形 $\widehat X_{\mathrm{red}}$ 是 $X$ 中以 $X'$ 为底空间的
唯一约化子预概形（\hyperref[I.5.2.1-fr]{5.2.1}）。
$\widehat X$ 为 Noether 的充要条件是 $\widehat X_{\mathrm{red}}$ 为
Noether 的；而 $X$ 为 Noether 的则是一个充分条件。
\end{corollary}

\oldpage[I]{196}

$\widehat X_{\mathrm{red}}$ 的确定是局部的
（\hyperref[I.10.5.4-fr]{10.5.4}），因此仍可假设 $X$ 是一个 Noether
环的仿射概形；沿用
\hyperref[I.10.8.5-fr]{10.8.5} 的记号，$\widehat A$ 的拓扑幂零元所成的
理想 $\mathfrak{T}$，是 $A/\mathfrak{J}$ 的幂零根在典范映射
$\widehat A\to\widehat A/\widehat{\mathfrak{J}}=A/\mathfrak{J}$ 下的
逆像（\hyperref[0.7.1.3-fr]{0, 7.1.3}），所以
$\widehat A/\mathfrak{T}$ 同构于 $A/\mathfrak{J}$ 对其幂零根的商。
因此第一个断言由 \hyperref[I.10.5.4-fr]{10.5.4} 和
\hyperref[I.5.1.1-fr]{5.1.1} 得出。若
$\widehat X_{\mathrm{red}}$ 是 Noether 的，则其底空间 $X'$ 也是
Noether 的，从而
$X'_n=\operatorname{Spec}(\mathscr{O}_X/\mathscr{J}^n)$ 都是 Noether 的
（\hyperref[I.6.1.2-fr]{6.1.2}），于是 $\widehat X$ 也是 Noether 的
（\hyperref[I.10.6.4-fr]{10.6.4}）；反命题由
\hyperref[I.6.1.2-fr]{6.1.2} 立即得出。

\begin{env}[10.8.7]
\label{I.10.8.7-fr}
各典范同态 $\mathscr{O}_X\to\mathscr{O}_X/\mathscr{J}$
（$\mathscr{J}\in\Phi$）构成一个逆系统，因而取逆极限便给出环层同态
$\theta:\mathscr{O}_X\to\psi_*((\mathscr{O}_X)_{/X'})
=\varprojlim_\Phi(\mathscr{O}_X/\mathscr{J})$；这里 $\psi$ 表示底空间的
典范嵌入 $X'\to X$。用 $i$（或 $i_X$）表示赋环空间的态射
（称为典范态射）
\[
  (\psi,\theta):X_{/X'}\to X
\]

对于每个凝聚 $\mathscr{O}_X$-模 $\mathscr{F}$，作张量积后，各典范同态
$\mathscr{O}_X\to\mathscr{O}_X/\mathscr{J}$ 给出
$\mathscr{O}_X$-模同态
$\mathscr{F}\to
\mathscr{F}\otimes_{\mathscr{O}_X}(\mathscr{O}_X/\mathscr{J})$；
这些同态仍构成一个逆系统，所以取逆极限便给出典范函子性
$\mathscr{O}_X$-模同态
$\gamma:\mathscr{F}\to\psi_*(\mathscr{F}_{/X'})$。
\end{env}

\begin{proposition}[10.8.8]
\label{I.10.8.8-fr}
\medskip\noindent
\begin{enumerate}
  \item[\rm{(i)}] 关于 $\mathscr{F}$ 的函子 $\mathscr{F}_{/X'}$ 是正合的。
  \item[\rm{(ii)}] $(\mathscr{O}_X)_{/X'}$-模的函子性同态
    $\gamma^\sharp:i^*(\mathscr{F})\to\mathscr{F}_{/X'}$ 是同构。
\end{enumerate}
\end{proposition}

\medskip\noindent
\begin{enumerate}
  \item[\rm{(i)}] 只须证明：若
    $0\to\mathscr{F}'\to\mathscr{F}\to\mathscr{F}''\to0$ 是凝聚
    $\mathscr{O}_X$-模的一个正合列，而 $U$ 是 $X$ 的一个仿射开集，
    其环 $A$ 是 Noether 的，则序列
    \[
      0\to\Gamma(U\cap X',\mathscr{F}'_{/X'})
      \to\Gamma(U\cap X',\mathscr{F}_{/X'})
      \to\Gamma(U\cap X',\mathscr{F}''_{/X'})\to0
    \]
    是正合的。此时
    $\mathscr{F}|U=\widetilde M$、
    $\mathscr{F}'|U=\widetilde{M'}$、
    $\mathscr{F}''|U=\widetilde{M''}$，其中 $M$、$M'$、$M''$ 是三个
    有限型 $A$-模，并且序列 $0\to M'\to M\to M''\to0$ 正合
    （\hyperref[I.1.5.1-fr]{1.5.1} 和
    \hyperref[I.1.3.11-fr]{1.3.11}）。取 $\mathscr{J}\in\Phi$，并取
    $A$ 的理想 $\mathfrak{J}$，使得
    $\mathscr{J}|U=\widetilde{\mathfrak{J}}$。于是
    \[
      \Gamma(U\cap X',
      \mathscr{F}\otimes_{\mathscr{O}_X}(\mathscr{O}_X/\mathscr{J}^n))
      =M\otimes_A(A/\mathfrak{J}^n)
    \]
    （\hyperref[I.1.3.12-fr]{1.3.12}）；因此按逆极限的定义，
    \[
      \Gamma(U\cap X',\mathscr{F}_{/X'})
      =\varprojlim_n(M\otimes_A(A/\mathfrak{J}^n))=\widehat M
    \]
    即 $M$ 关于 $\mathfrak{J}$-预进制拓扑的分离完备化；类似地，
    \[
      \Gamma(U\cap X',\mathscr{F}'_{/X'})=\widehat{M'},\qquad
      \Gamma(U\cap X',\mathscr{F}''_{/X'})=\widehat{M''}~;
    \]
    当 $A$ 是 Noether 的时，关于 $M$ 的函子 $\widehat M$ 在有限型
    $A$-模范畴上是正合的
    （\hyperref[0.7.3.3-fr]{0, 7.3.3}），我们的断言因而得证。

\oldpage[I]{197}
  \item[\rm{(ii)}] 问题是局部的，故可假设有正合列
    $\mathscr{O}_X^m\to\mathscr{O}_X^n\to\mathscr{F}\to0$
    （\hyperref[0.5.3.2-fr]{0, 5.3.2}）。由于 $\gamma^\sharp$ 具有函子性，
    而函子 $i^*(\mathscr{F})$ 和 $\mathscr{F}_{/X'}$ 都是右正合的
    （由 (i) 和 \hyperref[0.4.3.1-fr]{0, 4.3.1}），有交换图表
    \[
      \xymatrix{
        i^*(\mathscr{O}_X^m)\ar[r]\ar[d]_{\gamma^\sharp} &
        i^*(\mathscr{O}_X^n)\ar[r]\ar[d]_{\gamma^\sharp} &
        i^*(\mathscr{F})\ar[r]\ar[d]_{\gamma^\sharp} &
        0\\
        (\mathscr{O}_X^m)_{/X'}\ar[r] &
        (\mathscr{O}_X^n)_{/X'}\ar[r] &
        \mathscr{F}_{/X'}\ar[r] &
        0
      }
      \tag{10.8.8.1}\label{I.10.8.8.1-fr}
    \]
    其两行都是正合的。此外，两个函子 $i^*(\mathscr{F})$ 和
    $\mathscr{F}_{/X'}$ 都与有限直和可交换
    （\hyperref[0.3.2.6-fr]{0, 3.2.6} 和
    \hyperref[0.4.3.2-fr]{4.3.2}），因此归结为对
    $\mathscr{F}=\mathscr{O}_X$ 证明断言。这时有
    $i^*(\mathscr{O}_X)=(\mathscr{O}_X)_{/X'}=
    \mathscr{O}_{\widehat X}$
    （\hyperref[0.4.3.4-fr]{0, 4.3.4}），而 $\gamma^\sharp$ 是一个
    $\mathscr{O}_{\widehat X}$-模同态；所以只须验证
    $\gamma^\sharp$ 把 $\mathscr{O}_{\widehat X}$ 在 $X'$ 的一个开集上的
    单位截面送到它自身。这是立即可见的，并且表明在此情形
    $\gamma^\sharp$ 就是恒等映射。
\end{enumerate}

\begin{corollary}[10.8.9]
\label{I.10.8.9-fr}
赋环空间的态射 $i:X_{/X'}\to X$ 是平坦的。
\end{corollary}

这事实上由 \hyperref[0.6.7.3-fr]{0, 6.7.3} 和
\hyperref[I.10.8.8-fr]{10.8.8, (i)} 得出。

\begin{corollary}[10.8.10]
\label{I.10.8.10-fr}
若 $\mathscr{F}$ 和 $\mathscr{G}$ 是凝聚 $\mathscr{O}_X$-模，则存在
关于 $\mathscr{F}$ 和 $\mathscr{G}$ 的典范函子性同构
\[
  (\mathscr{F}_{/X'})\otimes_{(\mathscr{O}_X)_{/X'}}
  (\mathscr{G}_{/X'})
  \xrightarrow{\sim}
  (\mathscr{F}\otimes_{\mathscr{O}_X}\mathscr{G})_{/X'}
  \tag{10.8.10.1}\label{I.10.8.10.1-fr}
\]
\[
  (\mathscr{H}\!om_{\mathscr{O}_X}(\mathscr{F},\mathscr{G}))_{/X'}
  \xrightarrow{\sim}
  \mathscr{H}\!om_{(\mathscr{O}_X)_{/X'}}
  (\mathscr{F}_{/X'},\mathscr{G}_{/X'})
  \tag{10.8.10.2}\label{I.10.8.10.2-fr}
\]
\end{corollary}

这由 $i^*(\mathscr{F})$ 与 $\mathscr{F}_{/X'}$ 的典范等同得出：第一个
同构的存在性是对任意赋环空间态射都成立的结果
（\hyperref[0.4.3.3.1-fr]{0, 4.3.3.1}）；第二个同构的存在性是对任意
平坦态射都成立的结果（\hyperref[0.6.7.6-fr]{0, 6.7.6}），因而由
\hyperref[I.10.8.9-fr]{10.8.9} 得出。

\begin{proposition}[10.8.11]
\label{I.10.8.11-fr}
对于任意凝聚 $\mathscr{O}_X$-模 $\mathscr{F}$，由
$\mathscr{F}\to\mathscr{F}_{/X'}$ 导出的典范同态
$\Gamma(X,\mathscr{F})\to\Gamma(X',\mathscr{F}_{/X'})$ 的核，恰由
在 $X'$ 的某个邻域上为零的截面组成。
\end{proposition}

由 $\mathscr{F}_{/X'}$ 的定义，这样的截面的典范像为零。反过来，若
$s\in\Gamma(X,\mathscr{F})$ 在
$\Gamma(X',\mathscr{F}_{/X'})$ 中的像为零，则只须证明每个
$x\in X'$ 都有 $X$ 中的一个邻域，使 $s$ 在其中为零；于是可归结到以下
情形：$X=\operatorname{Spec}(A)$ 是仿射的，$A$ 是 Noether 的，
$X'=V(\mathfrak{J})$，其中 $\mathfrak{J}$ 是 $A$ 的一个理想，并且
$\mathscr{F}=\widetilde M$，其中 $M$ 是有限型 $A$-模。此时
$\Gamma(X',\mathscr{F}_{/X'})$ 是 $M$ 关于 $\mathfrak{J}$-预进制拓扑的
分离完备化 $\widehat M$，而同态
$\Gamma(X,\mathscr{F})\to\Gamma(X',\mathscr{F}_{/X'})$ 就是典范同态
$M\to\widehat M$。我们知道
（\hyperref[0.7.3.7-fr]{0, 7.3.7}），此同态的核是所有被
$1+\mathfrak{J}$ 中某个元素零化的 $z\in M$ 所成的集合。因此，对某个
$f\in\mathfrak{J}$ 有 $(1+f)s=0$；对于每个 $x\in X'$，由此得到
$(1_x+f_x)s_x=0$。由于 $1_x+f_x$ 在 $\mathscr{O}_x$ 中可逆
（$\mathfrak{J}_x\mathscr{O}_x$ 包含于 $\mathscr{O}_x$ 的极大理想），
故 $s_x=0$，命题得证。

\begin{corollary}[10.8.12]
\label{I.10.8.12-fr}
$\mathscr{F}_{/X'}$ 的支集等于
$\operatorname{Supp}(\mathscr{F})\cap X'$。
\end{corollary}

显然，$\mathscr{F}_{/X'}$ 是一个有限型 $(\mathscr{O}_X)_{/X'}$-模
（\hyperref[I.10.8.8-fr]{10.8.8, (ii)} 和
\hyperref[0.5.2.4-fr]{0, 5.2.4}），

\oldpage[I]{198}

所以其支集是闭的（\hyperref[0.5.2.2-fr]{0, 5.2.2}），并且显然包含于
$\operatorname{Supp}(\mathscr{F})\cap X'$。为证明它等于后一集合，立即
归结为证明关系 $\Gamma(X',\mathscr{F}_{/X'})=0$ 蕴涵
$\operatorname{Supp}(\mathscr{F})\cap X'=\varnothing$；这由
\hyperref[I.10.8.11-fr]{10.8.11} 和
\hyperref[I.1.4.1-fr]{1.4.1} 得出。

\begin{corollary}[10.8.13]
\label{I.10.8.13-fr}
设 $u:\mathscr{F}\to\mathscr{G}$ 是凝聚 $\mathscr{O}_X$-模的一个同态。
$u_{/X'}:\mathscr{F}_{/X'}\to\mathscr{G}_{/X'}$ 为零的充要条件是 $u$
在 $X'$ 的某个邻域上为零。
\end{corollary}

事实上，由 \hyperref[I.10.8.8-fr]{10.8.8, (ii)}，$u_{/X'}$ 与
$i^*(u)$ 等同；因此，若把 $u$ 看成层
$\mathscr{H}=\mathscr{H}\!om_{\mathscr{O}_X}(\mathscr{F},\mathscr{G})$
在 $X$ 上的一个截面，则 $u_{/X'}$ 就是在 $X'$ 上同它典范对应的
$i^*(\mathscr{H})=\mathscr{H}_{/X'}$ 的截面
（\hyperref[I.10.8.10.2-fr]{10.8.10.2} 和
\hyperref[0.4.4.6-fr]{0, 4.4.6}）。因此，只须对凝聚
$\mathscr{O}_X$-模 $\mathscr{H}$ 应用
\hyperref[I.10.8.11-fr]{10.8.11}。

\begin{corollary}[10.8.14]
\label{I.10.8.14-fr}
设 $u:\mathscr{F}\to\mathscr{G}$ 是凝聚 $\mathscr{O}_X$-模的一个同态。
$u_{/X'}$ 为单态射（相应地，满态射）的充要条件是 $u$ 在 $X'$ 的某个
邻域上为单态射（相应地，满态射）。
\end{corollary}

设 $\mathscr{P}$ 和 $\mathscr{N}$ 分别为 $u$ 的余核与核，因而有正合列
\[
  0\to\mathscr{N}\xrightarrow{v}\mathscr{F}\xrightarrow{u}
  \mathscr{G}\xrightarrow{w}\mathscr{P}\to0,
\]
由此根据 \hyperref[I.10.8.8-fr]{10.8.8, (i)} 得到正合列
\[
  0\to\mathscr{N}_{/X'}\xrightarrow{v_{/X'}}
  \mathscr{F}_{/X'}\xrightarrow{u_{/X'}}
  \mathscr{G}_{/X'}\xrightarrow{w_{/X'}}
  \mathscr{P}_{/X'}\to0.
\]
若 $u_{/X'}$ 是单态射（相应地，满态射），则
$v_{/X'}=0$（相应地，$w_{/X'}=0$），所以由
\hyperref[I.10.8.13-fr]{10.8.13}，存在 $X'$ 的一个邻域，使
$v=0$（相应地，$w=0$）。

\subsection{把态射延拓到完备化上}
\label{subsection:I.10.9-fr}

\begin{env}[10.9.1]
\label{I.10.9.1-fr}
设 $X$、$Y$ 是两个（通常的）局部 Noether 预概形，$f:X\to Y$ 是一个
态射，$X'$（相应地，$Y'$）是底空间 $X$（相应地，$Y$）的一个闭子集，
并且 $f(X')\subset Y'$。设 $\mathscr{J}$（相应地，$\mathscr{K}$）是
$\mathscr{O}_X$（相应地，$\mathscr{O}_Y$）的一个理想层，使
$\mathscr{O}_X/\mathscr{J}$（相应地，
$\mathscr{O}_Y/\mathscr{K}$）的支集为 $X'$（相应地，$Y'$），并且
$f^*(\mathscr{K})\mathscr{O}_X\subset\mathscr{J}$。注意，总存在这样的
理想层：例如，可取 $\mathscr{J}$ 为 $\mathscr{O}_X$ 中定义以 $X'$ 为
底空间的 $X$ 的子预概形的最大理想层
（\hyperref[I.5.2.1-fr]{5.2.1}）；此时假设 $f(X')\subset Y'$ 蕴涵
$f^*(\mathscr{K})\mathscr{O}_X\subset\mathscr{J}$
（\hyperref[I.5.2.4-fr]{5.2.4}）。因此，对于每个整数 $n>0$，都有
$f^*(\mathscr{K}^n)\mathscr{O}_X\subset\mathscr{J}^n$
（\hyperref[0.4.3.5-fr]{0, 4.3.5}）；所以由
\hyperref[I.4.4.6-fr]{4.4.6}，若置
\[
  X'_n=(X',\mathscr{O}_X/\mathscr{J}^{n+1}),\qquad
  Y'_n=(Y',\mathscr{O}_Y/\mathscr{K}^{n+1}),
\]
便由 $f$ 得到一个态射 $f_n:X'_n\to Y'_n$，并且各 $f_n$ 显然构成一个
归纳系。把它的归纳极限（\hyperref[I.10.6.8-fr]{10.6.8}）记作
$\widehat f:X_{/X'}\to Y_{/Y'}$，并且（滥用术语）称 $\widehat f$ 为
\emph{$f$ 到 $X$ 和 $Y$ 沿 $X'$ 和 $Y'$ 的完备化上的延拓}。立即可验证，
此态射与满足上述条件的理想层 $\mathscr{J}$、$\mathscr{K}$ 的选择无关。
事实上，只须在 $X$ 和 $Y$ 分别是 Noether 环 $A$、$B$ 的仿射概形时
验证这一点。此时
$\mathscr{J}=\widetilde{\mathfrak{J}}$、
$\mathscr{K}=\widetilde{\mathfrak{K}}$，其中 $\mathfrak{J}$（相应地，
$\mathfrak{K}$）是 $A$（相应地，$B$）的一个理想；$f$ 对应于满足
$\varphi(\mathfrak{K})\subset\mathfrak{J}$ 的环同态
$\varphi:B\to A$
（\hyperref[I.4.4.6-fr]{4.4.6} 和
\hyperref[I.1.7.4-fr]{1.7.4}）。此时，\emph{$\widehat f$ 是与连续同态
$\widehat\varphi:\widehat B\to\widehat A$ 相对应
（\hyperref[I.10.2.2-fr]{10.2.2}）的态射}；这里 $\widehat A$
（相应地，$\widehat B$）是 $A$（相应地，$B$）关于
$\mathfrak{J}$-预进制拓扑（相应地，$\mathfrak{K}$-预进制拓扑）的
分离完备化（\hyperref[I.10.6.8-fr]{10.6.8}）。我们还知道，若把
$\mathscr{J}$ 换成另一个

\oldpage[I]{199}

理想层 $\mathscr{J}'=\widetilde{\mathfrak{J}'}$，而
$\mathscr{O}_X/\mathscr{J}'$ 的支集仍为 $X'$，则 $A$ 上的
$\mathfrak{J}$-预进制拓扑和 $\mathfrak{J}'$-预进制拓扑相同
（\hyperref[I.10.8.2-fr]{10.8.2}）。

还要注意，按此定义，与 $\widehat f$ 对应的 $X_{/X'}$ 和 $Y_{/Y'}$
底空间之间的连续映射 $X'\to Y'$，正是 $f$ 在 $X'$ 上的限制。
\end{env}

\begin{env}[10.9.2]
\label{I.10.9.2-fr}
由前一定义立即可知，赋环空间态射的图表
\[
  \label{I.10.9.2.1-fr}
  \xymatrix{
    \widehat X\ar[r]^{\widehat f}\ar[d]_{i_X} &
    \widehat Y\ar[d]^{i_Y}\\
    X\ar[r]_f &
    Y
  }
\]
交换，其中竖直箭头为典范态射
（\hyperref[I.10.8.7-fr]{10.8.7}）。
\end{env}

\begin{env}[10.9.3]
\label{I.10.9.3-fr}
设 $Z$ 是第三个预概形，$g:Y\to Z$ 是一个态射，$Z'$ 是 $Z$ 的一个闭
子集，并且 $g(Y')\subset Z'$。若 $\widehat g$ 表示态射 $g$ 沿
$Y'$ 和 $Z'$ 到完备化上的延拓，则由
\hyperref[I.10.9.1-fr]{10.9.1} 立即得到
$\widehat{(g\circ f)}=\widehat g\circ\widehat f$。
\end{env}

\begin{proposition}[10.9.4]
\label{I.10.9.4-fr}
设 $X$、$Y$ 是两个局部 Noether $S$-预概形，且 $Y$ 在 $S$ 上为有限型。
设 $f$、$g$ 是 $X$ 到 $Y$ 的两个 $S$-态射，满足
$f(X')\subset Y'$、$g(X')\subset Y'$。则
$\widehat f=\widehat g$ 的充要条件是 $f$ 和 $g$ 在 $X'$ 的一个邻域上
重合。
\end{proposition}

此条件显然是充分的（无须对 $Y$ 作有限性假设）。为证明其必要性，首先
注意，假设 $\widehat f=\widehat g$ 蕴涵对于每个 $x\in X'$ 都有
$f(x)=g(x)$。另一方面，问题是局部的，故可假设 $X$、$Y$ 分别是 $x$
和 $y=f(x)=g(x)$ 的仿射开邻域，其环为 Noether 的；还可假设 $S$ 是
仿射的，并且 $\Gamma(Y,\mathscr{O}_Y)$ 是有限型
$\Gamma(S,\mathscr{O}_S)$-代数
（\hyperref[I.6.3.3-fr]{6.3.3}）。于是，$f$ 和 $g$ 分别对应于从
$\Gamma(Y,\mathscr{O}_Y)$ 到 $\Gamma(X,\mathscr{O}_X)$ 的两个
$\Gamma(S,\mathscr{O}_S)$-同态 $\rho$、$\sigma$
（\hyperref[I.1.7.3-fr]{1.7.3}），并且根据假设，这两个同态到
$\Gamma(Y,\mathscr{O}_Y)$ 的分离完备化上的连续延拓相同。由
\hyperref[I.10.8.11-fr]{10.8.11} 可知，对于每个截面
$s\in\Gamma(Y,\mathscr{O}_Y)$，截面 $\rho(s)$ 和 $\sigma(s)$ 在 $X'$
的某个邻域（依赖于 $s$）上重合；又因为
$\Gamma(Y,\mathscr{O}_Y)$ 是有限型 $\Gamma(S,\mathscr{O}_S)$-代数，
立即得知存在 $X'$ 的一个邻域 $V$，使得对于\emph{每个}截面
$s\in\Gamma(Y,\mathscr{O}_Y)$，$\rho(s)$ 和 $\sigma(s)$ 在 $V$ 上
重合。
% EGA1-ERR-0056: printed/diplomatic French has Gamma(Y,O_X); retained source-literally pending canon disposition.
若 $h\in\Gamma(Y,\mathscr{O}_X)$ 使 $D(h)$ 是包含于 $V$ 的 $x$ 的一个
邻域，则由上述结果以及 \hyperref[I.1.4.1-fr]{1.4.1, d)} 可知，$f$ 和
$g$ 在 $D(h)$ 上重合。

\begin{proposition}[10.9.5]
\label{I.10.9.5-fr}
在 \hyperref[I.10.9.1-fr]{10.9.1} 的假设下，对于每个凝聚
$\mathscr{O}_Y$-模 $\mathscr{G}$，存在典范的函子性
$(\mathscr{O}_X)_{/X'}$-模同构
\[
  (f^*(\mathscr{G}))_{/X'}\xrightarrow{\sim}
  \widehat f^*(\mathscr{G}_{/Y'})
\]
\end{proposition}

若把 $(f^*(\mathscr{G}))_{/X'}$ 典范地等同于
$i_X^*(f^*(\mathscr{G}))$，把 $\widehat f^*(\mathscr{G}_{/Y'})$ 典范地
等同于 $\widehat f^*(i_Y^*(\mathscr{G}))$
（\hyperref[I.10.8.8-fr]{10.8.8}），则命题立即由
\hyperref[I.10.9.2-fr]{10.9.2} 中图表的交换性得出。

\begin{env}[10.9.6]
\label{I.10.9.6-fr}
现设 $\mathscr{F}$ 是一个凝聚 $\mathscr{O}_X$-模，
$\mathscr{G}$ 是一个凝聚 $\mathscr{O}_Y$-模。若
$u:\mathscr{G}\to\mathscr{F}$ 是从 $\mathscr{G}$ 到 $\mathscr{F}$ 的一个
$f$-态射，则它对应于一个 $\mathscr{O}_X$-同态
$u^\sharp:f^*(\mathscr{G})\to\mathscr{F}$，因而经过完备化得到一个连续的
$(\mathscr{O}_X)_{/X'}$-同态
$(u^\sharp)_{/X'}:(f^*(\mathscr{G}))_{/X'}\to\mathscr{F}_{/X'}$；根据
\hyperref[I.10.9.5-fr]{10.9.5}，存在唯一的 $\widehat f$-态射
$v:\mathscr{G}_{/Y'}\to\mathscr{F}_{/X'}$，
\oldpage[I]{200}

使得 $v^\sharp=(u^\sharp)_{/X'}$。若把三元组
$(\mathscr{F},X,X')$（其中 $\mathscr{F}$ 是一个凝聚
$\mathscr{O}_X$-模，$X'$ 是 $X$ 的一个闭子集）视为一个
\emph{范畴}，并规定态射
$(\mathscr{F},X,X')\to(\mathscr{G},Y,Y')$ 由一个满足
$f(X')\subset Y'$ 的预概形态射 $f:X\to Y$ 以及一个
$f$-态射 $u:\mathscr{G}\to\mathscr{F}$ 组成，则可以说
$(X_{/X'},\mathscr{F}_{/X'})$ 是 $(\mathscr{F},X,X')$ 的一个
\emph{函子}，其取值于下述偶对 $(\mathfrak{Z},\mathscr{H})$ 所成的
范畴：$\mathfrak{Z}$ 是局部 Noether 形式预概形，
$\mathscr{H}$ 是一个 $\mathscr{O}_{\mathfrak{Z}}$-模；后一范畴的态射由形式
预概形的一个态射 $g$ 和一个 $g$-态射所成的偶对组成。
\end{env}

\begin{proposition}[10.9.7]
\label{I.10.9.7-fr}
设 $S$、$X$、$Y$ 是三个局部 Noether 预概形，
$g:X\to S$、$h:Y\to S$ 是两个态射，$S'$ 是 $S$ 的一个闭子集，
$X'$（相应地，$Y'$）是 $X$（相应地，$Y$）的一个闭子集，并且
$g(X')\subset S'$（相应地，$h(Y')\subset S'$）；置
$Z=X\times_S Y$，并假设 $Z$ 是局部 Noether 的；再置
$Z'=p^{-1}(X')\cap q^{-1}(Y')$，其中 $p$ 和 $q$ 是
$X\times_S Y$ 的投影。在这些条件下，完备化 $Z_{/Z'}$ 与
$S_{/S'}$-形式预概形的积
$(X_{/X'})\times_{S_{/S'}}(Y_{/Y'})$ 等同，其结构态射分别与
$\widehat g$ 和 $\widehat h$ 等同，而投影分别与
$\widehat p$ 和 $\widehat q$ 等同。
\end{proposition}

问题对 $S$、$X$ 和 $Y$ 是局部的，这是立即可见的；因此可归结到如下
情形：$S=\operatorname{Spec}(A)$、
$X=\operatorname{Spec}(B)$、$Y=\operatorname{Spec}(C)$、
$S'=V(\mathfrak{J})$、$X'=V(\mathfrak{K})$、$Y'=V(\mathfrak{L})$，其中
$\mathfrak{J}$、$\mathfrak{K}$、$\mathfrak{L}$ 是三个理想，满足
$\varphi(\mathfrak{J})\subset\mathfrak{K}$ 和
$\psi(\mathfrak{J})\subset\mathfrak{L}$；这里 $\varphi$ 和 $\psi$ 分别表示与
$g$ 和 $h$ 对应的同态 $A\to B$ 和 $A\to C$。此时我们知道
$Z=\operatorname{Spec}(B\otimes_A C)$ 且 $Z'=V(\mathfrak{M})$，其中
$\mathfrak{M}$ 是理想
$\operatorname{Im}(\mathfrak{K}\otimes_A C)
+\operatorname{Im}(B\otimes_A\mathfrak{L})$。结论源于
\hyperref[I.10.7.2-fr]{10.7.2} 以及以下事实：完备张量积
$(\widehat B\otimes_{\widehat A}\widehat C)^\wedge$（其中
$\widehat A$、$\widehat B$、$\widehat C$ 分别是 $A$、$B$、$C$ 对于
$\mathfrak{J}$-预进制、$\mathfrak{K}$-预进制和 $\mathfrak{L}$-预进制拓扑的
分离完备化），是张量积 $B\otimes_A C$ 对于
$\mathfrak{M}$-预进制拓扑的分离完备化
（\hyperref[0.7.7.2-fr]{0, 7.7.2}）。

还要注意，若 $T$ 是一个局部 Noether $S$-预概形，
$u:T\to X$、$v:T\to Y$ 是两个 $S$-态射，$T'$ 是 $T$ 的一个闭子集，并且
$u(T')\subset X'$、$v(T')\subset Y'$，则到完备化上的延拓
$((u,v)_S)^\wedge$ 与
$(\widehat u,\widehat v)_{S_{/S'}}$ 等同。

\begin{corollary}[10.9.8]
\label{I.10.9.8-fr}
设 $X$、$Y$ 是两个局部 Noether $S$-预概形，使得
$X\times_S Y$ 是局部 Noether 的；设 $S'$ 是 $S$ 的一个闭子集，
$X'$（相应地，$Y'$）是 $X$（相应地，$Y$）的一个闭子集，其在 $S$ 中的像
包含于 $S'$。对任意满足 $f(X')\subset Y'$ 的 $S$-态射
$f:X\to Y$，图态射 $\Gamma_{\widehat f}$ 与 $f$ 的图态射的延拓
$(\Gamma_f)^\wedge$ 等同。
\end{corollary}

\begin{corollary}[10.9.9]
\label{I.10.9.9-fr}
设 $X$、$Y$ 是两个局部 Noether 预概形，$f:X\to Y$ 是一个态射，
$Y'$ 是 $Y$ 的一个闭子集，$X'=f^{-1}(Y')$。则藉助交换图表
\[
  \label{I.10.9.9.1-fr}
  \xymatrix{
    X\ar[d]_f &
    X_{/X'}\ar[l]\ar[d]^{\widehat f}\\
    Y &
    Y_{/Y'}\ar[l]
  }
\]
形式预概形 $X_{/X'}$ 与积 $X\times_Y(Y_{/Y'})$ 等同。
\end{corollary}

只须在 \hyperref[I.10.9.7-fr]{10.9.7} 中把 $S$ 和 $S'$ 替换为
$Y$，并把 $X$ 和 $X'$ 替换为 $X$。
\oldpage[I]{201}

\begin{remark}[10.9.10]
\label{I.10.9.10-fr}
若 $X$ 是和 $X_1\amalg X_2$
（\hyperref[subsection:I.3.1-fr]{3.1}），$X'$ 是并集
$X'_1\cup X'_2$，其中 $X'_i$ 是 $X_i$ 的一个闭子集（$i=1,2$），则立即有
$X_{/X'}={X_1}_{/X'_1}\amalg {X_2}_{/X'_2}$。
\end{remark}

\subsection{对仿射形式概形上的凝聚层的应用}
\label{subsection:I.10.10-fr}

\begin{env}[10.10.1]
\label{I.10.10.1-fr}
在本节中，$A$ 始终表示一个\emph{Noether 进制环}，
$\mathfrak{J}$ 表示 $A$ 的一个定义理想。置
$X=\operatorname{Spec}(A)$、$\mathfrak{X}=\operatorname{Spf}(A)$；后者与 $X$ 的闭子集
$V(\mathfrak{J})$ 等同（\hyperref[I.10.1.2-fr]{10.1.2}）。此外，
\hyperref[I.10.1.2-fr]{10.1.2} 的定义和
\hyperref[I.10.8.4-fr]{10.8.4} 的定义表明，
\emph{仿射形式概形 $\mathfrak{X}$} 就是仿射概形 $X$ 沿其底空间的闭子集
$\mathfrak{X}$ 的\emph{完备化 $X_{/\mathfrak{X}}$}。因此，每个凝聚
$\mathscr{O}_X$-模 $\mathscr{F}$ 都对应于一个有限型
$\mathscr{O}_{\mathfrak{X}}$-模 $\mathscr{F}_{/\mathfrak{X}}$；它还是拓扑环层
$\mathscr{O}_{\mathfrak{X}}$ 上的拓扑模层。但每个凝聚
$\mathscr{O}_X$-模 $\mathscr{F}$ 都具有形式 $\widetilde M$，其中 $M$ 是一个有限型
$A$-模（\hyperref[I.1.5.1-fr]{1.5.1}）；我们置
$(\widetilde M)_{/\mathfrak{X}}=M^\Delta$。另外，若 $u:M\to N$ 是有限型
$A$-模的一个 $A$-同态，则它对应于一个同态
$\widetilde u:\widetilde M\to\widetilde N$，因而也对应于一个连续同态
$\widetilde u_{/\mathfrak{X}}:(\widetilde M)_{/\mathfrak{X}}\to
(\widetilde N)_{/\mathfrak{X}}$，我们把它记作 $u^\Delta$。显然有
$(v\circ u)^\Delta=v^\Delta\circ u^\Delta$；这样就定义了一个\emph{可加协变函子
$M^\Delta$}，它从有限型 $A$-模范畴取值于有限型
$\mathscr{O}_{\mathfrak{X}}$-模范畴。当 $A$ 是一个\emph{离散}环时，有
$M^\Delta=\widetilde M$。
\end{env}

\begin{proposition}[10.10.2]
\label{I.10.10.2-fr}
\medskip\noindent
\begin{enumerate}
  \item[\rm{(i)}] $M^\Delta$ 关于 $M$ 是正合函子，并且存在典范的函子性
    $A$-模同构
    $\Gamma(\mathfrak{X},M^\Delta)\xrightarrow{\sim}M$。
  \item[\rm{(ii)}] 若 $M$ 和 $N$ 是两个有限型 $A$-模，则存在典范的函子性同构
    \[
      (M\otimes_A N)^\Delta
      \xrightarrow{\sim}
      M^\Delta\otimes_{\mathscr{O}_{\mathfrak{X}}}N^\Delta
      \tag{10.10.2.1}\label{I.10.10.2.1-fr}
    \]
    \[
      (\operatorname{Hom}_A(M,N))^\Delta
      \xrightarrow{\sim}
      \mathscr{H}\!om_{\mathscr{O}_{\mathfrak{X}}}
      (M^\Delta,N^\Delta)
      \tag{10.10.2.2}\label{I.10.10.2.2-fr}
    \]
  \item[\rm{(iii)}] 应用 $u\to u^\Delta$ 是函子性同构
    \[
      \operatorname{Hom}_A(M,N)
      \xrightarrow{\sim}
      \operatorname{Hom}_{\mathscr{O}_{\mathfrak{X}}}
      (M^\Delta,N^\Delta)
      \tag{10.10.2.3}\label{I.10.10.2.3-fr}
    \]
\end{enumerate}
\end{proposition}

$M^\Delta$ 的正合性来自函子 $\widetilde M$
（\hyperref[I.1.3.5-fr]{1.3.5}）和 $\mathscr{F}_{/\mathfrak{X}}$
（\hyperref[I.10.8.8-fr]{10.8.8}）的正合性。根据定义，
$\Gamma(\mathfrak{X},M^\Delta)$ 是 $A$-模
$\Gamma(X,\widetilde M)=M$ 对于 $\mathfrak{J}$-预进制拓扑的分离完备化；但由于
$A$ 是完备的且 $M$ 是有限型的，我们知道
（\hyperref[0.7.3.6-fr]{0, 7.3.6}）$M$ 是分离且完备的，这就完成了
(i) 的证明。同构
\hyperref[I.10.10.2.1-fr]{10.10.2.1}（相应地，
\hyperref[I.10.10.2.2-fr]{10.10.2.2}）由同构
\hyperref[I.1.3.12-fr]{1.3.12, (i)} 和
\hyperref[I.10.8.10.1-fr]{10.8.10.1}（相应地，
\hyperref[I.1.3.12-fr]{1.3.12, (ii)} 和
\hyperref[I.10.8.10.2-fr]{10.8.10.2}）的复合得到。最后，由于
$\operatorname{Hom}_A(M,N)$ 是一个有限型 $A$-模，可以对它应用 (i)，从而把
$\Gamma(\mathfrak{X},(\operatorname{Hom}_A(M,N))^\Delta)$ 与
$\operatorname{Hom}_A(M,N)$ 等同；再利用
\hyperref[I.10.10.2.2-fr]{10.10.2.2}，便证明同态
\hyperref[I.10.10.2.3-fr]{10.10.2.3} 是一个同构。

由 \hyperref[I.10.10.2-fr]{10.10.2} 可以推出一系列结论，它们类似于从
\hyperref[I.1.3.7-fr]{1.3.7} 和
\hyperref[I.1.3.12-fr]{1.3.12} 推出的那些结论；我们留给读者自行表述。
\oldpage[I]{202}

注意，把 $M^\Delta$ 的正合性应用于正合列
$0\to\mathfrak{J}\to A\to A/\mathfrak{J}\to0$，可知这里记作
$\mathfrak{J}^\Delta$ 的 $\mathscr{O}_{\mathfrak{X}}$ 的理想层，与
\hyperref[I.10.3.1-fr]{10.3.1} 中以同一记号表示的理想层一致；这是
\hyperref[I.10.3.2-fr]{10.3.2} 的结果。

\begin{proposition}[10.10.3]
\label{I.10.10.3-fr}
在 \hyperref[I.10.10.1-fr]{10.10.1} 的假设下，
$\mathscr{O}_{\mathfrak{X}}$ 是一个凝聚环层。
\end{proposition}

若 $f\in A$，我们知道 $A_{\{f\}}$ 是一个 Noether 进制环
（\hyperref[0.7.6.11-fr]{0, 7.6.11}）；由于问题是局部的，可归结为
（\hyperref[I.10.1.4-fr]{10.1.4}）证明同态
$v:\mathscr{O}_{\mathfrak{X}}^n\to\mathscr{O}_{\mathfrak{X}}$ 的核是一个有限型
$\mathscr{O}_{\mathfrak{X}}$-模。此时有 $v=u^\Delta$，其中 $u$ 是一个
$A$-同态 $A^n\to A$
（\hyperref[I.10.10.2-fr]{10.10.2}）；由于 $A$ 是 Noether 环，$u$ 的核是有限型的，
换言之，存在一个同态 $A^m\xrightarrow{w}A^n$，使得列
$A^m\xrightarrow{w}A^n\xrightarrow{u}A$ 正合。因此由
\hyperref[I.10.10.2-fr]{10.10.2} 得到正合列
$\mathscr{O}_{\mathfrak{X}}^m\xrightarrow{w^\Delta}
\mathscr{O}_{\mathfrak{X}}^n\xrightarrow{v}
\mathscr{O}_{\mathfrak{X}}$，这就证明了 $v$ 的核是有限型的。

\begin{env}[10.10.4]
\label{I.10.10.4-fr}
使用上述记号，置 $A_n=A/\mathfrak{J}^{n+1}$，并设 $X_n$ 是仿射概形
$\operatorname{Spec}(A_n)=(\mathfrak{X},
\mathscr{O}_{\mathfrak{X}}/\mathscr{J}^{n+1})$，其中
$\mathscr{J}=\mathfrak{J}^\Delta$ 是与理想 $\mathfrak{J}$ 对应的
$\mathscr{O}_{\mathfrak{X}}$ 的定义理想层。设 $u_{mn}$ 是与典范同态
$A_n\to A_m$ 对应的预概形态射 $X_m\to X_n$，其中
$m\leq n$；形式概形 $\mathfrak{X}$ 是各 $X_n$ 对于 $u_{mn}$ 的归纳极限
（\hyperref[I.10.6.3-fr]{10.6.3}）。
\end{env}

\begin{proposition}[10.10.5]
\label{I.10.10.5-fr}
在 \hyperref[I.10.10.1-fr]{10.10.1} 的假设下，设
$\mathscr{F}$ 是一个 $\mathscr{O}_{\mathfrak{X}}$-模。下列条件等价：
\begin{enumerate}
  \item[a)] $\mathscr{F}$ 是一个凝聚 $\mathscr{O}_{\mathfrak{X}}$-模。
  \item[b)] $\mathscr{F}$ 同构于一个凝聚 $\mathscr{O}_{X_n}$-模序列
    $(\mathscr{F}_n)$ 的逆极限
    （\hyperref[I.10.6.6-fr]{10.6.6}），并且
    $u_{mn}^*(\mathscr{F}_n)=\mathscr{F}_m$。
  \item[c)] 存在一个有限型 $A$-模 $M$（根据
    \hyperref[I.10.10.2-fr]{10.10.2, (i)}，它在典范同构意义下唯一确定），使得
    $\mathscr{F}$ 同构于 $M^\Delta$。
\end{enumerate}
\end{proposition}

先证明 b) 推出 c)。有 $\mathscr{F}_n=\widetilde{M_n}$，其中 $M_n$ 是一个有限型
$A_n$-模，而假设导出
$M_m=M_n\otimes_{A_n}A_m$，其中 $m\leq n$
（\hyperref[I.1.6.5-fr]{1.6.5}）；因此，各 $M_n$ 对于典范双同态 $M_n\to M_m$
（$m\leq n$）构成一个逆系。由各 $A_n$ 的定义立即可知，这个逆系满足
\hyperref[0.7.2.9-fr]{0, 7.2.9} 的条件；因而它的逆极限 $M$ 是一个有限型
$A$-模，使得对每个 $n$ 都有 $M_n=M\otimes_A A_n$。由此得出，
$\mathscr{F}_n$ 由
$\widetilde M\otimes_{\mathscr{O}_X}
(\mathscr{O}_X/\widetilde{\mathfrak{J}}^{n+1})$ 在 $X_n$ 上诱导而来，因此按定义
（\hyperref[I.10.8.4-fr]{10.8.4}）有 $\mathscr{F}=M^\Delta$。

反过来，c) 推出 b)。事实上，若 $u_n$ 是浸入态射 $X_n\to X$，则
$u_n^*(\widetilde M)=(M\otimes_A A_n)^{\sim}$ 由
$\widetilde M\otimes_{\mathscr{O}_X}
(\mathscr{O}_X/\widetilde{\mathfrak{J}}^{n+1})$ 在 $X_n$ 上诱导而来，并且按定义
（\hyperref[I.10.8.4-fr]{10.8.4}）有
$M^\Delta=\varprojlim u_n^*(\widetilde M)$；由于当 $m\leq n$ 时
$u_m=u_n\circ u_{mn}$，各 $\mathscr{F}_n=u_n^*(\widetilde M)$ 满足 b) 的条件，从而得到所断言的结论。

现在证明 c) 推出 a)。事实上，按定义有
$\mathscr{O}_{\mathfrak{X}}=A^\Delta$；因为 $M$ 是某个同态
$A^m\to A^n$ 的余核，所以由
\hyperref[I.10.10.2-fr]{10.10.2} 可知，$M^\Delta$ 是某个同态
$\mathscr{O}_{\mathfrak{X}}^m\to\mathscr{O}_{\mathfrak{X}}^n$ 的余核；又因为环层
$\mathscr{O}_{\mathfrak{X}}$ 是凝聚的
（\hyperref[I.10.10.3-fr]{10.10.3}），所以 $M^\Delta$ 也是凝聚的
（\hyperref[0.5.3.4-fr]{0, 5.3.4}）。
\oldpage[I]{203}

最后，a) 推出 b)。把它们视为 $\mathscr{O}_{\mathfrak{X}}$-模，有
$\mathscr{O}_{X_n}=\mathscr{O}_{\mathfrak{X}}/
\mathscr{J}^{n+1}=A_n^\Delta$；
$\mathscr{F}_n=\mathscr{F}
\otimes_{\mathscr{O}_{\mathfrak{X}}}\mathscr{O}_{X_n}$ 是一个凝聚
$\mathscr{O}_{\mathfrak{X}}$-模
（\hyperref[0.5.3.5-fr]{0, 5.3.5}）；又因为它也是一个
$\mathscr{O}_{X_n}$-模且 $\mathscr{J}^{n+1}$ 是凝聚的，所以可知
$\mathscr{F}_n$ 是一个凝聚 $\mathscr{O}_{X_n}$-模
（\hyperref[0.5.3.10-fr]{0, 5.3.10}）。显然，当 $m\leq n$ 时有
$u_{mn}^*(\mathscr{F}_n)=\mathscr{F}_m$（要记住，底空间的连续映射
$X_m\to X_n$ 就是 $\mathfrak{X}$ 的恒等映射）。因此层
$\mathscr{G}=\varprojlim\mathscr{F}_n$ 是一个凝聚
$\mathscr{O}_{\mathfrak{X}}$-模，因为我们已经看到 b) 推出 a)。典范同态
$\mathscr{F}\to\mathscr{F}_n$ 构成一个逆系，取逆极限得到一个典范同态
$w:\mathscr{F}\to\mathscr{G}$；所有问题归结为证明 $w$ 是双射。现在这个问题是局部的，
因此可只考虑 $\mathscr{F}$ 是某个同态
$\mathscr{O}_{\mathfrak{X}}^p\to\mathscr{O}_{\mathfrak{X}}^q$ 的余核的情形。这个同态具有形式
$v^\Delta$，其中 $v$ 是一个同态 $A^m\to A^n$
（\hyperref[I.10.10.2-fr]{10.10.2}）；因而 $\mathscr{F}$ 同构于 $M^\Delta$，其中
$M=\operatorname{Coker}v$
（\hyperref[I.10.10.2-fr]{10.10.2}）。此时，由
\hyperref[I.10.10.2-fr]{10.10.2} 有
$\mathscr{F}_n=M^\Delta\otimes_{\mathscr{O}_{\mathfrak{X}}}A_n^\Delta
=(M\otimes_A A_n)^\Delta$；由于 $M\otimes_A A_n$ 上的 $\mathfrak{J}$-进制拓扑是离散的，有
$(M\otimes_A A_n)^\Delta=(M\otimes_A A_n)^{\sim}$（把它视为
$\mathscr{O}_{X_n}$-模）。我们在上面已经看到
$M^\Delta=\varprojlim\mathscr{F}_n$，因此在这种情形下 $w$ 确实是恒等态射。
\par\hfill C.Q.F.D.\par

\begin{corollary}[10.10.6]
\label{I.10.10.6-fr}
若 $\mathscr{F}$ 满足 \hyperref[I.10.10.5-fr]{10.10.5} 的条件 b)，则逆系
$(\mathscr{F}_n)$ 同构于由各
$\mathscr{F}\otimes_{\mathscr{O}_{\mathfrak{X}}}\mathscr{O}_{X_n}$ 组成的逆系。
\end{corollary}

\begin{env}[10.10.7]
\label{I.10.10.7-fr}
现设 $A$、$B$ 是两个 Noether 进制环，
$\varphi:B\to A$ 是一个连续同态；用 $\mathfrak{J}$（相应地，
$\mathfrak{K}$）表示 $A$（相应地，$B$）的一个定义理想，使得
$\varphi(\mathfrak{K})\subset\mathfrak{J}$，并置
$X=\operatorname{Spec}(A)$、$Y=\operatorname{Spec}(B)$、
$\mathfrak{X}=\operatorname{Spf}(A)$、
$\mathfrak{Y}=\operatorname{Spf}(B)$。设 $f:X\to Y$ 是与 $\varphi$ 对应的预概形态射
（\hyperref[I.1.6.1-fr]{1.6.1}），$\widehat f:\mathfrak{X}\to
\mathfrak{Y}$ 是它到完备化上的延拓
（\hyperref[I.10.9.1-fr]{10.9.1}）；它也是与 $\varphi$ 对应的形式预概形态射
（\hyperref[I.10.2.2-fr]{10.2.2}）。
\end{env}

\begin{proposition}[10.10.8]
\label{I.10.10.8-fr}
对每个有限型 $B$-模 $N$，存在典范的函子性
$\mathscr{O}_{\mathfrak{X}}$-模同构
\[
  \widehat f^*(N^\Delta)
  \xrightarrow{\sim}
  (N\otimes_B A)^\Delta.
\]
\end{proposition}

事实上，分别用 $i_X:\mathfrak{X}\to X$ 和
$i_Y:\mathfrak{Y}\to Y$ 表示典范态射。在典范函子性同构的意义下，根据
\hyperref[I.10.8.8-fr]{10.8.8} 有 $N^\Delta=i_Y^*(\widetilde N)$ 以及
\[
  (N\otimes_B A)^\Delta
  =i_X^*((N\otimes_B A)^{\sim})
  =i_X^*(f^*(\widetilde N))
\]
（\hyperref[I.1.6.5-fr]{1.6.5}）；因此，命题由图表
\hyperref[I.10.9.2-fr]{10.9.2} 的交换性得出。

\begin{corollary}[10.10.9]
\label{I.10.10.9-fr}
对 $B$ 的任意理想 $\mathfrak b$，都有
$\widehat f^*(\mathfrak b^\Delta)\mathscr{O}_{\mathfrak{X}}
=(\mathfrak bA)^\Delta$。
\end{corollary}

事实上，设 $j$ 是典范注入 $\mathfrak b\to B$；它对应于
$\mathscr{O}_{\mathfrak{Y}}$-模层的典范注入
$j^\Delta:\mathfrak b^\Delta\to\mathscr{O}_{\mathfrak{Y}}$。根据定义，
$\widehat f^*(\mathfrak b^\Delta)\mathscr{O}_{\mathfrak{X}}$ 是同态
$\widehat f^*(j^\Delta):\widehat f^*(\mathfrak b^\Delta)\to
\mathscr{O}_{\mathfrak{X}}=\widehat f^*(\mathscr{O}_{\mathfrak{Y}})$ 的像；但根据
\hyperref[I.10.10.8-fr]{10.10.8}，这个同态与
$(j\otimes1)^\Delta:(\mathfrak b\otimes_B A)^\Delta\to
\mathscr{O}_{\mathfrak{X}}=(B\otimes_B A)^\Delta$ 等同。由于
$j\otimes1$ 的像是 $A$ 的理想 $\mathfrak bA$，根据
\hyperref[I.10.10.2-fr]{10.10.2}，$(j\otimes1)^\Delta$ 的像因而是
$(\mathfrak bA)^\Delta$，由此得出结论。
\oldpage[I]{204}

\subsection{形式预概形上的凝聚层}
\label{subsection:I.10.11-fr}

\begin{proposition}[10.11.1]
\label{I.10.11.1-fr}
若 $\mathfrak{X}$ 是一个局部 Noether 形式预概形，则环层
$\mathscr{O}_{\mathfrak{X}}$ 是凝聚的，并且 $\mathfrak{X}$ 的每个定义理想层都是凝聚的。
\end{proposition}

由于问题是局部的，可归结到 Noether 仿射形式概形的情形；因此命题由
\hyperref[I.10.10.3-fr]{10.10.3} 和
\hyperref[I.10.10.5-fr]{10.10.5} 得出。

\begin{env}[10.11.2]
\label{I.10.11.2-fr}
设 $\mathfrak{X}$ 是一个局部 Noether 形式预概形，
$\mathscr{J}$ 是 $\mathfrak{X}$ 的一个定义理想层，$X_n$ 是局部 Noether 的（通常的）预概形
$(\mathfrak{X},\mathscr{O}_{\mathfrak{X}}/\mathscr{J}^{n+1})$；于是，$\mathfrak{X}$ 是序列
$(X_n)$ 对于典范态射 $u_{mn}:X_m\to X_n$ 的归纳极限
（\hyperref[I.10.6.3-fr]{10.6.3}）。使用这些记号：
\end{env}

\begin{theorem}[10.11.3]
\label{I.10.11.3-fr}
一个 $\mathscr{O}_{\mathfrak{X}}$-模 $\mathscr{F}$ 为凝聚的充要条件是：它同构于一个序列
$(\mathscr{F}_n)$ 的逆极限，其中 $\mathscr{F}_n$ 是凝聚
$\mathscr{O}_{X_n}$-模，并且当 $m\leq n$ 时有
$u_{mn}^*(\mathscr{F}_n)=\mathscr{F}_m$
（\hyperref[I.10.6.6-fr]{10.6.6}）。此时，逆系
$(\mathscr{F}_n)$ 同构于由各
$u_n^*(\mathscr{F})=\mathscr{F}
\otimes_{\mathscr{O}_{\mathfrak{X}}}\mathscr{O}_{X_n}$ 组成的逆系，其中 $u_n$ 是典范态射
$X_n\to\mathfrak{X}$。
\end{theorem}

由于问题是局部的，可归结到 $\mathfrak{X}$ 是 Noether 仿射形式概形的情形；此时定理是
\hyperref[I.10.10.5-fr]{10.10.5} 和
\hyperref[I.10.10.6-fr]{10.10.6} 的结果。

因此可以说，\emph{给定一个凝聚
$\mathscr{O}_{\mathfrak{X}}$-模，等价于给定一个凝聚
$\mathscr{O}_{X_n}$-模逆系 $(\mathscr{F}_n)$，使得当 $m\leq n$ 时
$u_{mn}^*(\mathscr{F}_n)=\mathscr{F}_m$}。

\begin{corollary}[10.11.4]
\label{I.10.11.4-fr}
若 $\mathscr{F}$ 和 $\mathscr{G}$ 是两个凝聚 $\mathscr{O}_{\mathfrak{X}}$-模，则（使用
\hyperref[I.10.11.3-fr]{10.11.3} 的记号）可定义一个典范的函子性同构
\[
  \label{I.10.11.4.1-fr}
  \operatorname{Hom}_{\mathscr{O}_{\mathfrak{X}}}
  (\mathscr{F},\mathscr{G})
  \xrightarrow{\sim}
  \varprojlim_n
  \operatorname{Hom}_{\mathscr{O}_{X_n}}
  (\mathscr{F}_n,\mathscr{G}_n).
  \tag{10.11.4.1}
\]
\end{corollary}

右端的逆极限是针对映射
$\theta_n\mapsto u_{mn}^*(\theta_n)$（$m\leq n$）而取的；这些映射从
$\operatorname{Hom}_{\mathscr{O}_{X_n}}
(\mathscr{F}_n,\mathscr{G}_n)$ 到
$\operatorname{Hom}_{\mathscr{O}_{X_m}}
(\mathscr{F}_m,\mathscr{G}_m)$。同态
\hyperref[I.10.11.4.1-fr]{10.11.4.1} 把元素
$\theta\in\operatorname{Hom}_{\mathscr{O}_{\mathfrak{X}}}
(\mathscr{F},\mathscr{G})$ 对应到序列 $(u_n^*(\theta))$。反过来，对逆系
$(\theta_n)\in\varprojlim_n
\operatorname{Hom}_{\mathscr{O}_{X_n}}
(\mathscr{F}_n,\mathscr{G}_n)$，把它对应到它在
$\operatorname{Hom}_{\mathscr{O}_{\mathfrak{X}}}
(\mathscr{F},\mathscr{G})$ 中的逆极限，再根据
\hyperref[I.10.11.3-fr]{10.11.3}，立即可见这就定义了前一同态的逆同态。

\begin{corollary}[10.11.5]
\label{I.10.11.5-fr}
同态 $\theta:\mathscr{F}\to\mathscr{G}$ 为满射的充要条件是，与它对应的同态
$\theta_0=u_0^*(\theta):\mathscr{F}_0\to\mathscr{G}_0$ 为满射。
\end{corollary}

由于问题是局部的，可归结到如下情形：
$\mathfrak{X}=\operatorname{Spf}(A)$，$A$ 是 Noether 进制环，
$\mathscr{F}=M^\Delta$、$\mathscr{G}=N^\Delta$ 且 $\theta=u^\Delta$，其中 $M$ 和 $N$ 是有限型
$A$-模，$u$ 是一个同态 $M\to N$。此外，此时有
$\theta_0=\widetilde{u_0}$，其中 $u_0$ 是同态
$u\otimes1:M\otimes_A A/\mathfrak{J}\to
N\otimes_A A/\mathfrak{J}$。结论来自：$\theta$ 和 $u$（相应地，$\theta_0$ 和
$u_0$）同时为满射
（\hyperref[I.1.3.9-fr]{1.3.9} 和
\hyperref[I.10.10.2-fr]{10.10.2}），而 $u$ 和 $u_0$ 同时为满射
（\hyperref[0.7.1.14-fr]{0, 7.1.14}）。

\begin{env}[10.11.6]
\label{I.10.11.6-fr}
定理 \hyperref[I.10.11.3-fr]{10.11.3} 表明，可以把每个凝聚
$\mathscr{O}_{\mathfrak{X}}$-模 $\mathscr{F}$ 视为一个\emph{拓扑}
$\mathscr{O}_{\mathfrak{X}}$-模，即把它视为各\emph{伪离散}群层 $\mathscr{F}_n$ 的逆极限
（\hyperref[0.3.8.1-fr]{0, 3.8.1}）。于是，由
\hyperref[I.10.11.4-fr]{10.11.4} 可知，凝聚
$\mathscr{O}_{\mathfrak{X}}$-模的每个同态
$u:\mathscr{F}\to\mathscr{G}$ 都自动地是\emph{连续}的。

\oldpage[I]{205}

（\hyperref[0.3.8.2-fr]{0, 3.8.2}）。此外，若 $\mathscr{H}$ 是凝聚
$\mathscr{O}_{\mathfrak{X}}$-模 $\mathscr{F}$ 的一个凝聚子 $\mathscr{O}_{\mathfrak{X}}$-模，则对每个开集
$U\subset\mathfrak{X}$，$\Gamma(U,\mathscr{H})$ 是拓扑群
$\Gamma(U,\mathscr{F})$ 的一个\emph{闭}子群。事实上，由于函子
$\Gamma$ 左正合，$\Gamma(U,\mathscr{H})$ 是同态
$\Gamma(U,\mathscr{F})\to\Gamma(U,\mathscr{F}/\mathscr{H})$ 的核；这个同态根据上面的结果是
\emph{连续}的，因为 $\mathscr{F}/\mathscr{H}$ 是凝聚的
（\hyperref[0.5.3.4-fr]{0, 5.3.4}）。因为 $\Gamma(U,\mathscr{F}/\mathscr{H})$ 是分离拓扑群，所以断言得证。
\end{env}

\begin{proposition}[10.11.7]
\label{I.10.11.7-fr}
设 $\mathscr{F}$ 和 $\mathscr{G}$ 是两个凝聚 $\mathscr{O}_{\mathfrak{X}}$-模。（使用
\hyperref[I.10.11.3-fr]{10.11.3} 的记号）可定义拓扑
$\mathscr{O}_{\mathfrak{X}}$-模的典范函子性同构
（\hyperref[I.10.11.6-fr]{10.11.6}）
\[
  \label{I.10.11.7.1-fr}
  \mathscr{F}\otimes_{\mathscr{O}_{\mathfrak{X}}}\mathscr{G}
  \xrightarrow{\sim}
  \varprojlim_n
  (\mathscr{F}_n\otimes_{\mathscr{O}_{X_n}}\mathscr{G}_n)
  \tag{10.11.7.1}
\]
\[
  \label{I.10.11.7.2-fr}
  \mathscr{H}\!om_{\mathscr{O}_{\mathfrak{X}}}
  (\mathscr{F},\mathscr{G})
  \xrightarrow{\sim}
  \varprojlim_n
  \mathscr{H}\!om_{\mathscr{O}_{X_n}}(\mathscr{F}_n,\mathscr{G}_n)
  \tag{10.11.7.2}
\]
\end{proposition}

同构 \hyperref[I.10.11.7.1-fr]{10.11.7.1} 的存在性由等式
\[
  \mathscr{F}_n\otimes_{\mathscr{O}_{X_n}}\mathscr{G}_n
  =
  (\mathscr{F}\otimes_{\mathscr{O}_{\mathfrak{X}}}
  \mathscr{O}_{X_n})
  \otimes_{\mathscr{O}_{X_n}}
  (\mathscr{G}\otimes_{\mathscr{O}_{\mathfrak{X}}}
  \mathscr{O}_{X_n})
  =
  (\mathscr{F}\otimes_{\mathscr{O}_{\mathfrak{X}}}\mathscr{G})
  \otimes_{\mathscr{O}_{\mathfrak{X}}}\mathscr{O}_{X_n}
\]
以及 \hyperref[I.10.11.3-fr]{10.11.3} 得出。如果把两端看成不带拓扑的模层，则同构
\hyperref[I.10.11.7.2-fr]{10.11.7.2} 来自以下两类层的截面定义：
$\mathscr{H}\!om_{\mathscr{O}_{\mathfrak{X}}}
(\mathscr{F},\mathscr{G})$ 和
$\mathscr{H}\!om_{\mathscr{O}_{X_n}}(\mathscr{F}_n,\mathscr{G}_n)$；还要使用同构
\hyperref[I.10.11.4.1-fr]{10.11.4.1} 的存在性，把它应用于 $\mathfrak{X}$ 的任意 Noether 仿射形式开集上诱导的预概形。剩下只须证明，在一个拟紧集上，同构
\hyperref[I.10.11.7.2-fr]{10.11.7.2} 是双连续的。因此可归结到
$\mathfrak{X}=\operatorname{Spf}(A)$、$A$ 为 Noether 进制环的情形；此时由
\hyperref[I.10.10.5-fr]{10.10.5} 有
$\mathscr{F}=M^\Delta$、$\mathscr{G}=N^\Delta$，其中 $M$、$N$ 是有限型 $A$-模。考虑到
\hyperref[I.10.10.2.1-fr]{10.10.2.1}、
\hyperref[I.10.10.2.3-fr]{10.10.2.3} 以及
\hyperref[I.1.3.12-fr]{1.3.12, (ii)}，问题归结为证明典范同构
$\operatorname{Hom}_A(M,N)\xrightarrow{\sim}
\varprojlim_n\operatorname{Hom}_{A_n}(M_n,N_n)$（其中
$M_n=M\otimes_A A_n$、$N_n=N\otimes_A A_n$）是连续的；这已经在
\hyperref[0.7.8.2-fr]{0, 7.8.2} 中证明。

\begin{env}[10.11.8]
\label{I.10.11.8-fr}
由于 $\operatorname{Hom}_{\mathscr{O}_{\mathfrak{X}}}
(\mathscr{F},\mathscr{G})$ 是拓扑群层
$\mathscr{H}\!om_{\mathscr{O}_{\mathfrak{X}}}
(\mathscr{F},\mathscr{G})$ 的截面群，它配有一个群拓扑。若
$\mathfrak{X}$ 是\emph{Noether}的，则由
\hyperref[I.10.11.7.2-fr]{10.11.7.2} 可知，这个群中 $0$ 的一个基本邻域系可以取为各子群
$\operatorname{Hom}_{\mathscr{O}_{\mathfrak{X}}}
(\mathscr{F},\mathscr{J}^n\mathscr{G})$（$n$ 任意）。
\end{env}

\begin{proposition}[10.11.9]
\label{I.10.11.9-fr}
设 $\mathfrak{X}$ 是一个 Noether 形式预概形，$\mathscr{F}$ 和 $\mathscr{G}$ 是两个凝聚
$\mathscr{O}_{\mathfrak{X}}$-模。在拓扑群
$\operatorname{Hom}_{\mathscr{O}_{\mathfrak{X}}}
(\mathscr{F},\mathscr{G})$ 中，满射同态（相应地，单射同态、双射同态）构成一个开子集。
\end{proposition}

根据 \hyperref[I.10.11.5-fr]{10.11.5}，
$\operatorname{Hom}_{\mathscr{O}_{\mathfrak{X}}}
(\mathscr{F},\mathscr{G})$ 中满射同态所成的集合，是离散群
$\operatorname{Hom}_{\mathscr{O}_{X_0}}(\mathscr{F}_0,\mathscr{G}_0)$ 的一个子集在连续映射
$\operatorname{Hom}_{\mathscr{O}_{\mathfrak{X}}}
(\mathscr{F},\mathscr{G})\to
\operatorname{Hom}_{\mathscr{O}_{X_0}}(\mathscr{F}_0,\mathscr{G}_0)$ 下的逆像，从而第一个断言得证。为证明第二个断言，用有限个 Noether 仿射形式开集 $U_i$ 覆盖
$\mathfrak{X}$。一个元素
$\theta\in\operatorname{Hom}_{\mathscr{O}_{\mathfrak{X}}}
(\mathscr{F},\mathscr{G})$ 为单射的充要条件是，它在所有（连续的）限制映射
$\operatorname{Hom}_{\mathscr{O}_{\mathfrak{X}}}
(\mathscr{F},\mathscr{G})\to
\operatorname{Hom}_{\mathscr{O}_{\mathfrak{X}}|U_i}
(\mathscr{F}|U_i,\mathscr{G}|U_i)$ 下的像均为单射。因此可归结到仿射情形，而这已经在
\hyperref[0.7.8.3-fr]{0, 7.8.3} 中证明。

\oldpage[I]{206}

\subsection{形式预概形间的进制态射}
\label{subsection:I.10.12-fr}

\begin{env}[10.12.1]
\label{I.10.12.1-fr}
设 $\mathfrak{X}$、$\mathfrak{S}$ 是两个\emph{局部 Noether}形式预概形；若存在
$\mathfrak{S}$ 的一个定义理想 $\mathscr{J}$，使得
$\mathscr{K}=f^*(\mathscr{J})\mathscr{O}_{\mathfrak{X}}$ 是
$\mathfrak{X}$ 的一个定义理想，则称态射
$f:\mathfrak{X}\to\mathfrak{S}$ 为\emph{进制的}；这时也称
$\mathfrak{X}$ 是一个\emph{进制 $\mathfrak{S}$-预概形}（对于 $f$ 而言）。
在这种情形下，对 $\mathfrak{S}$ 的\emph{任意}定义理想 $\mathscr{J}_1$，
$\mathscr{K}_1=f^*(\mathscr{J}_1)\mathscr{O}_{\mathfrak{X}}$ 都是
$\mathfrak{X}$ 的一个定义理想。事实上，问题是局部的，故可假设
$\mathfrak{X}$ 和 $\mathfrak{S}$ 都是 Noether 仿射形式概形；于是存在整数
$n$，使得 $\mathscr{J}^n\subset\mathscr{J}_1$ 且
$\mathscr{J}_1^n\subset\mathscr{J}$
（\hyperref[I.10.3.6-fr]{10.3.6} 和
\hyperref[0.7.1.4-fr]{0, 7.1.4}），从而
$\mathscr{K}^n\subset\mathscr{K}_1$ 且
$\mathscr{K}_1^n\subset\mathscr{K}$。第一个关系表明
$\mathscr{K}_1=\mathfrak{K}_1^\Delta$，其中 $\mathfrak{K}_1$ 是
$A=\Gamma(\mathfrak{X},\mathscr{O}_{\mathfrak{X}})$ 的一个开理想；第二个关系表明
$\mathfrak{K}_1$ 是 $A$ 的一个定义理想
（\hyperref[0.7.1.4-fr]{0, 7.1.4}），断言得证。
\end{env}

由上文立即可知，若 $\mathfrak{X}$ 和 $\mathfrak{Y}$ 是两个进制
$\mathfrak{S}$-预概形，则\emph{每个 $\mathfrak{S}$-态射
$u:\mathfrak{X}\to\mathfrak{Y}$ 都是进制的}。事实上，若
$f:\mathfrak{X}\to\mathfrak{S}$、
$g:\mathfrak{Y}\to\mathfrak{S}$ 是结构态射，$\mathscr{J}$ 是
$\mathfrak{S}$ 的一个定义理想，则 $f=g\circ u$，因而
$u^*(g^*(\mathscr{J})\mathscr{O}_{\mathfrak{Y}})
\mathscr{O}_{\mathfrak{X}}=f^*(\mathscr{J})
\mathscr{O}_{\mathfrak{X}}$ 是 $\mathfrak{X}$ 的一个定义理想；而根据假设，
$g^*(\mathscr{J})\mathscr{O}_{\mathfrak{Y}}$ 是 $\mathfrak{Y}$ 的一个定义理想。

\begin{env}[10.12.2]
\label{I.10.12.2-fr}
以下固定一个局部 Noether 形式预概形 $\mathfrak{S}$ 以及
$\mathfrak{S}$ 的一个定义理想 $\mathscr{J}$，并令
$S_n=(\mathfrak{S},\mathscr{O}_{\mathfrak{S}}/\mathscr{J}^{n+1})$。
进制（局部 Noether）$\mathfrak{S}$-预概形显然构成一个\emph{范畴}。
若局部 Noether 的（通常）$S_n$-预概形归纳系 $(X_n)$ 的结构态射
$f_n:X_n\to S_n$ 具有如下性质，则称它为一个
\emph{进制 $(S_n)$-归纳系}：当 $m\leq n$ 时，图表
\[
  \label{I.10.12.2.1-fr}
  \xymatrix{
    X_n \ar[d]_{f_n}
    & X_m \ar[l] \ar[d]^{f_m}\\
    S_n
    & S_m \ar[l]
  }
  \tag{10.12.2.1}
\]
交换，并且\emph{把 $X_m$ 等同于积
$X_n\times_{S_n}S_m=(X_n)_{(S_m)}$}。进制归纳系构成一个\emph{范畴}：
事实上，只须把这种系统的态射 $(X_n)\to(Y_n)$ 定义为一个
\emph{$S_n$-态射 $u_n:X_n\to Y_n$ 的归纳系}，使得当 $m\leq n$ 时，
$u_m$ 等同于 $(u_n)_{(S_m)}$。于是有：
\end{env}

\begin{theorem}[10.12.3]
\label{I.10.12.3-fr}
进制 $\mathfrak{S}$-预概形范畴与进制 $(S_n)$-归纳系范畴之间有一个典范等价。
\end{theorem}

这个等价按如下方式得到。若 $\mathfrak{X}$ 是一个进制
$\mathfrak{S}$-预概形，$f:\mathfrak{X}\to\mathfrak{S}$ 是结构态射，则
$\mathscr{K}=f^*(\mathscr{J})\mathscr{O}_{\mathfrak{X}}$ 是
$\mathfrak{X}$ 的一个定义理想；把 $\mathfrak{X}$ 对应到由各
$X_n=(\mathfrak{X},\mathscr{O}_{\mathfrak{X}}/\mathscr{K}^{n+1})$ 组成的归纳系，其中结构态射
$f_n:X_n\to S_n$ 对应于 $f$
（\hyperref[I.10.5.6-fr]{10.5.6}）。先证明 $(X_n)$ 是一个
\emph{进制归纳系}：若 $f=(\psi,\theta)$，则
$\psi^*(\mathscr{J})\mathscr{O}_{\mathfrak{X}}=\mathscr{K}$，故对每个
$n$ 都有
$\psi^*(\mathscr{J}^n)\mathscr{O}_{\mathfrak{X}}=\mathscr{K}^n$；并且由函子
$\psi^*$ 的正合性，当 $m\leq n$ 时有
$\mathscr{K}^{m+1}/\mathscr{K}^{n+1}
=\psi^*(\mathscr{J}^{m+1}/\mathscr{J}^{n+1})
(\mathscr{O}_{\mathfrak{X}}/\mathscr{K}^{n+1})$。因此，结论由
（\hyperref[I.4.4.5-fr]{4.4.5}）得出。此外，立即可以验证，对于进制
$\mathfrak{S}$-预概形的一个 $\mathfrak{S}$-态射
$u:\mathfrak{X}\to\mathfrak{Y}$，它对应于（采用显然的记号）

\oldpage[I]{207}

一个 $S_n$-态射 $u_n:X_n\to Y_n$ 的归纳系，使得当 $m\leq n$ 时，
$u_m$ 等同于 $(u_n)_{(S_m)}$。

上述构造确实定义了一个等价，这将由下面更精确的命题得出：

\begin{proposition}[10.12.3.1]
\label{I.10.12.3.1-fr}
设 $(X_n)$ 是一个 $S_n$-预概形归纳系；假设结构态射
$f_n:X_n\to S_n$ 使得图表
（\hyperref[I.10.12.2.1-fr]{10.12.2.1}）交换，并且当 $m\leq n$ 时把
$X_m$ 等同于 $X_n\times_{S_n}S_m$。那么归纳系 $(X_n)$ 满足
（\hyperref[I.10.6.3-fr]{10.6.3}）的条件 b) 和 c)。设
$\mathfrak{X}$ 是它的归纳极限，
$f:\mathfrak{X}\to\mathfrak{S}$ 是归纳系 $(f_n)$ 的归纳极限态射。
那么，若 $X_0$ 是局部 Noether 的，则 $\mathfrak{X}$ 是局部 Noether 的，
并且 $f$ 是进制态射。
\end{proposition}

由于在 $S_n$ 中定义子预概形 $S_m$ 的
$\mathscr{O}_{S_n}$-理想层是幂零的，根据
（\hyperref[I.4.4.5-fr]{4.4.5}），在 $X_n$ 中定义子预概形 $X_m$ 的
$\mathscr{O}_{X_n}$-理想层也是幂零的；所以
（\hyperref[I.10.6.3-fr]{10.6.3}）的条件确实满足。因此，问题在
$\mathfrak{X}$ 和 $\mathfrak{S}$ 上都是局部的，可假设
$\mathfrak{S}=\operatorname{Spf}(A)$、
$\mathscr{J}=\mathfrak{J}^\Delta$，其中 $A$ 是 Noether
$\mathfrak{J}$-进制环，并且 $X_n=\operatorname{Spec}(B_n)$。若
$A_n=A/\mathfrak{J}^{n+1}$，则假设蕴含 $B_0$ 是 Noether 的；若令
$\mathfrak{J}_n=\mathfrak{J}/\mathfrak{J}^{n+1}$，则
$B_m=B_n/\mathfrak{J}_n^{m+1}B_n$。同态 $B_n\to B_0$ 的核因而是
$\mathfrak{K}_n=\mathfrak{J}_nB_n$，而当 $m\leq n$ 时，同态
$B_n\to B_m$ 的核是 $\mathfrak{K}_n^{m+1}$。此外，由于 $A_1$ 是
Noether 的，$\mathfrak{J}_1$ 在 $A_1$ 上是有限型的，故
$\overline{\mathfrak{K}}_1=\mathfrak{K}_1/\mathfrak{K}_1^2$ 在 $B_1$ 上是有限型的，
并且\emph{当然}在 $B_0=B_1/\mathfrak{K}_1$ 上也是有限型的。于是由
（\hyperref[I.10.6.4-fr]{10.6.4}），$\mathfrak{X}$ 是 Noether 的。若
$B=\varprojlim B_n$，则 $\mathfrak{X}=\operatorname{Spf}(B)$；若
$\mathfrak{K}$ 是 $B\to B_0$ 的核，则
$B_n=B/\mathfrak{K}^{n+1}$。若
$\rho_n:A/\mathfrak{J}^{n+1}\to B/\mathfrak{K}^{n+1}$ 是与 $f_n$ 对应的同态，则
\[
  \mathfrak{K}/\mathfrak{K}^{n+1}
  =(B/\mathfrak{K}^{n+1})\rho_n(\mathfrak{J}/\mathfrak{J}^{n+1})
\]
由于与 $f$ 对应的同态 $\rho:A\to B$ 等于
$\varprojlim\rho_n$，$B$ 的理想 $\mathfrak{J}B$ 在 $\mathfrak{K}$ 中稠密；而
$B$ 的每个理想都是闭的
（\hyperref[0.7.3.5-fr]{0, 7.3.5}），故
$\mathfrak{K}=\mathfrak{J}B$。若 $\mathscr{K}=\mathfrak{K}^\Delta$，则关系
$f^*(\mathscr{J})\mathscr{O}_{\mathfrak{X}}=\mathscr{K}$ 由
（\hyperref[I.10.10.9-fr]{10.10.9}）得出，证明完毕。

\begin{env}[10.12.3.2]
\label{I.10.12.3.2-fr}
对于两个进制 $\mathfrak{S}$-预概形 $\mathfrak{X}$、$\mathfrak{Y}$，上述等价给出一个
\emph{典范双射}
\[
  \operatorname{Hom}_{\mathfrak{S}}(\mathfrak{X},\mathfrak{Y})
  \xrightarrow{\sim}
  \varprojlim_n\operatorname{Hom}_{S_n}(X_n,Y_n)
\]
这里的逆极限是针对当 $m\leq n$ 时的映射
$u_n\to(u_n)_{(S_m)}$ 而取的。
\end{env}

\subsection{有限型态射}
\label{subsection:I.10.13-fr}

\begin{proposition}[10.13.1]
\label{I.10.13.1-fr}
设 $\mathfrak{Y}$ 是局部 Noether 形式预概形，$\mathscr{K}$ 是
$\mathfrak{Y}$ 的一个定义理想，
$f:\mathfrak{X}\to\mathfrak{Y}$ 是形式预概形的一个态射。下列条件等价：
\begin{enumerate}
  \item[a)] $\mathfrak{X}$ 是局部 Noether 的，$f$ 是进制态射
    （\hyperref[I.10.12.1-fr]{10.12.1}），并且若令
    $\mathscr{J}=f^*(\mathscr{K})\mathscr{O}_{\mathfrak{X}}$，则由 $f$ 导出的态射
    $f_0:(\mathfrak{X},\mathscr{O}_{\mathfrak{X}}/\mathscr{J})
    \to(\mathfrak{Y},\mathscr{O}_{\mathfrak{Y}}/\mathscr{K})$ 是有限型的。
  \item[b)] $\mathfrak{X}$ 是局部 Noether 的，并且是一个进制
    $(Y_n)$-归纳系 $(X_n)$ 的归纳极限，其中态射 $X_0\to Y_0$ 是有限型的。

\oldpage[I]{208}

  \item[c)] $\mathfrak{Y}$ 的每一点都有一个 Noether 仿射形式开邻域
    $V$，它具有如下性质：
    \begin{enumerate}
      \item[(\textbf{Q})] $f^{-1}(V)$ 是有限多个 Noether 仿射形式开集
        $U_i$ 的并，并且 Noether 进制环
        $\Gamma(U_i,\mathscr{O}_{\mathfrak{X}})$ 在拓扑上同构于
        $\Gamma(V,\mathscr{O}_{\mathfrak{Y}})$ 上一个设限形式幂级数代数
        （\hyperref[0.7.5.1-fr]{0, 7.5.1}）模去一个理想（该理想必为闭的）所得的商。
    \end{enumerate}
\end{enumerate}
\end{proposition}

由（\hyperref[I.10.12.3-fr]{10.12.3}），a) 蕴含 b) 是立即的。为证明
b) 蕴含 c)，由于问题在 $\mathfrak{Y}$ 上是局部的，可假设
$\mathfrak{Y}=\operatorname{Spf}(B)$，其中 $B$ 是 Noether 进制环；令
$\mathscr{K}=\mathfrak{K}^\Delta$，其中 $\mathfrak{K}$ 是 $B$ 的一个定义理想。
由假设，$X_0$ 在 $Y_0$ 上是有限型的，所以 $X_0$ 是有限多个仿射开集
$U_i$ 的并，使得 $X_0$ 在 $U_i$ 上诱导的仿射概形的环 $A_{i0}$，是
$Y_0$ 的环 $B/\mathfrak{K}$ 上的有限型代数
（\hyperref[I.6.3.2-fr]{6.3.2}）。由
（\hyperref[I.5.1.9-fr]{5.1.9}），$U_i$ 也是每个 Noether 预概形
$X_n$ 中的仿射开集；若 $A_{in}$ 是 $X_n$ 在 $U_i$ 上诱导的仿射概形的环，
则假设 b) 蕴含：当 $m\leq n$ 时，$A_{im}$ 同构于
$A_{in}/\mathfrak{K}^{m+1}A_{in}$。因此，$\mathfrak{X}$ 在 $U_i$ 上诱导的
形式预概形同构于 $\operatorname{Spf}(A_i)$，其中
$A_i=\varprojlim_n A_{in}$（\hyperref[I.10.6.4-fr]{10.6.4}）；
$A_i$ 是 $\mathfrak{K}A_i$-进制环，而
$A_i/\mathfrak{K}A_i$ 同构于 $A_{i0}$，故是 $B/\mathfrak{K}$ 上的有限型代数。
由此从（\hyperref[0.7.5.5-fr]{0, 7.5.5}）得出：$A_i$ 在拓扑上同构于
$B$ 上一个设限形式幂级数代数的商（模去一个必为闭的理想，因为这样的代数是
Noether 的（\hyperref[0.7.5.4-fr]{0, 7.5.4}））。

为证明 c) 蕴含 a)，只须考虑 $\mathfrak{X}=\operatorname{Spf}(A)$ 也为仿射的情形；
这里 $A$ 是 Noether 进制环，同构于 $B$ 上一个设限形式幂级数代数模去一个闭理想所得的商。
于是由（\hyperref[0.7.5.5-fr]{0, 7.5.5}），
$A/\mathfrak{K}A$ 是 $B/\mathfrak{K}$ 上的有限型代数，并且
$\mathfrak{K}A=\mathfrak{J}$ 是 $A$ 的一个定义理想；故由
（\hyperref[I.10.10.9-fr]{10.10.9}），条件 a) 得到满足。

应当指出，若命题（\hyperref[I.10.13.1-fr]{10.13.1}）的条件满足，则性质 a)
对 $\mathfrak{Y}$ 的\emph{每个定义理想 $\mathscr{K}$} 都成立（由 c)），
因而在性质 b) 中，\emph{所有 $f_n$ 都是有限型态射}。

\begin{corollary}[10.13.2]
\label{I.10.13.2-fr}
若（\hyperref[I.10.13.1-fr]{10.13.1}）的条件满足，则
$\mathfrak{Y}$ 的每个 Noether 仿射形式开集 $V$ 都具有性质
（\textbf{Q}）；若 $\mathfrak{Y}$ 是 Noether 的，则 $\mathfrak{X}$ 也是 Noether 的。
\end{corollary}

这立即由（\hyperref[I.10.13.1-fr]{10.13.1}）和
（\hyperref[I.6.3.2-fr]{6.3.2}）得出。

\begin{definition}[10.13.3]
\label{I.10.13.3-fr}
若（\hyperref[I.10.13.1-fr]{10.13.1}）的等价性质 a)、b)、c) 满足，则称态射
$f$ 是\emph{有限型的}，或者称 $\mathfrak{X}$ 是一个
\emph{有限型 $\mathfrak{Y}$-形式预概形}，或称它是一个
\emph{$\mathfrak{Y}$ 上的有限型形式预概形}。
\end{definition}

\begin{corollary}[10.13.4]
\label{I.10.13.4-fr}
设 $\mathfrak{X}=\operatorname{Spf}(A)$、
$\mathfrak{Y}=\operatorname{Spf}(B)$ 是两个 Noether 仿射形式概形。
$\mathfrak{X}$ 在 $\mathfrak{Y}$ 上为有限型的充要条件是，Noether 进制环
$A$ 同构于 $B$ 上一个设限形式幂级数代数模去一个闭理想所得的商。
\end{corollary}

事实上，采用（\hyperref[I.10.13.1-fr]{10.13.1}）的记号，若
$\mathfrak{X}$ 在 $\mathfrak{Y}$ 上是有限型的，则由
（\hyperref[I.6.3.3-fr]{6.3.3}），$A/\mathfrak{K}A$ 是有限型
$(B/\mathfrak{K})$-代数，而 $\mathfrak{K}A$ 是 $A$ 的一个定义理想
（\hyperref[I.10.10.9-fr]{10.10.9}）。于是结论由
（\hyperref[0.7.5.5-fr]{0, 7.5.5}）得出。

\begin{proposition}[10.13.5]
\label{I.10.13.5-fr}
\medskip\noindent
\begin{enumerate}
  \item[\rm{(i)}] 两个有限型形式预概形态射的复合是有限型的。

\oldpage[I]{209}

  \item[\rm{(ii)}] 设 $\mathfrak{X}$、$\mathfrak{S}$、
    $\mathfrak{S}'$ 是三个局部 Noether（相应地，Noether）形式预概形，
    $f:\mathfrak{X}\to\mathfrak{S}$、
    $g:\mathfrak{S}'\to\mathfrak{S}$ 是两个态射。若 $f$ 是有限型的，则
    $\mathfrak{X}\times_{\mathfrak{S}}\mathfrak{S}'$ 是局部 Noether 的
    （相应地，Noether 的），并且在 $\mathfrak{S}'$ 上是有限型的。
  \item[\rm{(iii)}] 设 $\mathfrak{S}$ 是局部 Noether 形式预概形，
    $\mathfrak{X}'$、$\mathfrak{Y}'$ 是两个局部 Noether
    $\mathfrak{S}$-形式预概形，使得
    $\mathfrak{X}'\times_{\mathfrak{S}}\mathfrak{Y}'$ 是局部 Noether 的。
    若 $\mathfrak{X}$、$\mathfrak{Y}$ 是局部 Noether
    $\mathfrak{S}$-形式预概形，
    $f:\mathfrak{X}\to\mathfrak{X}'$、
    $g:\mathfrak{Y}\to\mathfrak{Y}'$ 是两个有限型
    $\mathfrak{S}$-态射，则
    $\mathfrak{X}\times_{\mathfrak{S}}\mathfrak{Y}$ 是局部 Noether 的，并且
    $f\times_{\mathfrak{S}}g$ 是一个有限型 $\mathfrak{S}$-态射。
\end{enumerate}
\end{proposition}

(iii) 由（\hyperref[I.3.5.1-fr]{3.5.1}）中的形式论证从 (i) 和 (ii) 得出，
所以只须证明 (i) 和 (ii)。

设 $\mathfrak{X}$、$\mathfrak{Y}$、$\mathfrak{Z}$ 是三个局部 Noether
形式预概形，$f:\mathfrak{X}\to\mathfrak{Y}$、
$g:\mathfrak{Y}\to\mathfrak{Z}$ 是两个有限型态射。若
$\mathscr{L}$ 是 $\mathfrak{Z}$ 的一个定义理想，则
$\mathscr{K}=g^*(\mathscr{L})\mathscr{O}_{\mathfrak{Y}}$ 是
$\mathfrak{Y}$ 的一个定义理想，而
$\mathscr{J}=f^*(g^*(\mathscr{L}))\mathscr{O}_{\mathfrak{X}}$ 是
$\mathfrak{X}$ 的一个定义理想。令
$X_0=(\mathfrak{X},\mathscr{O}_{\mathfrak{X}}/\mathscr{J})$、
$Y_0=(\mathfrak{Y},\mathscr{O}_{\mathfrak{Y}}/\mathscr{K})$、
$Z_0=(\mathfrak{Z},\mathscr{O}_{\mathfrak{Z}}/\mathscr{L})$，并设
$f_0:X_0\to Y_0$、$g_0:Y_0\to Z_0$ 是与 $f$、$g$ 对应的态射。
由假设，$f_0$ 和 $g_0$ 都是有限型的，故
$g_0\circ f_0$ 也是有限型的（\hyperref[I.6.3.4-fr]{6.3.4}）；它对应于
$g\circ f$。所以由（\hyperref[I.10.13.1-fr]{10.13.1}），
$g\circ f$ 是有限型的。

在 (ii) 的条件下，$\mathfrak{S}$（相应地，$\mathfrak{X}$、
$\mathfrak{S}'$）是局部 Noether 预概形序列 $(S_n)$
（相应地，$(X_n)$、$(S'_n)$）的归纳极限；由
（\hyperref[I.10.13.1-fr]{10.13.1}），可假设当 $m\leq n$ 时
$X_m=X_n\times_{S_n}S_m$。于是形式预概形
$\mathfrak{X}\times_{\mathfrak{S}}\mathfrak{S}'$ 是各预概形
$X_n\times_{S_n}S'_n$ 的归纳极限
（\hyperref[I.10.7.4-fr]{10.7.4}），并且有
\[
  X_m\times_{S_m}S'_m
  =(X_n\times_{S_n}S_m)\times_{S_m}S'_m
  =(X_n\times_{S_n}S'_n)\times_{S'_n}S'_m.
\]
此外，$X_0$ 在 $S_0$ 上是有限型的，所以
$X_0\times_{S_0}S'_0$ 是局部 Noether 的
（\hyperref[I.6.3.8-fr]{6.3.8}）。由此先从
（\hyperref[I.10.12.3.1-fr]{10.12.3.1}）得出
$\mathfrak{X}\times_{\mathfrak{S}}\mathfrak{S}'$ 是局部 Noether 的；又因为
$X_0\times_{S_0}S'_0$ 在 $S'_0$ 上是有限型的
（\hyperref[I.6.3.8-fr]{6.3.8}），所以由
（\hyperref[I.10.12.3.1-fr]{10.12.3.1}）和
（\hyperref[I.10.13.1-fr]{10.13.1}），
$\mathfrak{X}\times_{\mathfrak{S}}\mathfrak{S}'$ 在
$\mathfrak{S}'$ 上是有限型的，这就证明了 (ii)（关于 Noether 预概形的断言立即由
（\hyperref[I.6.3.8-fr]{6.3.8}）得出）。

\begin{corollary}[10.13.6]
\label{I.10.13.6-fr}
在（\hyperref[I.10.9.9-fr]{10.9.9}）的假设下，若 $f$ 是有限型态射，则它到完备化的延拓
$\widehat f$ 也是有限型的。
\end{corollary}

\subsection{形式预概形的闭子预概形}
\label{subsection:I.10.14-fr}

\begin{proposition}[10.14.1]
\label{I.10.14.1-fr}
设 $\mathfrak{X}$ 是局部 Noether 形式预概形，$\mathscr{A}$ 是
$\mathscr{O}_{\mathfrak{X}}$ 的一个凝聚理想层。若 $\mathfrak{Y}$ 是
$\mathscr{O}_{\mathfrak{X}}/\mathscr{A}$ 的（闭）支集，则拓扑赋环空间
$(\mathfrak{Y},(\mathscr{O}_{\mathfrak{X}}/\mathscr{A})|\mathfrak{Y})$ 是局部
Noether 形式预概形；若 $\mathfrak{X}$ 是 Noether 的，则它也是 Noether 的。
\end{proposition}

应当注意，由（\hyperref[I.10.10.3-fr]{10.10.3}）和
（\hyperref[0.5.3.4-fr]{0, 5.3.4}），
$\mathscr{O}_{\mathfrak{X}}/\mathscr{A}$ 是凝聚的，故它的支集
$\mathfrak{Y}$ 是闭的（\hyperref[0.5.2.2-fr]{0, 5.2.2}）。设
$\mathscr{J}$ 是 $\mathfrak{X}$ 的一个定义理想，并令
$X_n=(\mathfrak{X},\mathscr{O}_{\mathfrak{X}}/\mathscr{J}^{n+1})$。环层
$\mathscr{O}_{\mathfrak{X}}/\mathscr{A}$ 是各层
$\mathscr{O}_{\mathfrak{X}}/(\mathscr{A}+\mathscr{J}^{n+1})
=(\mathscr{O}_{\mathfrak{X}}/\mathscr{A})
\otimes_{\mathscr{O}_{\mathfrak{X}}}
(\mathscr{O}_{\mathfrak{X}}/\mathscr{J}^{n+1})$ 的逆极限
（\hyperref[I.10.11.3-fr]{10.11.3}），而这些层的支集都是
$\mathfrak{Y}$。层
$(\mathscr{A}+\mathscr{J}^{n+1})/\mathscr{J}^{n+1}$ 是凝聚
$\mathscr{O}_{\mathfrak{X}}$-模，因为 $\mathscr{J}^{n+1}$ 是凝聚的；所以
$(\mathscr{A}+\mathscr{J}^{n+1})/\mathscr{J}^{n+1}$ 也是凝聚
$(\mathscr{O}_{\mathfrak{X}}/\mathscr{J}^{n+1})$-模
（\hyperref[0.5.3.10-fr]{0, 5.3.10}）。若 $Y_n$ 是由这个理想层定义的
$X_n$ 的闭子预概形，则立即可见
$(\mathfrak{Y},(\mathscr{O}_{\mathfrak{X}}/\mathscr{A})|\mathfrak{Y})$

\oldpage[I]{210}

是各 $Y_n$ 的归纳极限形式预概形；而且由于
（\hyperref[I.10.6.4-fr]{10.6.4}）的条件满足，这就证明该形式预概形是局部
Noether 的；若 $\mathfrak{X}$ 是 Noether 的，它也是 Noether 的（因为此时由
（\hyperref[I.6.1.4-fr]{6.1.4}），$Y_0$ 是 Noether 的）。

\begin{definition}[10.14.2]
\label{I.10.14.2-fr}
% EGA1-ERR-0057: source-literal coherent-module wording retained; published EGA II errata says coherent ideal.
称形式预概形 $\mathfrak{X}$ 的任何如下形式预概形为它的\emph{闭子预概形}：
$(\mathfrak{Y},(\mathscr{O}_{\mathfrak{X}}/\mathscr{A})|\mathfrak{Y})$，其中
$\mathscr{A}$ 是一个凝聚 $\mathscr{O}_{\mathfrak{X}}$-模；称这个预概形为由
$\mathscr{A}$ 定义的闭子预概形。
\end{definition}

% EGA1-ERR-0057: source-literal coherent-module correspondence retained; published EGA II errata says coherent ideals.
显然，这样定义的凝聚 $\mathscr{O}_{\mathfrak{X}}$-模与
$\mathfrak{X}$ 的闭子预概形之间的对应是一一的。

拓扑赋环空间态射
$j=(\psi,\theta):\mathfrak{Y}\to\mathfrak{X}$，其中 $\psi$ 是嵌入
$\mathfrak{Y}\to\mathfrak{X}$，而 $\theta^\sharp$ 是典范同态
$\mathscr{O}_{\mathfrak{X}}\to\mathscr{O}_{\mathfrak{X}}/\mathscr{A}$；由
（\hyperref[I.10.4.5-fr]{10.4.5}），它显然是形式预概形态射，称为
$\mathfrak{Y}$ 到 $\mathfrak{X}$ 中的\emph{典范嵌入}。应当注意，若
$\mathfrak{X}=\operatorname{Spf}(A)$，其中 $A$ 是 Noether 进制环，则
$\mathscr{A}=\mathfrak{a}^\Delta$，其中 $\mathfrak{a}$ 是 $A$ 的一个理想
（\hyperref[I.10.10.5-fr]{10.10.5}）；由上文立即得出，此时在相差一个同构的意义下有
$\mathfrak{Y}=\operatorname{Spf}(A/\mathfrak{a})$，而 $j$ 对应于
（\hyperref[I.10.2.2-fr]{10.2.2}）典范同态
$A\to A/\mathfrak{a}$。

若局部 Noether 形式预概形的态射
$f:\mathfrak{Z}\to\mathfrak{X}$ 分解为
$\mathfrak{Z}\xrightarrow{g}\mathfrak{Y}\xrightarrow{j}\mathfrak{X}$，其中 $g$ 是
$\mathfrak{Z}$ 到 $\mathfrak{X}$ 的一个闭子预概形 $\mathfrak{Y}$ 上的同构，
$j$ 是典范嵌入，则称这个态射为\emph{闭浸入}。由于 $j$ 是赋环空间的单态射，
$g$ 和 $\mathfrak{Y}$ 必然是\emph{唯一的}。

\begin{proposition}[10.14.3]
\label{I.10.14.3-fr}
闭浸入是有限型态射。
\end{proposition}

立即可归结到 $\mathfrak{X}$ 是仿射形式概形
$\operatorname{Spf}(A)$，并且
$\mathfrak{Y}=\operatorname{Spf}(A/\mathfrak{a})$ 的情形；命题由
（\hyperref[I.10.13.1-fr]{10.13.1, c)}）得出。

\begin{lemma}[10.14.4]
\label{I.10.14.4-fr}
设 $f:\mathfrak{Y}\to\mathfrak{X}$ 是局部 Noether 形式预概形的态射，
$(U_\alpha)$ 是 $f(\mathfrak{Y})$ 的一个覆盖，其成员是
$\mathfrak{X}$ 的 Noether 仿射形式开集，并且各
$f^{-1}(U_\alpha)$ 都是 $\mathfrak{Y}$ 的 Noether 仿射形式开集。态射
$f$ 是闭浸入的充要条件是：$f(\mathfrak{Y})$ 是 $\mathfrak{X}$ 的闭子集，并且对每个
$\alpha$，$f$ 在 $f^{-1}(U_\alpha)$ 上的限制对应于
（\hyperref[I.10.4.6-fr]{10.4.6}）满射同态
$\Gamma(U_\alpha,\mathscr{O}_{\mathfrak{X}})\to
\Gamma(f^{-1}(U_\alpha),\mathscr{O}_{\mathfrak{Y}})$。
\end{lemma}

这些条件显然是必要的。反之，记
$\mathfrak{a}_\alpha$ 为
$\Gamma(U_\alpha,\mathscr{O}_{\mathfrak{X}})\to
\Gamma(f^{-1}(U_\alpha),\mathscr{O}_{\mathfrak{Y}})$ 的核，则可以如下定义
$\mathscr{O}_{\mathfrak{X}}$ 的一个凝聚理想层 $\mathscr{A}$：令
$\mathscr{A}|U_\alpha=\mathfrak{a}_\alpha^\Delta$，并且
% EGA1-ERR-0058: source-literal zero-on-complement rule retained; the ideal must be the unit ideal there.
在各 $U_\alpha$ 的并集的补集中把 $\mathscr{A}$ 取为零。事实上，由于
$f(\mathfrak{Y})$ 是闭的，并且
% EGA1-ERR-0059: source-literal support of the ideal retained; the quotient sheaf has this support.
$\mathfrak{a}_\alpha^\Delta$ 的支集是
$U_\alpha\cap f(\mathfrak{Y})$，问题归结为验证：在任一 Noether 仿射形式开集
$V\subset U_\alpha\cap U_\beta$ 上，$\mathfrak{a}_\alpha^\Delta$ 和
$\mathfrak{a}_\beta^\Delta$ 诱导相同的层。由于 $f$ 在
$f^{-1}(U_\alpha)$ 上的限制，是该形式预概形到 $U_\alpha$ 中的闭浸入，
$f^{-1}(V)$ 是 $f^{-1}(U_\alpha)$ 的 Noether 仿射形式开集，而且 $f$ 在
$f^{-1}(V)$ 上的限制是闭浸入；若 $\mathfrak{b}$ 是与这个限制对应的满射同态
$\Gamma(V,\mathscr{O}_{\mathfrak{X}})\to
\Gamma(f^{-1}(V),\mathscr{O}_{\mathfrak{Y}})$ 的核，则由
（\hyperref[I.10.10.2-fr]{10.10.2}）立即可知，
$\mathfrak{a}_\alpha^\Delta$ 在 $V$ 上诱导 $\mathfrak{b}^\Delta$。这样定义了理想层
$\mathscr{A}$ 之后，显然有
% EGA1-ERR-0060: source-literal composite retained; the type-correct order is j after g.
$f=g\circ j$，其中 $j:\mathfrak{Z}\to\mathfrak{X}$ 是由 $\mathscr{A}$ 定义的
$\mathfrak{X}$ 的闭子预概形 $\mathfrak{Z}$ 的典范嵌入，而 $g$ 是
$\mathfrak{Y}$ 到 $\mathfrak{Z}$ 上的同构。

\begin{proposition}[10.14.5]
\label{I.10.14.5-fr}
\begin{enumerate}
  \item[\rm{(i)}] 若 $f:\mathfrak{Z}\to\mathfrak{Y}$、
  $g:\mathfrak{Y}\to\mathfrak{X}$ 是局部 Noether 形式预概形的闭浸入，则
  $g\circ f$ 是闭浸入。

\oldpage[I]{211}
  \item[\rm{(ii)}] 设 $\mathfrak{X}$、$\mathfrak{Y}$、$\mathfrak{S}$ 是三个
  局部 Noether 形式预概形，$f:\mathfrak{X}\to\mathfrak{S}$ 是闭浸入，
  $g:\mathfrak{Y}\to\mathfrak{S}$ 是态射。那么态射
  $\mathfrak{X}\times_{\mathfrak{S}}\mathfrak{Y}\to\mathfrak{Y}$ 是闭浸入。
  \item[\rm{(iii)}] 设 $\mathfrak{S}$ 是局部 Noether 形式预概形，
  $\mathfrak{X}'$、$\mathfrak{Y}'$ 是两个局部 Noether
  $\mathfrak{S}$-形式预概形，使得
  $\mathfrak{X}'\times_{\mathfrak{S}}\mathfrak{Y}'$ 是局部 Noether 的。
  若 $\mathfrak{X}$、$\mathfrak{Y}$ 是局部 Noether
  $\mathfrak{S}$-形式预概形，
  $f:\mathfrak{X}\to\mathfrak{X}'$、$g:\mathfrak{Y}\to\mathfrak{Y}'$ 是两个
  本身为闭浸入的 $\mathfrak{S}$-态射，则
  $f\times_{\mathfrak{S}}g$ 是闭浸入。
\end{enumerate}
\end{proposition}

由（\hyperref[I.3.5.1-fr]{3.5.1}），仍只须证明 (i) 和 (ii)。

为证明 (i)，可假设 $\mathfrak{Y}$（相应地，$\mathfrak{Z}$）是
$\mathfrak{X}$（相应地，$\mathfrak{Y}$）的闭子预概形，由
$\mathscr{O}_{\mathfrak{X}}$（相应地，$\mathscr{O}_{\mathfrak{Y}}$）的凝聚理想层
$\mathscr{J}$（相应地，$\mathscr{K}$）定义。若 $\psi$ 是底空间的嵌入
$\mathfrak{Y}\to\mathfrak{X}$，则 $\psi_*(\mathscr{K})$ 是
$\psi_*(\mathscr{O}_{\mathfrak{Y}})
=\mathscr{O}_{\mathfrak{X}}/\mathscr{J}$ 的一个凝聚理想层
（\hyperref[0.5.3.12-fr]{0, 5.3.12}），从而也是凝聚
$\mathscr{O}_{\mathfrak{X}}$-模
（\hyperref[0.5.3.10-fr]{0, 5.3.10}）。所以同态
$\mathscr{O}_{\mathfrak{X}}\to
(\mathscr{O}_{\mathfrak{X}}/\mathscr{J})/\psi_*(\mathscr{K})$ 的核
$\mathscr{K}_1$ 是 $\mathscr{O}_{\mathfrak{X}}$ 的一个凝聚理想层
（\hyperref[0.5.3.4-fr]{0, 5.3.4}），并且
$\mathscr{O}_{\mathfrak{X}}/\mathscr{K}_1$ 同构于
$\psi_*(\mathscr{O}_{\mathfrak{Y}}/\mathscr{K})$；这证明
$\mathfrak{Z}$ 同构于 $\mathfrak{X}$ 的一个闭子预概形。

为证明 (ii)，显然只须考虑
$\mathfrak{S}=\operatorname{Spf}(A)$、
$\mathfrak{X}=\operatorname{Spf}(B)$、
$\mathfrak{Y}=\operatorname{Spf}(C)$ 的情形，其中 $A$ 是 Noether
$\mathfrak{J}$-进制环，$B=A/\mathfrak{a}$，$\mathfrak{a}$ 是 $A$ 的一个理想，
$C$ 是 Noether 进制拓扑 $A$-代数。问题归结为证明同态
$C\to C\widehat{\otimes}_A(A/\mathfrak{a})$ 是\emph{满射的}。然而，
$A/\mathfrak{a}$ 是有限型 $A$-模，它的拓扑是
$\mathfrak{J}$-进制拓扑；于是由
（\hyperref[0.7.7.8-fr]{0, 7.7.8}），
$C\widehat{\otimes}_A(A/\mathfrak{a})$ 等同于
$C\otimes_A(A/\mathfrak{a})=C/\mathfrak{a}C$，断言得证。

\begin{corollary}[10.14.6]
\label{I.10.14.6-fr}
在（\hyperref[I.10.14.5-fr]{10.14.5, (ii)}）的假设下，设
$p:\mathfrak{X}\times_{\mathfrak{S}}\mathfrak{Y}\to\mathfrak{X}$、
$q:\mathfrak{X}\times_{\mathfrak{S}}\mathfrak{Y}\to\mathfrak{Y}$ 是投影，
使得图表
\[
  \xymatrix{
    \mathfrak{X} \ar[d]_{f}
    & \mathfrak{X}\times_{\mathfrak{S}}\mathfrak{Y}
      \ar[l]_{p} \ar[d]^{q}\\
    \mathfrak{S}
    & \mathfrak{Y} \ar[l]^{g}
  }
\]
交换。于是，对每个凝聚 $\mathscr{O}_{\mathfrak{X}}$-模
$\mathscr{F}$，有一个 $\mathscr{O}_{\mathfrak{Y}}$-模的典范同构
\[
  \label{I.10.14.6.1-fr}
  u:g^*(f_*(\mathscr{F}))
  \xrightarrow{\sim}
  q_*(p^*(\mathscr{F})).
  \tag{10.14.6.1}
\]
\end{corollary}

为了定义同态
$g^*(f_*(\mathscr{F}))\to q_*(p^*(\mathscr{F}))$，已知这等价于定义同态
$f_*(\mathscr{F})\to g_*(q_*(p^*(\mathscr{F})))
=f_*(p_*(p^*(\mathscr{F})))$
（\hyperref[0.4.4.3-fr]{0, 4.4.3}）。
% EGA1-ERR-0061: source-literal equality retained; f_*(rho) is the adjoint correspondent of u, not u itself.
我们取 $u=f_*(\rho)$，其中 $\rho$ 是典范同态
$\mathscr{F}\to p_*(p^*(\mathscr{F}))$
（\hyperref[0.4.4.3-fr]{0, 4.4.3}）。为说明 $u$ 是同构，立即归结到
$\mathfrak{S}$、$\mathfrak{X}$、$\mathfrak{Y}$ 分别是 Noether 进制环
$A$、$B$、$C$ 的形式谱，并且满足上文
（\hyperref[I.10.14.5-fr]{10.14.5, (ii)}）中的条件的情形。此时
$\mathscr{F}=M^\Delta$，其中 $M$ 是有限型
$(A/\mathfrak{a})$-模
（\hyperref[I.10.10.5-fr]{10.10.5}）；
（\hyperref[I.10.14.6.1-fr]{10.14.6.1}）的两端分别由
（\hyperref[I.10.10.8-fr]{10.10.8}）等同于
$(C\otimes_A M)^\Delta$ 和
$((C/\mathfrak{a}C)\otimes_{A/\mathfrak{a}}M)^\Delta$，于是推论成立，因为
$(C/\mathfrak{a}C)\otimes_{A/\mathfrak{a}}M
=(C\otimes_A(A/\mathfrak{a}))\otimes_{A/\mathfrak{a}}M$ 典范地等同于
$C\otimes_A M$。

\oldpage[I]{212}
\begin{corollary}[10.14.7]
\label{I.10.14.7-fr}
设 $X$ 是局部 Noether 通常预概形，$Y$ 是 $X$ 的一个闭子预概形，
$j$ 是典范嵌入 $Y\to X$，$X'$ 是 $X$ 的一个闭子集，并令
$Y'=Y\cap X'$。那么
$\widehat{j}:Y_{/Y'}\to X_{/X'}$ 是闭浸入；并且对每个凝聚
$\mathscr{O}_Y$-模 $\mathscr{F}$，有
\[
  \widehat{j}_*(\mathscr{F}_{/Y'})=(j_*(\mathscr{F}))_{/X'}.
\]
\end{corollary}

因为 $Y'=j^{-1}(X')$，只须使用
（\hyperref[I.10.9.9-fr]{10.9.9}），并应用
（\hyperref[I.10.14.5-fr]{10.14.5}）和
（\hyperref[I.10.14.6-fr]{10.14.6}）。

\subsection{分离形式预概形。}
\label{subsection:I.10.15-fr}

\begin{definition}[10.15.1]
\label{I.10.15.1-fr}
设 $\mathfrak{S}$ 是一个形式预概形，$\mathfrak{X}$ 是一个
$\mathfrak{S}$-形式预概形，
$f:\mathfrak{X}\to\mathfrak{S}$ 是结构态射。称对角态射
$\Delta_{\mathfrak{X}|\mathfrak{S}}:\mathfrak{X}\to
\mathfrak{X}\times_{\mathfrak{S}}\mathfrak{X}$（也记作
$\Delta_{\mathfrak{X}}$）为态射
$(1_{\mathfrak{X}},1_{\mathfrak{X}})_{\mathfrak{S}}$。若
$\Delta_{\mathfrak{X}}$ 将 $\mathfrak{X}$ 的底空间映为
$\mathfrak{X}\times_{\mathfrak{S}}\mathfrak{X}$ 的底空间中的一个闭子集，
则称 $\mathfrak{X}$ 在 $\mathfrak{S}$ 上分离，或称
它是一个 $\mathfrak{S}$-形式概形，或称 $f$ 是一个
分离态射。若形式预概形 $\mathfrak{X}$ 在 $\mathbf{Z}$ 上分离，
则称它是分离的，或称它是一个形式概形。
\end{definition}

\begin{proposition}[10.15.2]
\label{I.10.15.2-fr}
设形式预概形 $\mathfrak{S}$、$\mathfrak{X}$ 分别是通常预概形序列
$(S_n)$、$(X_n)$ 的归纳极限，并且态射
$f:\mathfrak{X}\to\mathfrak{S}$ 是态射序列 $f_n:X_n\to S_n$ 的归纳极限。
那么，$f$ 分离的充要条件是态射 $f_0:X_0\to S_0$ 分离。
\end{proposition}

事实上，此时 $\Delta_{\mathfrak{X}|\mathfrak{S}}$ 是态射序列
$\Delta_{X_n|S_n}$ 的归纳极限
（\hyperref[I.10.7.4-fr]{10.7.4}），并且
$\Delta_{\mathfrak{X}|\mathfrak{S}}$ 所映出的 $\mathfrak{X}$ 底空间
（相应地，$\mathfrak{X}\times_{\mathfrak{S}}\mathfrak{X}$ 的底空间）
与 $\Delta_{X_0|S_0}$ 所映出的 $X_0$ 底空间
（相应地，$X_0\times_{S_0}X_0$ 的底空间）相同；结论随即得出。

\begin{proposition}[10.15.3]
\label{I.10.15.3-fr}
以下假设所考虑的所有形式预概形（相应地，形式预概形的态射）都是通常预概形
序列（相应地，通常预概形态射序列）的归纳极限。
\begin{enumerate}
  \item[\rm{(i)}] 两个分离态射的复合是分离的。
  \item[\rm{(ii)}] 若 $f:\mathfrak{X}\to\mathfrak{X}'$、
  $g:\mathfrak{Y}\to\mathfrak{Y}'$ 是两个分离的 $\mathfrak{S}$-态射，
  则 $f\times_{\mathfrak{S}}g$ 分离。
  \item[\rm{(iii)}] 若 $f:\mathfrak{X}\to\mathfrak{Y}$ 是一个分离的
  $\mathfrak{S}$-态射，则对于基形式预概形的每个扩张
  $\mathfrak{S}'\to\mathfrak{S}$，$\mathfrak{S}'$-态射
  $f_{(\mathfrak{S}')}$ 都分离。
  \item[\rm{(iv)}] 若两个态射的复合 $g\circ f$ 分离，则 $f$ 分离。
\end{enumerate}
\emph{（这里约定：若同一个形式预概形 $\mathfrak{Z}$ 在同一命题中出现多次，
则在它出现的每一处，都把它视为\emph{同一个}通常预概形序列 $(Z_n)$ 的归纳极限；
并把从 $\mathfrak{Z}$ 到一个形式预概形的态射（相应地，从一个形式预概形到
$\mathfrak{Z}$ 的态射），视为从 $Z_n$ 到通常预概形的态射
（相应地，从通常预概形到 $Z_n$ 的态射）的归纳极限。）}
\end{proposition}

采用（\hyperref[I.10.15.2-fr]{10.15.2}）的记号，事实上有
$(g\circ f)_0=g_0\circ f_0$ 和
$(f\times_{\mathfrak{S}}g)_0=f_0\times_{S_0}g_0$；于是
（\hyperref[I.10.15.3-fr]{10.15.3}）中的断言立即由
（\hyperref[I.10.15.2-fr]{10.15.2}）以及关于通常预概形的相应断言
（\hyperref[I.5.5.1-fr]{5.5.1}）得出。

对于（\hyperref[I.10.15.3-fr]{10.15.3}）中同一类形式预概形和态射，
我们把表述与（\hyperref[I.5.5.5-fr]{5.5.5}）、

\oldpage[I]{213}
（\hyperref[I.5.5.9-fr]{5.5.9}）以及
（\hyperref[I.5.5.10-fr]{5.5.10}）相对应的命题留给读者
（其中把“仿射开集”替换为“满足
（\hyperref[I.10.6.3-fr]{10.6.3}）中条件 b) 的仿射形式开集”）。

类似的论证还表明，每个 \emph{Noether 仿射形式概形}都是分离的；这也说明了
上述术语的合理性。

\begin{proposition}[10.15.4]
\label{I.10.15.4-fr}
设 $\mathfrak{S}$ 是一个局部 Noether 形式预概形，
$\mathfrak{X}$、$\mathfrak{Y}$ 是两个局部 Noether
$\mathfrak{S}$-形式预概形；假设 $\mathfrak{X}$ 或 $\mathfrak{Y}$ 在
$\mathfrak{S}$ 上为有限型（从而
$\mathfrak{X}\times_{\mathfrak{S}}\mathfrak{Y}$ 是局部 Noether 的），
并且 $\mathfrak{Y}$ 在 $\mathfrak{S}$ 上分离。设
$f:\mathfrak{X}\to\mathfrak{Y}$ 是一个 $\mathfrak{S}$-态射；那么图态射
$\Gamma_f=(1_{\mathfrak{X}},f)_{\mathfrak{S}}:\mathfrak{X}\to
\mathfrak{X}\times_{\mathfrak{S}}\mathfrak{Y}$ 是一个闭浸入。
\end{proposition}

可以假设 $\mathfrak{S}$ 是局部 Noether 预概形序列 $(S_n)$ 的归纳极限，
$\mathfrak{X}$（相应地，$\mathfrak{Y}$）是 $S_n$-预概形序列
$(X_n)$（相应地，$(Y_n)$）的归纳极限，$f$ 是 $S_n$-态射序列
$f_n:X_n\to Y_n$ 的归纳极限；于是
$\mathfrak{X}\times_{\mathfrak{S}}\mathfrak{Y}$ 是序列
$(X_n\times_{S_n}Y_n)$ 的归纳极限，而 $\Gamma_f$ 是序列
$(\Gamma_{f_n})$ 的归纳极限（\hyperref[I.10.7.4-fr]{10.7.4}）。
由假设，$Y_0$ 在 $S_0$ 上分离
（\hyperref[I.10.15.2-fr]{10.15.2}），所以空间
$\Gamma_{f_0}(X_0)$ 是 $X_0\times_{S_0}Y_0$ 的一个闭子空间；由于
$\mathfrak{X}\times_{\mathfrak{S}}\mathfrak{Y}$
（相应地，$\Gamma_f(\mathfrak{X})$）的底空间与
$X_0\times_{S_0}Y_0$（相应地，$\Gamma_{f_0}(X_0)$）相同，
已经可见 $\Gamma_f(\mathfrak{X})$ 是
$\mathfrak{X}\times_{\mathfrak{S}}\mathfrak{Y}$ 的一个\emph{闭}子空间。
现在注意，当 $(U,V)$ 遍历由 $\mathfrak{X}$（相应地，$\mathfrak{Y}$）
的 Noether 仿射形式开集 $U$（相应地，$V$）所组成且满足
$f(U)\subset V$ 的所有序偶时，这些开集
$U\times_{\mathfrak{S}}V$ 在
$\mathfrak{X}\times_{\mathfrak{S}}\mathfrak{Y}$ 中构成
$\Gamma_f(\mathfrak{X})$ 的一个覆盖；并且，若
$f_U:U\to V$ 是 $f$ 在 $U$ 上的限制，则
$\Gamma_{f_U}:U\to U\times_{\mathfrak{S}}V$ 是 $\Gamma_f$ 在 $U$ 上的限制。
因此，只要证明 $\Gamma_{f_U}$ 是闭浸入，就可由
（\hyperref[I.10.14.4-fr]{10.14.4}）得出 $\Gamma_f$ 也是闭浸入。
换言之，问题归结到
$\mathfrak{S}=\operatorname{Spf}(A)$、
$\mathfrak{X}=\operatorname{Spf}(B)$、
$\mathfrak{Y}=\operatorname{Spf}(C)$ 都是仿射的情形（$A$、$B$、$C$
为 Noether 进制环），且 $f$ 对应于一个连续 $A$-同态
$\varphi:C\to B$。此时 $\Gamma_f$ 对应于唯一的连续同态
$\omega:B\widehat{\otimes}_A C\to B$；它与典范同态
$B\to B\widehat{\otimes}_A C$ 和
$C\to B\widehat{\otimes}_A C$ 复合后，分别给出恒等同态与
$\varphi$。显然，$\omega$ 是\emph{满射的}，断言得证。

\begin{corollary}[10.15.5]
\label{I.10.15.5-fr}
设 $\mathfrak{S}$ 是一个局部 Noether 形式预概形，
$\mathfrak{X}$ 是一个有限型 $\mathfrak{S}$-形式预概形；那么，
$\mathfrak{X}$ 在 $\mathfrak{S}$ 上分离的充要条件是对角态射
$\mathfrak{X}\to\mathfrak{X}\times_{\mathfrak{S}}\mathfrak{X}$ 是闭浸入。
\end{corollary}

\begin{proposition}[10.15.6]
\label{I.10.15.6-fr}
局部 Noether 形式预概形的闭浸入
$j:\mathfrak{Y}\to\mathfrak{X}$ 是一个分离态射。
\end{proposition}

采用（\hyperref[I.10.14.2-fr]{10.14.2}）的记号，
$j_0:Y_0\to X_0$ 是闭浸入，因而是分离态射；只须应用
（\hyperref[I.10.15.2-fr]{10.15.2}）。

\begin{proposition}[10.15.7]
\label{I.10.15.7-fr}
设 $X$ 是一个局部 Noether（通常的）预概形，$X'$ 是 $X$ 的一个闭子集，
并令 $\widehat{X}=X_{/X'}$。那么，$\widehat{X}$ 分离的充要条件是
$\widehat{X}_{\mathrm{red}}$ 分离；而 $X$ 分离是一个充分条件。
\end{proposition}

事实上，采用（\hyperref[I.10.8.5-fr]{10.8.5}）的记号，
$\widehat{X}$ 分离的充要条件是 $X'_0$ 分离
（\hyperref[I.10.15.2-fr]{10.15.2}）；又因为
$\widehat{X}_{\mathrm{red}}=(X'_0)_{\mathrm{red}}$，这等价于说
$\widehat{X}_{\mathrm{red}}$ 分离
（\hyperref[I.5.5.1-fr]{5.5.1, (vi)}）。
\begin{thebibliography}{22}
\oldpage[I]{214}

\bibitem{I-1}
H.~\textsc{Cartan}、C.~\textsc{Chevalley},
\emph{Séminaire de l'École Normale Supérieure},
第 8 学年（1955--56），\emph{Géométrie algébrique}。

\bibitem{I-2}
H.~\textsc{Cartan}、S.~\textsc{Eilenberg},
\emph{Homological Algebra}, Princeton Math. Series (Princeton University
Press), 1956.

\bibitem{I-3}
W.~L.~\textsc{Chow}、J.~\textsc{Igusa},
\emph{Cohomology theory of varieties over rings},
\emph{Proc. Nat. Acad. Sci. U.S.A.}, t.~XLIV (1958), p.~1244--1248.

\bibitem{I-4}
R.~\textsc{Godement}, \emph{Théorie des faisceaux}, Actual. Scient. et Ind.,
n\textsuperscript{o}~1252, Paris (Hermann), 1958.

\bibitem{I-5}
H.~\textsc{Grauert},
\emph{Ein Theorem der analytischen Garbentheorie und die Modulräume komplexer
Strukturen}, \emph{Publ. Math. Inst. Hautes Études Scient.},
n\textsuperscript{o}~5, 1960.

\bibitem{I-6}
A.~\textsc{Grothendieck}, \emph{Sur quelques points d'algèbre homologique},
\emph{Tôhoku Math. Journ.}, t.~IX (1957), p.~119--221.

\bibitem{I-7}
A.~\textsc{Grothendieck},
\emph{Cohomology theory of abstract algebraic varieties},
\emph{Proc. Intern. Congress of Math.}, p.~103--118, Edinburgh (1958).

\bibitem{I-8}
A.~\textsc{Grothendieck},
\emph{Géométrie formelle et géométrie algébrique},
\emph{Séminaire Bourbaki}, 第 11 学年（1958--59），报告 182。

\bibitem{I-9}
M.~\textsc{Nagata},
\emph{A general theory of algebraic geometry over Dedekind domains},
\emph{Amer. Math. Journ.}: I, t.~LXXVIII, p.~78--116 (1956);
II, t.~LXXX, p.~382--420 (1958).

\bibitem{I-10}
D.~G.~\textsc{Northcott}, \emph{Ideal theory}, Cambridge Univ. Press, 1953.

\bibitem{I-11}
P.~\textsc{Samuel}, \emph{Commutative algebra} (D.~Herzig 记录), Cornell
Univ., 1953.

\bibitem{I-12}
P.~\textsc{Samuel}, \emph{Algèbre locale}, Mém. Sci. Math.,
n\textsuperscript{o}~123, Paris, 1953.

\bibitem{I-13}
P.~\textsc{Samuel}、O.~\textsc{Zariski},
\emph{Commutative algebra}, 2 卷, New York (Van Nostrand), 1958--60.

\bibitem{I-14}
J.-P.~\textsc{Serre}, \emph{Faisceaux algébriques cohérents},
\emph{Ann. of Math.}, t.~LXI (1955), p.~197--278.

\bibitem{I-15}
J.-P.~\textsc{Serre},
\emph{Sur la cohomologie des variétés algébriques},
\emph{Journ. de Math.} (9), t.~XXXVI (1957), p.~1--16.

\bibitem{I-16}
J.-P.~\textsc{Serre},
\emph{Géométrie algébrique et géométrie analytique},
\emph{Ann. Inst. Fourier}, t.~VI (1955--56), p.~1--42.

\bibitem{I-17}
J.-P.~\textsc{Serre},
\emph{Sur la dimension homologique des anneaux et des modules noethériens},
\emph{Proc. Intern. Symp. on Alg. Number theory}, p.~176--189, Tokyo--Nikko,
1955.

\bibitem{I-18}
A.~\textsc{Weil}, \emph{Foundations of algebraic geometry}, Amer. Math. Soc.
Coll. Publ., n\textsuperscript{o}~29, 1946.

\bibitem{I-19}
A.~\textsc{Weil}, \emph{Numbers of solutions of equations in finite fields},
\emph{Bull. Amer. Math. Soc.}, t.~LV (1949), p.~497--508.

\bibitem{I-20}
O.~\textsc{Zariski},
\emph{Theory and applications of holomorphic functions on algebraic varieties
over arbitrary ground fields}, Mem. Amer. Math. Soc., n\textsuperscript{o}~5
(1951).

\bibitem{I-21}
O.~\textsc{Zariski}, \emph{A new proof of Hilbert's Nullstellensatz},
\emph{Bull. Amer. Math. Soc.}, t.~LIII (1947), p.~362--368.

\bibitem{I-22}
E.~\textsc{Kähler}, \emph{Geometria Arithmetica},
\emph{Ann. di Mat.} (4), t.~XLV (1958), p.~1--368.

\end{thebibliography}
% EGA I 记号索引（大陆简体中文）。
% 法文 NUMDAM 印本第 217--218 页控制条目、次序和页界。
\clearpage
\oldpage[I]{217}
\section*{记号索引}
\phantomsection
\label{section:index-notation}
\begingroup
\small
\setlength{\columnsep}{2em}
\begin{multicols}{2}
\raggedcolumns
\newcommand{\EGAIndexLine}[2]{\par\noindent\hangindent=1em\hangafter=1 #1\ \dotfill\ \textup{#2}\par}

\phantomsection\label{section:index-notation-chapter0}
\medskip\noindent\textbf{第 0 章}\par\nopagebreak
\EGAIndexLine{$M_{[\phi]}$}{(0.1.0.2)}
\EGAIndexLine{$B\mathfrak{J}$, $\mathfrak{J}B$}{(0.1.0.3)}
\EGAIndexLine{$r(\mathfrak{a})$}{(0.1.1.1)}
\EGAIndexLine{$\mathfrak{R}(A)$}{(0.1.1.2)}
\EGAIndexLine{$S_{f}$}{(0.1.2.1)}
\EGAIndexLine{$S^{-1}A$, $S^{-1}M$, $m/s$, $i_A^S$, $i_M^S$, $i^S$}{(0.1.2.2)}
\EGAIndexLine{$A_{f}$, $M_{f}$, $A_{\mathfrak{p}}$, $M_{\mathfrak{p}}$}{(0.1.2.3)}
\EGAIndexLine{$S^{-1}u$}{(0.1.3.1)}
\EGAIndexLine{$\rho^{T,S}_{A}$, $\rho^{T,S}_{M}$, $\rho^{T,S}$}{(0.1.4.1)}
\EGAIndexLine{$\operatorname{Supp}(M)$}{(0.1.7.1)}
\EGAIndexLine{$\mathcal{F}|U$, $u|U$}{(0.3.1.5)}
\EGAIndexLine{$\mathcal{F}_{x}$, $s_{x}$, $\Gamma(U, \mathcal{F})$, $u(s)$, $\operatorname{Supp}(\mathcal{F})$}{(0.3.1.6)}
\EGAIndexLine{$\psi_{*}(\mathcal{F})$}{(0.3.4.1)}
\EGAIndexLine{$\psi_{*}(u)$}{(0.3.4.2)}
\EGAIndexLine{$\psi_{x}$}{(0.3.4.4)}
\EGAIndexLine{$\mathcal{G} \to \mathcal{F}$}{(0.3.5.1)}
\EGAIndexLine{$u^{\sharp}$, $v^{\flat}$, $\rho_{\mathcal{G}}$}{(0.3.5.3)}
\EGAIndexLine{$\psi^*(\mathcal{G})$, $\psi^*(v)$, $\sigma_{\mathcal{F}}$}{(0.3.5.5)}
\EGAIndexLine{$\mathcal{O}_{X}$, $\mathcal{O}_{X,x}$, $\mathcal{O}_{x}$, $1$, $e$}{(0.4.1.1)}
\EGAIndexLine{$\check{\mathcal{F}}$, $\bigwedge^{p} \mathcal{F}$}{(0.4.1.4)}
\EGAIndexLine{$\mathcal{J}\mathcal{F}$}{(0.4.1.5)}
\EGAIndexLine{$\Psi_{*}(\mathcal{F})$, $\Psi_{*}(u)$}{(0.4.2.1)}
\EGAIndexLine{$\Psi_{*}(\mathcal{C})$}{(0.4.2.4)}
\EGAIndexLine{$\Psi^*(\mathcal{G})$, $\Psi^*(v)$}{(0.4.3.1)}
\EGAIndexLine{$\Psi^*(\mathcal{C})$}{(0.4.3.4)}
\EGAIndexLine{$\Psi^*(\mathcal{J})\mathcal{O}_{X}$, $\mathcal{J}^{c}_{X}$}{(0.4.3.5)}
\EGAIndexLine{$\mathcal{G} \to \mathcal{F}$}{(0.4.4.1)}
\EGAIndexLine{$u_{\theta}^{\sharp}$, $u^{\sharp}$, $v_{\theta}^{\flat}$, $v^{\flat}$, $\rho_{\mathcal{G}}$, $\sigma_{\mathcal{F}}$}{(0.4.4.3)}
\EGAIndexLine{$u_{1} \otimes u_{2}$}{(0.4.4.4)}
\EGAIndexLine{$\mathcal{L}^{-1}$}{(0.5.4.3)}
\EGAIndexLine{$\mathcal{L}^{\otimes n}$}{(0.5.4.4)}
\EGAIndexLine{$\Gamma_{*}(\mathcal{L})$, $\Gamma_{*}(\mathcal{L}, \mathcal{F})$}{(0.5.4.6)}
\EGAIndexLine{$\mathcal{O}_{X}^{*}$}{(0.5.4.7)}
\EGAIndexLine{$\mathfrak{m}_{x}$, $\kappa(x)$, $f(x)$}{(0.5.5.1)}
\EGAIndexLine{$X_{f}$}{(0.5.5.2)}
\EGAIndexLine{$\operatorname{Im}(M' \otimes_{A} N')$}{(0.6.0)}
\EGAIndexLine{$\widehat A$, $\widehat M$}{(0.7.2.3) 与 (0.7.3.1)}
\EGAIndexLine{$A\{T_{1}, \cdots, T_{r}\}$}{(0.7.5.1)}
\EGAIndexLine{$A\{S^{-1}\}$}{(0.7.6.1)}
\EGAIndexLine{$\mathfrak{a}\{S^{-1}\}$}{(0.7.6.9)}
\EGAIndexLine{$A_{\{f\}}$, $\mathfrak{a}_{\{f\}}$}{(0.7.6.14)}
\EGAIndexLine{$A_{\{S\}}$}{(0.7.6.14)}
\EGAIndexLine{$(M \otimes_{A} N)^{\wedge}$, $M \widehat{\otimes}_{A} N$}{(0.7.7.1)}
\EGAIndexLine{$u \widehat{\otimes} v$}{(0.7.7.3)}

\oldpage[I]{218}
\phantomsection\label{section:index-notation-chapterI}
\medskip\noindent\textbf{第一章}\par\nopagebreak
\EGAIndexLine{$\operatorname{Spec}(A)$, $\mathfrak{j}_{x}$, $\mathfrak{m}_{x}$, $\kappa(x)$, $f(x)$, $M_{x}$, $r(E)$, $V(E)$, $V(f)$, $D(f)$}{(I, \hyperref[I.1.1.1-fr]{1.1.1})}
\EGAIndexLine{$\mathfrak{j}(Y)$}{(I, \hyperref[I.1.1.3-fr]{1.1.3})}
\EGAIndexLine{${}^{a}\phi$}{(I, \hyperref[I.1.2.1-fr]{1.2.1})}
\EGAIndexLine{$S_{f}'$}{(I, \hyperref[I.1.3.1-fr]{1.3.1})}
\EGAIndexLine{$\rho_{g,f}$}{(I, \hyperref[I.1.3.3-fr]{1.3.3})}
\EGAIndexLine{$\widetilde A$, $\widetilde{M}$, $\theta_{f}$}{(I, \hyperref[I.1.3.4-fr]{1.3.4})}
\EGAIndexLine{$\widetilde u$}{(I, \hyperref[I.1.3.5-fr]{1.3.5})}
\EGAIndexLine{$\widetilde{\phi}$}{(I, \hyperref[I.1.6.1-fr]{1.6.1})}
\EGAIndexLine{$A(X)$}{(I, \hyperref[I.1.7.1-fr]{1.7.1})}
\EGAIndexLine{$\mathcal{O}_{X/Y}$}{(I, \hyperref[I.2.1.6-fr]{2.1.6})}
\EGAIndexLine{$\operatorname{Hom}(X, Y)$}{(I, \hyperref[I.2.2.1-fr]{2.2.1})}
\EGAIndexLine{$\operatorname{Hom}_{S}(X, Y)$, $1_X$}{(I, \hyperref[I.2.5.2-fr]{2.5.2})}
\EGAIndexLine{$\Gamma(X/S)$}{(I, \hyperref[I.2.5.5-fr]{2.5.5})}
\EGAIndexLine{$X \amalg Y$}{(I, \hyperref[I.3.1.1-fr]{3.1.1})}
\EGAIndexLine{$X \times_{S} Y$, $X \times Y$, $(g, h)_{S}$, $u \times_{S} v$, $u \times v$}{(I, \hyperref[I.3.2.1-fr]{3.2.1})}
\EGAIndexLine{$X \times_{A} Y$, $X \otimes_{A} B$, $(g, h)_{A}$, $u \times_{A} v$}{(I, \hyperref[I.3.2.1-fr]{3.2.1})}
\EGAIndexLine{$X_{(S')}$}{(I, \hyperref[I.3.3.6-fr]{3.3.6})}
\EGAIndexLine{$f_{(S')}$}{(I, \hyperref[I.3.3.7-fr]{3.3.7})}
\EGAIndexLine{$\Gamma_{f}$}{(I, \hyperref[I.3.3.14-fr]{3.3.14})}
\EGAIndexLine{$X(T)$}{(I, \hyperref[I.3.4.1-fr]{3.4.1})}
\EGAIndexLine{$P \times_{R} Q$}{(I, \hyperref[I.3.4.2-fr]{3.4.2})}
\EGAIndexLine{$X(T)_{S}$}{(I, \hyperref[I.3.4.3-fr]{3.4.3})}
\EGAIndexLine{$X(B)$, $X(B)_{A}$}{(I, \hyperref[I.3.4.4-fr]{3.4.4})}
\EGAIndexLine{$X \otimes_{y} B$, $X \otimes_{\mathcal{O}_y} B$}{(I, \hyperref[I.3.6.3-fr]{3.6.3})}
\EGAIndexLine{$Z \leq Y$}{(I, \hyperref[I.4.1.10-fr]{4.1.10})}
\EGAIndexLine{$f^{-1}(Y)$}{(I, \hyperref[I.4.4.1-fr]{4.4.1})}
\EGAIndexLine{$\mathcal{N}_{X}$}{(I, \hyperref[I.5.1.1-fr]{5.1.1})}
\EGAIndexLine{$X_{\mathrm{red}}$}{(I, \hyperref[I.5.1.3-fr]{5.1.3})}
\EGAIndexLine{$f_{\mathrm{red}}$}{(I, \hyperref[I.5.1.5-fr]{5.1.5})}
\EGAIndexLine{$\Delta_{X/S}$, $\Delta_{X}$, $\Delta$}{(I, \hyperref[I.5.3.1-fr]{5.3.1})}
\EGAIndexLine{$\operatorname{rg}_{K}(X)$}{(I, \hyperref[I.6.4.5-fr]{6.4.5})}
\EGAIndexLine{$n(X)$}{(I, \hyperref[I.6.4.8-fr]{6.4.8})}
\EGAIndexLine{$\Gamma_{\mathrm{rat}}(X/Y)$}{(I, \hyperref[I.7.1.2-fr]{7.1.2})}
\EGAIndexLine{$R(X)$}{(I, \hyperref[I.7.1.3-fr]{7.1.3})}
\EGAIndexLine{$\mathcal{R}(X)$}{(I, \hyperref[I.7.3.1-fr]{7.3.1})}
\EGAIndexLine{$L(A)$}{(I, \hyperref[I.8.1.2-fr]{8.1.2})}
\EGAIndexLine{$\delta(f)$}{(I, \hyperref[I.8.2.1-fr]{8.2.1})}
\EGAIndexLine{$\mathcal{F} \otimes_{\mathcal{O}_{S}} \mathcal{G}$, $\mathcal{F} \otimes_{S} \mathcal{G}$}{(I, \hyperref[I.9.1.2-fr]{9.1.2})}
\EGAIndexLine{$\overline{\mathcal{G}}$}{(I, \hyperref[I.9.4.1-fr]{9.4.1})}
\EGAIndexLine{$\overline Y$}{(I, \hyperref[I.9.5.11-fr]{9.5.11})}
\EGAIndexLine{$\operatorname{Spf}(A)$, $\mathcal{O}_{\mathfrak{X}}$ ($\mathfrak{X}=\operatorname{Spf}(A)$)}{(I, \hyperref[I.10.1.2-fr]{10.1.2})}
\EGAIndexLine{$\mathfrak{D}(f)$}{(I, \hyperref[I.10.1.4-fr]{10.1.4})}
\EGAIndexLine{${}^{a}\phi$, $\widetilde{\phi}$}{(I, \hyperref[I.10.2.1-fr]{10.2.1})}
\EGAIndexLine{$\mathfrak{J}^{\Delta}$}{(I, \hyperref[I.10.3.1-fr]{10.3.1})}
\EGAIndexLine{$\mathfrak{X} \times_{\mathfrak{S}} \mathfrak{Y}$}{(I, \hyperref[I.10.7.3-fr]{10.7.3})}
\EGAIndexLine{$\mathcal{F}_{/X'}$, $\widehat{\mathcal{F}}$, $u_{/X'}$, $\widehat{u}$}{(I, \hyperref[I.10.8.4-fr]{10.8.4})}
\EGAIndexLine{$X_{/X'}$, $\widehat{X}$}{(I, \hyperref[I.10.8.5-fr]{10.8.5})}
\EGAIndexLine{$\widehat{f}$}{(I, \hyperref[I.10.9.1-fr]{10.9.1})}
\EGAIndexLine{$M^{\Delta}$, $u^{\Delta}$}{(I, \hyperref[I.10.10.1-fr]{10.10.1})}
\EGAIndexLine{$\Delta_{\mathfrak{S}|\mathfrak{X}}$, $\Delta_{\mathfrak{X}}$}{(I, \hyperref[I.10.15.1-fr]{10.15.1})}
\end{multicols}
\endgroup

% 以源文为先的 EGA I 术语索引大陆简体中文投影。
% 法文 NUMDAM 扫描件为控制性依据。各条目保留源索引顺序，
% 使每一行均可直接对照印刷本第 219--224 页。
\clearpage
\oldpage[I]{219}
\section*{术语索引}
\phantomsection
\label{section:index-terminology}
\begingroup
\footnotesize
\setlength{\columnsep}{2em}
\begin{multicols}{2}
\raggedcolumns
\newcommand{\EGATermIndexLine}[2]{\par\noindent\hangindent=1em\hangafter=1 #1\ \dotfill\ \textup{#2}\par}

\noindent\emph{条目顺序遵循法文源索引；因此中文译名并不总按汉语拼音顺序排列。}\par\medskip

\EGATermIndexLine{子预概形的闭包}{(I, \hyperref[I.9.5.11-fr]{9.5.11})}
\EGATermIndexLine{$\mathcal O_X$-代数}{(0.4.1.3)}
\EGATermIndexLine{进制环；$\mathfrak J$-进制环}{(0.7.1.9)}
\EGATermIndexLine{可容环}{(0.7.1.2)}
\EGATermIndexLine{完备分式环}{(0.7.6.5)}
\EGATermIndexLine{分式环}{(0.1.2.2)}
\EGATermIndexLine{有理函数环}{(I, \hyperref[I.2.1.6-fr]{2.1.6} 和 \hyperref[I.7.1.3-fr]{7.1.3})}
\EGATermIndexLine{仿射概形的环}{(I, \hyperref[I.1.7.1-fr]{1.7.1})}
\EGATermIndexLine{整环}{(0.1.0.6)}
\EGATermIndexLine{线性拓扑化环}{(0.7.1.1)}
\EGATermIndexLine{局部环}{(0.1.0.7)}
\EGATermIndexLine{支配的局部环}{(I, \hyperref[I.8.1.1-fr]{8.1.1})}
\EGATermIndexLine{$X$ 沿 $Y$ 的局部环；$Y$ 在 $X$ 中的局部环}{(I, \hyperref[I.2.1.6-fr]{2.1.6})}
\EGATermIndexLine{预进制环；$\mathfrak J$-预进制环}{(0.7.1.9)}
\EGATermIndexLine{预可容环}{(0.7.1.2)}
\EGATermIndexLine{约化环}{(0.1.1.1)}
\EGATermIndexLine{正则环}{(0.4.1.3)}
\EGATermIndexLine{同源的局部环}{(I, \hyperref[I.8.1.4-fr]{8.1.4})}
\EGATermIndexLine{$\mathcal O_X$-模的零化子}{(0.5.3.7)}
\EGATermIndexLine{与环同态相联系的谱映射}{(I, \hyperref[I.1.2.1-fr]{1.2.1})}
\EGATermIndexLine{有理映射；有理 $S$-映射}{(I, \hyperref[I.7.1.2-fr]{7.1.2})}
\EGATermIndexLine{在一点有定义的有理映射}{(I, \hyperref[I.7.2.1-fr]{7.2.1})}
\EGATermIndexLine{在开子集上诱导的有理映射}{(I, \hyperref[I.7.1.2-fr]{7.1.2})}
\EGATermIndexLine{在 $\operatorname{Spec}(\mathcal O_x)$ 上诱导的有理映射}{(I, \hyperref[I.7.2.8-fr]{7.2.8})}

\EGATermIndexLine{凝聚 $\mathcal O_X$-模}{(0.5.3.1)}
\EGATermIndexLine{凝聚 $\mathcal O_X$-代数}{(0.5.3.6)}
\EGATermIndexLine{$\mathcal O_X$-模以及 $\mathcal O_X$-模同态沿闭子集的完备化}{(I, \hyperref[I.10.8.4-fr]{10.8.4})}
\EGATermIndexLine{预概形沿闭子集的形式完备化}{(I, \hyperref[I.10.8.5-fr]{10.8.5})}
\EGATermIndexLine{不可约分支}{(0.2.1.5)}
\EGATermIndexLine{$\psi$-态射与 $\psi'$-态射的复合}{(0.3.5.2)}
\EGATermIndexLine{$\Psi$-态射与 $\Psi'$-态射的复合}{(0.4.4.2)}
\EGATermIndexLine{黏合条件}{(0.3.3.1 和 0.4.1.6)}
\EGATermIndexLine{代数预概形的基域}{(I, \hyperref[I.6.4.1-fr]{6.4.1})}
\EGATermIndexLine{几何点的剩余域}{(I, \hyperref[I.3.4.5-fr]{3.4.5})}

\EGATermIndexLine{$X\times_S X$ 的对角线}{(I, \hyperref[I.5.3.9-fr]{5.3.9})}
\EGATermIndexLine{双同态}{(0.1.0.2)}
\EGATermIndexLine{有理映射的定义域}{(I, \hyperref[I.7.2.1-fr]{7.2.1})}
\EGATermIndexLine{$\mathcal O_X$-模的对偶}{(0.4.1.4)}

\EGATermIndexLine{拓扑幂零元}{(0.7.1.1)}
\EGATermIndexLine{由一族截面生成（$\mathcal O_X$-模）}{(0.5.1.2)}
\EGATermIndexLine{截面的零点集}{(0.5.5.1)}
\EGATermIndexLine{整代数；有限整代数}{(0.1.0.5)}
\EGATermIndexLine{赋环空间}{(0.4.1.1)}
\oldpage[I]{220}
\EGATermIndexLine{局部赋环空间}{(0.5.5.1)}
\EGATermIndexLine{在开子集上诱导的赋环空间}{(0.4.1.2)}
\EGATermIndexLine{正规、约化或正则的赋环空间}{(0.4.1.3)}
\EGATermIndexLine{黏合所得的赋环空间}{(0.4.1.6)}
\EGATermIndexLine{Kolmogorov 空间}{(0.2.1.2)}
\EGATermIndexLine{不可约空间}{(0.2.1.1)}
\EGATermIndexLine{Noether 空间}{(0.2.2.1)}
\EGATermIndexLine{拟紧空间}{(0.2.2.4)}
\EGATermIndexLine{赋环空间的底空间}{(0.4.1.1)}
\EGATermIndexLine{拓扑赋环空间}{(0.4.1.1)}

\EGATermIndexLine{赋环空间上的代数层}{(0.4.1.3)}
\EGATermIndexLine{预层的层化}{(0.3.5.7)}
\EGATermIndexLine{$A$-模在 $\operatorname{Spec}(A)$ 上所确定的层}{(I, \hyperref[I.1.3.4-fr]{1.3.4})}
\EGATermIndexLine{取值于某范畴的层}{(0.3.1.1)}
\EGATermIndexLine{凝聚环层}{(0.5.3.5)}
\EGATermIndexLine{在一点正规、正规、在一点约化、约化、在一点正则或正则的环层}{(0.4.1.3)}
\EGATermIndexLine{分次环层}{(0.4.1.3)}
\EGATermIndexLine{有理函数层}{(I, \hyperref[I.7.3.1-fr]{7.3.1})}
\EGATermIndexLine{挠层}{(I, \hyperref[I.7.4.1-fr]{7.4.1})}
\EGATermIndexLine{理想层}{(0.4.1.3)}
\EGATermIndexLine{定义理想层}{(I, \hyperref[I.10.3.3-fr]{10.3.3} 和 \hyperref[I.10.5.1-fr]{10.5.1})}
\EGATermIndexLine{诱导层}{(0.3.7.1)}
\EGATermIndexLine{局部常值层}{(0.3.6.1)}
\EGATermIndexLine{黏合所得的层}{(0.3.3.1)}
\EGATermIndexLine{伪离散层}{(0.3.8.1)}
\EGATermIndexLine{常值层}{(0.3.6.1)}
\EGATermIndexLine{赋环空间的结构层}{(0.4.1.1)}
\EGATermIndexLine{仿射概形的结构层}{(I, \hyperref[I.1.3.4-fr]{1.3.4})}
\EGATermIndexLine{拓扑基上的层}{(0.3.2.2)}
\EGATermIndexLine{层的茎}{(0.3.1.6)}
\EGATermIndexLine{预概形态射的纤维}{(I, \hyperref[I.3.6.3-fr]{3.6.3})}
\EGATermIndexLine{有限代数}{(0.1.0.5)}
\EGATermIndexLine{有理函数}{(I, \hyperref[I.7.1.2-fr]{7.1.2})}

\EGATermIndexLine{点的一般化}{(0.2.1.2)}
\EGATermIndexLine{分次 $\mathcal O_X$-模}{(0.4.1.3)}
\EGATermIndexLine{态射的图}{(I, \hyperref[I.5.3.11-fr]{5.3.11})}

\EGATermIndexLine{拓扑环层的连续同态}{(0.3.1.4)}
\EGATermIndexLine{由截面定义的同态}{(0.5.1.1)}
\EGATermIndexLine{局部环的局部同态}{(0.1.0.7)}
\EGATermIndexLine{模的 $\phi$-同态}{(0.1.0.2)}

\EGATermIndexLine{可容环的定义理想}{(0.7.1.2)}
\EGATermIndexLine{$\mathcal O_X$-理想}{(0.4.1.3)}
\EGATermIndexLine{形式预概形 $\mathfrak X$ 的 $\mathcal O_{\mathfrak X}$-定义理想}{(I, \hyperref[I.10.3.3-fr]{10.3.3} 和 \hyperref[I.10.5.1-fr]{10.5.1})}
\EGATermIndexLine{素理想}{(0.1.0.6)}
\EGATermIndexLine{$\mathcal O_X$-模的直像}{(0.4.2.1)}
\EGATermIndexLine{预层的直像}{(0.3.4.1)}
\EGATermIndexLine{$S$-截面的像}{(I, \hyperref[I.5.3.11-fr]{5.3.11})}
\EGATermIndexLine{预概形在态射下的闭像预概形}{(I, \hyperref[I.9.5.3-fr]{9.5.3})}
\EGATermIndexLine{$\mathcal O_Y$-模的逆像}{(0.4.3.1)}
\oldpage[I]{221}
\EGATermIndexLine{$S$-态射的逆像}{(I, \hyperref[I.3.3.7-fr]{3.3.7})}
\EGATermIndexLine{预层的逆像}{(0.3.5.3)}
\EGATermIndexLine{$S$-预概形的逆像}{(I, \hyperref[I.3.3.6-fr]{3.3.6})}
\EGATermIndexLine{子预概形的逆像}{(I, \hyperref[I.4.4.1-fr]{4.4.1})}
\EGATermIndexLine{预概形的浸入、闭浸入或开浸入}{(I, \hyperref[I.4.2.1-fr]{4.2.1})}
\EGATermIndexLine{形式预概形的闭浸入}{(I, \hyperref[I.10.14.2-fr]{10.14.2})}
\EGATermIndexLine{预概形的局部浸入}{(I, \hyperref[I.4.5.1-fr]{4.5.1})}
\EGATermIndexLine{开子集上诱导赋环空间的典范嵌入}{(0.4.1.2)}
\EGATermIndexLine{子预概形的典范嵌入}{(I, \hyperref[I.4.1.7-fr]{4.1.7})}
\EGATermIndexLine{形式子预概形的典范嵌入}{(I, \hyperref[I.10.14.2-fr]{10.14.2})}
\EGATermIndexLine{几何单射（根态射）}{(I, \hyperref[I.3.5.4-fr]{3.5.4})}
\EGATermIndexLine{可逆 $\mathcal O_X$-模的逆}{(0.5.4.3)}
\EGATermIndexLine{可逆 $\mathcal O_X$-模}{(0.5.4.1)}
\EGATermIndexLine{与浸入相联系的同构}{(I, \hyperref[I.4.2.1-fr]{4.2.1})}
\EGATermIndexLine{预概形的局部同构}{(I, \hyperref[I.4.5.1-fr]{4.5.1})}

\EGATermIndexLine{$\mathcal O_{X_n}$-模的逆极限}{(I, \hyperref[I.10.6.6-fr]{10.6.6})}
\EGATermIndexLine{局部自由 $\mathcal O_X$-模}{(0.5.4.1)}
\EGATermIndexLine{取值于局部环的一点的所在点}{(I, \hyperref[I.3.4.5-fr]{3.4.5})}

\EGATermIndexLine{分式模}{(0.1.2.2)}
\EGATermIndexLine{忠实平坦模}{(0.6.4.1)}
\EGATermIndexLine{平坦模}{(0.6.1.1)}
\EGATermIndexLine{$A$-平坦模}{(0.6.2)}
\EGATermIndexLine{拟有限模}{(0.7.4.1)}
\EGATermIndexLine{$\mathcal O_X$-模}{(0.4.1.3)}
\EGATermIndexLine{赋环空间或拓扑赋环空间的态射}{(0.4.1.1)}
\EGATermIndexLine{忠实平坦态射}{(0.6.7.8)}
\EGATermIndexLine{平坦态射}{(0.6.7.1)}
\EGATermIndexLine{定义在拓扑基上的预层态射}{(0.3.2.3)}
\EGATermIndexLine{预概形的态射}{(I, \hyperref[I.2.2.1-fr]{2.2.1})}
\EGATermIndexLine{$S$-预概形或 $A$-预概形的态射}{(I, \hyperref[I.2.5.2-fr]{2.5.2})}
\EGATermIndexLine{双有理态射}{(I, 6.5.6)}
\EGATermIndexLine{有限型态射}{(I, \hyperref[I.6.3.1-fr]{6.3.1})}
\EGATermIndexLine{对角态射}{(I, \hyperref[I.5.3.1-fr]{5.3.1})}
\EGATermIndexLine{支配态射}{(I, \hyperref[I.2.2.6-fr]{2.2.6})}
\EGATermIndexLine{闭态射}{(I, \hyperref[I.2.2.6-fr]{2.2.6})}
\EGATermIndexLine{态射的图态射}{(I, \hyperref[I.3.3.14-fr]{3.3.14})}
\EGATermIndexLine{局部有限型态射}{(I, \hyperref[I.6.6.2-fr]{6.6.2})}
\EGATermIndexLine{可经由另一态射分解的态射}{(I, \hyperref[I.4.1.8-fr]{4.1.8})}
\EGATermIndexLine{开态射}{(I, \hyperref[I.2.2.6-fr]{2.2.6})}
\EGATermIndexLine{拟紧态射}{(I, \hyperref[I.6.6.1-fr]{6.6.1})}
\EGATermIndexLine{根态射}{(I, \hyperref[I.3.5.4-fr]{3.5.4})}
\EGATermIndexLine{约化态射}{(I, \hyperref[I.5.1.5-fr]{5.1.5})}
\EGATermIndexLine{分离态射}{(I, \hyperref[I.5.4.1-fr]{5.4.1})}
\EGATermIndexLine{$S$-预概形的结构态射}{(I, \hyperref[I.2.5.1-fr]{2.5.1})}
\EGATermIndexLine{满态射}{(I, \hyperref[I.2.2.6-fr]{2.2.6})}
\EGATermIndexLine{普遍单射的态射}{(I, \hyperref[I.3.5.4-fr]{3.5.4})}
\EGATermIndexLine{等价态射}{(I, \hyperref[I.7.1.1-fr]{7.1.1})}
\EGATermIndexLine{预概形的 $A$-态射或 $S$-态射}{(I, \hyperref[I.2.5.2-fr]{2.5.2})}
\EGATermIndexLine{形式预概形的态射}{(I, \hyperref[I.10.4.5-fr]{10.4.5})}
\EGATermIndexLine{形式预概形间的进制态射}{(I, \hyperref[I.10.12.1-fr]{10.12.1})}
\EGATermIndexLine{形式预概形的对角态射}{(I, \hyperref[I.10.15.1-fr]{10.15.1})}
\EGATermIndexLine{形式预概形的有限型态射}{(I, \hyperref[I.10.13.3-fr]{10.13.3})}
\oldpage[I]{222}
\EGATermIndexLine{形式预概形的分离态射}{(I, \hyperref[I.10.15.1-fr]{10.15.1})}
\EGATermIndexLine{形式 $\mathfrak S$-预概形的结构态射}{(I, \hyperref[I.10.4.7-fr]{10.4.7})}
\EGATermIndexLine{形式预概形的 $A$-态射或 $\mathfrak S$-态射}{(I, \hyperref[I.10.4.7-fr]{10.4.7})}
\EGATermIndexLine{预层的 $\psi$-态射}{(0.3.5.1)}
\EGATermIndexLine{从 $\mathcal O_Y$-模到 $\mathcal O_X$-模的 $\Psi$-态射}{(0.4.4.1)}

\EGATermIndexLine{环的幂零根}{(0.1.1.1)}
\EGATermIndexLine{$\mathcal O_X$-代数的幂零根}{(I, \hyperref[I.5.1.1-fr]{5.1.1})}
\EGATermIndexLine{预概形的几何点数}{(I, \hyperref[I.6.4.8-fr]{6.4.8})}

\EGATermIndexLine{仿射开子集}{(I, \hyperref[I.2.1.1-fr]{2.1.1})}
\EGATermIndexLine{仿射形式开集；进制仿射形式开集；Noether 仿射形式开集}{(I, \hyperref[I.10.4.1-fr]{10.4.1})}

\EGATermIndexLine{环的乘法子集}{(0.1.2.1)}
\EGATermIndexLine{饱和乘法子集}{(0.1.4.3)}
\EGATermIndexLine{$f$-平坦 $\mathcal O_X$-模}{(0.6.7.1)}
\EGATermIndexLine{预概形的取值于某环的点}{(I, \hyperref[I.3.4.4-fr]{3.4.4})}
\EGATermIndexLine{预概形的取值于某预概形的点}{(I, \hyperref[I.3.4.1-fr]{3.4.1})}
\EGATermIndexLine{$A$-预概形的取值于某 $A$-代数的点}{(I, \hyperref[I.3.4.4-fr]{3.4.4})}
\EGATermIndexLine{$S$-预概形的位于点 $s\in S$ 之上的点}{(I, \hyperref[I.2.5.1-fr]{2.5.1})}
\EGATermIndexLine{$S$-预概形的取值于某 $S$-预概形的点}{(I, \hyperref[I.3.4.3-fr]{3.4.3})}
\EGATermIndexLine{闭点}{(0.2.2.6)}
\EGATermIndexLine{泛点}{(0.2.1.2)}
\EGATermIndexLine{预概形的几何点}{(I, \hyperref[I.3.4.5-fr]{3.4.5})}
\EGATermIndexLine{位于 $s$ 之上的几何点}{(I, \hyperref[I.3.4.5-fr]{3.4.5})}
\EGATermIndexLine{位于 $x$ 上的几何点}{(I, \hyperref[I.3.4.5-fr]{3.4.5})}
\EGATermIndexLine{$K$-有理点}{(I, \hyperref[I.3.4.5-fr]{3.4.5})}
\EGATermIndexLine{常值预层}{(0.3.6.1)}
\EGATermIndexLine{拓扑基上的预层}{(0.3.2.1)}
\EGATermIndexLine{预概形}{(I, \hyperref[I.2.1.2-fr]{2.1.2})}
\EGATermIndexLine{Artin 预概形}{(I, \hyperref[I.6.2.1-fr]{6.2.1})}
\EGATermIndexLine{连通预概形}{(I, \hyperref[I.2.1.7-fr]{2.1.7})}
\EGATermIndexLine{基预概形}{(I, \hyperref[I.2.5.1-fr]{2.5.1})}
\EGATermIndexLine{模 $\mathfrak J$ 约化所得的预概形}{(I, \hyperref[I.3.7.1-fr]{3.7.1})}
\EGATermIndexLine{在开子集上诱导的预概形}{(I, \hyperref[I.2.1.8-fr]{2.1.8})}
\EGATermIndexLine{整预概形}{(I, \hyperref[I.2.1.7-fr]{2.1.7})}
\EGATermIndexLine{不可约预概形}{(I, \hyperref[I.2.1.7-fr]{2.1.7})}
\EGATermIndexLine{局部整预概形}{(I, \hyperref[I.2.1.7-fr]{2.1.7})}
\EGATermIndexLine{局部 Noether 预概形}{(I, \hyperref[I.6.1.1-fr]{6.1.1})}
\EGATermIndexLine{与预概形相联系的约化预概形}{(I, \hyperref[I.5.1.3-fr]{5.1.3})}
\EGATermIndexLine{$A$-预概形；$A$ 上的预概形}{(I, \hyperref[I.2.5.1-fr]{2.5.1})}
\EGATermIndexLine{$K$-代数预概形}{(I, \hyperref[I.6.4.5-fr]{6.4.5})}
\EGATermIndexLine{$S$-预概形；$S$ 上的预概形}{(I, \hyperref[I.2.5.1-fr]{2.5.1})}
\EGATermIndexLine{支配的 $S$-预概形}{(I, \hyperref[I.2.5.1-fr]{2.5.1})}
\EGATermIndexLine{$S$ 上的有限型预概形；有限型 $S$-预概形}{(I, \hyperref[I.6.3.1-fr]{6.3.1})}
\EGATermIndexLine{由基预概形扩张所得的预概形}{(I, \hyperref[I.3.3.6-fr]{3.3.6})}
\EGATermIndexLine{在 $S$ 上分离的预概形}{(I, \hyperref[I.5.4.1-fr]{5.4.1})}
\EGATermIndexLine{形式预概形}{(I, \hyperref[I.10.4.2-fr]{10.4.2})}
\EGATermIndexLine{进制形式预概形}{(I, \hyperref[I.10.4.2-fr]{10.4.2})}
\EGATermIndexLine{局部 Noether 形式预概形}{(I, \hyperref[I.10.4.2-fr]{10.4.2})}
\EGATermIndexLine{Noether 形式预概形}{(I, \hyperref[I.10.4.2-fr]{10.4.2})}
\EGATermIndexLine{$A$ 上的形式预概形；形式 $A$-预概形}{(I, \hyperref[I.10.4.7-fr]{10.4.7})}
\EGATermIndexLine{$\mathfrak S$ 上的形式预概形；形式 $\mathfrak S$-预概形}{(I, \hyperref[I.10.4.7-fr]{10.4.7})}
\EGATermIndexLine{$\mathfrak S$ 上的有限型形式预概形；有限型形式 $\mathfrak S$-预概形}{(I, \hyperref[I.10.13.3-fr]{10.13.3})}
\oldpage[I]{223}
\EGATermIndexLine{在 $S$ 上分离的形式预概形}{(I, \hyperref[I.10.15.1-fr]{10.15.1})}
\EGATermIndexLine{有限表示（模具有有限表示）}{(0.1.0.5)}
\EGATermIndexLine{有限表示（$\mathcal O_X$-模具有有限表示）}{(0.5.2.5)}
\EGATermIndexLine{Noether 归纳原理}{(0.2.2.2)}
\EGATermIndexLine{$S$-预概形的积}{(I, \hyperref[I.3.2.1-fr]{3.2.1})}
\EGATermIndexLine{形式 $\mathfrak S$-预概形的积}{(I, \hyperref[I.10.7.1-fr]{10.7.1})}
\EGATermIndexLine{集合的纤维积}{(I, \hyperref[I.3.4.2-fr]{3.4.2})}
\EGATermIndexLine{不同预概形上各层的张量积}{(I, \hyperref[I.9.1.2-fr]{9.1.2})}
\EGATermIndexLine{代数的完备张量积}{(0.7.7.5)}
\EGATermIndexLine{模的完备张量积}{(0.7.7.2)}
\EGATermIndexLine{乘积的典范投影}{(I, \hyperref[I.3.2.1-fr]{3.2.1})}
\EGATermIndexLine{子模的典范延拓}{(I, \hyperref[I.9.4.1-fr]{9.4.1})}
\EGATermIndexLine{态射到完备化的延拓}{(I, \hyperref[I.10.9.1-fr]{10.9.1})}
\EGATermIndexLine{$\mathcal O_X$-模的第 $p$ 外幂}{(0.4.1.4)}

\EGATermIndexLine{拟凝聚 $\mathcal O_X$-模}{(0.5.1.3)}
\EGATermIndexLine{拟凝聚 $\mathcal O_X$-代数}{(0.5.1.3)}

\EGATermIndexLine{理想的根}{(0.1.1.1)}
\EGATermIndexLine{环的 Jacobson 根}{(0.1.1.2)}
\EGATermIndexLine{局部自由 $\mathcal O_X$-模的秩}{(0.5.4.1)}
\EGATermIndexLine{无挠 $\mathcal O_X$-模的秩}{(I, \hyperref[I.7.4.2-fr]{7.4.2})}
\EGATermIndexLine{$K$ 上有限概形的秩}{(I, \hyperref[I.6.4.5-fr]{6.4.5})}
\EGATermIndexLine{$K$ 上有限概形的可分秩}{(I, \hyperref[I.6.4.8-fr]{6.4.8})}
\EGATermIndexLine{赋环空间在开子集上的限制}{(0.4.1.2)}
\EGATermIndexLine{态射在开子集上的限制}{(0.4.1.2)}
\EGATermIndexLine{预概形态射在子预概形上的限制}{(I, \hyperref[I.4.1.7-fr]{4.1.7})}
\EGATermIndexLine{预概形在开子集上的限制}{(I, \hyperref[I.2.1.8-fr]{2.1.8})}
\EGATermIndexLine{有理映射在开子集上的限制}{(I, \hyperref[I.7.1.2-fr]{7.1.2})}

\EGATermIndexLine{概形}{(I, \hyperref[I.5.4.1-fr]{5.4.1})}
\EGATermIndexLine{仿射概形}{(I, \hyperref[I.1.7.1-fr]{1.7.1})}
\EGATermIndexLine{局部概形；预概形在一点处的局部概形}{(I, \hyperref[I.2.4.1-fr]{2.4.1})}
\EGATermIndexLine{$K$-代数概形}{(I, \hyperref[I.6.4.1-fr]{6.4.1})}
\EGATermIndexLine{$K$ 上的有限概形；有限 $K$-概形}{(I, \hyperref[I.6.4.5-fr]{6.4.5})}
\EGATermIndexLine{$S$-概形}{(I, \hyperref[I.5.4.1-fr]{5.4.1})}
\EGATermIndexLine{仿射形式概形}{(I, \hyperref[I.10.1.2-fr]{10.1.2})}
\EGATermIndexLine{形式 $\mathfrak S$-概形}{(I, \hyperref[I.10.15.1-fr]{10.15.1})}
\EGATermIndexLine{层的截面}{(0.3.1.6)}
\EGATermIndexLine{$S$-预概形的 $S$-截面}{(I, \hyperref[I.2.5.5-fr]{2.5.5} 和 \hyperref[I.5.3.11-fr]{5.3.11})}
\EGATermIndexLine{$S$-预概形的有理 $S$-截面}{(I, \hyperref[I.7.1.2-fr]{7.1.2})}
\EGATermIndexLine{$\mathcal O_X$ 的单位截面}{(0.4.1.1)}
\EGATermIndexLine{设限形式幂级数}{(0.7.5.1)}
\EGATermIndexLine{预概形的不交并}{(I, \hyperref[I.3.1.1-fr]{3.1.1}；源索引印作 3.6.1)}
\EGATermIndexLine{由子模生成的 $\mathcal O_X$-子代数}{(0.4.1.3)}
\EGATermIndexLine{子预概形}{(I, \hyperref[I.4.1.3-fr]{4.1.3})}
\EGATermIndexLine{与浸入相联系的子预概形}{(I, \hyperref[I.4.2.1-fr]{4.2.1})}
\EGATermIndexLine{闭子预概形}{(I, \hyperref[I.4.1.3-fr]{4.1.3})}
\EGATermIndexLine{由理想层定义的闭子预概形}{(I, \hyperref[I.4.1.2-fr]{4.1.2})}
\EGATermIndexLine{闭形式子预概形}{(I, \hyperref[I.10.14.2-fr]{10.14.2})}
\EGATermIndexLine{点的特化}{(0.2.1.2)}
\EGATermIndexLine{环的谱}{(I, \hyperref[I.1.1.1-fr]{1.1.1})}
\EGATermIndexLine{可容环的形式谱}{(I, \hyperref[I.10.1.2-fr]{10.1.2})}
\EGATermIndexLine{模的支撑}{(0.1.7.1)}
\oldpage[I]{224}
\EGATermIndexLine{群层的支撑}{(0.3.1.6)}
\EGATermIndexLine{定义理想的基本系}{(I, \hyperref[I.10.3.6-fr]{10.3.6} 和 \hyperref[I.10.5.1-fr]{10.5.1})}
\EGATermIndexLine{预概形的进制归纳系}{(I, \hyperref[I.10.12.2-fr]{10.12.2})}

\EGATermIndexLine{$\mathfrak J$-进制拓扑；$\mathfrak J$-预进制拓扑}{(0.7.1.9 和 0.7.2.3)}
\EGATermIndexLine{谱拓扑}{(I, \hyperref[I.1.1.2-fr]{1.1.2})}
\EGATermIndexLine{平凡可逆 $\mathcal O_X$-模}{(I, \hyperref[I.2.4.8-fr]{2.4.8})}
\EGATermIndexLine{有限型 $\mathcal O_X$-模}{(0.5.2.1)}

\EGATermIndexLine{截面在一点的值}{(0.5.5.1)}

\medskip\noindent\textit{术语说明。} 源文仅将分离且完备的环或模称为
“进制的”；这与通常用法略有不同。
\end{multicols}
\endgroup
% EGA I 目录（大陆简体中文）。
% 法文 NUMDAM 印本第 225--228 页控制条目、次序和权威页码。
% 印本目录把参考文献列作第 215 页；本文件外交式保留该读法，
% 与实际扫描叶位置的差异另记于源勘误候选，不在正文中静默更改。
\clearpage
\oldpage[I]{225}
\section*{目录}
\phantomsection
\label{section:contents}
\begingroup
\small
\setlength{\parskip}{0pt}
\newcommand{\EGAContentsLine}[3]{\par\noindent\makebox[4.2em][l]{#1}\hangindent=4.2em\hangafter=1 #2\ \dotfill\ \textup{#3}\par}
\newcommand{\EGAContentsHead}[2]{\par\medskip\noindent\textbf{#1}\ \dotfill\ \textbf{#2}\par\nopagebreak}
\noindent\hfill\textbf{权威页码}\par\smallskip
\EGAContentsHead{引言}{5}
\EGAContentsHead{第 0 章——预备知识}{11}
\EGAContentsLine{§ 1}{分式环}{11}
\EGAContentsLine{1.0}{环与代数}{11}
\EGAContentsLine{1.1}{理想的根；环的幂零根与 Jacobson 根}{12}
\EGAContentsLine{1.2}{分式模与分式环}{13}
\EGAContentsLine{1.3}{函子性质}{14}
\EGAContentsLine{1.4}{乘法子集的变换}{15}
\EGAContentsLine{1.5}{环的变换}{17}
\EGAContentsLine{1.6}{模 $M_f$ 与一个归纳极限的等同}{19}
\EGAContentsLine{1.7}{模的支撑}{20}
\EGAContentsLine{§ 2}{不可约空间；Noether 空间}{21}
\EGAContentsLine{2.1}{不可约空间}{21}
\EGAContentsLine{2.2}{Noether 空间}{23}
\EGAContentsLine{§ 3}{关于层的补充}{23}
\EGAContentsLine{3.1}{取值于一个范畴的层}{23}
\EGAContentsLine{3.2}{拓扑基上的预层}{25}
\EGAContentsLine{3.3}{层的黏合}{28}
\EGAContentsLine{3.4}{预层的直像}{29}
\EGAContentsLine{3.5}{预层的逆像}{30}
\EGAContentsLine{3.6}{常值层和局部常值层}{33}
\EGAContentsLine{3.7}{群预层和环预层的逆像}{34}
\EGAContentsLine{3.8}{伪离散空间层}{35}
\EGAContentsLine{§ 4}{赋环空间}{35}
\EGAContentsLine{4.1}{赋环空间、$\mathcal O_X$-模、$\mathcal O_X$-代数}{35}
\EGAContentsLine{4.2}{$\mathcal O_X$-模的直像}{39}
\EGAContentsLine{4.3}{$\mathcal O_X$-模的逆像}{40}
\EGAContentsLine{4.4}{直像与逆像之间的关系}{42}
\EGAContentsLine{§ 5}{拟凝聚层与凝聚层}{44}
\EGAContentsLine{5.1}{拟凝聚层}{44}
\EGAContentsLine{5.2}{有限型层}{45}
\EGAContentsLine{5.3}{凝聚层}{47}
\EGAContentsLine{5.4}{局部自由层}{48}
\EGAContentsLine{5.5}{局部赋环空间上的层}{53}

\oldpage[I]{226}
\EGAContentsLine{§ 6}{平坦性}{54}
\EGAContentsLine{6.1}{平坦模}{55}
\EGAContentsLine{6.2}{环的变换}{55}
\EGAContentsLine{6.3}{平坦性的局部化}{56}
\EGAContentsLine{6.4}{忠实平坦模}{57}
\EGAContentsLine{6.5}{标量限制}{58}
\EGAContentsLine{6.6}{忠实平坦环}{58}
\EGAContentsLine{6.7}{赋环空间的平坦态射}{59}
\EGAContentsLine{§ 7}{进制环}{60}
\EGAContentsLine{7.1}{可容环}{60}
\EGAContentsLine{7.2}{进制环与逆极限}{62}
\EGAContentsLine{7.3}{Noether 预进制环}{66}
\EGAContentsLine{7.4}{局部环上的拟有限模}{68}
\EGAContentsLine{7.5}{设限形式幂级数环}{69}
\EGAContentsLine{7.6}{完备分式环}{72}
\EGAContentsLine{7.7}{完备张量积}{75}
\EGAContentsLine{7.8}{同态模上的拓扑}{77}
\EGAContentsHead{第一章——概形语言}{79}
\EGAContentsLine{§ 1}{仿射概形}{80}
\EGAContentsLine{1.1}{环的素谱}{80}
\EGAContentsLine{1.2}{环的素谱的函子性质}{83}
\EGAContentsLine{1.3}{模所确定的层}{84}
\EGAContentsLine{1.4}{素谱上的拟凝聚层}{90}
\EGAContentsLine{1.5}{素谱上的凝聚层}{92}
\EGAContentsLine{1.6}{素谱上拟凝聚层的函子性质}{93}
\EGAContentsLine{1.7}{仿射概形态射的刻画}{96}
\EGAContentsLine{§ 2}{预概形与预概形的态射}{97}
\EGAContentsLine{2.1}{预概形的定义}{97}
\EGAContentsLine{2.2}{预概形的态射}{98}
\EGAContentsLine{2.3}{预概形的黏合}{101}
\EGAContentsLine{2.4}{局部概形}{101}
\EGAContentsLine{2.5}{预概形之上的预概形}{103}
\EGAContentsLine{§ 3}{预概形的积}{104}
\EGAContentsLine{3.1}{预概形的不交并}{104}
\EGAContentsLine{3.2}{预概形的积}{104}
\EGAContentsLine{3.3}{积的形式性质；预概形的换基}{108}
\EGAContentsLine{3.4}{预概形的取值在预概形中的点；几何点}{111}
\EGAContentsLine{3.5}{映满和含容}{114}
\EGAContentsLine{3.6}{纤维}{117}
\EGAContentsLine{3.7}{应用：预概形的模 $\mathfrak J$ 约化}{118}
\EGAContentsLine{§ 4}{子预概形与浸入态射}{119}
\EGAContentsLine{4.1}{子预概形}{119}
\EGAContentsLine{4.2}{浸入态射}{122}

\oldpage[I]{227}
\EGAContentsLine{4.3}{浸入的积}{124}
\EGAContentsLine{4.4}{子预概形的逆像}{125}
\EGAContentsLine{4.5}{局部浸入与局部同构}{126}
\EGAContentsLine{§ 5}{约化预概形；分离条件}{127}
\EGAContentsLine{5.1}{约化预概形}{127}
\EGAContentsLine{5.2}{具有给定底空间的子预概形的存在性}{131}
\EGAContentsLine{5.3}{对角线；态射的图}{132}
\EGAContentsLine{5.4}{分离态射与分离预概形}{135}
\EGAContentsLine{5.5}{分离判据}{136}
\EGAContentsLine{§ 6}{有限性条件}{140}
\EGAContentsLine{6.1}{Noether 预概形与局部 Noether 预概形}{140}
\EGAContentsLine{6.2}{Artin 预概形}{143}
\EGAContentsLine{6.3}{有限型态射}{144}
\EGAContentsLine{6.4}{代数预概形}{147}
\EGAContentsLine{6.5}{态射的局部确定}{150}
\EGAContentsLine{6.6}{拟紧态射和局部有限型态射}{152}
\EGAContentsLine{§ 7}{有理映射}{155}
\EGAContentsLine{7.1}{有理映射和有理函数}{155}
\EGAContentsLine{7.2}{有理映射的定义域}{158}
\EGAContentsLine{7.3}{有理函数层}{161}
\EGAContentsLine{7.4}{挠层和无挠层}{163}
\EGAContentsLine{§ 8}{Chevalley 概形}{164}
\EGAContentsLine{8.1}{同源的局部环}{164}
\EGAContentsLine{8.2}{整概形的局部环}{165}
\EGAContentsLine{8.3}{Chevalley 概形}{168}
\EGAContentsLine{§ 9}{拟凝聚层补充}{169}
\EGAContentsLine{9.1}{拟凝聚层的张量积}{169}
\EGAContentsLine{9.2}{拟凝聚层的直像}{171}
\EGAContentsLine{9.3}{拟凝聚层截面的延拓}{172}
\EGAContentsLine{9.4}{拟凝聚层的延拓}{174}
\EGAContentsLine{9.5}{预概形的闭像与子预概形的闭包}{176}
\EGAContentsLine{9.6}{拟凝聚代数层；结构层的改变}{179}
\EGAContentsLine{§ 10}{形式概形}{180}
\EGAContentsLine{10.1}{仿射形式概形}{180}
\EGAContentsLine{10.2}{仿射形式概形的态射}{182}
\EGAContentsLine{10.3}{仿射形式概形的定义理想层}{183}
\EGAContentsLine{10.4}{形式预概形及其态射}{185}
\EGAContentsLine{10.5}{形式预概形的定义理想}{186}
\EGAContentsLine{10.6}{作为预概形之归纳极限的形式预概形}{188}
\EGAContentsLine{10.7}{形式预概形的积}{193}
\EGAContentsLine{10.8}{概形沿着一个闭子集的形式完备化}{194}
\EGAContentsLine{10.9}{把态射延拓到完备化上}{198}
\EGAContentsLine{10.10}{对仿射形式概形上的凝聚层的应用}{201}

\oldpage[I]{228}
\EGAContentsLine{10.11}{形式预概形上的凝聚层}{204}
\EGAContentsLine{10.12}{形式预概形间的进制态射}{206}
\EGAContentsLine{10.13}{有限型态射}{207}
\EGAContentsLine{10.14}{形式预概形的闭子预概形}{209}
\EGAContentsLine{10.15}{分离形式预概形}{212}
\EGAContentsHead{参考文献}{215}
\EGAContentsHead{记号索引}{217}
\EGAContentsHead{术语索引}{219}

\medskip
\begin{flushright}
1959 年 10 月 17 日收到。
\end{flushright}
\vfill
\begin{center}
1960 年——法国大学出版社印刷所，旺多姆（法国）\\
出版编号 25937\qquad 法国印制\qquad 印刷编号 16186
\end{center}
\endgroup
\end{document}
